% ============================================================
%  Harald Høffding, ETIK.  En Fremstilling af de etiske Principer og deres
%  Anvendelse paa de vigtigste Livsforhold.  4th ed., København og Kristiania:
%  Gyldendalske Boghandel, Nordisk Forlag, 1913.  vii + 606 pp.
%  "Ethics: An Exposition of Ethical Principles and Their Application to
%  the Most Important Relations of Life"
%
%  TRANSLATION (English), from transcription.tex (Danish) in this directory,
%  the text of record (complete 2026-10-04; see its header).  As far as is
%  known (2026-10-04) there is no earlier English translation; the book was
%  translated into German (Ethik, Leipzig: Fues/Reisland, 1888; 2nd ed. 1901)
%  and French (Morale, Paris: Alcan, 1903).  Those are not used as a base.
%
%  STATE (2026-10-05): COMPLETE -- works list, title page, both prefaces, the Contents
%  (pp. I--VII), chapters I--XLI (pp. 3--599) and the Index (Høffding's Register,
%  pp. 601--606).  No "% [text to be added ...]" markers remain.  Footnotes and
%  \opage marks checked page by page against transcription.tex for the whole book.
%  Translated against the transcription as repaired up to 2026-10-05 (postfix.tsv,
%  patches.py); if the transcription changes later, see TRANSLATION-PLAYBOOK.md §6a.
%
%  BOOK-SPECIFIC CONVENTIONS (beyond TRANSLATION-PLAYBOOK.md §2), from the
%  transcription's own:
%   * Pagination: "% --- p. N ---" + \opage{N} at the English word where
%     printed page N begins; front matter in roman (\opage{I}...).
%   * Headings as printed, mirroring the transcription's macros: \partitle,
%     \grouphead, \chapterhead{numeral}{title}; Høffding's paragraph numbers
%     (1., 2., ...) and lettered sub-heads (a. b. ...) as printed.
%   * \textsc = the print's small capitals (names in the text); \emph = its
%     letterspacing (names in the notes); \textit = its italics (titles,
%     emphasis, foreign words).  Carry all three over at the same words.
%   * Footnotes per page as printed, \footnote[k]{...}, at the same anchor.
%   * Quotations in Latin, German, French, English, Greek stay as printed;
%     what Høffding quotes in Danish is translated from his Danish.  Titles
%     cited stay as cited; "Smlgn." / "Se" / "Kfr." -> "Cf." / "See" / "Cf.";
%     citation numbers as printed.
%   * Misprints (PRINTED AS IS in the Danish): translate the sense, comment.
%   * Keep an odd-spots list while translating (playbook §6c).
%   * Names: English forms where the English reader has one (Platon -> Plato,
%     Aristoteles -> Aristotle), as in Spinoza's Ethica.
%
%  TERMS (proposed; settle them with chapter III and keep them):
%    Vurdering = valuation; Vurderingsmotiv = motive of valuation;
%    Handlingsmotiv = motive of action; Velfærd = welfare;
%    Velfærdsprincip(et) = the principle of welfare; Samvittighed = conscience;
%    Sindelag = disposition; Følelse = feeling; Drift = impulse;
%    Lyst/Ulyst = pleasure/pain; Selvhævdelse = self-assertion;
%    Hengivelse = devotion; Selvbeherskelse = self-control;
%    Personlighedsprincip(et) = the principle of personality;
%    Motivforskydning = displacement of motive; Værdiforskydning =
%    displacement of value; Humanitetens Rige = the kingdom of humanity;
%    Kultursamfund = cultural community; Ret = law / right (by context);
%    Retfærdighed = justice; Øjeblikkets Suverænitet = the sovereignty of
%    the moment; Grundværdi = fundamental value;
%    Aarsagsfrihed = freedom from causation; Menneskekærlighed = love of
%    mankind; Autoritetsprincip(et) = the principle of authority;
%    Opdragelse = education; Slægten = the race (as in ch. I).
%  Settled with chapters III--V (2026-10-05): Formaal = end; Maalestok =
%    standard; Grundlag = foundation; Vurderingsprincip = principle of
%    valuation; Forholdsfølelse = feeling of relation; Livshelhed(en) = the
%    whole of life; Grundvillie = fundamental will; Pligtfølelse = feeling of
%    duty; Retfærdighedsfølelse = feeling of justice; human Etik = humane
%    ethics (Høffding's term, as in his title "Om Grundlaget for den humane
%    Etik"); Humanitetsetik = ethics of humanity; Inderliggørelse = inward
%    deepening; Potensering = potentiation; Forskydning = displacement;
%    Nationaløkonomi = political economy; Retslære = jurisprudence;
%    Sanktion = sanction; Anger = remorse; Tilregnelighed = accountability;
%    Aarsagsfrihed = freedom from causation; Øjebliksstandpunktet = the
%    standpoint of the moment; Individets Suverænitet = the sovereignty of
%    the individual.  Ret = law (vs. Moral = morality) in ch. III, 14.
%  Hamlet (p. 119) is given in Shakespeare's English (III, 4), with
%  Høffding's Danish in a comment.
%  TRANSCRIPTION REPAIRS made while translating chs. III--V (2026-10-05):
%  66 rows in ebook-method/etik/postfix.tsv (read at the image, crop sheets
%  .parts/oddspots-ch3-5-*.png), among them: p. 50--51 the notes re-paired
%  (p. 51 has two notes; the second's mark was read as *)), p. 84 the Greek
%  (δήμου φάτις) restored, p. 105 two lost words in the Teresa sentence;
%  p. 84 "Odyss. I, 261---363" marked PRINTED AS IS (patches.py).
%  Chs. VI--X (2026-10-05): Velfærd/Lykke/Nytte = welfare/happiness/utility; where
%  Høffding argues from the Danish words (VII, 1 and 3: Lykke, Nytte, Gavn, Velfærd
%  as "fare vel") the Danish word is given in brackets.  Dyd = virtue; Ret (VIII, 1,
%  the fourth formal concept) = right; Selvhævdelse = self-assertion; Hengivelse =
%  devotion; Højmod (generositas) = high-mindedness; Aandshøjhed (sublimitas) =
%  loftiness of mind; Forbillede = model; Aladdinsnatur = Aladdin-nature.  Mill's
%  pig and Socrates (p. 140) in Mill's English; Dante (pp. 127, 141) translated
%  from Høffding's Danish, the Danish in a comment.  Further repairs to the Danish
%  (read at the image, sheets .parts/oddspots-ch6-10-*.png): 67 rows in postfix.tsv,
%  3 POST patches (the Dante verse of p. 127 set as one quote; the sheet signature
%  on p. 161 that had made a false paragraph break; an italic extent on p. 166).
%  The Contents' summaries are in sentence case.  Ch. II: "Retfærdighed" in the
%  Beatitudes (p. 21) = "righteousness", the biblical word, not "justice".
%  TRANSCRIPTION REPAIRS made while translating (2026-10-04, read at the image,
%  in ebook-method/etik/postfix.tsv and template.tex): p. 14 n. 1 "Kap. 111" ->
%  III; p. 16 ",Ja" -> „Ja"; p. 20 "ien" -> "i en"; p. 22 n. 1 closing quote;
%  Indhold: "Medfølelse" (XI, 12) and "c Frihed" (V, 2) marked PRINTED AS IS.
%  Chs. XI--XV (2026-10-05): Selvopholdelse = self-preservation; Selvbeherskelse =
%  self-control; Selvstændighed = independence; Makropsyki = macropsychy; Pietet =
%  reverence; Nødløgn = the white lie; intellektuel Redelighed = intellectual probity;
%  Motiv-/Værdiforskydning = displacement of motive/of value; Humanitetens Rige = the
%  kingdom of humanity; det frie Kultursamfund = the free cultural community; Slægter
%  (in marriage, XV) = kindreds; Hetærisme = hetaerism.  Coleridge (p. 256) and 2 John
%  10--11 (p. 258) in their English; Johnson (p. 258) translated from the Danish.
%  Misprints kept with a comment: Griesebach (p. 258), Ehrenfells (p. 271),
%  "écheation" for éducation (p. 229).  Repairs to the Danish (read at the image,
%  sheets .parts/oddspots-ch11-15-*.png, notes-ch11-12-*.png): 75 rows in
%  postfix.tsv; POST patches: five global rules (Il -> II for volumes, « and
%  \textasciitilde{} for an opening quote, "Lon ing" -> Löning), the p. 229 misprint
%  comment, and the heads on p. 281 (set as a quote with a stray "xv.").  Worst: the
%  long note of pp. 258--259 (a dozen faults).
%  Chs. XVI--XX (2026-10-05): Elskov(sfølelse) = (the feeling of) love; Troskab =
%  faithfulness; Ægtefælle = spouse; Skilsmisse = divorce; Vielse = wedding; Slægt =
%  family/kindred where it means a lineage, the race elsewhere; Kvindens Stilling og
%  Vilkaar = the position and circumstances of woman; Arbejdsdeling = division of labor;
%  Skøge = prostitute.  Kant, Burckhardt, Darwin, Mill quoted from Høffding's Danish,
%  marked in comments.  Kept as printed, with a comment: Paola Sarpi, Peter Plade
%  (= Peder Palladius), Marc Auréle, La Régime Moderne, Die Phänomene des sittlichen
%  Bewusstseins, History of Women Suffrage.  Repairs to the Danish: 74 rows in
%  postfix.tsv + 18 POST patches (sheets .parts/notes-ch16-20*.png, oddspots-ch16-20q-1.png),
%  among them the heads on p. 326 (again set as a quote), the Goethe line and two lost
%  words on p. 323, six words split at line ends.
%  Chs. XXI--XXIII (2026-10-05): Fostringsløn = return for rearing; Selvvirksomhed =
%  self-activity; Kultursamfund = cultural community; Frihed som Maal / som Middel =
%  freedom as goal / as means (Maal = goal; the Contents has Formaal = end); Livegenskab = serfdom;
%  Hoveri = compulsory labor; Stavnsbaand left in Danish with a gloss; Lav = guild.
%  Declaration of Independence in its English; Kant, Preyer, Lyell, the Henry George
%  letter, Falbe Hansen from Høffding's Danish.  Kept as printed with a comment: "paa de
%  aristoteliske Princip" (p. 333), "14. Juli 1791" (p. 354; the note says Juni).  Repairs
%  to the Danish: 33 rows in postfix.tsv + 2 POST patches (sheets .parts/notes-ch21-23*.png).
%  Chs. XXIV--XXV (2026-10-05): Modsætning = antagonism; Masse = mass; Arbejdsløn =
%  wages; Befolkningsspørgsmaal = population question.  Cicero's Latin kept; the Gotha
%  program and the coal-mine witness from Høffding's Danish.  Repairs to the Danish: 14
%  rows in postfix.tsv (sheets .parts/notes-ch24-25-*.png), among them Mill "III, 1, 2"
%  -> II (p. 372) and the superscript 2 read as a backslash.
%  Ch. XXVI (2026-10-05): Fagforening = trade union; Fællesforening = industrial
%  partnership; Produktionsforening = producers' association; Brugs-/Forbrugsforening =
%  consumers' association; Arbejdsgiver = employer; Levefod = standard of living; Samfund /
%  Selskab = community / society (Tönnies); Handel = commerce; Rovdrift = robber-farming.
%  English originals used where Høffding quotes English (Mill on the "tutelary" classes,
%  the Marx passage via the English Capital, the Labor Department act, Shaw roughly);
%  other quotations translated from his Danish and marked.  Kept as printed:
%  "gewerbschaftlichen" (p. 394), "kunde skaffe" (p. 420), "Das Kapital!" (p. 392).
%  Repairs to the Danish: 70 rows in postfix.tsv + 7 POST patches (sheets
%  .parts/notes-ch26*.png, oddspots-ch26-1.png).
%  Chs. XXVII--XXXI (2026-10-05): Fritid = free time (also in the Contents, XXVII, 2);
%  Aandsdannelse = culture of mind; Oplysning = enlightenment; Forsken = inquiry/research;
%  Lærerstanden = the teaching profession; Højskolebevægelsen = the folk high school
%  movement; Livsfølelse = feeling of life; kosmisk Livsfølelse = cosmic feeling of
%  life; Sjælesorg = cure of souls; Rensning = purification (Aristotle's catharsis).
%  Kant, Goethe, Michelangelo, Ibsen (Gengangere) and Julius Lange translated from the
%  Danish and marked.  Repairs to the Danish: 42 rows in postfix.tsv (sheets
%  .parts/notes-ch27*.png), among them "Isbs" -> 1885 (p. 444), "Sergei" -> Sergel
%  (p. 456), "Skoleplanen" -> Skoleplaner. (p. 445).
%  Chs. XXXII--XXXVI (2026-10-05): Religionsfrihed = religious freedom (also in the
%  Contents); Udskillelse = withdrawal (Contents XXXIII, 3), Adskillelse = separation,
%  Skilsmisse = divorce; Statskirke = state church; Ed = oath; Vielse = marriage ceremony;
%  Menneskekærlighed = love of mankind; Velgørenhed = beneficence; Naade = grace;
%  Fattigvæsen = poor relief; Folk = people; Flok = horde; Nationalfølelse = national
%  feeling; Ret (XXXVI, 5) = law.  KJV for 1 Cor. 15:32, 4:7 and Mark 3:24; Martensen,
%  Madvig, Eliot (via Bryce), the Württemberg judgment, Nielsen, Spencer, Gneist,
%  Aeschylus and the Danish verse line translated from the Danish and marked.  Kept as
%  printed: Durckheim (p. 479), Pater Matthew (p. 497), Armenflege (p. 511), Poul Louis
%  (p. 520).  Repairs to the Danish: 59 rows in postfix.tsv (sheets .parts/notes-ch32*,
%  ch33*, ch34*, ch36*.png), among them "Hundreåer" -> Hundreder (p. 490), "1874---1599"
%  -> 1899 (p. 506), "Ide"/"ide"/"iog" word splits.  Also: 30 lines where a text line
%  ended in a comment with no space before the %, which ran words together in the PDF
%  ("thing.''The"), now have the space.
%  Chs. XXXVII--XLI and Index (2026-10-05): Ret = law (right where it is a person's right);
%  Retsorden = legal order; Retsfølelse = sense of justice; Sæd og Skik = usage and custom;
%  Gengældelse = retribution; Afskrækkelse = deterrence; Straffetrusel = penal threat;
%  Forbedring = improvement; Formynderregimente = regime of guardianship; Selvstyrelse =
%  self-government; Folkeafstemning = popular vote (referendum); den evige Fred = perpetual
%  peace (as in the Contents, XLI, 2); Humanitetsriget = the kingdom of humanity.  English
%  originals for Hamilton (Federalist Nos. 1 and 51), Peel (29 June 1846) and the KJV
%  (Acts 5:29); Tocqueville's French as printed; Ørsted, Goos, Jhering, Gneist, Kant,
%  Bryce, Amos, Tocqueville (Démocratie), Bjørnson and others translated from the Danish and
%  marked.  Kept as printed: ``Majoriteterne'' (p. 583, for Minoriteterne; translated by the
%  sense), Bryce ``1885'' (p. 532), ``qvia'' (p. 563), ``Poul Louis'' (p. 520), ``bl. A.''
%  (p. 597).  Repairs to the Danish: 130 rows in postfix.tsv + 10 POST patches (sheets
%  .parts/notes-ch37*, ch38*, ch39*, ch40*, ch41*.png), among them a scrambled note on
%  p. 566 (Laas, Lotze, Westermarck) rebuilt from the image, ``p. ff. ff.'' -> p. 56 ff.
%  (p. 527), p. 452 -> 432 (p. 553), London 1574 -> 1874 (p. 595),
%  and a lost ``p. 8.'' (p. 556).  Also: two section junctions that had lost their
%  paragraph break (XXXIV, 3/4), and ``self- esteem'' split across p. 395.
% ============================================================
\documentclass[12pt,a4paper,openany]{book}
\usepackage[utf8]{inputenc}
\usepackage[T1]{fontenc}
\usepackage{libertinus}
\usepackage{libertinust1math}
\usepackage[english]{babel}
\usepackage{textalpha}
\usepackage[protrusion=true,expansion=true]{microtype}
\usepackage{setspace}
\usepackage{enumitem}
\usepackage[margin=1.2in]{geometry}
\usepackage{fancyhdr}
\setlength{\headheight}{28pt}
\addtolength{\topmargin}{-13.5pt}
\usepackage{hyperref}
\hypersetup{pdftitle={Ethics (4th ed., 1913)}, pdfauthor={Harald Høffding}, hidelinks}
\pagestyle{fancy}
\fancyhf{}
\fancyhead[LE,RO]{\thepage}
\fancyhead[RE]{\textit{Harald Høffding: Ethics}}
\fancyhead[LO]{\textit{1913}}
\fancypagestyle{plain}{\fancyhf{}\fancyhead[LE,RO]{\thepage}}
\setstretch{1.3}
\usepackage{marginnote}
\usepackage{multicol}
\usepackage{graphicx}
\DeclareUnicodeCharacter{0254}{\reflectbox{c}}
\renewcommand{\thefootnote}{\arabic{footnote})}
\newcommand{\opage}[1]{\marginnote{\footnotesize\textit{[#1]}}}
% the Contents, as printed: a chapter line, its summary, a sub-head, a group head
\newcommand{\ichap}[3]{\par\smallskip\noindent\hangindent=4.4em\hangafter=1\makebox[4.4em][l]{#1}#2\dotfill\,#3\par}
\newcommand{\idesc}[1]{\par{\small\leftskip=4.4em\noindent #1\par}}
\newcommand{\isub}[2]{\par\noindent\hspace*{4.4em}\textit{#1}\dotfill\,#2\par}
\newcommand{\igroup}[1]{\par\bigskip\begin{center}\large #1\end{center}}
\newcommand{\isubgroup}[1]{\par\medskip\begin{center}#1\end{center}}
% the Index: a letter heading, an entry
\newcommand{\regletter}[1]{\par\medskip{\centering\textbf{#1}\par}\smallskip}
\newcommand{\entry}{\par\noindent\hangindent=1em\hangafter=1 }
% the headings, as in the transcription
\newcounter{etikpt}
\newcommand{\partitle}[1]{\clearpage\stepcounter{etikpt}\pdfbookmark[0]{#1}{etikpt.\arabic{etikpt}}\vspace*{6em}\begin{center}{\Large #1}\end{center}\clearpage}
\newcommand{\grouphead}[1]{\bigskip\begin{center}{\large #1}\end{center}}
\newcounter{etikch}
\pdfstringdefDisableCommands{\renewcommand{\footnote}[2][]{}\renewcommand{\textit}[1]{#1}}
\newcommand{\chapterhead}[2]{\stepcounter{etikch}\bigskip\begin{center}\pdfbookmark[1]{#1 #2}{etikch.\arabic{etikch}}#1\\[3pt]{\large #2}\end{center}\medskip}

\begin{document}
\frontmatter
\pagestyle{empty}

% --- the list of the author's works facing the title page (scan PDF 6): not translated (titles stay
%     as printed); give it in Danish, as in the transcription, when the front matter is done.
% (as printed, in Danish: a list of titles)
\begin{center}
{\large HARALD HØFFDING}
\end{center}
\medskip
{\footnotesize\setstretch{1.05}\setlength{\parindent}{0pt}\setlength{\parskip}{2pt}%
\everypar{\hangindent=1.5em\hangafter=1}%
DEN ANTIKE OPFATTELSE AF MENNESKETS VILLIE. 1870.

PHILOSOPHIEN I TYSKLAND EFTER HEGEL. 1872.

DEN ENGELSKE PHILOSOPHI I VOR TID. 1874.

OM GRUNDLAGET FOR DEN HUMANE ETHIK. 1876.

PSYKOLOGI i Omrids paa Grundlag af Erfaring. (1882). Sjette paa ny gennemsete Udgave. 1911.

FORMEL LOGIK. Til Brug ved Forelæsninger. (1884). Sjette Udgave. 1913.

ETIK. En Fremstilling af de etiske Principer og deres Anvendelse paa de vigtigste Livsforhold. (1887) Fjerde, paa ny gennemsete og delvis ændrede Udgave. 1913.

PSYKOLOGISKE UNDERSØGELSER. 1889 (Videnskabernes Selskabs Skrifter).

ETISKE UNDERSØGELSER. Om Muligheden af en filosofisk Etik. --- Om Velfærdsprincipet. --- Forholdsloven i Etiken. 1891.

SØREN KIERKEGAARD SOM FILOSOF. 1892.

KONTINUITETEN I KANTS FILOSOFISKE UDVIKLINGSGANG. 1893. (Videnskabernes Selskabs Skrifter)

DEN NYERE FILOSOFIS HISTORIE. En Fremstilling af Filosofiens Historie fra Renaissancens Slutning til vore Dage. To Bind. (1894---95). Anden gennemsete Udgave. 1904.

JEAN JACQUES ROUSSEAU og hans Filosofi. (1896). Anden Udgave. 1912.

KORT OVERSIGT OVER DEN NYERE FILOSOFIS HISTORIE. (1898). Femte paany gennemsete Udgave. 1913.

MINDRE ARBEJDER. Første Række. 1899.

DET PSYKOLOGISKE GRUNDLAG FOR LOGISKE DOMME. 1899. (Videnskabernes Selskabs Skrifter).

VED AARHUNDREDESKIFTET. Tale holdt paa Københavns Universitet. 1901.

RELIGIONSFILOSOFI. (1901). Andet Oplag. 1906.

FILOSOFISKE PROBLEMER. 1902. (Universitetsprogram).

OM NOGLE RELIGIONSFILOSOFISKE ARBEJDER fra den nyeste Tid. 1903. (Universitetsprogram).

MODERNE FILOSOFER. 1904.

MINDRE ARBEJDER. Anden Række. 1905.

DANSKE FILOSOFER. 1909.

DEN MENNESKELIGE TANKE, dens Former og dens Opgaver. 1910.

RELIGION OG VIDENSKAB. 1910.

UDVALGTE STYKKER af dansk filosofisk Litteratur. Med Indledninger af \textit{Harald Høffding}. 1910.

PERSONLIGHETSPRINCIPEN I FILOSOFIN. Stockholm 1911.

MINDRE ARBEJDER. Tredie Række. 1913.\par}
\clearpage

% --- the title page (PDF 7) ---
\begin{titlepage}
\begin{center}
\vspace*{2em}
{\fontsize{48}{52}\selectfont E\,T\,H\,I\,C\,S}\\[18pt]
{\large AN EXPOSITION OF ETHICAL PRINCIPLES AND\\
THEIR APPLICATION TO THE MOST IMPORTANT\\
RELATIONS OF LIFE}\\[24pt]
{\scriptsize BY}\\[12pt]
{\LARGE HARALD HØFFDING}\\[10pt]
\rule{2.5cm}{0.4pt}\\[14pt]
{\small FOURTH, NEWLY REVISED AND PARTLY\\
ALTERED EDITION}

\vfill
{\small Translated from the Danish}

\vfill
{\small \textit{Etik. En Fremstilling af de etiske Principer og deres Anvendelse paa de vigtigste
Livsforhold.} Fjerde, paany gennemsete og delvis ændrede Udgave.\\
København og Kristiania: Gyldendalske Boghandel, Nordisk Forlag, 1913}
\end{center}
\end{titlepage}

\pagestyle{fancy}
\fancyhead[RE,LO]{\textit{Preface}}

% --- the preface to the first edition (PDF 9--11) ---
\chapter*{Preface to the First Edition}
\addcontentsline{toc}{chapter}{Preface to the First Edition}
When one catches sight of snow-covered mountains in the distance, they seem to float in the air.
Only as one draws nearer does one see clearly that they have a firm and solid ground to stand on.
So it is with ethical principles too. In the first enthusiasm one thinks one does them justice only
by giving them a place raised as high as possible above actual nature and actual life. On closer
reflection, and through longer and perhaps dearly bought experience, one discovers that they can
guide life only when they have themselves come forth from life. My task in this book has been to
show what the fundamental ethical ideas are, where they spring from, and how they find application
in the most important relations of life. Practical experience and theoretical inquiry have
confirmed in me more and more the conviction that ethical principles --- the foundation and the
standard of all judgments of good and evil --- have their source in the very nature and conditions
of man, independently of all authorities. This conviction I have here sought to establish and to
carry through.

I have wished to give not merely an abstract doctrine of ethical principles, but also a
demonstration of the applicability and the application of the principles set up. This task has
demanded a large and varied material, and difficulties have constantly arisen for me, partly in
finding the material, partly in limiting it. In order to compress my exposition within a frame of
tolerably limited size, I have taken pains to express myself as briefly and concisely as possible.
I am myself no lover of diffuse expositions, and I hope that no thoughtful reader will complain of
my brevity.

Ethical discussions are frequent in our day. It becomes ever clearer that when religious, social,
political and aesthetic questions stir minds so strongly, it is really the ethical side of these
questions that bears the sign of contradiction. There would seem, then, to be a rightful place for
an inquiry that sets out to bring to light the very foundation and standard of ethical judgments
and to draw their consequences for some of the most important relations of life. Such an inquiry
cannot, of course, treat exhaustively all the special fields it takes into account. It must find
its justification in the light that can be thrown on the particular relations of life by setting
them side by side with others under a common illumination. But it presupposes a need to go back to
the deepest-lying questions of principle, and one can politicize, aestheticize and
dogmatize\footnote[1]{I ought to have added: moralize! (Later note).}
with great zeal without feeling this need. I believe it would be of great importance if this need
could stir more than it does; it would give our intellectual life greater firmness and greater
depth; but how far my work will be able to satisfy it is, of course, quite another matter.
---

In the year 1876 I published a small work whose title was ``On the Foundation of Humane Ethics.''
It has often, wrongly but perhaps for brevity's sake, been referred to as ``humane ethics,'' though
it is not a treatment of the whole of ethics, but only of a few particular questions. In the present
work some of these questions are taken up for an entirely new inquiry; on other questions (notably
the doctrine of authority) I refer wholly to the older work; but the greater part of the contents of
the present work is neither touched upon nor treated in the older one. ---

As I now part from a work that has for a long time been the guiding thread of my reflection and my
studies, I feel all too well how little I have achieved of what I had in mind, let alone of what
might be demanded of a book that bears the venerable title: Ethics. I fear, too, that the book will
contain a good deal that must be put down to the author's personal account and that I have not
succeeded in giving scientific form and grounding. It cannot be otherwise in a work of this nature.
I must, moreover, here state something that will perhaps be a rank heresy in the eyes of eminent
men of science. I wish that philosophy, and ethics in particular, may acquire as scientific a
character as possible, and I have worked for that according to my ability; --- but the peculiar
attraction that philosophical studies have had for me from my earliest youth is nevertheless bound
up with the fact that the subjects they deal with stand, after all, in a closer relation to the
personality of the inquirer than do the subjects of other sciences. This does not, of course,
prevent my having felt at every point the obligation to ground my views in an objective way.

So I leave the book to its fate. Could it find a circle of readers who, after going through it with
attention and reflection, found that it did not wholly contradict the motto with which it ends, I
should have achieved all I dare wish.

\medskip
\noindent\hspace*{2em}\textit{April 7, 1887.}

\begin{flushright}HARALD HØFFDING.\hspace*{2em}\end{flushright}
\clearpage

% --- the preface to the fourth edition (PDF 12) ---
\chapter*{Preface to the Fourth Edition}
\addcontentsline{toc}{chapter}{Preface to the Fourth Edition}
In the second edition of this book some less essential corrections and additions were made. --- In
the third edition a section was inserted (III, 21) in which I gave a characterization of a series of
Danish works on the presuppositions of ethics. In the present fourth edition I have, besides some
minor corrections in various places, added some remarks, mostly of a literary kind.

\medskip
\noindent\hspace*{2em}\textit{October 25, 1913.}

\begin{flushright}HARALD HØFFDING.\hspace*{2em}\end{flushright}
\clearpage

% --- the Contents (pp. I--VII, PDF 13--19), set as printed ---
\chapter*{Contents}
\addcontentsline{toc}{chapter}{Contents}
\fancyhead[RE,LO]{\textit{Contents}}
{\setstretch{1.15}\setlength{\parskip}{0pt}
%
% [Contents p. I begins]
\igroup{ETHICAL PRESUPPOSITIONS.}
\hfill{\scriptsize Page}\par
\ichap{I.}{POSITIVE MORALITY AND SCIENTIFIC ETHICS}{3}
\idesc{1. The foundation and content of ethical judgments. --- 2. Misgivings about attempts at a scientific grounding of ethical judgments. --- 3. The possibility and necessity of scientific ethics. --- 4. Historical and philosophical ethics.}
\ichap{II.}{THEOLOGICAL AND PHILOSOPHICAL ETHICS}{13}
\idesc{1. The principle of authority. --- 2. Inner contradiction in the standpoint of authority. --- 3. Philosophical ethics independent of theology and metaphysics. --- 4. On the original Christian ethics.}
\ichap{III.}{THE PRINCIPLES AND METHOD OF ETHICS}{22}
\idesc{1. Ethics as a practical science. --- 2. Subjective and objective ethics. --- 3. Valuation presupposes a feeling of pleasure or pain. --- 4. The sovereignty of the moment. --- 5--7. The sovereignty of the individual. --- 8--10. The ethical feeling on the basis of sympathy. --- 11--12. The principle of welfare. --- 13--14. Its psychological-historical presuppositions. --- 15. The insufficiency of the subjective principle. --- 16. The relation between foundation and content. --- 17. Objective ethics is divided into individual and social ethics. --- 18. The relation between foundation (motive of valuation) and motive of action. --- 19. Ideal and reality. --- 20. Ethical idealization as negation, potentiation and combination. --- 21. Criticism of some other views.}
\ichap{IV.}{THE THEORY OF CONSCIENCE}{79}
\idesc{1. The psychology and history of conscience. --- 2. Individual differences. --- 3. The infallibility of conscience as authority. --- 4. The ethical sanction. --- 5. Its sufficiency. --- 6. Ethics and the hypothesis of evolution. --- 7. Will the time of the feeling of duty one day be past? --- 8. Can one do more than one's duty?}
\ichap{V.}{THE FREEDOM OF THE WILL}{106}
\idesc{1. Ethics and the law of causality. --- 2. Six different meanings of the expression ``freedom of the will'': a. Freedom from causation. b. Freedom from external %
% [Contents p. II begins]
compulsion. c. Freedom from internal compulsion. d. Ability, strength and competence. e. Freedom of choice. f. The will guided by ethical motives. --- 3. Determinism and indeterminism. a. Indeterminism makes a comprehensive view of existence impossible. b. Self-determination and freedom from causation exclude one another. c. Indeterminism makes the act of will a matter of chance. d. The ethical character of an action depends on its connection with the whole personality. e. The more modest form of indeterminism only increases the difficulties of indeterminism. f. The psychological nature and ethical significance of remorse. g. Ethics powerless without determinism. h. Theory and practice.}% the print has "c" without a stop (PRINTED AS IS in the Danish)
\ichap{VI.}{ETHICAL EVIL}{125}
\idesc{1. The psychology of evil. --- 2. Evil as resting on ignorance and blindness. The stupid devil. --- 3. Obduracy.}
\ichap{VII.}{THE THEORY OF WELFARE}{135}
\idesc{1. The general concept of welfare. --- 2. Two main problems. --- 3. Is not every complete satisfaction of equal worth? --- 4. Is culture a road to general welfare? --- 5. End and means.}
\ichap{VIII.}{INDIVIDUAL AND SOCIAL ETHICS}{155}
\idesc{1. The formal ethical concepts are not suited to furnish a principle of division for ethics; it must be divided into individual and social ethics. --- 2. Historical development of the relation between individual and society. --- 3. Individual ethics taken as ethics proper. --- 4. Absolute altruism. --- 5. Attempts to co-ordinate individual and social ethics. --- 6. The center of gravity of ethics is of a social nature. The principle of personality.}
\igroup{INDIVIDUAL ETHICS.}
\ichap{IX.}{THE DIVISION OF INDIVIDUAL ETHICS}{183}
\idesc{1. Justice the chief ethical virtue, comprising within itself self-assertion and devotion. --- 2. Individual differences.}
\ichap{X.}{THE PERSONAL FOUNDATION OF THE ETHICAL LIFE}{189}
\idesc{1. On the significance of practice. --- 2. Work on the strength, extent and purity of conscience. --- 3. Bodily and spiritual hardening. --- 4. The history of the virtues.}
\ichap{XI.}{SELF-ASSERTION}{207}
\idesc{1. Its three main forms.}
\isub{a. Self-preservation}{207}
\idesc{%
% [Contents p. III begins]
1--2. Self-preservation as instinct and as duty. --- 3. Bodily health and strength. --- 4. Suicide as an effect of insanity and as flight from obligations. --- 5. Suicide as an expression of enfeeblement of the will. --- 6. Can suicide be a right or even a duty? --- 7. The state and suicides.}
\isub{b. Self-control}{220}
\idesc{8. The psychology of self-control. --- 9. Individualistic and human-social view of self-control. --- 10. On various kinds of self-control, especially on the relation to the sexual instinct. --- 11. Going astray and getting stuck.}
\isub{c. Independence}{231}
\idesc{12. True self-esteem (macropsychy). --- 13. No isolation. --- 14. External, personal freedom. Honor and property. Civil rights.}% the print has "Medfølelse" (sympathy), a misprint for "Selvfølelse" (self-esteem), the word of the text, p. 231 (PRINTED AS IS in the Danish)
\ichap{XII.}{DEVOTION}{236}
\isub{a. Love of other beings}{236}
\idesc{1. Strength and extent. Christianity and Stoicism. --- 2. The various forms of love, especially high-mindedness. --- 3. On duties toward animals.}
\isub{b. Love of truth}{249}
\idesc{4. The grounding of the duty of truth. Honesty, personal truth and intellectual probity. --- 5. Pedagogical limitation of the duty of truth --- for the sake of truth itself. --- 6. Tolerance and reverence.}
\igroup{SOCIAL ETHICS.}
\ichap{XIII.}{INTRODUCTION AND DIVISION}{263}
\idesc{1. Ethics and sociology. --- 2. Historical relativity. --- 3. Humanization and emancipation. --- 4. The Aristotelian principle. Displacement of motive. Displacement of value. --- 5. Society and organism. --- 6. The kingdom of humanity. --- 7. Family, free cultural community and state.}
\isubgroup{\textit{A. THE FAMILY.}}
\ichap{XIV.}{THE ETHICAL SIGNIFICANCE OF THE FAMILY}{278}
\idesc{1. The family as the most intimate community. --- 2. The family as an all-round community of life. --- 3. The family and the larger communities.}
\isubgroup{1. MARRIAGE.}
\ichap{XV.}{SOCIOLOGICAL DATA}{281}
\idesc{1. Different forms of marriage. --- 2. The relation between the sociological and the ethical inquiry.}
%
% [Contents p. IV begins]
\ichap{XVI.}{FREE MONOGAMY}{286}
\idesc{1. The ethical grounding of monogamy. --- 2. Marriage and ``free love.'' --- 3. Fate and guilt in marriage. --- 4. The significance of common tasks and common fate. --- 5. Man and woman as equals. --- 6. Free monogamy and hetaerism.}% the print has "2" without a stop
\ichap{XVII.}{ON ENTERING INTO AND DISSOLVING MARRIAGE}{299}
\idesc{1. The influence of difference in views and character. --- 2. Responsibility in entering into marriage. --- 3. On marriage between near relatives. --- 4. External attestation of the entering into marriage. Equal rights for both parties. --- 5. Divorce.}
\isubgroup{2. THE POSITION AND CIRCUMSTANCES OF WOMAN.}
\ichap{XVIII.}{SOCIOLOGICAL DATA}{310}
\idesc{1. The position of woman is still a subject of discussion. --- 2. Interaction between nature and circumstances. --- 3. Various forms of division of labor between man and woman.}
\ichap{XIX.}{THE ETHICAL POSITION OF WOMAN}{315}
\idesc{1. Woman in the family. --- 2. The surplus of women. --- 3. Feminine peculiarities. --- 4. The testimony of experience. --- 5. Emancipation as a duty. --- 6. Political rights.}
\isubgroup{3. PARENTS AND CHILDREN.}
\ichap{XX.}{SOCIOLOGICAL DATA}{326}
\idesc{1. Unconditional parental authority at earlier stages of culture. --- 2. Motives that have brought about the recognition of the child's right.}
\ichap{XXI.}{ETHICS AND PEDAGOGY}{329}
\idesc{1. The development of family feeling. --- 2. The independent value of childhood. --- 3. Unconscious and conscious education. --- 4. Intellectual education. --- 5. Religious education.}
\ichap{XXII.}{THE STATE AND THE CHILDREN}{338}
\idesc{1. The duty of care for children in the physical respect. --- 2. The child's right to instruction.}
\isubgroup{\textit{B. THE FREE CULTURAL COMMUNITY.}}
\ichap{XXIII.}{FREEDOM AND CULTURE}{342}
\idesc{1. The principle of freedom in the family and in the cultural community. --- 2. Freedom as end and as means. --- 3. Unfreedom through division of labor. --- 4. The overvaluation of the principle of freedom in the 18th century. --- 5. Different kinds of culture.}
%
% [Contents p. V begins]
\isubgroup{1. MATERIAL CULTURE.}
\ichap{XXIV.}{SOCIAL ANTAGONISMS}{358}
\idesc{1. Older and newer valuation of material labor. --- 2. Possession and labor. The unfortunate effects of the division of labor.}
\ichap{XXV.}{THE SOCIAL QUESTION}{364}
\idesc{1. Its particular appearance in our time. --- 2. Its connection with the question of population. --- 3. The social question is an ethical question. The concept of the mass. --- 4. Two extreme conceptions. --- 5. Historically given points of departure necessary. The general discussion.}
\ichap{XXVI.}{POSSIBILITIES OF DEVELOPMENT}{373}
\idesc{1. The organization of the mass by free forces or by the intervention of the state. Ethics and political economy.}
\isub{a. The ordering of labor through free association}{375}
\idesc{2. Emancipation and association. --- 3. Trade unions. --- 4. Industrial partnerships. --- 5. Producers' associations. --- 6. Consumers' associations.}
\isub{b. The ordering of labor by the intervention of the power of the state}{386}
\idesc{7. The concept of socialism. Four kinds of socialism. --- 8. The justification of its fundamental idea and of its criticism of the present order. --- 9--12. Criticism of socialism. --- 13. The pedagogical significance of socialism. --- 14. Forms of the state's influence on the ordering of labor without conflict with the principle of freedom. --- 15. The ethical justification and significance of private property. --- 16. Commerce and its ethical significance.}
\isubgroup{2. IDEAL CULTURE.}
\ichap{XXVII.}{MATERIAL AND IDEAL CULTURE}{424}
\idesc{1. Their reciprocal relation. --- 2. The significance of free time. --- 3. Bright and dark sides of ideal culture.}
\isubgroup{\textit{A. INTELLECTUAL CULTURE.}}
\ichap{XXVIII.}{THE ETHICAL SIGNIFICANCE OF SCIENTIFIC KNOWLEDGE}{430}
\idesc{1. Fluctuations in the valuation of the ethical significance of intellectual culture. --- 2--3. The psychological and historical significance of these fluctuations. --- 4. The connection of knowledge with actual life. --- 5. The unity of scientific knowledge in spite of specialization; science as the common work of the race. --- 6. Schools and parties.}
\ichap{XXIX.}{THE FREEDOM AND INDEPENDENCE OF INTELLECTUAL CULTURE}{440}
\idesc{1. The freedom of science. --- 2. Science as an independent %
% [Contents p. VI begins]
contributing element in the life of the race. --- 3. The school and the political or religious parties.}
\isubgroup{\textit{B. AESTHETIC CULTURE.}}
\ichap{XXX.}{ART AND LIFE}{448}
\idesc{1. The relation between science and art. --- 2. Art as an ideal life. --- 3. Aesthetic and ethical valuation. --- 4. Art must not take the place of life. --- 5. National art, and art that answers to its time.}
\isubgroup{\textit{C. RELIGIOUS CULTURE.}}
\ichap{XXXI.}{ETHICS AND RELIGIOUS FEELING}{457}
\idesc{1. The points of contact of religious culture with ethics. --- 2. The cosmic feeling of life. --- 3. Its intellectual and ethical elements. --- 4. Individual differences. The cure of souls.}
\ichap{XXXII.}{THE SOCIAL-ETHICAL SIGNIFICANCE OF THE POSITIVE RELIGIONS}{467}
\idesc{1. Ethical elements in dogma and cult. --- 2. Positive religion a condensation of elements from all sides of spiritual life. --- 3. The contradiction in the position of positive religion under the division of labor in the spiritual domain. --- 4. Positive religion is social religion. --- 5. Contradiction in the idea of a universal religion that builds on a dogmatic confession. Faith and love. --- 6. Religious freedom and its consequences. --- 7. The different positions of believers and freethinkers toward these consequences.}
\ichap{XXXIII.}{STATE AND CHURCH}{483}
\idesc{1. The church as an educating cultural power. --- 2. State church and religious freedom. --- 3. The withdrawal of the church from the state is not a complete divorce between state and church. --- 4. Progressive loosening of the bond between state and church.}
\isubgroup{3. PHILANTHROPIC CULTURE.}
\ichap{XXXIV.}{THE NATURE AND SIGNIFICANCE OF PHILANTHROPY}{495}
\idesc{1. The relation of philanthropic culture to material and ideal culture. --- 2. Its independence. --- 3. The human rights of the suffering. --- 4. The relation between giver and receiver. --- 5. Misgivings about confessional philanthropy.}
\ichap{XXXV.}{THE ORGANIZATION OF PHILANTHROPY}{505}
\idesc{1. The necessity and the drawbacks of organization. --- 2. The intervention of the state. --- 3. Drawbacks of state philanthropy.}
\isubgroup{\textit{C. THE STATE.}}
\ichap{XXXVI.}{PEOPLE AND STATE}{513}
\idesc{1. The relation of the state to the family and the free cultural community. --- 2. A %
% [Contents p. VII begins]
people comes into being through common fate and common activity, which lead to common customs. --- 3. National feeling as a feeling of contrast. --- 4. National feeling as instinct. --- 5. The state is the organized people.}
\ichap{XXXVII.}{LAW AND MORALITY}{520}
\idesc{1. Original connection between law and morality. --- 2. The difference between law and morality. --- 3. The legal order as a part of the ethical order. --- 4. Legality and morality. --- 5. Two main points in the history of the development of law. --- 6. The relation of public opinion to law and morality.}
\ichap{XXXVIII.}{THE ETHICAL SIGNIFICANCE OF THE STATE}{534}
\idesc{1. Theological theory of the state. --- 2. The state conceived as realized ethics. --- 3. The state conceived as mere power. --- 4. The individualistic theory of the state. --- 5. The relation of the state to the ethical life of the people, to family, culture and law. --- 6. The task of the state is to order the given content of life in fixed forms. --- 7. The state as the expression of the people's unity and power.}% the print has no final stop
\ichap{XXXIX.}{THE PUNITIVE AUTHORITY OF THE STATE}{548}
\idesc{1. The instinct of revenge. --- 2. The limitation of retribution. --- 3. The passing of retribution to the state. The metamorphosis of the instinct of revenge. --- 4. Ethical grounding of the state's authority to punish. Education and deterrence. --- 5. Criticism of the doctrine of retribution. --- 6. On the measuring of punishment according to the doctrine of retribution and according to the pedagogical theory of punishment.}
\ichap{XL.}{THE CONSTITUTION OF THE STATE}{572}
\idesc{1. Given conditions and conscious intentions. --- 2. The double control as the great problem of politics. --- 3. The significance of political freedom. --- 4. Free constitution. --- 5. Its dangers. --- 6. Majority and minority. --- 7. Political parties. --- 8. Self-government.}
\ichap{XLI.}{CONCLUSION}{592}
\idesc{1. The kingdom of humanity and the smaller communities. --- 2. Tendencies in the direction of perpetual peace. --- 3. Enthusiasm for what is great and faithfulness in what is small.}
\par}

\mainmatter
\fancyhead[RE]{\textit{Harald Høffding: Ethics}}
\fancyhead[LO]{\textit{1913}}

% --- p. 1 ---
\opage{1}\partitle{ETHICAL PRESUPPOSITIONS}

% --- p. 3 ---
\opage{3}\chapterhead{I.}{POSITIVE MORALITY AND SCIENTIFIC ETHICS.}

1. Ethical judgments contain valuations of human actions. When in actual life we call an action good
or evil, we do not thereby explain how the action came about, but we declare what value it has in
our eyes. Every such valuation presupposes, on the one hand, that there is a need, a feeling, that
drives us to judge the action (whether the judgment is merely passed in our own thoughts or also in
words addressed to others), --- and, on the other, that we have a standard, an ideal, with which the
action can be compared and by which it can be tested. \textit{The motive of ethical valuation} I call
the \textit{foundation} of our ethics, or its \textit{motive of valuation.} Its character is
essentially determined by the \textit{fundamental value} on which the highest end of our life, and
with it the standard of our ethical valuation, depend. \textit{The standard of ethical valuation}
becomes determinative for the content of our ethics, since it decides which actions, directions of
life and forms of life can be called ethically good.

The character of an ethical view depends on what foundation it presupposes, what standard it
applies, and what relation holds between its foundation and its standard.

But ethical judgments are for the most part not passed with an express and clear consciousness of
the principles they presuppose. Such
% --- p. 4 ---
\opage{4}a consciousness, as it develops further, leads to \textit{scientific} ethics. But it awakens
only when reflection and doubt make themselves felt. The question then naturally arises here, what
justification and what significance the conscious bringing forward and discussion of ethical
principles has at all.

2. Not only are ethical judgments, judgments of good and evil, passed before thinking proper and
scientific inquiry awaken; but ethical judgments that are to have personal truth and practical
significance must always spring from a living feeling and impulse that gives us no rest until we
have spoken out. It is the condition of a sound and able life that we do not fetch the principles
of our most important decisions from afar and by laborious reflection, but that the decision
springs from what has become flesh and blood in us. In our ethical judgments we display our
personality; they must therefore be determined by our whole nature and not merely by reasonings
that we may carry on in our idle hours, if we have any. Life, besides, does not always give us time
to think things over, but demands that judgment be pronounced on the instant.

And even if we have the time and the ability for reflection, will not our feelings and our
impulses be so strong that through and through --- without our noticing it ourselves --- they
determine our thinking and lead it instead of being led by it? Is it not a prejudice, long since
overthrown by psychology, that reason should be the sovereign authority within us, and has it not
often enough been shown that what came forward as the utterance of reason was in reality an
outbreak of the heart's need? In logic and mathematics pure reason may perhaps speak through our
judgments; but its voice is too weak where the most concrete, the most personal thing of all is in
question, namely human actions.

Will not reflection and criticism here even be questionable, if not positively harmful? Reflection
always has a certain dissolving influence; it robs us of the instinctive certainty and security
with which we begin our life. It weakens our strength even where it does not paralyze us by the
doubts it awakens. We no longer act with concentrated energy and no
% --- p. 5 ---
\opage{5}longer pronounce our judgments with force and confidence; indeed we may in the end become
inclined to withhold our judgments, because it seems to us impossible to reach any sure decision.
But then life dries up and in the end comes wholly to a standstill. For ethical judgments are not
mere theoretical curiosities. They are not merely the expression of feeling, but also exert a
determining influence both on the feeling and action of the one who judges and on those of others.
They are not merely \textit{motivated,} but are themselves in turn \textit{motives,} practical forces of
the greatest significance.

The misgiving becomes greater still when we extend our view beyond the single individual. The
feelings and impulses of the individual are determined by the whole \textit{nature, conditions of
life and traditions of the race.} The individual receives some of his most far-reaching abilities
and instincts as an inheritance from the race. He is then led, by developing and being brought up
as a member of the family, the tribe and the state, into a certain spiritual atmosphere, in which
habits of life, ideas, incitements and tasks present themselves to him, not as something that can
be made the object of his deliberation and choice, but as something he involuntarily takes into
himself. As he has sprung from the womb of the race, so he drinks in the traditions of the race
with his mother's milk. His way of thinking, feeling and acting is an unconscious inheritance from
earlier generations. Through instincts of race and traditions, through involuntary imitation and
practice, the individual's ethics is founded for him before his conscious reflection can
intervene. The will is involuntarily led into definite courses, and the involuntarily formed
direction of the will determines the interests and ends of life, and these in turn the judgments
of good and evil. The virtues come into being --- as \textsc{Jacobi} says somewhere --- before they get
names and before they are made into demands: ``The book of life must be written before the index to
it can be drawn up; our theories of morals are such indexes, drawn up afterwards.'' And ---
Jacobi adds --- they are as a rule drawn up by people who understand nothing of the book itself! In
any case the book is more important than the index, and the index cannot replace the book. For in
ethical judgments, the strong outbursts of feeling in which a valuation of human actions is
expressed, we have not merely the disposition of the single individual
% --- p. 6 ---
\opage{6}expressed, but it is \textit{the result of the experiences of the race} that displays itself in
them. The individual does not make his ethics himself, does not ``hit upon'' it, does not
construct it absolutely from the beginning. And it is only so that it gets its proper strength. The
ethics living in the race is a condition of the health and strength of human life. He who by his
reflection and criticism would dissolve what nature has joined together has not only an enormous
resistance to overcome, but must also know well what he is doing, lest he end in a quagmire and
drag others down with him.

This actual ethics of the race, at work in life, has been called \textit{positive morality.} It
appears in \textit{current judgments and maxims,} which are often clothed in the form of the proverb,
and which may either be lasting expressions of the practical wisdom of a nation, a tribe or a
religious community, or have a more short-lived existence as the ``public opinion'' of a century or
an age. It appears also in the \textit{personal exemplars} (founders of religions, heroes,
lawgivers, etc.) to whom a single generation, or generation after generation, looks up as the
highest expression of what is truly human. It appears, finally, in the way in which the human
\textit{relations and forms of life} (the family, civil society, the state, the church) are ordered.
Positive legislation always contains a certain part of positive morality. Common to all forms of
positive morality is that an action that conflicts with them becomes the object of a reaction on
the part of society. A violation of the prevailing maxims, exemplars or forms of life brings with it
an outburst called forth by offense or indignation, the judgment of public opinion, on which the
agent's future honor depends. Or it does not stop at this moral censure, but the agent must make
amends and perhaps through pain or even death atone for what he has done and appease the
indignation he has aroused. This assertion of positive morality through social influence is
denoted by the expression \textit{sanction.} It need not be felt by the individual as something
merely external; by virtue of his original, involuntary growing into positive morality, the
individual may feel its power in his inmost self as one with the power of his own conscience, ---
if one is to speak of conscience
% --- p. 7 ---
\opage{7}here at all; conscience does not really exist here at all as a purely individual factor.

Positive morality is a bed into which the stream of human life has been led and which it has
deepened more and more, and which the stream goes on following as long as new channels or dams do
not prevent it. But will it be possible to intervene here with conscious thought? And if it is
possible, will it not bring with it a harmful weakening of the force with which the stream hurries
on?

3. Against such well-founded and very significant objections I shall try to show that scientific
ethics is so far from being either impossible, or harmful, or unnecessary, that it develops out of
positive morality itself, acts back on it as a clarifying and liberating force, and solves tasks
that could not otherwise be solved.

First of all it is clear that one cannot defend positive morality and maintain its significance as
a condition of the health and strength of human life without having a standard by which one judges
it. This standard one cannot, unless one is willing to move in a circle, fetch from positive
morality itself. One must then try to make clear to oneself what principle one lays at the basis of
one's valuation, --- but then one has already entered on the road that leads to scientific ethics.
As soon as one establishes that health and strength are better than sickliness and faintness, one
has already established a principle of valuation. The question will then naturally arise whether
positive morality really satisfies the demands of this principle everywhere, --- whether there are
not other principles that must be laid at the basis of valuation, --- and the discussion is then in
full swing!

It is self-contradictory to come forward as defender of positive morality if one sees a danger in
all reflection and criticism. Toward what is to act with the power of the unconscious, silence is
best. A serious doubt can be met only with decisive counter-reasons, or by showing that there are
questions that can be left unanswered without life suffering from it. It belongs to the noblest
needs of man to ask
% --- p. 8 ---
\opage{8}``why?'', and the suppression of such questioning easily becomes a suppression of spiritual
life as a whole. In questioning and doubt, too, the health and strength of life can show
themselves.

Life leads by its own road to questioning. Only where the horizon is narrow and the tasks are
limited is there full security. But with growing experience there begins the comparison between
the various laws, the various ideals, the various institutions that one comes to know among
various peoples at various times. Or the new experiences and situations set tasks that cannot be
solved according to the inherited ethics. Or one seeks to put in order the great multitude of
ethical judgments that one passes and hears others pass, so that the more significant can be
distinguished from the less. Ethical judgments are, to be sure, from the beginning involuntary
outbursts of feeling, and ``there is no disputing about feelings''; \textit{but one can dispute about
the validity of the} ideas \textit{to which the feelings are attached, and about the value of the
actions to which the feelings lead.}

It is certainly a difficult point in spiritual development when reflection and criticism awaken in
the single individual or in the people. Instinct and impulse lose their immediate strength and
certainty, and the great question is whether a substitute can be won for this
loss\footnote[1]{Cf. my \textit{Psykologi} VII B, 3.}. But where it is life itself that puts the
questions into our mouths, they must either be answered, or it must be decided why they cannot be
answered. And it must be well noted that certainty and strength are not absolute goods. The
sleepwalker walks more surely than the waking man could; but we wake him or stop him nevertheless,
to keep him from coming near the abyss. The greatest strength may display itself in a most
pernicious direction. When one then seeks to diminish this strength, one diminishes its pernicious
effects too. And one must then try to procure for the better insight that is won all the energy
that formerly stood in the service of immediate instinct. All sound spiritual development
consists, after all, in energy being transferred from narrower to larger ends. Nothing is produced
from nothing.

% --- p. 9 ---
\opage{9}Scientific ethics will not and cannot put itself in the place of positive morality. It will
only ground, develop and widen it. In scientific ethics we seek only to understand ourselves, to
become clear about the principles by which we lead our life, and to bring these principles to
greater clarity and mutual harmony. In human spiritual life there takes place a constant
interaction between the conscious and the unconscious, as well as between knowledge, feeling and
will among themselves. What is gained in one field can benefit all other spiritual fields.

4. Two tasks of scientific ethics must be distinguished. It is partly a historical, partly a
philosophical science. \textit{Historical or comparative ethics} seeks to present positive morality as
it appears in a given people at a given time, to show what development it undergoes under
different conditions, and to compare the different forms it may take among different peoples and
at different times. It seeks to find the causes of these different stages of development and forms
in definite physical, psychological and historical conditions. \textit{Philosophical ethics} has as
its task not the description and explanation of given ethical phenomena, but their
\textit{valuation.} By valuation is meant an inquiry into whether an action or an ordering of human
relations of life agrees with the end that we presuppose for all human striving, and so also with
the \textit{fundamental value,} the value to which we trace back all other values. Scientific ethics
may rightly be called a practical science, since it presupposes an end that is to be realized
through human actions. Every ethical judgment presupposes such an end, for feeling is set in motion
by the sight or the thought of an action only when this action either furthers or hinders something
whose existence and progress lie close to our heart. An ethical judgment presupposes an interest of
life that is either favored or injured by a human action or by an ordering of human relations.

What has developed historically is not, as a matter of course, ethically justified. We can
understand how positive morality has developed
% --- p. 10 ---
\opage{10}into the form it has taken in our country and in our time, and yet we may condemn it to a
greater or lesser extent. Historical ethics can teach us only what actually exists and how it has
come to exist; but unless one assumes that the actual as such is one with the rational, there still
remains a valuation of it to be made. Philosophical ethics is a valuation systematically carried
through on a definite foundation given in human nature and according to the standard that answers
to this foundation. Such an inquiry, carried through, may lead to the condemnation of certain
actions and forms of life and to the approval of others. It may be that \textit{inconsistencies are
shown} within positive morality, and since the presupposition that all ethical judgments must agree
with one another must necessarily lie at the basis, the task becomes to develop positive morality
so that the contradiction disappears. It may further appear, from \textit{the comparative
consideration} of various forms of positive morality, that the one form has the advantage over the
other; this presupposes a definite standard on which the comparison is grounded, and it may lead
one to seek means of developing the type of positive morality recognized as standing lower in the
direction of the one recognized as standing higher. Finally, it may often happen that the original
motives of certain actions and institutions can no longer be acknowledged, without these actions
and institutions themselves therefore losing their value and significance. Under the changed
conditions, and by being seen from another side than before, they may bring with them effects that
were formerly unknown. Traits of character such as knightly courage, chastity and self-control,
which were developed under raw and barbarous conditions, may be taken into the service of quite
other tasks than the original ones. Private property arose historically under the influence of the
lust for power and the acquisitive impulse and has been recognized by the power of society in the
interest of peace. But in the ethical discussion of its lasting significance and justification,
quite other considerations come to be taken into account. It is here as when, in a town, the old
ramparts that were originally built for defense are turned into promenades and pleasure grounds.
There takes
% --- p. 11 ---
\opage{11}place in such cases a \textit{displacement} in the ethical field, whereby the old motives and
values fall away and new motives and values take their place, while actions, habits and
institutions remain the same. It becomes evident from this that the historical explanation of an
ethical phenomenon is not, as a matter of course, decisive for its valuation, since the possibility
of such a displacement always stands open. Philosophical ethics respects nothing existing merely
because it exists; but in seeking to lead beyond what exists, it seeks to show how what has
developed historically can be applied and turned into new forms.

While historical and philosophical ethics, seen from this side, have a very different character,
seen from another side it will prove difficult to keep them sharply apart. --- The historical
ethicist will not be able to free himself from certain ethical presuppositions, which may come to
influence the light in which he sees the past. He must even have an eye for, and practice in,
making a systematic valuation, if he is to grasp the inner coherence (or want of coherence) in the
ethical views and customs of other times and peoples. And --- as has already been suggested --- as
soon as he makes a valuing comparison between different forms of positive morality and calls some
lower, others higher, he presupposes a definite standard (and a definite fundamental value) and has
thus already entered the field of philosophical ethics. --- The philosophical ethicist will of
course always be a child of his time and his people, and the positive morality prevailing in these
will in many ways come to influence his results. But this means only that the work of philosophical
ethics must always be taken up anew. We must separate out and test whatever in the work of our
predecessors is due only to passing conditions and influences. In Kant's ethics, for example,
one notices clearly that he is a German, that he was brought up under pietist influence, and that
as a mature man he lived in the age of rationalism. Some of the merits and some of the faults of his
ethics are to be explained by these causes. But such a dependence on people and
% --- p. 12 ---
\opage{12}time\footnote[1]{An interesting example of how the national direction of mind determines the work and
results of individual thinkers is offered by the history of Italian philosophy, especially as
regards ethics. Cf. Luigi \emph{Ferri}: \textit{National Character and Classicism in Italian
Ethics.} (International Journal of Ethics 1895). --- Cf., for ancient Greece: \emph{Leopold
Schmidt}: \textit{The Unity of the Ethics of ancient Greece.} (ibid. 1894).} does not prevent the
great philosophical ethicist from uttering thoughts that are valid for many generations.
Philosophical ethics is made impossible by historical limitation only if it wishes to be a pure
science of reason, an exposition of eternal truths. It is, and can never become anything other
than, the systematic carrying through of a historically given standpoint, and it can construct
ideals and norms only for certain definite conditions of life; it is in its way --- like positive
morality --- a means in the struggle for existence: a means of maintaining a stage of life
attained and of attaining a new stage of life that may deserve the name of a higher one. Its
scientific character depends on the clarity with which its foundation and its standard are fixed,
and on the consistency with which its principles are carried through in particulars. Its practical
significance depends on the degree to which the thoughts and results so won can pass into the flesh
and blood of larger or smaller circles of mankind, find points of attachment in strong feelings and
habits, and become guiding in the ordering of the various relations of life. Experience shows that
there may often be a long and intricate road from the general ethical thought to its practical
realization. ---

(1913). In recent times the French sociological school, whose most important representatives are
\textit{Durkheim} and \textit{Levy Bruhl,} has with great energy stressed the significance of the
historical study of positive morality at different times. They expect that on the basis of such a
comparative ethics there will one day be developed a ``moral art'' that stands in the same relation
to the ``science of morals'' (la science des mœurs) as the art of medicine stands to medical
science. When science has settled how every obligation of conscience has come into being,
% --- p. 13 ---
\opage{13}has been strengthened and recognized, --- what effects it has produced and what function it
performs in social life, --- then we shall also know to what degree it is possible and beneficial
to change it. Then a moral art can be practiced. In the meantime one neither can nor should sit
with folded hands, but must follow the best knowledge that our science makes possible at a given
time.\footnote[1]{\emph{Levy Bruhl}: \textit{La Morale et la Science des Mœurs.} p. 150. 223.}
Besides, it is only in appearance that there is so great a conflict between the ethical theories of
the philosophers. They are in reality different points of view on the fundamental moral fact (le
fait moral): social life, which stands at once as authority over against the individuals and as
the source of the highest goods to which they have access.\footnote[2]{Durkheim in \textit{Bulletin de
la Société Française de Philosophie.} 1906. p. 113---139.}

This significant direction of ethical inquiry, too, comes, partly against its will, to acknowledge
the difference between historical and philosophical ethics and the interaction between them. It is
not content to point out ``morality'' as a fact (le fait moral), but at once attributes value to
this fact. And it suggests that every obligation is to be valued according to its effects and its
significance for social life; but this at once presupposes a standard distinct from the facts
themselves. And the ethical problem concerns precisely the possibility and the grounding of such a
standard.

\chapterhead{II.}{THEOLOGICAL AND PHILOSOPHICAL ETHICS.}
% (chapter II begins on p. 13, after the end of chapter I)

1. Theological ethics stands in definite opposition both to historical ethics and to philosophical
ethics. --- It builds on tradition, on truth as given, on the ideal
% --- p. 14 ---
\opage{14}as historically revealed. To that extent it does, to be sure, fall under historical ethics: just
as positive jurisprudence sets out the system of the rules and commands laid down in a certain
country's juridical laws, so every theological ethics sets out the system of the ethical judgments
that follow from a certain ecclesiastical community's confession of faith. But theological ethics
does not acknowledge the method of scientific historical research, since the revelation on which it
builds is supposed to be due to a breaking through of supernatural powers that cannot be explained
by the physical, psychological and social laws on which historical science relies. It demands a
special position for its historical foundation and asserts that this must be regarded in an
entirely different way from the rest of history. --- It approaches philosophical ethics insofar as,
like the latter, it undertakes a valuation of historically given actions and forms of life. But it
does not undertake this valuation according to principles that can be shown to have their source in
human nature. It takes its point of departure not in human nature but above it, in the supernatural
revelation of an ideal. It builds on the absolute principle of
authority.\footnote[1]{On the concept of authority cf. my work \textit{Om Grundlaget for den humane
Etik.} Copenhagen 1876. Chap. III.}

Authority means the standing to command, and such standing can make itself felt in very different
ways. It can be grounded in physical power, so that it works through the fear it inspires. It can be
grounded in superiority of insight and ability, and work by furnishing a model for others, which
they seek in admiration to imitate. It can be grounded in the ability and the will to help and
protect, so that it works not merely by calling forth fear and admiration but by awakening
gratitude, trust and love. Human life cannot do without such authorities. The older generation is
the educator and the model of the younger. Parents are the authorities of their children, teachers
of their pupils; everywhere the one who has advanced further is an authority for the one who is not
so far along the road. The relation of authority is an important element in positive morality.

% --- p. 15 ---
\opage{15}It is always due, however, to a certain reflection, a beginning of thought, when the principle of
authority is expressly laid down and taken as unconditional. Men involuntarily follow what seems to
them great and eminent, and need no special command to do so. From the instinct of imitation and
from habit they follow even the insignificant and the indifferent, provided only that it works by
the power of repetition and that there are precedents. Only in times when criticism has awakened,
and when doubt replaces the secure resting in the half-unconscious, does one consciously set up the
principle of authority in its absolute form as the ultimate foundation of all ethics. When one
question calls forth another, so that it seems as if nothing stands fast, the need is awakened for
an absolute fixed point that is not drawn into the whirl. Human feeling and thought seem an
uncertain foundation for the ethical. The greatest question of human life: ``what is good and what
is evil?'' seems to need the strongest foundation, and what foundation can be stronger than the will
of an almighty being? When one can appeal to such a will, all discussion must cease, and the
clamorous doubts fall silent before the unconditional authority.

The good, then, becomes what agrees with God's will, and it is good only because it agrees with
God's will (bonum est, qvia deus vult). This was the doctrine that \textsc{Plato} had already had
occasion to combat (in the dialogue \textit{Euthyphro}), but which was first properly formulated
towards the close of the Middle Ages by \textsc{Duns Scotus}, who even stated outright that if God
had commanded murder, murder would be no sin. In so crude a form probably only few would acknowledge
this doctrine, but every theological standpoint contains an approximation to it. The unconditional
principle of authority is contained in the dogma of the infallibility of the Pope and in the narrow
literalism of Protestant orthodoxy. Through the living word of the head of the Church, or through the
written word of the books handed down by the Church, it is made known what is good or evil according
to God's will, and those to whom the word is addressed have only to receive it and follow it in
unconditional obedience.

Let us try to place ourselves at this standpoint, in order to see
% --- p. 16 ---
\opage{16}whether its presuppositions do not bring with them inner contradictions and insuperable
difficulties.

2. We ask first how one knows that something is God's will, and what guarantees one has of it. It
is of no use that I know that \textit{if} something is God's will it is good, when I do not know
with unconditional certainty \textit{what} is God's will in the particular case. This certainty must,
by the nature of the matter, be the most important thing to me at this standpoint. The faith on
which theological ethics builds therefore consistently becomes a faith in the Church, the bearer of
religious revelation and tradition. The Church guarantees me that its head is the rightful successor
of the Church's founder, or that the biblical books contain the true revelation. In the Catholic
dogma of infallibility, therefore, theological ethics finds its consistent conclusion. --- If, on
the other hand, I wish to learn God's will immediately, I must hold to the testimony of the Spirit
within myself; but how can I distinguish definitely between the divine voice and my own natural
feelings and thoughts? What standard have I for drawing the boundary between the divine and the
human? Here we lose the absolute firmness that was supposed to be precisely the advantage of the
principle of authority. We slide over from the theological standpoint to the psychological.

``Yes,'' it is said, ``we cling to the principle of authority precisely because we feel that we need
it. It is our own powerlessness and doubt that cast us into the arms of the unconditional
authority.'' --- But by stressing this personal need one raises a psychological problem in place of
the purely theological problem. Not authority itself but \textit{my need} of it then becomes the real
foundation. I justify my submission to authority by my need; but can this need itself, then, be
justified? Is it good to satisfy this need? How do I distinguish which of my many wishes and wants it
is good to satisfy, and which are to be suppressed? If, to decide this, one appeals to authority, one
goes round in a circle.

One comes to a similar circle in one's thinking when, as some theologians attempt, one wishes to
combine the two propositions:
% --- p. 17 ---
\opage{17}``the good is good because God wills it,'' and: ``God wills it because it is good.'' --- If the good
is what agrees with God's will, what does it mean, then, that God wills it because it is good? It can
evidently mean only that he wills it because it is his will. So by this we get no further. --- If,
on the other hand, God first recognizes something as good and then wills it, then there is a law
and rule to which his will bows, and his will is therefore not in and for itself the cause of the
good's being good. Even those who credit themselves with the most intimate acquaintance with the
psychology of the Godhead would yet become somewhat uneasy at the dividedness thereby posited in the
essence of the Godhead. Consistent theological thinkers have also seen that at the standpoint of the
principle of authority there can be no question of good or evil as regards the Godhead itself. God
can have neither conscience nor duty. Every transference of ethical ideas from men to God leads to
contradictions, and the only way out is to stop at the appeal to the divine will and to bow to it in
obedience. Otherwise one gives up the theological principle itself. As soon as one acknowledges that
there are other sources of knowledge of the good than the certainty itself of what God's will is,
one has really passed over from theological ethics to philosophical. The good is recognized --- that
becomes the main proposition; and that the good so recognized is God's will becomes an inference,
--- an inference of great practical importance, to be sure, since a supernatural sanction thereby
becomes possible, but an inference grounded on entirely different presuppositions from the main
proposition, and which could therefore very well fall away without the main proposition falling
away. ---

The unconditional authority utters a multitude of commands and prohibitions. As soon as reflection
awakens, we naturally try to bring connection and agreement among the particular pronouncements of
the authority. But this can be done only by our applying the fundamental law of our reason to them
and assuming that one pronouncement cannot conflict with another. But if we have the right to apply
this law of reason, why should other laws of our natural knowledge be excluded?
% --- p. 18 ---
\opage{18}And have we not, by that very application, already broken with the unconditional principle of
authority?

3. It is not, however, merely these inner contradictions that keep philosophical ethics from taking
theological presuppositions as its basis.

What called philosophical ethics forth in the first place, and ever anew awakens interest in it, is
the conviction that the ultimate ground of the ethical must lie in human nature itself. For all that
man is to be able to acknowledge as true and good, as beautiful and great, the standard must in the
end lie in himself. The principles of all understanding and all valuation, of all theoretical and
practical activity, must be found within him. However high the ideals stand above him, and however
mightily the majesty of the law makes itself known to his conscience, they are still ideal and law
\textit{for him} only through his free acknowledgment, through the assent he gives them by virtue of
his nature. Socrates therefore became the founder of ethics through his command: Know thyself! He
thereby set up the \textit{principle of personality}, or the principle of subjectivity, at once the
principle of free inquiry and of free conscience. Blind obedience as an abiding condition conflicts
with this principle; one thereby divests oneself of one's free personality and makes oneself an
impersonal machine.

Next, another motive also cooperates in the effort to develop a philosophical ethics, namely the
wish to make the ethical as \textit{independent} as possible of disputable presuppositions. Precisely
because judgments of good and evil intervene so powerfully in human life and condition the sharpest
oppositions within its domain, it is of great importance that no other disputed questions be mixed
up with the ethical ones. Ethical questions must therefore as far as possible be made independent of
religious and metaphysical problems. If ethics had to wait until agreement arose on dogmatic
questions, it would have to wait a long time. It is worth the trouble, then, to try whether agreement
in the ethical domain might not be greater than in the religious and metaphysical domain. And since
it is precisely the \textit{ethical} significance of the religious and metaphysical problems that is
to give them their great interest, it lies in the nature of
% --- p. 19 ---
\opage{19}the matter that one must have a foundation and a standard for the ethical which --- to a certain
degree at least --- must be independent of religious and metaphysical
assumptions\footnote[1]{A theological critic (in the journal \textit{Vidar} 1888) has found ``a
weakness'' in my standpoint, expressed in the phrases: ``as independent \textit{as possible} of
disputable presuppositions'' and ``to \textit{a certain degree} at least,'' which occur in this
passage. I chose these expressions deliberately, being convinced that here there can be question
only of an endeavor to reduce the disputable presuppositions to the least possible, an endeavor that
no single person can carry through completely. Every thinker has his limitation and must know that
he unconsciously operates with presuppositions that can perhaps only later be brought out and
tested. I have thus used those phrases out of scientific caution. They do not shake the principle:
that the ethical is to be investigated by itself alone and not \textit{derived from} religious or
metaphysical assumptions. This does not, of course, exclude that the ethical thinker may set up ideas
that one might perhaps call religious or metaphysical (cf. IV, 5 and XXXI, 3), provided only that he
does not make use of them in the \textit{grounding} of the ethical.}. For in attributing an ethical
significance to such assumptions one \textit{values} them, and the question must then be asked on
what foundation one takes one's stand in this valuation, and what standard one applies.

The general scientific maxim that principles must not be multiplied without necessity thus has a
practical significance here as well. In seeking to place itself outside the conflict between the
various religious confessions and between the various metaphysical theories (spiritualism,
materialism, etc.), philosophical ethics will not only attain a surer foundation for the ethical,
but will also work for the good of tolerance, freedom of religion and freedom of teaching. It will
counteract the hateful manner of discussion that consists in deriving immoral consequences from an
opponent's views.

The ethical problem therefore holds a very self-standing and independent position within philosophy
in relation to the other main philosophical problems. There is of course a connection between the theory
of knowledge, which investigates the foundation, presuppositions and limits of our knowledge,
metaphysics, which seeks to form a world-view, a general theory of the nature of existence,
% --- p. 20 ---
\opage{20}and ethics\footnote[1]{On the inner connection between the great problems cf. \textit{Den
menneskelige Tanke, dens Former og Opgaver} (1910). p. 246---254.}. But the connection is not so
direct and immediate as has often been assumed. And it will in any case be of importance to keep the
various problems apart from one another as much as possible; this does not exclude that they may
throw light on one another in many ways.

Philosophical ethics, which is to be a scientific doctrine of ethical valuation, must take its stand
on the same foundation as all other science. It does not seek to shake the principles, results and
hypotheses that other sciences have arrived at. It demands no special position. If it were to set up
presuppositions that conflicted with those on which other sciences build, it would come into the
domain of the disputable. It will go most surely when it builds on the common presuppositions. It is,
however, still a frequent assumption that ethics, at least at one point, namely in the question of
the application of the law of causality to the life of the will, should be compelled to set up
presuppositions that not merely deviate from, but directly conflict with, the general scientific
presuppositions. This question will be treated specially later on.

4. When theological ethics has been spoken of in the foregoing, it must be well noted that
\textit{Christian} ethics is not a theological ethics in all its forms and at all its stages of
development. Christianity did not begin with a theological system, any more than it began with an
ecclesiastical organization. Here, as everywhere where development takes place, a more indeterminate
form of life precedes system and organization. The ethics proclaimed in the first three Gospels is
in great part a doctrine of love of mankind. It is supported by theological ideas, and reference is
made to reward and punishment in a world beyond. But the God to whose requiting power reference is
made appears nevertheless chiefly as the ideal of perfection, which here, according to the context
(Matth. V. 43---48),\footnote[2]{In the parallel passage, Luke VI. 36, mercy stands in place of
perfection.} means:
% --- p. 21 ---
\opage{21}the ideal of love of mankind. Jesus founds the messianic kingdom, to which access stands open for
those who are pure in heart, who hunger and thirst after righteousness, who are peaceable and love
others as themselves. The spirit and the conditions of the new kingdom were entirely different from
those of the proud world-empires that history had hitherto known, and from those which the Jewish
people had imagined in its national fantasies.

What is significant in the oldest Christian ethics is the widening and inward deepening of love of
mankind that it inaugurates. Only through the dogmatic development of Christianity, founded chiefly
by Paul, did theological doctrines become of decisive importance for Christian ethics. Christianity
then received two principles, faith and love, which it has not succeeded in bringing into complete
harmony with one another. Differences of faith became barriers to love of mankind. The most decisive
standard of a human being's worth was now not whether he had an inward and loving mind, was pure in
heart and sought the truth, but whether he had the right faith. Only where differences of faith are
treated as Christianity itself had treated national differences and other external differences, only
there has the great ethical idea that was uttered in the oldest Christianity attained its full
development\footnote[1]{Cf. \textit{Om Grundlaget for den humane Etik.} p. 70 f. --- In that work,
in my discussion of this point, I did not take sufficient account of the oldest Christianity, which
lies before the more definite theological shaping.}.

Yet for a right understanding of the ethics of primitive Christianity it must be well noted that its
deepest motive was nevertheless not love of mankind itself, but the living, enthusiastic assurance
that the kingdom of heaven was at hand and that the end of all earthly things was drawing near.
Before this great hope all human ends and tasks lost their weight. There could be no question of
coherent ethical striving, and even the idea of love of mankind was not allowed to unfold its
consequences. Only when the Church pushed that expectation into the background did Christianity
become a cultural power. And the idea of love of mankind (which,
% --- p. 22 ---
\opage{22}incidentally, was not first conceived by primitive Christianity) has shown that it has vitality
beyond the limits both of that ecstatic expectation and of the ecclesiastical
dogmas\footnote[1]{Cf. my essay: \textit{Hedenske Sandhedssøgere} (in ``Mindre Arbejder'' I) and the
concluding section of my book \textit{Søren Kierkegaard som Filosof} (1892).}.

\chapterhead{III.}{THE PRINCIPLES AND METHOD OF ETHICS.}

1. Ethical judgments, judgments of good and evil, arise in the first instance without grounds. Just
as before a work of art or before nature we may relieve our heart by exclaiming: ``That is
beautiful!'', so before a human action we may relieve our heart by exclaiming: ``That is good!'' We
are guided by an immediate feeling that is awakened by the sight or the memory of the action. Which
actions we call good (or evil) will for the individual depend in the first instance on what
``positive morality'' (cf. I, 2) he has grown into. An ethical problem first arises when this
positive morality does not speak clearly, or when it proves to contain self-contradictions. Then
\textit{principles} are needed, \textit{guiding thoughts}, that can determine the choice between the
possibilities at hand.

One might doubt the possibility of setting up \textit{grounded} ethical judgments, and so the
possibility of a \textit{philosophical} ethics. Ethical judgments, after all, contain a demand; they
state \textit{what ought to be}, not merely \textit{what} is. Even when they pronounce praise or blame
of an action actually performed, they really mean a demand that such actions shall be performed or
left undone in the future. But does not all science aim at explaining, at finding the causes of what
is? Can there be a science of that which is not, but ought to be? Philosophical ethics goes
% --- p. 23 ---
\opage{23}a step further still: it wishes to be a doctrine of \textit{what ethical judgments we ought} to
pass! How is such a science possible? Add to this that judgments of good and evil spring from
feeling, --- but ``there is no disputing about feelings.'' And experience shows, too, that
discussions of ethical questions frequently end with each appealing to his own feeling.

To throw light on this problem it must first be remarked that ethical judgments aim at stating the
means and ways of attaining \textit{something that has value for us}, and which we therefore more or
less consciously make our \textit{end}. Ethical valuation regards the particular action, the
particular quality or the particular social institution in relation to an end. The more involuntarily
the judging and the valuing take place, the less consciousness we have of the end; we are seized
immediately and rest wholly in the thought of the object that is valued. Valuation can take place
quite instinctively. But when in cases of doubt the valuation is to be grounded on definite
principles, it must first be settled to what end we relate the action. The judgment on the action
will depend on whether, on closer examination, it proves to further or to hinder the end
presupposed. Every action will bring something about; otherwise it would be impossible and absurd.
This remark is not unnecessary, for great ethical thinkers, especially \textsc{Immanuel Kant}, have
held that it was a degradation of the will to lay particular weight on the effect and to grant the
end a determining influence on the valuation. Against this holds \textsc{Schleiermacher}'s apt
remark\footnote[1]{\textit{Ueber den Begriff des höchsten Gutes.} Werke zur Philosophie. 2 Band. p.
452.}: ``Why then do I act, if I wish to bring nothing about?''

We already have firm ground under our feet once it is established that every action must have an
end. For then it is possible to examine whether the action really serves this end or not. The
relation between the road and the goal, between the means and the end, can be scientifically
established. If the end stands fast, certain definite means must be employed. The relation
% --- p. 24 ---
\opage{24}between means and end is, after all, the same as that between cause and effect: what we call a
means is the cause that brings about the attainment of the end. Here, then, we see the possibility
of a science of what ought to be: the word ought expresses \textit{the demand for means to a given
end.} In every need or impulse lies such a demand, and what we call conscience is a deep inner need
to maintain or to produce something of value in the world. --- Through the necessary relation
between end and means a scientific ethics becomes possible. But the end is presupposed as given ---
and, as will appear, the problem arises again when it is asked \textit{which} end we presuppose as
the highest, --- which value is for us the fundamental value.

We arrive at the same result to which the consideration of the action in its relation to the end
has led us, by a closer consideration also of the proposition that judgments of good and evil are
utterances of feeling. The utterance of feeling here presupposes that the perception or the idea of
an action has moved us inwardly, has furthered or hindered an involuntary striving, a conscious or
unconscious interest in us. Feelings of pleasure and pain are connected --- as \textit{the biology of
the life of feeling} shows --- with the furthering or hindering of our life, or at any rate of a
part of our life. This determines what acquires value for us and can therefore be made an end. All
value depends on the conditions of life and on the stage of life. \textit{The human life of feeling
is therefore intelligible only in connection with the conditions of human life.} And only in its
simplest forms and in the most violent affects is feeling wholly blind. Its definite, fully developed
form it acquires only by being bound up with ideas. Joy presupposes the idea of an object that
arouses pleasure, sorrow the idea of the loss of one. Hope and fear presuppose ideas of the
possibility of certain events that arouse pleasure or pain.

Here too, then, we get firmer ground under our feet. ``There is no disputing about feelings''; but
one can dispute about the ideas to which the feelings are attached, since one can examine the
relation of these ideas to
% --- p. 25 ---
\opage{25}experience. And one can examine the very states that prove valuable, and seek to find under what
conditions they occur; from this one will then be able to draw conclusions about the way in which
such valuable states can be brought about in the future.

Not only in the relation between means and end, then, which plays a part in all human actions, but
also in the relation between experience and the ideas by which the feelings that find vent in
valuation are determined, we see \textit{a possibility of grounding ethical judgments, and so of
setting up a philosophical ethics.} And the two relations cannot be separated; for an end always
presupposes a feeling. Only what sets our feeling in motion can we make an end.

Ethical judgments concern human actions and human ordinances of life (institutions), the latter,
however, only insofar as they have proceeded from human actions and can in turn be altered by them.
For the most part, now, ethics has started from the action that is to be valued, and has not
examined valuation itself. It will appear that one must begin by discussing the condition of
valuation itself.

2. The action has its origin within the individual, in its instincts and impulses, its thoughts and
feelings. But after having developed there, it passes over as an effect into the external world.
Here, then, the question arises to which part of this whole course ethical valuation is chiefly to
keep.

A perfect valuation would require that the whole course of the action be taken in. The ethical
characterization of an action as good or evil would then presuppose that one could trace it all the
way from its first germs in the individual's inner life and through all its branching and distant
effects in the surrounding world. Ethicists have not laid equal weight on all parts of the action's
whole history. Some have taken their point of departure from the action's own origin, that is, from
the motives, the disposition, that give birth to it. The external effects of the action, they hold,
cannot be surveyed and are not wholly under the dominion of the will; it is the inner action in
which the will shows itself, and therefore it alone can
% --- p. 26 ---
\opage{26}become the object of ethical valuation. Others, on the contrary, start from the action's
consequences and effects in the external world. It is these alone, they hold, that give the action
practical significance. Valuation must keep to them first and foremost, especially since experience
shows that actions springing from the same motives can have very different effects, and actions
springing from different motives can have the same effect. Only secondarily does valuation extend
to the motives and the disposition, insofar as it can be shown that certain definite motives
generally lead to a certain definite manner of acting. Here, then, a distinction is made between
the valuation of the actions and the valuation of the acting person; the former is primary, the
latter secondary.

This opposition has played a great part in the history of ethics. One might describe it as an
opposition between \textit{subjective ethics} and \textit{objective ethics.} In the more recent history
of ethics subjective ethics is represented by Kant, objective ethics by \textsc{Bentham}, and the
debate in the field of philosophical ethics lies essentially between the tendencies that proceed
from these two thinkers. A decision of the dispute must depend on a closer determination of the
relation between subjective and objective ethics. Such a closer determination I seek to give in
this chapter. I shall rely especially on the consideration that, even if one agrees with Bentham
that an objective valuation of actions is possible only if one takes one's point of departure from
their effects and consequences, \textit{valuation} is nevertheless always \textit{a subjective
activity,} and every principle of valuation must presuppose certain subjective conditions. The
dispute between the schools of Kant and of Bentham has turned all too exclusively on the relation
between the motives of the action and its effects. Perhaps the problem will present itself more
clearly and become easier to solve if one goes further back and asks about the motives of
valuation itself, or about \textit{the psychological presuppositions on which the very circumstance
rests that a general and objective valuation of human actions is striven for.} Just as theoretical
knowledge leads back by analysis to certain principles, which it is the business of the theory of
knowledge to point out, and which stand in close connection with
% --- p. 27 ---
\opage{27}the nature of the human faculty of knowledge, so too practical valuation must prove to build on
principles that have their ground in the human mind, even if the question here should be somewhat
more intricate than in the theoretical field. Both the theory of knowledge and ethics lead us back
to subjectivity as the ultimate foundation. The task is precisely to show how, \textit{in spite of
the unavoidable subjective element in the foundation}, \textit{an objective knowledge and an
objective valuation} can be reached. Here, then, is a parallel between the problem of knowledge and
the ethical problem to which attention has not always been paid.

Just as \textit{the foundation (the motive of valuation)} is the subjective principle of ethics, so
the standard by which actions are valued from this foundation is the objective principle of ethics
and determines the content of ethics (cf. I, 1). For valuation presupposes two things: a motive that
drives one to exercise the valuing activity of mind, and a criterion according to which the
valuation is carried out\footnote[1]{In my work \textit{Om Grundlaget for den humane Etik} (1876) I
did not yet sufficiently distinguish between the concepts \textit{foundation} and \textit{motive},
or more clearly: between \textit{motive of valuation} and \textit{motive of action}, whereas the
separation between foundation and content is already made in the work named. --- In my essay
\textit{Filosofiske Problemer} (University programme, Nov. 1902) I introduced the concept of
\textit{fundamental value} as the presupposition of the motive of valuation. Every motive presupposes
experience of a value, and the motive of valuation presupposes experience of a value that for us
conditions all other values.}.

To show more fully the significance of the point of view brought forward here, I shall try to give
an account of how both the subjective and the objective principle of ethics (the foundation and the
content) spring from the nature and the manner of development of ethical valuation. The more
precise relation between them will thereby gain in clarity. --- But since the valuation of actions
and conditions of life not only takes place among different individuals and peoples on the basis
of different motives and according to different standards, but is also conditioned in each single
individual by composite motives, it will be necessary for us to think of the conditions as simpler
than they are in reality. No human being
% --- p. 28 ---
\opage{28}pursues just one single end; life is varied and manifold within us as outside us, and the
various ends contend for sole dominion within the consciousness of the single individual as well as
within society. The single moments as well as the single interests demand their right; the
individual feels itself a little world and yet a member of a larger world, or of several worlds,
each of which demands its right. Therefore human conscience, the feeling that determines
valuation, is as a rule a very composite product. It will then be expedient to examine particular
elements of this composite whole each by itself. The following account proceeds from the simplest
to the most comprehensive valuations; it is not a historical account, but orders the standpoints
according to the range of the ends that become determinative for valuation.

3. It is a fact that men value their own and others' actions and call them good or evil according
to the outcome of this valuation. How, now, is such a valuation possible?

We think first of the simplest case, namely where it is the acting subject that judges its own
action, without having any consciousness of other beings to whom regard might have to be paid. The
acting and valuing subject is thus thought of as a little world by itself. We make an abstraction
in order to get as simple an example as possible.

The first presupposition for an action's being valued is that it is remembered; it must not,
then, have vanished from consciousness once it has passed over from resolve into reality. The
image of the action must be capable of being called up again as an object of contemplation. But
this is not sufficient. Such an image in memory might in and for itself stand as something wholly
indifferent that did not set the mind in motion. Only if the action has in one way or another, at
one stage or another of its development, intervened in the individual's whole state and thereby
aroused either pleasure or pain, will the image of the action be able to arouse pleasure or pain.

% --- p. 29 ---
\opage{29}Here there appears in its very simplest form the important truth, already touched upon, that
\textit{all valuation of actions presupposes a subject with the capacity to feel pleasure or pain.}
Valuation presupposes that a demand is made on the actions, which they can satisfy in greater or
less degree. But such a demand stands wholly \textit{unmotivated} if the action is unable to arouse
pleasure or pain. This, as already remarked, is only another expression for the fact that valuation
presupposes an end by which the action can be measured. Only that can become an end which has value
for us, and which does not fall to our lot immediately but can be attained only through work.
Immediate value makes itself known through a feeling of pleasure, while a feeling of pain will
arise when the action hinders the realization of an immediate value.

In the simple case we imagined, the feeling that conditions the end, and which is thus to be
satisfied by the action, can only be the individual's own. The individual will then stamp the
action as good or evil according to the way in which it has intervened in its life. What character
and significance the valuation more precisely acquires will depend on whether the feelings of
pleasure and pain are determined only by, and answer to, the momentary states each by itself, or
whether they are determined by regard to the individual's life as a whole and to the conditions
under which it develops.

4. The lower the life of consciousness, the more isolated and self-standing are the single moments
in relation to one another, and the smaller the part played by memory and by the thought of the self
as a whole that comprises the single moments of life with all their content. It is then only a
half-unconscious instinct that keeps the individual from being wholly absorbed in the single moment.
The impulse of self-preservation leads the individual, in the single present moment, to take regard
also of the future and to make use of the experiences of the past. But the more the individual is
absorbed in the single moments, the less possibility there is of a valuation, since no comparison
and interaction between the different states can take place. The action itself is perhaps
forgotten at the moment when its effect makes itself felt
% --- p. 30 ---
\opage{30}in consciousness. Each moment is filled in its own way, with its own state of feeling, which has
no influence on the other moments. The single moment stands as an absolute egoist over against the
other moments, and will not give up its claim to satisfaction for their benefit.

Here there appears the possibility of a standpoint at which all valuation falls away, because,
though feelings of pleasure and pain stir, they answer merely to the state of the moment, not to
life as a whole. No end is set up; even the instincts with their unconscious ends are pushed back;
and when every regard to an end falls away, valuation becomes impossible. It is a childlike
standpoint; but it is possible to attempt to make it the definitive standpoint. In the history of
ethics this has been attempted by Aristippus of Cyrene\footnote[1]{On the relation between
Aristippus's standpoint and Søren Kierkegaard's ``aesthetic stage'' see my book \textit{S.
Kierkegaard som Filosof.} p. 84. Note.}. It is the principle of \textit{the sovereignty of the
moment} that he maintains. This is the most radical ethical standpoint that can be conceived. It has
the fewest possible presuppositions, so few that all valuation falls away. For the good, for
Aristippus, was the purely momentary feeling of pleasure. Why --- so the thought runs --- should one
moment be sacrificed or subordinated to another? The one has in and for itself just as much right to
be as the other. --- At this standpoint --- and only at this --- the good coincides wholly with the
feeling of pleasure, evil with the feeling of pain. --- It follows of itself that such a standpoint
cannot be a wholly primitive and naive standpoint. It will be carefree and thoughtless, as the child
is, --- but the very express willing of it shows that it is only an imitation of the childlike
standpoint. A certain art and self-control are needed to check the involuntary craving beyond the
present moment. The instincts themselves already establish a connection between the single moments
that the carrying through of the standpoint of the moment on principle cannot bear. Herein lies
precisely a great difficulty for the practical maintenance of
% --- p. 31 ---
\opage{31}this standpoint. But it is of great significance for us to see that its justification has been
maintained in theory. By the nature of the case this could happen only within an advanced culture.
The simplest standpoints are by no means always those that appear earliest in history. But there is
a natural tendency to seek out limiting cases, to see how our concepts are formed when we take up
the extreme points of life; and the attempts at an ethics of the moment that exist are very
instructive. Perhaps still more perfect expositions of this standpoint could be given than exist
historically.

Such a standpoint does not lack justification. It might be objected that it abolishes all ethics,
since it excludes valuation, and since all ethical endeavor presupposes that a more limited regard is
subordinated to a more comprehensive one. Of such a subordination there can be no question here, since
life is thought of as consisting of absolutely sovereign moments! --- But to this it must be
answered that ethics must itself prove its justification for setting up the demand that a
satisfaction be given up in one moment for the benefit of other moments. The possibility presents
itself that every moment might be filled by an immediate value of one and the same degree, or that
there might at any rate be a harmony between the values that fill the single moments. The burden of
proof lies on him who demands sacrifice and resignation. This can be denied only by an absolutely
ascetic view, i.e.\ a view for which asceticism is an end, not merely a means. Every moment has a
natural right to be, and has as well, if one may so put it, its instinct of self-preservation, since
the need for full satisfaction in the moment yields to other promptings only under a certain
resistance. Even at a purely instinctive stage of life this resistance is felt.

5. If the principle of the sovereignty of the moment could be carried through in practice, no
reasoning could overthrow it. But there is hardly any conscious individual in whom there do not stir
instincts \textit{and impulses} that lead beyond the moment. And in human individuals \textit{memory and
expectation} will constantly assert themselves and become points of attachment for feelings of
pleasure and pain that are determined by the individual's
% --- p. 32 ---
\opage{32}lasting or constantly recurring conditions. In every reflection on himself and his conduct the
individual will rise above the single moments in their diversity and isolation, and his feeling will
(in the moments of reflection at least) answer to the way in which \textit{the whole of life}, and
not the single moment, stands before consciousness. A \textit{real} self first forms when there is a
circle of feelings and ideas that make up a firm core in consciousness, even if they do not assert
themselves at every moment. This firm core is made up above all of the values the individual knows,
the ends it sets itself. These values and ends govern its whole life, make up the interest that
gathers life into a whole and lets it stand as more than a series of isolated moments. Every strong
feeling has a tendency to spread and to govern the whole life of consciousness, to color all that
makes up the content of consciousness. Through this expansion of feeling, not only can a mood
awakened at a moment extend its effects over several moments, but a real self can also form, a
ruling fundamental mood that marks the level of the individual life of feeling as against the
momentary fluctuations.

Now where the state of feeling of the single moment, regarded as the effect of the individual's own
action, collides in consciousness with the state of feeling determined by the idea of the whole of
life, there will arise a new feeling, determined by the relation between them, a relation that may be
either harmonious or disharmonious. In this feeling of the relation between the momentary state and
the state determined by regard to the whole of life consists \textit{valuation.} It expresses the
relation between two values. The capacity for such valuation is conscience, as it might be thought
to express itself in a wholly isolated individual. In the widest sense conscience is a
\textit{feeling of relation}, and presupposes only that there asserts itself in consciousness a
relation between something central and something peripheral, something more and something less
comprehensive. The single moment, and the single action that has given it its character, are valued
--- at the standpoint here presupposed --- according to the way in which they can be fitted in as
members of the
% --- p. 33 ---
\opage{33}whole of the individual life. Here there is at work not only the tendency of the real self and the
ruling ends to determine everything in consciousness, but also a natural need, deeply grounded in
the general nature of consciousness, for unity and coherence in the states of the soul. We wish our
life to be a whole, not only because only so can our highest interest in life prevail, but also
because unity and coherence are goods in and for themselves. It is \textit{the formal side of
personality} that expresses itself in the need to will one thing, so that everything else we desire
besides this one thing is only a means to it. It is a spiritual impulse of self-preservation that can
never be wholly lacking, and rarely as one finds full consistency in reality, it is just as impossible
that an endeavor to be consistent can ever be wholly lacking\footnote[1]{With regard to the
psychological ideas and laws used here cf. my \textit{Psykologi} V B, 5; VI E, 1 and F, 4.}.

Here the task presents itself to the individual of bringing about harmony between the single parts
of its life. It is a task that is hardly solved involuntarily, without any conscious striving, in any
human being. Here, therefore, the valuation of earlier actions according to the way they have
contributed to the solution of this task also acquires significance for the individual. Valuation is
thus not only \textit{made possible} by the central feeling that answers to the whole of life, but is
also \textit{motivated} by it. A fine and developed sense for what serves the individual life, whose
single members the moments are, is a condition of its existence and development. It is a kind of
higher instinct of self-preservation and need not confine itself merely to what the existence of
physical life requires, but can comprise ideal needs no less than physical ones. --- In the ethics of
\textsc{Plato} and \textsc{Aristotle} this harmonious individualism has found its most characteristic
expression. Ethical virtue is regarded by them as a spiritual health and harmony. \textsc{Aristotle}'s
definition of virtue (as the mean determined by the nature of each single individual) is of
particular interest in this respect. Yet the whole ethics of the thinkers named is not exhausted by
this individualistic doctrine.

% --- p. 34 ---
\opage{34}To designate this standpoint I use the word ``individualism'' and not ``egoism,'' since by the
latter word one should preferably think of a conscious setting aside and subordination of others'
welfare to one's own. The individualist need not be an egoist, but he can become one.

6. At such an individualistic standpoint the task will be not only to determine how \textit{much}
energy must be spent in the single moments, but also to employ the energy at one's disposal in as
\textit{varied and diverse} a way as is compatible with the interest of the whole of life. By a
psychological law of nature the life and freshness of the feelings are conditioned by the different
states standing in a certain relation of contrast to one another, while uniformity and repetition
have a damping or dulling effect. --- Besides, the natural course of life will also bring with it that
the individual is gradually no longer absorbed merely in physical self-preservation, but acquires
needs of a more ideal and more composite character. The individual's life will attain the greater
fullness, the more different directions it can spread into, the richer and more manifold a content it
can comprise, without its unity and collected strength being thereby weakened.

Since the individual is here thought of as having set itself ends that point beyond the moment and
beyond the satisfaction of a single side of its nature, there can be question of a demand, an
``ought,'' a law, and the single moments and the single actions are valued according to how far they
serve the end set up. It might, for example, be intellectual or aesthetic enjoyment and development
that is set as the end, and with that the demand will be given that is to be made on the single
moments and actions.

The \textit{ethical law} at the standpoint of individualism is found by formulating what the
harmonious relation between the interest of the whole of life and the need of the single moments
requires. It will consist of two main commands, a negative and a positive: 1) the single moment must
not have greater independence than answers to its significance within the whole of life; 2) but on
the other hand one is to live in the single moment as richly
% --- p. 35 ---
\opage{35}and fully as is compatible with the preservation of the whole of life and with the attainment of the
end that stands as the highest.

Good, then, here becomes the action that preserves the whole of life and gives the content of life
fullness and life; evil becomes the one that has a more or less marked tendency to burst or to narrow
the whole of life and its content. Thus the single moment and the single impulse are evil in their
rebellious isolation from the rest of life; and evil will lie all the deeper and become the object of
all the stronger rejection in the valuing feeling, the more it is willed, i.e.\ has come about as
the fruit of deliberation and choice, not merely of momentary prompting. The strength with which
valuation asserts itself will depend on the strength with which the interests of the whole of life
assert themselves in the individual's central feelings.

7. What held of the sovereignty of the moment holds of individualism, or the principle of the
\textit{sovereignty} of the individual: it cannot be overthrown by reasoning. An absolute individualist
or egoist would be absolutely unassailable. When the assertion of his own life, its existence, its
fullness, is the only end he acknowledges, there is no \textit{logical} transition from this
standpoint to another. With sufficient consistency and self-control, and with sufficient power over
circumstances, an individual will then be able to carry through the maxim: Thou shalt act so that
thou thyself always standest as the only end, as the absolute center of the world. --- If a change
is to take place, the feeling that determines the fundamental value, and with it the valuation, must
be altered by being attached to a more comprehensive circle of ideas than the one that concerns only
the individual's own life. Until that has happened it is of no use either to appeal to conscience;
for this is, as we have seen, a feeling of relation, the expression of the relation between the
central and the peripheral in the individual's life of feeling, and is thus itself determined by
what the central is.

Philosophical ethics has often, especially in earlier times, believed that it must lay claim to being
a pure science of reason that appealed to no presuppositions other than those lying in the nature of
human reason. But this conflicts
% --- p. 36 ---
\opage{36}with the character of ethics as a practical science. Action can be valued only according to ends,
and the fixing of an end presupposes \textit{feeling of pleasure and pain} in the subject that fixes
it. Apart from a consciousness with the capacity to feel pleasure or pain, therefore, ethical
valuation has no meaning. But on the other hand the mere capacity to feel pleasure and pain as yet
implies nothing about \textit{how comprehensive the circle of ideas is to which the feeling of
pleasure and pain is attached.}

From the standpoint of the moment to the individualistic standpoint the transition took place
through the forming in consciousness of feelings that were determined by the interests of the whole
of life, and which could therefore become strong links between the changing moments. Supported by
them the individual feels its unity in spite of the constant change. If there is a wider standpoint
than that of individualism, there must be feelings that, by stirring in the single individual, can
bind him to a more comprehensive whole in a way similar to that in which the single moments and
promptings in him are bound to the whole of the individual life. There must be a power that binds
the single individuals to one another and abolishes the isolation between them.

8. Only approximately can individualism be carried through in practice. The sovereignty of the
individual, the conception of the individual as a closed and wholly independent whole, proves to rest
on a violent and artificial abstraction. The individual comes into being out of the race and lives
its whole life as a part of the life of the race, with an organization in which it inherits the
consequences of earlier generations' doing and suffering, under conditions of life and in a spiritual
atmosphere that the development of the race has brought about. And just as the instinct of
self-preservation abolishes the isolation of the single moments of life and thereby becomes the
foundation of feelings determined by the interests of the whole of life, so in \textit{the
sympathetic instincts} there stir forces that abolish the isolation of the single individuals and
assert the conditions of life of the race within them. In the most primitive form the sympathetic
instincts appear in the founding of the family relation. However much the forms and institutions of
the family relation vary among different peoples and at different times, and however loose and
% --- p. 37 ---
\opage{37}arbitrary in particular the relation between the two sexes often proves to be, there is yet one
relation that by the nature of the case cannot be abolished or essentially altered, namely the
relation between mother and child. In this relation sympathetic feeling grows immediately out of
natural instinct. It forms the firm core by which the higher forms of family life become possible.
Here the ground is laid for a community of life within which the feelings of sympathy can be
cultivated and attain such strength that they can gradually comprise larger circles. Mother-love
nevertheless continues always to stand as the model and standard of all sympathy, both as regards
strength and purity. It is a reminder that the family relation is the constant source of the
sympathetic feelings, when the general love of mankind finds its most striking expression in the
proposition that all men are brothers.

This is not the place to enter upon a closer psychological inquiry into the nature of the sympathetic
feelings and into their different character according to the elements of which they consist and the
range in which they express themselves (cf. on this my \textit{Psykologi} VI C)\footnote[1]{As may be
seen not only from the text here but also from my \textit{Psykologi} (VI C), I use the word
\textit{sympathy} in the sense of a fellow-feeling in which one immediately makes others' pleasure or
pain one's own, or, as one may also put it, puts oneself in others' place and feels life in the way
they must feel it. The word is used in this sense not only in the older Danish psychology
(\emph{Treschow}: \textit{Om den menneskelige Natur.} Copenhagen 1812. p. 355. F. C. \emph{Sibbern}:
\textit{Psykologisk Pathologie.} Copenhagen 1828. p. 27), but also in later works on the psychology of
feeling (\emph{Domrich}: \textit{Die psychischen Zustände.} Jena 1849. p. 218. \emph{Bain}: \textit{Emotions
and Will.} 3. ed. London 1875. p. 111) and in ethical works (Hume, A. \emph{Smith}, and in our day
\emph{Wundt}: Ethik\textsuperscript{1} I. p. 208). --- Apart from this sense the word sympathy is used
in three other senses. --- It can denote a passive, reflex-like mirroring of others' states. The sight
of others' enjoyment or suffering makes us ourselves involuntarily attuned in a similar way, so that
we may even, for example, feel pain in our own leg when we see another's leg maimed. Suffering here
as a rule works more strongly than enjoyment. --- It can next denote a capacity for imagining others'
feelings, whether we share them or not. Here \textit{tact and inference} play a larger part than an
immediate feeling of unity. The one with whom we sympathize in this way still continues to stand as
an object. --- Finally it can mean a feeling of \textit{satisfaction in finding again} in others
feelings of pleasure or pain that we ourselves would have in similar cases. It is a strengthening, a
self-confirmation, that we thereby experience. --- It is idle to dispute about which of all these
senses is the ``right'' one. --- Sympathy in the three last senses of the word can come into
opposition to sympathy in the first sense of the word. But the different kinds of sympathy are
related, and in particular the three last kinds (the reflex, the intellectual and the interested
sympathy) can pave the way for the first through psychical transitions and displacements.}. Where
sympathy has attained its full purity,
% --- p. 38 ---
\opage{38}it is a feeling of pleasure or pain determined merely by other beings' feeling pleasure or pain.
As regards its range, the most significant point in the history of its development was when it
widened from merely comprising family, nation and race to holding for the whole of mankind. Greek
philosophy already (the Peripatetic and the Stoic schools) led to the idea of a general love of
mankind, grounded on all men's natural belonging together in one great community\footnote[1]{Cicero:
\textit{De finibus} V, 65 cf. \textit{De officiis} I, 16---17; III, 6. --- Cf. my essay:
\textit{Hedenske Sandhedssøgere} (Mindre Arbejder I. p. 159 f.; 167---174). --- Several theological
critics have complained that it is not shown whence the sympathetic feelings arise. From my
\textit{Psykologi} they could have seen that I assume a physiological point of departure and a
psychological-sociological course of development for these feelings. As regards degree, kind and
range they are subject to definite conditions, which I have indicated briefly. It is remarkable how
frequent is the opinion that self-love is the only natural feeling. This opinion is found especially
at two opposite standpoints: on the one side among skeptics and blasé men of the world, who appeal
here to their supposed experiences, and on the other side among orthodox theologians, according to
whom all true love comes from a supernatural source. Indeed, even an enlightened and critical
theologian such as \emph{Adolf Harnack} denies that love of mankind can be a ``product of nature,''
and holds that only self-love is natural. ``Enlightenment left to itself'' leads in his opinion to
egoism and ``general dissolution.'' See his lecture at the Evangelical-Social Congress in Frankfurt
a. M., 17 May 1894.}. But this idea first acquired greater historical significance by becoming, as a
chief command of Christianity, one of the guiding thoughts of a great world religion.

9. Ethically regarded, the great significance of this is that the horizon
% --- p. 39 ---
\opage{39}has been widened, so that the feeling of pleasure and pain is now determined not merely by the
individual's own fate, but by \textit{the conditions of life of a community of which the individual is
only a single member.} When such feelings come to prevail, actions will be valued according to the
relation in which the states they produce stand to the interests of this community, the larger
whole. The standard is now not taken merely from the individual's isolated life; this stands as
belonging to a larger order of the world, and the question is whether the actions hinder or further
that order.

It would not be correct to say that at this standpoint sympathy becomes one with the ethical feeling
or conscience. Conscience appears here too as a feeling of relation, determined by the relation
between what is the prevailing or central feeling in the individual and the results of the actions.
The difference between this standpoint and the individualistic one is that the foundation is more
comprehensive; the valuation may therefore turn out differently: one and the same action may be
called good from the individualistic standpoint, evil from that of the race. When the individual
immediately feels its own personal interests as subordinate to regard for the welfare of the great
whole, of which in sympathy it feels itself a single member, the ethical feeling appears as
\textit{the feeling of duty.} Even at the standpoint of individualism one could speak of duty. For the
concept of duty need contain only the formal relation between a more limited and a more
comprehensive regard, that the latter is to have precedence over the former. The individual can in
the single moment feel itself bound by regard for the whole of its own life. The feeling of duty
springs --- where duty is more than something imposed from without --- from within the personality,
and is conditioned by the demand for, and the need of, preserving the coherence of one's life through
the assertion of the highest ends and interests of life in the single moments and under the special
circumstances, --- of finding oneself again, according to one's central nature, in the single state
or action. In duty, then, a fundamental will expresses itself: we \textit{ought} because in our
inmost ground we \textit{will.} Our central willing becomes an ``ought'' because peripheral tendencies
of will can rise up
% --- p. 40 ---
\opage{40}in opposition to it, and because we nevertheless wish to remain ``true to ourselves and to our best
striving.'' But at the standpoint we are now describing, regard for the individual's whole of life
will itself again be subordinate to regard for the whole of life of the race, and likewise the
sympathy that has a narrower community as its object will be subordinate to the sympathy directed
toward a more comprehensive community. Or rather: the individual has here made the interest of life
of society or of the race its own highest interest of life, the object of its own central willing.

Seen from another side, the ethical feeling appears in its further development as \textit{the feeling
of justice.} For the sympathy that lies at the foundation will, when it is more than a blind
instinct, be guided in its manner of expression (as regards both kind and degree) by regard for the
peculiarity of the beings to which it applies. And when its range is widened to all beings with the
capacity to feel pleasure and pain, it must in each single case express itself in such a way that the
single being to which it is attached is satisfied according to its peculiarity, without other beings
being thereby injured in their equally marked peculiarities. Every difference and inequality that is
made to count must be grounded in what the welfare of the realm to which both the individual himself
and the other beings belong requires. Sympathy, regarded from the active side, is a need to impart;
this imparting must, if it is not to happen blindly, be determined by principles, and these must be
taken from the nature of sympathy itself. When sympathy is universal, the difference it makes in the
distribution of its goods can be due only to the reason that these goods, if distributed in another
way, would not really, or not in so high a degree, be goods for those to whom they fell, or would not
make for so great an advance of all the interests of life. On the ground of sympathy the ethical
feeling thus develops into a feeling for distributive justice. Just as the individual feels itself
as one among many, so it also regards every other individual as one among many, insofar as no
special reasons motivate special regards. --- At the standpoint of individualism an analogy to
distributive justice could be found, since the relation between the
% --- p. 41 ---
\opage{41}different moments and promptings corresponds to the relation between the different individual
beings at the standpoint we are speaking of here. \textsc{Plato} used the expression justice
precisely of the harmonious relation between the different sides or parts of the single
individual's life of consciousness.

It is this logical character of the ethical feeling, the strictness with which it demands grounds
for every difference and inequality that is laid down, that has chiefly led people to think that the
ethical law is an expression of pure reason. What \textsc{Kant} formulated as the content of the
categorical imperative expresses in reality only the impartiality, the disregard of irrelevant and
accidental considerations, which the ethical feeling built on sympathy demands. Kant demanded that in
valuing an action one should take one's stand at a universal standpoint, and he held that every other
standpoint contained a contradiction. But for such a contradiction to come forth, quite definite
psychological conditions must be present; a collision between two interests must be felt. It holds
here again that conscience presupposes a relation in which neither of the terms is vanishingly small.
That demand is thus by no means self-evident. It came forth, after all, only at a definite stage of
historical development. In Kant himself the idea of this demand arose under the influence of an
inquiry into the development of human social life\footnote[1]{Cf. my work: \textit{Den nyere
Filosofis Historie}\textsuperscript{2}, II p. 72---76 and my essay: \textit{Rousseaus Indflydelse paa
den definitive Form for Kants Etik} (Oversigt over det kgl. danske Videnskabernes Selskabs
Forhandlinger 1896).}, an inquiry in which, however, he overlooked precisely the development of the
sympathetic feelings under the influence of social life, which is why the categorical imperative too
came to stand as a mystical power in human nature. Kant regarded all feeling of pleasure and pain
(sympathetic feeling too) as egoistic, and thereby cut himself off from the possibility of
understanding how interests and ends can develop in man that lie beyond his own existence and
enjoyment.
The ethical law therefore comes, in
% --- p. 42 ---
\opage{42}Kant, in a mysterious way, as a proclamation from another world, into man's empirically given
nature.

10. An \textit{ethical} law comes into being when the conditions of life of a more comprehensive whole
are formulated in definite thoughts. At the standpoint we have before us here, \textit{the standpoint
of the ethics of humanity}, its content can be nothing other than the maxim that actions are to lead
to \textit{as great welfare and progress for as many conscious beings as possible.} And this in turn
comprises two main commands, a negative and a positive: 1) to no single individual ought more to be
allotted than is its due according to the place that, by its peculiarity, it occupies within the
race; 2) but on the other hand the faculties and impulses of every individual are to be developed and
satisfied as fully and richly as is compatible with what the life of the race as a whole requires.
These two commands follow with logical necessity from the concept of society as a multiplicity of
conscious beings bound into a unity. It conflicts with the unity of society that a single individual
or single individuals be arbitrarily preferred or set aside in favor of others: every special
position must be grounded in what the common conditions of life require; but on the other hand a
society is the more perfect the more freely and independently its single members move, the more
different possibilities they realize, while at the same time the unity is preserved and acquires an
ever more inward character and a higher validity.

When the ethical feeling on the basis of sympathy develops into the feeling of duty and of justice,
it will be the principle stated in this law that is in the end the standard for the ethical
judgments that are passed\footnote[1]{Whereas I have here gone from the concept of welfare to the
concept of justice, I take the opposite road in an outline of lectures from 1901 which is reprinted
under the name ``Human Etik'' in \textit{Mindre Arbejder, Tredie Række.} There I set up the
hypothesis that justice is the supreme ethical quality of character (``virtue''), and then examine
--- after having given a description of it --- how far and in what way such an assumption can be
grounded. (1913).}. At this standpoint it also follows of itself that ethical judgments concern not
merely the individual's own actions but also those of other individuals, and that the
% --- p. 43 ---
\opage{43}standard must be the same. Only, with regard to its own actions the individual will be better able
to follow the action back to its origin, and will thus be able to value longer stretches of its
history than when it stands before the actions of other individuals.

An action, then, will be good when it preserves and develops the welfare of conscious beings, evil
when it tends to the opposite. Evil will here, as at the standpoint of individualism, be what
dissolves and isolates. When a single individual either makes itself or is made by others an
absolute end, the community of conscious beings is dissolved. Evil is therefore egoism in its various
degrees and under its various forms. And the judgment on it becomes the sterner the more conscious it
is; for the deeper and firmer the root the action has in the whole disposition, and the harder it is
to remedy the isolation and dissolution. The stern \textit{valuation} itself finds its motivation and
its limitation in the wish for such a restoration. It is therefore \textit{more than an aesthetic
judgment; it has a definite end,} --- \textit{the same end as that by which the action is valued.}

11. The objective principle, the principle for determining the content of ethics and for valuing
human actions, thus becomes here the principle of \textit{the general welfare.} According to it no
action, and no institution or ordinance of life founded by action, has value except insofar as it
serves to further the life and happiness of conscious beings. There are many other things besides
human actions that work to further or to harm these. Through the very course of unconscious natural
life things work for or against the welfare of conscious beings; but the valuation of such effects
has no ethical character, because they do not spring from any consciousness and so cannot be
influenced by any judgment. The judgment on what happens outside the domain of human actions
acquires an aesthetic or religious character, that is, it has the character of an outburst of mood
before what must essentially be taken as it is. The ethical judgment, on the other hand, can itself
be motivated by the principle according to which it is passed, and must serve the definite
% --- p. 44 ---
\opage{44}end of bringing about greater welfare. Therefore it can concern only such events as can be
motivated by a judgment, that is, actions. This comes out most clearly where the valuing individual
is the same as the acting one, or where the ethical judgment is acknowledged by the acting
individual; in other cases it will be a special task how the acting individual can be brought to
acknowledge the judgment, a task of a psychological-pedagogical nature. But it follows also that
ethical valuation can itself be valued --- and that according to its own principle. If ethical
judgments can be motivated only by their practical influence on the acting beings they concern, it
is clear that in the end they must be regarded as pedagogical means. Of the two concepts, valuation
and education, education is the higher, the more comprehensive. A valuation that could not itself
serve as a means in the service of the end that forms the foundation of the valuation itself would
be a palpable inconsistency. In this way, as the following inquiries too will often show in
particulars, ethics by its own principle points beyond itself.

I use the word \textit{welfare} rather than words such as happiness or utility, because these can
easily give rise, and have given rise, to misunderstandings. By the word welfare I mean everything
that leads to the satisfaction of the needs of human nature in its whole range. Ethics is to take
account of all the stages of life and cannot therefore start from a distinction between external and
inner, lower and higher welfare. Such a distinction is, after all, already a valuation, and can thus
first take place when the principle of valuation is given. The misunderstandings of which the
principle of welfare (or, as its most important champion, \textsc{Bentham}, calls it: the principle
of the greatest possible happiness of the greatest possible number of
men)\footnote[1]{Bentham is not the author of this formula. He has himself declared that he has it
either from \emph{Priestley} or from \emph{Beccaria}. But both these writers are in turn influenced by
\emph{Hutcheson}, in whom the formula frequently occurs. (See Den nyere Filosofis
Historie\textsuperscript{2} I. p. 400). Hutcheson again is influenced by his zealous study of the later
Stoics, in whom the real origin of the famous principle must be sought. This has been shown by W. R.
Scott: Francis Hutcheson. Cambridge 1900. p. 272---279. Cf. also above p. 38 Note.} has been the
object derive in great part from not having kept this in view. --- There is, however, also another
point that is more clearly designated by the word
% --- p. 45 ---
\opage{45}welfare than by other expressions one might use. For it points to a total state. Momentary feelings
of pleasure and pain are no sure criterion of life as a whole. According to a psychological
hypothesis pleasure and pain are, to be sure, as a rule the expression of the advance or decline of
life: pain is a sign of a beginning dissolution of life, pleasure of its normal and harmonious
unfolding. (Cf. my \textit{Psykologi} VI D, 2. 3). But from single isolated feelings of pleasure and
pain nothing certain can be inferred; nor would a mere summing up of them lead to the goal. Rather, a
connection must be sought between the feelings and the whole character, the real unity in
consciousness to which they belong, as well as with the whole state of the organism. Likewise the
feeling of pleasure or pain of single individuals must be regarded in connection with the whole
social state. So-called utilitarianism, the ethical school founded chiefly by Bentham, which has the
great merit of having asserted the principle of welfare with energy, has harmed its cause by taking
as its basis a psychological theory that dissolves consciousness into a multitude of sensations and
feelings and society into a multitude of individuals. The significance that the feeling of pleasure
and pain has for lasting and comprehensive welfare cannot be exhausted by a simple sum. --- It follows
of itself that the burden of proof lies on him who demands the giving up of a feeling of pleasure in
the single moment or in the single individual: the interests must be shown for whose benefit the
sacrifice is to take place (cf. § 4).

12. The principle of welfare by no means gives us a magic key that can open everything without more
ado. It is a misunderstanding of what a principle is to think that it settles all the particular
cases for us at a stroke. We look for principles precisely because the particular cases can be so
intricate and
% --- p. 46 ---
\opage{46}complex that we gain some sort of survey only by starting from certain definite presuppositions.
And, as already indicated, a principle essentially has the significance of \textit{distributing the
burden of proof.} The principle of welfare demands a proof from him who would hinder life in its
involuntary unfolding and who would inflict pain, while no special proof of justification is needed
in cases where one furthers the development of life and procures increased happiness. And although
philosophical ethics is a practical science, it nevertheless satisfies a theoretical interest as
well: namely that of grounding our valuation of actions. Such a grounding becomes impossible unless
--- through many intermediate links perhaps --- we go back from the intricate historical context in
which each single action presents itself to a general principle that in and for itself contains
nothing but the spirit and direction in which ethical valuation is to proceed. The principle of
welfare stands in ethics as the principle of causality stands in the theory of knowledge. Neither
theoretically nor practically do the general principles give us complete solutions of the special
problems. The doctrine of principles is, to be sure, placed first in the synthetic or systematic
exposition of a science, but that does not mean that they are found first; on the contrary, they are
found in the first instance by the analysis of certain prominent phenomena and are then set up
hypothetically as presuppositions of the special explanations.

The reasoning of philosophical ethics must not be confused with the way in which, in practical
deliberation, we discuss whether an action is to be undertaken. In practical deliberation we are
guided by instincts and impulses, by motives that for a great part do not come fully forth in our
consciousness, by thoughts and feelings whose first origin we do not discuss. We are guided by the
``positive morality'' into which we have grown, and which is for a great part an inheritance of the
race. The ethical art precedes the ethical science; the latter seeks partly to point out, partly to
correct the principles by which the former is unconsciously guided. Theoretical reflection will of
course often act back on practical deliberation, correcting and purifying it; therefore complete
freedom of discussion of ethical questions
% --- p. 47 ---
\opage{47}is of such extraordinarily great importance. But only in cases of doubt does conscious reflection
come to play a direct part in ethical decisions\footnote[1]{Cf. my \textit{Etiske Undersøgelser.}
Copenhagen 1891. p. 26---29.}.

13. If regard for the greatest and most comprehensive welfare is to be the principle of all ethical
grounding, it cannot of course itself be ethically grounded. If it really is the right ethical
principle, that is: if it really designates the decisive standard for the valuation we make from our
standpoint, --- then the ethical way of looking at things first comes into force with its ---
conscious or unconscious --- acknowledgment, and no ethical discussion can be carried on with him who
denies it. Only by an \textit{indirect} way can there be question of a grounding of the principle
itself. For insofar as he who denies the validity of our principle himself sets up other principles,
one may seek to show that these presuppose it and are, logically regarded, derived in relation to
it. This would be the only fruitful way in which the debate could be carried on between the schools
of Kant and of Bentham. This way \textsc{Henry Sidgwick} has taken with great skill in his work
\textit{The Methods of Ethics.} He goes even further there and seeks to show how the moral tradition
prevailing in modern Europe finds its full explanation only when one thinks of it as having arisen
under the influence of an unconscious utilitarianism. Scattered commands and decisions can be brought
into connection with one another, and the doubt and indeterminacy that may prevail as to the decision
of particular questions can be overcome, when one takes the principle of welfare as the basis.
Likewise we thereby gain an explanation of why certain virtues have played a particularly prominent
part at certain times and among certain peoples.

But against him who resolutely takes his stand at the standpoint of individualism (let alone at that
of the sovereignty of the moment) this, as was said, is of no help. These standpoints are logically
unassailable so long as they are consistent. The same holds against those who extend the principle of
valuation only to a narrower social group, making the family, the caste, the nation,
% --- p. 48 ---
\opage{48}the race or the sect all in all. The narrower community can thereby come to stand as an
individualist or egoist in relation to the more comprehensive one. In this way we can think of as
many ethical systems as there are larger or smaller wholes. The ethical world extends from the moment
to mankind, but between these extremes there are many points where one can stop, and beyond which
only the psychological-historical development leads, through the changes it brings about in the
human life of feeling. A great example of this is the way in which, as a consequence of the
development of culture after the conquests of Alexander the Great, there arose a common feeling of
humanity that became the foundation of a new ethics. It becomes the business of psychology and of
history to show how the different ethical standpoints and motives of valuation come into being. One
standpoint is replaced by another through a course of development that is subject to psychological
and historical laws. There is a continuous connection between the standpoints or motives of
valuation, insofar as they all spring from human nature and come into being in the course of human
history\footnote[1]{Cf. \textit{Etiske Undersøgelser.} p. 7---14.}. Historical or comparative ethics
(cf. I, 4), which relies precisely on history and psychology, takes hold where philosophical ethics,
which seeks the grounding of ethical judgments, stops. The limit of the scientific character of
philosophical ethics is thus not the limit of all scientific treatment of ethical judgments. In
history the different ethical systems contend and seek by way of education and agitation to displace,
or rather to absorb, one another. The way in which this education and agitation is practiced must of
course consistently be determined by the individual system's own principles: one educates others not
according to their points of view but according to one's own. What matters is to bring about a
psychological foundation hitherto lacking, to awaken a certain definite motive of valuation. One
wishes to raise the level of the others, not to lower one's own. But before the education is
completed the others cannot acknowledge the principles according to which they have been educated.
---

% --- p. 49 ---
\opage{49}In short: \textit{every standard for the valuation of actions rests on definite
psychological-historical presuppositions.} He who is to acknowledge and seriously apply the
principle of the greatest possible welfare of the greatest possible number of conscious beings must
not be an egoist or individualist, a fanatical patriot or sectarian, but must be able to follow human
actions with disinterested and universal sympathy. This is \textit{the subjective presupposition of
the objective principle.} Apart from it, applying the principle in judgment and drawing consequences
from it is merely intellectual curiosity.

It was \textsc{Bentham}'s chief fault that he did not see clearly that a subjective principle formed
the presupposition of the objective principle. He demanded an objective principle to control and
regulate the various subjective views and assertions in the ethical field, and such a principle he
found in the principle of the greatest possible happiness of as many as possible. But \textit{the
subjectivity (the conscience) that is to be regulated by the objective principle always} itself
\textit{lies at the foundation of this principle's} acknowledgment. An ethics that does not take this
into account always acquires a dogmatic character, however excellently its principle may be suited
to serve as a guide in the special discussion. \textit{Philosophical ethics must expressly establish
on what standpoint the valuing subjectivity stands.} The road to this standpoint may, historically
regarded, be long and winding\footnote[1]{Cf. \textit{Om Grundlaget for den humane Etik} (1876).
Chap. III: Authority. It is there sought to show how the relation of authority is the most important
source of the development of moral feeling in history. I would now, however, stress somewhat more the
significance of fellow-feeling and the feeling of community in this respect.}. ---

A few words on how the whole conception of the significance and limitation of ethics developed in
this chapter came into being will be in place here. --- I started from the indubitable fact that in
life ethical judgments are passed, judgments that approve or disapprove of human actions or ordinances
of life, --- judgments of good and evil. The first presupposition we bring to such judgments is that
they must not
% --- p. 50 ---
\opage{50}conflict with one another. The standard applied in approval or disapproval must always be the same;
the valuation must be carried through consistently from the given standpoint. It is a standard we
apply to all ethical systems: there must be a necessary connection between the end taken as the basis
and the means and ways employed. When Socrates laid such great weight on the command ``Know
thyself!'' he set up a principle whose validity and applicability are independent of the real
difference between ethical systems. It must be possible to agree whether, where an end of definite
content and range has been set, it is consistent to pronounce a certain approval or disapproval of a
given action. But formal consistency does not settle everything. What judgments are passed depends
not only on the consistency of thinking, but also on what end is presupposed. When the end (and with
it the foundation, the motive of valuation) is changed, the whole system of ethical judgments is
changed too. Two individuals who take wholly different highest ends as the basis of their valuation
can therefore continue the discussion with one another only up to a certain point. The different end
brings with it a different valuation. And men are brought to set themselves new ends through
practical, personal influence (education and agitation), not through mere theoretical inferences. Man
grows with his greater ends, says \textsc{Schiller}. But man must already have grown in order to be
able to set himself greater ends, and on such growth, therefore, it also depends what ethical
propositions can come to stand as grounded\footnote[1]{This whole line of thought I first developed
in an essay on \textit{Den filosofiske Etiks Principer} (Det kgl. danske Videnskabernes Selskabs
Oversigter 1886), which was taken into the first edition of the present work (1887), where it made up
the greater part of chap. III. Later I treated the same subject in the first chapter of
\textit{Etiske Undersøgelser} (1891) and in my Zurich lectures (\textit{Etische Principienlehre.}
Bern 1896), and in \textit{Den menneskelige Tanke, dens Former og Opgaver} (1910) I came back to it
once more (p. 346---359). The guiding thought is: how far is grounding possible in this domain? ---
In \textit{Grundlaget for den humane Etik} (1876) I did not yet see this problem sharply.}.

% --- p. 51 ---
\opage{51}Good and evil are concepts of relation, --- as all concepts are\footnote[1]{Cf. \textit{Psykologi}
V D, 5.}. They express a relation to a conscious or
unconscious end and through it to a will. It sounds idealistic, but it is at bottom meaningless, to
say that there is anything good or evil in and for itself --- i.e.\ independently of every end. If
one only goes back far enough, one will find an end and a will behind every ethical judgment.

(1913). Quite briefly, the sting of the ethical problem can be expressed by saying that value can be
measured only by value. In the end a fundamental value must be presupposed that determines the
validity of all other values and their place relative to one another. Every valuation (or
``revaluation'') presupposes, consciously or unconsciously, such a fundamental value. A closer
examination will show that every such fundamental value is conditioned by regard for an individual
or social whole and the conditions of its existence and development. The rational character of ethics
depends on the possibility of finding and formulating these conditions. But ever anew the opposition
and conflict between different individual and social wholes will pose the problem. The concepts of
value are the most special, the richest in content and the most conditioned of all our fundamental
concepts (categories)\footnote[2]{See on this \textit{Den menneskelige Tanke.} p. 237---245;
346---349.}.

14. Ethical judgments, then, always point back --- however much one may ground them by an objective
principle of valuation --- to a \textit{subjective foundation.} They are in the end the expression of
a feeling. --- But there is no disputing, it is said, about feelings. This is true insofar as a
feeling is a psychological fact and must be taken as such. But into every feeling there enter certain
definite ideas on which its character and direction depend, and the coherence and validity of these
ideas can be discussed. Such a discussion will, though slowly, act back on the foundation.
Thus it will be possible to show how little right the individualist has
% --- p. 52 ---
\opage{52}to regard himself as an isolated and unique being. It will be possible to point out the illogicality
of restricting one's love of mankind to comprise only human beings of a certain color, a certain
descent or a certain faith. The complete shifting of the center of gravity, however, will first be
brought about when personal and historical experiences cooperate.

Such is the position of every practical science. The significance of its propositions depends in
the end on an interest of feeling that dictates the end for which the means are sought.
\textit{Political economy} presupposes the acquisitive impulse as given, and its doctrines stand and
fall with the working of this impulse. It examines the ways and means by which the end set by this
impulse can be reached. And insofar as other needs stir in men that come to influence their conduct,
the propositions of political economy acquire an abstract and hypothetical character, hold only under
a presupposition that is not completely confirmed in reality. Indeed, one must go a step further.
Political economy is not a science wholly independent of ethics. It is not only a doctrine of
production but also a doctrine of distribution. It does not regard the economic domain as wholly
isolated, but sees it in its connection with general culture, and in the end with the ethical ends of
human social life. The task of political economy becomes in the end that of ordering and guiding
production and distribution so that a human existence becomes attainable for as many as possible in
society. Political economy is thus an ethical science, and rests in the end on the same foundation as
humane ethics. The same holds of \textit{jurisprudence.} When one distinguishes between morality and
law, the essential difference is that the content of law can be carried through and realized by
external force, while what is properly moral demands inner acknowledgment and the assent of the will.
But only he bows to force who either acknowledges the use of force as justified or perceives his own
powerlessness; and on the other hand force is used to enforce the law only where there is ability and
will for it in the persons to whom it belongs. The propositions of jurisprudence therefore also
acquire an abstract and hypothetical character.
% --- p. 53 ---
\opage{53}In general jurisprudence it is even disregarded whether the means of enforcing the law are in fact
present. For in the widest sense legal relations are understood to be such relations of life as are
\textit{suited} to legal ordering. One builds, then, on the idea of a society where everything suited
to it is legally ordered, just as political economy in many of its propositions builds on the idea of
a society where the acquisitive impulse is at work, and humane ethics on the idea of a society where
love of mankind (disinterested and universal sympathy) rules and the general welfare is striven for.
And in the last instance the definite difference from ethics falls away for jurisprudence too, as for
political economy, for the grounding of the use of force in jurisprudence is in the end of an ethical
nature, since it is sought in the necessity of forming a fixed, external order of life with definite
limits to the freedom of action, so that the ethical life and the higher human endeavors may unfold
under its shelter. Jurisprudence too, then, is an ethical science and rests on the same foundation as
humane ethics. --- In any case, by discussing the ultimate foundation of political economy and
jurisprudence one comes back to problems analogous to those to which the question of the grounding of
ethical judgments leads.

Both political economy and jurisprudence have, like ethics, often suffered from the misunderstanding
that they were to be pure sciences of reason. They have then acquired a dogmatic and unhistorical
character that conflicts with their nature as practical sciences. They presuppose a subjectivity with
certain definite interests, and the possibility is not excluded that these interests may change.

15. If in ethics one wished to start only from the subjective principle, the foundation, ethics would
become merely a doctrine of the ethical feeling. Since every ethics must nevertheless teach something
about \textit{what} ought to be done, one would be obliged to derive the content from the foundation.
But if this is not done in such a way that an objective principle is first sought, ethics becomes
merely a series of postulates which it becomes impossible to ground even for those who stand on the
same foundation.

% --- p. 54 ---
\opage{54}I have sought to show how the subjective and the objective principle develop together and answer to
one another. A disinterested and universal sympathy cannot --- without giving itself up --- take any
other principle as the basis of the valuation of action than the principle of the general welfare.
From shortsightedness and impatience it may come to act against this principle; but it then also acts
against its own fundamental value, its end, and the more it becomes conscious of this, the more it
will also follow the principle of welfare in particulars. Conscience is not infallible: therefore it
needs to be guided by an objective principle. At any definite moment the valuation must of course be
made according to the insight at hand; there is no other tribunal than conscience as well informed as
possible. But for all that the decision may nevertheless be objectively wrong; a more consistent
application of the principle of welfare, built on more extended experiences, may be able to show
this. Yet perhaps the greater consistency and the more extended experience themselves became possible
only because the decision in the given case was made according to the insight one had. Conscience is
the highest authority, but an authority that can constantly make itself more perfect. The objective
principle not only makes negotiation between different consciences (on the same foundation)
possible, but is also the condition of the single individual's conscience being able to judge itself.

In this way, then, the content acts back on the foundation, the objective principle on the
subjective. Subjectivity controls itself by the principle contained in the end it sets itself. In a
corresponding way it controls itself in the theoretical field by the principle of identity and the
principle of causality, although these principles too are in the last instance posited by a
subjective activity\footnote[1]{Cf. my \textit{Psykologi} VB, 11; D, 3.}.

16. When one considers more closely \textit{the relation between foundation and content,} it may
easily seem as if the inference from foundation to content has greater certainty than that from
content to foundation. History shows us that one and the same ethical content,
% --- p. 55 ---
\opage{55}the same maxims, commands and laws, have been set up on very different foundations. The mere
circumstance that certain commands of duty recur within very different religions is enough to show
this. Ethical systems on different foundations may very well present common elements, though it will
appear that these elements occupy a somewhat different place in the different systems. The egoist's
system, for example, does not exclude regard for family and state, though these stand there only as
means. National and humane ethics may lay extraordinarily great weight on the single individual's
self-assertion, though this is not, as in the egoist's system, the absolute end. --- If we keep to the
principle of welfare as the most significant example, could it not be motivated in another way than
has been done here? --- To this the answer is that historically regarded it does indeed seem so, but
that inconsistencies and imperfections can always be pointed out when content and foundation do not
answer to one another completely and reciprocally.

There are especially two ways, besides the one used in this inquiry, in which a foundation has been
sought for the principle of welfare.

It has been held that the endeavor to further the general welfare is \textit{the road by which the
single individual best furthers his own welfare,} since for his own person he thereby succeeds in
carrying through ends that he could otherwise not reach, at any rate not so easily and so surely. The
endeavor for the general welfare thus becomes merely a detour that must be taken so that the single
person may reach his own satisfaction. Every ethical command indicates such a detour, and ethics
becomes the systematic doctrine of such detours. This was the conception of the foundation of ethics
to which the founder of modern utilitarianism, \textsc{Bentham}, was most inclined, and which he
expressed in a particularly crude way in a work published after his death
(Deontology)\footnote[1]{In \textit{Principles of Morals and Legislation,} which is Bentham's chief
work, this doctrine does not come out so definitely. He does not there express himself clearly on the
foundation of ethics. In one single passage it is said: ``The dictates of utility are neither more
nor less than the dictates of the most extensive and enlightened benevolence'' (X, 36). But no use
whatever is made of this proposition in setting up the ethical principle of valuation itself (ch.
I---II). --- Cf. my characterization of Bentham in Den nyere Filosofis Historie\textsuperscript{2} II.
p. 363---368.}. In my essay ``Om Grundlaget for
% --- p. 56 ---
\opage{56}den humane Etik'' I criticized Bentham from this side, and only from this side. Here I will add
that it is a great and very precarious hypothesis that Bentham uses as the foundation of his ethics.
It is nothing less than the bold presupposition that there is a harmony between the egoistic interests
of all individuals, provided only that each of them is clear about what these really are. He who
works for his own well-understood interest thus works, on this assumption, for those of all others
too. Therefore no motives other than egoistic ones are really needed, and the very setting up of the
principle of welfare can also be motivated egoistically. Every valuing subjectivity judges actions
according to the way they further the general welfare, because the valuing individual finds itself
personally best served by such actions.

It will easily be seen that this is a hypothesis that stands on very weak feet and constantly
challenges doubt. It leads into very intricate inquiries that can never be concluded. It is therefore
a very questionable foundation on which to erect ethics. And besides, even if one assumes this
foundation, the content built on it will be only an external means, and will not immediately in and
for itself answer to what lies in the foundation. The content of ethics becomes a system of
concessions that the individual, who really stands at the standpoint of absolute individualism, must
allow to be extorted from him. For the straight road is after all always the best, and a detour only
a makeshift.

Another foundation on which the principle of welfare has appeared is \textit{the theological.} Here too
one relies on far-reaching and disputable presuppositions. Ethical valuation is thereby made dependent
on the solution of the religious problems. But the endeavors to set up a philosophical ethics have
proceeded precisely from the wish to make the ethical problem
% --- p. 57 ---
\opage{57}independent of the religious problems. If ethics is built on a theological foundation, these
endeavors must be said to have failed. Besides, theologically grounded ethics is struck by an
objection similar to that against the ethics built on the presupposed harmony between well-understood
egoistic interests. For the whole of ethics becomes the means to the satisfaction of an interest that
is in itself different from --- at any rate does not wholly coincide with --- the ethical interest
itself. The individual's relation to what is assumed to be the will of the Godhead is a relation by
itself, and it has first to be proved that there is agreement between the demands of the Godhead and
those of the principle of welfare. In reality, too, every theological ethics, even if it places the
principle of welfare very high, has a class of special duties that are due to the theological
foundation but are not grounded by the principle of welfare, and may even come into conflict with
it\footnote[1]{See \textit{Om Grundlaget for den humane Etik.} p. 68---72.}.

In setting up the principles of a science the law of parsimony should prevail. In the ethical feeling
conditioned by disinterested and universal sympathy we have a foundation whose reality is
indubitable, though there may be dispute as to how great a part it plays in the world. Love of
mankind has, as was indicated in the foregoing, slowly worked its way forward and is still a
struggling power, --- must struggle for dominion both in the outer world and within the single person
himself. Whether it will become the motive of valuation that will one day prevail in all, we do not
know, and we need not know\footnote[2]{Axel \emph{Hägerström} (Undersökning \textit{af den empiristiska
Etikens} möjlighet med särskild hänsyn till dess moderna hufvudformer. Upsala 1895. p. 20 f.) holds
that I must assume universal sympathy to be equally great in all men, since otherwise ethics gets no
universally valid foundation and thereby loses the character of a science. My opinion, however, is
that universal sympathy develops both sporadically and successively, and that the same holds of any
other motive of valuation one may name, --- and precisely therein I see a limitation of the scientific
character of ethics (not merely of my ethics, but of every ethics). --- A theological critic (N.
Tejsen in ``Vidar'' 1888) asks what philosophical ethics will do if the feeling it presupposes as the
motive of valuation is lacking. Here one may answer with a counter-question: what does theological
ethics do with him who believes neither in God nor in the Devil?}. It
% --- p. 58 ---
\opage{58}already plays such a part in the world that it is of very great interest to examine what ethical
judgments consistently spring from it. And it will prove to be an extraordinarily expedient choice to
take precisely this motive of valuation as the basis. For there is hardly any other motive that in
such a degree as this gives birth to an interest in following the effects of actions in their whole
course, as far as is at all possible. The individualist stops his inquiry when he has made sure that
the effects of an action do not touch \textit{his} welfare; he who stands on the ground of national
ethics is unconcerned so long as the effects of an action do not touch \textit{his} people. Not even
the hypothesis of a harmony between the different individual interests (of single persons or of
nations) becomes of decisive importance for these standpoints. What does it matter to the
individualist whether others' interests agree with his, provided only that he has the means to secure
his own interest? At most he may feel a certain security at the thought of that harmony. Universal
sympathy as the motive of valuation, on the contrary, leads to tracing the way in which the effects
of actions intervene in the life of all personal beings, and thereby gives birth to an interest in
inquiries of a scientific, psychological, historical, social and statistical kind, which can go
entirely together with the scientific interest in such inquiries. By virtue of the motive of valuation
chosen we are thus led beyond ethics, or ethics is led to rely on natural science, history,
psychology, statistics and political economy, so that valuation can become as \textit{objective} and as
\textit{realistic} as possible. And this transition takes place with inner consistency and in a more
direct way than at other standpoints that attempt to take up and ground the principle of welfare. To
universal sympathy as motive of valuation answers the principle of welfare as an objective principle
for the valuation of human actions; it merely unfolds more fully the logical consequences of the
% --- p. 59 ---
\opage{59}psychological foundation and makes possible a consistent grounding of its single expressions. The
foundation and the content of ethics have here moved as close to one another as is at all possible.

Every ethical conception that starts from an authority as its foundation --- whether this be
supernatural or natural, consist in the will of those in power, in the attraction of models or in the
compulsion of public opinion --- must yet grant that as soon as the individual can immediately make
the ends of the authority his own ends, the only sure and simple foundation of the ethical will thereby
have been reached, and the possibility of a constant correction of the ethical at the same time most
surely given. It is a detour when I enter into relation with the highest ends that my will
acknowledges only through the will of an authority. In energetic and intellectually developed love of
mankind we have the only foundation where the individual and the universal ends can immediately be
identified. Over against a force such as this all authority can have only provisional and pedagogical
significance. All authority may vanish without the abiding foundation of the ethical vanishing with
it. It is one thing that the ethical law (when we disregard the standpoint of the moment and the
standpoint of absolute individualism) is the conditions of life of a comprehensive whole formulated
in definite thoughts, another how these thoughts are acknowledged by the single individual. Even where
those conditions of life are enjoined by external laws and by public opinion, the individual may, after
all, have come to them by way of his own conviction. Perhaps the single person even has a far clearer
and more far-reaching insight into the conditions of life of society than the laws and public opinion
have so far been able to express, --- and should then the disposition with which he follows this
insight in his conduct not deserve the name of ethical, because no authority sanctions it? The motive
of valuation that it will be of greatest interest to test must surely be the one that is most directly
able to value every authority by the standard that all authority must satisfy in order to be regarded
as justified. If one builds all ethics on authority, where then will one find a standpoint from which
the ethical valuation of authorities can take place?

17. The difference between foundation and content does not
% --- p. 60 ---
\opage{60}coincide with the difference between \textit{individual} and \textit{social} ethics. To the content
of ethics belong both the individual qualities of character and the social forms of life that agree
with the principle of valuation; it thus comprises both individual and social ethics. The question
whether individual and social ethics are independent of one another, or whether one of them is
determinative for the other, must be decided according to the principle of welfare. Within special
ethics an individualism may assert itself of another kind than the one spoken of earlier, since it
does not build on the sovereignty of the individual but is grounded by regard for the general welfare,
which requires as many independent and peculiar points of departure for action as possible. It goes
in a corresponding way with the grounding of the existence and significance of the smaller communities
within the larger. The general principle of welfare demands regard for all conscious beings whose
feelings of pleasure and pain can be set in motion by the actions. But precisely in accordance with
this regard it may become of the greatest importance that interest concentrates itself on a narrower
circle. Through life in the family and in the definite activity of a calling one perhaps works in the
best way for the advance of the nation, and through taking part in the national life for that of
mankind. All these questions are of a more special kind than the one it is the task of this chapter
to discuss.

18. As we saw in § 2, the opposition between ``subjective'' and ``objective'' ethics arose in the
first instance from the different weight that was laid on the inner and on the outer side of the
action. I have sought to show that this opposition goes further back and is closely connected with
the character of ethics as a practical science. I now return in conclusion to those two sides of the
action in order to have occasion to determine the relation between what I have here called the
ethical \textit{foundation} or \textit{the motive of valuation} and the \textit{motives} of the single
actions.

The history of ethics shows us that valuation first concerns \textit{the external action and its
effects,} but that it gradually widens to apply also to \textit{the motives, the disposition, the
character of the agent.} It is, after all, natural that one keeps first to what seems to lie plain
to view and can become
% --- p. 61 ---
\opage{61}the object of sensuous observation. The spiritual eye is directed to the external before it is able
to apprehend the internal. At the more primitive stages, moreover, actions have essentially the
character of reflex movements and expressions of instinct; the motives do not particularly come
forward, and interest does not particularly dwell on them. Nor is any distinction made between
morality and law, between inner and outer agreement with the prevailing law or statute. Ethical law,
custom and juridical law have not yet separated out from one another. A significant side of ethical
development consists precisely in these differences making themselves felt. This presupposes a
capacity to understand the dependence of the action on definite motives and their tendency to work in
a certain definite direction. Ethical valuation is then extended to the internal too. The great
revolutions in the ethical field appear essentially as advances in the \textit{inward deepening} of
ethical valuation. Such an advance we have where ethics separates itself from jurisprudence by
stressing the inner source of the action as the essential thing. What matters is not merely to perform
or to refrain from the external action, but agreement between the whole inner state and the demand of
the ethical law. Such an inward deepening is at the same time a \textit{generalization.} For when a
motive (e.g.\ egoism, hatred) is rejected, all the single actions that spring from it are rejected
with it, and these need not be specially enumerated. Likewise the approval of a motive is the approval
of all actions that spring from it. The transition from external to internal valuation is therefore
finally a great \textit{simplification} of the ethical law. It is not the multiplicity of commands,
but the qualities of character that are to rule, and the direction in which the forms of society are
to be developed, that have to be laid down.

Examples of such inward deepening, generalization and simplification we have in Christianity's
break with Judaism and in Protestantism's break with Catholicism, and in general wherever a purely
ethical way of looking at things separates itself from a view that relies on custom, authority and
external force.

The ethical law demands the presence of the disposition
% --- p. 62 ---
\opage{62}that carries it in the race. This is what Kant expressed in the proposition that it is a duty to
have a conscience. Since the acknowledgment of duties presupposes conscience, it might seem as if one
were here going round in a circle. But this is only apparently so; for the foundation and the motive of
action \textit{need} not coincide, and it need not be an imperfection either when they do not coincide.
According to the principle of welfare it may very well be necessary that motives other than the
feeling of duty itself make themselves felt. It may, for example, be necessary or healthy that the
instinct of self-preservation, or again immediate fellow-feeling, lead man to undertake actions for
his own and others' welfare, and that conscience does not stir at every single action. It may even be
a sign of perfection when actions that demand effort and sacrifice are undertaken without its
happening from a feeling of duty. Spiritual practice leads, after all, to this: that what at first
could be performed only through several psychological intermediate links, through express reflection
and exertion of will, can at last be performed directly and without special consciousness of why one
really acts so.

19. All ethics is practical idealism. It presupposes that we set ourselves ends, but an end is not
something that is, but something that \textit{is to} be. All ethics therefore presupposes an impulse, a
living feeling and craving bound up with an idea of what is craved. If reality satisfied every need,
there would be no ideal and therefore no ethics. But an ideal that ethics is to be able to use must not
merely stand above reality; it must also have points of attachment in the given reality, so that from
it at least an approximation to the ideal becomes possible. It must --- subjectively regarded --- be
able to find acknowledgment, and it must --- objectively regarded --- be capable of being carried
through more and more in the world of experience. It is real men who are to will and act, and a real
world that is the scene of their willing and acting. What ethics demands must be physically,
psychologically and historically possible, must not conflict with the laws of the real world. He who
rejects all ethics because it sets up ideals
% --- p. 63 ---
\opage{63}can only mean by this that its ideals lack points of attachment in reality. He will not be able to
avoid setting up ideals himself, so that he overthrows ethics by a new ethics. When, for example, he
demands that one be honest enough to admit that ethics is an impossibility, or when he perhaps even
rejects ethics because all the misfortunes of human life derive from ``morality,'' he speaks in the
name of honesty or of love of mankind, and so appears as an ethical idealist. The discussion can turn
only on which ideals one is to set up. No one escapes ideals, however he twists and turns. But ideals
can be of very different worth, and it is about the worth that there is dispute.

Theoretical science, too, always has an idealistic character. It operates with principles and
presuppositions so simple that real experience never completely covers them. All coherent,
constructive thinking becomes possible only by our holding fast a single side, a single element of
the given, and drawing all the consequences from it. Thus geometry regards things merely as extended
in space and derives the general laws of extension without taking account of the other properties of
things. In all fields, in order to be able to think sharply, we must think of things as simpler than
they are in reality. This does not exclude that the results of thinking can hold of things, provided
only that it is essential and significant sides of them that we treat in our constructive thinking.
Just as geometry constructs an ideal space, ethics constructs on the ground of psychology and history
an ideal conscience. In reality every man's conscience is composed of different, sometimes even
conflicting, elements. Imitation and tradition, egoism and ambition, regard for others' opinion and
fear of punishment, reverence and love of mankind, religiosity and social feeling --- these and still
more motives stir within the ordinary conscience, which is not so simple a feeling as is often
supposed. The different possible ethical systems characterized in the foregoing (4---13) contend not
only in history and in the race but also \textit{within} the single
% --- p. 64 ---
\opage{64}human being's consciousness. Everyone has --- or can have --- several interests and ends that contend
with one another as to which of them shall in the end determine the valuation a man makes of his own
and others' actions. Anyone who would attempt to develop a coherent ethics must take a single motive
of valuation as the basis, as ruling and determining. And in taking a single motive as the basis an
idealization has already been undertaken. The history of philosophical ethics shows us how men have
tried to carry through the different possible motives of valuation. Thus \textsc{Hobbes} and
\textsc{Bentham} started from egoism or well-understood self-interest, \textsc{Hutcheson},
\textsc{Hume} and \textsc{Adam Smith} from men's sympathy grounded in common interests, \textsc{Hegel}
from the idea of the state and of the cultural community. Philosophical ethics has here the
significance of bringing the different ends men can set themselves to full consciousness and of
drawing all the consequences from them. The different ideals are thereby so illuminated that the
struggle between them can be carried on with greater clarity; it is one of the most important
contributions theoretical ethics can make to practical ethical striving. The worth of the
contribution will depend on how significant the motive of valuation chosen for examination is. ---
From the foundation chosen, all motives and tendencies are valued according to how far they lead,
directly or indirectly, in the direction of the ideal that the foundation leads one to set up. The
difference between theoretical and practical science is here this: that in theory our ideals are
approximations to reality, which is the goal thought seeks to reach, while in practice it is reality
that is to be changed and brought nearer to the ideal.

20. In three respects ethical ideals depart from what is actually given.

First, there is much in actual willing and acting that directly \textit{conflicts with} the ethical
demand. Toward this ethics takes up a \textit{checking and negating} position, demands that it be
excluded from the ethical world. In the practical life of the will there corresponds to this the
checking activity by which involuntary promptings and inclinations, original or acquired, can be
pushed back\footnote[1]{\textit{Psykologi} II, 4 e; 6 e. --- IV, 4.}.

% --- p. 65 ---
\opage{65}Next, actual willing and acting often presents only \textit{weak and incomplete} realizations of
what the ethical demand contains. Here what is needed is increase both in degree and in range; what
is needed is potentiation and \textit{widening.} Ethics here pronounces an Excelsior! The direction is
good, but the movement lacks strength and fullness. In the practical life of the will there
corresponds to this the exertion of attention, the summoning of energy, in general the will's capacity
to act back on itself through its influence on ideas and feelings\footnote[1]{\textit{Psykologi} VII
A, 5; B, 2.}. --- Or the direction is good, but the \textit{motive} cannot be acknowledged, so that an
entirely different grounding is necessary if the will or the action is to become perfect. A direction
of the will that has arisen from one motive can later (through \textit{displacement,} see I, 4) perhaps
be acknowledged only by virtue of an entirely different motive\footnote[2]{Psykologi VI C, 2.}.

Finally, actual willing and acting may \textit{lack coherence,} unity and harmony. Different and
mutually conflicting tendencies and impulses may make themselves felt. What is needed then is
\textit{combination and concentration,} the working together and reconciling of the conflicting and
scattered forces. To this there corresponds in the field of the practical life of the will the
practice of composite movements and series of ideas, and the capacity to make a resolve and thereby
put an end to deliberation\footnote[3]{Psykologi VII B, 1 c: 3; 5 a}.

In the checking, in the heightening or motive-changing, and in the uniting activity alike, the
principle of welfare is the guide. Only where it is satisfied do these activities acquire value. One
can just as well check something good as something evil, just as well heighten something evil as
something good; the ethical value therefore lies not in these activities in and for themselves, but in
the direction in which they are exercised.

The processes spoken of find application not only in the ethical development of the single
individual but also in that of society and of the race. There may be individual expressions of life
that conflict with the welfare of the race, though they do not conflict with the individual's own
% --- p. 66 ---
\opage{66}isolated welfare; they must be checked. An endeavor for social welfare may lack strength or range
or right motivation; here, then, there is occasion for potentiation or widening or displacement.
Finally there may be lacking the coherence and harmony between the different expressions of life and
endeavors within society and the race that is necessary if the general welfare is to be reached; what
is needed then is to work together and organize the scattered activities. ---

From this side, too, it appears that philosophical ethics is both conservative and radical. (Cf. I,
4). It takes its points of departure in the given. On bare ground it can accomplish nothing. But this
given it then seeks to develop further in the three ways stated.

21. In more or less sharp opposition to the theory of ethical principles and of the possibility of
grounding them that is put forward in this chapter, several Danish writers have sought to find what is
in their opinion a more objective and rational foundation for ethics. Their endeavors aim at asserting
the scientific character of ethics; some of them even think they have shown a scientific foundation
for ethics for the first time. In particular their endeavors aim at freeing it from the empirical and
subjective limitation to which, on my view, it is necessarily subject. Now it is not probable that
every such limitation could fall away. For it appears everywhere, where one goes back to the
foundation in principle of our knowledge, that grounding and limitation go together. And it cannot,
after all, be otherwise, since every grounding rests on certain presuppositions. To overlook this
leads to dogmatism. It has been my endeavor to point clearly and definitely to the points where a
strict, rational ethics rests on psychological and historical presuppositions that do not hold
everywhere and at all times. I grant, however, that the works referred to may serve to supplement my
exposition by throwing light on sides of the problems that it has perhaps not allowed to come fully
into their own. And since they also bear witness to the interest in ethical problems that has stirred
in this country in recent
% --- p. 67 ---
\opage{67}years, I shall go somewhat more closely into these works as regards the decisive main point.

The difficulties that make themselves felt in every attempt at a scientific ethics have their ground
in the fact that all ethical concepts and propositions concern relations of value. Ethical judgments
state the relation between human actions and the ends they are to serve, and the concept of end in its
turn presupposes the concept of value, unless one takes the word ``end'' in a figurative sense. But
value again can belong only to what is thought of in relation to a living being's feeling. A subjective
element thereby enters. And besides, all living beings do not necessarily know the same values, and
therefore do not set themselves the same ends either. Here, then, value can stand against value, at any
rate when one goes back to the fundamental values on which the special values depend.

In the foregoing I have shown how I myself stand toward these difficulties. Every ethics must, from
whatever fundamental value it starts, seek means and ways to produce what is valuable in the world of
reality. It must take its stand on the ground of experience, both when it is a matter of showing in
what the valuable consists and when it is a matter of finding means and ways to produce it. In both
respects ethics is grounded, on my view, on biology, psychology and sociology. All human values are
bound up with life, the existence and unfolding of personal life and of social life, and the
conditions of life must be learned from life itself. This is the fundamental thought of my ethics,
and through it it becomes possible for ethics, in spite of its subjective point of departure, to
become an objective science (see above § 11 and 14---16). One will, I hope, not find in my exposition
a single attempt at grounding that does not rely on biology, psychology or sociology. Just as little as
the subjective character of the qualities of sense makes physics impossible as an objective science,
just as little does the subjectivity of values make ethics impossible as an objective science. In
both cases the task is to find the objective relations to which the subjective determinations answer.
But an essential difference between the position of ethics and that of physics arises from this,
that while the qualities
% --- p. 68 ---
\opage{68}of sense are tolerably the same for all, experiences of value vary among different individuals and
peoples and at different times. I am too much of an empiricist and historian to be able to overlook
this fact or to think I can escape it by abstract constructions. The task of scientific ethics then
becomes for me to examine what special values can consistently be derived from the different
fundamental values or types of value. (See § 3---13). The possibility of different ethical systems
thereby appears. These will not, however, fall wholly outside one another; they may meet in particular
propositions, which are then set up with different groundings. (Cf. § 13 and 19). Even if they contain
the same elements, these will yet be combined in different ways, and the timbre will therefore be
different. Or perhaps what stands as fundamental value in the one system will stand as derived value
in the other. Such an inquiry will in any case have its significance, even if it is regarded only as a
preliminary discussion (a ``Phänomenologie des Geistes''). There is a constant development in the
ethical field, both as regards single individuals and as regards the race. But we are unfortunately
in the midst of this development and cannot rise above the imperfections that follow from it. The
dispute still stands, and I cannot see otherwise than that the road I have taken must in any case
bring greater clarity about the object of the dispute. In the long run, clarity about the consequences
of the different standpoints will of course affect the course of development\footnote[1]{It has been
of very great interest to me to see that the French mathematician and philosopher Henri Poincaré took
up quite a similar position on the ethical problem to the one I have put forward in this chapter. See
his \textit{Dernières Pensées.} p. 236---238 (also p. 228 f.). (1913).}.

If one does not find the way in which I stand toward the difficulties named satisfactory, there will be
chiefly three roads one can take.

a. One can try to reduce the concept of value, which is, after all, what causes the difficulties, to
its most elementary forms. About the most elementary values there will be least dispute, and they will
most easily be brought into definite connection
% --- p. 69 ---
\opage{69}with objective relations. This road is taken in our literature by Dr. \textsc{Starcke}
\textit{(Samvittighedslivet.} 1894---97) and Dr. N. H. \textsc{Bang} \textit{(Begrebet Moral.} 1897).

b. One can go a step further and try to eliminate the whole concept of value, strike out all psychical
elements and presuppositions, and ground ethics on biology as a purely objective science. Ethical
propositions then become biological deductions. Such a ``biological'' ethics has been maintained by
Dr. S. \textsc{Hansen} \textit{(Etikens Begrundelse.} 1903).

c. Finally, one can seek to reduce the oppositions between the different fundamental values to purely
logical contradictions, which fall away when one seriously conceives man as a thinking being. This
road is taken by Professor \textsc{Kroman} \textit{(Begrebet ``Det Etiske''.} University programme, Nov.
1903).

It is interesting that all three roads have been tried by Danish writers. I shall now try to show that
none of these three roads leads to Rome.

a. Dr. \textsc{Starcke}'s conception of ethics has evidently developed under the influence of his
studies of primitive forms of society, studies that have borne fruit, among other things, in his
excellent critical work on the primitive family. For an inquirer who has immersed himself with great
energy in such studies, it is natural to apply relations from the more primitive societies to social
life as a whole. The clan, accordingly, stands in his account of ``the life of conscience'' as the
frame that encloses the single person's life, in that the respect and esteem of the others become the
most important condition of his secure conduct of life. ``The demands of morality\dots{} come to man
from without\dots{} The moral requirements are now of such a kind that the power of society seeks to
carry them through by force, so that he who transgresses them is punished; now they are referred to the
sanction that public opinion, the good repute and kindly disposition of other men can give them\dots{}
In the right of others to hold the individual responsible for his actions lies the center of gravity
of morality.'' (Samvittighedslivet. I. p. 120). This conception of the ethical consistently leads to a
limiting of the content of ethics. Everything not needed to bring about the
% --- p. 70 ---
\opage{70}feeling of security that arises in the single person from knowing himself and his conduct to be in
agreement with the punishing authority and public opinion falls outside ethics. When this limit is
reached, the single person can go his own ways. The whole of social life stands as merely a means for
the single person to be able to feel secure. And when this end is attained, the single person can seek
his happiness as he will. --- The question is whether individual and society can be set over against
one another in this external way. Even if the individual thinks he has settled accounts with society,
the way in which he seeks his happiness may yet be of the greatest social significance. It cannot be a
matter of indifference how single persons conduct their inner life, since it is in the inner, personal
life that the values come into being that can later change and develop society. The single person
therefore has, even in his most intimate relations, a social responsibility that only he himself can
make good. Here he is himself guardian and judge. It is placed in his hand whether at this central
spot in the world new values are to arise, or whether the germs of such are to perish. Dr. Starcke
does in fact teach that it is the conditions of social life themselves with which the single person
must feel himself in agreement in order to be able to feel secure. But if this is so, it is clear that
reward and punishment, esteem and contempt --- the only sanctions acknowledged in Dr. Starcke's
account --- can be only educative means. Why, now, should precisely the feeling of security be the
motive of valuation? And may it not be a matter of indifference to the single person who has
sufficient power whether the others feel secure, provided only that he does himself? The elementary
life of valuation that lies at the foundation of Dr. Starcke's account does not, then, lead with
necessity to the results he wishes to ground.

Dr. Bang holds that usage has already settled what morality is. ``A judgment that declares an action
moral or immoral states that the action, as a general practice, will further or hinder the good of
society, and that it is in the interest of the general good that society, through its organs, the state
or public opinion, further or hinder actions of that kind.'' (Begrebet Moral. p. 113). To this I
remark that,
% --- p. 71 ---
\opage{71}however the case may stand with usage, the grounding given in this statement holds only for him who
takes an interest in the general good; for anyone else it becomes merely a thought-experiment.
Experience shows that men can, with greater or less approximation, make their own good, not that of
society, their end; for such men a deduction from the good of society will have nothing really
convincing. They may grant that certain demands must consistently be made when one starts from the good
of society; but they do not themselves start from it unless they are compelled to, and compelled they
cannot be if they have courage and power enough to defy the state and public opinion. And according to
Dr. Bang the content of ethics does not extend beyond what can be maintained through the laws of the
state and the verdicts of public opinion. Outside the relatively narrow circle of duties that these
authorities rightly demand to be fulfilled, the single person can lead his own life and seek happiness
by his own ways.

Thus the application of the most elementary points of view necessarily leads to a dualism or a purely
external relation between individual and society. And in Dr. Bang this dualism appears in a
particularly striking way, in that he declares that ``the most admirable qualities of soul, the
qualities of the great man'' perhaps come forth precisely in him who breaks with ``morality.'' ``We have
never asserted that moral action is the most perfect action'' (p. 145). ``Morality is one perfection,
one virtue, among others; it is the fundamental social virtue, but not the only one, not the
all-comprehensive one, and not the highest either'' (p. 190). But where, then, is the standard by
virtue of which ``morality'' is compared with other virtues? There must surely be such a standard, for
otherwise it would be meaningless to speak of higher and lower qualities. Unfortunately this standard
is not stated --- though it must be precisely its nature that a treatment of the ethical problem ought
to concern. Why is the ``moral'' good, and when does it cease to be good, to give place to ``higher''
and ``highest'' qualities? That is the kernel of the problem, and it is this that I have tried to bring
forward.

% --- p. 72 ---
\opage{72}As regards the appeal to ordinary usage, Dr. Bang really takes it back himself, since he admits with
regret that usage is indefinite. The word ``moral'' is used, after all, of purely individual rules of
prudence too, and noble or meritorious action that goes beyond the minimum of what is ``dutiful'' is
also called moral action (p. 30. 35). It seems to me that if one regards usage as an authority, one
must also follow its fluctuations. Living usage is always on the move, and it may be guided by just as
right an instinct when it moves as when it stands still. How right the instinct is that guides it in its
variation one will discover when one lets the words go and keeps to the observation of what different
ends men can set themselves, and what influence this consistently has on their valuing judgments. The
questions that the scientific inquiry into valuing judgments has to treat are then how men come to set
themselves ends, and how these ends determine their valuation of their own and others' conduct. The
fact that usage varies in the way spoken of shows that ordinary consciousness itself is already aware of
the problem in its larger range\footnote[1]{In his book \textit{Naturlig Ret} (1907) \textit{Severin
Christensen} has given a sketch, interesting in several respects, of an ethics confined to formulating
the elementary conditions of human living together. An attempt is made here to isolate the most
elementary part of ethics and to treat it by itself. It recalls Fichte's attempt to set up a doctrine of
right independent of all ethical presuppositions. (See on this below chap. XXXVII, 4). Just as little
as the attempts spoken of in the text does this conception escape the difficulties that follow from the
concept of value. Here again it must be asked whether it may not be a matter of indifference to the
single person who has sufficient power what the conditions of living together are. He transgresses them
when he can, and he observes them only when he can use them for his own ends. (1913).}.

b. The ``biological'' ethics, as Dr. S. Hansen has sought to found it, would dispense with the concept of
value and thinks thereby to escape the difficulties it brings with it. It is granted, to be sure, that
all ethical inquiry turns on whether a certain action is a means to the attainment of an end; but about
% --- p. 73 ---
\opage{73}the end lying at the foundation there is here to be no doubt; it is an end given beforehand, and one
that ethics does not itself set up (p. 107). The end set by nature itself is the existence and advance
of the race, or the greatest possible unfolding of life (p. 24).

Now is this really an end that is absolutely necessary? Dr. Hansen himself grants that other ends
\textit{can} be set. That end is, he says, ``as little arbitrarily chosen as may well be possible''
(p. 28), set up ``with least arbitrariness'' (p. 30); another goal cannot be set up ``with greater
logical right,'' and in any case no goal can be set up that conflicts with this one (p. 234). What
becomes, then, of the logically compelling? Different ends, after all, can be set up, and no single one
has beforehand the absolute ``logical'' right. That is in reality all I have asserted, and it is therein
that for me the difficulties about the presuppositions of ethics lie. Dr. Hansen grants me (with whom
on this point he thinks he is ``as much at odds as may well be possible'') everything I wish. For him
too several consistent ethical systems must be possible. I have come to this recognition by the way of
experience, supported by what biology, psychology and sociology have taught me. And if there is anything
the biology of the last half-century has impressed upon us, it is surely above all the struggle for
existence --- a struggle between living beings over whose life is to endure and unfold. It is a
struggle between individuals among themselves and between groups of individuals. We still stand in the
midst of this struggle; it is waged in the human world as well as in the rest of living nature. It is in
it that ethical ideas are to prove their validity, not in the world of abstractions, not for ``life in
general,'' but for quite definite, individually and historically formed stages of life.

According to Dr. Hansen the compelling end is really given with the definite outcome of the struggle.
He grants that he uses the word ``end'' in a transferred sense (p. 24). It really means only a tendency
of movement in a certain direction, so that one may with the same right say that the rolling avalanche
has its ``goal'' in the valley, or the river its ``goal'' in the sea.
% --- p. 74 ---
\opage{74}But can one then everywhere set down the actual result as the goal? It is precisely here that the
concept of value announces itself as unavoidable. For if we could, we would stop the avalanche when it
threatens to destroy the town down in the valley, and divert the river's water when it can be used to
make our fields fruitful. The unavoidable we must put up with; but it does not always stand as something
that can become a goal for us. Thereby ethics becomes something other than the acknowledgment of the
right of the stronger and submission to blind, brutal facts. Here the proud word holds that was said of
Cato the Younger, that the gods held with the victorious cause, Cato with the vanquished. The victorious
does not always have value on its side. Experience shows us both development and dissolution. Which of
them is really nature's ``goal''? Mephistopheles holds that everything that comes into being is worthy
to perish, and that it would therefore be best that nothing came into being. How will one be able to
refute him? A consistent advocatus diaboli would in any case hardly give in before lyrical phrases such
as, for example, ``the eternal right of life to assert itself'' (p. 111). How can there be question of
right before ethics has come into force? The ``biological'' ethics expressly maintains, after all, that
the ``goal'' itself cannot be grounded ethically.

When experience shows us different tendencies of movement, there is nothing else to do but to hold fast
and favor the tendency that has the greatest value for us, and value is experienced through the feeling
of pleasure. Here, then, we come into the subjective world, which one cannot get round, either in ethics
or in logic, unless one wishes to dogmatize. But there exists no gulf between the subjective and the
objective. All thinking about values leads to examining the conditions of their arising, and thereby the
possibility of producing them anew and developing them further. No tendency in ethics has impressed this
with greater energy than the morality of utility or of welfare. According to it ethical work consists in
converting the natural energies into such forms as enable them to work toward valuable results. Human
impulse and passion do not --- any more than the avalanche and the river --- always work in the direction
of valuable ``goals''; therefore ethics demands a
% --- p. 75 ---
\opage{75}metamorphosis of them, whereby they can --- with the conservation of energy --- work in other
directions. What is demanded is a conversion from the real to the ideal. Plato was the first to set this
task in the grand style (in the ``Symposium''). This conversion is to take place not in a world of
fantasy but in the real world itself. That is why I have also always maintained a close connection
between the problem of valuation and the other main problems. ``What is for me the highest end, which
determines the value I attribute to everything else, must surely be so constituted that it can be
asserted and striven for within the world whose reality can be shown by the most consistent criterion I
am able to apply [namely the causal connection shown by science]. Something can prove itself valuable
only through the way in which it intervenes in the real world. And on the other hand this real world
itself acquires greater or less value for me according as it proves to further or hinder a striving
toward the end that determines my judgments of value''\footnote[1]{\textit{Det psykologiske Grundlag
for logiske Domme.} (Vid. Selsk. Skr. 6. Række) p. 65.}.

In all fields the setting up of a principle has first and foremost the significance of distributing the
burden of proof. The significance of the causal proposition is that the burden of proof lies on him who
would maintain that an event has no cause. And in an analogous way the principle of welfare means that
the burden of proof lies on him who maintains that a pain is better than a feeling of pleasure. When I
strike the cup out of the hand of the thirsty man, I must show that it is harmful to his health to
drink, or that the water was poisonous. A pain or a disappointment may, according to the principle of
welfare, be more valuable than a feeling of pleasure, if it is a necessary purgatory, or if it is a
contributing element in a more valuable state. The ethical consideration here relies in the end on the
biology of feeling. Therefore one may very well stress the subjective element in ethical
judgments\footnote[2]{It has escaped Dr. Hansen's attention that by ``ethical judgments'' I mean from
the first (see I, 1) the judgments of good and evil that men pass in real life. The ethical problem is
posed when it is asked whether such judgments can be grounded. Scientific ethics is a system of grounded
ethical judgments. With the grounding a closer determination and further development will follow.}
without giving up the possibility of a grounding.

% --- p. 76 ---
\opage{76}c. Professor Kroman agrees with me that one cannot get round the concept of value, and that in all
ethical judgments passed in real life an element of feeling of decisive significance cooperates: ``The
expression `for the sake of' can be used with meaning only with regard to a feeling being'' (p. 52). But
he is convinced that man ``in the last resort'' must combine the different elements of feeling in a way
free from contradiction. The right action, the one that the ideal conscience demands and acknowledges,
will be ``the one agreeing with man's inmost abiding nature, the one free from contradiction in the last
resort'' (p. 37). The ethical is no historical accident, but ``something necessarily connected with
man's inmost abiding nature, an outcome, constant in all essential features, of the foundation common to
all individuals'' (p. 44).

It is indisputable that when a contradiction makes itself felt between different interests or
tendencies, a need must arise to get it out of the way, either by a choice between the opposites or by a
harmonizing of them. In this way contradiction drives forward. But the great question is precisely
whether the contradiction is present in the particular case, and in what \textit{degree} it is felt. And
it is here that experience, psychological and historical experience, shows us extraordinarily great
differences. These differences are by no means historical accidents. They must be assumed to have their
definite causes, which it is the task of psychology and history to find. One has no right to start from
the assumption that all tendencies of feeling are found in every man in such a degree of development that
the only thing needed now would be to unite them in a way free from contradiction. Therefore I cannot get
over the different stages and standpoints as easily as my colleague. I have, to be sure, sought to show
that Lombroso's well-known theory of the criminal man cannot be grounded by
experience\footnote[1]{Cf. \textit{Etiske Undersøgelser} (1891). p. 73---81.};
% --- p. 77 ---
\opage{77}but the opposite assumption, that all the elements an ideal conscience demands are present in
everyone, in every race and at every time, cannot be grounded either.

Even if all the elements were present, it would still be of importance in what relation they stand to
one another. For a contradiction to be present --- a contradiction of such a kind that it can work as a
motive (both motive of valuation and motive of action) --- the different elements must each come forward
with great energy. No contradiction comes forth when one element sounds only as a faint tone, a distant
ring, while the opposite deafens the ear of the soul with its pealing. A certain development is thus
required for the contradiction to be able to assert itself, and especially for it to become a driving
and motivating force. So long as development is still going on, in the single individual and in the
race, ``the ideal conscience'' therefore stands as an abstraction not covered by any actual standpoint.
Each of us, every generation and every people, stands at a certain definite stage of development, and by
this definite stage of life all values, ends, ideals and norms are determined. If one would build ethics
on the foundation of experience, one cannot pass by the difficulties that lie herein. \textit{When} are
we at ``the last resort''? And how long must a tendency express itself and prevail in order to deserve
the predicate ``abiding''? Here there is always only the possibility of approximation, in life and in
thought.

For the rest, I believe that one may not regard the contradiction between the elements of human nature as
the only condition of development. Often biological-psychological development takes place as a positive
growth, in that an element unconsciously and involuntarily attains such a degree of development that it
gains the upper hand in the soul and determines the timbre with which the whole interplay of the elements
of the soul appears. Where in this way a single element becomes prevailing through expansion,
contradiction plays no part. And it is perhaps the greatest things in the field of the soul that take
place in this way. It is by no means so that we always first have a series of single values and ends,
and then get rid of the contradiction in one way or another.
% --- p. 78 ---
\opage{78}There are values that involuntarily come to prevail, and there are ends we set ourselves because we
have grown without knowing it. Not only can it go as in the parable, that we sow the seed, and it then
grows, becomes stalk and sets ear, while we sleep; but even the seed may have been sown without our
suspecting it. As I have said elsewhere, it holds not only that man grows with his greater ends, but
also that he must have grown in order to be able to set himself greater ends. And besides this
unconscious growth there may take place displacements of value and of motive whose results cannot always
be foreseen. (Cf. \textit{Psykologi} VI C and the present work XIII, 4). Nor in such transitions is
contradiction the decisive thing. A psychological process is, after all, something other than a formal
logical train of thought, and man is, to be sure, ``a thinking being,'' but his thought unfolds at each
given stage under the influence of quite definite historical conditions.

The psychological-historical foundation whose consequences and special applications I examine in my
ethics appeared in Stoicism in the last centuries before Christ; then for the first time the idea of a
general feeling of humanity (caritas generis humani) was set up as the most comprehensive motive of
valuation. Deepened and made inward by Christianity, and clarified through the intellectual and social
development of modern times, it poses for ethical thinking the most comprehensive tasks any system can
wish to treat. Among these tasks is also that of deciding how one is to stand, in thought and in life,
toward those in whom that foundation is not developed, or who do not acknowledge it. A demonstration of
self-contradictions may here have its significance. But man's nature is too composite and too changeable
for the mere striving to be freed from self-contradiction to be able to furnish a sufficient foundation
for the ethical.

% --- p. 79 ---
\opage{79}\chapterhead{IV.}{THE THEORY OF CONSCIENCE.}

1. In the preceding chapter (§ 5 and 9) conscience was described chiefly as a feeling. But no feeling
is wholly without connection with the life of consciousness as a whole, and it is important to stress
this in particular in order to understand the nature of conscience.

It is interesting that the concept of conscience is from the first closely connected with the concept
of consciousness. The corresponding word (συνείδησις) was first used among the Greeks in the general
sense of self-consciousness, the capacity to apprehend one's own inner states, before it was used in
the more special sense of ethical self-valuation. And the German Gewissen (= Mitwissen) likewise first
had a more general theoretical meaning before it acquired the special ethical
meaning\footnote[1]{Rudolf \emph{Eucken}: \textit{Geschichte der philosophischen} Terminologie.
Leipzig 1879. p. 175. --- The Danish ``Samvittighed'' was taken over from Low German (probably in the
14th or 15th century). The Low German samwitticheit had (as a translation of conscientia) the meaning
``consciousness'' and the meaning ``conscience.'' Older than it is Old High German gewissen or
gewizzida, which meant knowing, knowledge, and only later (after the Latin conscientia) took on the
meanings consciousness and conscience. I owe this information to colleagues in linguistics.}. The
expression conscience implies that man has the capacity as it were to divide himself into two persons,
one of whom is the object of the other. We are thereby referred to the general mark of consciousness:
a combining activity by which the single experiences are placed in relation to one another. All life
of consciousness is the expression of a need for unity and coherence, for a synthesis. And in
conscience this peculiarity expresses itself in a particularly inward and concentrated way. However
great the opposition may be between what man recognizes as the ideal and his own actual
% --- p. 80 ---
\opage{80}willing and acting, conscience nevertheless leads man to own both: ``That is my ideal, however much
it conflicts with my conduct!'' or: ``That I have willed and done, however much it conflicts with my
ideal!'' The arising of conscience presupposes that within consciousness a difference makes itself
felt between a central content, what prevails and rules, and a peripheral content, the single passing
thoughts, moods and utterances of will. Where an idea of a chief end, an ideal, a model, has developed,
it must --- together with the feelings and impulses attached to it --- belong to what is central in
consciousness. Conscience is then a form or a part of the reaction of the central against the
peripheral, a feeling of relation that acquires a different character according as there is a
harmonious or a disharmonious relation between the central and the peripheral in the personality.
There may be great resistance to overcome before the immediate holding together and comparing of the
highest and the lowest that consciousness contains can come about. The descent into hell of
self-knowledge, to use Kant's expression, requires the summoning of great spiritual energy. But it
comes about according to general psychological laws and only expresses the general nature of
consciousness in a particularly marked way. There is therefore no reason to take conscience for a
special faculty or sense, implanted once for all and in all\footnote[1]{I do not know on what
\emph{Victor Cathrein} \textit{(Philosophia moralis,} Friburgi 1895. p. 60) relies when he attributes
to me the assumption of a special ``sensus moralis.''}. The possibility of conscience lies in the
general nature of consciousness; but whether this possibility is developed, and in what direction it
is developed, --- what special content conscience thus acquires, --- depends on the social conditions
under which the individual comes into being and develops (see I, 2). Since the relation of contrast
between the different elements plays so great a part in the arising of conscience, it is
psychologically intelligible, as experience seems to show, that the ``bad'' conscience arises before
the ``good'': the contrast is naturally
% --- p. 81 ---
\opage{81}greatest where there is present not merely a difference but a sharp conflict. Personal ethical life
therefore often first really comes into being through one's feeling oneself in actual conflict with
what one takes to be right.

The special character of conscience will depend on what ends or interests of life are the highest for
the single person and make up his real self. On this, after all, depends also, as we have seen, the
difference between the different possible ethical systems, all the way from the egoistic to the
universally humane\footnote[1]{Conscience, when the different fundamental values are disregarded, has
a purely formal character. But even the purely formal has its value. That is why I have distinguished
in \textit{Den menneskelige Tanke} (p. 241---245) between formal and real differences of value. To the
first belong the differences between elementary and ideal values, between immediate and mediate values
and between potential and actual values. In its purely formal character conscience is a guardian over
a harmonious relation between these opposites. To another group of differences of value belong partly
the difference between biological (economic), intellectual, aesthetic, ethical and religious values,
partly the difference between individual, social and cosmic values. --- Cf. also my \textit{Psykologi}
VI C, 8 a. (1913).}. But however comprehensive the interests of life to which the single person has
attached himself, he nevertheless always feels himself in conscience as a harmonious or disharmonious
whole. His real I is the one that lives in those interests; his outer I is the one that appears in his
actual willing; and in conscience he perceives the relation between his real and his outer I, both of
which, after all, belong to his personality. If he denies one of them, conscience does not come into
being; but then he also deceives himself. There expresses itself in conscience, as F. C.
\textsc{Sibbern} has shown, a deep care for oneself, a spiritual impulse of self-preservation. And it
expresses itself first and foremost in one's wishing to be honest with oneself, --- even if one has not
been true to oneself.

It is not always, and not from the first, that conscience has this free, personal character. It may be
bound to authority, and its task is then to hold the pronouncements of authority up against the
individual's own willing and acting. The real
% --- p. 82 ---
\opage{82}self is then determined by obedience and reverence, and what matters is whether these feelings are
injured or satisfied by a man's conduct. The spiritual position of power of the Catholic Church rests on
its having as it were branched out through the confessionals into the personal centers of its
adherents. Great firmness and boundness in certain things is thereby brought about, but at the same
time great freedom in other things. For where authority does not speak clearly, I have according to
Catholic ethics\footnote[1]{Victor Cathrein: \textit{Philosophia} moralis. p. 152. --- Alfonso de'
\emph{Liguori}, one of the authorities of Catholic moral theology, enjoined a certain caution in
weighing the probability that a law might be the right one. He demanded that the possibilities between
which the individual had free choice should not merely each have a serious ground at all, but also that
on express holding together and comparing they should stand as equally or almost (sic) equally
probable. This somewhat stricter doctrine is not, however, the prevailing one in the theory and
practice of the Catholic confessionals. On this question see J. P. \emph{Gury}, S. J.:
\textit{Compendium theologiæ} moralis\textsuperscript{6}. Lugd. 1899. I. p. 133---140.} --- by
virtue of the principle that an uncertain law is no law (lex non certo promulgata non obligat), ---
the right to follow whatever opinion I will, even the \textit{least} probable --- provided only it is
still probable! --- One reasons here: whether I drown in 10 or in 3 fathoms of water is of no
consequence! Once I am outside the reliable pronouncement of authority, I need not be so particular!
With boundness to authority probabilism consistently follows.

It is otherwise where the individual does not have the standard outside itself, but where the decision
is placed in its own hand. The free conscience does not find it so easy to feel secure; it does not
know the abrupt boundary between certainty and uncertainty. It has both a more inward character ---
since the single person is here his own judge --- and a more comprehensive task --- since it must
constantly seek to learn from life.

It is not always that the end or interest of life with which the actual willing and acting is held up
for comparison itself comes forth in consciousness. Approval or disapproval may be felt toward an
action without one's being able to account for what it is that has been satisfied or injured. We
% --- p. 83 ---
\opage{83}have here an analogy to such cases where an effect of contrast sets in against a sensation that cannot
itself be shown to have been in consciousness, or that was at any rate there so briefly and fleetingly
that it was not noticed and remembered\footnote[1]{\textit{Psykologi} III, 6. --- \textit{Det psykol.
Grundlag for log.} Domme. p. 35---39; 47.}. In such cases conscience acquires a mysterious character,
which is stressed especially by those who deny the natural origin of conscience and regard every
attempt to explain its arising by a psychological-historical way as ``materialism.'' The explanation
lies in many cases in this, that a constant interest can, through repetition, sink beneath the
threshold of consciousness, but that it does not for that reason cease to work. There is a constant
relation of interaction between the conscious and the unconscious.

In such cases conscience takes on the character of an instinct. Instinctive working is everywhere
characterized by serving an end that is not itself an object of consciousness. The obscure and
unconditional way in which conscience often works can perhaps be explained only by the assumption that
inherited elements cooperate. One may rightly speak of a conscience of the race, and the common
tradition will not always be a sufficient explanation. But it must nevertheless --- even if one accepts
the assumption of hereditary elements in conscience --- be firmly held that inherited dispositions are
always indeterminate, and that the way in which they are developed will depend on the social conditions
and traditions in which the individual grows up. The single individual involuntarily appropriates the
positive morality and through it acquires the ends and interests of life that determine its life of
conscience. The custom and usage prevailing in society are immediately felt as binding. Just as the life
of feeling in general develops under the constant but quiet power of the conditions of
life\footnote[2]{Cf. \textit{Psykologi} III, 7 and VI F, 4 c.}, so the life of conscience in the single
person cannot always be traced back to quite definite and single experiences either. The mysterious
character of the ``unwritten laws''
% --- p. 84 ---
\opage{84}rests in part on this. The child hears that those around it approve or disapprove of its conduct,
that it is the object of ethical judgments, and in its great urge to imitate it repeats these judgments
on occasion and so involuntarily appears as its own judge, regards itself from the standpoint of those
around it. At the same time it is led --- likewise involuntarily --- to imitate those around it in its
conduct; it thus practices the virtues that are particularly required in the society in which it lives.
It acts --- as \textsc{Aristotle}\footnote[1]{Eth. Nic. II, 1.} says --- well before it becomes good,
just as only by playing the cithara can one become a cithara-player. The action thus far precedes the
capacity. Through both passive and active imitation the child thus grows into the practical ethics of
society and thereby acquires the first form of the life of conscience. Conscience does not yet exist
here as a free, individual capacity for valuation, but is an echo that --- perhaps deeply and inwardly,
but yet from without --- resounds within the single person. As it always goes with the child, so it
goes at earlier stages of culture --- and everywhere where positive morality rules alone --- with men
in general. Their ethical ideas and their conduct of life are determined first and foremost through the
social milieu. Fear and awe of the gods of the tribe, who represent, after all, the experiences of the
past, respect for custom and usage and for public opinion (δήμου φάτις) determine --- as the oldest history of
the Greek people in particular teaches us --- the judgments and actions of men at this stage. Awe of the
gods, for example, holds men back from using poisoned arrows (Odyss. I, 261---363), % the print has 363 for 263 (PRINTED AS IS in the Danish)
 and a son who in
anger would kill his father is held back from it by the thought of ``the talk of the people and the many
reproaches of men'': he feared to be called a parricide! (Iliad. IX, 460 f.). In its noblest form
ethical conduct appears at this stage as determined by a considerateness that can rise to modesty and
reverence (αἴδως), and which is felt not only toward the aged and the mighty, but also toward equals,
indeed particularly toward the poor and the suffering, where it becomes one with pity. In a
characteristic way we see here sympathy as a motive of valuation
% --- p. 85 ---
\opage{85}and of action gaining influence on positive morality and ennobling it, keeping it from becoming mere
regard for what others say. The Greeks distinguished between that noble modesty and the feeling of shame
(αἰσχύνη) that is only a shrinking from incurring blame. The transition to conscience proper in a free,
individual form is marked by the demand that \textsc{Democritus} set up: ``Have modesty before thyself,
and thou wilt not come to be ashamed before anyone!''\footnote[1]{P. \emph{Natorp}: \textit{Die Ethica des
Demokritos.} Marburg 1893. p. 10. 102. Cf. also the interesting inquiry into the concepts of modesty
and shame in L. \emph{Schmidt}: \textit{Die Ethik der alten Griechen.} I. p. 168---186.}.
\textsc{Socrates}'s demand for self-knowledge as the condition of all right action was the most
important step toward moving the center of gravity from the outer to the inner, and the emancipation of
conscience is completed by \textit{the Stoics}, who also fixed the expression conscience in its ethical
sense. --- With the dissolution and downfall of ancient culture this development was broken off, and
with the ecclesiastical dominion of the Middle Ages the principle of authority came forward in the
sharpest and most conscious form it has ever had. \textit{Protestantism} and \textit{modern philosophy}
again asserted in practice and theory the cause of the free conscience\footnote[2]{See my work:
\textit{Den nyere Filosofis Historie}\textsuperscript{2}. I. p. 8---74.}. So long as conscience is only
``the echo of the authorities,'' so long there is an opposition between freedom and conscience: for so
long the appeal to conscience will be one with the giving up of every attempt to ground one's judging and
acting. The peace of conscience then easily becomes a pillow to sleep on, while conscientiousness ---
``obedience to the inner Thou shalt!'' --- becomes a finer kind of bondage. If by conscience one
understands only a social echo, it is consistent when it is proclaimed that ``the free man of reason''
must be without conscience\footnote[3]{Bruno Wille: \textit{Die Philosophie der Befreiung durch das
reine Mittel.} Berlin 1894. p. 243 ff.}. But, as we have seen, conscience is closely connected with the
general nature of consciousness. And only when self-knowledge has awakened, so that man can become
conscious of his highest ends, will a clear and thorough
% --- p. 86 ---
\opage{86}examination of the justification of judgments and actions be possible, and not merely possible, but
a spiritual necessity. In particular, as has been shown above (III, 16), the conscience whose motive of
valuation is universal sympathy will necessarily be led to examine, with all the means at hand, the
influence of the action or the judgment on life within us and outside us. How far the examination can be
carried in the single case, nothing can of course be said beforehand. It will go here as with all
knowledge: where one stops, another will perhaps be able to continue. Limited (``borné'') we all are,
more or less.

The free conscience will have a tendency to develop to a higher degree of express consciousness than the
instinctive conscience bound to positive morality possesses. Ideal images will involuntarily form that
depart more or less from the traditional models. In place of the mere instinct that works involuntarily
there then steps an impulse with an idea of an end. Experiences and reflections determine the
development of the ideal images. The individual tries to assert his ideal against resistance within
himself or from without; it thereby stands forth more clearly before him, but at the same time there
arises the possibility of altering and reshaping it: the resistance experienced may, after all, at times
be recognized as justified, or at any rate lead to renewed self-examination. And even if man from
weakness is not always able to assert his ideal against the resistance, the attempt is not without
effect; strength is exercised and may perhaps grow strong for the next struggle. It is through such
trials that conscience gains its experiences. In particulars it unfolds by virtue of general
psychological laws. Now it is by virtue of the laws of the association of ideas that thoughts and
misgivings come up; now it is to an effect of contrast that the awakening of conscience is due. Through
fusions of ideas and blendings of feelings, through displacements and deferments, as well as through
express concentration of the will in definite directions, there forms a resultant of the individual's
experiences that determines the way in which it stands toward new calls to judge and to act.

% --- p. 87 ---
\opage{87}Without wholly losing its instinct-like or impulse-like character, conscience can develop into what
has aptly been called practical reason. A need will arise for consistency and totality in ethical ideas.
Actions --- not only the actual ones but also those merely thought of --- are followed out in their
consequences; ideas form of all the beings with the capacity to feel pleasure and pain that are touched
by the actions and their effects. Here a developed historical experience will be of particular
importance. With its help later Greek philosophy was led beyond the opposition between Hellenes and
barbarians, and with its help modern philosophy was led to a consciousness of humanity that rose above
differences of class and religion. Both historical experience and logical consistency contribute to the
constant self-correction of conscience. Valuation gradually takes place more consciously and according to
definite principles. A philosophical ethics is only a methodical exposition of the content of conscience
as it takes shape at such a stage. What Kant depicted in the categorical imperative and in the formal and
universal moral law of reason was a result of a long historical and psychological course of development.
---

Where conscience works as instinct, the individual does not know what it is doing. Where it works as
impulse, it has a dawning idea of it. But where it works as practical reason, there arises a clear
consciousness of ideals and rules.

2. Conscience appears in the most various forms and in the most various degrees in the single
individuals. Not only can it come forth now as a half-unconscious instinct, now as impulse, now as
practical reason. But now immediate sympathy, now the feeling of duty, now the feeling of justice may be
the chief element in it. It may express itself now predominantly negatively and checkingly, now more
positively, partly heightening, partly uniting. (Cf. III, 20). It may express itself now as enthusiastic
devotion, now as a more quietly and steadily working need. It may now assert with inwardness the
interests of life of a small community, now be determined by a comprehensive view of the welfare of large
circles. It would be impossible to depict even merely the chief forms of its appearance in the single
% --- p. 88 ---
\opage{88}individuals. But it is of great importance to be attentive to these differences, since we still
suffer from a dogmatism that cuts everyone to the same pattern. It will follow precisely from the
principle of welfare that every individual ought as far as possible to live and work according to its
peculiarity, and nowhere can this hold in a higher degree than when it is a question of the ethical
foundation itself. Precisely where the greatest responsibility rests on the individual, since it is to
assert the highest interests it knows, the individualization of conscience will be a necessary demand:
for only by virtue of it can the individual work with full strength. Where the individual has not acted
according to its full peculiarity, it has (as \textsc{Schleiermacher} has said) not been wholly active
either. And the more the individual in its most significant action can work out of its inmost
peculiarity, the more it stands not merely as a means but also as an end; its own deepest satisfaction
then accompanies its action. ---

One must surely go a step further still. There may be men who possess no ethical feeling proper, and
who do not need any either. What they can give they give with all their heart, without making any
valuation of their own or others' actions. They are wholly absorbed in an activity that agrees entirely
with their faculties and impulses, without ever doubting its justification and significance, and without
their zeal cooling. It may be the cultivation of art or science, activity in the service of society or
for the support of their family. Or they belong to the happy natures that spread light and joy about
them merely by existing. Whether there are many men of that kind is not easy to say; but there is
nothing to prevent there being such at all. From the heart they do the works of the law without having
the law, and what should ethics have to object to that? Ethics exists for the sake of life, and not life
for the sake of ethics.

There may thus be extraordinarily great differences between individuals in respect of conscience, and it
would be dogmatic to lay down here one definite type as the only right one. These differences come out,
of course, especially
% --- p. 89 ---
\opage{89}where one moves away from positive morality; for this has in all its forms a tendency to regard all
individuals as alike. The attempts to develop a philosophical ethics bear witness to the differences
that are actually present, in that the individual philosophers each describe and set out the kind of
conscience they know best. A ``personal equation'' makes itself felt here\footnote[1]{Cf. Frank Chapman
Sharp: \textit{The personal equation in Ethics.} (Transactions of the Wisconsin Academy.
1894---1895). --- Sharp applies the concept of personal equation at another point than I do in
``Etiske Undersøgelser.''}, and the history of philosophical ethics has, among other things, the great
significance of drawing attention to the typical differences in the ethical field. And yet most
philosophers hold the assumption that the ethical demands, the content of the ethical law, must be the
same for all, so that conscience can only become a formal capacity to acknowledge and fulfill the law.
It might, however, be that the differences concern not merely the form, manner or degree of
consciousness of conscience, but also become determinative for the content of the ethical law in both
qualitative and quantitative respects.

If by an ethical law one understands a sum of demands laid down by an external authority or by absolute
reason without regard to the individuals for whom they are to hold, then it is clear that individual
differences cannot come to influence the content of the law. ``The objectively right'' is demanded ---
and demanded in equal degree of all. But when one acknowledges the principle of welfare, one must also
hold firmly that the right ethical law cannot have been found so long as the individual's personal
peculiarity does not come into its own. Neither in respect of the kind of faculties and motives nor in
respect of their degree are the different individuals alike, and a law that in qualitative and
quantitative respects would demand ``the same'' of all would thus in fact demand a most different labor,
lay most different burdens on the shoulders of the single persons.
If there is to be equality before the ethical law, then ``the same'' must not be demanded of all, and
the single person's conscience must
% --- p. 90 ---
\opage{90}therefore make demands that in qualitative and in quantitative respect would be different from those
that others' conscience makes. Now something other and more, now something other and less, is
demanded.

In the qualitative respect this will most readily be granted. Every individual has --- on account of
its special endowment --- its calling, its special direction in which it is to work. Only where a
caste system in one form or another exists is the same direction of activity demanded of the most
different individuals. Under freer forms of society it is recognized that man's tasks are determined not
merely by general laws and external circumstances, but first and foremost by his own peculiar nature.
This was something people already had an eye for in Greek antiquity. \textsc{Socrates} enjoined the
necessity of knowing one's nature and one's faculties in order to find one's tasks. \textsc{The Stoics}
carried this thought further: he who does not follow his own nature will not be able to preserve the
personal coherence and unity of his life; therefore one should do as the good actors do, who do not
choose the greatest parts, but choose the parts that best suit their gifts!\footnote[1]{Cicero: De
officiis I c. 31. --- \emph{Cicero} here follows \emph{the Stoic Panaetius} (of the 2nd century B.C.).
See \emph{Schmekel}: \textit{Die Philosophie der mittleren Stoa.} p. 41; 219.} --- Conscience has here
to guard the individual peculiarity, but at the same time to keep self-delusion out.

It will be harder to grant that the ethical law also varies in quantity for the different individuals,
and yet it is not possible to see how this consequence can be avoided if one holds firmly enough to the
difference between an ethical and a juridical law. Ethically regarded, a far greater labor may have been
performed, and a far greater self-control and self-sacrifice shown, by a man who nevertheless does not
come up to the average level demanded from without, than by many a one who easily satisfies the
``social'' demand but who has on the other hand done no work to develop his faculties beyond the
average level, although there was possibility of it. And should the former, ethically regarded, then
nevertheless stand below the latter? Such an absurdity
% --- p. 91 ---
\opage{91}every ethics falls into that takes up a purely objective standpoint. The principle of welfare
demands, as will be shown later, as a necessary consequence, that a personal being must at no point be
treated as a mere means. But that would be the case if ethics had said its last word with the
formulation of a universal law that drew a line through all personal differences. The true ethical law
must individualize, must state for the single person the work he can do. The virtues and duties
demanded of the single person must be means and ways to the development of \textit{his} personal
peculiarity. --- Ethical laws are therefore not so easy to find as is often thought. The task must be to
find out how the single person, by following his law, satisfies the general law in the best way. And
this task is solved when the general law, valid for all, states only the direction in which one is to
strive, the type that is to be held fast. Only with dogmatic arbitrariness and by overlooking the
deepest-lying psychological and ethical problems can one fix ``the qvantum satis of manly will'' for all
at once. At any given time a society will, through its leaders or through public opinion, fix certain
qualities in a certain quantity as those that can be demanded and as standing within the power of every
``normal'' man. What is understood in the particular cases by words such as courage, self-control,
etc., is due to the average experiences one thinks one has had concerning the resistance that can be
overcome. One here applies an estimate that one trusts holds for all individuals. One does not see that
a personal equation holds here, which sets a limit to universally valid ethical reasoning. If men were
more alike, it would be easier to be a moral philosopher. But to make the matter easier than it is
surely cannot be good philosophy\footnote[1]{Cf. \textit{Etiske Undersøgelser.} p. 51---83, where I
have set out and illustrated in detail the individual relativity or the personal equation in ethics.
--- In antiquity the doctrine of the personal equation was hinted at by Aristotle, in modern times by
Hutcheson, Jacobi, Schleiermacher, Beneke, and in our own literature by A. S. Ørsted (Om Grænserne
mellem Teori og Praxis i Sædelæren. 1811. Reprinted in Eunomia I) and by S. Kierkegaard
(Uvidenskabelig Efterskrift. 1846. p. 376.). (Samlede Værker VII. p. 426). (``Where is the boundary for
the single individual in his concrete existence between what is lack of will and what is lack of
ability, what is sluggishness and earthly selfishness and what is the limitation of finitude?'') Cf.
already Holberg: Moralske Tanker. 1744. p. 463.}. It is no solution of the
% --- p. 92 ---
\opage{92}problem to distinguish between valuation and imputation and to leave it to the latter to take account
of individual differences in disposition and faculties. This way out Alexius \textsc{Meinong} has
attempted: ``Moral valuation in its generality cannot enter into personal particulars, while the
imputation directed at the single personality can hardly be personal enough''\footnote[1]{\textit{Psychologisch-etische
Untersuchungen zur Werth-Theorie,} Graz 1894. p. 201. (Cf. my review of this work in Göttinger gelehrte
Anzeigen. 1896. Nr. 4). --- On the juridical and the theological conception cf. my \textit{Etiske
Undersøgelser.} p. 63.}. This, however, is an all too juridical way of looking at things. Jurists have
been accustomed to distinguish sharply between the general law and the individualizing treatment in
punishment, and between sentence and pardon. Ordinary opinion and public opinion also apply this
distinction; so too the theologians, when they distinguish between law and grace. But the law and the
sentence are unjust if they do not take the agent's individual presuppositions into account, --- and no
less are the execution of punishment, pardon and grace indefensible if they are not the expression of a
higher justice.

3. Since all ethical judgments presuppose conscience as their psychological foundation, it is \textit{the
highest authority and lawgiver,} and every other authority, of whatever kind may be named, is
subordinate and derived in relation to it. To wish to go beyond one's conscience is to wish to go beyond
oneself. Where I bow to another man whose judgment I trust more than my own, this is justified only if it
happens by virtue of my conscience. Conscience is infallible, if by its infallibility one means this,
that at every single moment it is
% --- p. 93 ---
\opage{93}the highest judge. Conscience, to be sure, where it stirs in earnest, must lead precisely to seeking
an objective decision with the help of the best grounds and the most exact experiences that can be
procured; but the application of the objective criterion cannot always be carried through before the
moment for action has come, and we are then compelled to act according to the best conviction we have
won up to that very moment. There is thus not, as has been objected against me\footnote[1]{Giacomo
Laviosa: \textit{La filosofia del diritto in} Inghilterra. I Torino 1897. p. 629---631.}, a
contradiction in declaring conscience infallible and yet holding firmly that it is to be guided by an
objective criterion.

But this infallibility is not of an objective kind. The highest ethical authority may itself prove to be
on a false track. It would be dogmatic to assert that conscience never deceives us. By such an assertion
one arbitrarily shuts one's eyes to some of the most tragic conflicts of life. The purest and most
earnest conviction may rest on a complete error. We have an altogether immediate and infallible mark of
truth just as little in the ethical field as in other fields\footnote[2]{On the criterion of truth in
other fields cf. \textit{Psykologi} V D, 3; VII B, 4. --- \textit{Den menneskelige Tanke.} p.
99---108.}. Every more serious conviction comes forward in the form of conscience; but the mere form is
no guarantee of absolute truth. \textsc{Fichte} grounded the assertion of the infallibility of conscience
on this, that the highest goal for us is to act independently, out of our own personality; if we really
do that, the perfect is reached, as it can be realized at this definite point of existence\footnote[3]{Cf.
e.g.\ \textit{Appellation.} 1799. p. 32.}. But the tragic thing is that the most inward personal
conviction may prove to have been in conflict with what an objective valuation of the circumstances
enjoins. Could not those who crucified Jesus have acted according to their best conviction? Should Kant
be right that an inquisitor cannot have a good conscience? Was it not in good faith that Aristotle
defended
% --- p. 94 ---
\opage{94}the justification of slavery, that Calvin, with Melanchthon's approval, had Servetus go to the stake,
and that Sand murdered the ``traitor to his country'' Kotzebue? ---

No less dogmatically does one cut the knot from the opposite side. Where one regards ethics merely from
the objective side, as the doctrine of forms of society and external actions, one declares without more
ado that subjective feeling is unjustified over against the objective relations and their demands.
Ethicists of standpoints as different as \textsc{Hegel} and \textsc{Bentham} meet in distrust of
conscience. According to \textsc{Hegel} it even stands on the very point of being the evil principle in
the world, since subjective conviction can just as well aim at something objectively reprehensible as at
something objectively right\footnote[1]{Cf. \emph{Hegel}: \textit{Philosophie des Rechts.} § 139---140. ---
\emph{Bentham}: \textit{Principles of Morals and Legislation.} ch. II. § 11---19. --- \emph{James} Mill
blamed a bad action equally severely whatever its motive had been. Praise and blame were for him
motivating forces; therefore they must not be influenced by the motives of the action. (\emph{Stuart}
Mill's \textit{Levned.} Danish transl. p. 52). There was in his opinion no motive that under all
circumstances works in the right way; not even the social feelings; therefore it is of constant
importance to calculate the effects of the action according to the principle of utility. \emph{James}
Mill: \textit{Fragment on Mackintosh.} London 1835. p. 176 f).}. --- It is easy enough, as these
ethicists do, to warn against self-confidence and self-will and to require that one bow to an objective
law. The law we follow must, after all, always make itself known to us in the form of conscience. The
light that makes everything else clear to us must in the end be found in ourselves, and who guarantees
us in the particular case that it is not a will-o'-the-wisp we let ourselves be led by?

Here there lies the possibility of a conflict between the two principles on which ethics builds, the
foundation and the content (the motive of valuation and the principle of valuation). There can hardly be
any other solution of it than the one already indicated, that at the moment conscience must be obeyed,
since there can be no higher authority. It is presupposed, of course, that it speaks clearly and after
due deliberation. But to this it must be added that conscience can control and correct itself; the
later, more
% --- p. 95 ---
\opage{95}practiced and experienced conscience judges the earlier. Precisely by following his best conviction at
the moment man can make his way to a more correct conviction, which teaches him that what at the moment
of resolve stood before him as the best was nevertheless wrong and pernicious\footnote[1]{The appeal to a
good conscience is often made in such a way that one precisely bars oneself from the possibility of such
a correcting. Cf. \emph{Albrecht Ritschl's} apt remark: ``Man kann beobachten, dass wer im Falle des
Streites sein Gewissen als die höchste sittliche Instanz für sich ausspielt, und sich auf keine weitere
Ueberlegung sittlicher Regeln einlässt, hiebei meistens eine Gereiztheit hervorkehrt, welche kein ganz
gutes Gewissen verräth.'' \textit{Ueber das Gewissen.} p. 17.}. So long as one acts according to one's
best conviction, a man's inner being is sound, however the case may stand with the objective character of
the action; and this inner soundness was the germ out of which the more correct insight could develop.
Therefore, from an ethical point of view too, a pernicious action performed in the conviction that it was
good can be placed above a good action performed in the conviction that it was bad. In the first case the
source itself was pure, in the last it was corrupt. In the personal inwardness with which what is
regarded as true and good is held fast and practiced lies the point at which one individual best
understands another, one generation can best sympathize with another. In the field of external actions
and forms of life the differences and oppositions are often so great that understanding is impossible.

Only he who has the courage to err can accomplish anything great. Willingness to expose oneself to such
error is willingness to suffer for the truth. Nor is it, then, the cold and the narrow-hearted, but those
who have enthusiasm for the true and the good, who err in this way. Nothing venture, nothing win.

But the capacity of conscience to correct itself is developed only when a \textit{definite principle or
criterion} for valuation can be found. Such a principle conscience can of course not let be prescribed to
it from without; it must spring from conscience's own nature. I have already sought to show (see
% --- p. 96 ---
\opage{96}III, 10---15) that when conscience is determined by a disinterested and universal sympathy, the
principle of welfare is the only criterion that can come into question. It expresses nothing but what
lies in the end that a personality animated by sympathy must set itself.

Ethical development, however, will hardly ever go forward without conflicts between subjective conviction
and the objective connection of things, repeated again and again. Just as it is an infinite task that is
set our theoretical science, namely to explain reality to us with the help of the principle of
causality, so too it is an infinite task that is set us in the ethical field, namely to value all actions
and conditions of life according to the demands of the principle of welfare.

4. In close connection with the concept of authority stands the concept of \textit{sanction.} Authority
commands or forbids, but it is the sanction that makes the command or prohibition hold good. The sanction
consists in the pleasure or pain attached to the observance or transgression of the command, the reward
or punishment one incurs by one's action, the heaven or the hell one draws nearer to by the action.

Only where the authority itself is an external one does the sanction stand in this external relation to
the action. And in this external form it has no \textit{immediate} ethical significance. The ethical
character of the action rests, subjectively regarded, on its springing from the inmost disposition,
objectively regarded, on its agreement with the principle of welfare. What ethical significance could it
have if to this there were added a pleasure or pain that did not spring from the action itself? Because I
act well now, why should I then at a following moment obtain a feeling of pleasure? Because I act evilly
now, why should I then at a later moment suffer a pain? There is, ethically regarded, nothing
self-evident in this. The whole idea of reward and punishment belongs to pedagogy, not to ethics. It may
be needful for education that reward and punishment increase the power that the feeling of the action's
worth and significance does not always possess in sufficient degree. But it is no immediate ethical
necessity. When it is asserted that the thought of retribution is a thought clear
% --- p. 97 ---
\opage{97}in itself, this is because one is guided by an instinct (to revenge or to gratitude) whose ethical
justification and necessity is not given from the first. The ``self-evident'' is very often the blind.
The involuntariness of the prompting makes one not ask for grounds, and then one thinks that no grounds
are needed either.

The external sanction, which consists in reward and punishment, can, as was said, only be educative. The
ethical sanction must be an \textit{inner} one. It can consist only in the agent's feeling of harmony with
his own highest conviction, of agreement between his ideal (his fundamental will) and his actual willing.
There thereby arises an inner peace that can be stronger than all contradiction and resistance from
without. The individual may have the feeling that there prevails in it a force that could transform the
world if it prevailed in all. It feels its being widened and heightened. The men it most admires it could
confidently imagine as spectators of its action, even if all the motives lay uncovered before them. It
is not pride that stirs here, but a quiet feeling of being, in all lowliness, akin to what is great. The
counterpart to this is the feeling of forsakenness that may seize him who sees that his conduct must
isolate him from all. When Richard III wakes after his dream before the battle, he is in despair because
he knows that no one can love him. Not even the world of dream and of memory can for him enclose the
image of anyone who approves his conduct.

Such an inner sanction is not merely the effect of the action, but consists in a feeling that may already
be present before the action is performed. It is then only the continuation and full development of what
led to the resolve and made it possible. ``Blessedness,'' says Spinoza, ``is not the reward of virtue, but
is virtue itself.'' It is of the same kind as the satisfaction that accompanies the fulfillment of a
deep-lying need. The source of the action and the effect of the action are here closely connected.

This satisfaction can be so great and strong that everything else loses its value in relation to it. A
great and beautiful
% --- p. 98 ---
\opage{98}moment can stand above a long but insignificant life. Self-sacrifice thereby becomes psychologically
intelligible. Between the satisfaction felt in performing a self-sacrificing action and every other
possible feeling of pleasure there may be such a distance that the latter as good as vanishes from
consciousness. The image of life as it takes shape when the self-sacrificing action belongs to it may be
so radiant that life as it would be without this action loses all value. This is only a single example of
a well-known psychological law. Already in the old Buddhist didactic poem \textit{Dhammapada} it is said
(v. 111---112): ``Better than to live a hundred years in blindness and distraction is it to live a single
day in understanding and enthusiasm. Better than to live a hundred years in torpor and feebleness of will
is it to live a single day in the unfolding of strength of will.'' \textsc{Aristotle} too has seen this.
``The good man,'' he says\footnote[1]{Eth. Nic. IX, 9.}, ``will do much for his friends and for his
country; if it is necessary, he will even die for them. He will sacrifice goods and honors, and in general
all the contested goods, in gaining for himself the beautiful; for he will rather feel a great joy for a
short time than a slight joy for a long time, rather live one year beautifully\footnote[2]{By the beautiful
Aristotle understands that which we will for its own sake, not as a means to something else.} than many
years in an indifferent way, rather perform one great and beautiful action than many small ones. So,
surely, it goes with those who die for others: they win for themselves something very beautiful.'' ---
\textsc{Immanuel Kant} expresses a related thought in the following words\footnote[3]{\textit{Kritik der
praktischen Vernunft.} Kehrbach's ed. p. 106 f.}: ``Is not an upright man, in the greatest misfortune of
his life, which he could have avoided if only he had been able to set himself above his duty, held up by
the consciousness that he has upheld humanity in his person in its dignity and has honored it, so that he
has no cause to be ashamed before himself and to shun the inner gaze of self-examination?''

5. How great a part this inner sanction plays in
% --- p. 99 ---
\opage{99}actual life is beside the point here. There are perhaps very few human actions that need no support
other than it. But it is of importance to hold firmly that it can be sufficient to carry the ethical
life. On this fact depends the independence of ethics over against dogmatics and metaphysics, and through
it ethics becomes independent of postulates of faith and of hypotheses.

From an ethical standpoint one must have great misgivings about the way in which the ethical is so often
made dependent on certain definite religious or speculative assumptions. In the first place the
consideration then lies near that he who no longer holds these dogmas has thereby also emancipated
himself from the ethical, and will perhaps even act consistently when he follows the principle: ``Let us
eat and drink, for tomorrow we shall die!'' But next, one robs the action of its proper ethical character
by turning attention to what lies outside the nature and source of the action, and declaring regard for
reward or punishment a necessary motive.

There are those who cannot hold fast to the validity of the ethical except with the help of faith in a
higher order of things, where everything is perfect and complete that in the reality we know is imperfect
and fragmentary. But the need to live in such a faith must not be regarded as a universal human
necessity, even if it is a need that has stirred in some of the most excellent men who have lived.

Not even to faith in a progress \textit{within} the world of experience ought one to attribute an absolute
value for ethics. Such a faith may, theoretically regarded, have its difficulty in being maintained. But
\textit{even if} it should not stand its test, \textit{even if} the victorious tendency in the course of
the world went against ethics, the ethical principles would not therefore be shaken. The tasks would
become different. There would be greater occasion for pity and resignation; but where the ethical
disposition was present, one would hold with the vanquished cause, though the gods held with the victors.
Ethical value does not rest on mere power. He who follows his conscience can of course do so only because
it is the strongest thing in himself; but he does not do so
% --- p. 100 ---
\opage{100}because it is the strongest thing in the world.
The legend of Christopher, who wished to serve the strongest, is therefore unethical, as so many legends
are. Whether the development of the world ends (if it ends at all) as a tragedy or as a comedy, we can
know nothing about, and in any case he who obeys the law of conscience does not on that account change
anything in his part. Our ethics is the ethics of wayfarers (Ethica viatorum). We see the road open
before us, but its end we do not know. We do know, on the other hand, that we can get further, even if
we should not be able to go it to the end; and we know that standing still is decline and death.

In the very fact that conscience has arisen and developed within the human world we have, however, a
testimony that valuable forces are at work in existence. Here, even if it is in a vanishing corner of
the universe, a power has come into being that can carry life and is not afraid to pronounce its
judgment on it. We therefore live not merely ``on possibilities,'' but on a real foundation, which we
seek to make ever broader and stronger. The ethical life is a struggle for a kingdom that is coming into
being, that can grow, and whose fate stands in the closest connection with everything that gives life
value and significance. To this consideration ethics withdraws, and lets the dogmatic and speculative
discussions go their own way.

6. The hour of the birth of conscience is the moment when a feeling arises that is determined by the
difference between ideal and reality. It arises as every instinct and every impulse arises. The hour of
its death would be the moment when this difference fell away for good. By the nature of the case this
can happen in two ways, either by reality getting the better of the ideal, or by the ideal getting the
better of reality.

If for the present we keep to the first possibility, there may here again be different ways in which
reality can, so to speak, stifle the ideal. This may happen through spiritual enfeeblement, through
blasé indifference or through surrender to purely animal impulses. But such cases have an individual
and special character and belong to special ethics. Here, where we are discussing the general
presuppositions of ethics, we ask, on the other hand,
% --- p. 101 ---
\opage{101}whether there might not be \textit{such a scientific conception of reality that there would no longer
be any room for the ideal.}

Such a conception is present, in the opinion of many, in \textit{the modern hypothesis of evolution.}
According to it there prevails, after all, throughout nature, the human world included, an inexorable
law of selection, by virtue of which what cannot adapt itself to the given conditions perishes, and only
what suits the circumstances endures. Physical force, then, seems to be set on the throne, and dominion
in the world to be handed over to brutality. Can our ideals then become more than pious wishes and sighs
that change nothing in the course of things? Through natural selection and in the struggle for existence
even the noblest and best is trampled down, after all, if it cannot or will not meet force with force!

It is not the business of ethics to teach us to know reality. That it leaves to other sciences. It is
able neither to confirm nor to refute any of the hypotheses that have proved justified and grounded on
purely theoretical grounds. But it naturally has a great interest in examining whether such a hypothesis
can be reconciled with the ethical principles.

If we now look more closely at the hypothesis of evolution, it rests, when applied to human life, on a
fundamental thought that is so far from conflicting with the presuppositions of ethics that it must be
said to be precisely one of them, the thought, namely, that human development, as soon as it is more
than mere physical growth, becomes possible only through the awakening of definite wishes and impulses
that can work as moving forces. The greatest hindrance to development, indeed even to mere existence, is
torpor and indifference. In the purely elementary forms of human development it is need that drives
forward: but in the course of development needs arise that are directed toward more than mere
self-preservation. Only where no need, no impulse stirs, and where none can be awakened, is hope gone.
Only he who strives can be saved. This truth, which idealistic conceptions have so often proclaimed in
more or less vague forms, has by
% --- p. 102 ---
\opage{102}the modern hypothesis of evolution only been impressed on us more closely and more urgently --- more
brutally, if you will.

It is not a matter of indifference whether or not we work along with development consciously and
willingly. The very impulse to ethical valuing and ethical acting is one of the forces that cooperate
in determining the course and direction of development. Precisely when we assume that this impulse has
developed according to definite laws of nature, it must be clear to us that it is not in conflict with
the hypothesis of evolution. It has, to be sure, awakened late in its freer and higher forms and has not
yet exerted its full influence. It is in the process of becoming, but precisely therefore it has the
future before it.

Although the modern hypothesis of evolution has so strongly stressed the influence of external
circumstances, it nevertheless also maintains that this influence is in each single case more closely
determined and limited by the inner conditions with which every being steps out into the struggle for
existence. The higher a living being stands, the more it is able to intervene \textit{actively} in this
struggle. And the hypothesis of evolution even leads directly over into ethics by showing us how the
struggle for existence becomes at higher stages a \textit{common} struggle for the preservation and
development of human life.

When the struggle for existence is spoken of, one is inclined to think only of the most elementary and
brutal forms. But life has many forms and stages, and the struggle for existence therefore also acquires
in particulars a very different character. He who works in the service of science and art, or he who
struggles to keep himself up and not to deny what he holds to be right, struggles for life just as much
as he who clings to bare existence and follows the physical instinct of
self-preservation\footnote[1]{Cf. \textit{Om Grundlaget for den humane Etik.} p. 30---32. ---
\textit{Etiske Undersøgelser,} p. 17---23.}. It is the business of special ethics to show what value the
different stages and forms of life have in relation to one another.

7. The hypothesis of evolution not only leads \textit{over into} ethics; with some of its most optimistic
representatives it even leads
% --- p. 103 ---
\opage{103}beyond ethics. According to \textsc{Spencer} the ethical feeling has its place only in an intermediate
period of mankind's development, an intermediate period that we are of course far from having passed
through. All practice brings it about that actions can be performed with less and less resistance and
difficulty and with less and less express attention and exertion of will. If this is applied to the
ethical life, we are compelled to assume that so much ``organic morality,'' or involuntary capacity and
need for right action, gradually accumulates in human nature that a special ethical feeling no longer
makes itself felt and will no longer be necessary either, because harmony will have been brought about
between man's instincts and what regard for the welfare of the race demands. What formerly had to be
made the object of an express command of duty is now practiced instinctively and without deliberation
being needed\footnote[1]{\textit{Data of Ethics.} p. 128. 275.}. Actual human nature would then have
acquired so ideal a character that the time of conscience would be past.

Ethics in and for itself can have nothing to object to the assumption that its time will one day be
past. If it dies a happy death, that is, is replaced by a perfect and harmonious state, that is, after
all, a sign that it has done its work. Perfect beings have no ethics and can have none, since ideal and
reality coincide. Just as the relation of authority (in its various shapes) is replaced by the free
conscience, so this in turn can be thought of as replaced by an immediate living in the good and the
right. There are (as mentioned earlier) perhaps already now some natures of whom this may be said to
have happened.

But even if this possibility must be granted, we are in reality so far removed from such a state that
it becomes a matter of faith or a purely speculative concern whether one will assume that it ever wholly
sets in, and ethics, which is loath to enter into far-reaching speculations, cannot lay great weight on
it. So long as we are still in ``the hobbledehoy age,'' if not in ``the embryonic state'' (to use
Sibbern's
% --- p. 104 ---
\opage{104}expression), one must not weaken attention to the great obstacles that hinder progress by such
fantasies of the future. We can strengthen our hope and courage by the thought that human nature is
undergoing a slow but constant change, and that there may be ground for believing that this change is an
advance. But this assumption will not be able to gain any essential influence on the manner of treating
ethics.

We must surely even assume that progress itself will bring with it new ideals and new tasks. Just as, by
comparing the ethics of civilized man with the ethics of the savage, we find that there are great fields
that in the former have been drawn under the ethical point of view while in the latter they still lie
wholly unheeded, so it will probably go with our present ethics when it can one day be compared with a
higher stage of development. There will surely present themselves many relations whose ethical
significance we have not yet gained an eye for on account of our thick-skinnedness, our ignorance and our
egoism. The ideal will rise anew each time we think we have reached the end.

\textsc{Spencer} assumes (like \textsc{Kant}) that the feeling of duty is necessarily bound up with a
feeling of compulsion and therefore of pain. But the concept of duty need imply, as remarked earlier
(III, 9), only that a more limited regard is subordinated to a more comprehensive one, without this
difference or opposition between a ``lower'' and a ``higher'' \textit{needing} to be felt as compulsion.
The feeling of compulsion may fall away without the time of the feeling of duty being thereby past.

8. Can one do more than one's duty? Can one earn merit by exceeding the demand that can ethically be
made? --- This question can be answered yes only from a standpoint that thinks of the ethical demand as
made on man from without, either from a supernatural authority or from other men; such an external demand
or expectation one can exceed, but not the inner demand of conscience. Catholicism answers yes. It
distinguishes between commandments and counsels. By not only obeying the commandments but also following
the counsels,
% --- p. 105 ---
\opage{105}man earns a special, surplus merit\footnote[1]{The grounding of the difference between praeceptum and
consilium is, however, sometimes carried out in a way that leads precisely to the opposite result from
the one intended. That it is nobler (nobilius) to sacrifice temporal happiness is grounded by
\emph{Cathrein} (Philos. mor. p. 208) thus: ``The giving up of external goods not only removes many
hindrances and dangers to salvation, but also affords constant occasion for perfect exercise in
virtue.'' --- And it should not be a \textit{duty} to make such a sacrifice, when one is able to? It must
be a duty to realize everything of value that one is able to think and will. \textit{It must be a duty to
earn merit,} when one is able to. When one can avert ``hindrances and dangers'' of essential
significance for oneself or for others, it would conflict with a developed feeling of duty to omit doing
so. One of my hearers, a Catholic who has studied at a theological institution abroad, told me that his
teacher of moral theology added to his exposition of the difference between commandment and counsel an
exhortation not to make use of this distinction in practice. --- Saint Teresa found no rest in her soul
until she had earned all the ``merit'' she was able: it was her ``duty'' to earn ``merit''! --- Cf. my remarks in
Göttinger \textit{gelehrte Anzeigen.} 1896. p. 305.}. Popular ethics, which is inclined to think of the
ethical law in greater or less likeness to a juridical law, likewise holds fast to the difference between
the dutiful and the meritorious. And the same is the case with such ethicists as are inclined to make the
ethical demand one with what society or public opinion can demand (Mill and Bain)\footnote[2]{Cf. my work:
\textit{Den engelske Filosofi i vor Tid.} Copenhagen 1874. p. 82 f.}.

It is not easy to see how, if one does not stop at such an external conception of the ethical law, one
will be able to hold fast to the difference between duty and merit. He in whom conscience is sufficiently
developed must feel it his duty to do everything he can to further welfare in the world, and when he goes
beyond the goal that others or even he himself had set, this is only a sign that the goal was not set
high enough. One must not confuse what certain men expect or demand with what can really be reached. Even
the highest sacrifice a
% --- p. 106 ---
\opage{106}man can make is his simple duty when under the circumstances at hand it is really beneficial and
possible. Popular ethics keeps to a certain customary average, is glad when it is reached, and admires
what goes beyond it. It may be pedagogically right to attach special approval (``merit'') to actions
that go beyond the usual; but he who performs such an action has nevertheless the right disposition only
when he feels that he has merely done his duty, and that he stood as one who could not do otherwise. To
the external spectator who comes with his average it may be admirable to see an action that shoots far
beyond it. But such an action may for all that very well have proceeded from the feeling of duty in the
agent himself. The level from which the different individuals start in their judging and acting does not
always lie at the same height; and from a highly developed feeling of duty and from a nature endowed with
great possibilities actions may proceed that for other natures stand as more than duty. --- On this point
too it is the acknowledgment of the great personal differences that decides the dispute.

Even if for pedagogical reasons one would hold fast to that difference between commandment and counsel,
the difference will always be relative: what at a lower stage of development stands as something that is
counseled but not demanded is perhaps at a higher stage of development one of the most elementary ethical
demands. For the savage it may be a meritorious action to spare the vanquished and disarmed enemy; but it
belongs to the elementary commandments in the ethics of civilized nations.

\chapterhead{V.}{THE FREEDOM OF THE WILL.}

1. Ethical valuation arises in the first instance as an outburst of feeling. But it has its abiding
significance through its
% --- p. 107 ---
\opage{107}motivating, will-determining force. It thereby itself becomes a contributing cause of development
toward the highest welfare. It is of great importance to hold firmly to this, that conscience is itself
a cause. For one thereby takes up the right standpoint toward the so often discussed question whether
ethics can consent to acknowledge the validity of the law of causality for the human life of the will.
One will then be able to see that ethics not only need not demand any restriction of the validity of
the law of causality, but that such a restriction would be harmful to ethics itself, indeed would even
make it impossible.

It is, after all, a frequent assertion that all true ethics stands and falls with the assumption that
the human will is not, or at any rate not through and through, subject to the law of causality. There
can, it is held, be question of moral responsibility only if man can begin a causal series absolutely
afresh, i.e.\ can be a cause without being an effect. It is such an absolute beginning that is meant
by the infelicitous expression ``freedom of the will.'' This expression is infelicitous already for the
reason that it can be, and is, used in a whole series of different meanings, which are often confused
with one another.

2. Applied to the human will the word ``freedom''\footnote[1]{This word is now used predominantly in a
negative sense, so that as soon as one hears it one must ask: ``Freedom from what?'' This negative
sense is, however, hardly the original one. \emph{Steinthal} \textit{(Allgemeine Ethik.} Berlin 1885. p.
356 f.) remarks of the word ``frei'': ``It comes from a root that means to tend, to spare, to show
affection, to love and to make oneself loved. Free, then, is dear.'' The word is related to words such
as Freund, Freude, Friede.} is used in at least 6 (six) different meanings.

a. The meaning of the word freedom with which alone the debate about ``freedom of the will'' is
concerned is the one according to which a ``free'' will is not subject to the law of causality, is not
a member of the causal series like other phenomena, but is a cause without being an effect. One might
call freedom in this meaning of the word ``freedom from causation.'' It is about this that the dispute
stands, determinism denying it, indeterminism assuming it. What
% --- p. 108 ---
\opage{108}one is free from, then, is here: everything. To will freely is to will \textit{without cause,}
independently of everything that has gone before.

b. Freedom can next mean merely \textit{the absence of external compulsion.} Here, then, not all causes
are excluded, but only those that lie outside the willing personality. He is free whose resolve is not
hindered by external force from passing over into action. It is thus freedom of action rather than
freedom of will that is in question here. I have freedom to go out of my room when I have the key in my
pocket and the door is not barred; in the opposite case I am unfree and must stay where I am.

c. Freedom can also be a \textit{freedom from internal compulsion.} One often calls unfree the willing
that springs from pain or from fear, as opposed to the willing that springs from pleasure or hope. Even
if the door stands open, I perhaps do not go out, because I fear being attacked by a lurking enemy. In
such cases I do not do what I have a mind to; I feel a hindrance that makes me not act with the whole
and full assent of my mind. ``You build yourself a house,'' says Augustine, ``because otherwise you
would have no dwelling. It is then necessity, not your free will, that makes you will.'' The ``free''
will in this sense we often call in daily speech our ``good'' will. That I do not do something with my
``good'' will means that I do it reluctantly, do it, but with a strong feeling of pain, do it because I
prefer a lesser evil to a greater evil. He who hands over his money to the robber who threatens his life
does it deliberately and by a conscious act of will; but he can say: ``I have willed it --- but under
compulsion'' (coactus volui, as an old Roman jurist already aptly expressed it)\footnote[1]{Cf.
\emph{Jhering}: \textit{Der Zweck im Recht.} I. 2. Aufl. p. 17.}.

d. In the fourth place the ``freedom'' of the will can mean the will's \textit{ability, strength and
competence.} The question here is \textit{how much} the will can accomplish, not whether it is dependent
on or independent of what has gone before. One can be an indeterminist, that is, hold that the will is
not conditioned by anything preceding, and yet assume that this causeless will plays
% --- p. 109 ---
\opage{109}a very small part in the world. And one can be a determinist, that is, hold that the will is through
and through conditioned by what goes before, and yet assume that this causally determined will plays an
extraordinarily great part in the world. In the famous dispute between \textsc{Augustine} and
\textsc{Pelagius} the question was rather of the strength and competence of the will (the natural human
will)\footnote[1]{Augustine maintained that the human will is at variance with itself and is not able
to let itself be lifted up by the truth. ``The soul commands itself to will, and it would not command it
if it did not already will! And yet what it commands does not happen; but neither does it will with its
whole being (ex toto), and therefore it does not command with its whole being either\dots{}. If the will
were complete (plena), it would not command that it should be, for then it would already be! . . . . It
is a sickness of the soul; for the soul does not rise in its wholeness: it is supported by truth, to be
sure, but it is weighed down by habit. There are, then, two wills, since neither of them is complete;
what the one has, the other lacks\dots{}. It was I myself who willed, and it was I myself who did not
will, and I was torn asunder in my inmost being.'' Confessiones. VIII, 9---10.} than of the will's
relation to the proposition that everything has its cause. The question was, after all, to decide
whether man needed supernatural help in order to act well, --- that is, whether human nature had
\textit{sufficient} motives at its disposal, not whether motives were needed \textit{at all}. The same
point of dispute was taken up again, as is well known, at the time of the Reformation. In the Augsburg
Confession (I, 18), accordingly, the words strength or ability (vis) and freedom (libertas) are used as
synonymous. --- How easily one can confuse this meaning of the word with the first-named (freedom from
causation) is seen from the use indeterminists frequently make of the word ``ability.'' They then speak
of freedom as the ability to begin absolutely afresh. But if they distinguish between the actual
beginning itself and the ability to begin, then the will, the actual willing, becomes dependent after
all, namely on its ``ability.'' According to the strict indeterministic line of thought nothing may go
before the ``free'' willing, not even the ability. For if any meaning is to be attached to the word
ability, it means the conditions in our nature that must be present for us to be able to perform a
certain activity. Ability to will can
% --- p. 110 ---
\opage{110}reasonably mean, after all, only the inner conditions for a willing to be able to occur.

e. Very frequently by the ``freedom'' of the will one understands \textit{freedom of choice,} the ability
to choose. A choice by no means presupposes the invalidity of the law of causality. It presupposes only
that one has ideas of different possible actions which one deliberates, that is to say,
\textit{compares} with one another. The decision as to which of the actions is to be performed then
depends not on momentary prompting or isolated affects, but is determined through an inner debate between
a multiplicity of ideas and feelings. ``Freedom'' here means not the opposite of necessity, but
\textit{the opposite of blindness.} Through freedom of choice a deeper, more composite necessity shows
itself than in the action called forth by momentary affects and impulses. The ``free'' will here means
the mature, self-conscious will, which yet in the moment of choice says: ``I can do no other!'' To the
external spectator it may look as if man at such a moment could ``just as well'' have willed the
opposite, --- as if, for example, Luther at Worms could ``just as well'' have pronounced a retraction.
But for one who knew Luther's character and everything that moved within him, it must have stood as an
impossibility that he could have willed otherwise than he did. The external spectator does not know the
inner conditions that turn the scale. And it is in a purely external way that the popular conception as a
rule regards the will.

He who knows only one possibility cannot choose, ``has no choice.'' The cause of his knowing only one
possibility may be very various. It may be limited experience, neglect to use the experiences he has had,
momentary indisposition, passion, enfeeblement, or downright insanity. By such causes it can also be
explained when the different possibilities do not present themselves in deliberation with the degree of
clearness they would otherwise have. But all this precedes the choice. Whether a choice can take place
depends, then, on what has gone before. Freedom of choice does not conflict with determinism, and it is
not about it that the dispute stands.

f. Finally the word freedom can designate the will as the will that
% --- p. 111 ---
\opage{111}is guided by ethical motives. In this meaning of the word only the good man is free. What is
presupposed here is such a spiritual development and practice that conscience can come to play a
decisive part in every deliberation and resolve. But this obviously presupposes in its turn that there
exists a psychological causal connection. On a call or occasion for action, conscience must then --- by
virtue of the laws of the connection of ideas with one another and with the feelings --- be capable of
being awakened. Freedom here means that certain thoughts and feelings have become dominant and have
displaced others. Freedom here stands \textit{in opposition to bondage under the dominion of sensuous and
egoistic impulses and passions.} This is sometimes called the true or the higher freedom. This meaning of
the word is the oldest; it derives from Socrates (or Antisthenes?) and is used by Augustine, Spinoza and
others. It has nothing whatever to do with the dispute between determinism and indeterminism.

It is only the first-named meaning that we have use for in this connection, however significant the
others may also be in other respects.

3. It is now granted, I suppose, by most that it is not theoretical but practical-ethical motives that
lead to the assertion that our will is not subject to the law of causality. If one regarded the will
merely as a psychologist or historian, one would hardly hit upon setting up such an assertion. There
would always be many utterances of will for which one could find no explanation; but this would contain
no ground for asserting that they had no cause. On the other hand it is held that such an assertion is a
necessary presupposition of ethics. Should this really be the case? Should we be obliged to acknowledge
such a disharmony between our intellectual and our ethical nature that, in order not to deny the validity
of the ethical, we must deny the principle by virtue of which alone existence can become intelligible to
us? --- We must at least look carefully before we enter upon so desperate a conception. It is not
everyone who finds it so very easy to push aside the postulate of causality and set up other
% --- p. 112 ---
\opage{112}postulates in its place, much as one takes off one's everyday coat to put on an evening dress.

I shall then try to show the untenability of indeterminism\footnote[1]{In our own literature
indeterminism has been defended by \emph{Heegaard} \textit{(Om Intolerance.} 1878), \emph{Wilkens}
\textit{(Sociologi.} 1882), \emph{Kroman} \textit{(Vor Naturerkendelse.} 1883), H. \emph{Scharling}
\textit{(Kristelig Sædelære.} 1884). --- When one (like the last author) holds that there is a ``higher
mediation'' of determinism and indeterminism, I am unable to discuss this question, simply because I
cannot attach any meaning to it. \textit{Either} the law of causality holds for the will, \textit{or} it
does not hold. What third can there be? --- On Prof. Kroman's standpoint on this problem in the 2nd
edition of his Tænke- og \textit{Sjælelære} I must refer to my \textit{Psykologiske Undersøgelser.}
(Vidensk. Selsk. Skrifter 6. Række) p. 99---103. In his most recent statement on this question Prof.
Kroman finds the possibility of a partly causeless willing by distinguishing between the central and the
peripheral in man (or between man's ``essence'' and his ``character''). \textit{Begrebet ``det}
Etiske''. (University programme. Nov. 1903.) p. 87. But here the problem comes again: Is the central or
the ``essence'' always the same? Did ``I myself'' produce it in the first place? Can it be released into
activity equally easily at every moment? And is it always good if it is released?} and to show that
determinism is even an indispensable presupposition of ethics.

a. Indeterminism tears apart the bond between the individual and the race, indeed between the
individual and the whole of the rest of existence. The individual no longer stands as a peculiar member
of the great connection of existence, but is torn out of it precisely at the most decisive points. It
therefore becomes impossible for indeterminism to conceive existence as a whole. Every deeper-going
philosophical or religious view becomes impossible. The only religious view compatible with indeterminism
is polytheism; for every being that can absolutely begin a causal series is a little god (or a little
devil), an absolute being, and we thus get as many gods as there are ``free'' men. Perhaps one does not
care so very much about such a comprehensive view either. But the consideration put forward here has its
significance nevertheless, especially when determinism is attacked as a godless or antireligious
doctrine.
% --- p. 113 ---
\opage{113}If one conceives the Godhead as an absolute and almighty being, the assumption of a causeless will in
finite beings becomes downright self-contradictory. And if one protests against this and maintains that
we here stand before a ``mystery,'' then one must see to it how one can distinguish between mystery and
self-contradiction\footnote[1]{It is hardly historically correct when Kroman \textit{(Vor
Naturerkendelse.} p. 256) holds that antireligious feeling has as a rule taken the side of determinism.
One might rather assert the opposite; at least the rationalistic tendencies have usually taken the side
of indeterminism (cf. Pelagius against Augustine, Erasmus against Luther, the Arminians against the
orthodox Calvinists). When Kroman himself declares that a mathematical world-formula is impossible if man
has free will (p. 257), he must also grant that it is self-contradictory at once to believe in an
omniscient God and to assume the will's freedom from causation. For the omniscient God must be summus
mathematicus and be able to foretell everything with absolute exactness.}.

b. Indeterminists commit a significant confusion of two different meanings of the word freedom when they
express themselves now so that the will has \textit{no cause}, now so that we ourselves are the cause of
our acts of will. The second, fourth and fifth meanings are here mixed up with the first. For when ``we
ourselves'' are the cause of our acts of will, our will is not free from causation, but has its cause in
our nature. ``We ourselves'' are at every moment not something ``abstract,'' but something quite
definite, namely the definite thoughts and feelings, dispositions, instincts and impulses that stir in us;
in them, then, the origin of the acts of will must be sought. The two concepts, self-determination and
freedom from causation, which are often regarded as identical, in reality cancel one another as soon as
one attaches a definite meaning to the word ``self.''

The indeterminist perhaps answers to this: ``I mean precisely that in the causal series we must stop at
man himself. If we go further, and let his will [and nature?] too have a cause, then the last cause of his
resolve comes to lie far beyond himself; therefore he himself is not the cause.'' To this one need only
answer that in this way there would be no causes at all in the world. It would then not be the lightning
that
% --- p. 114 ---
\opage{114}killed the man, but the causes (the electricity of the air, etc.) that called forth the lightning; but
these again would not really be causes of the lightning --- and so on in an infinite series. This
consideration has its interest; but why does one apply it only to the will? --- Every thing and every
being in the world is both effect and cause. The more peculiar and the more richly endowed a being is,
the more conditions it presupposes, and the more effects can proceed from it. It will appear that ethics
can very well, indeed that it must, adopt this way of looking at things.

When self-determination is set in particular opposition to inertia, it must be well noted that precisely
the most significant acts of will of all are performed in enthusiasm, in devotion, in absorption in great
ideas, and in immersion in thoughts that the individual has not itself consciously produced, at any rate
not at the very moment when the act of will takes place. It is, after all, never the naked and abstract
self that wills something. The self that wills has been torn out of its state of equilibrium by something
that finds points of attachment in it, to be sure, but of whose working the self is not the only cause.

c. If the act of will, or a part of it, is not subject to the law of causality, it stands as something
\textit{isolated and accidental} in relation to the whole personality. It does not become flesh of its
flesh, blood of its blood. It is not only existence that indeterminism cannot conceive as a whole; in the
single personality either the coherence cannot be preserved. It is a strange fate that indeterminism here
suffers. It thinks it is asserting the dignity of man against determinism, which according to its
assertion makes man a mere machine, --- and then it comes itself to make man into something far meaner
than a machine, into something incoherent and accidental. When one and the same motive, without any
change having taken place in the inner or outer circumstances under which one acts, is followed now by
one resolve, now by another, in what then does this ``freedom'' really differ from capriciousness and
chance? What value will one be able to attribute to it at all? And how can one, under such
% --- p. 115 ---
\opage{115}presuppositions, go on speaking of a man's having a certain character?

d. The ethical judgment on my conduct builds on the action's really being my action. Therefore the
ethical judgment is clear and definite only where the psychological connection between motives and
resolve lies plain to view. The less my action can be understood from knowledge of my character and my
circumstances, the sooner shall I be regarded as not accountable, and the less shall I be able to hold
myself responsible. By giving up the causal connection in the life of the will I give up precisely the
mark of accountability!

The juridical way of looking at things agrees here with the ethical. The penal code regards agitated state
of mind and lack of premeditation as extenuating circumstances, repetition, relapse and clear
premeditation as aggravating circumstances. At the basis of this way of looking at things lies the
conception that a criminal act of will must be punished the more severely, the greater the connection in
which it proves to stand with the man's whole character. Although the juridical judgment is concerned
essentially with the external action, not with the inner disposition, it is nevertheless necessary for
the judge to gain insight into the motives and origin of the action. And that not only in order to be
able to determine the degree of the punishment, but also in order to be able to have a completely sure
conviction that the action has really been committed. The better the judge understands how the action
sprang necessarily from the inner and outer circumstances at hand, the more sure he can be that the
individual has really committed the action\footnote[1]{Cf. my \textit{Psykologi} VII B, 4. ---
\emph{Anselm} v. \emph{Feuerbach} \textit{(Aktenmässige Darstellung merkwürdiger Verbrechen.} Giessen
1829. II. p. 502 f.) maintains that a judge ought not to pass a sentence of death when he has not
understood \textit{how} the action came into being.}.

e. When several recent writers (Heegaard, Kroman) hold that the postulate of freedom from causation can be
restricted to the least possible, to ``a very small quantity, a trifle,'' it is not easy to understand
what significance it can have that so small a part of the act of will is free from causation, since it
is, after all, \textit{the act of will as a whole} that the ethical judgment concerns.
One sees that
% --- p. 116 ---
\opage{116}indeterminism, which first dissolves the coherence of existence and next the coherence of the
personality, ends by not even conceiving the single act of will as a whole. Certain per cent of its
elements are to be causally determined, others free from causation. We are not told whether the
causeless component is equally large in all acts of will of the same individual and in acts of will of
different individuals, nor how he who holds himself responsible, let alone he who is to judge others,
can discover how large the component is, and whether it is there in the single case. But even apart
from this the assumption stands as something most strange.

If, as has been said\footnote[1]{Kroman: \textit{Vor Naturerkendelse.} p. 246.}, it is ``the Achilles'
heel of determinism'' that the individual after every base action would have the right to say: ``It was
necessary, therefore it is not my fault!'' --- then it is not easy to see how the indeterminism that
assumes only an insignificant part of the action to be free from causation (and which to that extent
can also be called, and calls itself, a relative determinism) can avoid being struck by the same
difficulty, if indeed it is present at all. The individual will then rightly say: ``Only a
\textit{thousandth} of my action am I to blame for, since only so much of it is free from causation! Why,
then, am I held responsible for the \textit{whole} action? I have committed no murder, but only the
thousandth part of a murder!'' --- The most consistent thing, however, would be to admit the
impossibility of determining how far a man was accountable for a single definite action. This
consequence \textsc{Kant} draws in one passage, where he declares that ``how \textit{large a part} of our
character is the pure effect of freedom, and how large a part can be ascribed to mere nature and to the
unmerited faults or fortunate constitution of temperament, no one can fathom''; therefore ``the real
morality of actions remains wholly hidden from us.'' \textit{(Kritik der reinen}
Vernunft\textsuperscript{1} p. 551). ---

That modest indeterminism is prevented precisely by its modesty from attaining what it wishes to attain,
and it attempts
% --- p. 117 ---
\opage{117}the impossible thing of making a weight of a hundredweight hang from a cobweb. If responsibility really
stands and falls with freedom from causation, then we must assume that the agent could just as well have
omitted the \textit{whole} act of will, the whole action.

That kind of modesty usually enters into the history of theories at the point where their vitality and
self-confidence have vanished, and where the enemy has already been let into one's own camp without one's
being conscious of it. The old assertions then hold on for a time still, like rudimentary organs; but it
is an illusion to think that they really perform any function.

f. The words \textit{responsibility, guilt} and \textit{accountability} are, like so many other ethical
expressions, carried over from the juridical to the ethical field, or rather, derive from a time when the
difference between these two fields did not yet properly make itself felt. That I am held responsible or
regarded as accountable for an action means that I stand as the one to whom reward or punishment for the
action is allotted. \textit{Why} the action is rewarded or punished is a question by itself. It is, as
already shown (IV, 4), not at all self-evident that an action is to be rewarded or punished. But even if
we accept the theory of retribution and regard it as self-evident, it is not at all necessary that the
will held responsible should be the \textit{absolutely first} cause, should begin a causal series
absolutely afresh. Retribution demands only that there be \textit{someone to lay hold of,} and there is,
once the origin of the action in a human will has been established. The primitive instinct of retribution
does not even go so far back; it is satisfied if only a man of the same family or tribe as the
perpetrator suffers for the action, and it demands satisfaction whether there is present a really
\textit{willed} action or not. The idea of individual guilt and individual accountability has slowly worked
its way forward. But the demand that the individual will shall have no cause derives from a metaphysical
speculation into which neither the primitive instinct nor the ethical feeling of guilt enters.

Why, then, do we not in our ethical feeling of guilt go further
% --- p. 118 ---
\opage{118}back than to the will? Why does ethics stop here and leave everything that goes before to psychology?
--- This question is here to be answered from a purely ethical standpoint. Juridical responsibility can
first be treated when, in the doctrine of the state, we come to the grounding of the punitive authority
of the state.

Ethically regarded, the feeling of guilt, responsibility or accountability can mean only that I feel my
action subject to the judgment of conscience. The inner sanction (IV, 4) comes into force. The
presupposition of this is that I am conscious of having really willed the action. I impute it to myself
because it is nothing foreign to me, but would not be in the world without my resolve. I bear the guilt
because I have willed. When I now perceive the conflict between the act of will I own and what I
recognize as right, there arises \textit{the feeling of remorse,} a feeling of inner disharmony,
unworthiness and self-contempt, which can rise to the highest spiritual pain. It arises as a feeling of
contrast, called forth by the opposition between my actual willing and my acknowledgment of the ideal.
This feeling arises under the conditions indicated with psychological necessity.

From a deterministic standpoint too one may very well define remorse as ``dissatisfaction with oneself at
not having acted otherwise''\footnote[1]{\emph{Kroman}: \textit{Vor Naturerkendelse.} p. 243.}. Such
dissatisfaction is only another form of the lively wish, arising in the moment of remorse, that one had
after all acted otherwise, a wish that, as I shall show directly, has its great practical significance.
Such a wish arises naturally when in \textit{the moment of reflection} one submits one's willing and acting
to a serious test and finds that it cannot stand it. But from this wish and the dissatisfaction it
awakens one cannot infer that in \textit{the moment of action} one could just as well have acted quite
otherwise than one actually did. In many perhaps this illusion arises; but they then merely date the
dearly bought experience they have gained through error and remorse back to an earlier
moment\footnote[2]{Psykologi VII B, 5 b (conclusion). --- A more detailed demonstration of the ways in
which the idea of a causeless willing can arise I have given in my \textit{Psykologiske Undersøgelser.} p.
92---99.}.

% --- p. 119 ---
\opage{119}But with the psychological explanation of remorse its ethical value and grounding are not yet given.
Remorse is a feeling of pain, and it is only from an absolutely ascetic standpoint that pain is in and for
itself something good. Why should I not seek as quickly as possible to forget my reprehensible action,
just as I naturally seek to keep unpleasant memories away? When Hamlet has awakened his mother's
conscience, she says: ``O Hamlet, thou hast cleft my heart in twain!'' And he answers:
% Shakespeare's English (Hamlet III, 4), where Høffding quotes a Danish translation:
%   „O Hamlet, du har søndersprængt mit Hjerte!“ / Saa kast den slette Halvdel bort, og lev /
%   des renere da med den anden Halvdel!

\begin{quote}
O, throw away the worser part of it,\\
And live the purer with the other half!
\end{quote}

Why does ethics hold that it cannot go so briskly as Hamlet thinks, but that the pain of remorse has its
great significance?

It is not because it holds the principle ``like for like,'' or ``an eye for an eye, a tooth for a
tooth.'' The doctrine of retribution is no ethical doctrine. Ethics cannot see in remorse ``God's
terrible and unspeakable wrath,'' as \textsc{Melanchthon} calls it, nor, with the Catholics, see a penance
or atonement in it. --- A more thorough criticism of the doctrine of retribution than is contained in the
remarks already put forward will find a better occasion in the inquiry into the punitive authority of the
state. Here I shall confine myself to a single point that is of particular interest in this connection.
--- According to the doctrine of retribution, remorse (where it is the only punishment) would have to be
strongest in him who has committed the greatest crime. But in fact it is not, and this finds its
psychological explanation in the fact that remorse rests on a relation of contrast between the ideal and
the actual willing, a relation of contrast that of course cannot come into being except where the idea of
the ideal becomes living. The more living and developed this idea is, the more actions will be regarded in
its light, and the more sharply they will be judged. Therefore the feeling of remorse is often strongest
in the purest and best characters. There is a psychological truth in the tale of Queen Dagmar, who found
no rest in death because she had laced her sleeves on a Sunday. Against a snow-white background spots are
seen that would otherwise not be noticed.

Psychologically regarded, remorse appears --- like duty, seen from another
% --- p. 120 ---
\opage{120}side (III, 9) --- as the expression of the personality's striving to preserve its unity and
coherence. Remorse presupposes that consciousness of the good has been awakened, which it perhaps was not
at the moment of action, or else that it has been heightened in kind or range or both, a higher ethical
stage having been reached. Pity is felt for the victim of the action, and one feels the value of the
interests that have been injured. Remorse is thus already a sign that the good is about to prevail in man.
But the pain of remorse is caused by the consciousness of the actual willing and acting standing in the
sharpest conflict with the consciousness of the ideal. Therefore \textsc{Dante} says
\textit{(Purgatorio} XXXI) that Beatrice, i.e.\ the ideal thought, was the cause of the pain that ``the
nettles of remorse'' caused him. In spite of this painful inner opposition, which shows man his own soul
in an entirely different light from before, he strives to preserve the connection with his past. He can
neither let go of the new I nor deny the old I: both together make up his own I. He recognizes --- in
spite of the new I --- honestly and without reserve himself as the cause of the earlier acting. Through
this constant striving to identify himself with himself and to make the old and the new I cover one
another by trying to think them together, the contrast is felt all the more strongly. Only in the form of
pain can the personality here assert its unity. But the attempts at covering may lead to a purification,
a fusion, whereby the new I prevails, but is at the same time enriched by the experience man has had of
his limit. The consciousness of the ideal becomes all the more inward as the oppositions of life have been
tried at a depth to which those whose development runs more in a straight and unbroken line do not reach
down\footnote[1]{This development is illustrated by the laws of the connection and fusion of feelings
developed in my \textit{Psykologi} VI B, 2 d and e, and might itself have been cited there as an
example.}. What is acquired through such a purgatory has a firmness that is not easily won in any other
way. The energy that blindness formerly took into its service may now, through a metamorphosis, come to
work in an entirely different
% --- p. 121 ---
\opage{121}direction. In this way remorse bears witness to the causality of the personality and to the
conservation of psychic energy. Already \textsc{Bröchner} has (in his work \textit{Det Religiøse i dets
Enhed med det Humane.} 1869. p. 157. 172) expressed this conception of remorse: ``The preceding life is
joined together through remorse with the present into a continuous, intensive whole of life.'' ``In the
transformation of force, here as in the domain of the physical, the sum of force is abiding.''\footnote[1]{Cf.
already \emph{Giordano Bruno} \textit{(Den nyere Filosofis} Historie\textsuperscript{2} I. p. 141 f.).}

Only in the express consciousness of the character of the earlier state does one have the sign that one
is beyond it. If the will to bear the opposition between the old and the new slackens, the metamorphosis
and the firm founding of the new character do not take place. He who cannot bear to think of his earlier
conduct, --- he who cannot pass through the purgatory of self-knowledge, --- is still entangled in his
past, just as the insane man is not fully cured if he will not admit that he has been insane. There is of
course question here only of the judgment the single person holds within himself; it is not said that
anyone else has the right to settle his accounts for him. --- Where the boundary lies between the
necessary self-knowledge and morbid brooding over what cannot be otherwise belongs to the very hardest
questions in life; no general formula can instruct us about it. Melancholy natures easily come --- by
virtue of the expansion of feeling\footnote[2]{Psykologi VI F, 4.} --- to regard themselves as more
guilty than they are. Activity is thereby paralyzed, and the impulse to progress is cowed\footnote[3]{There
is often an egoistic sentimentality in brooding remorse and petty self-accusation. \emph{Goethe} tells (in
one of his conversations with Eckermann) of a little boy who could not be calmed over a little fault he
had committed. ``It was disagreeable to me,'' says Goethe, ``to notice this; for it testifies to an all
too tender conscience, which values its own self so highly that it will not forgive it anything. Such a
conscience produces hypochondriac men if it is not counterbalanced by a great activity.'' In
\textit{Wilhelm Meisters Lehrjahre} (IV, 12) it is said: ``Die Eigenliebe lässt uns sowohl unsere
Tugenden als unsere Fehler viel bedeutender, als sie sind, erscheinen.'' --- Cf. also my remarks in
Rousseau og hans Filosofi, p. 13 f. ---}.

% --- p. 122 ---
\opage{122}For the ethical significance of the penetrating consciousness of the reprehensibility of the action,
which the pain of remorse partly presupposes, partly awakens, is that of acting as a spur and driving
forward. An attention is called forth for things one formerly passed by blind and hard. There rises an
impulse to reach a higher stage, a craving that springs naturally from the painful contradiction that is
felt. Only by being a \textit{motive} is remorse of an ethical nature. It is a means for future
development. Ethics looks forward, not back, or rather, it looks back only in order to be able to look
forward all the better. Ethical judgments are, as has been remarked several times in the foregoing, in the
first instance outbursts of feeling at the sight or memory of an action. This they continue to be, however
well they are grounded. But the pronouncing of such a judgment is justified only where it can be a cause,
a motive to continued or changed willing. Neither toward ourselves nor toward others have we the right to
judge when the judgment cannot act as a motive. A great part of the difficulties many find in reconciling
ethics and determinism derives from the Old Testament way in which they conceive concepts such as
responsibility and accountability. Only he ought to be held responsible in whom the settling of accounts
can become a motive --- a deterring or an encouraging motive. If the popular conception means something
else by the words remorse, responsibility, accountability, etc., then they are untenable concepts that it
has formed.

In the wish to have acted otherwise and in the pain of not having done so one perhaps acknowledges
something that in its time made itself felt during deliberation, but could not then gain power. We see
here the great significance of the inner debate before the resolve. Even if the impulses and promptings
that had right on their side had to succumb, this struggle has perhaps not been in vain. They may have
made the equilibrium less stable, even if they could not wholly
% --- p. 123 ---
\opage{123}overcome the resistance to movement in the direction in which they point. But the work they have done
need not go to waste. They rise again when the mind has become calmer, and by contrast they may now perhaps
gain the power they formerly lacked, or they get help and supplementing from new promptings and impulses
that on their side would not alone be able to determine the will.

g. If the law of causality did not hold in the spiritual domain, ethical striving would be hopeless. Only
where conformity to law prevails can my will intervene in such a way that something is accomplished. That
we can intervene in external nature and make it obey our ends is made possible by our knowing the laws of
nature and knowing what conditions we must procure in order to attain what we will. We stand in a similar
way toward human nature. Only if conformity to law prevails can we change it at the points where it
conflicts with our ideals. What matters is to procure motives of the right kind and strength, and this must
happen by virtue of the psychological laws of nature. But what would every possible effort avail if one
and the same motive under the same circumstances were followed now by one resolve, now by quite another?
\textit{What has no cause, I cannot prepare;} toward it I stand as toward the accidental and capricious. I
have command of my future willing only if there is a causal relation between my present and my future
willing. Then what I now sow as a seed may grow into a vigorous plant.

We now understand why responsibility does not go further back in the causal series than to the will. It is
because it is the will that is to be changed. What lies before the act of will interests us, ethically
regarded, only insofar as it has influence on the will. Only indirectly, therefore, are we held responsible
for our thoughts and feelings, namely insofar as they are will-determining forces. The extent to which they
are to be drawn in will of course be very different in the different cases.

Out of nothing comes nothing. This proposition ethics must for its own sake let hold. What matters is to
awaken and nourish wishes and impulses of the right kind in myself and others. How
% --- p. 124 ---
\opage{124}far the single person can get in the single case, only experience and trial can convince him.

h. A few words should still be given to throwing light on the strange assertion that the determinist must
consistently stand as a mere spectator of life, both of his own and of others' life. As if one cannot feel
joy or sorrow at life and desire to intervene in it because one assumes that it is subject to definite
laws! If my responsibility gives me no time at all to think, I cannot of course be a determinist. But then
I cannot be an indeterminist either. For one becomes that, too, only by thinking. Indeterminism is just as
much a theory as determinism, and there are those for whom it demands the greatest intellectual effort to
make out what indeterminism really means. If I am at every moment absorbed in practice, I cannot
theorize. The tiller of the soil gets no time to study chemistry when he is to plow and harrow without
interruption. It is not, of course, everyone whose lot it is to study psychology and history. There may be
natures in whom the constant interest in finding the causes of their own and others' actions can weaken
the practical interest. But then it is only the general opposition between theoretical and practical
natures that we stand before. Everyone must here see to how much room he gives theory in his spiritual
household. It is of course unethical to appear as a mere inquirer where action is called for. Who, however
zealous a physicist he may be, would study the nature of the effects of light while he stood at the
deathbed of one of his dear ones? But will one therefore say that physics and love exclude one another? I
cannot (whether I am an indeterminist or a determinist) study the psychology of the will at the same moment
as I am to act; but because I cannot do two different things at the same time, these two things need not
therefore logically conflict with one another. I cannot at one and the same time stand on my feet and on my
head; but I can (perhaps) do the one first, the other later.

One presupposition, to be sure, lies at the basis of the whole ethical consideration, whether it is
deterministic or indeterministic, namely that existence is not completed,
% --- p. 125 ---
\opage{125}finished or perfect, but that there is a work to be done for human willing. The human will must be one
of the causes on which the development of the world depends. What weight it has in relation to other causes
only experience can teach; but it must not be nil if an ethics is to be possible. Every ethics that stands on
the ground of reality is a doctrine of how tendencies of development that express themselves in more
elementary forms in the external life of nature can be carried further in the human mind. The course of the
world must in part take place through human willing, and every attempt at a conclusive world-view must also
take this factor into account\footnote[1]{Cf. below VII, 1 and 4 and XXXI, 3, as well as my
\textit{Filosofiske Problemer} (University programme Nov. 1902). III, 5 and IV, 2.}.

\chapterhead{VI.}{ETHICAL EVIL.}

1. What ethical evil is follows already from what has been developed earlier (III, 6; 10; 20). It
consists in a more or less conscious isolation of the single moment from the whole of the individual's
life, and of the single individual from the whole of the life of the race. From this there follows not
merely a resistance to what the welfare of the individual or of the race requires, but also a slackening
of energy and a dissolution of coherence in the individual or in the race. --- We shall here go a little
more closely into the psychological causes of this isolation.

Two causes explain the tragic disproportion that is designated by the concept of evil. First,
\textit{the sporadic character of development,} which consists in this, that the single faculties in the
single personality and the single individuals within the race do not develop simultaneously, nor equally
strongly and quickly, but in such a way that now the one element, now the other, comes forward with the
% --- p. 126 ---
\opage{126}greatest energy; thereby the possibility arises that they do not come to fit harmoniously into the
larger whole, but even offer definite resistance to it. This disharmony can arise in two forms. It may be
that the single element (the single faculty in the individual, the single individual in the race) carries
on with energy a tendency that has developed beforehand. By virtue of a kind of inertia it goes on
existing and working as before, although other demands must now be made on it from the side of the whole
than before. Or it may be that the single element involuntarily varies, strikes into new directions,
precisely because life stirs in it so vigorously that it must have a more ample outlet than before, ---
but that no possibility now appears of finding a connection into which the new might enter and gain
value. What at first seemed to make an entirely new stage of life possible now becomes a break in the
coherence of life, something twisted and distorted. --- Neither the conserving nor the varying tendency
can be done without; but both contain the possibility of a break with the coherence of life and of an
injury to the fundamental will that valuation presupposes. They are forces that involuntarily will the
good, but often work evil. They can lead to isolation, so that the part sets itself in the place of the
whole, because it will not change its state --- whether this be essentially continuation or constant
variation.

The second cause lies in this, that \textit{the involuntary precedes the voluntary,} the unconscious or
dimly conscious precedes the clearly conscious. The thought that can guide must first be won, and is
often won only through error. One often becomes wise only through harm. Zeus, says Aeschylus, has laid
down the law that wisdom is won only through suffering (πάθει μάθος). The individual weaves itself into
life without knowing into what connection it is entering. And the first thought itself arises
involuntarily: the transition from the involuntary to the voluntary must of course take place
involuntarily. The first ends man sets himself he sets without deliberation:

% --- p. 127 ---
\begin{quote}
\opage{127}How the first knowledge comes to be awakened\\
No man can tell, nor know whence it wells up,\\
Nor how his desires first came to life.\\
As the honey-making urge shows itself in a bee,\\
So came they to him; that \textit{first thought of the will}\\
Deserves no praise, and no blame either.\\
(Dante: \textit{Purgatorio} XVIII).
\end{quote}
% Høffding's Danish (Dante, Purgatorio XVIII, 55--60):
%   Hvordan den første Kundskab vakt kan blive, / Véd ingen Mand, vèd ej, hvorfra den vælder,
%   Og hvordan først hans Lyster kom til Live. / Som Honningdriften hos en Bi sig melder,
%   Kom de til ham; hin Villiens Førstetanke / Fortjener ingen Ros, og Last ei heller.

And yet this ``first thought'' can become decisive for the direction of the whole of life. It becomes the
point of departure of the life of thought, the first center of association, which determines the course
and kind of the following thoughts. But there is no certainty that the direction thereby given agrees with
the highest interests of life.

Isolation and blindness, then, are the two features that evil presents, when we regard it in its
possibility. Evil is the inertia that will not change its state, the heaviness that will not let itself be
lifted, the narrowness that will not let itself be widened. It is the full opposite of the motive of
valuation that was presupposed in the foregoing: for sympathy is a need to impart, to make oneself one
with others, to open oneself to the world. --- Greek ethics laid most weight on the blindness that awakens
arrogance and insolence (ὕβρις). \textsc{Homer} and \textsc{the tragedians} conceived the origin of sin as
infatuation (ἄτη), \textsc{Socrates} downright as ignorance. The ethics of modern times stresses most the
isolation proceeding from inertia. So \textsc{Spinoza} and J. G. \textsc{Fichte}. The latter thinker shows
in his interesting psychological development\footnote[1]{J. G. \emph{Fichte}: \textit{Sittenlehre.}
1798. p. 262 ff. See my book \textit{Den nyere Filosofis Historie}\textsuperscript{2}. II. p. 150 f. ---
For \emph{Spinoza} sin was a spiritual impotence \textit{(Tract. polit.} II, 7. 11. 20), which formed the
opposite of the way in which the impulse of self-assertion is able to carry man's development forward
from lower to higher forms \textit{(Ethica.} III, 6 ff.). See \textit{Den nyere Filosofis
Historie}\textsuperscript{2}. I. p. 325---327. --- F. C. \emph{Sibbern} has in his \textit{Psychologie}
(1856) p. 119 and in his \textit{Moralfilosofi} (1878) p. 7---11 applied his doctrine of the sporadic
course of all development to the problem of evil.} how out of inertia, as resistance to going beyond the
given state and direction, there can arise cowardice and falseness, since one would rather give up one's
freedom than take up the struggle, and since one, when there must nevertheless
% --- p. 128 ---
\opage{128}be a struggle, would rather use cunning than honest weapons. For both thinkers ethical evil is the
opposite of spiritual strength and freedom. This spiritual freedom was for them closely connected with
social life, and thereby evil became indirectly also a break in social harmony, an injury to the common
striving for freedom. But at the standpoint on which we have here taken our stand (III, 8---11), ethical
evil must be characterized directly by its antisocial character, by the absence of the motives and
qualities that make spiritual community possible, not merely of those that make spiritual freedom
possible.

When the consciousness of the opposition to a higher stage of life arises (which, however, one of course
does not acknowledge as higher), and the stage once entered upon is nevertheless held fast, inertia rises
to \textit{defiance.} The contrast here works to fix and to provoke, and increases the resistance. The
standpoint is shaped into a system that can be held fast with strict consistency. When such an isolating
stage is called evil, this valuation is made from another ethical system than the one acknowledged by him
who stands at that stage. Every standpoint asserts its right to be, sets its valuation against every other
valuation. Here, then, we have in the sharpest form the struggle between ethical systems that we earlier
(III, 13) considered from a purely theoretical point of view. It is the great struggle in the world, in
which what on one side is stamped as inertia, isolation, blindness and defiance is declared on the other
side to be the outcome of natural and justified self-assertion, or perhaps even the germs of new advances.
In this struggle it is not merely thought that stands against thought, but passion against passion, and
will against will. At the foundation of all ethical valuation there lies, after all, a fundamental will,
an activity that sets ends. Where valuation stands against valuation, it is the most fundamental forces of
spiritual life that have risen up against one another.

Philosophical ethics stands toward this struggle otherwise than theological ethics does. It learns from
psychology and history how standpoints and systems come into being, and how they successively unfold out
of elements that have their justified place in the connection of life. This successive unfolding of
% --- p. 129 ---
\opage{129}the standpoints that are rejected and are to be overcome, consistent theological ethics cannot
accept. It does not acknowledge the simple truth that the involuntary precedes the voluntary, and that the
parts do not always join together into a whole. Sin is for it \textit{pride,} --- just as obedience is the
sum of all virtue. When man isolates himself, tears himself loose from the source of all good, the cause
is to be sought in his arrogance. Whence arrogance and pride in turn arise we are not told; but from
Augustine\footnote[1]{\textit{De civitate dei.} XIV, 12---13. (Malæ voluntatis initium quæ potest esse
nisi superbia?) \textit{De moribus eccl. cath.} c. 12. --- See \textit{Om Grundlaget for den humane Etik.}
p. 137---141. --- Beside the strong stress on pride as the origin of the evil will there goes in Augustine
the strong emphasis on sensual desire (concupiscentia), which reveals itself especially in the sexual
impulse. Since this stirs independently of the will, it is clear that nature is corrupt, and since the
propagation of the race takes place precisely through the act of generation accomplished with sensual
desire, it is no wonder that mankind has become a chaos of sin (massa peccati). Cf. A. \emph{Harnack}:
\textit{Lehrbuch der} Dogmengeschichte.\textsuperscript{3} II. p. 29. But it is nevertheless Augustine's
express doctrine that it is the proud will that has corrupted nature, so that evil thereafter arises
involuntarily and rules the will. De civ. dei XIV, 3---6.} onward they stand at the head of all vices.
While the development of sin according to the philosophical conception proceeds from the unconscious and
involuntary to the conscious and voluntary, as experience shows it to us in all development of the soul,
it becomes the task of theological ethics to show how the less conscious and voluntary forms of evil
develop out of the most conscious and voluntary. In medieval ethics an attempt was made to explain the
seven capital sins genetically, so that out of pride there successively develop envy, hatred, sloth
(acedia), avarice, intemperance and lust\footnote[2]{Hugo a St. \emph{Victore} derives acedia from the
first three sins, which lead man to disgust with himself; out of acedia in turn spring the last three
sins, since man, having lost inner joy, turns outward to find comfort (\emph{Liebner}: \textit{Hugo von
St. Victor.} Leipzig 1832. p. 278---280). --- In \emph{Dante} the same course of development is hinted at:
\textit{Purgatorio.} XVII. Cf. \emph{Ozanam}: \textit{Dante et la philosophie} catholique \textit{au 13e}
siècle. Paris 1843. p. 94. --- \emph{Andreas Sunesøn}: \textit{Hexaëmeron.} ed. Gertz. p. 153 has the same order
(which seems to derive from St. Gregory). --- According to Thomas Aquinas all seven capital sins can be
the point of departure (radix); but since at the foundation of them all there lies a craving for one's own
eminence and loftiness, pride (superbia), the unbridled impulse to selfish loftiness (inordinatus amor
propriæ excellentiæ), may rightly be called the beginning of all sin, just as it is the queen of all sins.
Thomas sets up the following order, which, however, he does not regard as the only possible one: pride
(superbia), avarice (avaritia), lust (luxuria), envy (invidia), gluttony (gula), wrath (ira), sloth
(acedia). A parallelism between the capital sins and the cardinal virtues Thomas does not regard as
possible. Summa theologica. Pars II. Qvæstio 84, 2---4.}. Even if interesting observations are woven
into this account,
% --- p. 130 ---
\opage{130}it is nevertheless built on a psychological impossibility. F. C. \textsc{Sibbern} has rightly said that
theological ethics derives evil from the highest evil. That is to turn the matter upside down. Only
through a series of stages of development does unconscious inertia, which may be beneficial, rise to a
conscious rejection of every attempt to draw the individual out of its isolation. Only at the culmination
of the resistance begun by inertia does the individual make everything else a means to its enjoyment or
its feeling of power. Instead of being one among several, it has now attempted to make itself an absolute
end. In particulars it is so far from being the case that only one single series of development can be
run through here, that a great multiplicity of ways presents itself to psychological inquiry.

The more special forms of evil will depend on the general dispositions of feeling or temperaments that
the single individuals bring with them from the first. The dark temperament leads, in passive natures, to
dejection, despondency and unreceptiveness. The individual withdraws from everything that would lay claim
to it, especially also from enlivening and awakening influences. In active natures, given this
temperament, there arise malice and envy, harshness and cruelty. One feels others' pleasure and happiness
as a jarring disharmony, or regards it as the expression of stupidity and shortsightedness. There is an
inclination to be touched by everything in a painful way, and the frequency of the feelings of pain easily
produces
% --- p. 131 ---
\opage{131}egoism through the attention the individual is compelled to direct upon itself. Pain hinders it from
devotion, hinders it from forgetting itself. --- The bright temperament leads, in the form of passivity, to
a craving for enjoyment, to seeking satisfaction for the wishes and impulses rising at the moment, or to
wallowing in radiant fantasies and thereby losing the capacity to work at the tasks life sets. In active
natures it leads to reckless urge for activity and craving for independence, to lust for power or
fanaticism. Here too harshness and cruelty can develop, especially when one's own capacity to suffer pain
is slight. Great sacrifice and effort may often be shown here in order to satisfy the need to feel one's
powers, the egoistic lust for power. This becomes the heroic form of evil. --- Under such different forms
the individual entrenches itself in itself, makes itself the center and makes full devotion impossible.
Still more forms will be found when one takes account of the manifold displacements of motive that become
possible from the original disposition, and also of the way in which the general psychological elements
appear under the different historical and social conditions. Evil is not merely a psychological but also
a sociological phenomenon.

That evil can thus appear in many different forms is of great practical importance. When the one does not
have his imperfection where the other has his, mutual criticism and ethical influence can become possible.
The single person can become a model in certain respects, even if he cannot become so in all respects. If
such a diversity --- differentiation or ``division of labor'' --- did not take place, the possibilities of
development and progress would be far smaller than they are.

2. I shall go a little more closely into the question: how high a degree of consciousness can be present
in the states now depicted? --- --- The predicate ``evil'' is pronounced from a standpoint that he to whom
it is applied does not share. The man of the moment, after all, does not share the standpoint of the
individualist, and the individualist does not share the view of him who takes his stand at the standpoint
of the family, the tribe, the state or the whole race. If
% --- p. 132 ---
\opage{132}he who stands at the narrower or lower standpoint had complete and clear consciousness of the validity
of the predicate and of its applicability to himself, then the corresponding feelings and impulses must
also be set in motion, and his conduct would have to become different. There must always be ignorance and
blindness in evil: ignorance, insofar as the ideas are not clear and complete; blindness, insofar as it is
certain ruling feelings and passions that hinder the right development of the ideas. There may be an idea
of another stage of life, but this idea is heard only as a distant voice, gains no power, awakens no
strong feeling, does not enter into the real self. \textsc{Socrates} was in the main right that sin is
ignorance, if only by ignorance one thinks of more than a lack of enlightenment of the understanding, or
if to ignorance one adds blindness, in order to express that not only is right knowledge lacking, but
there are also positive hindrances to its being able to develop and gain power over the mind. Those who
act evilly ``know not what they do.'' The reason they do not know it lies, to be sure, deeper than
\textsc{Socrates} supposed, or at any rate deeper than he knew how to express. But the self-knowledge that
Socrates demanded would, if carried through, also bring to light the hidden foundation of the whole
standpoint. All ethics must therefore in the end agree with Socrates.

But surely a real act of will requires that one be clearly conscious of what one is doing! --- Here,
therefore, a distinction must be made between consciousness of the actual circumstances that the action
concerns and consciousness of the ethical character of the action. Ignorance and blindness need not
concern the actual circumstances. I can know definitely what I am doing, as, for example, the inquisitor
knows what he is doing when he passes the sentence of death on a man on account of his faith, without yet
being clear that, seen from a higher ethical standpoint, it is something reprehensible. Such clarity would
make the action impossible. For it is psychologically impossible to act against one's definite and full
conviction when this is not obscured and displaced by other promptings. One can ``have'' a conviction
% --- p. 133 ---
\opage{133}in the sense that there is a possibility of its being called forth and becoming living in
consciousness when there is sufficient call or occasion for it. But the question is whether it presses
forward in the mind spontaneously (or at any rate with a minimum of external occasion) at the moment of
deliberation and of action; for only then is it really a definite and full conviction. Here there are
infinitely many degrees and shades of clarity, and it is dogmatic to lay down the highest degree as given
without more ado.

One might think that such a willing against better knowledge is nevertheless found when man seeks to
obscure the ethical consciousness in himself in order to escape the discord and dividedness that arise
where it cannot immediately become dominant, or when man seeks to prevent the arising of thoughts that he
fears would change his whole present valuation, his present system\footnote[1]{Cf. \textit{Psykologi} VI F,
2.}. But in such cases the higher ethical consciousness has, after all, not gained a firm footing;
otherwise the painfulness of the divided state could not become the prevailing motive. Evil is not done
with full and clear consciousness that it is evil.

It is only in the most recent times that great poets have depicted evil with a finer and deeper psychology
than the one at the basis of the usual type of the devil, which even Shakespeare followed all too much in
his Richard III and his Iago. In the first rank here stands the beginning of Dostoevsky's Raskolnikov. But
I must also refer to a single speech in \textsc{Ibsen}'s ``Rosmersholm.'' The portrayal of Rebecca's
character I cannot undertake to defend; but one single speech is masterly. ``There are two kinds of will
in a human being, I should think! I wanted to have Beata away. In the one way or in the other. But I never
believed that it would come all the same. At every step that I ventured and risked forward, it seemed to
me as if something cried out within me: Now no further! Not a step further! --- And yet I \textit{could}
not let it be. I \textit{had} to venture one tiny little bit more. Only a single one. And then one more ---
and always one more. --- And then it came. --- That is
% --- p. 134 ---
\opage{134}the way such things happen.'' A lesser poet would have depicted her as a devil who with full
consciousness laid her plan and cleverly calculated the many small pinpricks that were to lead to the goal.

The theological idea of the ``sin against the Holy Ghost'' is likewise a psychological impossibility.
\textsc{Luther} describes it in a sermon on this subject as ``devil's sin, which, though manifest and
convinced, will not be convinced.'' He says: ``This sin I would never believe in former times, when I was
still a learned doctor, to exist in the world.'' Later he came to fear that the papists, who could not
refute the right doctrine, would stick fast in this sin! As a learned doctor \textsc{Luther} was a better
psychologist than when he had become the great agitator. If one will not be convinced, how can one then
\textit{be} convinced? So long as one will not be convinced, there must surely be a part of the mind left
where conviction has not penetrated with its full effect; how otherwise could such a will arise? ---
Luther says: ``When anyone goes so far that he will hear or see nothing, indeed will even defend his
blasphemy, then he is no longer to be counseled and helped.'' But he who defends himself appeals, after
all, to grounds; he thus still acknowledges the power of truth, only he does not acknowledge as truth what
his opponents call so. Here, then, something good is still left. The wholly devilish would be to
acknowledge no truth at all; but that would in the end lead to complete meaninglessness and
thoughtlessness and to complete impotence. Absolutely consistent wickedness would thus be absolute
stupidity, and it is a right instinct that led popular belief to transform the evil devil into the stupid
devil\footnote[1]{Already in the old Persian theology this feature appears: Ahriman in his ignorance does
not suspect that Ormuzd, the good power, exists, and when he has discovered this power, he underrates it,
likewise from ignorance. Edv. \emph{Lehmann}: \textit{Zarathustra.} En Bog om Persernes gamle Tro.
Copenhagen 1899---1902. II. p. 161.}.

3. A step further on than blindness lies
% --- p. 135 ---
\opage{135}obduracy. But it too is not one with conscious wickedness. Obduracy is blindness that has become
fixed, has passed into flesh and blood, so that a change of character is extraordinarily difficult. It
would, however, be unjustified to make such obduracy one with incorrigibility. Our art of education, both
the individual and the social, is hardly so far advanced that we may declare a man absolutely
incorrigible. Our art and our means are perhaps not sufficient. But in this field neither the strength nor
the skill has so far been put in that mankind has applied in so many other fields. One has been content to
put such individuals in a ``house of correction'' and to refer them to a hell in the other world. The art
and the great, patient feeling of humanity that alone could overcome the obstacles that may be present here
are still only in their becoming\footnote[1]{Cf. more fully on this my \textit{Etiske Undersøgelser} p.
73---81}.

\chapterhead{VII.}{THE THEORY OF WELFARE.}

1. In the three immediately preceding chapters we have been occupied with the \textit{foundation} of
ethics. We now pass on to a closer determination and illustration of the principle of the content of
ethics. To the foundation on which we have here taken our stand there answers, as shown earlier (III,
10---15), only one such principle, namely the principle of welfare. It gives the general answer to the
question: \textit{what is the good?}

From a misunderstood idealistic interest it has often been overlooked that the concept of the good is a
purely \textit{formal} concept, and a certain cult has been made of it without any closer determination of
it being given. The true good, it was held, must be the good in and
% --- p. 136 ---
\opage{136}for itself, whose value does not depend on anything else. But what is it that has its value in
itself? What content satisfies this condition?

As developed earlier, this question must be answered differently at the different ethical standpoints.
When a disinterested and universal sympathy is determinative for ethical valuation, only that can be good
which preserves and heightens the welfare of conscious beings, increases their pleasure or diminishes
their pain. Every action that works in this direction without bringing with it more remote effects in the
opposite direction is thereby justified; every action of which the opposite holds is reprehensible.

\textit{Ethics here only works further on what has been begun by the hand of nature.} For on the whole,
feeling of pleasure is bound up with the sound and natural use of the powers and with what preserves and
furthers life, feeling of pain with the opposite. Pain is a sign of a beginning dissolution of
life\footnote[1]{\textit{Psykologi} VI D.}. When ethics values our willing and acting by the degree in
which pleasure thereby comes to prevail over pain, it works --- by virtue of \textit{the close connection
between pleasure and life, pain and death} --- for the development of human life, of human faculties and
powers, in as rich a measure and in as harmonious a way as possible. Ethical problems therefore also
concern the valuation of the pleasure and pain of single moments in relation to the whole of the
individual's life, and of the pleasure and pain of single individuals in relation to the whole of the life
of the race. Only so far as this valuation is possible can the valuation of acts of will and of actions be
grounded, and ethical action have a definite content. Every act of will and every action is like a stone
thrown into the water; the movement propagates itself in larger or smaller circles. And the valuation
depends on how the action intervenes in the life and feeling of conscious beings. Just as theoretical
science explains one natural phenomenon by another, so ethics values the one feeling by other feelings:
the satisfaction the agent feels in, or obtains by, his
% --- p. 137 ---
\opage{137}willing and acting is to be called unconditionally good only when the effects of the action do not
intervene disturbingly in the feeling of pleasure of other beings. Every such intervention requires
justification by being shown to be a necessary means to a correspondingly greater welfare for the sufferer
himself or for others. So, according to the principle of welfare, the burden of proof must be distributed.
It is the infliction of pain that is to be grounded, while the feeling of pleasure has its provisional
justification in itself. This maxim was first set up by \textsc{Plato} \textit{(Protagoras.} Heise's
transl. p. 198 f.), who taught that when sensual enjoyment is to be blamed, it is on account of its
consequences, not on account of the pleasure it gives at the moment\footnote[1]{Among modern thinkers
\emph{Spinoza} has maintained it \textit{(Ethica} IV, 43. 60. App. c. 30); likewise \emph{Malebranche}
\textit{(Recherche de la vérité} IV, 10). \emph{Arnauld}, on the other hand, in his polemic with
Malebranche, had great misgivings about it. Cf. \emph{Bayle}'s \textit{Dictionnaire.} Art. Épicure.}.

Welfare expresses itself as a lasting state of the feeling of pleasure. It may perhaps be won only by a
course of development that leads through pain and privation; but it is the goal for whose sake pain and
privation are endured. One need not be an optimist to acknowledge welfare as the objective principle of
ethics. For nothing whatever is thereby said about how easy it is to reach the goal, or in what degree it
is at all possible to reach it. Should pessimism be right, should all feeling of pleasure, all happiness,
be an illusion, and painful burning and need the inmost thing in us all, --- the principle of welfare
would not thereby be overthrown. The task would then be to bring as much relief as possible, and such
relief would then be all the welfare that could be attained.

Only from an absolutely ascetic standpoint, that is, one that would maintain that pain in itself and to
all eternity is to be preferred to pleasure, can the principle of welfare be disputed here. But such a
standpoint cannot be carried through either; even the ascetic torments himself and others only in order
to produce a greater feeling of pleasure thereby. Every ethical line of thought leads out in the end to a
point where the ultimate ground of the approval of an
% --- p. 138 ---
\opage{138}action proves to be that at some point or other in the world pleasure is produced or pain annihilated.

As already remarked earlier (III, 11), the resistance to this principle is due in part to
misunderstandings that the words have caused, and I prefer to use the word ``welfare'' instead of
``happiness'' or ``utility'' in order to avoid such misunderstandings. --- By ``happiness'' [\textit{Lykke}]
one often thinks of something accidental and external. ``It is called happiness,'' says P. E.
\textsc{Müller} in his ``Synonymik,'' ``when a good befalls us that we cannot ourselves bring about.'' It is
in this sense that the poet speaks of ``the capricious play of fortune.'' Often the word happiness is also
used of sated contentment, rest in what has been attained. Whichever of these meanings one takes, it is
clear that there can be something better than ``happiness.'' But why should the word happiness not be
usable of the feeling of pleasure that is grounded in and bound up with self-activity and can be held fast
only by constantly continued development? --- Here, then, \textit{provisionally}, no kind or degree of
feeling of pleasure is to be excluded; the order of rank among the different feelings of pleasure can be
fixed only by the following more special inquiries. But in and for itself every ethics that acknowledges
sympathy as the motive of valuation must also find it valuable that one confers on others goods that they
cannot themselves bring about; if they could bring about these goods themselves, they would, after all,
need no help. And satiety is, after all, in and for itself nothing slight, whether it follows on bodily
or on spiritual hunger and thirst. --- By ``utility'' [\textit{Nytte}] one easily thinks only of what is a
mere means to something else, and thus has no value in itself. It is a business expression and grates on
many ears merely because it is not solemn enough\footnote[1]{For solemn use we have the old Danish word
``Gavn.'' The word Nytte came from Germany. ``Since \textit{Gavn} was more rarely heard in daily speech, it
acquired a more general and more honorable meaning, and the priest used it in the pulpit'' (P. E.
\emph{Müller}). When Jesus says: ``What \textit{gavner} it a man if he gains the whole world but suffers
harm to his soul,'' it would, according to P. E. Müller, not have been right to use the word ``nytter''
instead of ``gavner''! --- (Addition 1913). \emph{Kant} would not call respect a feeling; it was to be
neither pleasure nor pain. Fichte distinguished sharply between ``Genuss'' and ``Würde.'' According to
Hegel the happy times are the blank pages of history, because opposites are lacking. Carlyle distinguishes
sharply between ``happiness'' and ``blessedness.'' --- On the other hand Rousseau maintained that there is
no happiness without courage. Grønbech teaches (in his book Lykkemand og Nidding) that the Germanic word
Lykke contains a potential moral valuation; Lykke is the will-side in man, the last and deepest expression
of his being; it gathers in itself the conditions of life of the kin. And Jhering declares: ``Wohlsein ist
aufgespeicherte Arbeitskraft.'' --- One sees how the ceremonial of words can gain influence on the fate of
theories.}. But in turning away with overstrained indignation
% --- p. 139 ---
\opage{139}from this plain prose, one easily comes to overlook that ``having value in itself'' can mean nothing
other than being the immediate object of a need of feeling. Value ``in itself'' is the same as immediate
value. As the ultimate standard we always find again the pleasure and pain of one or more conscious
beings. --- By welfare I mean the ``true'' happiness or the ``true'' utility, leaving it to the following
sections to show its conditions more closely.

2. But the general concept of welfare contains some difficulties that we must bring forward
here\footnote[1]{In my \textit{Etiske Undersøgelser,} p. 24---42, I have discussed the various objections
that might be raised, or have been raised, against the principle of welfare, in part somewhat more fully
and from a somewhat different side than in what follows.}.

Can a lasting state of the feeling of pleasure not be reached at very different stages of development and
by very different roads? ``True,'' or rather \textit{actual}, happiness can, after all, mean nothing other
than full satisfaction of all the need that can be felt. And the need one feels conforms to the nature one
happens to have. Must not, then, the welfare or happiness of the one be just as good as that of the other,
since each of them can get no further than to obtain satisfaction for his need, and the one naturally
does not need satisfaction for a need that not he but another has? How, then, can the principle of welfare
furnish a standard for the valuation of the feelings of pleasure at the different stages and in the
different beings?

And second. Does the development of human faculties and powers, does the whole development of culture
really lead to
% --- p. 140 ---
\opage{140}welfare? Through culture, to be sure, the means at the disposal of the race are increased, and faculties
and skills are trained that were not present before. But are not the needs, both the material and the
ideal, increased at the same time, and with them the possibilities of privation, pain and disharmony? Even
the development and widening of sympathy and of the feeling of duty threaten welfare, since one thereby
becomes vulnerable at far more points, makes greater demands on oneself and others, and does not so easily
find peace and rest in one's mind!

These two questions I shall examine each by itself.

3. \textsc{Stuart Mill} has said: ``It is better to be a human being dissatisfied than a pig satisfied;
better to be Socrates dissatisfied than a fool satisfied.'' % Mill's English (Utilitarianism, ch. II), where Høffding gives it in Danish
He grounds this assertion by saying that even if the pig and the fool should think the opposite, their
judgment would have to be rejected, because they do not know the higher standpoint from which man and
Socrates regard life, whereas man knows the needs of the pig, and Socrates sees through the fool. One must
keep to the judgment of those who know both kinds of needs and can therefore make a valuation of
them\footnote[1]{\textit{Moral grundet paa Nytte- og Lykkeprincipet.} Danish transl. p. 10---12. ---
A similar line of thought is found already in \emph{Plato} \textit{(Statens 9de Bog.} Heise's transl. p.
86---92). \emph{Voltaire} too has discussed the problem. See \textit{Histoire d'un bon brahmin.}}.

I must, however, take the pig and the fool somewhat under my protection. The difficulty is greater than
Mill supposes. Man, to be sure, knows the needs of the pig, and Socrates will have no difficulty in
entering into those of the fool. But man does not have the needs of the pig, Socrates does not have those
of the fool, as \textit{sole rulers}, and it is, after all, on this that everything depends in this
connection. Man cannot make himself a pig without ceasing to be man\footnote[2]{If it succeeded, the
enjoyment would perhaps become all the greater, provided, that is, the Leipzig students in Goethe's Faust
are right when they sing: Uns ist so kannibalisch wohl, als wie fünfhundert Säuen.}, and Socrates will
hardly be able to make himself
% --- p. 141 ---
\opage{141}so one with the fool that the Socratic needs wholly fall away. Now if the pig can obtain full
satisfaction of all its need, is its happiness not greater than man's, who hardly gets all his wishes and
all his craving fulfilled? And the fool, who has not many thoughts and makes no great demands on life, is
he not happier than Socrates, who spends his whole long life learning to know himself and awakening
others, only to declare at last that death is at bottom better than life?

What is in question here is, after all, not a definite sum of external goods, but goods that are really
goods for this definite being. One and the same sum of money may be a great gain for the one and a sheer
trifle for the other. It is a chief psychological law that the degree of a feeling is determined by the
change produced in the individual's whole state.

Finally, he who has won \textit{full} satisfaction for his need has no motive to compare his state with
that of others. That the pig is not bothered by man, nor the fool by Socrates, comes of course from lack of
capacity. But even where the intellectual conditions for making such a comparison are present, one will
not be able to feel any sting in the comparison, any need to get beyond one's state. It will go as in
\textsc{Dante}'s Paradise, where there are different degrees of blessedness, but where those placed on the
lower degrees feel no unrest at the thought of the higher, since their ``will is brought to rest, so that
only in what it has it finds its joy.'' The poet continues:

\begin{quote}
Then it became clear to me that everywhere\\
in heaven, where I wander, there is Paradise,\\
though God lets grace rain down in different measure\footnote[1]{Dante: \textit{Paradiso.} III v.
88---90. --- Cf. \emph{Augustine}: \textit{De civitate dei} XXII, 30.}.
\end{quote}
% Høffding's Danish: Da blev det klart mig, at der alle Vegne / er Paradis i Himlen, hvor jeg vanker, /
%   skønt Gud forskelligt lader Naaden regne.

In \textsc{Dante}'s Paradise, therefore, all striving and development is over. But will it not go so
everywhere? And how can the principle of welfare be taken as the basis for a comparison of the states of
different
% --- p. 142 ---
\opage{142}beings, when each of these has the character of complete contentment?

Having now posed the problem that presents itself here as sharply as possible, I shall try to answer it.

Just as in external nature there is no absolute rest, but everywhere movement under different shapes, so
too welfare as the absolute satisfaction of every need for every being is possible only approximately, and
every approximation to it will be subject to attack. This lies in the fact that every being is subject to
external and internal changes. Nature and life do not stand still, but go on working both in the
individual and around it, and these changes will awaken new need and set new tasks.

As regards the external changes, no closer demonstration is needed. We always seek in vain to withdraw
into an idyllic existence. Sooner or later the idyll is interrupted by changes in the surrounding world
that make new efforts and labors necessary.

And with the external changes follow the internal. Ideas of the new and the strange will not in the long
run be kept away, and they will set feeling and will in motion. If we may apply our psychological laws to
the conditions in Dante's Paradise, then in the blessed, together with the ideas of the higher stages of
life, there will also arise a craving and longing to attain them. --- And even if new ideas do not arise
and work, feeling itself will need contrast and variation in order to keep itself fresh and living.
Uniformity and repetition dull and weaken. That is why the idyll as a genre of art needs as its
background a hint of a larger, manifold and agitated life in order to produce its effect. Thus in
\textsc{Goethe}'s ``Hermann und Dorothea'' the revolutionary war and its effects form the restless
background\footnote[1]{After Schiller, in the magnificent poem \textit{``Das Ideal und das Leben,''} had
depicted the transition from the heavy struggle with reality to the ideal Olympian harmony, he conceived
the plan of writing an idyll, a picture of a blessed existence raised above all strife. He regarded this as
the highest poetic task. But he saw the necessity that the strife and struggle of human life should yet be
recalled: ``The chief persons would indeed be gods, but in Hercules I can still find a point of attachment
between them and mankind and bring a movement into the picture.'' (Briefwechsel zwischen Schiller und
Wilhelm v. Humboldt. Stuttgart und Tübingen 1830. p. 328). --- The plan was not carried out.}. ---
Finally, when
% --- p. 143 ---
\opage{143}uniformity and repetition do not lead to decline and weakening, there will naturally arise a need for
movement and activity. We move, after all, not merely in order to attain definite ends, but also in order
to use our powers and find an outlet for accumulated energy. This need for movement will grow the more
powers are accumulated, the longer the rest lasts. The activity it leads to will lead to new experiences
and thereby to new struggles.

Even if these causes do not work in the single person who has reached a certain state of equilibrium, they
will nevertheless work in the race after the lapse of not many generations. It will appear that, if
there is no decline or dissolution, every attainment of an end will be only the beginning of a new end's
being set. Welfare is therefore an illusion if by it is meant \textit{a passive state, produced once for
all.} It must consist in \textit{activity,} in work, in development. This lies already in the word: welfare
[\textit{Velfærd}] is to \textit{fare} well --- that is, to be in movement, not at rest. When the word
``fare'' can also mean ``to be in a certain state,'' this is, according to Molbech, a ``figurative and
improper'' meaning. Welfare as rest can mean only a \textit{provisional} conclusion, the attainment of a new
level on which a new development can be begun. We need not, however, take back on that account our
definition of welfare as a lasting state of the feeling of pleasure. What must be rejected is merely the
idea of a \textit{passive} state. Activity, after all, is a state too, just as rest is. By a natural
illusion we like to think of our more distant ends as states of rest. Because of the distance we cannot see
the multiplicity they contain and the tasks they set. From this illusion actual experience constantly tears
us. We learn to understand that only the happiness found in work for the advance of life in ourselves and
others is built on firm ground.

There is here a significant parallel between the theory of welfare
% --- p. 144 ---
\opage{144}and the theory of conscience. (Cf. IV, 3---4). Neither of them can assume any absolute conclusion, but
both lead us to the thought of a course of development that cannot be surveyed. And there is also an
interaction between them. For the richer an idea one has of what belongs to welfare, the more conscience is
developed and exercised, since it gets the more tasks to treat. And the finer and more sharply developed
conscience is, the greater the demands it makes on what (in others and in ourselves) is to deserve the name
of welfare. He whose eye is dulled and whose horizon is narrow is easy to satisfy.

When great minds often judge disparagingly of the happiness that can be attained in the world, they are
thinking more of the passive states that fill the intervals between the times of great activity than of
the very working that has made them the great minds. It is psychologically intelligible that in the active
times, when life stirred strongly in them, they did not think of their own state, but were absorbed in the
work for their cause. It is always a bad sign when attention glides over from the thought of the end to
the feeling of pleasure its attainment can bestow\footnote[1]{Psykologi VII, B, 1 a; 3.}. One thereby comes
to separate what is not to be separated, and the state is thereby transformed from an active into a passive
one. At the same moment that one asks oneself whether one is happy, one easily isolates oneself from the
content in and for which one lives. That question does not come up at all so long as one is absorbed in
this content; that it comes up points to a beginning doubt. It holds, after all, of all ethical reflection,
of all knowledge of good and evil, that it \textit{presupposes} (not, as in the old story, \textit{brings
about}) the end of the paradisal state. Reflection can of course be awakened from without --- by him who
brings the serpent into Paradise. But then it is presupposed that it has been awakened in others, and it
will not work where there was not a beginning disharmony beforehand. Happiness is expansion, not
reflection, or --- to use Ludvig \textsc{Feilberg}'s expression --- ``running straight on,'' not
``circling.''
In a certain sense one may therefore say that we only rightly know happiness when it has
% --- p. 145 ---
\opage{145}vanished. The striving for happiness then becomes a striving to reproduce at a higher stage what we have
known at an earlier stage. And the more activity itself is bound up with the feeling of pleasure, the less
means and goal --- the striving for happiness and happiness itself --- fall apart, the more firmly welfare
is grounded. It is by no means a matter of indifference how means and goal, cause and effect, here stand
to one another, --- whether they stand in a purely external connection, or whether there is an inner
coherence between them, so that the end already stirs in the means, and the effect is only an unfolding of
what lies in the cause. Two states may apparently be quite alike, and yet the one may, on account of its
manner of coming into being, acquire greater value than the other through the hidden elements (the
potential energy) it contains, namely partly the possibility of rich memory, partly the possibility of
future development. The value 2 may here be greater when it is reached by $7 - 5$ than when it is reached by
$1 + 1$. The welfare that is reached through pain and struggle may be greater than the apparently same
welfare that is reached by quiet and gentle unfolding, --- may contain more possibilities than this. --- An
end completely and once for all \textit{attained} is no longer able to set the impulse to activity in
motion, and the feeling attached to the idea of it can therefore not become so strong and living either
as the one conditioned by the idea of an end that can constantly be realized in a higher degree and to a
greater extent. The highest we can conceive is an advance in which each single step is felt as a good,
because it sets our powers in motion without yet straining them beyond their capacity.

We are thereby led over to the second of the questions raised, that of the ethical significance of the
development of culture.

4. Welfare consists in activity, but it does not follow that welfare is attached to every activity. By this
simple proposition the relation between \textit{ethics} and \textit{culture} is provisionally designated.
Culture is the wider concept, ethics the narrower. All ethical striving is cultural activity, but not all
cultural activity is ethical.

% --- p. 146 ---
\opage{146}By culture we understand the working up of the material given by the hand of nature. Yet no fixed
boundaries can be drawn between nature and culture. What is culture seen from one point of view is nature
seen from another. For what the one generation laboriously produces is a given foundation for the
following ones. Not even in the animal world do we always find pure and bare natural life. The bird has
its nest, which it has built itself, and the fox its den, which it has dug out itself. And besides, every
generation has in its nature an inheritance from the experience and activity of earlier generations.

The concept of culture is in and for itself a purely formal concept. Working up and activity tell us
nothing about whether the work and activity are for benefit and advance. The so-called consciousness of
culture can often be a heartless joy at seeing the many wheels in motion, the many threads crossing one
another that are knotted together into an ingenious web. One then does not think of what is attained by
the whole, or of those who are crushed under the rolling wheels. What value, really, has the whole
restless labor continued from generation to generation? Our skills and our needs rise; but would it not be
better to be spared these new skills and these new needs, if welfare could be attained on simpler and
easier terms? --- By this question we are led back again to the first difficulty that lay in the concept of
welfare. The two difficulties into which this concept leads us are so closely connected that the
overcoming of the second forms the continuation of the overcoming of the first. From welfare as a passive
state we had to pass on to welfare as activity; but here it appears again that a closer limitation is
necessary.

This limitation is given by the general theory of \textit{the connection of the feeling of pleasure with
the sound, natural and life-developing use of the powers.} Activity ceases to be welfare when it either
overstrains the powers, or when the powers are scattered and split up, or when the use of power takes
place one-sidedly in a single direction at the expense of other important directions. Where the
\textit{resistance} is greater than is necessary for
% --- p. 147 ---
\opage{147}the unfolding of the powers, pain and suffering will arise; and such a painful activity can be
justified only as a transition to a state in which a more harmonious relation has been brought about
between the external and the internal conditions. Where the powers are not in the end gathered toward a
common goal, life dissolves into fragments, and no collected feeling of life becomes possible. And
finally, where some powers are exercised at the expense of others, a disproportion between the inner and
the outer will sooner or later appear; the \textit{one-sidedly} developed inner powers will not answer to
the many-sided demands that the external conditions of life make. In these three respects the
development of culture can go astray. It has a tendency to strike into every new direction of activity
that offers itself, being guided by the momentary satisfaction that can thereby be won. There is something
blind and reckless in cultural activity that is connected both with its merits and its faults. It is loath
to leave any possibility whatever untried; it has its feelers out in all directions; thereby the
significant experiences are made; but thereby it can also lead to disharmony and suffering for the single
individuals and for the race.

Culture is not an object of choice. It is the continuation of the development of nature and can just as
little be stopped as this. The need for restless development, for trying new ways and for training
specialties and producing diverse forms, is the great means in the hand of nature for making possible the
continued existence of living beings when these increase greatly within a relatively small area.
Differentiation and restlessness spring from the density of population\footnote[1]{Cf. Ch. Darwin:
\textit{The Variation of Animals and Plants under Domestication.} London 1868. I. p. 5---7. --- E.
\emph{Durkheim}: \textit{De la division du travail social.} Paris 1893. p. 282---294.}. But because the
development of culture, as it actually takes place, is a necessity of nature, it is not said that it is
admirable. Ethics feels no special admiration for an order of the world under which development is often
not possible without overstraining of the powers and without one-sidedness and dispersion in the use of
the faculties. It is
% --- p. 148 ---
\opage{148}not so thick-skinned that over the apparent brilliance of the external results it can forget the
anguish and the pain, the sweat and the blood, that they cost. Therefore it demands that the heavy
burdens be lightened, that the split powers be gathered, and that all valuable faculties be developed. The
demands it thus makes on the development of culture are related to those it makes on human willing and
acting in general (III, 20). Through insight into the conditions of the development of culture it seeks to
find means to alter, to mitigate and to humanize, so that culture and welfare may be brought into as close
a connection as possible. It is not sentimental and shortsighted, so that it should forget that advance can
take place only through effort and suffering. Only in taking part in the work of culture does it find a
real task and a real content for the life of the will. And in particular it sees that only through the
development of culture and the differentiation this brings with it do personal peculiarities come to be
unfolded and employed. Each single person is, to be sure, made more dependent on the others; but with this
mutual dependence follows a mutual supplementing and a more inward union than is possible so long as the
division of labor has not come into force. Both the multiplicity and the coherence can now become greater
and firmer. This culture can bring about --- but it has side effects that prevent it from unfolding all its
happiness-bringing effects.

The chief ethical task with regard to culture is to enjoin that life is not to be made a mere means to the
solution of impersonal tasks. Culture is to reach its completion in personality, so that this itself comes
to stand as the highest culture. Personal life must neither be wasted in fruitless struggle against
insuperable resistance nor be scattered and narrowed by too great multiplicity or one-sidedness. Culture
contains a possibility that men may come to set themselves higher ends and to grow with their higher ends;
it is this possibility that ethics demands be used. Enthusiasm for the state of nature is unethical when it
leads to giving up work. A meager content is easy enough to order; where there are no differences and
oppositions, it is easy enough
% --- p. 149 ---
\opage{149}to keep spiritual disharmony away. Often, no doubt, a restriction of needs may be the only possibility
of winning victory over the obstacles. But often precisely the feeling of privation and want may awaken
new powers and thereby make possible a satisfaction that would otherwise not have been possible. Neither of
the two assertions, neither that needs are to be restricted nor that they are to be increased in order that
welfare may be attained, is therefore unconditionally true.

This finds its application when the question arises with what right one abolishes the harmony and
happiness that are present in order to lead men into new ways in the field of spiritual or material
culture. So when savage nations are torn out of their ``state of nature,'' --- when one robs the child of
its illusions or awakens new needs in it, --- when one awakens doubt where a blind but secure faith reigned
before. The decision in such cases can be made only with the help of the principle of welfare. If one is
sure that by the new road a new harmony, a new satisfaction, can be won that has a firmer ground and richer
content than the earlier one, then it is justified to intervene. And it becomes a duty where one stands
before a sleepwalker's state that may end in misfortune for the sleeping individual himself and for others.
A man's idyll may bring harm and pain to others, or at any rate diminish their happiness. There may be unused
powers that the race needs, although the individual has so far felt no need to use them. Even if they
should be awakened by a painful shock, it may be a duty to administer it. It may in the single case be
extraordinarily difficult to decide whether such a painful shock is to be administered. The ethical
calculus will never be able to lead to full certainty: one must \emph{venture}, and a leap must be made out
into the dark. The principle of welfare can here serve only as a guiding point of view. It does not conflict
with venturing; but it conflicts with venturing the leap into the dark when one has no hope that there is
light on the other side of the dark. The relation may be illuminated by an analogy that is, however, more
than an analogy. It does not conflict with the method of empirical science to set up
% --- p. 150 ---
\opage{150}hypotheses, provided only one holds fast to the endeavor to win the confirmation of experience for these
hypotheses. Every hypothesis is a venture, but the history of science shows us that without such ventures
our knowledge would not have advanced. So too it will be a consequence of the principle of welfare that one
must venture, that one cannot always wait until full guarantee of happiness-bringing effects has been
attained. The principle of welfare must expressly lead to misgivings about letting reflection and the
weighing of grounds for and against play too great a part, since energetic and resolute action, which in
fact proves to be of such great significance, is thereby excluded. \textsc{Paul Hensel}\footnote[1]{\textit{Ethisches
Wissen und ethisches Handeln.} Ein Beitrag zur Methodenlehre der Ethik. Freiburg i. Br. 1889. p. 31.}
says in his criticism of utilitarianism: ``The cautious utilitarian must in any case take a dismissive
attitude toward an attempt to lead society into new and untried ways; for every such attempt undoubtedly
disturbs much happiness, and it is problematic whether it is able to create new happiness.'' But precisely
the principle of welfare leads one not to be all too cautious. Caution is a burgomaster's virtue, but it
would not agree with the principle of welfare if we were all burgomasters, and even burgomasters can be too
cautious. The adherents of utilitarianism, especially in earlier times, often --- like the empiricists in
general --- underestimated the significance of the venture, of the hypothesis. But that was a decided
inconsistency on their part, and no decisive objection to the principle of welfare can be built on it. When
Hensel declares besides: ``Where great ideas enter into struggle with one another, at the turning points of
history, where a new time comes into the world in thousandfold throes, there utilitarianism must fall
silent,'' --- he overlooks that the great ideas in history have always pointed precisely to a higher stage of
life, a kingdom of personalities of a more perfect kind than the existing societies made possible, and that
it was only in confidence in the greater value of the goods this kingdom would bring that the struggle was
ventured. But he also overlooks that history in fact reaches its
% --- p. 151 ---
\opage{151}results neither by the shortest nor by the cleanest roads, and that one need not approve of these roads
because one holds fast to the value of what has been attained, through many displacements, by following
them. The result does not sanctify the method.

If we would always direct our conduct by the need and the wants that men had at the moment, we would
thereby lower the level of human life and human culture. The principle of welfare demands that we do not
shun the struggle against prejudice and inertia. The best one can do for others is often simply to let them
feel that in their wishes and wants they still stand too low and do not make great enough demands. Thus ---
to take a single example --- the great artist often goes his lonely way, not understood or even misjudged by
the great multitude. And yet he follows the principle of welfare, perhaps without thinking of it, by
strictly asserting the demands of art. He increases the spiritual capital of the race and bestows on it a
power that may in the future come to work in large circles. Only a shortsighted conception and application
of the principle of welfare keeps to the need prevailing at the moment, instead of looking to the lasting
conditions of life of the race and to the constant sources of new life and new activity. \textsc{Goethe}
said (in a conversation with Eckermann): ``In my calling as an author I have never asked: what does the
great mass want, and how can I benefit the whole? I have always striven only to make myself more insightful
and better, to raise the worth of my own personality, and then always to utter only what I had recognized
as good and true.'' Had he wished to conform to the taste and capacities of the public, he could have made a
greater momentary effect, but his works would then also have had a shorter life, and he would not have
become the great teacher of many generations. --- How large a horizon is to be applied in the single case in
discussing the value of an action or an endeavor, nothing can be said in general. With a vigorous conviction
of having the conditions for accomplishing something there is as a rule bound up a more or less conscious
and confident appeal to the judgment
% --- p. 152 ---
\opage{152}of the future. One must venture every time something is to be put into the great connection of the
world, which we cannot survey. And even as regards the subsequent confirmation by experience, which
historically decides the value of the venture, it holds that it can never be wholly concluded so long as
historical development lasts. There will always be room for a revaluation. Nietzsche has even demanded a
``revaluation of all values.'' That is to take one's mouth rather full; but a limit to how far the
revaluation that constantly goes on, and must go on, is to be extended cannot be set beforehand. It will
therefore be impossible to form an absolutely concluded ethical system. Every science that holds fast to the
necessity of confirmation by experience (verification) must at every time bear the stamp of incompleteness.
It is therefore no decisive objection to the conception of ethics put forward here that it leads only to
``empirical certainty''\footnote[1]{Axel \emph{Hägerström}: Undersökning \textit{af den empiristiska
etikens} möjlighet. Upsala 1895. p. 16. 129.}. Just as human motives of valuation are in constant
development, so a more extended application of the principles of valuation will always be possible. We are
constantly on the way --- and not we alone: the world, existence as a whole, is perhaps not finished either,
but in constant development. The ethical striving itself is, after all, a sign that the world is not
finished. There is a part of the world's work that only men can do. A part of the development of the world
takes place through human action itself.

5. The roads that human action follows in order to reach its ends are by no means indifferent. As we have
seen, it is, according to the principle of welfare, an imperfection when road and goal, cause and effect,
stand in an external relation to one another. These roads will be the more perfect the more the means
employed are in the spirit of the end itself, or, as \textsc{Bruno Wille} has expressed it, are \textit{pure
means.} The well-known proposition that the end sanctifies the means starts from a purely external relation
between means and end. The more the means is in the
% --- p. 153 ---
\opage{153}spirit of the end, and thus is ``holy'' in itself, the more the end is already attained, although we are
still on the way: it is attained in the disposition that leads to working in the spirit of the end and is
thus an anticipation of the end. The proposition that the end sanctifies the means may for that matter
have its justification. There may be cases in which in the service of an end we do, and must do, things we
would not otherwise do --- not because they are bad, but because they are indifferent: what under other
circumstances has no value for us now acquires value as a means. Modern Jesuit ethics restricts the
proposition to such cases\footnote[1]{Cathrein: \textit{Philosophia} moralis. p. 80.}. But to these must
be added such cases in which a lesser evil must be inflicted in order that a greater evil may be avoided.
When a man who is dying of hunger smashes a pane in a baker's shop and takes a piece of bread, he performs
an action that is sanctified only by the intention. Here belongs also all education through suffering
punishment. According to the principle of welfare, after all, the infliction of pain is justified and
necessary when it is the only means to an indispensable good; but that is also the only ground that can
make the infliction of pain permissible. The decisive thing will then of course be whether the end itself
is really ``holy''; if it is not, it cannot sanctify the means. If, then, it is egoistic ends or party ends
(ecclesiastical, political, literary, etc.) that lie at the basis, the proposition becomes questionable.
And even if the end is really valuable, it is not said that it is able to sanctify the means. The
employment of a means may bring with it side effects that cause greater corruption than giving up the good
intention would. He who would murder with a noble intention would surely, by such a breach of the
elementary laws of human living together, produce a greater evil than he could remedy by attaining his
intention. Assassinations as a rule strike only the external consequences or symptoms of a bad social
state, not the real cause, and they often help to make bad worse. Significant ends require a long course of
development for their realization and are often made possible precisely only
% --- p. 154 ---
\opage{154}when there is no breach of the elementary ethical laws that one thinks oneself entitled to break in the
consciousness of one's good intention. Better social conditions are attained by working on the character and
way of thinking of the people, by developing its powers and resources. Finally, it stands as a
psychological problem how the noble intention and the crude or cruel means can dwell side by side in the
same soul. It can perhaps be explained only by the deep indignation that resistance to the great end breeds:
but actions that come into being in this way are rather to be conceived as symptoms, as outcomes of an
unhappy social state and of the tormented mood such a state has awakened in idealistic natures, than as
utterances of will that could have an exemplary character and lay claim to ``holiness.'' And precisely here
side effects lie so extraordinarily near: through displacement of motive what was at first only reluctantly
chosen as a means can spread in the soul and at last become the chief end. What was begun in a purely ideal
form can, through increasing bitterness, end in blind hatefulness and lust for destruction. It is a
dangerous road one treads when one begins to choose means that must first be ``sanctified.'' The history of
coups d'état and assassinations bears witness to this. --- Besides the involuntary displacement of motive
there is also a voluntary displacement, namely when one ``directs the intention'' in order to be allowed to
perform actions one has a mind to. It was the permissibility of this that \textsc{Pascal} accused the Jesuits
of teaching: although Christ forbids repaying evil with evil, we should nevertheless be allowed to persecute
our enemies, provided only we do it --- not in order to inflict evil on them, but --- in order to assert our
honor!\footnote[1]{\textit{Lettres écrites à un provincial.} VII.}

% --- p. 155 ---
\opage{155}\chapterhead{VIII.}{INDIVIDUAL AND SOCIAL ETHICS.}

1. There are some concepts that will unavoidably be formed whenever an ethical theory is developed in its
full range, since they designate different sides of, or points of view for, all ethical action.

An action must have an end, which consists in the valuable thing that is striven for. The concept of it is
the concept of \textit{the good.} Every end stands as a good, and this concept is in its generality common
to all practical sciences, to political economy and jurisprudence as well as to ethics. (Cf. III,
1---14). The whole character and direction of an ethics are determined by the ends that are taken as its
basis from the first, and therefore the greatest difficulties in every ethical discussion arise when the
dispute comes to concern the fundamental ends (the fundamental values, the motives of valuation). (Cf.
III, 12---13). What acquires value and can therefore become an end depends on the need and striving of
life, especially of the life of the soul, and will be different at the different stages of life. (Cf. I,
4; III, 13. 21).

The concept of end already implies the presupposition of a certain distance between what is striven for
and the given state. What is immediately given as valuable need not be set up as an end, except insofar as
it is a matter of preserving it in the future. Now when the factor of distance is especially stressed, the
difficulties of attaining the end being sharpened, the concept of \textit{duty} arises. Bringing about what
is acknowledged as a good then becomes a conscious task; man feels himself bound to work for it and to push
back opposing tendencies. The presupposition of this is that man wishes to hold fast to coherence and
agreement with himself, to remain identical with himself, true to himself through the changing moments and
circumstances. Only then will he maintain the conviction of the validity of the acknowledged value beyond
the moment of acknowledgment.
% --- p. 156 ---
\opage{156}From the point of view of self-preservation it becomes a duty to subordinate the satisfaction of the
moment to the assertion of the more comprehensive values of life; from the point of view of the race it
becomes a duty to subordinate the purely individual interests to the common ones. (Cf. III, 5---7.
9---10).

The concept of duty concerns especially the single actions that stand before self-reflection as a
necessary consequence of the acknowledged end. If, on the other hand, the weight is laid on the abiding
qualities of character that by their nature will lead to endeavors to realize the end, the concept of
virtue arises. Between duty and virtue there is in ethics a difference similar to that between emotion and
disposition in psychology. (Cf. \textit{Psykologi} VI E, 5). Duty presupposes a concentration of the ethical
feeling on a single point, with a conscious effort at the single moment; virtue presupposes a lasting mood
and striving, without the consistency with which the single action follows from the acknowledged end
needing to be an object of full consciousness. Just as the transition from the immediate feeling of value
to the feeling of duty takes place through an increasing stress on the factor of distance, so the
transition from duty to virtue takes place through increasing practice in the moods and endeavors that the
acknowledged end calls forth, until at last the single action proceeds from an involuntary need.

The concepts of duty and virtue spring as a direct consequence from the concept of the good. The concept
of right, on the other hand, is an indirect consequence of it. Ethically justified is everything that does
not conflict with or hinder the fundamental value that my highest end presupposes, and therefore does not
conflict with any duty or virtue either. I am within my good ethical right wherever my conscience can make
no objection, and where I therefore also demand that others shall acknowledge my conduct. If this
acknowledgment is not forthcoming, I confidently appeal from the imperfect conscience to a perfect
conscience, which I hope will develop and which I strive to have developed in them. The ethically justified
is the free unfolding of life, the free use of the powers. Life does not exist for the sake of conscience,
but conscience for the sake of life; therefore
% --- p. 157 ---
\opage{157}conscience upholds me in my conduct when this does not injure the values of life. First and foremost the
fundamental value itself stands as justified, subject only to the condition that it is to be realized in a
way agreeing with its own nature. (Cf. VII, 5). --- The ethically justified is not merely of a negative
nature. It is wrong to call it the permissible. That would imply that it is merely tolerated. But where
life has unfolded freely, without conscience needing to intervene regulatively, there the goal toward which
all ethical development strives has already been reached. It is of decisive importance that the domain free
from conscience be not narrowed. Ethics here points beyond itself, to an order of things where it itself
need not exist. Since it is a duty to work toward such an order of things, it is also a duty to assert one's
right where it coincides with the right of life itself. It is here right that grounds duty. The ethically
justified springs from the same source as the concepts of duty and virtue, but nearer the spring of the
source than these. ---

The character of every ethics is determined by the relation in which these four fundamental concepts stand
to one another. They are so closely connected that one will always, in a more or less natural way, be able
to find one's way from any one of them as a point of departure to the other three. They state the most
important formal points of view from which every action and every relation can be regarded in an ethical
respect. It will of course be of decisive importance for the characterization of an ethical theory which of
these concepts it takes as its point of departure. If one would ground ethics on experience, there can,
however, be no doubt that the concept of the good must be the first concept of ethics. An action, a striving
and a valuation necessarily presuppose an end, and an end in turn presupposes an immediate value. This point
of departure was accordingly taken by \textit{Greek ethics,} and its significance for all time rests not
least on this. In sharp opposition to this Kant would take the concept of duty as the basis and determine the
other fundamental ethical concepts from it. Only in an artificial way could he derive special
% --- p. 158 ---
\opage{158}ends by this road\footnote[1]{Cf. below VIII, 6, and \textit{Den nyere Filosofis
Historie}\textsuperscript{2}. II. p. 81---84.}; the chief end for him became the will of duty itself, as he
expressed it in his famous proposition: ``Nothing can be conceived in the world or outside the world that
could be called good without limitation, except solely the good will.'' To the question when, then, a will
is good, Kant answers: when it is in agreement with the nature of a rational being, and therefore in its
principles the same in every rational being. In this way duty itself becomes the highest good. Virtue
consists for Kant in the energy with which man practices his duty (die Stärke des Menschen in Befolgung
seiner Pflicht. Tugendlehre § IX). There are therefore many duties, but only one virtue and one good. ---
Against Kant \textsc{Schleiermacher} stressed that one cannot will and act without having an end. (See above
III, 1). He showed more closely in a series of interesting essays (which are found in the second volume of his
philosophical writings) how the concepts of the good, duty and virtue lead over into one another. In his own
ethics he took the concept of the end as the basis and derived the two other fundamental concepts from it.
\textit{(Philosophische Sittenlehre.} § 113---121). The concept of right he does not set up in this
connection. He has, however, indirectly prepared it through his criticism of the concept of the permissible,
which --- as he aptly shows --- belongs to jurisprudence, not to ethics, and comes into ethics only when one
confuses this with jurisprudence\footnote[2]{\textit{Ueber den Begriff des Erlaubten.} (Werke zur Philos.
II.) p. 443---445.}. In the same direction as \textsc{Schleiermacher} utilitarianism, and in general all
ethics of welfare, must of course go. --- The third fundamental concept, virtue or the ethical quality of
character, is taken as the basis in modern times in \textsc{Herbart}'s ethics\footnote[3]{Cf. \textit{Den
nyere Filosofis Historie}\textsuperscript{2}. II. p. 254 f.}. Ethics here becomes first and foremost a
description of qualities, a setting up of models. Besides \textsc{Herbart}, \textit{Swedish idealism} and,
with us, F. C.
% --- p. 159 ---
\opage{159}\textsc{Sibbern}'s moral philosophy also go in this direction\footnote[1]{\emph{Biberg} criticized
Schleiermacher's doctrine of the three fundamental concepts and demanded that the concept of virtue be taken
as the basis, because it states the source of ethical action (fons actionis), while the concept of duty
states the single action (actio) that derives from this source, and the good the result that arises through
the action (opus morale). (Cf. Axel \emph{Nyblæus}: \textit{Den filosofiska forskningen i Sverige.} III, 2.
p. 228). ``All morality,'' says \emph{Boström} \textit{(Föreläsningar i Etiken.} Ed. Ribbing. p. 70),
``rests on the will as principle or point of departure\dots{} It can be regarded 1) as resting: as a
constitution, but at the same time the ground of activity; 2) as becoming: whereby it has the idea as its
law; 3) as having come to be or attained, or as the highest good.'' --- F. C. Sibbern takes the doctrine of
the disposition (the fundamental virtues) as the basis, because it depends on the way in which man is fitted
into the great whole of life. Then follows the doctrine of ends and last the doctrine of duties.
\textit{Moralphilosophie som Retsindigheds- og Tilbørlighedslære.} p. 20.}. The difficulty of
carrying it through lies in this, that the value of a quality of character can be grounded only by its
containing a tendency to work in the direction of a valuable goal. The concept of virtue, like the concept
of duty, depends on the concept of the good.

When in the year 1896 I lectured in Zurich on the doctrine of ethical principles, I kept essentially to the
concept of end and discussed the problems contained in it. Since the concept of duty arises only where a
relation of distance enters between value and reality, and since the value of duty depends on the value of
the end, I did not speak of it in particular. In a discussion held after the end of the lectures it was
objected against me that the concept of duty was missing from my account. I answered: ``The concept of duty
does not belong to the doctrine of principles. For duty first arises when what is acknowledged in principle
is to be held fast and carried through in reality. I admire Kant's portrayal of the consciousness of duty,
which presents this significant side of ethical life in great, fundamental features. But Kant went wrong in
conceiving the concept of duty purely formally and not as conditioned by real ends. His psychological
explanation is not as good as his psychological description. But every ethical standpoint can take up Kant's
description of the consciousness
% --- p. 160 ---
\opage{160}of duty, as soon as the question is how the ends and principles set by the will are to be held fast and
carried through in the struggle of life.'' ---

The four concepts spoken of will come into application again and again in the particular questions within
special ethics. But it would not be expedient to take them as the basis for the division of ethics. They are
purely formal concepts, and the same real content would recur under all four points of view. Besides, the
relation between them depends in each single case on individual and historical conditions. What has
immediate value for the one man has perhaps for the other only mediate value and therefore requires a special
effort if it is to be held fast: what is a good for the former becomes a duty for the latter. Perhaps in a
third the character is so developed that a striving to assert what is valuable is released without effort.

Instead of the purely formal division that might be built on the fundamental concepts spoken of, it is more
expedient to follow a more real division\footnote[1]{On formal and real differences of value see \textit{Den
menneskelige Tanke.} p. 241---248. Cf. above IV, 1.}. \textit{The ethical world} extends from the
moment to mankind. What fills the moments is valued from the point of view of the whole of the individual
life, what fills the individual's life is valued also from the point of view of society, and every larger or
smaller society is valued in the end from the point of view of mankind. The most far-reaching opposition here
will be that between the single individual and society. The different kinds of society present common
features, and they all consist of individuals as their ultimate elements. The question then becomes that of
the relation between \textit{individual and social ethics.} Does the one of these comprise the other, or do
they stand side by side with equal right?

2. If we look first \textit{historically} at the relation between individual and society, we do not find at
earlier stages of development the individual as the bearer of special duties and rights. It is regarded only
as a social element. The weight that at different
% --- p. 161 ---
\opage{161}times and among different peoples is laid on individual virtues such as self-control, moderation and the
like finds its explanation in this, that these were qualities the tribe needed at the given time. The
admiration felt for the individual virtues is thus originally due to their social significance. Primitive
ethics is social ethics, and individual ethics develops out of it. As regards Greece, the philosophical schools
represent the development of an individual, in part an individualistic, ethics. Increasing weight is laid on
harmony in the individual's own nature, on the individual's disposition and character, and society is at
last regarded as a merely external means for the individual's self-development. In the Christian idea of the
single soul that is to be saved at any price there lies the same pointing toward the single personality. The
modern endeavors at emancipation set the individual more and more free over against the social forms, and
regard what goes on in it and with it as what is really decisive. And yet it will appear that social ethics
is always presupposed by individual ethics, not merely historically but also conceptually.

3. One might think that \textit{individual ethics} must be \textit{ethics proper,} since it is, after all,
always single, definite individuals that we have to do with in the ethical domain. Ethical judgments concern
willing and acting, and it is always single, definite individuals that will and act. A society acts only
insofar as the single individuals act. And each of these acting individuals seems in the end to have to do
only with itself. It is to act according to its own conscience; it cannot directly enter into what goes on
within others; it has only to see to it, then, that it does not itself suffer harm to its soul; and when
everyone does so, all will be well with society!

This way of looking at things has exercised great influence. It comes forth, as already touched upon, first
in Greek philosophy. It was especially the Cynics and the Stoics who maintained such a regard for self, a
higher self-preservation, as the most important thing in the ethical respect. In this respect they are the
forerunners of Christianity. In our
% --- p. 162 ---
\opage{162}modern time S. \textsc{Kierkegaard} is the great representative of the ethics of the single individual.
Under his influence the same line of thought is found in \textsc{Heegaard}'s book ``Om Intolerance,'' where the
need for intercourse with others is even characterized as a need ``to see, hear, enjoy, or to be admired,
flattered and worshipped.'' In \textsc{Henrik Ibsen}'s ``A Doll's House'' the same tendency comes out when
Nora, by virtue of the principle that we have duties to ourselves before duties to others, withdraws from her
very nearest relations --- in order to become a human being!

This conception, for all its one-sidedness, has proceeded from a significant ethical line of thought. The
single personalities are the living material in which the ethical has its reality. But it does not follow
that the whole of existence, society and all other individuals are to be made means to the single person's
self-development. That would bring with it an isolation, a one-sided absorption in oneself, that would easily
pass over into egoism. There is just as much an ascetic egoism as there is an enjoying or a power-seeking
egoism. The single person loses the capacity for self-surrender. The gaze is fixed on one's own welfare, not on
the great common welfare of which one's own is a part.

Anxious ideas are entertained here about what the self or the personality can bear and needs. One thinks one
can preserve oneself only by keeping one's attention uninterruptedly directed on oneself. But a self is not so
fragile a thing. It becomes fragile precisely through much care for itself. A sound and vigorous self is
reached by brisk and vigorous work on real tasks, and a large and comprehensive self is reached only when one
works on common tasks, when one makes oneself one with something greater than the self. Only through
self-forgetfulness does the self grow and develop. By living in others and in something other than itself the
individual acquires new and vigorous impulses for its own individual life too. It is raised above its
original narrow circle and filled with a richer content. Everything a man lives in and for through sympathy is
an enrichment and widening of his personality. It is here at every point both end and means. H. C.
\textsc{Ørsted}'s maxim: ``Forget your self, but do not lose it!'' aptly expresses what is at stake here.
% --- p. 163 ---
\opage{163}For the sake of individual ethics itself, then, it must not be made the only ethics.

The overvaluation of individual ethics is connected with an abstract way of thinking that draws man out of the
definite, historically given relations and thinks it has the true man only in the being thus isolated. But I
am not \textit{first} a man in general, and \textit{then} a man in these special relations. I am not
\textit{first} a man, and then a Dane, a citizen, the father of a family, etc. Individualism comes forth
naturally in times when the historically given conditions of life lose their absolute authority and do not
stand as the only possible and justified ones. It then has the great significance of nourishing a
consciousness of freedom that can become a fruitful source of new development. But not only does it bring with
it, as I have sought to show, great misgivings from the standpoint of individual ethics itself; --- it will not
be possible from such a way of looking at things to ground a conception of human social life by which this
becomes more than a mere means. Its interests and tasks could at most become exercises for the single
individuals, exercises that in and for themselves could just as well be exchanged for others.

An entirely different kind of individualism appears in the most recent times, especially in literary and
aesthetic form. Fear of leveling and uniformity and of the spiritual tyranny of the crowd has especially worked
to call it forth, and it appears as a cult of the strange, the solitary, the desert-like, the willful, the
daring. In the practical domain this leads to a cult of ``the superman,'' and the superman is thought of
especially as the superior power that is raised not only above prejudices and customs, but also above the
oppositions within which the life of the mass moves. What seems especially to awaken admiration of the
imagined superman is the contempt he is able, from his standpoint, to feel for the great crowd. It is this
factor of distance (``das Pathos der Distanz'') that is placed in the foreground. With it there is then often
bound up the longing and admiration that the weak and sick naturally feel for great strength and great health.
These different
% --- p. 164 ---
\opage{164}tendencies can be shown in \textsc{Friedrich Nietzsche}. Thus he demands \textit{(Jenseits von Gut und
Böse.} Aphor. 258) of a ruling class that it should not ground its value on the function it performs in
society, but that ``its fundamental conviction should be that society does not exist for the sake of
society, but only as a substructure and foundation on which a select kind of beings is able to raise itself
to its higher task and in general to a higher being.''\footnote[1]{I cannot here discuss whether in such
utterances we have Nietzsche's real fundamental thought. On this I must refer to my work \textit{Moderne
Filosofer} (1904).} Yet the problem of the relation between individual and social ethics comes back within
that select kind of beings, namely in the question of their relation to one another and to the society they
form. Of one of the supermen Nietzsche most admired, Napoleon, it has been said that he could not bear proud
souls. That became his nemesis. Many of Nietzsche's followers escape this nemesis for good reasons. There
betrays itself in them, on the whole, a fairly disinterested feeling, insofar as they themselves seem very
far from occupying the high station from which the great contempt could be motivated.

In a more philosophical way individualism appears in idealistic anarchism, as it is taught, for example, by
the German writer \textsc{Bruno Wille}. Bruno \textsc{Wille}'s book \textit{Die Philosophie der Befreiung durch
das reine Mittel} (Berlin 1894) gives an interesting and instructive account of this standpoint. It is
demanded here that individuality shall develop from within, and that all ideals shall be self-created.
Protest is made against all spiritual and physical compulsion, and such compulsion Wille finds also in moral
influence on the part of others. Since he conceives conscience as a merely external, social product, he
declares that the free man of reason must be without conscience.

Bruno Wille builds on a thought that has also played a significant part in the history of Danish
thought\footnote[2]{See my book \textit{Søren Kierkegaard som Filosof.} Copenhagen 1892. p. 22---27;
98---107.}. The personality must itself produce --- or produce anew
% --- p. 165 ---
\opage{165}--- the truth that is to be truth for it. Ethical judgments must spring from the single person's own
experience and conviction; otherwise they are only an echo of the judgments of those around him. Everyone is
to be his own inner judge, his own authority. All influence from without is to aim at producing the capacity
for ethical independence. But what idealistic anarchism overlooks is the pedagogical significance that
imitation and obedience, in general the relation of authority, can have. The single person needs help from
others in order to become free and independent. This help can very often best be given by another way than by
moral judgment. Direct moralizing often hinders the awakening of self-activity, since the ethical valuation is
conveyed to the individual wholly finished and formulated instead of being produced by it. Yet direct
moralizing really is a ``pure means'' (cf. VII, 5) when it aims at awakening self-knowledge and the impulse to
development. This awakening may often be painful; but there are pains that are necessary transitions to a
higher stage.

In vain will one strive to reach the highest by the way of isolation. Not only does the single person receive
his education in society, but his content and his tasks too he receives only by devotion to the life of the
race. Only here does he come to share in the higher ends with which he can grow. In the cult of contempt to
which aesthetic individualism in the most recent times is inclined, the impossibility of isolation comes to
light through a peculiar nemesis. When it is what Nietzsche has called \textit{the pathos of distance} that
determines the essential content this standpoint acquires, then it leads precisely to dependence on the great
crowd; for there must be someone one can despise, someone from whom one can feel oneself at a distance. The
exalted station is used essentially to look down from it on those who have not reached it. One forgets that
true spiritual power and superiority show themselves in the capacity to carry not merely one's own life but
the lives of many others. There is, says \textsc{Spinoza}, no way in which a man can better show his strength
of mind than by developing others to live according to their own free knowledge.

% --- p. 166 ---
\opage{166}4. As a counterpart to this conception there appears the one that demands \textit{complete giving up of
oneself} and regards \textit{self-forgetfulness as the only virtue.} Within modern philosophical ethics such a
standpoint has been asserted by J. G. \textsc{Fichte} and A. Comte. True life consists, according to Fichte,
in sacrificing oneself for the race: ``There is only one virtue, to forget oneself as a person, and only one
vice, to think of oneself\dots{} He who at all thinks of himself as a person and desires life, existence and
enjoyment otherwise than in and for the race, is --- by whatever good deeds he may otherwise seek in other
fields to hide his deformity --- nevertheless at bottom a common, small, bad and therefore unblessed
man''\footnote[1]{\textit{Die Grundzüge des gegenwärtigen Zeitalters.} Berlin 1806. p. 70 f.}. Comte
designates by the word altruism, which he introduced\footnote[2]{From alter (the other, i.e.\ the
neighbor), as egoism from ego (I).}, complete absorption in others, the absolute opposite of egoism. Comte
declares, to be sure, that altruism does not suppress self-assertion; but the latter's justification is to
rest only on its being sanctified by working in the service of the former: both duty and happiness consist in
\textit{vivre pour autrui}\footnote[3]{Politique \textit{positive.} IV. p. 45---49. --- There is on this
point a certain opposition between Comte's earlier and his later standpoint (\textit{Den nyere Filosofis
Historie}\textsuperscript{2} II. p. 346 f. and 357), which forms an analogy to the opposition between Fichte's
two stages (ib. 142 f. and 151 f.).}.

This conception too was called forth by definite historical conditions, in the case of Fichte and Comte by
the opposition to the individualism and the morality of happiness of the 18th century. In its way it starts,
like the individualistic one, from a sharper opposition between regard for self and regard for others than
need exist according to the real relations. From the principle of welfare it follows that the single person
is only one among many; but it follows also that the single person really is one among the many. Everyone
counts, and the giver is not to blot himself out for the sake of the receiver, any more than the receiver is
to be made a mere means to the development of good qualities in the giver. Only isolated self-satisfaction,
entangled in itself, conflicts with the principle of welfare. If the impulse
% --- p. 167 ---
\opage{167}to self-preservation, self-assertion and self-development were of evil, our inmost nature would be evil,
and all ethics would then be self-contradictory and impossible.

We gain powers and means to work for what is different from ourselves and greater than ourselves only by
preserving and developing our own being. If ethics condemns the instinct of self-preservation, it condemns
its own means. Precisely when one takes one's stand on the ground of altruism, it becomes a special task for
the single individual to strengthen and develop his powers and faculties. Self-assertion and self-development
work in the same direction as altruism. Indeed, as \textsc{Spencer} remarks, the happiness that can radiate
from a sound and fresh individual over others is often worth more than that which one can consciously and
with great sacrifice prepare for them. ``Though our ethical reasonings take little account of it, it is yet
evident that since happiness and sorrow are infectious, the care for self that leads to health and high
spirits is a benefaction to others''\footnote[1]{\textit{Data of Ethics.} London 1879. p. 194.}. % translated from Høffding's Danish
And this holds not least of the higher, ideal development of the personality through work for interests that
perhaps do not directly benefit others, as for example in artistic and scientific endeavors.

In another respect too ethics, by taking absolute altruism as its basis, would come into contradiction with
itself. All sacrifice for others aims, after all, at supporting their self-preservation, self-assertion and
self-development, and how can these be justified in the receivers if they are not so in the givers? Absolute
sacrifice for others can even have a harmful influence on them. He who constantly receives sacrifices on the
part of others easily becomes an egoist, if he does not go quite to sleep without knowing what is going on.
He is made a passive being, who takes what is given without entering into what is sacrificed for him. His
capacity for work is weakened, he becomes a parasite, not an independent, self-active being. The right
sacrifice must be the one that robs neither giver nor receiver of their independence,
% --- p. 168 ---
\opage{168}but makes them both independent members of the kingdom of personalities.

Such standpoints as Fichte's and Comte's come near to conceiving society or the race as a mystical whole that
has an existence apart from the single individuals. A similar tendency appears in \textsc{Hegel}\footnote[1]{See
\textit{Den nyere Filosofis Historie}\textsuperscript{2} II. p. 181.} and in the most recent times in
\textsc{Wundt} and Paul \textsc{Carus}. In \textsc{Wundt}'s conception\footnote[2]{Wundt: Ethik. Eine
Untersuchung der Thatsachen und Gesetze des sittlichen Lebens.\textsuperscript{3} II. p. 116.} ``the public
good'' does not consist in the sum of welfare of as many single persons as possible, and general progress does
not consist in the progress of as many individuals as possible. For every individual is a creature of a day:
``However richly blessed and however perfect the individual existence may be, it is yet only a drop in the
ocean of life. What can individual happiness and individual pain signify for the world?'' And
Carus\footnote[3]{P. Carus: \textit{The Ethical Problem.} Chicago 1890. p. III cf. p. 33, 38, 40. --- Cf. the
discussion between Carus and me in The Monist. July 1891. I cannot see that Carus by his closer explanations
robs his idea of its unclearness, so long as he (like Wundt) denies that in ethics one must always think of a
sum of individuals (of greater or less extent, perhaps of an extent that cannot be surveyed) when one wishes to
operate with the concept of society.} designates social life as ``superindividual''; the end of ethics is
according to him ``neither the welfare of the self nor that of other individuals, but welfare for the
superindividual interests.''

If the concept of society is to be usable scientifically, it must be so in such a way that at every point one
can decide what group of individuals it represents. The great significance of this concept rests on its
expressing the common and abiding interests of the individuals existing simultaneously and successively, as
opposed to the interests of single individuals or to the interests of a small group or a single generation.
The ethical consideration leads (at the standpoint where we have taken our stand here) to regarding human
actions not merely with regard to the welfare of the single person or of a limited circle or age, but with
regard to the welfare of society or of the whole
% --- p. 169 ---
\opage{169}race, so far as we can trace its conditions. But as soon as it proves impossible to convert the concept of
a society into the concept of a group of individuals living under certain definite conditions, that concept
gives us no ethical instruction; no ethical norms can be derived from it. Mysticism steps into the place of
thought.

Such a mysticism may have its significance. It may prove impossible to form a clear concept of the great mass
of human interests in which an action or an ordinance of life may under certain conditions intervene in a
decisive way. Expressions such as society and race designate very well what is unconcluded and cannot be
surveyed in many of these series of effects. And they designate also the possibilities or dispositions, the
potential energy, that human work can bring about and accumulate, and which perhaps only in distant times can
be converted into actual values that are felt by real, personal beings. But when one declares such a
conversion impossible or unnecessary, one gets stuck in mysticism. A welfare that is not at some stage or other
the welfare of definite individuals is a self-contradiction. And an action or an ordinance of life that does
not at some point of time or other lead to welfare for definite individuals has no ethical value. A possibility
that can never become reality is an impossibility.

The single individual is only a drop in the ocean. And the ocean does not exist for the sake of the single drop
alone. But what is an ocean that does not consist of drops? And will not the whole ocean be clear when each
single drop is clear? And only then is it perfectly clear.

Just as there are those who cannot see the forest for the trees, so there are those who cannot see the trees for
the forest. In ethics this shows itself in human striving being regarded as a means to superhuman ends. Every
ethics that will not give up the possibility of progressive verification must hold fast that society is always
represented by definite individuals. It need not therefore overlook that ethical activity, like all
manifestation of force, is connected with the whole process of the world.

% --- p. 170 ---
\opage{170}5. One might try to take a middle road and set \textit{individual and social ethics side by side as two
independent domains.} So \textsc{Bentham} and \textsc{Stuart Mill} assume that there is a domain where the
individual rules alone, and another where social regards rule. \textsc{Mill} even holds that such faults as
concern only ourselves and have no influence on the welfare of others are not really ``moral'' faults. Foolish
conduct, bad taste, the dominion of low and sensual impulses, lack of personal dignity and of self-respect may,
according to \textsc{Mill}, be present without the interests of others being in the least touched thereby. No
one, Mill holds, is in the least responsible to his fellow men for his self-development, ``since it is not for
the sake of the good of society that one is to answer for it''\footnote[1]{\textit{Om Friheden.} Danish transl.
p. 141 f. --- \emph{Bentham}'s statements about ``private ethics'' \textit{(Principles of Morals and
Legislation.} XVII, 3, 20) are not quite clear, and there appears in him a tendency to subordinate social ethics
wholly to individual ethics. Cf. above III, 16. --- In some Danish writers of the most recent times (Starcke
and Bang) a dualistic theory similar to Mill's has appeared. See above III, 21.}.

What Mill combats is the intervention of state power and of public opinion (the ``moral police'') in the
individual's inmost concerns. He fears spiritual tyranny and unhealthy compulsion if one starts from the
assumption that there is nothing in a man's conduct that concerns him alone. But he mixes up two things here. It
is one thing that the individual through and through, even in his inmost being and with all his qualities, must
be regarded, and regard himself, as a member of the race, and that his willing and acting are therefore subject
to an ethical valuation whose standard is given in the principle of the general welfare, --- quite another thing
\textit{who} in the single case is to appear as the proclaimer and enforcer of the ethical judgment. Ethically
regarded, everyone stands and falls first and foremost before his own inner tribunal in conscience, and ethics
is a doctrine of what judgments this inner tribunal must consistently pass. The degree in which a man develops
his faculties and powers must above all else fall under this tribunal. A. S. \textsc{Ørsted}
% --- p. 171 ---
\opage{171}has rightly said that worse than committing a breach of the right of property is squandering, through
self-interest, vanity or sluggish inactivity, the strength with which we could work in the service of
progress\footnote[1]{\textit{Om Grænserne mellem Teori og Praxis i Sædelæren.} (Eunomia. Første Del.
Copenhagen 1815). p. 105.}. It is not therefore said that it is expedient to allow other men a right to meddle
in the single person's inner ethical household. One can grant the one and deny the other. Mill, on the other
hand, denies the ethical significance of self-development because he fears thereby to give others the right to
penetrate into the individual's private conduct. But he could reach his goal, namely the assertion of individual
freedom against social pressure and moral police, without excluding self-development from the ethical domain.
Nothing is more questionable than the great misunderstanding that ethics and public opinion are the same, and
that a relation ceases to be subject to ethical judgment because no external authority has the right to call it
to account\footnote[2]{In his excellent work \textit{The Utilitarians} (III, p. 286---299) \emph{Leslie
Stephen} has, in a direction similar to mine, criticized Mill's doctrine of the boundary between individual and
social ethics.}.

There is nothing at all in the free development of the individual that cannot acquire ethical significance.
What appears in the single person's inner life may perhaps be the first germ of great social processes. Every
individual development can produce new forms and models; the lonely paths the single person walks may later
become highroads. The energy, the fidelity and the wisdom he uses to seek out and clear these paths may become
of decisive significance. It is not possible to separate the self-development and the inner household of single
persons from the life of society and of the whole race. Even in play, in art, in faith and in the free life of
mood of the single person the ethical is potentially present. Just as in external nature there is only
apparently an absolute rest, while energy is always present, though often in so-called potential form, so there
is in human life no point that may not possibly
% --- p. 172 ---
\opage{172}be subject to an ethical judgment. Some have particularly wished to maintain the erotic relation and
religion as private matters. This has no doubt been done for tactical reasons, namely in order not to be
molested by morality in aesthetics, and in order to be able to move in social policy without colliding with
ecclesiastical faith. In reality both the sexual relation and the religious relation are of very great ethical
significance. Precisely in these domains qualities such as courage, fidelity, honesty and love of truth have the
greatest value; precisely here baseness of disposition, wavering and timidity can become fateful. And the mere
establishing of the presence of such qualities is, after all, an ethical judgment.

6. The inner relation between individual and social ethics is really given with the setting up of the principle
of welfare itself. When the foundation of valuation is universal and disinterested sympathy, and the standard
the general welfare, it follows that individual ethics must be subordinate to social ethics, without therefore
vanishing in relation to it. The full grounding of the individual virtues, duties and rights can be reached only
by regarding the individual in the social milieu, as a member of the struggling and striving race, whose
inheritance it has received in its nature and in its external conditions of life, and whose development it is
to carry further so far as it is able. Through taking part in the life and work of the race it will develop its
own being.

The closer determination of the relation between individual and social ethics can be won with the help of some
propositions that can be derived from the general principle of welfare.

The simplest form in which this principle appears is that the infliction of pain must always be grounded, while
the production of pleasure and joy is justified in itself. The application of this proposition is simple and
easy only in the very simplest cases, but on closer consideration one will find the proposition as the ultimate
presupposition in composite cases too. The more remarkable is it that the proposition has been found
questionable; it really says only what everyone --- with the exception of the absolute ascetics, if there are
any --- must grant. Even the ascetics have as a rule regarded pain as a means of penance and practice, not as an
end in and for itself.

% --- p. 173 ---
\opage{173}In this simple proposition, on closer psychological consideration, another very important proposition can
be found enclosed.

In pain there makes itself known a hindering or a dissolution of life. This is clear in bodily pain, e.g.\ the
pain of wounds, toothache, the pains of stone, cancer. Pain here answers to a transitional stage between
unhindered, undivided life and the complete hindering and dissolution that comes with death. A Greek writer has
said that if man were an absolutely simple being, he would not be able to feel pain: for then no dissolution
could take place. But the converse holds too: only because the elements of life are always in the most intimate
unity and coherence in the living being and do not form a mere aggregate, only for that reason can pain be felt.
This holds for the life of the soul as well as for bodily life. Spiritual pain too expresses a hindering or a
dissolution. Grief at a loss appears now as a hindering, the thought of the irrevocability of what is lost
displacing all other thoughts, now as a dissolution, the opposition between the value of what is lost and the
loss threatening to burst consciousness. Doubt is a state of pain, consciousness being drawn in opposite
directions by conflicting possibilities and tendencies. The pain of remorse is conditioned by the violent
opposition between our acknowledged ideal and the thought of our actual willing and acting. On the other hand
everything that furthers the unity and harmony of the life of consciousness is bound up with the feeling of
pleasure. Our life of consciousness involuntarily strives to overcome the sufferings and pains that cannot be
removed by summoning higher, comprehensive powers. A new feeling forms, in which the bitterness is taken up as an
element in a more composite or mixed feeling. Most of our so-called higher feelings have this composite
character. But at all stages our life of feeling stands as the expression of the advance or decline of our life.
The difference between the feelings rests in this respect on the more or less central position they occupy
within the whole of our life\footnote[1]{Cf. on the psychological points of view used here my
\textit{Psykologi} VI B, C and D.}.

% --- p. 174 ---
\opage{174}By virtue of this close connection between feeling and the individual's being as a whole, the
principle of welfare leads directly over to \textit{the principle of the free personality,} i.e.\ to the
maxim that no personal being may be treated and regarded merely \textit{as a means,} but must always at the
same time be an end. In personal beings we have, after all, the central places to which the effects of our
actions propagate themselves, and where their most important significance appears. At these points in the
world the value of life is felt, --- and the value of life exists (as actual value) only when it is felt. An
injury or a hindering here will therefore necessarily diminish the value of life, and if no suffering is felt
at the hindering because this produces dullness, then that is precisely the greatest injury to the principle
of welfare, since the very capacity to feel value is abolished. --- The principle of welfare presupposes that
we are able to transport ourselves to the strange central places and to consider what influence the actions
that are to be judged exert there. On this capacity, after all, rests also the possibility of sympathy proper.

This principle, which thus follows from the principle of welfare, was set up by \textsc{Immanuel Kant} with the
help of an entirely different grounding, which it will be of interest to consider. Since Kant regards the
ethical law as purely formal and a priori, he cannot really derive definite ends for ethical action from it.
Every end that was to be set would lead beyond the pure form of the law. And yet Kant sees that every action
must have an end. The difficulty that thereby arises he thinks he solves in the following way. Since the human
personality has the capacity to subject itself with freedom and inwardness to a universal law, such a
personality stands as an absolute end and ought never to be degraded by being treated as a mere means. It is
the ethical in its inwardness and sublimity that is man's patent of nobility, which under no circumstances may
be taken from him by treating him as a mere means. In \textit{Kritik der praktischen Vernunft} (Kehrbach's ed.
p. 106) Kant expresses himself on this as follows: ``In the whole world everything one wills, and over which
one has power, can also be used \textit{merely as a means;} only man, and with him every
% --- p. 175 ---
\opage{175}rational being, is \textit{an end in itself.} For by virtue of the autonomy of its freedom it is the
subject of the moral law, which is holy. Precisely on account of this autonomy of freedom every will, even the
single individual's own will directed upon itself, is limited by the condition of agreement with the autonomy of
the rational being, --- namely so limited that this being cannot be subjected to any purpose that is not
possible according to a law that could spring from its own will, --- thus so that it is never treated merely as
a means, but at the same time as that which is itself an end.''

It is Kant's merit to have deepened the concept of personality, since he stresses at once the inwardness and
the universality of personal life. The personality is an inner world that follows its own law and yet can
stand in connection with the great world. But the question is whether Kant has here applied this significant
thought in a justified way. Personal beings we learn to know only by experience, and only by experience ---
through the interaction that social life brings with it --- do we learn to put ourselves in their place and to
feel how life takes shape in them. But from the purely formal law neither the capacity for this nor the interest
in it can be derived; no definite content whatever can be derived from it. Perhaps that principle of
personality stands as all the more significant for having developed on the ground of real life and of history,
even before we can assign it its systematic place in ethics. That it follows directly from the principle of
welfare there can be no doubt.

It was further a one-sidedness in Kant that the concept of personality was restricted to the strictly moral.
In all spiritual fields and in all relations of life the personality can appear, and the law of its working can
spring from itself. Everywhere --- not only in the moral field --- self-activity is the mark of personality; it
carries on the struggle for life in its own way, from its own center. In the rhythm of pleasure and pain there
expresses itself at all stages of life the favoring
% --- p. 176 ---
\opage{176}or the hindering of this activity, perhaps awakened from without, but determined by an inner law.

But although Kant's grounding cannot be acknowledged, it is nevertheless a great merit he has earned by setting
up the principle of the free personality. And even if one grants the validity of the grounding he has
attempted, it is of no slight importance that the important proposition can be reached by different roads. Both
according to Kant's grounding and according to the one I have attempted, a personal being stands both as an
end and as a means. On this rests the significance of the free development of the personality.

Freedom is an \textit{end:} for the unhindered unfolding of faculties and powers is a good, since it is bound up
with the most inward satisfaction, while the feeling of compulsion and hindrance is a feeling of pain. Of
freedom we ask: \textit{why not?} Of compulsion we ask: why? The burden of proof lies on those who would use
compulsion. Only as a means and preparation for the free unfolding of the powers can compulsion and barriers be
justified, as for example when the boundless unfolding of one individual would make the development of one or
more other individuals impossible.

But freedom is also a \textit{means:} the free unfolding of individuals produces new possibilities and leads to
the discovery of new directions for the life of the whole race. The new always begins at single points and then
spreads to larger circles. Personal beings are not merely the central places where the value of life is felt,
but also the places from which the movements of life ever anew go out. Through freedom new centers of
independent activity are created, and since the race consists of personal centers, the life of the race too
thereby becomes richer and more vigorous. Herein lies the great significance of the modern movements of
emancipation: that the slave, the peasant, the worker and woman come to stand as free and independent members of
the race can only make its life richer and more fruitful. Emancipation is demanded by society's own interest.
This thought \textsc{Stuart Mill} enjoins precisely in his book \textit{Om Friheden} (p. 115---119), in
remarkable contradiction to his assertion of individual self-development as indifferent in the moral respect.
% --- p. 177 ---
\opage{177}He demands the protection of the independence and originality of single personalities, because the impulse
to everything wise and noble comes from single individuals, and because the spiritual dominion of the masses
would lead to general mediocrity. Remarkably enough, at the same time as the work named, an analogous thought
came forth in \textsc{Charles Darwin}'s doctrine of the development of organic life through the new means and
ways for the struggle for existence that individual variations make possible. Only where such variations
appear (from one cause or another), only there can the natural selection exercised by the conditions of life
come to work. Both biologically and sociologically, then, the free unfolding of the nature of single individuals
proves to be a significant means of progress.

We can here throw light from a new side on the relation between motive and action. As we have seen (III, 18),
valuation consistently goes back from the action to the motive, because it must be the main thing to hit the
source from which the action springs. We started there from the practical significance of valuation and of
valuing judgments. Here we can supplement this consideration with the help of the principle of the free
personality. For the inner life of a personal being ought not --- by virtue of this principle --- to be made a
mere means to the production of external effects. It must not be a mere wheel in a machine. Even where the
individual is regarded as a member of a larger whole, its action is to spring from its peculiar being, from its
inner disposition. Only then can the action be called perfect, when it at once springs from personal
peculiarity and intervenes in the great whole in a furthering and happiness-bringing way. The duties and virtues
that society demands of the single person are to be at the same time means and forms of the individual's own
development.

It is the great pedagogical (or perhaps better: psychological) art, which there may be occasion to practice both
on the part of single individuals and on the part of the power of society: to call forth motives that make
self-activity possible. If by compulsion one would understand every intervention on the part of others, such a
motive-awakening activity would be a compulsion, even if it
% --- p. 178 ---
\opage{178}were exercised ever so indirectly and Socratically. It is more natural, however, to understand by
compulsion merely an external influence that calls forth pain or fear (cf. V, 2 c). And the motive-awakening
intervention does not necessarily presuppose compulsion in this narrower sense of the word. A responsibility
follows with all such intervention; but when it consists in compulsion proper, it needs a quite special
justification. There is much to learn here from the theoretical anarchists, who reject all authority and all
spiritual and physical compulsion. When one asks how such an anarchist can bring up his children, he will answer
that he seeks indirectly to awaken the involuntary feelings that in his opinion make up the ethical foundation,
so that the children come to do what is right without command and without compulsion. He will thus try to carry
through what Rousseau called ``negative education.'' The question (which the anarchist settles rather quickly)
is how far all motive-awakening influence can take place in this way, and in how high a degree compulsion proper
can be restricted. But here we have to do only with the general principle, that is, with the distribution of the
burden of proof. ---

But here the question arises again whether in every single individual there is the same possibility of the
arising of motives for the actions that must be demanded on the part of society. When one does not dogmatically
start from the assumption that all men are wholly alike in the respect named, there arises here, as shown
earlier (see IV, 2), the unavoidable demand for the individualization of the ethical law. The ethical law cannot
then be completely given by a determination of ``the conditions of social life,'' regardless of how the single
individuals are placed with regard to satisfying these conditions. Or rather: when one does not start from a
purely external relation between individual and society, one cannot fix the conditions of social life without
including among them the peculiarities of the single individuals in qualitative and quantitative respect. The
condition of a sound social life is, after all, first and foremost that there be a harmonious relation between
the development of the single persons and the demands of the whole. And such a relation presupposes in turn that
the task set the single
% --- p. 179 ---
\opage{179}person answers to the qualitative and quantitative presuppositions he brings with him. Since society
consists of all the single persons, this is not an external demand made on it. Only there is a real society
present where the single person can at every point be both end and means. A conflict between the demand of
society and individual nature therefore points to an imperfection in society itself, not merely --- or not
always --- in the individual.

Not only the usual moral sermon overlooks this. But every conception that lets the ethical come to man ``from
without'' --- whether one builds on the principle of authority, or on a conception of the race as something
other and more than an organic whole of individuals, or on public opinion as the sum of all ethics --- every such
conception must of course reject this consideration. One thereby cuts oneself off from treating some of the
deepest-lying personal problems. An illusory and purely formal conclusion of the theory is bought by setting
aside tasks that can, to be sure, only be solved approximately, but which it is nevertheless of great importance
to hold fast. As at its first foundation (the motive of valuation), so too in the consistent carrying through of
its principles ethics stands before the irrational (to use an analogy taken from arithmetic). There can here be
question only of finding as many decimals as possible. ---

Only by the consideration carried through here can there be a harmonious relation between individual and social
ethics. They point mutually to one another, the individual being regarded as an organic member of society, and
society as an organic whole of individuals. We set out individual ethics first because it presents the simplest
relations; but we constantly hold fast to the social point of view as the one lying at the foundation, just as in
social ethics we shall always take as the basis the point of view that the highest end and the best means of
social life is the free unfolding of the peculiar powers of the single individuals.

% --- p. 181 ---
\opage{181}\partitle{INDIVIDUAL ETHICS}

% --- p. 183 ---
\opage{183}\chapterhead{IX.}{THE DIVISION OF INDIVIDUAL ETHICS.}

1. The relation fixed above between individual and social ethics leads naturally to the division that must be
followed within individual ethics.

Since from the very first we regard the individual as a member of society, and within society as both end and
means, two tendencies appear in personal life that must be acknowledged and furthered: self-assertion and
devotion.

The single individual is an end, since it is the very representative of society; in its life the life of
society stirs, and by injury to it the life of society is injured. It is one of the end points in which
development runs out. The more it is able to assert itself and to unfold its faculties and impulses, so that
life stirs in it with full strength and harmony, the higher a stage is reached --- not merely for the single
person himself, but for society. Society consists of individuals, and its life must therefore be the fuller and
more vigorous the more each single person is able to develop what lies in his nature. The single person is a
little world whose existence and development is an end valid in itself. \textit{Self-assertion} --- in its
different forms: self-preservation, self-control and independence --- thereby comes forward as an essential
virtue. It is a duty the single person has at once toward himself and toward society. Regard for self and regard
for society here go so immediately into one that it
% --- p. 184 ---
\opage{184}may become indifferent which of the points of view one stresses. No appeal to a conscious interest in
society or to the abstract principle of welfare is needed to show that self-assertion must be regarded as
valuable by others too than by him who practices it. A vigorous and harmonious unfolding of the personality is
the object of immediate sympathy and admiration. As \textsc{Hume}\footnote[1]{\textit{Enquiry concerning the
Principles of Morals} VI, 1. Cf. \textit{Treatise} III, 3, 1.} already remarked, it is still harder to explain
from egoistic interest the admiration in others for the virtues that benefit an individual himself than it is to
explain the acknowledgment of social virtues in this way. We do not necessarily regard others' self-assertion as
a mere means. It is an unnecessary detour to take when Hobbes\footnote[2]{\textit{De cive} III, 25.} explains
the abhorrence of drunkenness by its easily leading to breaches of the peace and in general making it
impossible to rely on the individual. At the sight of a lack of self-control in others we feel immediately a
flaw in the vigorous and harmonious self-unfolding of which we wish others' conduct of life to give us a
picture. We recognize that where such a picture meets us, there at this point of the world the highest has been
reached. The tacit presupposition, of course, lies at the basis that the strength and harmony of the single
person do not conflict with the conditions under which others can strive to reach a similar goal for their
part, but are perhaps even a means to it. There is, after all, as developed earlier (VIII, 5), hardly anything
within individual development, even in its freest peculiarity, that cannot become of social significance; in all
self-assertion a potential social ethics is present. But in acknowledging self-assertion as an end ethics points
beyond itself. If ethics exists for the sake of life, and not life for the sake of ethics, there will be a
limiting point where the validity of the norms ceases, because the immediate unfolding of life needs no
justification. Where this limiting point lies, and how near one is to it in the single case, may be difficult
to
% --- p. 185 ---
\opage{185}decide when doubt arises. But immediate self-assertion has value as an end, not merely as a means; about
this there can be no doubt on the whole point of view that is adopted here.

Since the individual, however, is always only one among many, and since the interests of the life of society
make up a larger whole than the sum of the interests of his isolated life, \textit{devotion} may come into
opposition to self-assertion. Devotion to a larger circle of interests may bring with it a widening of interest
that hinders the individual harmony for which self-assertion strives. Self-assertion has to do with the
individual's own little world and regards it as a closed whole, but devotion demands that this little world be
brought into connection with a larger world, and this may bring with it a provisional disturbance of the inner
harmony of the little world. It is not said, after all, that the ordering of the little world fits without more
ado into the great order of the world. Closing off and widening may stand in opposition, even in conflict, with
one another. Therefore self-assertion and devotion must be regarded as two different tendencies in character, two
different virtues, and we find accordingly that different weight has been laid on them in the history of ethics.

Yet it is not necessary that they come into conflict with one another. The need for unity and continuity that
already expresses itself in self-assertion, insofar as the single personality is regarded as a closed whole, can
lead beyond the individual totality and appear as a need for connection with a more comprehensive whole. The
single person can perhaps gain coherence within his own little world only by attaching his interest to something
lasting, something ever advancing, something that cannot be surveyed. He can perhaps be true to himself only by
being true to something greater than himself. In this way devotion would be the continuation of self-assertion.

It follows from this, too, that it would be wrong to set up devotion as the passive virtue in opposition to
self-assertion as the active one. Sympathy, the feeling that leads immediately to devotion, is not a tendency to
passive sinking down that would form an opposition to the energetic impulse of self-preservation. In the most
recent times there has been proclaimed
% --- p. 186 ---
\opage{186}the right of the stronger, egoistic ruthlessness, as the highest thing, as the mark of character of ``the
superman,'' and it has then been regarded as a slave revolt that love of mankind was proclaimed as an ethical
principle. In opposition to this assertion, made by \textsc{Nietzsche}, it must be maintained that sympathy or
love of mankind, where it is genuine and original, is precisely an expression of strength, of spiritual power. It
presupposes that not all energy is consumed in the service of one's own purely individual needs, but that there
is a surplus that makes it possible to feel pleasure or pain at others' fate, even if this does not intervene in
one's own individually closed existence. The mind then has a greater fullness at its disposal than in isolated
self-assertion. On account of this surplus fullness of strength and interest the individual has in genuine
sympathy a true superiority. Its conduct in devotion is determined from within and is independent of others'
hatred or love, of their contempt or admiration. It is --- to use an image of \textsc{Marcus Aurelius} --- like
the pure and strong spring that lets its stream well forth even if one throws dirt and stones into it; it washes
them away, and its purity is not changed thereby.

Already \textit{primitive Christianity} described love as a power that is independent of whether others love us
or not. The love that is only ``reactive,'' called forth only by others' services, it declares insufficient. True
love makes no distinctions, comprises all --- just as the sun and rain of heaven fall to the share of all. Love
bears all things, endures all things, overcomes all things! The long-suffering of love lets it endure in spite of
the greatest resistance. (The Gospel of Matthew V. --- 1 \textit{Corinthians} XIII. --- \textit{De imitatione
Christi} III, 5).

In modern times this conception of sympathy as an expression of strength is found in Spinoza and Rousseau. For
Spinoza the spiritual strength that makes up the essence of virtue appears partly as self-assertion or courage
of life (animositas), partly as high-mindedness (generositas), and the high-minded man seeks to help others and
to join them to himself in friendship, for ``souls are conquered not by arms, but by love and
high-mindedness.'' Rousseau explains love of others from the overflowing strength of self-assertion,
% --- p. 187 ---
\opage{187}which involuntarily spreads and comprises others, provided only they resemble us, because it is too great
for its own needs. The difference between ourselves and others is not heeded; the fullness of life goes beyond
all barriers. ``La force d'une âme expansive m'identifie avec mon semblable!'' Love is, as Rousseau expresses
himself, a consequence of love of oneself (amour de soi), which he distinguishes from self-love (amour propre).
Self-love sets barriers, since man compares himself with others and separates himself from them. But thereby we
become precisely dependent on others, which we are not in love. Later a similar conception has been put forward
by \textsc{Jouffroy} and other French philosophers, especially Guyau, to whom \textsc{Krapotkin} attaches
himself\footnote[1]{Cf. my \textit{Psykologi.} VI C, 3. 7. --- \textit{Den nyere Filosofis
Historie}\textsuperscript{2}. I p. 326; 494 f. --- \textit{Jean Jacques Rousseau og hans
Filosofi}\textsuperscript{2}, p. 101---105. On \emph{Jouffroy} see \emph{Bouillier}: \textit{Du plaisir et de la
douleur.} Paris 1877. p. 169 f. --- Guyau: \textit{Esquisse d'une Morale sans obligation ni sanction.} Paris
1885. --- \emph{Krapotkin}: \textit{Anarkistisk Moral.} (Danish transl. in ``Proletaren'' 1896). --- Cf. also
\emph{William James}: \textit{Varieties of belief.} p. 279---284.}.

There is of course a soft and passive kind of feeling that also goes under the name of sympathy and love. And
there is a sentimentality that enjoys one's own supposedly great sympathy, just as another kind of sentimentality
enjoys one's own supposedly great pains. But these forms may not be regarded as typical, and the last-named is
even rather a kind of egoistic self-enjoyment. ---

A characteristic difference between self-assertion and devotion consists in this, that in the former the
individual stands as an end and yet works as a means, since society is strengthened and developed by the
self-assertion of its single members, --- in the latter the individual works as a means and yet becomes an end,
since its personal life becomes richer and fuller by getting a larger field to work in. This difference also
shows the inner connection between them. If we conceive \textit{justice} in the way indicated earlier (III, 9),
we can find in it the harmonious unity of self-assertion and devotion, of the tendency to close oneself off in
oneself and to open oneself to a larger connection.
% --- p. 188 ---
\opage{188}As regards psychological origin, justice may have its root in self-assertion, namely when this is bound
up with the acknowledgment of others' equal entitlement, an acknowledgment that is perhaps first extorted by
power and authority, but which through displacement of motive can pass into the individual's flesh and blood.
But it may also have its root in devotion, namely when this is bound up with an understanding of the peculiarity
and significance of the single personalities (one's own personality too). Self-assertion and devotion are each
by itself more elementary tendencies. Justice is the more comprehensive tendency that sets the crown on the
development of ethical character. Justice presupposes that the individual asserts and develops himself not merely
without hindering the development of others, but in such a way that this is thereby furthered, --- and that the
individual devotes himself to comprehensive interests of life in such a way that his own personality is thereby
developed and asserted\footnote[1]{Already in Greek ethics justice appears in close connection with love of
mankind. Cf. \emph{Aristotle}: \textit{Ethica Nic.} VIII, 1. 13. 14. \emph{Cicero}: \textit{De finibus.} V, 23, 65.}.

Not only self-assertion and devotion are here in harmony, but also knowledge and feeling. A great end stands
before thought, and this seeks by the way of experience and inquiry the means to realize it. It is a great work
of distribution that is to be done in the service of the race, and the distribution is determined by a clear
understanding of the need and welfare of the single personalities. ---

Individual ethics, according to this whole consideration, is naturally divided into the doctrine of
self-assertion and the doctrine of devotion. When these two tendencies are depicted in such a way that a
harmonious union between them becomes possible, the chief ethical virtue, justice, will thereby also have been
depicted.

2. In the different individuals these tendencies play a different part. There are natures that achieve the
highest they are capable of without the ethical feeling stirring in them as a special feeling beside the others.
Self-assertion and devotion unfold in them without special motives and without
% --- p. 189 ---
\opage{189}conscious effort. The possibility of such ethical Aladdin-natures cannot be denied. There may be those in
whom such an Aladdin-nature is present only as regards some, not as regards all, tasks and relations. As easily
and unconsciously as they solve some of the ethical tasks, so much reflection and effort must they apply to
others. But there are also natures for whom it is of the very greatest importance that the ethical feeling be
awakened in them and come to speak a decisive word in the single cases. The immediate impulse to self-assertion
and to devotion either does not stir strongly enough in them, or else there is not the right relation between the
two tendencies. --- It is of ethical-social significance that such differences exist. A larger and richer work is
thereby made possible than if all came with the same conditions. Both the multiplicity itself and the relations
of opposition it makes possible work to release and to develop.

In the systematic exposition, of course, all these individual differences cannot be exhausted. We must confine
ourselves to giving an ordered and grounded account of the most important of the qualities of character that the
principle of welfare demands.

\chapterhead{X.}{THE PERSONAL FOUNDATION OF THE ETHICAL LIFE.}

1. For various reasons it has been denied that preparation and practice have any value whatever in the
ethical respect.

The reaction against the ascetic line of thought has especially contributed to this assertion. Asceticism ---
by which is here understood freely chosen effort and suffering that serves merely for practice, not for the
attainment of demonstrable ends --- might, it is held, be consistent at the standpoint of Neoplatonism or of
medieval
% --- p. 190 ---
\opage{190}Christianity. For the Neoplatonists the task was to free the soul from the impurity it had incurred by
entering the earthly world. From the strict Christian standpoint abstinence, mortification and
self-humiliation were sacrifices of atonement to an angry God. But where such views no longer prevail, can
asceticism there be granted any justification and significance?

With the transition to a new view of life there is easily bound up an inclination, for the sake of a
superficial ``consistency,'' to throw away everything to which value was attributed at the earlier standpoint.
What has developed historically under the influence of certain motives may very well keep its value, because
these motives fall away and must be replaced by others. (Cf. I, 4). \textit{Every} ethical standpoint must make
demands for bearing effort and suffering, and practice will therefore always be necessary. It is, to be sure,
a maxim of great importance that all infliction of pain must be justified, while the production of pleasure
has its provisional justification in itself. But where one strives toward a distant and comprehensive end, much
preparation and many privations may become necessary. Every serious view of life will therefore be able to
learn from the heroes of asceticism. These too, after all, really regarded suffering and privation as a good
only because they were thereby strengthened and practiced for attaining the great and distant goal they
craved. Their fault was that they set the goal in another existence, whose conditions were to be \textit{wholly}
different from those of the present life, so that the goal could not be reached through positive work in the
given world. Their asceticism rested on a distrust of natural life and its powers. But precisely at a
standpoint where one does not get the rules for one's conduct through supernatural revelation or mystical
promptings, but where weight is laid on the independent authority of human thought and feeling, precisely there
great demands are made on the individual's powers, and it becomes of decisive importance that these powers be
exercised, so that the foundation of the ethical life may become firm and secure. There are, besides, natures
that in the ethical respect know only an Either---Or, either strict discipline or complete
% --- p. 191 ---
\opage{191}slackness. They can keep themselves up only by submitting to a thoroughgoing discipline. The great
spread of the temperance movement shows that there are not a few such natures. Often these commit the error
of believing that what is necessary for them must be a law for all.

It has been objected that it is a sign of imperfection when actions are undertaken merely for practice, since
means and end then fall apart. Such things may, it is said, be necessary when one wishes to acquire mechanical
skills; but the ethical life is to be the life of the whole personality. Such things may also suit the
standpoint of mystical and supranaturalistic asceticism, where one shunned real life and exercised the powers
of the soul in solitude, so that one perhaps even found something especially meritorious in the
insignificance of the tasks, as when Egyptian monks watered sticks that were planted in sand.

This objection rightly starts from the assumption that life itself is the best school of all, even if it is
the most troublesome and the most costly. It is inadmissible to withdraw from its demands in order to work on
oneself in solitude. One cannot step out of the ranks to become a proper human being and then begin again.
All true education takes place through interaction with real conditions. What the art of self-education can do
here is to choose precisely the tasks and the relations that are suited to exercise and strengthen the powers
before one proceeds to greater works. One ought always to exercise one's powers on something that has real
significance. The definite calling that is chosen will be the natural school of character. At first imitation
and the mechanical performance of prescribed functions are perhaps the highest one can attain; but through
repetition and practice dispositions and aptitudes may be trained that make it possible to perform the work of
one's calling with freedom and supported by one's own strength. Then a virtue answering to the calling has
formed. But it is not at all times that the real tasks lay claim to all our strength. The time can then be
used to prepare oneself for the possibilities. Practice then becomes like a game or an art. In art we have to
do not with life itself but with
% --- p. 192 ---
\opage{192}a reflection of it, and yet art and life have the same content; the use of the powers in art can
therefore be a preparation for their use in life. Besides, the domain of possibilities is greater than that of
reality. In external action often only a slight part of personal life gets occasion for development and
formation. On the stage of the inner life, in the clash of moods and thoughts and through the interaction of
budding motives, an essential part of the work on the development of character often takes place.
\textsc{Schleiermacher} has particularly expressed this in his ``Monologues'' (chap. 4). ``From how many
sides,'' he asks, ``would man not remain undetermined and unformed if his inner acting aimed only at what
actually befalls him from without?\dots{} I know that my outer life will never be able to present and complete
my inner being from all sides\dots{} I live in quiet concealment --- and yet on the world's great stage, rich in
deeds!'' Later S. \textsc{Kierkegaard} expressed a similar thought: ``Only he who is formed by possibility is
formed according to his infinity.'' He adds that the real situations are not so clear and unconditional as the
possibilities can be thought, and that the test through possibilities can therefore be the stricter; but he
also has an eye for how easy it is to deceive oneself when one is tested merely by possibilities. That the
possibilities can also suck the strength out of the will is a side of the matter on which he does not
dwell\footnote[1]{\textit{Søren Kierkegaard som Filosof.} p. 50---51.}.

2. The development of all spiritual faculties can serve to fortify and develop the ethical feeling. The
\textit{strength, range} and \textit{purity} of conscience depend not merely on an education of the life of
feeling, but also on the development of knowledge and of the will. At the standpoint of the free conscience, after
all, it is not merely a matter of blind obedience, of the command being heard in consciousness and the action
following mechanically. The free conscience seeks to set its tasks itself and to see their justification. It is
not only the philosopher who turns testingly against the traditional and instinctively approved rules; the task
may be set for anyone to tear himself loose from what has hitherto held and
% --- p. 193 ---
\opage{193}try untrodden ways. The clarity and consistency of knowledge let the single person catch sight of the
cases in which from weakness or egoism he evades carrying through his own acknowledged law, and it also teaches
him to find holes and unmotivated barriers in the current morality. What he has thus recognized in the calm and
bright times, it is then a matter of holding fast in the dark and agitated moments; only through an exertion of
the will can fidelity to oneself and to the truth one has recognized be preserved.

\textit{Strength} and \textit{range} conscience acquires only where thought and imagination are alive and active.
Thought must follow actions in their effects, which often branch out into many relations, and imagination must
give vivid images of these effects. Narrowness and prejudice often come not from a lack in the life of feeling,
but from a lack of imagination or of logical consistency. Mere knowledge is of course not enough; it must be
practiced and made a constant thought, that is, a thought that by virtue of the laws of the connection of ideas
asserts itself again easily and quickly where there is use for it. It thereby becomes a part of our real self,
of the circle of thoughts and feelings that forms the center of gravity of our spiritual nature, and to which we
always return, however much momentary influences and promptings draw us away from it. Neither the development of
ideas nor that of feelings takes place wholly without our self-activity, and once attention has been awakened,
many occasions will present themselves for intervening regulatively. --- One thing shall only be especially
stressed. It is not only theological ethics that has its books of edification. There can be books one seeks out
ever anew in order to strengthen one's thoughts and one's feeling. Newspaper reading and entertaining reading
ought not to be all the spiritual food one gets, if one does not wish to weaken one's inner life. What each will
choose here depends on his character and turn of mind; but no one can do without the collectedness and
concentration that serious reading brings with it. We are thereby led into a realm of ideas and possibilities
from which we can go with widened and strengthened feeling to the practical tasks.

% --- p. 194 ---
\opage{194}As regards the purity of conscience, self-examination is of the greatest importance, so that we may become
clear about our own motives and learn to know our powers. Self-knowledge always has its limits, since much of
what moves within us comes forth only very indistinctly before clear consciousness. There may, besides, lie
possibilities and aptitudes in our individuality that have not yet made themselves felt either inwardly or
outwardly. We can therefore never be sure that the center of our consciousness is also the center of our
individuality\footnote[1]{\textit{Psykologi} VII B, 5 a.}. But as regards the conscious motives, --- what we
have really thought, felt and resolved, --- we must be able to get clear. There is, no doubt, in most people an
inclination to self-deception, to wish to look better in one's own eyes than one really is, to embellish for
oneself the purity of one's motives and the seriousness of one's resolves. But it is of decisive importance that
one be honest with oneself and learn to tear apart the illusions in which one so easily spins oneself. This
inner truth is the condition of every vigorous and sound activity. On a lie no lasting building can be erected.
--- Self-deception occurs especially easily in such fields where no exact rules can be set up, and ethical
problems very quickly become so intricate that no general principle can untie the knot for us. Even if we are
not in doubt about the quality of what is right, insoluble doubts may arise as to what degree and what range our
activity is to have. If only those were right who hold that one has done one's duty when one has done what
``society'' requires, and that it is not a duty to earn ``merit''! But there are those whose inner spectator is
not so easily satisfied, and in whom the fear of self-deception in fixing ``the qvantum satis of manly will''
will therefore always stir anew.

Modern civilized life has an outward-turned tendency, easily leads to an unhealthy self-forgetfulness. Absorbed
as one easily becomes by the changing multiplicity of external events and
% --- p. 195 ---
\opage{195}circumstances, one comes to regard what goes on within oneself as indifferent. That outer world one
regards as the real reality, the inner experiences as a fleeting reflection of it. But the real life is always
the inner life; it is that which gives the outer its significance. From it, besides, forces can go out that can
intervene to change the outer world; but such forces are not fostered when one lets oneself drift with the
stream. Every personal inner being is one of the spiritual hearths in the world, and the task is to keep the
fire pure and vigorous. There can be a kind of spiritual suicide that consists precisely in letting the fire on
the inner hearth go out from torpor. The ideal is let go, no motive of valuation stirs, and in the end not even
instinct holds one up. In suicide of this kind not even a splash is heard, as from one who drowns himself.
\textsc{Sibbern} has rightly described conscience as a spiritual impulse of self-preservation, the expression of
a kind of care for oneself. One shows care for oneself here by not lowering the demands one makes on oneself. The
level of the self must not be lowered.

Self-knowledge may sometimes be made very difficult by the composite in human character and the sporadic in
human development. Different, even mutually conflicting motives and tendencies may press forward; now the one
tendency prevails, now the other, --- and in which of them, then, is man to recognize his real I? What stirs in
us in the enthusiastic, the ``best'' moments may be most different from what comes forth in the evil times, and
yet both belong to our I. Or we judge and act rightly and surely in some relations, while in other relations we
let ourselves be blindly led by selfish passions. Here the charge of hypocrisy lies so near that popular morality
does not care to undertake, nor indeed understands how to undertake, any thorough psychological inquiry --- and
still less understands how to hold back its judgment. In such problematic states there often expresses itself a
fermentation out of which a clear and harmonious character can later develop. But there are also natures that are
and remain problematic, or in whom
% --- p. 196 ---
\opage{196}the conflict between the elements of the soul at last bursts the unity of character\footnote[1]{Cf.
\textit{Jean Jacques Rousseau og hans} Filosofi. p. 11. --- \emph{Joseph Butler} has a series of interesting
remarks on self-deception (Sermon X).}.

3. But there can also be question of asceticism in a narrower sense, of hardening and of practicing the capacity
to endure effort and suffering. Among savage and warlike peoples a great part of education, self-education too,
consists in the development of hardiness, the capacity to bear cold, hunger and physical pain. In civilized life
such physical hardening is no doubt less necessary, but many of the drawbacks of civilized life surely stem from
one's having lost sight of it too much. Physical hygiene is for a great part also a school for the will. The
self-conquest and patience that must be applied for the benefit of physical health can benefit spiritual health.

But there is also a direct spiritual hardening. There can be inner cold, darkness and impotence just as well as
outer. Even without a definite external cause heavy and dark times may come, when nothing will succeed, when
hardly a train of thought can be held fast, and when everything loses its luster and its freshness for us. Even
the hold one otherwise has in the inner sanction of conscience may fail. One then easily reaches for external
means to keep oneself up, to comfort or to stimulate oneself\footnote[2]{``When a man begins to grow lukewarm, he
fears a slight effort and willingly accepts external comfort.'' \textit{De imitatione Christi.} II, 4, 3.}, and
with that the first step to decay is often taken. Where it is purely physical states of depression that set in,
it may of course be entirely right to use external means. But it may often be the only right thing to offer
resistance, to put one's will into keeping oneself up, in order not to come to suffer from spiritual softness. It
may, as an old mystic says, be necessary ``to do without all comfort.'' The will is thereby exercised in standing
on its own feet.

Among the dark and anxious thoughts that in such times
% --- p. 197 ---
\opage{197}may haunt us, one likes in modern times to count also the thought of death. In antiquity this was not
so common. It has become common only since Christian theology specially sought to sharpen the fear of death,
``like the mothers who rule their children by making them believe that the darkness is full of ghosts that
seize the disobedient''\footnote[1]{Lecky: \textit{History of European Morals from Augustus to Charlemagne.} I,
p. 221 f.}. % translated from Høffding's Danish
In opposition to the predominant place and significance that theological ethics gave to the thought of death,
\textsc{Spinoza} declared that the free man, that is, the man who is guided by right knowledge, thinks of nothing
less than of death; his inquiry is a meditatio vitae, has life as its object, not death\footnote[2]{Ethica IV,
67.}. It is not thoughtlessness that is recommended here. Thoughtlessness can often cover a secret fear. No
serious view of life can avoid struggling with the thought of death. But the difference between theological and
philosophical ethics shows itself here in a characteristic way. For theological ethics the whole of life is a
preparation for another life. We live in front of a curtain behind which true life is hidden, and everything we
do and leave undone is in the end to serve to make us ready when the curtain is at last drawn aside. The thought
of death therefore becomes here an all-deciding, a guiding thought. Philosophical ethics does not build on
assumptions about what lies beyond experience, but maintains that life must first and foremost have its end, its
value, in itself. We attribute value to life not because it is a preparation for another life that we do not
know, but because it contains in itself something that is beautiful and good, something that deserves that one
let oneself be filled by it and strike a blow for it. We live on realities, not on possibilities. Death is the
limit of life, and he is best prepared to approach this limit who has really \textit{lived,} that is, who has had
a share in the best in the life we know.

About the realm of possibilities that begins where our world of experience ends, views will probably always
% --- p. 198 ---
\opage{198}differ. A scientific decision cannot be conceived. We stand here before one of the open questions, and
everyone may here, on his own account, have his faith or his hope. What philosophy here draws our attention to is
merely the fact that all the features and qualities with which the images we can form of another world are
furnished are taken from the world of experience, only that they are imagined in far higher degrees and freed
from all imperfections. If one did not know from experience something that was true, good and beautiful, one
would not be able to form the idea of an absolute truth and goodness, a complete blessedness. But then the true,
the good and the beautiful in the world of experience are a common ground where we can all meet, whether we
assume another world or not. It is a fact that a life can be lived that is devoted to work in the service of the
true and the good without any need being felt to believe in another existence. And no really ethical quality can
be named that would become possible only through such a faith. Even if faith in another existence is a personal
necessity for the single person, he has nevertheless no right to make it an ethical necessity for all.

The counterpart to the struggle against the dark possibilities is the thought of the shining models. Such
personalities, in whom spiritual life stirs with strength and harmony, self-assertion and devotion working
together at every single point, stand as models for those who must struggle laboriously to rise above what
drags them down, or to hold together what threatens to dissolve in strife. An apparent paradox may arise here
from the fact that the personalities who stand as models have not always themselves won perfection through toil
and struggle. Energy and harmony may have been gifts of nature, and yet they can serve as signs that show the
direction and strengthen hope. They show us what possibilities lie in human nature. Had those personalities had
their hands full struggling within themselves, their lives would perhaps not have presented us with the picture
of strength and harmony that now strengthens and guides us. In Greek
% --- p. 199 ---
\opage{199}antiquity \textsc{Socrates} stood, with his wonderful harmony and energetic clarity of mind, as the great
model. Powers that otherwise so easily come into greater or less opposition to one another, --- thought and will,
--- were in him, so far as one can see, united in an original unity. It was this original strength of character
in \textsc{Socrates} that made Socratic philosophy so one-sidedly intellectualistic and lay at the basis of the
doctrine that virtue is knowledge. He did not know from his own experience the difference between knowing and
willing\footnote[1]{Karl \emph{Joël}: \textit{Der echte und xenophontische Sokrates,} Berlin 1893. I. p. 256.
--- The well-known passage in the ``Phaedrus'' where \emph{Plato} has Socrates say that he does not know whether
he is a simple and gentle or a manifold and flaming being must be put to Plato's account; it was he, not
Socrates, who first distinguished different parts of the soul. Cf. Aristotle: \textit{Magna Moralia.} c. 1.
\emph{Joël} op. cit. p. 277.}. But what made his philosophy so enigmatic in its one-sidedness has precisely made
him a great model. Later, when Greek philosophy in \textsc{Neoplatonism} acquired a mystical-religious character,
and virtue was conceived as a striving to become like God, the proposition was enjoined with particular emphasis
that what becomes virtue in the copy is not virtue in the model; for as absolute being God can have no
virtue\footnote[2]{Plotinus: \textit{Ennead.} I, 2, 2. Cf. already Plato: \textit{Theaetetus} p. 176 B.}. And
in \textit{the Christian Church} Christ was set up as the model, although according to the doctrine of the Church
he was by his nature the sinless one. There is this truth in it, that development in the ethical field, as in the
organic field, is made possible by involuntary variations presenting new possibilities and leading into new
directions. Already at low human stages of development not only the physically superior but also the spiritually
superior will be able to come to stand as models for imitation. There has revealed itself in them, in a happy and
natural form, a feature of character that is felt in its value, and that one believes will enrich life if it can
be propagated further. Ethical Aladdins have been the lights on the often dark road of ethical development. The
involuntary always precedes the voluntary, and the qualities that were raised to virtues had
% --- p. 200 ---
\opage{200}first to reveal themselves involuntarily before ethical selection could take place. The impulse to imitate,
admiration and reverence have exercised this selection, together with which the strict discipline on the part of
ruling authorities and of social opinion has constantly worked; yet an interaction has taken place, since the
holders of power and popular opinion could not escape the influence of the models either. \textsc{Fichte},
\textsc{Goethe} and \textsc{Carlyle} have pointed to the great significance of personal models. ``The good,''
says \textsc{Goethe} in a conversation with Eckermann (1 April 1827), ``is no product of human reflection, but is
inborn, innate beautiful nature. More or less it is innate in all men, but in a high degree in certain quite
preeminently gifted spirits. These have revealed their divine inner being through great deeds or teachings, and
this has then, through the beauty of its appearance, won men's love and mightily led to reverence and imitation.
But through experience and wisdom the value of the morally beautiful and good could come to consciousness, the
bad proving in its consequences to be something that disturbed the happiness of the single person and of the
whole, while the noble and right showed itself as something that brought about and secured particular and general
happiness.'' There lies a great biological and ethical truth in these words of Goethe\footnote[1]{Cf. my
\textit{Etiske Undersøgelser.} p. 72. 82 f. --- On Fichte and \emph{Carlyle} see \textit{Den nyere Filosofis
Historie}\textsuperscript{2}. II. p. 151. 385; cf. also p. 358 (Comte's positivist calendar).} --- a truth that is
not diminished by our having to leave the cause of the spontaneous variations (what Goethe calls the innate or
inborn in the eminent men) standing as a problem.

What no moral law and no ethical deduction can give, the great models give in living and vivid forms. Without a
relation to such models the personal foundation of the ethical cannot develop into inwardness and firmness. Which
models are chosen will depend on the need and circumstances of the choosing individuality.
Perhaps the same model will not be able to guide us
through the whole of life, just as it is not the same stars that guide the seafarer in
% --- p. 201 ---
\opage{201}the northern and the southern hemisphere. There is a constant interaction between the ethical models and
experience, as Goethe has seen so clearly.

4. History shows that ethical striving, in single persons as in society, in every period and in every people,
aims at calling forth and developing certain qualities of character that are then regarded as virtues. An
ethical standpoint is designated not merely by the foundation of feeling that determines valuation, nor merely
by the objective standard that is applied, or by the motives of action that are demanded, but also by the
qualities of character that are admired and favored. A virtue, like a duty, can be grounded in different ways;
but it is of special interest, both practical and theoretical, to see which virtues are particularly brought
forward at different times.

It was \textit{the Greeks} who first (in European mankind) set up an individual ethics. Early there formed a
series of qualities that the Greek people admired and demanded. First in the series came to stand the four chief
virtues, wisdom, courage, self-control and justice. \textsc{Plato} \textit{(Republic} Book 4, p. 427 E) takes it as
given that goodness consists in these four qualities, and then tries, partly by a psychological, partly by a
sociological way, to show that all virtue can be traced back to these four, later so-called cardinal virtues
(from cardo, a door-hinge: thus the qualities on which everything turns). The remarkable thing now is that
\textsc{Plato}'s psychological derivation is wholly independent of his sociological derivation. It was thus his
conviction that the single individual in and for itself has a task to solve. This task consisted for Plato in
making oneself a personal work of art, within which the different faculties and promptings work harmoniously
together. The weight laid in the whole of Greek ethics, not merely among the philosophers, on harmony,
limitation and moderation becomes properly intelligible only when we remember that Greek culture had as its
historical presupposition a passionate and chaotic past. There were titanic forces that had to be formed before
Olympian calm and loftiness could come to prevail. The different cardinal virtues express, each from its side,
the ideal thus won.
% --- p. 202 ---
\opage{202}Wisdom lets thought be the guiding power, courage holds up personal dignity against external resistance
and inner slackening, self-control brings it about that the sensual impulses are limited without being killed,
and justice stands as the expression of the interplay of all the elements of the soul: it designates the state
of the soul in which each single part of the soul exercises its own function without hindering the others'; it
is thus a harmony of the soul. Although \textsc{Aristotle} did not take up this doctrine of the four chief
virtues, he nevertheless conceived virtue as an individual harmony that the single person, in and for himself,
apart from his relation to society, strives to realize. Classical Greek ethics ended at this point in the great
problem whether the demands society makes on the individual, and those the individual makes on itself for its
full development, can be brought into agreement. Plato's and \textsc{Aristotle}'s individual ethics suited,
besides, only those who could live for their personal development independently of material labor. And for both
Plato and Aristotle there was a virtue that in the end stood above all and threatened to burst the inner personal
harmony: in the capacity for thinking and inquiry, especially for speculative contemplation, the great Greek
philosophers saw man's real patent of nobility. Theoretical wisdom became for them the highest.

The four cardinal virtues nevertheless went on following ethical thought on its way. The Stoic
\textsc{Panaetius} took them up (in the 2nd century B.C.), without its being clearly seen how he brought them into
connection with his doctrine of love of mankind, and from him in turn \textsc{Cicero} got them and rendered them in
his well-known work ``on duties'' \textit{(De officiis).} If a certain displacement probably took place already
in the classical Greek list of virtues when it passed over from Platonism to Stoicism, then a displacement quite
certainly took place when it passed from the Greek thinkers into the hands of the Roman rhetor. The citizen's
feeling of honor and patriotism came to play a particularly great part; honestum, what a good citizen is to
fulfill, stepped into the place of ``the beautiful,'' social duty into the place of personal harmony. Only in
political conduct, not in private life, could greatness of soul be shown according to the Roman conception.

% --- p. 203 ---
\opage{203}In \textit{primitive Christianity}, as we have already seen (II, 4), love of mankind is crossed by the
ecstatic expectation. In close connection with this last motive stands the relation to God in faith and obedience
as the condition of the soul's blessedness. And just as in Plato and Aristotle wisdom tended to raise itself
above personal harmony, so within Christian ethics faith has a constant tendency to displace love from precedence
among the virtues. Consistently, the more Christian ethics stressed the principle of authority, the obedience of
faith had to become the highest virtue, --- just as the first and highest sin was pride and arrogance (VI, 1).
When \textit{the Church Fathers} traced all virtue back to love, they usually meant by it love of God, which is one
with faith. And since we can learn to know God only with the help of the Church, the proposition: ``Without faith
no virtue!'' is soon replaced by the proposition: ``Outside the Church no virtue!'' The virtues of the heathen
were then to be regarded only as splendid vices: for it is pride and conceit to wish to practice virtue for its
own sake!\footnote[1]{Augustine: \textit{De civitate dei} XIX. 25. --- Cf., for the rest, De Wette:
\textit{Lærebog i den christelige Sædelære og sammes Historie.} Danish transl. § 148. 156. 171.}

Nevertheless on the side of the Church one felt in the field of ethics, as in that of dogmatics, a need to take
up elements and forms from Greek thought. In \textsc{Ambrose}'s work \textit{De officiis ministrorum,} which became
the Middle Ages' handbook of ethics, Plato's and Panaetius's old division of the virtues was transferred, with
Cicero as the intermediate link, onto ecclesiastical ground. This could not of course happen without various
reinterpretations and displacements. In place of the individual harmony that Plato called justice there steps
living love of mankind, which leads to devotion --- in place of wisdom there steps faith, obedience to authority,
--- courage shows itself especially as patience, and in place of self-control purity is
set\footnote[2]{An interesting comparison between Panaetius, Cicero and Ambrose is given by R. \emph{Thamin}:
\textit{St. Ambroise et la morale chrétienne au 4e} siècle. Paris 1895.}. Later \textsc{Augustine} tried to show
that the four traditional
% --- p. 204 ---
\opage{204}virtues appear in the Christian as four different forms under which love of God expresses
itself\footnote[1]{\textit{De moribus ecclesiæ catholicæ.} c. 15; 25. --- Augustine refers to ``the Book of
Wisdom'' (VIII, 7), where the four virtues are named.}. In a more external way ancient and Christian ethics were
combined in the Middle Ages, for example in Thomas Aquinas, the three ``theological'' virtues, faith, hope and
love, being placed above the four ``philosophical'' virtues. But in spite of all reinterpretation and
revaluation it is a sign of the continuity of ethical development that one has for so long been able to use, and
still sometimes\footnote[2]{Victor Cathrein: Philos. mor. p. 98---111. --- S. Alexander: \textit{Moral Order and
Progress.} 2. ed. London 1891. p. 250 f. --- P. \emph{Natorp}: \textit{Grundlinien einer Theorie der
Willensbildung} (Archiv für system. Philos. 1895).} to this day thinks one can use, the old Greek fourfold
division of the ethically valuable qualities of character.

With \textit{the Renaissance} vigorous self-assertion comes forth as an individual tendency that is cultivated and
admired. Joy of life, power and enjoyment of a more or less ideal kind are placed in the first rank. In
philosophical theory, in a whole series of thinkers from the 16th to the 19th century, \textit{loftiness of
mind} (sublimitas)\footnote[3]{This concept had its model in a concept that occurs under different names in
\emph{Plato}, \emph{Aristotle} and the Stoics, without properly finding its place in their division of the
virtues. Cf. Plato: \textit{Republic} Book 6, p. 486 A (διανοίας μεγαλοπρέπεια, loftiness of mind). Aristotle:
\textit{Eth. Nic.} IV. 7 (μεγαλοψυχία, having a great soul). \emph{Seneca}: \textit{Dialogi} I, 4---6; II, 11.
14 (magnanimitas). Both Aristotle and Seneca regard this quality as the highest.}, \textit{courage of life}
(animositas) and \textit{high-mindedness} (generositas) come to stand as the greatest virtues. They are derived
from the impulse of self-preservation and regarded as its highest and noblest forms. Telesio, Bruno, Campanella,
Descartes and Spinoza are especially to be named here. Kant and Fichte join this series in their doctrine of the
ethical qualities, since for Kant personal dignity asserted with strength of soul, as opposed to servility, for
Fichte unconditional self-activity, as opposed to inertia, stood as the highest in the domain of character. In
this valuation of character the reaction is at work against Christian ethics'
% --- p. 205 ---
\opage{205}tendency to make obedience and humility the highest virtues.

In another series of thinkers the memory of the Stoic doctrine of love of mankind works together with the
command of love of Christianity, so often displaced by faith, to lay the chief weight on \textit{fellow-feeling}
and the \textit{feeling of humanity} as grounded in human nature itself and as qualities of character of the
highest value. This tendency is represented by Shaftesbury, Hutcheson, Hume, Adam Smith, Comte, Schopenhauer and
Herbert Spencer.

For this last tendency it is evidently easiest to get a natural connection between individual and social ethics.
In strength of mind and loftiness of mind the individual stood as free and independent, borne by the
consciousness of having the conditions of its life predominantly in itself. A chief feature especially stressed
in these virtues is that they lead to despising external honor in the consciousness of inner worth. As in the
ancient doctrine of harmony, the individual stands here separated from society. In high-mindedness, as Descartes
and Spinoza depict it, it is nevertheless an essential element that the individual has the consciousness of
belonging to a larger whole whose interest is preferred to the purely individual ones. It is the great horizon,
reaching beyond the individual personality, that conditions high-mindedness. But for that very reason this forms
the transition to the qualities of character on which the other series of philosophers laid the chief weight. And
it presupposes them. To high-mindedness one comes only by walking not only the ways of self-assertion but also
those of devotion. The assertion of personal dignity described by Kant is conditioned not only by man's having
the law for his action within himself, but also by this inner law of his being universal, holding for all
rational beings. Nor is this tendency of character, then, psychologically and historically possible except when
the individual is able to feel solidarity with a larger whole.

As we have sought to show earlier (see IX, 1), self-assertion and devotion do not necessarily stand in conflict
with one another. They express two tendencies in human nature whose mutual relation may be most different in the
different
% --- p. 206 ---
\opage{206}individuals, but which may very well stand in a harmonious relation to one another, and which must both
enter as elements into the perfect character. Devotion develops and widens the individual's being, and
self-assertion fits it to fill its place in the whole to which it belongs. The Platonic concept of justice as
personal harmony we can take up to express the inner connection, the higher unity of self-assertion and devotion.
The ancient concept of harmony must be widened and deepened with the help of the ethical experiences expressed in
the valuations of character of Stoicism, Christianity, the Renaissance and modern humanity. But such an alteration
of the content of the concept is very well possible without bursting its frame. Our ethics is Greek ethics, which
later ethical development, through its manifold fluctuations, has been able to alter and correct, but which can
never be displaced so long as there is to exist anything that can rightly be called human ethics.

By a special inquiry into self-assertion and devotion each by itself we shall show the possibility and necessity
of uniting them in harmonious cooperation and thereby of working for the realization of the highest human quality
of character. It is, to be sure, only an indirect description of justice that is reached by examining how
self-assertion and devotion each from its side find their completion in it. But the question is whether it is
possible to get further, if the ideal is not to become an abstraction floating in the air. In each single case
justice develops either predominantly out of self-assertion or predominantly out of devotion. It thereby gets its
different timbre in the different personalities. Over against this richness an ideal depiction of the higher unity
of the two directions in our nature, which stand in constant tension, would acquire a dull tone. It is better,
then, that we sharpen the problem and keep to the different roads by which development can approach the highest
quality of character. Perhaps the time will come when it becomes possible for an ethicist to give a more perfect
and more positive description of it.

% --- p. 207 ---
\opage{207}\chapterhead{XI.}{SELF-ASSERTION.}

1. The virtues and duties that belong to self-assertion can be referred to three main forms:
self-preservation, self-control and independence. In all of them it will appear that the individual's
striving to assert itself is, ethically regarded, conditioned and limited by regard for the existence and
development of the whole race. As already remarked, there is in modern times, on account of the zealous and
significant endeavors at emancipation, a certain tendency to overlook this, while it is easier to show at
more primitive stages of development. But when one goes back to the foundation from which we here regard the
relation between individual and society, it will nevertheless appear that the proposition is right.

\begin{center}a. Self-preservation.\end{center}

2. If it is not in itself a good to live, all ethics is unreasonable. Directly or indirectly all ethics aims at
preserving life, protecting it and developing it further. It continues what nature itself begins. Even where a
clear consciousness has not awakened, an instinct of self-preservation asserts itself in all living beings,
which makes the individual seek what is beneficial to its existence and shun the opposite. \textsc{Kant} even
held that the instinct or impulse here was so strong from the very hand of nature that one could not speak of
a duty of self-preservation: for what everyone unavoidably wills of himself, he cannot be obliged
to\footnote[1]{\textit{Metaphysische Anfangsgründe der Tugendlehre.} Einleitung § 4.}. This view is
connected in part with \textsc{Kant}'s mistaken assumption that we can be obliged only to what we will
\textit{unwillingly}. But it also rests on the no less mistaken assumption that we really always and
unavoidably seek what serves our self-preservation. That is
% --- p. 208 ---
\opage{208}not so even with the animals, in whom the instincts work, after all, to a far greater extent and with far
greater power than in men. Even there instinct can lead astray, or it can be weakened and altered. But for man,
at any rate, self-preservation becomes the business not merely of instinct but also of the will proper.

3. Bodily health and strength are the necessary foundation of all further and higher development of life. At
any given time we have at our disposal only a limited sum of energy, and it therefore becomes an important task
to see to it that it is preserved and increased and does not go to waste. We cannot produce something out of
nothing; even the hero is nothing without food and drink. Asceticism with its self-inflicted pains and
privations, every kind of self-torment and groundless despondency, levity and dissolute sensuality can each in
its way waste the capital that could be used for valuable ends. Prejudice, shortsightedness and willfulness lead
either to rejecting the necessary means or to consuming them uselessly. When it is clear to the single person
that the consequences of his conduct in this respect can strike not only himself but also others, partly those
he might himself bring into the world, and to whom his weakened nature may pass by inheritance, partly also
those who must do without his help and his work, then care for physical self-preservation will stand as a
significant duty.

But herein also lies its limitation. It would be meaningless to spend life in gathering a capital that is never
used, to make life one single great means without an end. The ascetic is more reasonable even so, since he
trains himself to attain a distant end. To live is to use the powers, and they are used only by being used up.
A great physiologist has even given the paradoxical definition: ``Life is death!'' because every manifestation
of life brings with it the dissolution of organic tissue. A narrow conception of self-preservation would make it
one with the accumulation of matter and would not see that consumption is just as essential a thing in life as
the taking in of matter, even if it is this that in the end puts an end to life. It is not even always that we
can preserve the harmonious relation between intake and consumption.
% --- p. 209 ---
\opage{209}A disproportion may come forth with ethical necessity. A task that it is my duty to solve may make
necessary the setting aside of what the duty of self-preservation would demand under ordinary circumstances.
Döbeln (at Juthas) had no time to be ill. Life must even be wholly sacrificed when, to preserve it, we would
have to deny what gives it its real worth.

To live for one's health is an absurdity. It may happen that out of regard for one's failing powers one must
withdraw to more limited tasks; but to limit one's ends is not the same as to make an end of what is to be only a
means. \textsc{Plato} mocks the conduct of the physician Herodicus during an incurable illness: ``Since he had no
other occupation, he lived merely to doctor himself and was most unhappy at the slightest departure from his
accustomed way of life; in this way, by his wisdom, under a hard death struggle, he reached a great age.''
\textsc{Plato} remarks that in the case of the poor it follows of itself that they cannot live to doctor
themselves; it is only the rich who can set themselves this as a task; and he holds that ``he who cannot follow
the natural way of life ought not to be cured either, since it is of benefit neither to him nor to the
state''\footnote[1]{\textit{Staten.} 3rd Book. Heise's transl. p. 178---182.}. We would not proceed so strictly.
We have learned that under and through suffering qualities may unfold that perhaps were hitherto pushed back, and
that could give life a new value for the sufferer himself and for those who are influenced by his example. His
character perhaps only now gets its completion. We attribute value to the life of thought and feeling apart
from the external effects. We see the single personality not merely as a member of the state, but also as a
member of the family, of the circle of friends, and perhaps see him even in his outwardly inactive state being
the greatest joy and strengthening to others. Of the last suffering days of the famous freethinker \textsc{David Strauss}, which were so edifying for his friends, his biographer E. \textsc{Zeller}
writes\footnote[2]{\textit{D. Fr. Strauss in seinem Leben und seinen Schriften geschildert,} p.
122.}: ``He did not demand of nature that it
% --- p. 210 ---
\opage{210}should spare him anything of what its course and its laws brought with them; but all the more did he
demand of himself that he should know how to accept these laws, to make suffering too the material for a moral
and spiritual activity, and to find what is beneficent in the midst of pain.''

4. We now come to the much discussed question of the justification of \textit{suicide.} By the nature of the
case this question can be raised only insofar as suicide is a real action, the fruit of deliberation and
resolve. If suicide always springs from decided insanity, it does not fall under ethical consideration. The
suicide that springs from a desperate situation into which the individual has brought himself by his own conduct
may also be of such a character that it is no premeditated action, but the effect of a sudden mood overwhelming
all reflection. Here, however, ethical valuation does not fall away, but is directed (as with actions performed
in drunkenness) at the conduct by which the individual has brought himself into his desperate state. Thus in him
who has plunged himself into misfortune by fraudulent or at any rate indefensible financial speculations and then
taken his life, the real guilt lies before the suicide. So too in \textsc{Musset}'s ``Rolla,'' who takes his own
life after he has broken down his physical and spiritual personality by debauchery. Debauchery can weaken the
power of resistance, even if it does not directly produce desperate situations. Alcoholism and suicide rise, as
statistics show, side by side.

In such cases where suicide clearly springs from the wish to evade definite obligations, there is no ground for a
long debate either. It is an open desertion, an action of the same kind as self-mutilation in order to escape
military service. Under Roman military criminal law attempted suicide was punished by beheading, being regarded
as attempted desertion. He who murders himself in order to escape a painful or dangerous situation into which he
and his family have come, and who thereby withdraws his help from those nearest to him, undoubtedly performs a
reprehensible action. But the unethical thing then lies not in and for itself in his putting an end to his life;
it
% --- p. 211 ---
\opage{211}lies in his evading his obligations, something he could also do in another way.

The question is whether all cases of suicide can be referred to the kinds named here. In that case the problem
would fall away. Before I go more closely into examining this, I shall touch on some circumstances connected with
the, as it seems, strong increase of suicides in modern times.

5. In his book \textit{Om Selvmord i Kongeriget Danmark} \textsc{Kayser} communicates a series of letters written
by suicides just before death. The common tone of these letters is a lack of capacity to bear one's fate, to
continue the struggle for existence under the given conditions. Even if one does not desert, one capitulates.
One feels oneself facing an insuperable resistance. It is neither courage nor cowardice that shows itself in such
suicides, but a weakening of the strength of will, a cowing of the impulse of life and the joy of life. From an
inquiry published by the State Statistical Bureau, \textit{Selvmordene i Danmark i Tiaaret} 1886---1895, it
appears that the most important cause has been melancholy, weariness of life and insanity, more than
\textsuperscript{1}/\textsubscript{4} of all the suicides, and in the case of women even
\textsuperscript{2}/\textsubscript{5} of the suicides, having been committed from these motives. This agrees with
the impression one gets from the letters published by Kayser. As regards the different age groups it is seen
that in youth (from the 15th to the 25th year) unhappy love and remorse or fear of punishment play a considerable
part as motives of suicide; in manhood drunkenness and economic worries come forward in the first rank; in old
age bodily illness and melancholy, weariness of life or insanity become more and more the predominant motives,
and in the case of women earlier than in that of men. The weakening of the power of resistance and the blindness
to the given values of life or to brighter possibilities to which the frequency of suicides (--- in Copenhagen a
suicide almost every third day ---) bears witness have no doubt for a great part causes that lie far back, in
temperament and in the involuntary life of mood. But circumstances can surely also be pointed out that come with
the whole development of culture in modern times. Modern culture,
% --- p. 212 ---
\opage{212}with its one-sided development of the intellectual faculties at the expense of the other faculties, with
its wide horizons and its restless but also often planless and aimless spirit of enterprise, makes it harder to
hold fast the compactness and concentration, the limitation as regards ends and means, that is a necessary
presupposition of a sound and vigorous life of the will. Wide horizons are attractive to thought and imagination;
but the will, if it is to be more than mere impulse and wish, must be directed to something quite definite, and
to something that is within reach. Modern times are a time of doubt, not particularly favorable to harmony
between the powers of the soul. It is the task of spiritual hygiene to remedy the disharmony and morbidity that
so easily develop under these conditions. The educator and the physician must work together. Everything that
serves the more vigorous development of the active faculties --- to which belong not only physical capacity for
work but also the capacity for independent observation and thought and the artistic imagination --- will work
healingly against the disease of the will. Not by breaking off the work of culture and cowing doubt --- which is,
besides, an impossibility --- but by strengthening self-activity and awakening confidence in the human powers,
where these work in their right field, does one help to overcome the unsteady, brooding or blasé doubt that
follows the magnificent development of modern times as the shadow follows the body.

To this there is added yet another thing. Modern culture makes individuals at once more dependent --- by
awakening more needs in them --- and more isolated --- by ``emancipating'' them and setting them on their own
feet. And the isolation reaches its culmination when the individual with his many unsatisfied needs stands alone
amid the great bustle around him. The single individuals must fight more each for himself than under the old
order of things. The new associations no doubt counteract this isolation; but they are still in their childhood
and have hardly yet had sufficient influence on the disposition of single persons. The unhappy and suffering man
feels himself left to himself. A herd of animals maltreats or kills
% --- p. 213 ---
\opage{213}the individuals that cannot keep up with the pace of the herd; in the human world they must often shift for
themselves, even if they are not maltreated and trampled down. A living relation of sympathy, in which the single
person feels the sympathy of others and himself follows their conduct with sympathy, will see to it that his
courage of life is not so easily broken, and his horizon not wholly darkened. He who does not let go of others
will not so easily feel himself forsaken by all either. Blasé indifference and weariness will not so easily step
into the place of freshness and interest. --- It seems as if suicides decrease during great political crises,
perhaps because minds are then occupied by the great events and the common interests. If that is so, constant
absorption in common tasks must be a power that can prevent despondency and self-abandonment. The resignation that
is to be able to carry one beyond a great disappointment or a great distress becomes psychologically possible only
where there is a great interest of life in which one can be absorbed and for which one can live, when that in
which one has hitherto found one's real I must be given up. Suicide springs from regarding that in which one has
hitherto found one's I as the only thing. This earlier I must often be wholly renounced if suicide is to be
prevented. The impulse to suicide is the expression of an impulse of self-assertion that yields only to the
higher kind of self-assertion that appears as renunciation of oneself. It is in this sense that \textsc{Mme de Staël} has said: Le renoncement à soi même est en tout l'opposé du suicide.

He who has lived as a lonely stranger in a great city knows what an oppressive feeling of forsakenness can seize
the mind, and what an encouragement it can be merely to get a friendly word from a passer-by whom one asks the way
or to whom one does a service. So a sympathetic utterance, a friendly address, will often be able to stop the
suicide on his way. He will feel that the bond that ties him to the race is not broken, and may feel a need on his
side too to tie it fast again. In the old days the monasteries were places of refuge for unhappy and dispirited
natures, who could find support and tasks there in a common life. The homelessness of natures of that kind has
become greater since the monasteries were closed. It is now placed in their own hands to find tasks and a home.
--- What
% --- p. 214 ---
\opage{214}matters, then, is to set in motion forces that can abolish the isolation of the single faculties within
man and between men in society. Herein this special question is connected, in its ethical ground, with the great
general ethical and social questions.

6. Is suicide, now, always due to sickness of the will, downright insanity, or else a conscious wish to evade
definite obligations? Or can there be a right, perhaps even an obligation, to commit suicide?

Where the individual's full freedom to dispose of his life has been maintained, one has always started from a
purely individualistic standpoint. --- Thus \textit{the Stoics} taught that the wise man, who is sufficient to
himself and the absolute master of his fate, will not only give up his life when the good of his country requires
it, or when he can preserve it only by committing an unethical action, but that he will leave life when on account
of external dispensations he cannot lead it as he wills. The later Stoics extended this right ``to lead oneself
out'' of life very far. ``Not only in the extremity of need,'' says \textsc{Seneca} (Epist. 70), ``but as soon as
fortune begins to become somewhat questionable,'' the wise man will take his own life. \textsc{Marcus Aurelius}
demands merely that one shall leave life without anger and sorrow, as one goes out of a room where there is
smoke \textit{(Comment.} V, 29; X, 8). To end one's life from cowardice and softness the Stoics regarded as
immoral, --- although on the other hand they held that cowardly and soft men were unworthy guests at the feast of
life, and that they would therefore perhaps do best to withdraw. --- In modern times suicide was defended from an
extremely individualistic standpoint by a series of writers in the 18th century. \textsc{Montesquieu}, for
example, grounds the justification of suicide on this, that since society is grounded on mutual advantage, one
must be allowed to withdraw from it when one has only trouble from living in it; one thus transgresses no duty by
putting an end to one's life when one finds it unbearable \textit{(Lettres persanes.} Nr. 76).

Even if one were now to assume that the single individual had a free right to leave life when this was
unbearable, it would
% --- p. 215 ---
\opage{215}be very difficult to show that the unbearableness was really present in the single case. For how can the
individual know that it has reached the limit of its strength and has used all means? --- To this it has been
answered that the very fact that the natural instinct of self-preservation is cowed and the natural fear of death
overcome bears witness that man must have been brought to the extremity. Thus \textsc{Holberg} says in his 135th
Epistle: ``Experience teaches that men hold life dear; from which one may conclude that those who commit such a
deed [suicide] must be overwhelmed by anguish, suffering and adversity, indeed in such a state as rather awakens
pity than anger.'' But because I am overwhelmed by anguish and pain, so that my power of resistance is broken, I
may for all that well be in an illusion. Instinct does not always guide me surely, and the cessation of instinct
can therefore not always be a sure mark either. A momentary feeling of anguish and distress is no reliable
expression of the nature of the real circumstances. Perhaps the individual is precisely at the lowest point of the
rhythmic swing of fate, so that things will go upward for him again if only he can hold out a little longer. It is
one thing that the individual despairs, another whether the state of affairs is really desperate. A fully grounded
assurance of this can never be had. When in the opinion of some a decrease in the number of suicides can be traced
when emigration increases\footnote[1]{Legoyt: \textit{Le suicide ancien et moderne.} Paris 1881. p. 257
f.}, an example is thereby given of how new possibilities and ways out can keep courage up. A slackening of energy
and courage may bring it about that one does not discover and use the ways out; but an objective proof that they
are not there will be difficult, if not impossible, to give.

If the individual is absolutely sovereign (cf. III, 5---7), it will not matter so much here either whether such an
objective proof is possible. But when the highest standard of ethical valuation is sought in the single
individual's connection with the race, its significance for the race, the demand for strict proof cannot be
abated.
It was
% --- p. 216 ---
\opage{216}therefore a right thought when in some Greek states (Massilia and Ceos) the blame of suicide fell away if
the suicide first convinced the authorities that there was sufficient ground for taking such a step. In Thomas
More's \textit{Utopia} priests and magistrates urge the incurably sick who can no longer work and are a burden to
themselves and others to end their lives voluntarily, whereas a suicide that lacks such authorization brings
disgrace. The presupposition lies at the basis here that the individual does not live for himself alone, but is
placed before tasks that are of significance for others than himself. It may be wholly right that life is
unbearable in the single case when the individual knows no other standard than his own egoistic enjoyment. An
approximation to such a narrowing of interest will be found in all suicides who act with deliberation. They have
staked everything on one card, have made their life a game of hazard. Their passion is concentrated on a single
point, with which the whole of life stands and falls for them. This holds, for example, of the admired suicides of
antiquity. The ``last Romans'' killed themselves because their horizon did not reach further than the
aristocratic-republican life of Rome. We know human relations both more comprehensive and more inward that can
make claims on us. And even Cato acknowledged this when (according to Plutarch) he advised his son, whom he did not
wish to take with him into death, to keep away from politics in the times that must now come. If the son's honor
did not suffer from following this counsel, the father himself could have followed it without denying himself.
The strength and defiance that make up the heroic in deliberately chosen suicide often have something theatrical
about them. They would be better used if directed toward the goal: to make life more bearable, if not for
oneself, then for others. --- Where suicide springs from the weakening of the will described above (5), too, the
individual succumbs to resistance at a single point, often insignificant objectively regarded. But there this
concentration on a single point is due to a preceding weakening of energy. In the cases described here, on the
other hand, interest has beforehand been gathered with the whole
% --- p. 217 ---
\opage{217}energy of passion on one point, and everything therefore falls together with it. Two wholly opposite
psychological processes can thus lead to the same outcome.

The hope the single person has lost for his own part he will be able to hold fast for the race, especially when he
grasps the great doctrine of the significance of small effects. Modern science has opened our eyes to what the
uninterrupted accumulation of small effects can accomplish. The mightiest coral reefs have arisen by the deposit
of the shells of countless tiny creatures. What the greatest revolution could not accomplish is brought about by the
quiet working of everyday, inconspicuous forces. The most modest activity in the narrowest circle contributes a
mite to the great common life of the race. Nothing need go to waste. Every expression of inward feeling, of an
ideal direction of life, of firmness of character can acquire its significance and intervene in a beneficial way,
even if we do not get the theatrical feeling of doing something great. \textsc{Rousseau} rightly says (in \textit{La
nouvelle Heloïse),} that if he who feels himself tempted to suicide would examine whether there were not still
some good actions he could perform, a poor man he could help, an unhappy man he could comfort, an oppressed man he
could defend, he would surely be held back from carrying out his purpose.

Just as little as it can be shown that all possibilities are exhausted, just as little, then, can it be shown that
all obligations are fulfilled. The more seriously and deeply life is conceived, the more possibilities and
obligations present themselves. This is the chief consideration that ethics must assert in this matter. For the
rest, ethical endeavors with regard to suicide must aim rather at preventing it by strengthening the power of
resistance, the will to live and sympathy with the living, than at formulating stern judgments on suicides
committed. What matters is to bring it about that every individual feels itself both as an end --- by being an
independent member of the race --- and as a means --- by being a member whose activity can become of significance
for other members. Fortifying the foundation of life is better than moralizing, and especially when this comes
afterward, and with regard to an action that is as a rule
% --- p. 218 ---
\opage{218}the effect of violent blindness. In all actions that are ethically reprehensible blindness is at work; but
none of these actions is so frequently motivated by real misfortune and distress as the one by which the
individual puts an end to his life. ---

It must still be asked whether under certain circumstances it may not become downright a \textit{duty} to take
one's own life. Suicide and self-sacrifice may border so closely on one another that it may be difficult, if not
impossible, to draw the line between them. The individual may think that he stands in the way of others' happiness,
or that by leaving life he frees others from a danger. Thus when someone fears that under a sharp interrogation he
will come to betray secrets of great importance, and therefore kills himself in order not to cause irreparable
harm. Or when someone who has been bitten by a mad dog and notices the outbreak of the disease kills himself in
order not to come to bite others. Or when a poor father of a family takes his life because he knows that if he
dies (and only then) his family, which is now exposed to hunger and distress, will be provided for. If such cases
occur, there can be no question of the individual's evading obligations by taking its own life. Its action
expresses precisely an acknowledgment that it does not feel itself and its existence as the only or highest end,
but regards itself as a member of a larger whole for which it must sacrifice itself. On the other hand the
difficulty still remains whether really all possibilities are exhausted. Where a human life is at stake, where the
annihilation of one of the living elements of the race is at stake, a definite demonstration of the necessity of
the action must be required if it is to be ethically acknowledged.

In the ancient Church, which was otherwise so stern toward suicide, an exception was nevertheless made of the women
who killed themselves in order not to be violated by the persecutors; they were honored as
saints\footnote[1]{Cf. \emph{Barbeyrac}: \textit{Traité de la morale des pères de l'église.} Amsterdam 1728. p.
242 ff. --- Only \emph{Augustine} had some misgivings about the matter; but he helped himself with a hint that
those women might have acted in accordance with an immediate divine revelation (De civitate dei I, 26).}. Wholly
apart from the overstrained
% --- p. 219 ---
\opage{219}ascetic presuppositions from which the Church started, we can here find an example where suicide is an
ethical action. Suicide here establishes that there is an abuse of human beings that in brutality stands on a
level with murder, and it thereby contributes mightily to sharpening and raising the ideas of the significance
of feminine purity.

How frequent suicides from noble and high-hearted motives are is difficult to decide, as it is in general bound up
with great difficulties to get information about the motives that lead to suicide in the single
cases\footnote[1]{Cf. \emph{Marcus Rubin}'s critical remarks on the statistics of suicide in ``Tilskueren''
1884. p. 466 ff.}. \textsc{Morselli}\footnote[2]{\textit{Der Selbstmord.} Leipzig 1881. p. 270 ff. Cf. also
Oettingen: \textit{Moralstatistik.} 3. Aufl. p. 780 ff.} disputes the assertion made by some writers that in men
who do not suffer from insanity it is for the most part noble motives that bring about suicide. In our day at any
rate suicide is, he holds, essentially an outcome of egoism. ``Yet,'' he adds, ``outcomes of the better part of our
nature are not lacking either, and that especially in the female sex\dots{} In men one's own interest appears as
the wholly ruling motive, and since only a fourth or a fifth of the suicides are committed by women, the rarity of
noble motives in suicide is already secured by this numerical relation.'' It must be left to the statisticians to
judge this assumption.

7. Suicide can be regarded only from a medical, a psychological and an ethical standpoint. It cannot rightly be
drawn under juridical consideration. Of the state the individual is a member only so long as it lives, and the
punishment of the state cannot strike the suicide himself. The punishment here in reality comes to fall on others
than the individual himself, whether it consists in the confiscation of the inheritance (as in the Roman imperial
period, when the motive was the wish to escape punishment for lese-majesty) or in contemptuous treatment of the
corpse (as has been common in all countries from the Middle Ages on). --- Even
% --- p. 220 ---
\opage{220}Bishop \textsc{Martensen}\footnote[1]{\textit{Individuel Etik} p. 461.} found in what is called ``the
laxity the state shows in the burial of suicides'' a sign of the degree to which ``the spirit of Christianity has
been pushed back.''

\begin{center}b. Self-control.\end{center}

8. While self-preservation rests especially on overcoming external resistance, self-control consists in overcoming
resistance that stems from man's own inner nature. Self-control presupposes that two or more different tendencies
are given in man, and that one of them must be subordinated to or displaced by the other. If there could not be
several mutually opposed impulses or passions in the human self, self-control would be impossible. In self-control
it is feeling that yields to feeling, impulse that yields to impulse, passion that yields to passion. One does not
become master of oneself by one's ``reason'' alone, but by gathering the energy of consciousness into a strong
feeling or passion directed toward what stands before us as the highest. Self-control is in and for itself a purely
formal virtue, whose value depends on the cause in whose service it works. He who loves money above all is able by
this passion to hold other impulses and promptings down. In others the joy of knowledge and the passion of inquiry,
in others again the feeling of love or love of country, is the ruling force. By these and many more ways the
capacity for self-control is developed. Historically the feelings determined by the relation to authorities (in
family, state and church) have especially had great influence on the development of the capacity for self-control.
The relation to an authority that man feels far surpasses his own power and ability awakens a feeling of
boundness, directs all powers to one goal, subordinates all regards to a single one, and so works as a
concentrating power in consciousness. Under the dominion of authorities, with obedience as the chief virtue and
with hope and fear or --- at higher stages --- enthusiastic reverence as the ruling passions, mankind has gone
through a school of self-control
% --- p. 221 ---
\opage{221}and gained a practice that need not be lost because the original motives are replaced by others.

Only through self-control is the inner freedom and unity won that makes full and collected activity possible. It is
said to have been the Cynic Antisthenes who first enjoined self-control as a virtue of its own, teaching that he
who does not master himself is a slave, whether he is outwardly free or not. It was an ethical point of view from
which the difference between the free man and the slave lost its significance\footnote[1]{Cf. Karl \emph{Joël}:
\textit{Der echte und der xenophontische Sokrates.} II, 1. p. 39---47. --- With regard to the theory of
self-control one must note especially the strong stress of several older thinkers on the point that an impulse can
be checked only by another impulse --- an analogy to the principle of inertia or to the law of energy in the
physical domain. Cf. \emph{Plato}: \textit{Phaedo.} Heise's translation. p. 23---25. Bacon: \textit{De augmentis
scientiarum.} VII, 3. \emph{Spinoza}: \textit{Ethica} IV, 7 (cf. V, 20. Schol.).}. --- Inner freedom is the
condition of outer freedom and independence. The different impulses and promptings in our nature are just so many
ways by which the outer world can become master over us. If it can, without more ado, awaken any promptings whatever
in us at any moment, it plays on us as on an instrument whatever melody it will, and we stand powerless before it.
What matters, then, is that we have a center in ourselves, a circle of thoughts and feelings that always determine
the main direction of our life. Self-control is not to make us without feeling, but our feelings are not to sway
hither and thither with every wind. --- It needs no closer demonstration how great a significance self-control has
for the development of conscience and for self-preservation. In the foregoing we have therefore already presupposed
it at many points.

9. Spiritual life develops sporadically, begins from different dispositions and points of departure, and only
gradually can harmony be attained between the different tendencies that stir. Different powers and impulses move
consciousness in different, often opposite, directions. Only through a period of fermentation and struggle does one
often succeed in bringing about order and unity. It is a condition of the higher development of life that
% --- p. 222 ---
\opage{222}a rich material be given to work up. But the rich material may be refractory and difficult to unite.
Asceticism cut the knot by cowing and pushing back the natural impulses. It took up a suspicious attitude not only
toward the sensual instincts but also toward the impulse to knowledge and the aesthetic needs. The great art, on
the contrary, is to give the immediate and involuntary promptings the free course that is a condition of a sound
and fresh life, and yet to be master of the direction of the course, so that this does not conflict with what one
regards as one's real goal. Dead and mechanical order and self-tormenting anxiety are signs of imperfection. It is,
besides, not always healthy to draw the half-conscious promptings and fantasies wholly forth into the daylight of
consciousness; one gets rid of them again more easily when they remain on the outskirts of consciousness. Much
moralizing directs attention precisely to what is to be displaced.

Self-control takes on a different character not only according to the nature of what it serves, but also according
to what it displaces. It may be enjoyment, advantage, honor, revenge that is the object of the flaring affect that
is to be counteracted. Every affect or emotion hinders clear and calm deliberation and does not allow ideas to
combine in the way that would otherwise be natural to them. Special dispositions in the single individual can make
the struggle hardest for it in the one or the other direction. There may be innate dispositions to despondency,
irascibility, vanity and sensuality. The work the single person must do here, what he must suffer in order to
attain inner freedom and harmony, will therefore be most different and cannot proceed according to any definite
recipe. The demand for self-control sounds alike to all, but means for each something most different both in
degree and in kind. Just as one and the same burden weighs differently on different shoulders, so very different
work is required of the different individuals in order that they may keep themselves at the same stage of
development. Many a man must struggle hard merely in order not to sink below the current average of ethical
development. To the external observer a struggle
% --- p. 223 ---
\opage{223}of this kind does not easily become visible or admirable. But Leonidas was nevertheless no less a hero than
Alexander the Great, although the latter could use his courage to conquer the world, while the former had to fight
to the death in a vain attempt to keep the enemy away from his country.

Already \textsc{Aristotle}, in his famous definition of virtue as the right mean, by which he means a harmonious
relation between the human impulses, clearly acknowledged these individual differences in the kind and degree of
self-control. For the right mean, according to Aristotle, does not lie at the same point in all
individuals\footnote[1]{Eth. Nic. II, 5---9}. Since their inclination to the various affects is different, what in
one individual is a sign of great self-control may in another individual be something that happens quite
involuntarily. What in one is a sign of avarice may be thrift in another and prodigality in a third. Yet Aristotle
speaks more as a psychologist than as an ethicist; at any rate he does not see the ethical consequences of his
doctrine. He is entirely right that the conditions for reaching inner harmony are most different in the different
individuals as a consequence of their different nature. But he overlooks that this inner harmony does not always
exhaust what is demanded. The significance of self-control rests, after all, on our being to work with collected
strength in the work for our tasks. But what if my task demands a greater degree of self-control than agrees with
``the right mean,'' the inner harmony? There may be equilibrium in my nature without my nature therefore satisfying
the demands that circumstances make. I may perhaps combat my irascibility to such a degree that I show greater
self-control than he who is by nature gentle and peaceable; and yet it is not said that I am secured against being
carried away by the flaring affect in a reprehensible way. There asserts itself here in Aristotle the individualistic
tendency spoken of earlier (III, 5), which stands in remarkable opposition to the great weight Greek ethics
otherwise lays on state and society.

% --- p. 224 ---
\opage{224}Both self-control and self-preservation look different from a purely individualistic standpoint than from
the standpoint of humane ethics (III, 10). The individualist must, to be sure, strive for order and harmony between
the different sides of his being (III, 6), but from the standpoint of social ethics demands are made on self-control
that the individualist would have no reason to make. If he only preserves calm and harmony within himself, why should
he not follow his acquisitive impulse, his impulse to revenge? From the relation to others he always turns back to
himself and has no occasion to follow up the effects of his actions when these do not fall on himself.

There lies, however, an abiding truth in the Aristotelian doctrine. It is incontestable that one and the same ethical
demand can necessitate a most different inner work in different individuals. The problem that Aristotle passed
lightly over, because he did not bring individual ethics into connection with social ethics, keeps its significance
from the standpoint of a social ethics too, --- provided, that is, that this standpoint is really to be
\textit{ethical} and not coincide with the juridical standpoint or with that of public opinion. When there is an
inner relation between individual and society, it cannot be a matter of indifference with what presuppositions in
the ethical respect the single individual comes. Only a dogmatism that would assert that all men are equally
endowed with the capacity for self-control could maintain that the right ethical law had been found when the same
external quantum of self-control was demanded of all. (Cf. IV, 2 and VIII, 6). Rather than grant this, which would
lead to barbarism and pharisaism, ethics must acknowledge its limit in the scientific respect. Wherever individual
differences come into force in a decisive way, there is a limit to science. What can be demanded of everyone of
whom self-control can be demanded at all (by himself or others) is that his action lie in a series that would,
continued, lead to perfect mastery over himself; but at what point in the series the single individual's action is
to lie in the single case is something that ethical
% --- p. 225 ---
\opage{225}thought, with the means at its disposal, is unable to determine in such a way that both individual and
social regards come into their own. And that for the simple reason that the different individuals do not begin
from the same point.

10. The difference between the standpoint on which we have here taken our stand and the individualistic one is seen
also from the way in which the different kinds of self-control are valued. We do not regard the single individual as
a little closed world by itself, but always in its relation to others. If it is asked, then: is it most important
to master one's lust for revenge, or one's covetousness, or one's sensual impulse to enjoyment? --- an answer can of
course be given according to the way in which these impulses intervene in the individual's own inner being; but this
answer will not be sufficient. Sensuality and craving for enjoyment make man dependent on the external and thereby
hinder his free development. They also hinder the development of the ideal feelings and weaken the energy of the
will, perhaps even come into conflict with what physical self-preservation requires. Of greatest importance in the
ethical respect is that devotion to sensual enjoyments narrows the mind. Sensual enjoyment often cannot be shared
with others, indeed can perhaps only be attained at others' expense. It is perhaps at first a thoughtlessness called
forth by the affects that lets the enjoying individual disregard the consequences its enjoyments may have for
others; but through habit this thoughtlessness turns into heartlessness. Dissolute men are often more hard-hearted
and cold than men who offend against others' property or let themselves be carried away by flaring anger. Ethical
and juridical consideration here lead to different results. He who steals for the sake of his starving children
stands, ethically regarded, far above him who in debauchery destroys the material or spiritual capital that could
have benefited many. Positive morality and especially public opinion here come, as so often, too near the merely
juridical way of looking at things. The thief is excluded from ``good society,'' but the dissolute egoist perhaps
occupies the place of honor.

As regards the sexual relation in particular, it can first
% --- p. 226 ---
\opage{226}find its full judgment within social ethics. Within individual ethics the most important rule is that to the
pure all things are pure. The impure arises when consciousness is entangled in sensual ideas and cannot free itself
from them. Shyness and prudery may be questionable signs, because they bear witness that such ideas very easily
present themselves to the individual. Not by cultivating a mystical shyness of the phenomena of sexual life, but by
regarding them like all other natural phenomena with an unprejudiced eye and clear understanding, does one work for
a sound and natural purity of mind. There are only two sound ways of standing toward the sexual relation: partly in
the way indicated, as a physiological and psychological observer, partly as one who experiences its power through
inward devotion, that is, through more than merely sensual enjoyment. Different from both is dalliance, coarseness
and lewd fantasizing.

The relation to sexual life poses one of the very hardest ethical problems. Nature and culture collide here in the
sharpest way. As a natural being man is a mature man or woman long before his existence as a cultural being is yet
so established that he can provide for his offspring. Strong demands are constantly raised here on the side of
nature that cannot be satisfied to the extent to which they press forward without bursting cultural and social life.
Here, as so often, the natural need stirs with unruliness and prodigality, just as if nature wished at any price to
secure the existence of life on earth, however it goes with the single individual. It is nature's carefree
scattering of the germs of life that shows itself here from its most glaring side and sets the beings in whom the
life of consciousness has awakened before the greatest tasks. Nature uses the individual organism as a means; but
when a personal life has developed, it will not be a mere means. There is a struggle between the instincts we have
in common with the animals and the interests that lead to a higher development of human life. \textsc{Kant} even
drew from this the conclusion that man was not by nature disposed for civilized life, but that nature's intention
was only the preservation of the human race as an animal species. He found here a deep disharmony and a source of
misery
% --- p. 227 ---
\opage{227}that can be overcome only when men reach a stage of culture from which they are still far
removed\footnote[1]{\emph{Kant}: \textit{Muthmasslicher Anfang der Menschengeschichte.} (1786). \textit{(Vermischte
Schriften.} Halle 1799. III, p. 48 f.). --- Cf. with this \emph{Elie Metchnikoff}: \textit{Études sur la nature humaine.} Paris 1903. chap. 5 (Désharmonies dans l'organisation et le functionnement de l'appareil de la
reproduction).}. The great thinker has here pointed at once to the great difficulty and to the only way in which
one can, ethically regarded, stand toward it. It is easy enough to cut the knot and demand without more ado the
unbound satisfaction of the natural instincts. Such satisfaction is possible precisely only if we would go back
again to the animal stage. And even at the animal stage complete unboundness does not prevail, as the choice of
mates and the struggle for mates show plainly enough. Even if asceticism in its hatred of nature and official
morality in its prudery end in untruth and unnaturalness, the reaction against them in its way goes too far when it
would maintain that the sexual relation is neutral, or at any rate of subordinate significance, in the ethical
respect. It is by virtue of a right instinct that at all times the sexual relation and its ordering have been
regarded as one of the very most important ethical questions. Here we have to do with it as yet only from the
standpoint of the single (though not sovereign) individual and can therefore give no complete treatment of this
question; such a treatment the doctrine of the family can first give. But so much is clear, that since the sexual
instinct can be satisfied only with the help of another individual, there can be no question here of any
unlimited individual freedom. If the satisfaction of the instinct brings misfortune and degradation to another
individual, or even presupposes the existence of a whole class of human beings who live by being means for others'
sensual pleasure, then no closer demonstration is needed that we by no means stand here before something
indifferent, a purely private matter. The sexual instinct hides in itself the germ both of the highest and noblest
and of the lowest and basest in human nature. It can lead the individual beyond itself in enthusiasm and joy over
another
% --- p. 228 ---
\opage{228}human being and thereby show itself related to all devotion and sacrifice; but it can also lead the
individual, entangled as it is in its promptings, to make others mere means to its enjoyment. Between these poles
the instinct swings, and a remarkable blindness is needed to think that self-control should have no significance
here. Self-control proves precisely here to be man's patent of nobility, the condition of his being able to
continue the course of development by which he has raised himself out of an animal existence.

Self-control is also possible in a far higher degree than is often thought. One often speaks here of physiological
necessities that are not present, but which one can, to be sure, foster by constantly talking of them. There are
individuals who preserve their purity without harm to the soundness and strength of their spiritual and physical
life. How far the single individual's capacity for self-control goes, only it itself can try; but just as little as
anyone has the right to cast the first stone, just as little has anyone the right to call his own weakness a
universal human necessity\footnote[1]{There is no pedagogical reason to lower the level. \emph{Esquirol} remarks:
``Si la continence, dans quelques cas très-rares, a causé l'aliénation mentale, le libertinage est une cause plus
fréquente.'' \textit{Des Maladies mentales.} Bruxelles 1838. I. p. 24. --- In very strong terms the English
physician \emph{Lionel Beale} \textit{(Our Morality and the Moral Question,} chiefly from the medical side. London
1887) attacks those who assert that physiological necessity. He cites (p. 99) a statement in the same direction by
Sir James Paget. --- Cf. also S. \emph{Ribbing}: \textit{Om den sexuela hygienen} (Stockholm 1888). (Lectures for
Swedish students). --- \emph{Albert Heim}: \textit{Das Geschlechtsleben des Menschen vom Standpunkte der natürlichen Entwickelungsgeschichte.} Zürich 1900. (A lecture for Swiss students). --- \emph{Erik Pontoppidan}:
\textit{Kønssygdommenes Betydning og Forebyggelse.} Copenhagen 1903. (Lectures for Danish students). In this last
work it is especially stressed (p. 65 f.; 98) how ethical and sanitary regards here coincide: ``At the head of all
good counsels to youth concerning these diseases and their prevention is placed an exhortation to the education of
the personality and its subordination to firm moral principles, in order to strengthen the character and set a
manly will against the sexual temptations.''}. He who must struggle hard to keep his head above water can do so with
the
% --- p. 229 ---
\opage{229}consciousness that he is taking his share of the suffering the race must go through to reach a higher
stage, and that he is contributing his mite to furthering the development toward it. We are all members of the
great series of development that leads from the animal to man. And in no respect, perhaps, is it so important as
here that the need be kept from sinking back to the purely animal form. The mighty need in question here expresses
itself in the animal as a rule blindly, in isolation and indifferent to its object, although features are found in
the animal world that bear witness that a higher form is already possible there. The need becomes humanized the less
it is separated from other elements in the personality, the more it is bound up with and determined by the special,
individual qualities of its object, and the more it becomes the supporting foundation for sympathetic feelings and
ideal interests\footnote[1]{In an interesting discussion on ``éducation sexuelle'' % PRINTED AS IS: écheation (read at the image)
 in the French philosophical
society both the opposition between the biological and the sociological consideration of the sexual relation and
the opposition between the theological and the philosophical consideration of this relation made themselves felt.
Against the purely biological consideration it was asserted from the sociological side that the sexual act must be
seen in its connection with family life, and that it stands as properly human only when it is the expression of an
inward and lasting community. Against the theological consideration it was asserted from the philosophical side that
religious motives are not necessary for an ethical attitude toward the sexual relation, but that self-control is
here motivated partly by human dignity, partly by the feeling of justice. \textit{(Bulletin de la Société Française de Philosophie.} 1911). (1913).}.

Just as little here as with regard to the motives that lead to suicide can the struggle always be carried on
directly. Self-control cannot always be exercised at the single, extreme moment. By bodily exercise, by sound and
vigorous nourishment for the imagination and by enthusiasm for ideal ends the instinct is kept from spreading and
calling forth the fixed ideas and the giddiness that lead to a fall.

11. While positive morality does not lack a pharisaic attitude toward the aberrations that come about through the
unfolding of luxuriant instincts and promptings, it has not so much eye for the fact that self-control can become of
very great
% --- p. 230 ---
\opage{230}significance under wholly different circumstances too. If in the period just spoken of there is danger of
\textit{going astray,} there may in a later period be danger of \textit{getting stuck.} And many who go astray in
youth get stuck all the more later. Many young stormers of heaven end as Philistines; many a life divides into a
period of debauchery and a period of blasé indifference. He who remains true to himself through the struggles into
which the restless promptings lead will also more easily be kept from stagnating.

The influence of habit and of regular conditions of life often makes no slight demands on self-control if spiritual
freedom and freshness are to be preserved and a new growth constantly be possible. The new must find us new; but the
presupposition of this is a constant endeavor to keep the eye open and clear. Our growth cannot, to be sure, go on
to infinity. Care has been taken that the trees do not grow up into heaven. But what matters is to get all valuable
possibilities of growth developed. There is no more beautiful sight than a vigorous and lively development at an age
when all development is otherwise wont to be concluded. Constant spiritual work and warm sympathy with everything
great and beautiful is the best means of avoiding rigidity\footnote[1]{Cf. my characterization of F. C. Sibbern's
personality and course of development in \textit{Mindre} Arbejder.}. The Emperor \textsc{Marcus Aurelius} addresses
the following exhortation to himself in his notes (VI, 30): ``See that you are not Caesarized (ὅρα μὴ
ἀποκαισαρωθῇς)! Take care that you do not take on the color of your position! That happens, after all, so easily.
Keep yourself therefore simple, good, pure, serious, without pomp, just, pious, full of good will and love,
steadfast in fulfilling your duties. Struggle so that you may remain such a man as your striving for wisdom would
make you!'' --- The danger of becoming a Philistine that the philosophical emperor feared threatens everyone,
whether he is of low or of high degree, and the remedy against it is everywhere the one he recommended: to be true
to oneself and to one's best striving.

% --- p. 231 ---
\begin{center}\opage{231}c. Independence.\end{center}

12. Both self-preservation and self-control are conditions of personal independence. The single individual is to
assert itself outwardly and inwardly in order to come to stand as a peculiar and independent member of the race. It
has something in it that does not recur in quite the same way in any other individual. The need for independent
development, for making one's dispositions and faculties count, is one of the most important forces that lead to the
advance of human life. Without faith in oneself and confidence that there is something in one's self that deserves
to live and unfold, nothing can be accomplished. Despondency and distrust of oneself paralyze all willing and all
work, are indeed often the sign of a weakening of the will that can approach idiocy. Already \textsc{Aristotle}
enjoined that he who has too little self-esteem does not come to accomplish all the beautiful and good he otherwise
could, because he does not regard himself as worthy of it. Such ``micropsychy'' (as Aristotle calls it) hinders the
unfolding of what lies in the individual's nature. True self-esteem (macropsychy) is, to be sure --- as Aristotle
also shows --- difficult and rare; it presupposes that one not merely regards oneself as worthy of something great,
but also really is worthy of it. He who possesses this feeling has a fixed standard in himself. He regards external
honor and external goods only as something subordinate. He knows his limits, seeks to accomplish everything he can
within them, tolerates no encroachment on the part of others, but commits no encroachment himself either. He goes
through life secure and confident\footnote[1]{Aristotle's doctrine of macropsychy and micropsychy is found in
\textit{Eth. Nic.} IV, 5---9. --- Cf., from a more recent time: \emph{Hume}: \textit{Treatise} III, 3, 2 (Of
greatness of mind), and Adam Smith: \textit{Theory of moral sentiments.} VI, 3 (Conclusion).}.

Humility and modesty are not such great virtues as they are often made out to be. The great praise that has fallen
to their share can be explained only partly as a reaction against the exaggerated and ruthless self-esteem that is
so frequent, partly
% --- p. 232 ---
\opage{232}as an after-effect of the ascetic line of thought, for which obedience is the highest virtue. What matters
is to know oneself and then strive to bring forth everything one can do. If only the ideal is set high enough, no
special praise of humility is needed. We shall then feel our barriers well enough and need not weaken our courage by
dwelling too much on our impotence\footnote[1]{Spinoza defines humility (humilitas) as ``the sorrow that arises from
man's considering his impotence or weakness,'' a feeling that easily paralyzes and leads to contempt of oneself
(abjectio). Eth. III. Aff. Def. 26---29. --- \emph{Aristotle} too will not regard bashfulness (αἴδως) as a virtue
proper. \textit{Eth. Nic.} IV, 15.} instead of gathering and employing the power we really have. He who knows his
limits has at once faith in his peculiarity and an eye for his smallness\footnote[2]{Cf. \emph{Goethe}'s valuation
of himself (in a conversation with Eckermann, 30 March 1824). People had wished to set Tieck beside him as a poet. To
this \emph{Goethe} says that it would be the same as if he himself would place himself beside Shakespeare.}. He feels
himself at once great and small. True self-esteem thereby acquires a stamp of humor.

13. One does not assert one's independence by anxiously keeping all foreign influence away. The more material a man
can treat and appropriate, the more he will be able to show his peculiarity. True independence shows itself
precisely in being able to immerse oneself in a rich content. The most original men are usually the most willing to
acknowledge what they have learned from others. \textsc{Marcus Aurelius} begins his work by enumerating all those
who have had influence on his development. \textsc{Goethe}, who like few others had the right to feel his
independence, says (in a conversation with Eckermann): ``People speak of originality; but what does it mean! As soon
as we are born, the world begins to work on us, and so it goes right to the end. And what, after all, can we call
our own except energy, strength, will! If I could say all that I owe to great predecessors and contemporaries, there
would not be much left.''

Often the individual thinks it can assert its independence
% --- p. 233 ---
\opage{233}only by withdrawing to feeling and appealing to it. Feeling has something more individual about it than
thought and action. But if individuality is not to be the same as isolation, and if the one-sidedness and limitation
from which every single individual always suffers are to be supplemented, it must not stop at the mystical appeal to
feeling, but the peculiarity must show itself in the work of thought and will. --- An overstrained feeling of one's
own personality without the power of understanding as a counterweight is a frequent trait of character in the
insane\footnote[1]{\emph{Maudsley}: \textit{Pathologie de l'esprit.} Trad. de l'anglais. p. 228 ff.}.

14. External independence is necessary for the internal to be rightly present. Dependence on others brings with it
barriers to the free movement and unfolding of the powers. A mature ethical personality will therefore strive for
\textit{external freedom.} We see, too, that the history of the world is really a great history of liberation. What
is struggled over in history is not merely the means to live, but also the means to be able to live as one will.
\textsc{Hegel} has aptly said: in the Orient one was free, in Greece some were free, in modern times all are free. We
are now all called ``Herre'' [master]; what first designated the superiority of the one over the other now
designates the single person's personal independence. It was an ascetic tendency that asserted itself when
Buddhism, Stoicism and the oldest Christianity alike maintained freedom as something purely inward and declared it
indifferent whether one was outwardly free or a thrall. The slave may, no doubt, have his inner sanctuary where no
one can penetrate; but the complete opposition between inner and outer is not only painful, but can harm ethical life
itself by hindering the inner powers from being applied outwardly. When history, seen from one side, is a great
process of liberation, its ethical significance is that independent, personal points of departure for the life of
the race are created. The race lives only in the single individuals; the more vigorously and freely these move, the
fuller and richer will the life of the race become too. Individual and social ethics thus meet here in the same
% --- p. 234 ---
\opage{234}demand. (Cf. VIII, 6). The limit of the single person's personal freedom must be determined by what the
equal freedom of others requires, and it will be the task as far as possible to get these limits to coincide with
the limits of the individual's ability and impulse.

Honor and property can be regarded as external extensions of personal freedom.

In order to be able to reach our goal it is not enough that we have faith in ourselves; we also need a certain
acknowledgment on the part of others. Others too must regard us as a being with a justified peculiarity.
\textit{Honor} comes under sound conditions of itself, as a presupposition that need not draw special attention to
itself, but which nevertheless gives our appearance before others security and firmness. It is unhealthy when
attention is directed too much to the image of us that is outlined in others' consciousness. What is to be a
condition of our activity then easily becomes its end, and appearance then comes to prevail.

Personal development needs, next, material means too. So long as the individual lives from hand to mouth, no higher
development can come about; and as soon as it has more than from hand to mouth, it has \textit{property.} Here in
individual ethics we regard property only as a means for personal independence, as an extension of the personality.
In every ordering of society that does not make the single person a mere machine, a will-less tool, he must have
disposal of the means he can procure under certain conditions. However these conditions are to be formulated, and
however great a control one would have exercised over his disposal of the means, there must always remain a domain
where he himself makes the decisions. And this domain ought to be as large as is compatible with regard for the
welfare of others. Every limitation of disposal is in itself an evil. Why should I not use my powers to work up the
material I find, or my power to lay hold of as much as I am able, when I harm no one whatever thereby? In and for
itself it is a good to be able to spread
% --- p. 235 ---
\opage{235}oneself in existence; evil first comes when it comes to a collision between several who all have this need
to spread themselves. The burden of proof therefore lies on him who would restrict my sphere of power, just as the
burden of proof in general (by virtue of the principle of welfare) lies on him who would have us feel pain instead of
pleasure. This need for property and the unfolding of power is not necessarily egoism. It is not said, after all,
that one makes possession and power one's highest ends. They can be used, like health, physical and spiritual
strength, in the service of mankind.

At all stages of human existence we find this need for property, at any rate movable property. The community of
property, especially of land, which is so common at more primitive stages of development, disappears everywhere with
the more vigorous development of the single personalities that civilization brings with it. It proves incompatible
with a more independent and freer development of personal life. Under community of property, as for example in the
South Slavic zadrugas, everyone's fate is fixed and cannot become essentially different from the others'. But there
naturally arises a need to ``live as one pleases, work for oneself, drink from one's own glass,'' --- to ``enjoy the
charms of the independent life and defy its dangers''\footnote[1]{Laveleye: \textit{Om Ejendomsretten og dens
primitive Former.} Danish transl. p. 120. --- On the history of property cf. \emph{Spencer}: \textit{Political
Institutions.} Chap. 15. --- Cl. \emph{Wilkens}: \textit{Sociologi} p. 284---295.}. It is very possible that the
right of property, its grounding and its extent, will and must undergo manifold changes. But whatever ordering is
then applied, an essential part of the value that can be ascribed to it will depend on the extent to which it
satisfies the need for independence in as many as possible. In independence lies one of the most important
conditions for the use's being good\footnote[2]{V. \emph{Falbe Hansen} \textit{(Stavnsbaandsløsningen og
Landboreformerne.} Set fra Nationaløkonomiens Standpunkt. I. Copenhagen 1888. p. 69 f., 73) remarks: ``The
community, like all such old institutions, had had its natural justification\dots{}. But in the course of time it
lost its justification and became in a high degree harmful, especially because it hindered change and progress. The
single person had to follow the others in cultivating his land; he not only had to keep to the same system of
husbandry, but he could not either without difficulty make greater changes in the details\dots{}. The enclosure broke
old custom and routine; it compelled the peasant, who had formerly followed thoughtlessly and will-lessly in the
inherited common husbandry of the village, to independent, individual thinking and acting. It made it possible for
the single person to go forward without being hindered by others.''}. ---

% --- p. 236 ---
\opage{236}Civil rights must be regarded in a similar way as honor and property. History (especially that of Rome and
England) shows us of how great significance it has been for the development of a people, indeed of the whole of
mankind, that the single person stands firm on his good right. ``Complete surrender of right,'' says
Jhering\footnote[1]{\textit{Kampen for Retten.} Danish transl. p. 20. --- Cf. also \textit{Der Zweck im Recht.} I,
2. Aufl. p. 74 f. 259. --- Steinthal \textit{(Allg. Ethik} p. 154) remarks: ``In Germany, where the development of
law was for centuries so wholly hindered, characters are sparsely sown and the feeling of right weak.''}, ``is a moral
suicide.''

\chapterhead{XII.}{DEVOTION.}

\begin{center}a. Love of other beings.\end{center}

1. Love works as directly as possible for what is the final goal of all ethics. Were there only this one feeling
in man, and were it not so often blind, neither practical nor theoretical ethics would be necessary or possible.
But with so simple and perfect a compass men are, after all, not equipped.

Love presupposes a capacity to recognize a life of feeling in other beings. It extends as far as it
% --- p. 237 ---
\opage{237}is possible for us to put ourselves in others' place, to feel and suffer with them. General love of mankind
has slowly developed through the successive widening of the circle of beings with whom one could sympathize. The
range of the feeling is not, however, without influence on its kind and strength. In the narrowest relations of
family and friendship it appears otherwise than in the looser and more distant relations. In the question of the
right kind and strength of love in the different relations it becomes clear, as has been shown earlier, that love
passes over into justice. Justice is love ordered according to its own principle. General love of mankind (the
disinterested and universal sympathy) leads, when it becomes clear about itself, to the demand that precisely for
the sake of the welfare of the whole race the greatest love shall be shown in the narrowest circles, where men can
be and do most for one another. Through an immediate and constant community of life, as in the family or in a close
relation of friendship, mutual access is opened to as deep an insight into the personalities as is at all possible.
The feeling of unity can here become as great as it can at all become. The strength of love, whether it has the
character of vehemence or of inwardness, cannot become so great in the wide circles as in the narrow ones. No
distinction is made here between mine and thine, between giving and receiving, and yet precisely the peculiarities
of the personalities come to light, since nowhere can they unfold so freely and confidently as in such a relation.
In such narrow relations sympathy was born in the human race in the first instance, and it is born ever anew. Even
if, therefore, ethical valuation, in order to take in all regards, takes its stand at the standpoint of the general
welfare, it nevertheless demands precisely from this standpoint the existence of the small circles. In them the
highest that can be reached is in many respects reached. They become models for the larger circles. And not only
that, but from them flows also the force that can animate the larger circles.

There can often be a disproportion between the strength and the range of love. It may limit itself so that a group
egoism arises, when the interests of the family, the class, the nation
% --- p. 238 ---
\opage{238}or the race are cultivated at the expense of great universally human interests. Just as the single
person's impulse of self-preservation must be subordinated to regard for the welfare of all, so the limited sympathy
too must sometimes be sacrificed to the more comprehensive one. Our duties as members of families or as friends do
not stand above our duties as patriots and as men. But the disproportion may also be of the opposite kind. There may
be a sacrifice and enthusiasm for great national and universally human interests that leads to ruthless and brutal
conduct in the smaller circles. ``There is no lack,'' says Maudsley\footnote[1]{\textit{Pathologie de l'esprit.} p.
260.}, ``of examples of the martyrs of the cause of mankind themselves producing martyrs among those to whom they
stand in an intimate and daily relation. The humble and tedious duties and self-denials of daily life demand a calm
and unostentatious self-control\dots{} but they make no claim on public attention and sympathy.'' There is rightly
indicated in this utterance of Maudsley that the sympathy that leads to activity in larger circles very easily gets
an admixture of ambition and vanity. Often too it may be restlessness that leads beyond the narrower but more inward
relations.

When devotion and love in general are set up in ethics as a virtue to be practiced, or a duty to be fulfilled,
peculiar psychological difficulties are bound up with this, which have often been stressed from the philosophical
side --- in antiquity especially by the Stoics, in modern times especially by Kant and his disciples.

The idea of general love of mankind appears, as we have seen (III, 8; IX, 1), in Greek philosophers in the last
centuries before the appearance of Christianity. This feeling takes on a particularly inward form in the Stoics of
the imperial period. \textsc{Marcus Aurelius} thus demands that we shall feel ourselves not merely as a part (meros)
but as a member (melos) of the realm of men. One is thus to stand toward others not as a stone in a heap stands
toward other stones in the heap, but as an organ stands toward other organs in the same organism. And with this is
connected the inward absorption by which love
% --- p. 239 ---
\opage{239}becomes an enrichment of one's own inner life. ``If you are merely a part, not a member,'' says Marcus
Aurelius, ``then you do not yet love men from the heart: doing good to others does not yet so gladden you that it
wholly absorbs you, but you do it as mere duty, not yet with the feeling that you thereby do good to yourself.'' ---
What nevertheless prevented love from winning the acknowledgment of the Stoic philosophers that it won in
Christianity was another tendency that was peculiar to Stoicism --- and still more peculiar to it than the tendency
to devotion, namely the endeavor for spiritual independence and independence of everything external. He who devotes
himself makes himself dependent, after all; he no longer disposes of his life of feeling; his joy and sorrow are
conditioned by what happens outside him, and his feelings may be set into stronger oscillations than is compatible
with the preservation of the inner harmony that was the highest thing for Greek individualism. Stoicism has
misgivings about all strong emotion, but especially when this is called forth from without. In \textsc{Epictetus}
this side of the matter appears most clearly. We are, to be sure, to have fellow-feeling with others' sorrow and to
express this fellow-feeling, to sigh with them: ``but,'' adds \textsc{Epictetus}, ``take good care that you do not
come to sigh in your inmost being!''\footnote[1]{According to the Stoics pity belonged to the affects (τὰ πάθη)
and was a vice. The wise man feels no pity. \emph{Diogenes} Laërt. VII, 111. 123. --- While --- says
\emph{Seneca} --- vigorous natures are inclined to anger, gentle natures are inclined to lesser faults, such as pity.
Just as it is unworthy of the wise man to become dependent on others' baseness by becoming angry, so it is unworthy of
him to become dependent on others' misfortune by feeling pity. \textit{De ira} II, 7. 15 (cf. I, 7---8). --- For
Indian wisdom too pity contained a danger. The royal sage Bharata was just about to attain the redeeming knowledge,
but missed it when he tended a suffering young gazelle. (R. \emph{Garbe}: \textit{Die Sankhyaphilosophie.} p. 143).
--- See my comparison between Buddha and Jesus: Religionsfilosofi p. 270 and 272.} --- In opposition to this fear of
becoming dependent through love stands Christianity's faith that spiritual freedom can subsist with the strongest
oscillations of the life of feeling, indeed even first properly comes into being in
% --- p. 240 ---
\opage{240}him who lives through wholly and fully the great oppositions of sorrow and joy, fear and hope. Christianity
here stands with a realistic stamp in relation to Greek philosophy. It does not fear that devotion to the reality of
life, the individual's full entering into the fate of the race, will rob the life of the soul of its freedom and
strength. It stakes life in order to win it. Here is an element in the ethics of primitive Christianity that
philosophical ethics must appropriate. It is an ethical experiment that has an abiding significance, independent of
the special circumstances under which it was made\footnote[1]{See the comparison between Stoicism and Christianity
in my essay \textit{Hedenske Sandhedssøgere} (Mindre Arbejder. I. p. 174---179).}.

\textsc{Immanuel Kant} and his school did not see the matter in this way. When Kant thought he was in agreement with
Christian ethics, it was because he conceived its command of love as the setting up of an ideal of perfection, to
which it was possible to approach, but which no finite being could completely realize. Love is a feeling, --- and a
feeling is not produced on command, says Kant; therefore what is commanded by that command must be the practical
performance of duty in the ideal form, where the performance springs from the inmost heart and disposition. But in
man, Kant continues, other motives than that which consists in respect for the law always express themselves, and
self-compulsion is therefore always necessary. The ideal is holy, but man can bring it only as far as
virtue\footnote[2]{\textit{Kritik der praktischen Vernunft.} Kehrbach's ed. p. 100---102. Cf. \textit{Tugendlehre}
§ 25---26 (love of mankind, not as ``Liebe des Wohlgefallens,'' but as ``Maxime des Wohlwollens'').}. One of Kant's
Danish disciples, M. G. \textsc{Birckner}, had the particular misgiving about the Christian command of love, as it was
uttered by Christ, that it leads to grounding morality on self-interest: ``The qualities that awaken the feeling of
love in us are precisely only such as either immediately contribute or at any rate could contribute to gladdening and
benefiting ourselves.... Love is thus at bottom, according to its origin, a self-interested feeling\dots{}. As love
has its origin in self-interest,
% --- p. 241 ---
\opage{241}so on the other hand respect for persons has its origin in the pure, disinterested respect we feel in them
for the moral law itself.'' \textsc{Birckner} solves the difficulty in a similar way to Kant. What is meant in
Christianity by the word love is, he holds, ``desire for the moral good for its own sake''\footnote[1]{M. G.
Birckner: \textit{Efterladte Skrifter.} Copenhagen 1800. p. 175---184.}.

The psychology of feeling proves the Kantians both right and wrong in their misgivings. It is in a certain sense right
that love cannot be commanded forth. It is of no use to demand love where there is no love. But if one starts from the
assumption that ethics is more than law that has been ``turned inward,'' and that the ethical can keep its validity
even if the tone of command is laid aside, one will see that feelings can very well be called forth, and that a
striving to bring about the conditions for the arising of love can therefore very well be demanded. Often a feeling
can be called forth only indirectly and by detours; that makes the task intricate, but does not make its solution
impossible. It is an important consideration in the ordering of the relations of living together, in education and
self-education, that care be taken that the conditions for loving devotion come to be present. Involuntary living
together and working together will, when it is allowed to unfold in peace, work quietly and slowly in this direction.
And the regard of great models' inward and loving life will awaken to enthusiastic imitation of the strength of mind
that expresses itself in it. The Kantians are surely right, too, that there can here be question only of
approximation. True love is such a force that everyone in comparison with it will have constant occasion to feel his
impotence and his distance from the ideal. But this holds, after all, of the relation to all ideals, to the ideals of
self-assertion no less than to those of devotion. Of how many ethical qualities of character does it hold at all that
they can be brought about immediately, on categorical command?

It rests on false psychology when the Kantians think that all feeling, with the exception of respect for the moral
law,
% --- p. 242 ---
\opage{242}is of a self-interested kind. Self-interest or egoism requires a conscious referring of everything to
oneself as end. It is not enough that a feeling is bound up with self-satisfaction for it to be called egoistic
(unless one gives this word a very wide meaning). From the fact that devotion brings with it, or is conditioned by,
pleasure in the object to which we devote ourselves, it cannot be inferred that we make ourselves the end and treat
the object only as the means to our enjoyment. Devotion would be psychologically impossible if the object could not
awaken our pleasure\footnote[1]{Cf. my \textit{Psykologi} VI C, 7.}. And the feeling of respect that Kant describes
is a psychological impossibility if what is respected cannot awaken any pleasure whatever in us. If love is
self-interest, so is respect.

2. In each single case love may take on different forms, each of which may have its ethical
interest\footnote[2]{Aristotle in the 8th and 9th books of the Nicomachean Ethics submitted these different forms to
an inquiry that still has classical significance. (Danish translation of these books by E. \emph{Bojesen} in the Sorø
school programme 1859).}. It cannot always have the character of an immediate feeling of unity, even if the
individuals it concerns do not stand far from us. The one party may stand so far above the other that the relation
on the part of the latter acquires a stamp of admiration and reverence, which yet need not change anything in the
equality of the mutual love. Community and belonging together may be felt in spite of the difference in character and
faculties. In the superior party love expresses itself as high-mindedness. What is high in high-mindedness rests not
on the individual's feeling himself standing at a higher stage than others, but precisely on his not feeling himself
standing at a higher stage, although he really does. Inward devotion lets all such differences vanish.
High-mindedness may appear as humor, when there is indeed an eye for others' pettiness and limitation, their
self-contradictions and imperfections, but this does not hinder sympathy with them. What is high in high-mindedness
rests further on fellow-feeling's
% --- p. 243 ---
\opage{243}subsisting independently of others' conduct, in spite of their ingratitude and indifference. And finally
what is high in high-mindedness rests on the eye's being fixed on great, comprehensive interests of life, so that
sympathy is not limited to such narrower circles where the single person stands immediately face to face with others.
The high-minded man's love of mankind is determined by a far-seeing knowledge of the conditions of life of the human
race; his striving is to assert these, even if in particulars he must injure what a sympathy bound to narrower circles
would demand. Wrongly, therefore, does \textsc{Adam Smith}\footnote[1]{\textit{Theory of Moral Sentiments.} IV, 2.}
hold that high-mindedness (generosity) and love of mankind (humanity) are two wholly different feelings. He infers
this from the fact that he who stakes his life to save another would perhaps feel no real sorrow if this other had
died from a cause whose effect it was not possible to prevent. But this inference is not right. It is, after all,
quite natural that love of mankind stirs when there is a possibility of action, even where there exists beforehand no
closer relation to the other individual. It is only in the narrower circles that it frequently has occasion to appear
as actual; toward beings standing further off it keeps itself, apart from special occasions, in potential form.
\textsc{Schopenhauer} (\textit{Neue Paralipomena} § 244) has aptly remarked that there is a kind of courage that
springs from the same root as goodness of heart, namely from man's feeling his existence as almost identical with that
of others. ``This consciousness produces courage in that man clings less to his individual existence, since he lives
almost as much in the existence of all beings and therefore cares less for his own life and what concerns it. This is
not always the source of courage; for courage can spring from different causes. But that courage is the noblest kind
of courage, which shows itself in its being here bound up with great gentleness and patience.'' --- There is in
high-minded love such a strength and fullness that it belongs in itself to the things that give life value. It is
valuable not merely through its effects,
% --- p. 244 ---
\opage{244}but by its very existence, and its effects are only its involuntary expansion. ---

True love cannot be shown from above downward. This becomes especially significant for charity. Often the gift is
handed down from the clouds, and pity is frequently disguised pride. High-minded love wipes out the difference
between giver and receiver, acknowledging the other's right to take as well as its own duty to give. It works toward
making the suffering and needy man an independent member of the race. Nor does it treat the receiver as a mere means
of getting an outlet for blind sympathetic instincts or for a sentimental inclination. We may often follow the blind
instinct here just as little as anger, ambition and other promptings, even if in cases of doubt we do best to let the
heart speak. We do not give, after all, merely because we feel a need to, but because we wish to bring real help.
\textsc{Malthus} was the first to show how difficult an art charity is. Here again the close connection of love with
justice appears. In the working of the sympathetic instinct regard for the welfare of larger circles must be held
fast.

In him who fixes his eye on the great gulf between the needy and suffering on the one side and those who have
relatively easy access to the material goods of life on the other, a giddiness easily arises that can weaken both
sympathy and energy. One must not, of course, shut one's eyes to this gulf. But the great problem of poverty and
suffering is far more intricate and ramified than that it could be solved by direct intervention. It would be of no
use if we all gave all our goods to the poor. We must seek our justification in this, that by the work we have chosen
we make the best contribution we can to the welfare of the race. But we can never feel completely secure. When have
we done enough? Is our zeal sufficient? Is our more favorable position not due in part to injustices in the
distribution of goods in society? Are we really working in the direction of a better distribution? --- Questions such
as these will keep their sting. The physical and moral misfortunes
% --- p. 245 ---
\opage{245}in the world are sometimes felt most strongly precisely by those who do not themselves directly suffer
under them. ---

There is not only a physical but also a spiritual charity, and the latter is just as important as the former. Here
too what matters is to help others and to assert oneself. Spiritual help is still harder to bring than physical.
\textsc{Kant} even denied that one could work for the perfection of other men, while he would indeed grant that one
could work for their happiness. His thought was that a man's most essential qualities must be his own work, spring
from his own willing. Herein lies the truth that all spiritual help can consist only in the awakening of
self-activity. But this holds --- though Kant would not grant it --- of ``happiness'' just as well as of
``perfection.'' Happiness has firm ground only when it is self-acquired. Yet it holds quite especially of inner,
spiritual goods. The giver must here, even more than with physical help, make himself a means. Without impatience and
without egoism he must enter into the receiver's state in order to be able to give him, or be for him, what is
needed. From above downward nothing can be accomplished here either. There must, as \textsc{Socrates} expressed it, be
given a spiritual midwifery; the germs the individual hides in itself are to be coaxed forth. This Socratic method is
the only fruitful one. It is love's own method. --- Spiritually degraded men are helped already by showing them that
one believes them able to accomplish something. The confidence one shows them can become the germ of their own
self-confidence. ---

High-minded love is not conditioned by what the individual has itself received or suffered. It is the full and rich
inner life that makes the high-minded man not change his conduct on account of others' coarseness, ingratitude or
enmity. (Cf. IX, 1). It is the expression in feeling of the unity of the race; but this unity cannot be abolished by
anything that may happen. It was an important breakthrough in the ethical development of the race, the most important
step on the road to the general love of mankind, when men began to rise above the proposition that expresses the
standpoint of the barbaric stage of development: thou shalt love thy
% --- p. 246 ---
\opage{246}friend and hate thine enemy! That was a dualism of love and hate. Within the European world we find in
\textsc{Plato}'s writings (Crito, Gorgias, the first book of the Republic) for the first time\footnote[1]{There is on
this point a conflict between \emph{Xenophon}'s and \emph{Plato}'s accounts of \emph{Socrates}. According to the
former Socrates held the old standpoint that it is the sign of an able man to love his friends and harm his enemies.
\textit{(Memor.} II, 3, 14; IV, 2). --- The solution might perhaps lie in this, that Socrates himself did not express
himself definitely on this point, and that the dispute about it first arose in his school, where Xenophon (in
agreement with Antisthenes) then maintained the current conception, while \emph{Plato}, in opposition to this and from
his ideal ethics, interpreted the Master's more indefinite utterance so that the opposition between friend and enemy
does not hold at the highest standpoint. This hypothesis has been set up by Karl \emph{Joël}: \textit{Der echte und
der xenophontische Sokrates.} I. p. 396 f. (Cf. II, 2. p. 1008 f.). --- \emph{Democritus} had for that matter already
declared (fragm. 48 \emph{Natorp}) that he who does wrong is unhappier than he who suffers wrong. This utterance is
closely connected with Democritus's laying the weight on the inner disposition, the inner state of the mind, in
opposition both to external fate and to external action. Herein the Democritean and the Platonic ethics meet. Cf.
Natorp: \textit{Die Ethica des Demokritos.} p. 102 f.} an entirely different maxim set up: ``one must never inflict
evil except as educative punishment, and it is better to suffer wrong than to do wrong.'' The proposition is
introduced with great emphasis and with full consciousness of its opposition to the current way of thinking. ``This I
know well,'' says \textsc{Socrates} in Plato's \textit{Crito} (p. 49 D), ``that both now only a very few think this,
and that it will always be so. Those, then, who hold this view and those who do not hold it cannot take counsel
together; no, they must unavoidably despise one another when they see each other's way of thinking.'' The new thing is
that even the relation to the enemy now falls under the relation of love, the instinct of revenge being subordinated
to regard for a higher justice. There does not follow from this an ascetic passivity, such as would follow from the
well-known command of the Sermon on the Mount if one took it according to its wording. One does not make oneself a
mere object for others because one does not let blind lust for revenge prevail, which knows no limits but its own
satiety, and which does not have the other's good as its final goal. --- In the doctrine of the
% --- p. 247 ---
\opage{247}punitive authority of the state we shall meet again this opposition between the maxim of revenge
(retribution): ``Evil for evil!'' and the maxim of educative love: ``no evil except for the sake of the good!''

3. When we define sympathy as pleasure in others' pleasure, pain at their pain, nothing more is presupposed about
these others than that they have the capacity to feel pleasure and pain. Disinterested and universal sympathy must
therefore extend beyond the human world, and the principle of welfare will be able to find application also in the
relation to other than human beings; in its generality it demands, after all, merely that the feeling of pleasure shall
increase, the feeling of pain decrease in the world as far as possible. One may therefore very well speak of duties
toward the animals and of rights for the animals. It is not, after all, high reason and great power of comprehension,
but merely the capacity to feel and suffer, that is the condition for being the object of ethical duty and
\textit{to that extent} the subject of ethical right. As soon as pity acquires a decisive significance as an element in
the motive of valuation, it will follow of itself that duties toward the animals are acknowledged, and that these
cannot be treated as mere means, so that any treatment of them whatever should be defensible. After \textsc{Clarke}
first pleaded the cause of the animals from his peculiar moral principle, which demanded that every being be treated
according to its nature, and after with us Fr. \textsc{Eilschow}, from the idea of a great society of all living
beings, had maintained that there were duties toward the animals, it was \textsc{Hutcheson} and Rousseau who derived
the duties toward the animals from the feeling of pity: since the capacity to suffer is common to animals and men,
unnecessary infliction of suffering must be just as reprehensible toward the former as toward the latter. And later
Bentham declared: ``The question is not: can the animals think? or can they talk? but: can they
suffer?''\footnote[1]{Fr. \emph{Eilschow}: \textit{Philosophiske Breve.} Copenhagen 1748. p. 90 f. (Eilschow polemizes
especially against \emph{Pufendorf}, who in his work \textit{Et Menneskes og en Borgers Pligter efter Naturens Lov.}
Danish transl. 1742. p. 257 had declared that ``since the beasts had no reason, without which one can form no concept
of any so-called right and obligation properly speaking, there is consequently no common right for men and beasts'').
--- \emph{Hutcheson}: \textit{System of Moral Philosophy.} London 1755. I. p. 314. --- \emph{Rousseau}: \textit{Discours sur l'inégalité.}
Amsterdam 1755. p. LXII f. --- \emph{Bentham}: \textit{Principles of Morals and Legislation.} XVII, 1, 4. Note. --- On \emph{Clarke}'s
doctrine see \emph{Ueberweg-Heinze}: \textit{Grundriss der Geschichte der Philosophie.} 8 Aufl. III, 1. p. 159.} % Bentham's words: "the question is not, Can they reason? nor, Can they talk? but, Can they suffer?" -- here from Høffding's Danish
--- When we nevertheless do not count the animals
% --- p. 248 ---
\opage{248}as belonging to the ethical world, which comprises only the human race, this is for various reasons.

In the first place the animals are, to be sure, objects of duty, but not themselves subjects of duty in the same sense
as men are. They cannot, then, fill a place in the ethical world as independent members, and their ethical right can
therefore not be the same as man's either. In the second place, the animal's fate, on account of its lower and more
limited nature, does not have the great significance that a man's fate has. And in the third place the pain animals
feel can hardly reach such a degree of strength as human pain can\footnote[1]{See \emph{Panum}: \textit{Indledning til
Fysiologien.} 2nd ed. p. 7. --- Cf. my \textit{Psykologi} I, 6. --- This is granted even by the zealous friend of
animals \emph{Schopenhauer} \textit{(Die beiden Grundprobleme der Ethik.} 2. Aufl. p. 245).}. When man makes use of
the right of the stronger toward the animal and uses it as a means for his ends, this must be justified by the
significance of human ends in comparison with the pleasure or pain the animal can feel. But that the animal can at no
point be regarded as a mere means comes precisely from its being able to suffer. (Cf. VIII, 6). The burden of proof
must lie on him who would have the animals treated otherwise than men.

When one will not ground the duties toward the animals directly, from the principle of welfare, one attempts it by a
detour. One then says, for example, with \textsc{Kant}\footnote[2]{\textit{Tugendlehre.} § 16---17.}, that cruel treatment of the
animals conflicts with man's duty toward himself, ``because fellow-feeling with their fate is thereby dulled, and so a
natural disposition very beneficial to moral conduct toward other men is weakened and gradually rooted out.'' This
consideration has its significance as a subordinate motive. But fellow-feeling
% --- p. 249 ---
\opage{249}with the animals has not merely the significance of being a preparation and school for fellow-feeling with
men; it also has an immediate value in relieving a part of the pain that is felt in the world. He who helps a bird out
of a trap does a good deed, not merely because he strengthens his capacity to have fellow-feeling with other men. From
the side of the philosophy of law the prohibition of cruelty to animals is usually grounded on regard for human feeling,
which is corrupted or injured thereby. It is here rather the feeling of others than the agent's own that is
regarded\footnote[1]{\emph{Jhering}: \textit{Der Zweck im Recht.} II. p. 138 ff. \emph{Goos}: \textit{Retslære.} I. p. 169.
--- The German Imperial Penal Code forbids only public and offensive cruelty to animals (§ 360); the Danish Penal Code
forbids coarse cruelty or other cruel and revolting cruelty to animals (§ 297). --- On the other hand \emph{Löning} stresses
as the purpose of the prohibition of cruelty to animals not merely the protection of the moral feeling of the
population, but also ``the protection of the animals themselves against useless and therefore immoral cruelties.''
\textit{Sittlichkeitspolizei} (Schönberg's Handb. der polit. Oekon. 1st ed. II). p. 638. Some German states have
taken this point of view into their penal codes}. But this too is a detour or a side consideration that must not step
into the place of the direct grounding.

\begin{center}b. Love of truth.\end{center}

4. Love can be directed not merely toward single definite beings, but also toward ideas that captivate us and lay claim
to our sacrifice. A conflict can arise between the two kinds of love. There can be a passionate devotion to ideas
(scientific, artistic, political or religious ideas) that does not hesitate to trample underfoot the welfare of real,
living and feeling beings, or perhaps does not heed it at all. On the other hand the need to help in the single,
personal relations can also blind one to the significance of the general interests and endeavors on which, after all,
all higher and freer development of human life depends. Absorption in general ideas has been called
\textit{Alsind} [public-mindedness] and set
% --- p. 250 ---
\opage{250}in opposition to \textit{humanity} in the sense of care for each single human
existence\footnote[1]{Gabriel \emph{Sibbern}: \textit{Om Humanitet og Alsind.} Copenhagen 1855. Cf. F. C. Sibbern:
\textit{Moralfilosofi} § 22.}.

The single man's need for happiness, peace and comfort is the first thing and in the end the most important. All
progress can, after all, in the end be measured only by whether such need can be satisfied better and more than before.
The significance of public-mindedness, or, as we shall here call it, love of truth, does not therefore vanish. Truth
designates, after all, nothing but the connection of existence, so far as we can grasp it. Into this connection we can
and should grow; if we do not, we build on sand. What the ideas named above show us is, after all, only that we live
under more comprehensive and differently constituted conditions than we had imagined.

The demand: seek the truth! springs from the impulse of self-preservation. Man's first encounter with truth as the
expression of the fixed connection of things takes place when he sees for the first time that he unavoidably needs
certain means to reach his ends\footnote[2]{Psykologi VI F, 3 (cf. V D, 4).}. Here, then, there is not yet any
opposition between truth and feeling, between theory and practice, even if the wish to reach the goal at once stirs with
impatience. But with the further development of knowledge a relation of hostility may set in. It is not said that the
results of knowledge immediately satisfy feeling. New ideas can therefore often be carried through only with a certain
ruthlessness. In order to bring forward new points of view and the wider horizons, one must often provisionally awaken
dejection and ill will, suffering and strife. It goes here as when a statesman plunges his people into a war: he knows
that he will cause thousands of men pain, death and sorrow; so sure he cannot be that by the war he will secure and
heighten the future happiness of the people; and yet it may be necessary to venture it, if greater misfortunes would
otherwise break in. Even he who founded the religion of love of mankind said: ``Suppose ye that I am come to give peace
on earth? I tell you,
% --- p. 251 ---
\opage{251}Nay; but rather division: for from henceforth there shall be five in one house divided, three against two, and
two against three''\footnote[1]{Luke XII, 51---52.}. % King James wording, where Høffding quotes the Danish Bible
--- No wonder that he who in the name of truth would raise a great and bitter strife may for a while be held back by the
thought of the happiness and security of life he will disturb by the sharpened demands he feels obliged to assert! His
own happiness he could sacrifice, and perhaps he does not even have any happiness to sacrifice; but the suffering of
others becomes a difficulty that may hold him back for a while. In Søren Kierkegaard's journals before his violent
struggle against the usual conception of Christianity this difficulty is depicted in an inward and lofty
way\footnote[2]{See my book \textit{S. Kierkegaard som Filosof.} p. 138---140.}.

It has been held that the unconditional search for truth could not be grounded by the principle of welfare. ``We
owe,'' it has been said\footnote[3]{\emph{Lecky}: \textit{History of European Morality.} I. p. 52.}, ``more to our illusions
than to our knowledge. Our ever-working imagination probably contributes more to our happiness than reason, which in
the speculative field is chiefly critical and destructive.'' --- To this it must be answered that imagination can very
well work without our being entangled in illusions. Illusions first arise when we confuse images of fancy and reality.
All truth is in the end practical, is a light that is to guide the will. If illusion can guide the will just as well as
truth, then we lack any distinguishing mark between illusion and truth. For the only mark we have of truth is that we
can draw all consequences from it, can build on it both in our acting and in our thinking\footnote[4]{Psykologi V D, 1.
5. (Cf. V B, 4).}. The boundary between illusion and truth cannot be drawn once for all. It is an infinite task that is
set our science: to form a faithful and coherent picture of the world as the expression of our observations. The work
on the solution of this task is --- even apart from the practical yield --- bound up with a feeling of satisfaction
that can far surpass the joy one feels in rocking oneself in illusions.

% --- p. 252 ---
\opage{252}The duty to seek truth springs from the psychological relation between knowledge and will. Even the most
primitive utterances of will, instinct for example, which has so often been praised for its sureness, can err.
Instinct is always lured forth by sensations; but these sensations can occur without the conditions for the instinct's
successful working being present. The animal that is lured into the trap by the smell of the bait comes to suffer
because instinct hinders it from examining all the circumstances. We are in general inclined to attribute greater
validity to our ideas than they have, and it is only the practical consequences that lead to the needful critical
limitation. He who does not watch over the truth and clarity of his thoughts does not watch over the direction his will
takes either; and since his willing brings consequences not only for himself but also for others, the lack of love of
truth can turn into hard-heartedness and indifference to the welfare of others.

More closely regarded, there are three relations in which love of truth can show itself: the relation between thought
and word, the relation between thought and personality, the relation between the content of thought and its objective
validity.

The word is to answer to the thought, when one gives one's thought words at all. This excludes meaning one thing and
saying another. The difficult point lies here, as at so many points in ethics, in the pedagogical regard that it may
be necessary to take. Possibly a word that fully expressed the thought would not lead to so full and independent an
understanding as a half-sung song or an ironical contradiction. Such a pedagogy springs precisely from love of truth.
What conflicts with it is only that by accommodations and reinterpretations one goes on using words that have meaning
only when quite other thoughts lie behind them than those one actually has. In the religious, the political and the
social domain there is ample occasion to show this kind of love of truth --- the love of truth of honesty --- in
opposition to dishonorable compromises.

But deeper than this relation lies the relation between
% --- p. 253 ---
\opage{253}thought and personality. What matters here is whether the thought is really the property of the personality.
Truth here means \textit{personal truth,} agreement between thought and personality, an agreement that arises from the
thought's meeting with living assent from what is central in the personality, and from its acknowledgment being the
fruit of self-activity, not of blind, passive reception. \textsc{Poul Møller} even went so far (in his fragments on
affectation) as to say: ``No utterance of life has truth unless creative self-activity lies in it.'' Affectation
consists for him in taking thoughts into oneself without making them one's real property, or rather one's own work. But
affectation is, after all, excluded where the assent is self-active, although self-activity cannot always be called
``creative.'' There must, besides, often be times of transition when the new thoughts (whether they are produced by the
personality itself or by others) have not yet become personal property, and it would be unjust to stamp the state in
such fermenting and seeking times as affectation. \textsc{Søren Kierkegaard} continued Poul Møller's fragmentary
thinking by his passionate stress on the subjective relation to the thoughts that determine the view of life. On these
points he demanded a thinking ``with the pit of the heart,'' --- a ``subjective thinking,'' a ``thinking in
existence,'' a thinking that goes into one with a willing. He thereby introduced a great measure of value: views of life
are to be tested not merely by their objective content, but especially by the root and hold they have in the
personalities who hold them; --- and conversely, holding a view of life now no longer consists in rocking oneself in
certain thoughts or articles of faith, but in living and dying on them as the inmost hold one can find. And in his own
life and his own thinking \textsc{Kierkegaard} himself gave a great model of such existential
thinking\footnote[1]{\textit{Søren Kierkegaard som Filosof.} p. 24---26; 44---53; 127---129; 146.}. He thereby rendered
mankind the greatest service he could render it according to his nature and the conditions of his life: to undertake
the great attempt to test a traditional --- often in a very external way
% --- p. 254 ---
\opage{254}traditional --- view of life's relation to personal life in our day, to examine its carrying capacity and the
possibility of fulfilling its real demands. An experiment unique of its kind, which perhaps gives more important
contributions to the valuation of a view of life than most learned inquiries into its objective content!

But the validity of the objective content of views too must be tested without reserve. The subjective thinkers, who lay
so great weight on the relation of the personality to ideas, are as a rule\footnote[1]{Cf., besides Kierkegaard:
Pascal, Rousseau, Jacobi, Carlyle. (See the relevant sections in ``Den nyere Filosofis Historie'').} inclined to
underestimate the significance of objective criticism and inquiry. In their zeal to make the truth personal they forget
the duty of examining whether it is really the truth to which they relate themselves. Against their will they thereby
easily come to reassure others in their habitual relation to what is traditional. They overlook that self-activity is
to show itself precisely in testing and finding the truth. As the third kind of love of truth there must therefore be
named \textit{intellectual probity,} which prevents one from accepting and pronouncing something as objective truth
without sufficient testing. It is through this virtue that love of truth can lead over from the single person's
striving to preserve the coherence between word and thought and between thought and personality to scientific inquiry.
This virtue has rightly been called ``the youngest and hardest of all virtues.'' It could not be enjoined before
definite scientific methods had formed, and before several ``truths'' had transformed themselves into just as many
problems. In its sharpest form it is set out by W. K. \textsc{Clifford} in his essay \textit{The Ethics of Belief}
(Contemporary Review 1877). He maintains that a man's opinions or beliefs are by no means a private matter, so that he
should be entitled to take account only of what appeals to him or comforts him. Truth is a sacred inheritance from
generation to generation, and everyone has a contribution to make to it. To make oneself credulous is therefore a crime
against
% --- p. 255 ---
\opage{255}mankind. If one cannot win the truth itself, one is to test the probability, and if one cannot oneself test
the truth, one is to test the authorities one relies on, --- test them not merely as regards their honesty, but also as
regards the possibility of their having found the truth. Tradition does not thereby lose its significance: it is to
give us questions to ask and means of trying to answer them, but not finished results; it is to make it possible for us
to ask new questions. And if in our assumptions we go beyond what experience teaches us, this may happen only by virtue
of the presupposition that what we do not know resembles what we know. --- These demands have been found too strict when
they are made of all men, and it has rightly been maintained that involuntary faith and trust do, after all, precede
doubt, and that long experience may be needed before an independent doubting can arise\footnote[1]{These misgivings have
been put forward by a man who is himself in practice one of the finest examples of love of truth (both as honesty, as
striving for personal truth and as intellectual probity) in our time: \emph{Christoph Schrempf}. See his criticism of
Clifford in his journal \textit{Die Wahrheit.} Stuttgart 1894. I. p. 163 ff.}. From an entirely different side too
Clifford's assertion has been criticized, it being stressed that the courage to venture is a condition of finding the
truth. Better to be fooled many times than to lose the possibility of guessing right! Does not Clifford behave like a
general who tells his soldiers that it is better to keep out of the battle than to be wounded? And an analogy has been
drawn between the involuntary promptings that often lie at the basis of discoveries and the spontaneous variations that
play so great a part in the struggle for existence\footnote[2]{William \emph{James}: \textit{The Will to believe.} New
York 1897. p. 18 f.; 93; 249. --- Cf. with this \emph{Leslie Stephen}'s remarks on Mill's book ``On Liberty,'' in which a
view related to Clifford's is developed. \textit{The Utilitarians.} III. p. 255---259. Stephen stresses especially that
inquiry is a great common task, in whose treatment the one must have help from the other. Mill isolates the single
individuals too much and lays too great weight on even well-grounded assumptions finding contradiction.}.
--- But not only with regard to the uncritical faith which the Church
% --- p. 256 ---
\opage{256}has enjoined as the highest virtue, but also with regard to the great popularization of scientific results
in narrative and descriptive form, it is of great importance that intellectual probity come forward into the front rank
of the virtues. That is the condition for intellectual work not lapsing or being underrated. And especially of those
who are to be public teachers in church and school this virtue ought to be required to a quite different extent than
has hitherto been the case, --- and required in the name of love of truth.

Intellectual progress not only hangs together very closely with progress in all other domains, but is itself the form
of progress that is easiest to demonstrate. Even those who otherwise do not assume that any progress takes place admit
it in the domain of knowledge. Feeling and will change and develop more slowly, and their advance is to a very great
extent bound up with that of knowledge. The proposition, current for a time, that ideas rule the world is not
psychologically correct; the play of history is driven by men's feelings and passions; --- but the clarity and truth
of thoughts is one of the most important of all conditions for the healthy development of feelings and passions. ---

The interest in truth has this in common with ethical feeling, that it leads us to look away from our accidental
individuality and our accidental circumstances in order to ask about what is universally valid. He who cannot bear the
truth cannot see beyond his own isolated individuality; therefore shrinking from the truth leads to egoism, if it does
not spring from it in the first place. Religious ethics has often had a pernicious effect by setting up faith as the
only meritorious thing and branding doubt as sin. It has not seen that an honest and reasonable doubt is both healthy
and good. It easily comes to set a premium on stupidity; for not all simplicity is ``holy simplicity.'' If only one
barrier is set to the striving for truth, it will soon bring more in its train. \textsc{Coleridge}, who was himself a
believing Christian, has said: ``He, who begins by loving Christianity better than Truth, will proceed by loving his
own Sect or Church better than Christianity, and end in loving himself better than all.''% Coleridge's English (Aids to Reflection, Moral and Religious Aphorisms XXV); Høffding: „Den, som begynder at elske Kristendommen mere end Sandheden, vil siden elske sin egen Sekt eller Kirke mere end Kristendommen og vil ende med at elske sig selv mere end Alt“.

% --- p. 257 ---
\opage{257}It is not always glowing faith that leads to regarding doubt as something evil. In our days it is very
often fear of the unknown, of lacking a foothold when one lets go of the accustomed ideas, of the unrest and toil that
it costs to win a new view. With regard to him who cannot give up an opinion because he finds in it his only comfort,
his only refuge, no one will raise doubt and discussion. One should not rob the poor man of his one lamb. % cf. 2 Sam. 12:1--4
But from there it is a great leap to the assumption that what is truth depends on our need and our wishes, --- a leap
which is nevertheless made by many with astonishing ease.

5. From the fact that the truth is always to be sought it does not follow, as already indicated, that it is always to
be \textit{spoken.} The duty of truth aims at bringing truth to rule; but this end can often be hindered by speaking the
truth. By uttering the truth too early one behaves like the child who pulls the plant up out of the earth to see how
deep its root has gone down. One establishes what has been reached --- but may thereby prevent more from being reached!
Self-appointed apostles of truth often sin in this way. The duty to speak the truth is pedagogically limited. Truth is
so great and comprehensive that man can be initiated into it only step by step. The part of the truth that he cannot
yet grasp he easily misjudges and despises. What matters, then, is to speak the truth in such a way that this remainder
becomes in each single case as small as possible. And what matters is to find the form in which truth can best gain
entrance, and the gradual manner of disclosure that can secure its being understood. This holds especially when the
talk is of ethical ``truths.'' There is much uncalled-for moralizing in the world, and it is constantly overlooked that
the uttering of ethical judgments is itself ethically conditioned, and that all valuation stands in the service of an
end, is a means of education (see III, 11). To utter ethical judgments where they are not needed is barbarism or
Pharisaism. Indirect education is often the right one. Yet it can sometimes be precisely pedagogically right to arouse
offense and exasperation and to set up a sign of contradiction. % Modsigelsens Tegn: cf. Luke 2:34
It is often only through a passionate struggle for and against that the truth properly comes forth before us.

% --- p. 258 ---
\opage{258}With this the question of the justification of the white lie % Nødløgn
is essentially settled. Just as little as one is obliged to say of one's own accord everything one thinks, just so
little is one obliged to answer all questions. The murderer who is looking for his prey, or him who would coax from me
a secret that it is my duty to keep, I am both entitled and obliged to lead astray. The same can hold in many smaller
relations. He who declines an invitation cannot always give the real reason, because the one inviting would not
understand it, or because it would wound him. But we surely take the matter all too lightly in this respect in daily
life. Partly from untimely politeness, partly from indolence, partly from lack of courage, we extend the use of the
white lie far further than is necessary and justified.

6. When we hold fast to the view that truth is in the process of becoming, tolerance and reverence % Pietet
become easily compatible with love of truth. We all find ourselves, each for himself, at different stages of
approximation to the truth. A sharp boundary line between some who have the truth and others who do not have it cannot
be drawn. Truth is so mighty and so closely bound up with reality that no one, even with his best will, can completely
isolate himself from it. It is, after all, reality, contact and interaction with real circumstances, that educates to
truth, and no one can wholly escape having a share in this education. But the education can be long and laborious, and
how far the individual gets depends on manifold conditions.

\textit{Tolerance} (to keep this ugly word for a beautiful thing) does not bargain away anything of the truth, but
springs from insight into the hard conditions under which knowledge comes into being. He who reasons thus: ``there can
be only one truth; if I put up with others' thinking and teaching what conflicts with my conviction, I admit that what I
assume is \textit{perhaps} not truth!'' --- he ends as a persecutor. It was by virtue of this line of thought that the
ancient Church early entered upon the ways of persecuting heretics\footnote[1]{See already \textit{the Second Epistle of
John} 10---11: ``If there come any unto you, and bring not this doctrine, receive him not into your house, neither bid
him God speed: For he that biddeth him God speed is partaker of his evil deeds.'' % KJV; Høffding quotes the Danish Bible
--- In the 2nd century orthodox Christians were not allowed to converse with the excommunicated and with heretics.
(\emph{Lecky}: \textit{History of European Morality} I, p. 451). --- Gregory of Nazianzus (in the 4th cent.) declared
that one could not permit the heretics to hold assemblies without admitting the truth of their teaching.
(\emph{Barbeyrac}: \textit{Traité de la Morale des Pères.} p. 170). --- In the 17th century many zealous Lutherans
shunned the Reformed to such a degree that they would not sit at table with them or live in the same house with them.
(\emph{Gass}: \textit{Die Lehre vom Gewissen.} p. 164). --- Dr. Johnson was much offended on hearing that one of Hume's
theological opponents had dined with Hume, shaken his hand and exchanged visits with him. ``An infidel,'' he said,
``ought not to be treated handsomely by a Christian merely because he endeavors to commit his robbery with a certain
sincerity.'' % translated from Høffding's Danish
(This utterance, as well as a whole series of utterances by Johnson in the same tone about Rousseau, Voltaire, Hume and
A. Smith, is found in the \textit{Edita und Inedita Schopenhaueriana} edited by \emph{Griesebach}. % PRINTED AS IS: Griesebach for Grisebach (read at the image)
Leipzig 1888. p. 127---132). Examples from our own day could also be cited.}.

% --- p. 259 ---
\opage{259}One can very well be fully convinced of a truth and yet see why others do not accept it. And if one does
not see the cause of their not accepting it, one will not be able to find the right way of convincing them. By
searching for this cause, one will at last get beyond the domain of thought. One will then find, not only that there
are feelings which can hinder the development of clear knowledge, but also that there are feelings which can subsist to
a certain degree independently of the development of knowledge. Precisely the most important feelings --- love of
mankind, conscience, etc. --- are compatible with the most diverse theoretical and religious assumptions. What the
freethinker regards as a meaningless dogma is for the believer an indispensable thought, the expression of a need of
feeling which the serious freethinker perhaps also knows in his own way, though he gives it vent in another manner. By
seeking out the life of feeling that lies behind knowledge, one will find what is common and pave the way for a
sympathy that from a distance might seem impossible. Tolerance rests, where it is more than indifferent passivity, on a
faith in the human kernel in all, a faith that is one with love itself.

Just as tolerance is the feeling toward those who, contemporaneously with us, fight for truth in another camp, so
\textit{reverence} ties a bond between us and the representatives of the past.
% --- p. 260 ---
\opage{260}We stand in a relation of reverence to that which has been an authority for us, and whose school we have
left without bitterness. We have an understanding of its significance which he cannot have who comes as a stranger and
looks at it from outside. We relate ourselves to it as to our own past; even if we have broken with it, we do not
wholly give it up, so long as we hold fast to the continuity of our development. In a similar way reverence is the
feeling in which the continuity of the race lives; it forms the bond between the changing generations. Often it is
precisely the next generation that has difficulty in preserving this bond. The 18th century had no great reverence for
the Middle Ages; that came only (and in excess) with the Romanticism of the 19th century, which in return had no great
reverence for the 18th century; and in our time, which is hostile to Romanticism, one in turn feels more reverent
toward the 18th century. --- Reverence can be felt both toward single persons and toward a whole tendency or a whole
generation. He who in his relation to the authorities whose dominion he has withdrawn from has not yet got so far that
he can recognize and sympathize with their historical significance, is not yet free, but is dependent on his opposition
to the past. Reverence is the true historical feeling. Everything that has had value lives on in it. It is one of the
most important forms under which the individual lives himself into the race.

% --- p. 261 ---
\opage{261}\partitle{SOCIAL ETHICS}

% --- p. 263 ---
\opage{263}\chapterhead{XIII.}{INTRODUCTION AND DIVISION.}

1. Social ethics differs from sociology (the history of civilization) in a way similar to that in which individual
ethics differs from psychology. While psychology investigates human ideas, feelings and expressions of will as purely
natural-historical phenomena and seeks only to understand them in their origin and development, it is the business of
individual ethics to \textit{evaluate} them according to their relation to the ideal for the life of the individual. And
while sociology studies human social life in its various forms, at the various times and among the various peoples, and
stands toward these forms as natural science stands toward natural phenomena, social ethics will evaluate them according
to their relation to the ideal for human social life. Ethics everywhere uses the given as the point of departure for a
new development, and when it undertakes a valuation of what is astir in society, this is because this valuation
exercises a determining influence on the future course of development. How far this influence can extend, only
experience can teach. As soon as we evaluate, we presuppose the possibility of intervening to bring about change. This
holds not only for ethics, but for every practical science. When general jurisprudence investigates not only which
relations are legally regulated, but also which are \textit{suited} to it, there lies in this an ideal which may need
much struggle to be realized. Political economy
% --- p. 264 ---
\opage{264}investigates not only actual production and distribution, but criticizes these insofar as they do not
accord with the conditions of a healthy and happy social life. Both jurisprudence and political economy lead back to
ethics as their general foundation. (Cf. III, 1. 14). ---

Just as every trait of character and every act of will that individual ethics requires must be psychologically
possible, so too every social order that social ethics requires must be sociologically (historically) possible. Yet an
ethical ideal, even if it cannot be realized for long ages or perhaps ever, can have great significance through the
direction it gives to the human mind and to human powers, provided only that fantastic overexcitement and misjudgment of
real circumstances are kept away. New possibilities of development are often discovered only by him whose gaze is not
fettered to what is momentarily given, but who is able to rise above it. It is an important turning point where the
conscious setting of ends and ideals replaces blind, instinctive action. This turning point is reached through
historical development, and the conscious valuation which thereby becomes possible acts back in turn upon the further
development.

2. Seen from a twofold side, all ethics, and especially all social ethics, acquires a historical character. It has
already been stressed earlier (III, 13) that evaluating subjectivity always proceeds from a psychological foundation
which is not the same at the various times. But here we now see that the material and the objects too, with which
ethical valuation and ethical work are concerned and which they treat, present different properties at the different
times. From the same foundation, then, one will reach a different result at different times. What is possible at one
stage of development is not possible at another. One cannot construct from the general principle of welfare an order of
family or state which could without further ado be introduced everywhere. Any social order whatever presupposes certain
inner and outer conditions given in those who are to make up the members of the society. No social building can be
erected on bare ground. It is, as history amply shows, easier to overturn old
% --- p. 265 ---
\opage{265}institutions and ``introduce'' new ones than to alter the dispositions and impulses in human nature which
gave the old institutions their vitality. In the spiritual and social domain we always have old wineskins for the new
wine, and the wine always acquires a certain taste from them. Often the new wine must burst the old wineskins before
they can be replaced by new ones. % cf. Matt. 9:17

The historical relativity which thus necessarily marks every one of the results of ethics has often, once men had had
their eyes opened to it, been exaggerated and misunderstood. It does not mean that everything can be equally good or
equally bad. It means that what is to be good under certain definite circumstances must correspond to them. The ethical
ideals are individualized in different ways at different times and in different places; and this is so far from
conflicting with their nature that an ideal which could not be individualized would be empty fantasy. But the ethical
ideals have more the character of general tendencies and directions than of formulas which could without further ado be
applied under any circumstances whatever. What we can attain is that our actions, and the forms in which social life is
ordered, come to lie in a series that runs in the direction of the ideal, and as far forward in this series as
possible. We do not get beyond constant striving. Relativity means precisely that only a constant approximation is
possible. (Cf. III, 12; IV, 3---4 and VII, 3). Here it must not be overlooked that the ideals themselves undergo changes
in the course of progressive historical development, since there are discovered both new possibilities, new demands
--- and new limits\footnote[1]{Cf. further on historical relativity: \textit{Etiske Undersøgelser.} p. 47---51.}.

3. The ethical process of development is partly a process of humanization, partly a process of emancipation. --- At the
lower stages of development the animal instincts rule, to which the preservation of the individual and of the race is
tied. And they do not cease to operate because ideal and sympathetic feelings develop. They even have in their nature a
violence against which the higher and gentler feelings are often powerless. When higher interests or
% --- p. 266 ---
\opage{266}values --- that is to say: such as point beyond the moment, indeed perhaps beyond the whole life of the
individual and of the generation --- assert themselves, this does not happen through new energy being created. Ethics
has no reason to deny the law of energy. But energy can be converted from its more elementary forms of expression into
forms in which it can serve higher values, and moreover energy may be released for which there was no use in the
service of more limited interests. What matters is to transfer the energy which at first is seized upon by the animal
instincts to the ideal, more human feelings, and thereby occasion will also arise for hitherto slumbering powers to be
awakened. Strict asceticism is here a revolutionary principle, since it demands outright that those instincts be
killed. It is only through a slow course of development that the instincts, under the influence of the conditions of
life --- which by no means makes conscious exertion of will superfluous --- can be softened and ennobled. Besides this
\textit{process of humanization,} by which we have gradually become and are becoming men from animals, there also goes
on a constant \textit{process of emancipation,} by which we gradually become free men. Bodily and spiritual compulsion
is characteristic of the lower stages of ethical and social development; the relation of authority rules in its
stricter or milder forms. Nor can this either --- as the 18th century in its revolutionary zeal believed --- be
abolished suddenly. Through the educating influence of authority, and through the free choice of authorities as
intermediate links, the development of the independent personality takes place.

Every individual, every generation and every people finds itself at a certain stage of this process of development (in
its twofold form), and all institutions, customs and forms of life will bear the stamp of it, however one goes about
things. What matters is which institutions and forms of life can best take the powers and impulses at hand into their
service, in such a way that they at the same time carry development upward toward a higher stage.

4. From an absolutely individualistic standpoint society can only be conceived as having arisen through the union of
individuals. In the more or less conscious thoughts and feelings of individuals
% --- p. 267 ---
\opage{267}one then finds the explanation of the kind and form of social life. One presupposes independent individuals
given prior to the rise of society. But every individual in fact \textit{finds} a society already there, within which
it develops, and only after the individual's development has received its definite stamp from the circumstances and
traditions of this given society can the individual in turn act back determiningly upon society. Before it becomes
clearly conscious of itself, it lives and acts, thinks and feels only as a member of society. The involuntary always
precedes the voluntary. There takes place an education of the will, conditioned by the unconscious influence of social
circumstances, which is every bit as important as the education that goes on through conscious intervention and
influence.

This too belongs to the psychological and ethical truths for which \textsc{Aristotle} already had a clear eye. His line
of thought is roughly the following:\footnote[1]{Eth. Nic. II, 2.} Qualities (e.g. ethical virtues) are produced by
performing the corresponding actions. By building we become skilled builders, and by practicing justice we become just.
What we are to carry out consciously with skill, we must first be trained in. Therefore it is of great importance
whether the social order from the very first leads the young person in the right track, guides him to activity in the
right direction. The ethical significance of the social order depends on what habits it imparts to individuals. --- In
what follows we shall very often have use for this proposition: that involuntary training precedes fully conscious
practice, and for brevity's sake we shall call it \textit{the Aristotelian principle.} This principle finds application
prior to the apparently opposite principle, that what is first practiced with consciousness can later be practiced
unconsciously, as habit, or as instinct. For at the outset the individual does not know why it is training itself in
certain ideas or actions; this happens involuntarily and partly unconsciously. The involuntary always precedes the
voluntary, and the very transition from the involuntary to the voluntary takes place involuntarily. ``The first thought
of the will'' (see VI, 1) emerges without intent. We
% --- p. 268 ---
\opage{268}have no control over the rise of our first ideas, on which in turn our conscious willing, our voluntary
action, depends.

Just as there is thus a constant interaction between the conscious and the unconscious, and between the voluntary and
the involuntary, so too there goes on an interaction between the will on the one side, and its means and its results on
the other. What the will first seized upon as a mere means can later become an end for it, and the result in which it
thought to come to rest can become the point of departure for new striving. There can take place displacements of
motive, and there can take place displacements of value.

According to \textit{the law of the displacement of motive,} what is first done from one motive can later be done from
an entirely different one, the original means having become an end. Figuratively this can be expressed by saying that
the individual's interest has shifted. (Cf. I, 4). Now since it is not said that the original motive wholly
disappears, there can be bound up with the displacement a fusion, a kind of chemical compound of different motives,
through which there arises a new motive, a new feeling, whose properties cannot be immediately derived from the earlier
feelings from which it has arisen. By the help of these psychological laws\footnote[1]{Cf. on these psychological laws
my \textit{Psykologi} VI B, 2 d; C, 2 and 5; E, 4---5 (cf. II, 6 d; V B, 2). --- The law of the displacement of motive
has been demonstrated especially by Spinoza, Hartley and James Mill, the phenomenon of fusion by Hartley and F. C.
Sibbern.} it is understood how the social interests of life can stand directly before the individual as the highest,
and become the central point of view to which he finally returns in his valuation and in his plan of life. It is such a
standpoint that we have here chosen as the foundation for our systematic valuation. (Cf. III, 9---13). Only on such a
standpoint does it follow of itself that social ethics includes individual ethics. The process by which this standpoint
comes into being is now of great interest for social ethics itself, since it must naturally be a main task of the
latter to understand the foundation on which its own position of principle depends.

% --- p. 269 ---
\opage{269}Experience shows further that from the involuntary or voluntary, self-loving or self-sacrificing actions of
individuals there spring far more and often wholly different effects than were foreseen or could be foreseen, let alone
intended, by those who acted. From the interaction of individual interests and passions there can spring results of
lasting and comprehensive significance. Often something stands as a last, concluding end that bounds the horizon of the
will, and it may yet afterward prove that it can be used as a means to still more distant ends, or that it possesses
values of another kind and of greater range than had been imagined beforehand. There are, after all, both powers that
would do harm and yet come to do good, and powers that would do good and yet come to do harm. But the main thing is the
widening of the horizon of values which can take place, and which constantly takes place in history. An example is
offered by the development of the feeling of humanity after Alexander's conquests. For Alexander the founding of a
world empire was the final goal; but it proved that when this goal was reached, there had thereby been brought about
conditions for such an interaction, such an intercourse between the various peoples hitherto kept apart, that a
general feeling of humanity could develop. The values (potential or actual) that are produced in the course of human
development rest to a very great extent on this \textit{objective displacement of value} (as one might call it in
distinction from the displacement of motive as a \textit{subjective displacement of value).} In the brilliant
introduction to his \textit{Philosophie der Geschichte} \textsc{Hegel} has pointed out this phenomenon, ``that in world
history, through the actions of men, something other comes about besides what they intend and attain, indeed besides
what they know and will at all.'' He calls it \textit{the cunning of reason} that it [world-reason] lets the passions
work for it: ``The idea [world-reason] pays the tribute to existence and transience not out of its own pocket, but out of
the passions of individuals.''\footnote[1]{This unconscious production of more comprehensive values \emph{Wundt} calls,
with an expression borrowed from zoology, ``die Heterogonie der Zwecke.'' Ethik.\textsuperscript{3} I p. 275.} The
results that are thus
% --- p. 270 ---
\opage{270}realized right through men's conscious work concern ethics, however, only insofar as they can really be
called values and as such enter into relation to will and feeling (III, 1 and 14). If one stretches the consideration so
far that the last results, the final goals of all human development, are to lie beyond everything that could possibly
become an end for human willing and stand as valuable for human feeling, then one has overstepped the boundary of
ethics and passed over to the philosophy of history and the philosophy of religion (or, if one will, metaphysics). It is
then a transcendent result, overstepping all experience, that stands as the final goal of human development. In this
direction goes \textsc{Hegel}'s conception and, in most recent times, \textsc{Wundt}'s; the latter, however, concedes that
that transcendent final goal is not and cannot become known to us. But even if one will not leave the domain of ethics
and pass over into the regions of speculative thought, the law of the displacement of value contains a thought of great
significance, especially as against a short-sighted and short-winded conception of the principle of welfare.

In the displacement of motive the connection between the old and the new is closer than in the objective displacement
of value. In the former there is a mode of action or an order of life which in essentials remains as it was, but is
maintained and acknowledged for new reasons; in the latter it is an entirely new mode of action or order of life that
arises, and the reasons which make it valuable may then be such as were formerly not known at all. The social result
which Alexander's conquests brought about by breaking down the national barriers around the eastern Mediterranean not
only could not be foreseen, but could not be acknowledged either from the classical Greek standpoint. Common to the two
kinds of displacement is this, that the acknowledgment of the new values does not necessarily entail acknowledgment of
the way in which they came into being. Acknowledgment of the significance of private property does not entail
acknowledgment of the way in which private property has historically developed, and the feeling of humanity need not
lead to unconditional admiration for Alexander's lust of conquest. A sharp opposition can arise here between the
ethical and the historical mode of consideration. The
% --- p. 271 ---
\opage{271}solution of this conflict must be sought in this, that ethical valuation is by its nature directed toward the
future, toward what is to happen. It takes life up on the foundation which, as things are, has actually been brought
about. Just as little as the end always sanctifies the means, just so little need any luster shine back from the new
values attained upon the ways by which they were reached. When Mephistopheles wills evil but always ends up producing
good, we do not thank him for it; at most we can put up with there being fellows of that sort. But ethics does not care
to sit in judgment on the past either; it regards what has been produced as the soil on which work is now to be done.
After revolutions of the earth that have led to subsidences and upheavals of the strata, plant growth and animal life
spread over the new terrain, however uneven and topsy-turvy it may be. It is not in the power of life to undo the
revolution (even if perhaps much that was glorious has perished); but life is able by its quiet power to bring forth a
new world on the ruins of the old --- if only the germs of life have survived the catastrophe. The relation between
ethical life and the historical world processes must be conceived by analogy with this.

The displacement of motive speaks rather in favor of the conservative conception, the objective displacement of value
in favor of the revolutionary conception, the former showing the possibility of new reasons for the acknowledgment of
the old, the latter the possibility that the suddenly arisen new will also find new motives of valuation.

There can be an inner connection between the two kinds of displacement, inasmuch as the condition for a means becoming
an end can consist in the result's being \textit{anticipated} during the work and thereby giving the work a share in the value
it itself has. A road we travel to a joyful goal can itself become joyful. But this does not hold for all cases. Often
it is the work \textit{itself} that becomes attractive through the use of powers that it entails. It is then not the
result that gives the work a share in its value, but it is the work that itself wins a new value\footnote[1]{\emph{Ehrenfells} % PRINTED AS IS: Ehrenfells for Ehrenfels (read at the image)
\textit{(System der Werttheorie.} I. p. 137 ff.) designates what I call displacement of motive as ``Wertbewegung nach
abwärts,'' i.e. a retrograde movement of value, in contrast to the objective displacement of value as ``Wertbewegung
nach aufwärts.'' But not all displacement of motive takes place ``abwärts.'' --- Cf. moreover my review of
\emph{Ehrenfells}' book in ``Göttinger gelehrte Anzeigen.'' 1900. No. 9.}. But of
% --- p. 272 ---
\opage{272}both these subspecies of displacement of motive it holds that continuity is greater in them than in the
objective displacement of value.

Common to the Aristotelian principle and the two forms of displacement is this, that conscious acknowledgment comes only
afterward. Narrow limits are thereby set to every deductive ethics; the necessity of experience is enjoined with
emphasis. And at the same time all three processes entail a decisive difficulty for negotiation between standpoints of
which the one lies before, the other after, one or more of the metamorphoses described. (Cf. ch. III).

5. It has often been thought that guidance for understanding the relation between individual and society could be found
in the analogy between society and an organism. In various forms this analogy is carried through and applied by Plato,
Hobbes and \textsc{Spencer}. Society is then regarded as a great man, or as an organism whose various organs or cells
the single individuals are. Many interesting considerations can also, as \textsc{Spencer}'s sociology in particular
shows, be undertaken with the help of this analogy. But it has its definite limit, at which we shall here
dwell\footnote[1]{It is moreover already brought forward by Spencer: \textit{Principles of Sociology.} I, p. 478 ff.},
since the relation between individual and social ethics is thereby more closely illuminated. (Cf. III, 17 and VIII).

In the organism the single elements are unconscious, and only to the nervous central organs is consciousness attached.
This centralization is the more prominent the higher the organism stands. The state and activity of the single cells or
organs are in the end evaluated according to their influence on the state of these central parts, which appears in
consciousness as pleasure or pain. Since the central parts can be said to represent the whole organism, all cells and
organs thus become subordinate means for the welfare of the whole organism. --- In society it is
% --- p. 273 ---
\opage{273}on the contrary precisely all the elements, the single members, that have the capacity to feel pleasure and
pain; this capacity is not confined to one center. Only by fantastic mysticism can one ascribe consciousness to society
as a whole, apart from the single individuals. The welfare of society is the welfare of the single individuals. Society
is a connection of personal beings, but not itself a personal being.

Only apparently does this view contain a contradiction of the subordination of individual ethics to social ethics
(VIII). Although society and the race consist of individuals and in any given place come forward before us only through
certain definite individuals, the concepts society and race have nevertheless the great significance that, over against
the single individual, they represent the sum of \textit{all} individuals (itself included as one of the many). The
concepts society and race have, as was remarked earlier (VIII, 4), the ethical significance of pointing partly to the
incalculable range of the interests of life that must be taken into account, partly to the necessity of bringing about
possibilities, conditions, potential values, and not stopping at what has merely actual value and perhaps has this only
for a short time. They thus make possible the most comprehensive point of view for the valuation of the single
individual's willing and acting. They require that the individual's willing and acting be regarded not merely from its
own standpoint, but also from a higher one, just as Copernicus required that one should regard the earth not merely from
its own standpoint, where everything revolves around it, but from the standpoint of the sun, from which it is itself
seen to share in the motion.

6. The highest idea of social ethics is the idea of \textit{the kingdom of humanity,} a society of harmoniously and
richly developed personalities. Such a society is the more perfect, the more distinctive and independent each single
personality is, and the more intimate and firm the connection between the single personalities is, --- the more, that
is, both multiplicity and unity at once come into their own. (Cf. III, 10). This idea springs from the principle of
welfare. For the greatest welfare must be present where each single individual develops in an independent way and
thereby at the same time, both consciously and unconsciously, helps others
% --- p. 274 ---
\opage{274}to a similar development from their own position. The chief virtues of individual ethics: justice,
self-assertion and devotion, spring directly from this. And not only are the development of character of individuals
and their mutual harmony made possible hereby. There is also made possible the production of results which the
individuals each by himself would not be able to reach, and which are yet of significance for the development of the
race as a whole. The order of society and the works of culture are such results. They become intelligible only through
this, that --- by virtue of the Aristotelian principle, of the displacement of motive and of the displacement of value
(4) --- something other and more can come out of the work of individuals than they think of in their philosophy.

The idea of the kingdom of humanity is formed by simple combination of the two concepts society and welfare. But the
more special forms of ethical social life cannot be derived from it. We can use it as a standard in valuation; but since
the possibilities of development are different at the different stages, the demands in particulars too may turn out
very different.

The various stages of human development differ, as regards social life, not least in this, that at lower stages only
slight difference is found between various kinds of society; here as in other domains it is differentiation that is
especially characteristic of the higher stages.

7. \textit{The various kinds of society} between which one can distinguish at higher stages of development differ from
one another partly by the nature of the \textit{forces} that bind the individuals together in the society, partly by the
\textit{ends} that are striven for in the society, partly by the greater or smaller \textit{circle} of individuals that
the society embraces.

In \textit{the family} it is instinctive sympathy that is the binding force. In the relation between mother and child,
the kernel of the family, it shows itself most clearly and most strongly. It is the bond of blood that here ties human
beings together. But motives other than the elementary sympathetic instincts also operate in the family at its various
stages of development. At primitive stages of development women and children are regarded predominantly as the property
of the men,
% --- p. 275 ---
\opage{275}and as subject to their dominion; and insofar the primitive family recalls the state. It recalls the
church, insofar as it has its cult and its sacrifices to the spirits of the ancestors. And it is a society for the
furtherance of culture, insofar as it provides not only for physical care and preservation, but in addition for the
development of capacities and skills. At higher stages of development too the family is not only the constant source of
nourishment for the sympathetic feelings, but also the first power that leads man into culture and awakens the
consciousness of a fixed social order.

Another kind of society arises when common, or at least connected, interests and ends unite human beings. The one often
attains his ends only when he helps the other to attain his. Such is the relation between seller and buyer, in all
exchange. To the offer of the one corresponds the demand of the other; the one has in surplus what the other lacks and
desires. Here, then, the individuals supplement one another. Their interests are connected, although there need not
otherwise be any closer or further understanding and sympathy between them. --- But often the ends are common and can
be reached only by united forces. Security, for example, is a common good which is reached through union. --- The
cooperation of one or another of these motives gradually produces (by virtue of what was remarked in 4) a feeling of
community, a fellow-feeling, which can become disinterested sympathy, even if at first it was egoism that drove the
individual to seek intercourse with others. Intercourse resting on the mutual supplementing of interests, or on the
community of interests, can thus have an educating effect. It is of great importance that the outer conditions of life
already have a binding effect. --- It is not only egoistic ends, however, that can in this way lay the foundation for
intercourse and society. There are also such ends as do not directly play a part in the egoistic struggle for
existence. Science and art, religion and love of mankind set human beings in common motion and bring them together
around common and unselfish ends. --- In both cases, both where it is at first the individuals' struggle for
self-preservation that brings them
% --- p. 276 ---
\opage{276}together, and where it is ideal interests that are the binding force, there can arise a feeling of
brotherhood, determined partly by the personal union, partly merely by the consciousness of standing in the service of
the same great ends. --- This kind of society, where the binding force lies in the ends, and where neither elementary
sympathetic instincts nor outer compulsion rule, we may call \textit{the free cultural community.} Its boundaries are
not so narrow as those of the family; indeed, they extend much further, just as far as culture is cultivated; and that
is to say, again, that in the end they coincide with the boundaries of the human race.

Finally there is a society which differs from the two preceding ones in that it builds not only on natural sympathy or
on the binding force of interests, but in the last instance rests on the use of force and compulsion. As regards ends,
\textit{the state} has, or can have, much resemblance to the family and the free cultural community. It will not only
preserve and secure the life of its members, but also work for their advance in culture, both material and ideal
culture. But wherever it acts, power stands in the background. In the law that holds in the state we have the sum of
the rules laid down for the use of force. Within the family and the free cultural community too we can speak of law, if
we understand by it usage and custom, such as have imperceptibly formed themselves in the course of development and
lead involuntarily to treating new relations and cases in a manner similar to earlier ones of a similar kind.
Historically we scarcely find a family or a cultural community without outer compulsion of one kind or another, and the
state on the other hand regards it as one of its most important tasks to protect culture and the family. The three
societies cannot be kept wholly apart from one another, but designate different ways of viewing human society. The
state, to be sure, has power as its last argument, but also only as the last. It seeks to draw advantage from the same
forces that rule in the family and the cultural community; precisely for that reason it is not indifferent to it how
these develop, and its own existence becomes all the more firmly founded the less it is mere power that rules. In
extent the state differs both from the family and from the free cultural community:
% --- p. 277 ---
\opage{277}it is larger than the former, but smaller than the latter. A large family can be called a small state. It
has been said that a family is a state when it cannot be subdued without war\footnote[1]{Thomas Hobbes:
\textit{Leviathan.} Cap. 20.}. Therein lies that it is power that is peculiar to the state, and that it finds
application especially against outer enemies. And therein lies at the same time the difference between the state and
the cultural community. The state presupposes a people, a group of human beings that feels itself a unity in contrast to
other groups. But the cultural community can spread out over all such groups and connect them by harmonizing or common
ends, even if no power encompasses them all. ---

These three forms of social life we shall now consider more closely, with special regard to their ethical significance
and to the spirit and direction in which their further development ought to proceed, when we apply the standard set up
in the foregoing.

% --- p. 278 ---
\opage{278}\grouphead{A. THE FAMILY.}

\chapterhead{XIV.}{THE ETHICAL SIGNIFICANCE OF THE FAMILY.}

1. Nothing gives a better example of how nature can pave the way and form a foundation for what ethics demands or
approves than the circumstance that the most intimate and most perfect of all human societies is founded by some of the
strongest instincts in human nature. The kingdom of humanity, which is the highest ideal of ethics, not only has its
first germ and its constant source in the family relation, but is realized in family life, when this has reached its
highest form, in a way to which other forms of society can show no counterpart. The high points of all other forms of
society are measured by the degree to which they recall the intimacy and strength of the family relation. The kingdom
of humanity would have reached its completion if a universal brotherhood united all; and to designate an intimate
relation between master and servant, master and apprentice, authorities and people, the most apt expression we have is
the relation of parents to their children.

From several different sides the great ethical significance of the family relation can be demonstrated.

2. In other societies man takes part with only a part of his being, but in the family he can find nourishment for all
sides
% --- p. 279 ---
\opage{279}of his nature. Only there does he properly live as a \textit{whole human being.} The most primitive instincts
and the most ideal feelings find their satisfaction there. The community of life extends (or can extend) from the pure
natural instincts, through whose working the individual can often seem to be a mere means for the race's will to live,
through all material and spiritual domains. All interests can find their first cultivation in the family. The family
is a little world that sets all powers in activity. With this, again, it hangs together that the family more than any
other society unites (or can unite) the \textit{independence} of personalities with their intimate \textit{connection.}
Precisely because the family is not a special society but a general community of life, the greatest distinctiveness and
independence in its single members can stir freely and yet meet with understanding and sympathy. In other societies
deeper-going independence and distinctiveness are more or less compelled to hide themselves, or else to keep others at
a distance. In the family, because of the complete community of life, even apparently odd and paradoxical qualities can
meet with sympathy, because they are apprehended and understood in connection with the individual's whole nature.
Strangeness is abolished. Therefore there is rest and strength in no relation as in this one.

In the family relation life is lived not only with full consciousness, but a multitude of unconscious or half-conscious
influences continually make themselves felt. In the other social relations conscious observation and reflection and
conscious decision and action play a far greater part. Demands are made on strained attention and on self-conscious
conduct. But in family life the involuntary and the unconscious rule to a far higher degree. Through countless
influences, memories and moods the feeling for the home grows to such a height that, when it is attacked, it can appear
as an affect. Of the feelings tied to home and family it holds in a higher degree than of others that they are
nourished and increased by small accretions, which only by summing themselves up come to make themselves felt in
consciousness. The life of feeling thereby acquires a firmness which it
% --- p. 280 ---
\opage{280}lacks where it appears as a series of violently flaring but equally quickly vanishing
affects\footnote[1]{Cf. \textit{Psykologi} III, 7; VI E, 4---5}.

Finally, the family connects \textit{different generations} by the bond of nature. It throws the bridge between the
race's past and its future. It makes an understanding possible where between those standing further apart it would be
impossible. The conflict between old and new is subdued within the family by the deeper-lying sympathy which here has
the first word.

3. One might object against the family that it is too narrow a society, that it concentrates feeling and interest on a
small circle, so that what stands outside this becomes a matter of indifference. A family egoism can develop, of wider
extent, to be sure, than purely individual egoism, but no less than this a hindrance to the development of the general
love of mankind. To this the answer is that sympathy must first develop in narrow circles before it can spread in wide
ones. In the feelings that develop and find constant nourishment within the family we have the first and strongest
means of subduing and educating individual egoism. The general love of mankind is only an extension of a feeling that
arises within the family, an extension which does not always take place without resistance, but which nevertheless
always presupposes that the germ has been laid in the small. There is therefore no necessary contradiction between
family love and the general love of mankind. And to this is added (cf. XII, 1) that there is an inverse relation
between the strength and the extent of sympathy. If the extent is widened, it easily comes at the cost of the strength.
Very few people are (at any rate as yet) able to embrace distant and wide relations with such intimacy and force as what
lies near them. If, then, the strength of the feelings is not to suffer, there must be narrow circles in which they grow
strong. --- The family is here at once end and means. It affords the highest satisfaction for the needs of individuals,
and it is the home of powers that become of significance for the whole race. The
% --- p. 281 ---
\opage{281}more homes there are, the more hearths for the fire that keeps the life of the race going. ---

Another objection against the family is this, that by its limitation and by its life built on tradition and repetition
it easily leads to stagnation, monotony and dullness. Thereby what should be precisely one of the advantages of family
life is often prevented: that the individual could here unfold himself on all sides of his being. Many-sided unfolding
is hampered by the narrow content and by a peculiar shyness about appearing unreservedly in word and deed before those
nearest. This makes it intelligible that many individuals are valued less highly in their home than outside it ---
although the reverse would be most natural. --- These misgivings, however, stem only from the isolation of the home or
the family from the cultural community and the state. From these larger worlds come the fresh currents that can drive
out the stale air which, through the home's isolation, can gather within its walls\footnote[1]{Cf. my introductory
essay to the work ``Vort Hjem.''}.

The difficulties are wholly overcome only where the family has reached its highest form. Reality here, as everywhere,
falls far short of the ideal. But it offers approximations which it is our task to develop further. --- To throw closer
light on the questions that present themselves here, we shall first consider marriage, then the position and condition
of woman, and finally the relation between parents and children.

\grouphead{\textit{1. MARRIAGE.}}

\chapterhead{XV.}{SOCIOLOGICAL DATA.}

1. At the lowest stages we know, the union between the two sexes has the character of a relation of power and
compulsion, in which the
% --- p. 282 ---
\opage{282}weaker party is merely a means for the satisfaction of the other's lusts, or must slave for his needs. Some
recent writers (\textsc{Bachofen}, \textsc{Mc Lennan}, \textsc{Lubbock}) have thought that originally there reigned
complete freedom of mating between the two sexes (promiscuity, general hetaerism), so that every woman belonged to every
man who could for a moment take possession of her. This is surely, however, an exaggeration of the great freedom that
prevails in sexual relations among savage peoples round about the earth. There has scarcely been any time when complete
hetaerism reigned. Even where no custom or law prevented the constant change of sexual unions, individual preference as
well as the desire for property and power would still lead to the unions' becoming more than merely
momentary\footnote[1]{Cf. \emph{Spencer}: \textit{Principles of Sociology.} I. p. 662 ff.}. Woman, moreover, was early
regarded not merely as a means for the satisfaction of the sexual instinct; she also became the first slave, and
thereby acquired a value which was bound to bring it about that complete freedom could not be permitted. One cannot
prove that all other forms of marriage have developed out of an original hetaerism, even if one finds great
approximations to such a state round about the earth\footnote[2]{That the proofs of the existence of an original
hetaerism or promiscuity are untenable has been shown from various sides, notably by C. N. \emph{Starcke} \textit{(Die
primitive Familie.} Leipzig 1888) and \emph{Edward Westermarck} \textit{(The History of Human Marriage.} I. Helsingfors
1889. --- II. London 1891). (German translation 1902). --- Cf. my review of Starcke's book in ``Tilskueren'' 1888.}.
Since monogamy is found among animals too, there is in itself nothing to prevent its being found among human beings
also at purely primitive stages, and this is indeed the case. However attractive it might seem, one cannot demonstrate
a coherent course of development that leads from hetaerism as the lowest stage through polygamy --- partly polygyny,
partly polyandry, partly group marriages (where several men have several women in common) --- to monogamy. The
development has surely had a very different character among different tribes. The ordering of marriage everywhere hangs
together with so many other social relations that scarcely any single, clear and simple course of development can be
demonstrated.
% --- p. 283 ---
\opage{283}But one thing can be shown to have exercised an essential and necessary influence, namely the degree to which
the right of the individual personality has been felt and acknowledged. Before this right both the sensual impulse and
the impulse to dominate have had to give way. It is this that has brought it about that monogamy has, in the course of
development, been more and more acknowledged as the highest and the true form of union between man and woman.

Even if times or places could be shown where the horde or the tribe had women and children in common, as was the case
with other property, this community must nevertheless fall away as the individual began to feel his independence. The
need to have something for oneself was bound to express itself in this domain too. It is at first a purely selfish
need, which expresses itself, or at any rate wins acknowledgment, only in the case of the strong, and that is to say,
of the men. Polygyny therefore occurs more frequently and maintains itself longer than polyandry and group marriage,
which occur especially where circumstances are straitened. Monogamy too is often only an outcome of the need for
property of one's own, thus an outcome of power, except that this, for one reason or another that is not always easy to
point out, restricts its domain. The relation between man and woman in marriage here still has predominantly the
character of a relation of compulsion. It is perhaps even founded by violence, robbery or purchase. Purchase marks an
advance relative to violence and robbery; for by it it comes clearly to light that the woman has an economic value, and
she is spared like any costly ware. But even where purchase no longer takes place, the spouses (especially the woman)
are often given away without free choice; marriage is regarded as an affair, not between the two individuals, but
between the two kindreds which are joined through them, which are propagated through their children, and whose property
passes by inheritance to these. Marriage is here not a personal compact, but a ``league of
kindreds''\footnote[1]{Cf., for the North: R. \emph{Keyser}: \textit{Efterladte Skrifter.} II. Kristiania 1867. p. 306.
--- \emph{Michel de Montaigne} (Essais III, 5) expresses himself (at the end of the 16th cent.) thus on marriage: ``On
ne se marie pas pour soy, quoy qu'on die; on se marie autant, ou plus, pour sa posterité, pour sa famille; l'usage et
l'interest du mariage touche nostre race, bien loing pardelà nous: pourtant me playt cette façon qu' on le conduise
plustost par main tierce que par les propres.'' --- In circles other than the noble ones too this conception is found
right up to our own day.}. In the old marriage the personal
% --- p. 284 ---
\opage{284}relation between man and woman was not the main thing, but the founding of a new
household\footnote[1]{\emph{Leist}: \textit{Alt-arisches Jus civile.} I. Jena 1892. p. 154. 166.}. That adultery is
punished so severely in earlier times hangs together with the woman's being regarded as the property of the kindred or
of the husband\footnote[2]{A. H. \emph{Post}: \textit{Die Grundlage des Rechts und die Grundzüge seiner
Entwickelungsgeschichte.} Oldenburg 1884. p. 374 ff. --- \emph{Fustel de Coulanges}: \textit{La cité antique.} 4. éd.
p. 109.}. --- Only where the two individuals choose one another from free impulse and by free decision do we have the
highest form of marriage: \textit{free monogamy.}

2. In many people a certain dizziness arises when they first become aware of the great diversities in the ordering of
marital relations round about the earth and at the various times. Chance seems to rule, and the one arrangement seems
able to be as good as the other; at any rate, what has once been practiced seems thereby alone to acquire a certain
authority. If, for instance, hetaerism has been general, or at least approximations to hetaerism, might it not then be
conceived that we should return to it again? To this the answer is that the sample card which sociology shows us of the
various orderings of marriage is motley indeed; but that there nevertheless appears, as already indicated, a clear
connection between the development of free monogamy and the acknowledgment of the significance of the individual
personality. This connection will be demonstrated more closely in the following chapter. It comes into being originally
through displacements of motive and of value. In the being that the man attaches to himself for the sake of his lust he
discovers an economic value, and later he discovers that his female slave has a soul like himself, and that she must
stand for him not merely as a means, but as an end. --- Sociology can give ethics valuable pointers, but it is not
absolutely
% --- p. 285 ---
\opage{285}determinative for ethics. Even if one could demonstrate a course of development that began with hetaerism and
ended with free monogamy, there would still remain for ethics the task of demonstrating the value of the latter by
virtue of the principle of welfare. Historical development might very well have led in a pernicious direction, so that
ethics would have to demand a change, however difficult such a change might be to carry through. What is real is
therefore not on that account rational. % Hvad der er virkeligt, er derfor ikke fornuftigt: against Hegel's "Was wirklich ist, das ist vernünftig"

The real state of affairs, however, is this, that free monogamy is indeed officially acknowledged in the most civilized
countries as the only right form of marriage, but that the other forms have not thereby vanished from life. Hetaerism
spreads around us, even if it is forced to keep to the dark. And among a large part of the population it does not even
keep to the dark, but unfolds freely and involuntarily. The lowest forms of sexual union constantly appear anew, just as
the lowest animal forms can go on existing even after higher types have long since developed. What is designated as
``moral laxity''\footnote[1]{\emph{Hafstrøm}: \textit{Om Sædelighedsforholdene i det danske Folk.} Copenhagen 1888.
p. 67.} is very often old forms of life that have not vanished. Yet here nature and overcivilization meet in a
remarkable way. The constantly flourishing hetaerism is due not only to the impulses of fresh, unbridled nature that are
not yet softened and harmonized. It is due in part also to a one-sided training and refinement of some faculties of
mind at the expense of the whole character. Intellectual and aesthetic culture can be united with great looseness in
sexual respects. One has one's interests, perhaps serious ones, in other domains, and behaves essentially in an
enjoying and playful way in the sexual domain, without its ethical significance being made an object of reflection.
With one-sided intellectual and aesthetic culture there follow, moreover, not seldom a blasé and refined way of thinking
and a taste for the piquant, which finds nourishment especially in the sexual domain. Two other, opposite motives too
can meet: bitter want and overflowing vitality.
% --- p. 286 ---
\opage{286}Want leads to giving oneself up, and the craving for enjoyment leads to using those who are forced to sell
themselves as means to a moment's lust. --- From causes so different there can constantly proceed that isolation of the
sexual impulse from the other elements of the human personality which is the sting of the sexual problem. That this
impulse can appear in isolation and yet with such great power hangs together with the surplus of energy which nature
everywhere stakes on maintaining the existence of the species. The isolation becomes the more serious the more the life
of culture replaces the life of nature. The sexual problem is therefore closely interwoven with the whole problem of
culture and with the whole social problem, and can be properly understood only when this connection is not forgotten.

\chapterhead{XVI.}{FREE MONOGAMY.}

1. A community between human personalities can be perfect only when none of the participants stands as a mere means
for the others, and when no part of the single participant's being is one-sidedly drawn forward or pushed aside. All
polygamous forms of marriage conflict with this principle. In polygamy one man divides himself among several women, or
one woman among several men. Full devotion cannot then be present toward any one: the relation is always divided, and
the full union of the personalities does not come about. On the one side one parcels oneself out, on the other one makes
do with a fragment. ``In polygamy,'' says \textsc{Kant} \textit{(Rechtslehre} § 26), ``the person who gives herself away
wins only a part of the person to whom she wholly surrenders herself, and thereby makes herself a mere thing.'' % translated from Høffding's Danish
The many are made means for the one, or the one for the many. And this in turn entails that within the single
personality there comes about a separation
% --- p. 287 ---
\opage{287}of what should work together undivided. Where the feeling of love is more than animal heat, a physical and an
ideal element work inseparably together, and devotion embraces the whole personality. (Cf. XI, 10). But in polygamy the
purely physical element must necessarily be especially drawn forward. The fragment of oneself that one can give to many
is here only the purely physical element. As a purely physical relation, as a purely animal act, the sexual relation
also has only physical limits. The physical side of the sexual relation is the least individualized; the more it is
drawn forward, the more the difference between the objects vanishes; and conversely, the more different objects the
sexual instinct is directed toward, the more its physical side must become predominant. The sexual life of flies is a
readily accessible example. The purely physical union is momentary, and the halt in the mutual struggle for existence
that it brings with it often lasts only for the moment of mating itself. In some species of spider the nature of the
beast of prey awakens again immediately after mating, so that male and female now see only a possible prey in one
another. This shows how isolated and momentary the sexual instinct can appear. In the feeling of love proper, on the
other hand, there expresses itself, besides the elementary instinct, the image of the other individual in its
distinctiveness and an intimate joy in it. The instinct then works not as a purely blind power, but opens the eye to
the other individual's nature and thereby paves the way --- through a displacement of motives and values --- for a
devotion that is more than mere physical union. The ethical significance of the sexual instinct will at higher stages
be this, that personal peculiarities are discovered and valued which would otherwise remain unnoticed. An attention is
awakened which works with the power of a force of nature. The intimacy of the union then rests not merely on the race's
urge for self-preservation binding one individual to an individual of the other sex, but also on qualities and
expressions of life that can appear only in a lasting life together exercising their influence. Not only love, but
faithfulness too is an organ for making spiritual discoveries.

Monogamy alone, therefore, can be the form of sexual union between independent personalities who enter into the
% --- p. 288 ---
\opage{288}relation each with his whole nature, not in order to satisfy a single, isolated need.

2. It belongs to personal devotion that it embraces the personality not only in its whole extent, but also the whole of
life, as long as this lasts for both parties. A feeling that does not involuntarily believe in its own continuance is
not genuine. It would be unnatural if, when the feeling toward another individual is in its strongest movement, a
reservation were nevertheless made and there were an express consciousness that the whole was to hold only for a
certain time. If the feeling is genuine, such a possibility will not be able to come before thought. Because the
feeling believes, and must believe, in its own continuance, it need not, to be sure, continue. The superlative lies in
the nature of feeling; when it is strong, it excludes limitation and comparison. But even if what was meant to have
eternal existence proved to be transient, the feeling would still not be able to fill the time it lasted in the right
way if it were not bound up with the belief in its continuance. Illusions are frequent; but a relation that is founded
with express consciousness of, or even calculation on, one's always being able to withdraw from it --- such a relation
is not only as a rule a deception of the other party, but can hardly become a whole and full relation.

A personal being cannot be \textit{parceled out} into detached and changing moments. Personality is present only where an
inner unity and coherence runs through the single moments, and where there are certain definite feelings that make up
the firm kernel. But neither can a personal being be \textit{exhausted} in a single moment. It encloses within itself
such riches that, where real sympathy is present, there can be enough to discover through a whole life. It is a great
illusion to think that a series of changing sexual unions should furnish rich material for knowledge of men and for
human development\footnote[1]{In that case a Roman woman mentioned by Jerome must have carried knowledge of men to a high
degree: she was married to her twenty-third husband and was his twenty-first wife.}.
% --- p. 289 ---
\opage{289}Momentary unions do not lead into the innermost sanctuary; this opens itself only to constant, faithful
sympathy.

What is called, with an often misused word, the doctrine of ``free love'' is an attempt to put inconstancy into a system
and to proclaim it an essential moment in the feeling of love. \textsc{Enfantin}, St. Simon's well-known disciple, held
that marriage as it now exists suited only constant and jealous natures, not the inconstant and capricious, --- the
Othellos, not the Don Juans. But since all kinds of natures have a claim to satisfaction --- including ``les personnes
vives, coquettes, séduisantes, attrayantes, changeantes'' --- he demands divorce as an institution, as an acknowledgment
of the right to inconstancy, not merely as a way out in cases of conflict\footnote[1]{Cf. \emph{Paul Janet}: \textit{St.
Simon et le St. Simonisme.} p. 138---142.}. A similar line of thought is found in the anonymous author of ``Elements of
Social Science.'' When this latter author nevertheless will not abolish the concept of self-control, and when he
maintains that the sexual instinct must be satisfied in such a way that it does not become harmful to others, he
demands a subordination of momentary impulses to more comprehensive considerations, which conflicts with the line of
thought that leads to the doctrine of free love. For passion is absolutist, and its sole dominion is abolished when
other feelings, e.g. sympathy, assert themselves beside it. There is scarcely any question that if a lively sympathy
with other beings is awakened before the sexual urge can announce itself in its full force, this urge will undergo a
considerable metamorphosis. It is on isolation that an essential part of its strength rests. --- For the rest, so as
not to do this latter author an injustice, one must note that by marriage he means indissoluble marriage. Where there is
reasonably easy access to divorce, as in Germany, there --- he thinks --- marriage is in fact
abolished\footnote[2]{\textit{Elements of Social Science.} By a Doctor of Medicine. 13th ed. London 1875. p. 371.}. One
sees how easily the whole thing becomes a quarrel about words. Many of the attacks which, especially by aesthetes, are
directed against marriage in
% --- p. 290 ---
\opage{290}general are surely due only to reminiscences of French and English writers who polemicize against
indissoluble marriage, or marriage dissoluble only on quite a few grounds.

The inconstancy that is sometimes reckoned an essential moment in the relation between man and woman cannot avoid
leading to pain and unhappiness, as long as it does not lie in the nature of all, and --- it must be added --- as long
as the urge for change does not set in in both parties (by a remarkable harmonia praestabilita) at the same time. A
Danish author\footnote[1]{The anonymous author of \textit{En Livsanskuelse, grundet paa Elskov} (Copenhagen 1881) and
\textit{Forholdet mellem Mand og Kvinde belyst gennem Udviklingshypotesen} (Copenhagen 1884). Cf. also his essay in
``Tilskueren'' (1885): \textit{Om en Reaktion mod den moderne Stræben efter større sexuel Sædelighed.}}, who has
discussed the relation between man and woman from a purely psychological and ethical standpoint with great ability and
seriousness, has shown that as long as there are, besides the polygamous natures for whom inconstancy and change belong
to the essence of the feeling of love, also monogamous natures who cannot conceive of true love without constancy and
faithfulness, so long the urge for change on which the champions of free love lay so much weight cannot be satisfied
without bringing pain and sorrow to others. It will often go as with a couple in \textsc{Daudet}'s ``Sappho.'' The
relation was founded with the mental reservation that it was to be temporary; but for the young woman the relation has
become serious, and she takes her own life when the separation is to come: On meurt donc quelquefois de ces ruptures!
--- If it is reckless to play with fire, it is a thousand times more reckless to play with a human being's welfare.

3. The individual personality is as a rule not fully finished in the period of life in which the feeling of love plays
the greatest part. When this feeling has now led to marital union, what matters is whether the continued development of
the two personalities can proceed in mutual harmony. This is of all the greater significance since it is not only the
character of the two individuals as a whole, but quite especially the feeling that binds them, that can and will in
the course of time
% --- p. 291 ---
\opage{291}undergo changes. Where development proceeds happily, the feeling of love will turn from a flaring affect into
an intimate feeling that --- when it counts --- can have the same strength, if not the same momentarily overpowering
violence, as at its first outbreak.

The development that the individuals, and especially the feeling that binds them, can undergo in their continued life
together is not independent of their will. It is one of the current psychological misunderstandings that the will is
something wholly different from thought and feeling, a power that must first be specially appealed to in order to step
forth like a deus ex machina and resolve the difficulties. Under sound and natural conditions the will develops and
fixes what thought has grasped and feeling has attached itself to. The will must be there from first to last; only so
can it also help in the single, fateful moment. It is not a blind fate to which the two individuals are subject. The
matter is to a great extent placed in their own hands; it depends on the seriousness with which they conduct their
life. Marriage, like every life in common, requires self-control and effort in order to last. If every ill humor, every
contrast of character, is allowed to become decisive for the future of the relation, this future will be short. If no
other satisfaction is known than that which the flaring affect and the influence of the new and surprising give,
marital love proper will not be able to develop. There is a form of aesthetic view of life for which the highest thing
is the changing play of feelings, and which therefore always seeks change and novelty. (Cf. III, 4). Such a view of life
cannot enter into marriage; but consistently it cannot enter into any fixed relation whatever either. Firmness comes
only when the will is put to work, and when another form of the life of feeling is known than the affects of the
moment. What matters is that the erotic feeling be taken up as a member of the whole personal life and not go on
standing as an accidental and detached element. And for such an intimate union with life's other elements it has
perhaps a greater possibility than other feelings. There was something apt in it when one of the later Stoics (among
% --- p. 292 ---
\opage{292}whom, as has been remarked\footnote[1]{\emph{Bonhöffer}: \textit{Die Ethik des Stoikers Epiktet.} Stuttgart 1894.
p. 88.}, one finds the most ideal conception of marriage that antiquity knows) describes the marital union as a chemical
union, a ``complete mixture,'' in contrast to the more mechanical union in other relations of kinship and friendship. He
who wants to preserve the sexual element in its isolation, in order thereby to secure himself a source of enjoyment of
its own, will fear the metamorphosis which that element undergoes in true marriage. And he who does wish for such a
metamorphosis, but fears that it will not succeed, will also have misgivings about marriage. But it can bring with it a
wholly different conception of the relation between the two sexes whether one knows it only without such a metamorphosis
or after one. An experienced observer of convicts and their circumstances\footnote[2]{\emph{Hafstrøm}: \textit{Om
Sædelighedsforholdene i det danske Folk.} p. 82.} remarks: ``Although convicts as a rule come from the poorest strata
of society, their correspondence regularly shows how faithful and enduring most wives are while their husbands are
serving sentences for crimes\dots{} Most of the men too, however deeply sunk they may otherwise be, cling with love to
their wives and children, and it is quite strange to see how the married prisoners as a rule have an entirely different
view of woman from the unmarried. In marriage they quickly learn that woman has quite another significance and value
than merely being an instrument for the man's sensuality.''

It would be unpsychological to assert that the will can do everything. In the relation of life here in question very
many elements work together, and not all of them can be mastered even by the most earnest will. Without real guilt on
either side the development can take an unfortunate direction. Disharmony can arise in the characters and in the views
of life; the relation between the physical and the ideal element in the feeling of love itself can develop in a
different, even conflicting, way in the two individuals; toward the new circumstances and questions that life brings
with it they can find themselves in quite opposite moods. Indeed, when one rightly considers how many
% --- p. 293 ---
\opage{293}outer and inner conditions work together here, it must be regarded as a great piece of good fortune when the
original feeling stands its test and undergoes its metamorphosis in the course of the individuals' continued development
without running aground on any of the many reefs; the Stoic Antipater was to that extent right to call entering into
marriage a heroic deed (ἡρωϊκόν). That is why marriage contains the possibility of so many dramas and leads to so many
tragedies. The involuntary and the voluntary, fate and guilt, can here less than at any other point be kept sharply
apart. It would be doctrinaire to assert that if only the feeling was serious from the first, it would also last. It may
have been serious from the first, but the later development may have withdrawn from it a nourishment without which it
cannot subsist.

4. Of essential support in the metamorphosis that the feeling of love undergoes in the continued life together are the
common tasks set the two individuals and the common activity they carry on. New values arise which were not
anticipated beforehand. --- They share the work for their material livelihood. They have to struggle with the same outer
dispensations. They can work in common on their further spiritual development. When they keep together in their
seeking, they will more easily keep together in their finding. --- But the greatest significance belongs to common care
for the children. Common anxiety and common sacrifice let them come still closer to one another than in the light and
glad time of the new feeling. The sympathy between them is deepened and heightened. The responsibility for preserving
the relation is felt all the more strongly when not merely the fate of a single individual, but the future of several
helpless beings depends on what seriousness and self-control are shown\footnote[1]{Charles \emph{Darwin} tells of his
father, who was a physician: ``Owing to his power of winning confidence he received many strange confessions of misery
and guilt. He often remarked how many miserable wives he had known. In several instances husbands and wives had gone on
fairly well together for between twenty and thirty years, and then hated each other bitterly; he attributed this to
their having lost a common bond in their children having grown up.'' % Life and Letters, ch. 1 (Darwin's English as remembered; Høffding gives it in Danish)
(Life and Letters. I. p. 13). --- Other common bonds than this, then, are necessary.}.
% --- p. 294 ---
\opage{294}Here is at the same time a point where the advantage of monogamy over polygamy clearly shows itself. Only in
monogamous marriage can the child receive the whole and full love of the parents. Only there can full harmony and unity
reign in the home. And even where the feeling of love has not been the first motive for entering into marriage, but
where ``considerations of reason,'' especially economic considerations or the wish to leave offspring, have been
decisive, the most intimate union may arise through a displacement of motive. The common tasks, and the mutual
understanding of the personalities that life together can bring, can make possible a sympathy that is often every bit
as intimate as that which arises from the feeling of love and its metamorphoses. --- Where the wish for offspring has
not been the motive that led to the marriage, the offspring stand as a new value which the relation entered into
brings forth. --- The course of development of marriage in individual cases can thus offer examples of all three kinds
of metamorphosis mentioned above (XIII, 4).

5. The characterization of marriage given in the foregoing excludes not only polygamy but also unfree monogamy, in which
the woman does not stand fully as the man's equal in marriage, partly because she does not choose in freedom but is
given away by her family, partly because within marriage she does not have the same rights as the man, and partly
because she is not his peer in spiritual respects. The last-named point contains the reason for the first two. When the
woman essentially leads an obscure life of feeling and instinct, and when her calling is thought to be exhausted by
housekeeping and the care of children, she can be only an appendage to the man, fill out a single side of his life. To
be more she lacks the independence in the life of thought and will which is the condition for being able to choose and
to assert herself in accordance with her choice. Marriage then becomes on the most essential points a relation between
an active and a passive party, and does not attain the perfection it has where it is a relation between two parties
who are each active in their own way and can follow one another's activity with understanding sympathy. The division of
labor between them can then conform to their individuality, and superiority will not be as good as
% --- p. 295 ---
\opage{295}self-evident on the one side. Who is to be first must not be due to a legalized fiat of power. But this can
be attained only through a more many-sided development of the feminine gifts than has until the most recent times been
thought possible and justifiable. Only then will monogamy reach its completion. The work that nature has once assigned
to woman she will naturally not push away from herself; a \textit{real} natural vocation does not let itself be denied.
--- A closer consideration of this question will be undertaken in a later section.

The different position that man and woman occupy in marriage according to the hitherto current way of thinking comes to
light especially in the different judgment passed on their unfaithfulness. That the wife's unfaithfulness is judged
most severely might seem to find its justification in the fact that it entails far more drastic consequences for the
family than the man's, whose adventures need not entail any consequences felt in the family at all. But this loses its
significance in the face of the demand for perfectly equal rights that free monogamy makes. Where this demand is not
satisfied, there is in reality a polygamy, and the different judgment of male and female unfaithfulness also stems from
a stage of development at which woman was regarded as the man's property. (See XV, 1). The different judgment has
persisted after the official acknowledgment of free monogamy has taken place.

An Italian writer, Mantegazza\footnote[1]{\textit{Die Physiologie der Liebe.} (Translated from the Italian). Jena 1877.
p. 351 ff.}, fights stoutly for the unequal judgment. His line of thought is built chiefly on two grounds. First, on
this, that woman's whole ethical task lies in the family relation. While --- he says --- a hundred virtues are required
of the man, who must work and struggle in such manifold relations, faithfulness is the only [!] virtue demanded of the
woman: ``is that perhaps too much?'' --- But secondly: man is by nature inclined to polygamy, he is more faithless,
brutal, capricious and sensual than woman, who is not so easily overcome by an intoxication of the senses. Mantegazza
paints this difference in such strong colors that it becomes incomprehensible
% --- p. 296 ---
\opage{296}how so pure a being can ever fall. In any case it would be reasonable to turn the severest judgment not
against the fallen angel but against the devil who causes her fall, especially if those are right who think that
woman's fall is far more often than man's due to a higher motive than mere sensuality. --- That the difference is present
to a certain degree can hardly be denied. But here it is screwed up to an unnatural height, at the same time as woman's
whole task as a human being is restricted to the family relation. And even within this there is surely occasion enough
for woman to display more virtues than that one, exceedingly important as it is.

Both the effect that laxity in marital relations has on family life as a whole, and the way in which unfaithfulness
intervenes in the individual's personal life, will naturally differ according to psychological and social
circumstances. Thus \textsc{Burckhardt} remarks \textit{(Kultur der Renaissance.} 4. Aufl. II. p. 183 ff.): ``What stands
out as peculiar to Renaissance Italy is that marriage and its rights are trodden underfoot to a higher degree, and at
any rate with greater consciousness, than elsewhere, and that the principle is openly pronounced that marriages are to
be concluded only for a time, namely as long as the wife pleases the husband. The young girls of the higher classes do
not come into consideration; all passion has married women for its object. But it is remarkable that marriages are
nevertheless not seen to have declined, and that family life by no means suffered the disturbance it would suffer under
similar conditions in the North.

The more highly cultivated, individually developed woman disposes of herself with
quite another sovereignty than in the North, and unfaithfulness does not make any terrible break in her life, if she
can secure herself against the outward consequences. The husband's right to her faithfulness has not the firm foundation
which it acquires among the Northerners through the poetry and the passion that courtship and betrothal bring with
them: after the most fleeting acquaintance and directly from the guardianship of her parents or the convent the young
woman steps out into the world, and only now does her individuality develop, and that with marvelous rapidity. For this
reason above all the husband's right is very conditional.'' % translated from Høffding's Danish
--- In the southern
% --- p. 297 ---
\opage{297}countries conditions have in this respect not wholly changed since the days of the Renaissance.

6. Free monogamy as it has been described in the foregoing is indeed an ideal, but not one that floats in the air. Every
happy and serious marriage shows us an approximation to it. It fulfills in the best way the ends of the family, not only
to afford full satisfaction to its members, but also to be the hearth of sympathy and the best school for the new
generation. Its maintenance and development therefore becomes a significant task.

But although free monogamy, insofar as it really exists, stands as one of the most important fruits of historical
development, we nevertheless find beside it everywhere the lower, indeed the very lowest, forms of sexual union in full
bloom. Hetaerism, momentary unions and loose unions have not vanished because free monogamy has won acknowledgment.

It would be an unhistorical view to take free monogamy not only for the \textit{ideal} form of marriage, but also for
the \textit{original} form, and to regard hetaerism and loose unions as having arisen through decline from it and denial
of it. One would thereby not only come into conflict with sociology, but also be at a loss to explain the possibility of
such a fall; and one would stand hopeless before the real circumstances. If, on the other hand, free monogamy has
developed historically, then we have only to go further in the same direction, to maintain and develop what has been
begun. Hetaerism and loose unions are then signs that human nature is not yet adapted to higher social conditions of
life, --- that it still has too much of the animal in it to be able to realize the properly human form of sexual life.
It will always show itself here that the two endeavors, humanization and emancipation (XIII, 3), belong closely
together. The more brutal and low the character of sexual life, the further behind it stands also as regards the
liberation of personalities. Only spiritually immature individuals let themselves be used as mere means for the lust of
others.

There is no lack of sufferings and misfortunes on account of this unfinished development. A desperate struggle is often
waged here
% --- p. 298 ---
\opage{298}between the different tendencies in human nature. And the struggle will be long. In the ethical domain
development moves forward slowly, not least where it is a matter of subordinating one of the mightiest impulses of
nature to ethical-social laws that correspond to a higher stage of life than the one at which man parted from the
animal. But the worst solution of all would be to bargain down the ideal. That would only be the beginning of a
regression. --- Legal prohibitions, compulsion and pharisaic moralizing will accomplish little. A sound and harmonious
education, a lively sense for the right of personal life in all, a warm love of mankind will give the best support in
the struggle. What matters above all is to prevent sexual life from being isolated from the rest of the life of the
soul. It is at once to exercise its attracting and inspiring power and to have its proper limitation as a single one of
life's forms. Everyone who in his conduct of life strives in this direction is working at an essential part of the
welfare of the human race. ---

The persistence of hetaerism, however, is, as already remarked, to be explained not only by the instinct's having still
kept its unbridled character, but also by material want and bodily and spiritual degradation preventing a purely human
life. If it is true that in London every seventh and in Hamburg every ninth (younger) woman is a
prostitute\footnote[1]{\emph{Oettingen}: \textit{Moralstatistik.} 3. Aufl. p. 197. --- In Copenhagen a little over 1
per cent of women between 20 and 30 years of age are prostitutes. M. \emph{Rubin} in Nationaløkonomisk Tidsskrift 1887.
p. 40. --- In Paris there are about 5000 registered prostitutes, but the number in clandestine prostitution is given as
50,000. E. \emph{Pontoppidan}: \textit{Kønssygdommene.} p. 71.}, then it is not the natural instinct that must bear
the whole blame, but it is the effects of material and social conditions that we have before us here. We stand here
before a branch of the great social question, that abyss into which we look down from so many sides. It is above all
bitter want that drives so many to throw themselves away\footnote[2]{\emph{Prosper Despine}: \textit{Psychologie
naturelle.} Paris 1868. III. p. 215. (Often it is not their own want, but the want of parents or children, that makes
women seek a living as prostitutes). Cf. \emph{Edv. Ehlers}: \textit{Bidrag til Diskussionen om
Prostitutionsspørgsmaalet.} Copenhagen 1896. p. 9: ``He who believes that prostitution can be eradicated must also be
able to believe in the eradication of poverty.''}.
% --- p. 299 ---
\opage{299}A better and more independent development of the feminine faculties will both give woman greater power of
resistance by strengthening her self-esteem, and procure her better prospects of making her way in the world. --- In the
man's case too social and economic conditions intervene here, since the age at which his income permits him to enter
into marriage is, in the middle class, becoming later and later\footnote[1]{\emph{Rubin} and \emph{Westergaard}:
\textit{Ægteskabsstatistik paa Grundlag af den sociale Lagdeling.} Copenhagen 1890. p. 47 f., 51---53.}.

\chapterhead{XVII.}{ON ENTERING INTO AND DISSOLVING MARRIAGE.}

1. Everything that divides human beings into different camps can become a hindrance to marriage between individuals who
would otherwise feel inclination for one another. But even in such cases the real reason lies in the differences of
character that condition, or are conditioned by, the conflicting tendencies. For experience shows that neither
religious nor national nor political oppositions are absolute hindrances to happy marital union. What matters is not so
much agreement in views as agreement in characters, so that these either resemble or supplement one another. From a
strictly theological point of view ``mixed'' marriages are naturally to be condemned. Old teachers of the Church
regarded them as fornication; Catholicism forbids them; modern Protestant theologians content themselves with
designating them as ``monstrous.'' Nor can it be denied that the possibility of marital union between individuals of
different faith points to a weakening of
% --- p. 300 ---
\opage{300}confessional passion. But humane ethics sees no misfortune in such a weakening, provided only that spiritual
lukewarmness and blasé indifference do not follow with it. It regards such unions, where they succeed, as the
expression of nature's victory over unnatural barriers. It assumes that behind different confessions of faith there may
very well lie a deep spiritual kinship, and it finds in the power of the feeling of love to let this deeper-lying
spiritual kinship come into its own one of the most significant sides of this mighty force. Two striving natures will
understand one another\footnote[1]{Schleiermacher's wife, who stood in so intimate a relation to him, nevertheless
needed a more orthodox Christianity than the one her husband preached. This gave occasion to a beautiful exchange of
thoughts between them. \textit{(Aus Schleiermachers Leben.} In Briefen. II. p. 434---436). On the other hand it is
related of the wife of the ultramontane Romantic Görres that she never went to church, and yet here too the marriage was
extraordinarily happy. (\emph{Sepp}: \textit{Görres und seine Zeitgenossen.} p. 68). Cf. with this Bjørnson's
wonderfully moving portrayal of the relation between the pastor Sang and his wife \textit{(Over Evne.} I).}. But on the
other hand, as was said, different faith can hang together with differences of character, or the differences of faith
can be so dogmatically pointed that a mixed marriage becomes a risky venture\footnote[2]{According to experience from
Switzerland, mixed marriages are more exposed to dissolution than marriages in which both parties are Catholics or
Protestants. (H. \emph{Westergaard} in Nationaløkonomisk Tidsskrift. 1887. p. 30). --- One of the things that brought
unhappiness into the relation between Carlyle and his wife is said to have been that the highly gifted woman agreed with
her husband in his criticism of the positive religions, but not in the faith he himself held to. (\emph{Froude}:
\textit{My relations with Carlyle.} p. 7 f.).}. Marriage attains its most intimate form only where approximation and
union are possible in the whole way of conceiving and conducting life.

2. Entering into marriage is naturally ethically justified only where there is a prospect of being able to support a
family. But this ethical condition cannot be maintained by way of legal compulsion, as was attempted in earlier times in
some German states. In the first place it is questionable and intolerable to let it depend on the judgment of a public
authority whether two individuals who
% --- p. 301 ---
\opage{301}wish to marry have a prospect of getting by in economic respects. The individual can best decide this
himself, and it is beneficial that he bear the whole responsibility himself. There opens up, moreover, the possibility
of all kinds of abuse of authority. And it has proved that by such measures one only succeeded in making bad worse, the
number of loose unions and of illegitimate children increasing\footnote[1]{Cf. E. \emph{Löning}: \textit{Armenpflege und
Armenpolizei.} (Schönberg's Handbuch der politischen Oekonomie. 1st ed. II), p. 581. 595. --- \emph{Stuart Mill}
remarkably sympathizes with such prohibitions and finds nothing in them that encroaches on the freedom of the
individual, although he hints that they may perhaps, owing to local conditions, not always prove practical. \textit{(On
Liberty.} Danish transl. p. 196).}.

Not only for entering into marriage, but also for the children who are brought into the world in it, the individuals
have an ethical responsibility. Inner and outer circumstances can make it impermissible to leave offspring. So when one
oneself suffers from hereditary diseases (such as insanity, leprosy, etc.). Or when the outer circumstances are so
hopeless that the children are condemned to suffering and ruin. The great mortality among newborn children, which has
been called the modern worship of Moloch\footnote[2]{\emph{Rümelin}: \textit{Reden und Aufsätze.} Freiburg u. Tübingen.
1875. p. 331 f. --- \emph{Rubin} and \emph{Westergaard} \textit{(Befolkningsstatistik} etc. p. 100 ff.) have shown that
the more children are born in a marriage, or the more quickly the children's births follow one another, the greater is
mortality in proportion.}, can be prevented only when no more children are brought into the world than can find care and
nourishment under the given conditions. Children do not, after all, come into being of themselves without the parents'
knowledge and will. How the individual can fulfill the responsibility that rests on him here is a question that belongs
rather to medicine's decision than to that of ethics.

3. A deep-lying feeling, whose first source and motive it is not easy to point out, now fills most people with
loathing and horror at the thought of a sexual union between parents and children or between brothers and sisters. Among
savage peoples, on the other hand, such unions are frequent\footnote[3]{\emph{Spencer}: \textit{Principles of
Sociology.} I. p. 636 f. 680.};
% --- p. 302 ---
\opage{302}yet different customs appear early among different tribes\footnote[1]{\emph{Snorre Sturleson}: \textit{Ynglingesaga.} Ch. 4: ``When Njord was among the Vanir, he had his sister to wife; for so the law there
permitted\dots{}. But among the Æsir it was forbidden to have such near kin to wife.''}. Severe punishments for marital
union with mother, daughter and daughter-in-law are laid down by the Babylonian lawgiver Hammurabi (2200 years B.C.);
likewise in the Law of Moses, where the relation to the sister is also forbidden\footnote[2]{Cf. Fr. \emph{Buhl}:
\textit{Kong Hammurabis Lovsamling.} Nordisk Tidsskrift 1903. p. 468.}.

Sociologically that feeling can perhaps be connected with the custom prevailing among so many peoples that the wife is
to be fetched from another tribe than the man's own; but the rise of this custom itself (exogamy) is not yet properly
explained. In Indian and in medieval legislation, as is well known, there are strict prohibitions of marriage even with
more distant relatives, down to the sixth or seventh degree. According to the custom prevailing in Protestant countries
only marriage between the very nearest relatives is impermissible. Whatever the sociological explanation of this may
be\footnote[3]{It is scarcely likely that experiences of the weakening of offspring through marriage between near
relatives have been decisive. The necessity of such a weakening has not so far been demonstrated. --- If morbid
dispositions are once present in a family, they will naturally be propagated and increased by repeated unions between
members of the family. --- An attempt at a sociological explanation is given by \emph{Starcke} \textit{(Die primitive
Familie.} p. 238---247). The old prohibitions are derived from the wish to avoid the confusion in the legal relations
resting on kinship that would arise from marriage between relatives. This does not, however, explain the constant
after-effect and persistence of the feeling which is supposed originally to have been due to purely economic
circumstances.}, the ethical justification will be found in the following consideration.

The relation between brothers and sisters and between parents and children would lose its free and secure character if
the possibility of a sexual relation, and of the passions bound up with it, were present. The complete trust on which
that relation must rest if it is to have its full value would be lacking if the erotic impulse forced its way in here.
In the close and constant
% --- p. 303 ---
\opage{303}life together there would frequently be occasion for passionate relations that could have a harrowing and
dissolving effect\footnote[1]{Cf. already \emph{Maimonides}, a Jewish philosopher of the 12th cent. (\emph{Stäudlin}:
\textit{Geschichte der Vorstellungen und Lehren von der Ehe.} Göttingen 1826. p. 457). Cf., among later writers,
\emph{Grotius}: \textit{De jure belli et pacis.} II, 12---13. \emph{Hume}: \textit{Enquiry concerning the principles of
Morals.} Sect. IV.}. A psychical impossibility will here be the best safeguard. And only through such an impossibility
does the distinctive, rich and intimate relation between brother and sister come into being; a new value is won for
life. The relation between father and daughter, mother and son, too, is first shaped in its distinctiveness when it
appears in decided difference from the sexual relation. --- And on the other side there is something natural in the
erotic affect's not arising amid the familiarity and intimacy of home, the feeling of original belonging together that
reigns in the relation between parents and children and between brothers and sisters. In the erotic affect there
expresses itself a need for supplementation which is best satisfied by union with an individual of another family. An
affect arises from what is hitherto unknown, from what brings something hitherto unexperienced. That is why the erotic
affect is excluded under ordinary circumstances between brother and sister, who do not feel themselves to be so
different, and whose relation does not first have to be initiated and founded by the awakening of a special impulse.
It is characteristic of the feeling of love that a strange personality is discovered; therefore the love relation is
founded by declaration and choice. The relation between brothers and sisters and that between parents and children, on
the other hand, are not due to discovery and choice, but are given as a natural relation.

4. Entering into marriage naturally entails new duties and rights for the parties in their relation to one another and
to others. The community founded embraces all sides of life, and it is therefore of importance that it be established
on record. The formation of a new social cell interests more than the two individuals themselves. Not only their
families, but also the state, whose task it may become to maintain rights founded by the marriage, is interested in the
union's being officially declared. Nor does this violate the freedom of marriage. It gives it
% --- p. 304 ---
\opage{304}only an outer frame. Where personal feeling has spoken and tied the union, there it will be only a matter of
drawing all the consequences when the outer social and legal side is put in order, and no sense of freedom, however
romantic, can be violated by it. The opposition to a public declaration of entry into a marital relation stems chiefly
from its seeming to lie in the usual forms --- both in civil and in church ``wedding'' --- that the union should have
its sanction from outside and not be hallowed by the very intimate feeling that has founded it.

The ethical validity of marriage does not rest on the ``sanction'' of it by society and the state. Such a view can be
maintained only from the standpoint of the principle of authority. Marriage receives its ethical validity through the
free declaration between the two persons themselves, and the task of the state can only be to lay down and maintain the
rights and duties they thereby take upon themselves. It is a confusion of shell and kernel to regard the official
establishment of marriage as its ethical founding. This holds whether the official establishment takes place in the
church's way or in the civil way. There can be a marriage in spirit and in truth without these outer forms, and it would
be barbarism if --- as has been proposed --- one were to force such forms on anyone by the power of the police. The need
to establish the entering into the relation was originally bound up with the social significance of marriage as the
founding of a new household. With the development of civil law the public declaration in oral or written form became the
main thing, and the bride's transfer to the bridegroom's house, which in older times was the sign of the new household's
founding and the main point of the wedding, gradually became a mere festive custom\footnote[1]{\emph{Leist}:
\textit{Alt-Arisches Jus Civile.} II. Jena 1896. p. 120 ff.}. The civil-law act that took the place of the
involuntarily developed form for establishing the entry into marriage did not at once receive a purely civil form, but
was surrounded with religious symbols. The Christian Church, however, long retained the memory of the more primitive
form of entering into marriage, inasmuch as for a long time it did not
% --- p. 305 ---
\opage{305}regard the church ceremony as a necessary condition of valid marriage. The marriage of Adam and Eve, the model
of all other marriages, was after all concluded without witnesses! According to Catholic doctrine marriage was a
sacrament administered by man and woman, even when they enter into their union without priest and without witnesses;
and it was therefore an inconsistency when the Council of Trent laid down the necessity of a church
wedding\footnote[1]{This inconsistency is laid bare in a very clear way by \emph{Paola Sarpi} % PRINTED AS IS: Paola for Paolo (read at the image)
in his \textit{History of
the Council of Trent.} (French transl. by Amelot de la Houssaie. Amsterdam 1686. p. 765). Church consecration of
marriage early became custom in the ancient Church; it was laid down by law first by the East Roman emperors for the
Greek Church and by the Council of Trent for the Roman Church. --- That secret marriages were regarded as valid in the
first Protestant church too is seen from \emph{Peter Plade}'s \textit{Visitatsbog} % Peder Palladius
 (ed. by S. Grundtvig, p. 82 f.),
where the public establishment in church is recommended merely by the importance, in legal respects, of there being
witnesses to the entry into marriage. --- Cf. on this and related questions: J. \emph{Nellemann}:
\textit{Retshistoriske Bemærkninger om kirkelig Vielse.} (Historisk Tidsskrift. Fifth Series. First Volume).}.

The provisions that the laws of the state contain concerning marriage concern not only the conditions under which legal
effects can flow from it, but also aim at laying down the relation between the spouses themselves. They still bear in a
high degree the stamp of the old notion that the man is the woman's master, and that she is subject to his correction.

It lies in the nature of free monogamy that the wife's legal position in marriage must be coordinate with the husband's.
It ought not to be possible to dispose of her inheritance or acquired property without her consent. And she ought to
have a right of co-determination in all important decisions concerning common interests. In fact she does as a rule
exercise a great influence and bears an ethical responsibility for the decisions that are taken; but this responsibility
will be sharpened when she has a legal right of co-determination. Where she cannot approve the husband's dispositions,
she ought at least to be able to secure her children and herself against threatening ruin\footnote[2]{How much woman can
accomplish even without legal right is seen from the way in which the workers' wives in Rochdale made use of the famous
cooperative store there. ``Many married women become members because their husbands will not take the trouble, and
others in self-defence, to prevent the husbands from spending the money in drink. The husband cannot withdraw the
savings standing in the wife's name unless she signs the order. Of course by the existing law the husband could, by
legal process, possess himself of the money; but a legal process takes time, and before it can be set in motion the
husband becomes sober and thinks better of it.'' % after Mill's English, as far as Høffding's Danish allows
(\emph{Stuart Mill}: \textit{Principles of Political Economy.} IV, 7, 6).}. The
% --- p. 306 ---
\opage{306}more spiritually developed she becomes through an improved manner of education, the greater the independence
and practical experience she wins through richer opportunity to come to know the life around her, --- the better will
such an arrangement be able to subsist without hindering reasonable and justified undertakings. --- When the wife's
rights are laid down at the entry into marriage as something self-evident, there will be nothing in this that can
violate the intimate relation between the spouses, and yet a safeguard will have been provided should disagreement and
ill feeling arise. The thoughts expressed here are often dismissed with the remark that marriage is to be an intimate
relation of trust, not a legal relation; unity is to be the decisive thing, not difference. But it is a legal relation
all the same when all authority is assigned to the one party. To fix this authority on the husband is, after all, to
stress the difference between the two parties to the marriage far more strongly than would happen if equal authority
were allotted to them in economic and legal respects. It may be that such an arrangement has its technical
difficulties; but the future will surely succeed in getting past them\footnote[1]{Among a Russian sect, the Molokans,
there is community of property between husband and wife, and the husband cannot mortgage or sell the possessions without
the wife's consent. The Molokan woman is accustomed to independence from earliest childhood, and her spiritual life is
developed (within the sect's puritan frame) in an independent way. Cf. N. \emph{Tsakni}: \textit{La Russie sectaire.}
Paris 1888. p. 153 f.}. As for spiritual authority, it can naturally be fixed neither by tradition nor by law, and here
the distribution of authority then often stands (and will more and more come to
% --- p. 307 ---
\opage{307}stand) in a comic disproportion to the traditional and legal distribution of authority toward the outside
world. --- For the rest, in the most recent times in Germany --- the country where the greatest resistance is offered to
woman's independence --- some not inconsiderable advances have been made with regard to woman's legal position within
marriage in the drafting of the new civil code for the German Empire. A proposal was indeed rejected that in differences
of opinion on economic matters the decision should belong to that one of the spouses from whose fortune the greater
part of the family's expenses was met, and the same fate befell a proposal that the wife should have the right to keep
her family name. But through like-minded members of the law commission and the Reichstag the German women nevertheless
won for themselves the right to be appointed guardians and to be members of the family council. It was also laid down
that the husband can forbid the wife to practice a calling or a trade only with the approval of the court of
guardianship. And of the very influence that the German women's associations have exercised through their meetings and
petitions on the form in which the Empire's new civil code was adopted, a German woman who is herself a doctor of law
writes: ``Never before have women taken so lively a part in legislation as on the occasion of the debates on the draft
of the civil code. Never before have they exercised so great an influence on the formulation of the legal provisions that
concern them\dots{}. It is surely a unique case in modern cultural history that every single one of the demands of the
women and their spokesmen was debated and tested in the commissions and found expression in the law, insofar as they
could be reconciled with the views of the members''\footnote[1]{Dr. jur. \emph{Emilie Kempin}: \textit{Die deutschen
Frauen und das bürgerliche Gesetzbuch.} (Schweizerische Blätter für Wirtschafts- und Socialpolitik. 1896. p. 679).}. It
is, however, only a small step forward that has been taken here in the direction of what justice, and the growing
development of woman's capacity for practical and theoretical work, require. German women's associations have
energetically protested against the limitation
% --- p. 308 ---
\opage{308}of woman's rights and duties which the new German code retained.

5. That monogamy prevails in Europe is due chiefly to Greek philosophy, Roman law and Christianity. The influence of the
last, however, has been more indirect; for nowhere in the New Testament is the polygamy of the Old Testament abolished.
Luther and Melanchthon were also clear that as theologians they could not oppose bigamy; in their eyes it was the
business of the secular authorities, not of the Church, to order marital relations\footnote[1]{Cf. \emph{Hausrath}:
\textit{Kleine Schriften religionsgeschichtlichen Inhalts.} p. 239 f.}.

Our marriage is the Roman one, with the abolition of the arbitrary right to divorce which from the end of the
Republic's time played so great a part\footnote[2]{Cf. \emph{Henry Maine}: \textit{Early History of Institutions.}
p. 60. \emph{Lecky}: \textit{History of European Morals.} II, p. 316. \emph{Renan}: \textit{Marc Auréle.} % PRINTED AS IS: Auréle (read at the image)
p. 547.
\emph{Leist}: \textit{Alt-ar. Jus civ.} I. p. 438. \emph{Robert Bartsch} \textit{(Die Rechtsstellung der Frau als Gattin
und Mutter.} Leipzig 1903. p. 39) stresses the Greek influence: ``Es ist sicher, dass das auf hellenistischem Milieu
entstandene Christenthum wesentliche Stücke seines familienrechtlichen Inhalts der griechischen Philosophie verdankt.''
--- Among the Jews polygamy is said to have persisted until about the year 1000. It was then given up --- ``as a
concession to the Occident.''}. According to strict Christian doctrine divorce is impermissible where there has been no
physical unfaithfulness. But here an overstrained idealistic tendency fares, as so often, in such a way that it comes
precisely to lay exaggerated and unnatural weight on the purely physical side of the matter. There can be an
unfaithfulness in mind and will that is more pernicious for the whole relation than physical unfaithfulness. Marriage can
be dissolved in its inner spirit and kernel by the feeling that was to bear and animate it being gone. When the life is
gone, the outer form can indeed be kept up by compulsion, but it then has no value. Honesty and sincerity are conditions
of life in the marital relation; where they are lacking, free, unconditional devotion cannot take place. Resignation and
self-control can accomplish great things; but there is a limit to their power, and when the animating feeling itself has
wholly dried up, it cannot be replaced by anything else. There need, as was already shown earlier
% --- p. 309 ---
\opage{309}(XVI, 3), be no guilt on either side, and yet the relation can become profoundly unhappy. The unhappy
consequences extend over the whole home. Distance and coldness, or even bitterness and enmity, between the parents will be
the most pernicious atmosphere for the children to grow up in.

If the legal order is to be a support for ethical family life, it must therefore not set too great obstacles to divorce.
It would thereby degrade marriage into an institution of compulsion. Out of regard for the interests of the various
individuals there must be prescribed definite forms for the dissolution of marriage as well as for its entry. By these
forms the spouses are protected not only against each other's encroachments but also each against his own rashness. A
certain time ought to pass before complete divorce takes effect; and this time ought to be longer when it is only one of
the parties who wishes the divorce. Faithfulness has its right and ought to have time to test its power. It was an
individualism carried too far when \textsc{Wilhelm} v. \textsc{Humboldt}, in his youthful work on the limits of the
state, considered the declared wish of one party sufficient to dissolve the marriage.

In the suffering that a married life ending in divorce brings with it there will lie warning enough for those who, owing
to the ease of obtaining divorce, would show recklessness in entering into marriage, so that there is no need to fear
that a liberal divorce legislation should weaken the general respect for marriage. In the state of Indiana, which has
for many years had liberal divorce laws, family life is just as fine and lasting as in New York and South Carolina,
where it can be dissolved for one cause only. Among some Russian sects (the Doukhobors, the Molokans), where divorces are
frequent and very easy to carry out, high and noble ideas of marriage prevail, and on the whole the moral level is said
to be higher among them than among their orthodox neighbors\footnote[1]{\textit{History of Women Suffrage.} Boston 1878.
I. p. 742. --- \emph{Lecky}: \textit{History of European Morals.} II. p. 325 f. --- \emph{Tsakni}: \textit{La Russie
sectaire.} p. 143 f. 152.}. Strict marriage laws (especially:
% --- p. 310 ---
\opage{310}prohibitions of or obstacles to the remarriage of divorced persons) have never proved able to maintain and
promote morality. They crush and bind the pure and noble natures; the frivolous do not let themselves be bound, but
always find their ways out. In Catholic countries such as Italy, where divorce is not permitted, marital relations in
many circles are in the highest degree degrading, so that the right to divorce is demanded in order to protect the
sanctity and dignity of marriage\footnote[1]{\emph{Mantegazza}: \textit{Die Physiologie der Liebe.} Ch. 20 and 21.}.
The frequency of divorce depends, in the opinion of a French statistician, not on the laws but on customs, religion,
race, class. Divorces seem moreover to be most frequent in the countries where suicide too is most frequent, and where
there are thus most people who have difficulty in keeping their balance in the various relations of life (individus mal
équilibrés). It is most often the wife who wishes divorce; and divorce is comparatively rare where there are
children\footnote[2]{J. \emph{Bertillon} in ``Dictionnaire des sciences anthropologiques.'' p. 384---385. --- Cf. with
this H. \emph{Westergaard}'s essay in ``Nationaløkonomisk Tidsskrift,'' 1887. --- W. F. \emph{Willcox}: \textit{The
Divorce Problem.} New York. 1891.}.

\grouphead{\textit{2. THE POSITION AND CIRCUMSTANCES OF WOMAN.}}

\chapterhead{XVIII.}{SOCIOLOGICAL DATA.}

1. An inquiry into woman's ethical position forms a natural supplement to the foregoing doctrine of free monogamy. It
appeared that this is first properly realized when not only the man but also the woman has the capacity and the right to
lead
% --- p. 311 ---
\opage{311}an independent human existence. She can fill her place in marriage only when she has the possibility of being
able to fill a place outside it as well. Only then can her choice become free, when both possibilities present
themselves to her. Whatever she then chooses, she chooses not merely as a means of securing her existence, but as a task
of life. --- The time will surely come when it will be unnecessary to enter into a closer justification of woman's right
to independent development and to free choice of a task in life. Just as little as we now have in ethics a special
section on the position and circumstances of man, will the ethics of the future need a section on the position and
circumstances of woman. But the notion has not yet disappeared that woman is by her nature disposed only to be wife and
mother, and that every endeavor to make possible an activity of another kind for her rests on a dubious error.

2. With regard to woman's nature two views stand sharply opposed: the one (held in antiquity by Plato, in the present by
Stuart Mill), that the difference between the nature and capacities of man and woman is only a difference of degree, if
it is present at all, --- the other (held among others by Spencer), that the difference from nature's hand is so deep
and fixed that it must always condition a difference in kind between the position and activity of the two sexes.

One must always be careful in speaking of natural peculiarities and natural differences as if they were eternal and
unchangeable\footnote[1]{The most recent investigations of the nature of man and woman are made use of in A.
\emph{Clod-Hansen}'s interesting book \textit{Mand og Kvinde.} Copenhagen 1895. And it is stressed precisely here, on
several points, how difficult it is to decide what is due to nature and what to education and practice.}. Nature is in
constant development, especially in the domain of living beings. The nature we stand before has itself always
\textit{come into being,} and a new nature will develop out of it in turn. But on the other side this development goes
on slowly, and we must take good heed how far we have come at the given time.

% --- p. 312 ---
\opage{312}Besides, woman's nature and her circumstances stand in interaction with one another. Her circumstances are
determined by her nature, but also act back upon it. Therefore her nature depends to a great extent on what is required
of her and what rights she is given. What matters here, as everywhere, is to adjust demands and rights in such a way that
they do justice to nature and at the same time develop it in the most fortunate direction.

3. When we consider woman's nature and position at different times and in different places, we find different forms of
the division of labor between man and woman.

At the most elementary stage of development known, there is no division of labor. Man and woman each provide for their
own subsistence, and only the sexual relation brings them into a shorter or longer union\footnote[1]{Cf.
\emph{Westermarck}: \textit{Geschichte der menschlichen Ehe.}\textsuperscript{2} p. 218 f.}. When more varied work is
needed in order to live, the strong one naturally loads the most disagreeable work onto the weak one, and woman then
becomes the first slave. She must do all the work that the man will not take upon himself. She must carry and drag,
provide for dwelling, clothing and the tilling of the soil. She is valued first and foremost by her capacity for
work\footnote[2]{See e.g. Th. \emph{Waitz}: \textit{Die Indianer Nordamerika's.} p. 99. --- Fr. \emph{Müller}:
\textit{Allgemeine Ethnographie.} p. 161. 288 ff. --- \emph{Spencer} \textit{(Principles of Ethics} § 428) asserts that
the treatment of women is the saddest subject in history, and that the unwritten history here --- if we knew it --- would
prove to be still sadder. Cannibalism and the maltreatment of prisoners have taken place only occasionally; but the
brutal maltreatment of women is universal!}. Already as a child she is set aside for her brothers. Although she is the
weak sex, she must still use her strength in good earnest. But she also commands more strength than in civilized life.
To cite just one example: an Indian woman frequently delivers herself and immediately after her confinement takes up her
hard work again. In the opinion of some investigators, among the first human beings whose traces we find the difference
in build and strength was not
% --- p. 313 ---
\opage{313}so great as it later --- through altered division of labor --- has become.

A new division of labor comes into force when the hardest work can be loaded onto the conquered but spared enemy. Woman
then gets only the work within the house. She then gets her own domain, which until the most recent times she has kept
under different conditions, according as polygamy or monogamy was the prevailing form of marriage.

Finally, in our day, access has been opened to women in several countries to a whole series of activities which formerly
could be imagined only as carried on by men. They work as physicians, lawyers, clergy, scientific investigators,
engineers, clerks, stationmasters, etc., and they have gained the vote for municipal and political assemblies. The way is
thereby paved for an entirely different development and entirely different circumstances for woman from those which
older times considered the only natural ones. Yet doubt is still raised by many about the necessity and benefit of this
last ordering of labor.

It will be of all the greater interest to examine this question since woman's position can be said to furnish a standard
for ethical progress in the human race. The Greek boasted before the barbarian of his treatment of woman; the Roman
boasted of the same before the Greek, and the Christian in turn before the Roman.

Yet history shows that woman's position at the various times was determined not only by the economic and social
division of labor, but also by the whole character of general culture and of spiritual life.

Within Greek culture it was especially the Stoic philosophers who maintained the ethical equality of women with men.
Even if some tasks lay nearer to the man, others to the woman, there was nevertheless --- they held --- no justified
striving from which the one sex or the other should be excluded. The later Stoics wanted the young women, not only the
young men, instructed in philosophy. It was these same philosophers who had the noblest conception of marriage as an
intimate community of life; a further-reaching spiritual development was in
% --- p. 314 ---
\opage{314}their eyes a condition for woman's being able to fill her position in her domestic circle in the right way.
In this lay a considerable advance, when one remembers that the Athenian women were excluded from spiritual education,
so that the men could satisfy the need for the company of educated women only in the circles of the hetaerae.

Christianity too contributed to a higher position for woman through a deeper-going spiritual community. Faith and hope
were, after all, common, and woman's life as well as man's was borne by the expectation of the great things that were to
happen, and to happen soon. The expectation of the speedy end of ``this world'' gave every individual's life a high goal
wholly independent of their particular lot. Woman could then reach her personal destination even if she was not wife
and mother. For ascetic reasons primitive Christianity even had a preference for the unmarried state, especially for
virginity. Quite apart from the motives underlying this, it was of great significance that woman's life received a great
task independent of her position in the family. Over against this conception of principle it was of lesser weight that
woman's subordination was strictly maintained in oriental fashion, and that woman was to keep silence in the congregation
and let herself be taught at home by her husband. The main thing was the great principle, the imperishable glory of
primitive Christianity, that the highest stood open to all, without distinction of sex, class and race.

In the Renaissance the new, great stream of culture fell to the share of women as well as men, and the strong
development of individualities went on in the feminine world as well as in the masculine. \textsc{Burckhardt}
\textit{(Kultur der Renaissance in Italien.} 4. Aufl. I. p. 124) says of this: ``There is no question of a special,
conscious emancipation, since the thing was a matter of course. The woman of rank had at that time, just like the man,
to strive after a finished personality, complete in every respect. The same process in mind and heart that makes the man
perfect was to make the woman perfect too.'' % translated from Høffding's Danish
It was peculiar to this time that this development of personality first became possible for the married woman. The young
girls were brought up to a great extent in the convents, and their
% --- p. 315 ---
\opage{315}active education began only after marriage. There reigned in this respect --- and there still reigns in the
Romance countries --- a remarkable dualism in woman's position.

This much we see from these historical notes, that when great ideas are astir which reach deep into human life,
there is also astir a tendency to regard men and women as equals. In the face of a great widening of the spiritual
horizon the great difference that had otherwise been made between the position of the two sexes disappears.

\chapterhead{XIX.}{THE ETHICAL POSITION OF WOMAN.}

1. Nature has entrusted to woman the bearing within herself of the germ of the new generation and its nourishment until
it can lead its own life. That this natural office is decisive for her whole organization can scarcely be doubted. There
cannot then be left over for other functions all the energy that the man can have at his disposal. This makes woman the
weak sex. And this also makes her nature purer and better than the man's. Even if animality and coarseness rule in all
other relations, the first germ of humanity is nevertheless found in the relation between mother and child. In the cult
of the mother with the child (Kourotrophos; Madonna) a memory was preserved of this germ of all that is good in the human
world.

The acknowledgment of woman's right had to begin with men's coming to see the great significance of the function which
she, as a natural being, has to perform for the race. It was a great advance when (with the second division of labor)
she was recognized as the weaker sex, and the heaviest burdens were taken from her shoulders. Even this has not yet
completely happened. In the poor classes she must often work in such a way that it comes at the cost of
% --- p. 316 ---
\opage{316}her duties as a mother, and it is an important side of the great social question to bring things to the
point where she can devote herself to her home and her children\footnote[1]{Cf. \emph{Jevons}' essay: \textit{Married
Women in Factories.} (Methods of Social Reform. London 1883).}. --- With the acknowledgment of woman's great significance
in the home a step has been taken that cannot and must not be retraced. The place that woman fills in the family cannot
be replaced by any advance whatever in another direction. The great significance of family life for the development of
the race rests first and foremost on woman's position in it as mother, wife or sister. There is a force here that cannot
be dispensed with. But precisely because the office that woman exercises here has a deep and firm foundation in nature,
we need not anxiously watch that it be taken up. Nature does not need our protection; it will announce itself well
enough. And the differences that really are present between man and woman will always assert themselves, when they are
not due to passing, more or less artificial conditions. \textsc{Stuart Mill}'s contention (in his work on ``The
Subjection of Women'') was that the notions people had formed of woman's nature and capacities rested only on habit and
tradition, not on real experience; what we call ``feminine nature'' was in his conviction only an artificial product.
Precisely for that reason, in his conviction, the barriers hitherto set to the development and use of feminine
capacities had to be removed, so that one could come to know the real feminine nature. It was now a matter of --- for the
first time in the history of the world --- letting experience bear its testimony in this matter! --- \textsc{Mill}
expected that the feminine gifts would show themselves in extraordinary splendor under freer conditions. It is doubtful
whether experience justifies such an expectation, and it is not needed either to ground the demand Mill wished to make.

It is, as already remarked, precisely the wish that woman may be able to fill her place as wife and mother in a still
fuller and more independent way than hitherto that leads first and foremost to the demand for a more many-sided
development of
% --- p. 317 ---
\opage{317}her capacities. Even where she does not take part in productive work proper, her position in the house as
``director of consumption'' will require a training of capacities of various kinds. And in order that she may follow her
husband's activity with understanding and sympathy and be his counselor, as well as that she may direct her children's
education, she must be as well oriented in the spiritual and the social domain as possible. The great influence that
woman has in the family, and that she thereby exercises on the whole ethical, social, religious and political
development, makes necessary as rich a training of her capacities as possible.

To this is added, further, that domestic activity and productive work do not always exclude one another. There may be
use for feminine labor in production too, and it will be no evil if it does not come at the cost of family
life\footnote[1]{Cf. \emph{Marcus Rubin}: \textit{Om Kvindens Adgang til Erhverv.} Copenhagen 1886. (Published by Dansk
Kvindesamfund). p. 10 ff. --- Mary \emph{Gilliland}: \textit{Women in the Community and in the Family.} (International Journal of
Ethics. 1895).}. --- Even where woman strikes out on paths that hitherto only man has walked, she will --- if there
really is the great difference in the capacities and nature of the two sexes that is often maintained --- not deny her
peculiarity under the new conditions. How the peculiarity will express itself, only experience can show.

2. Circumstances over which women have no control bring it about that they cannot always reach the sole destination
that the opponents of women's emancipation assign them. Partly the conditions of mortality, partly emigration, bring it
about that in Europe and in the eastern American states there are more women than men. In Germany there were some years
ago 800,000 more women than men; in the state of Massachusetts there were in the year 1880 over 66,000 more women than
men. In the \emph{country} the number of unmarried women is proportionally somewhat greater than that of unmarried men,
because ``death carries off proportionally more men, so that the chances for women of being married become smaller.''
The great surplus of women in the towns finds its explanation for the most part in
% --- p. 318 ---
\opage{318}more women than men migrating to the towns from the country. The fact is that not only are there on the
whole more women than men, but that this is also especially the case for the age groups in which it is most natural to
enter into marriage (the time between 20 and 40 years). There lived in Denmark in the year 1901 over 60,000 more women
than men; between the ages of 20 and 40 there were 30,000 more women than men\footnote[1]{\emph{Rümelin}:
\textit{Bevölkerungslehre.} (Schönberg's Handbuch der politischen Oekonomie I). p. 1207---1209. --- \emph{Rubin}:
\textit{Om Kvindens Adgang til Erhverv.} p. 21. --- The journal ``Kvinden og Samfundet.'' 1886. p. 126. --- \emph{Rubin}
and \emph{Westergaard}: \textit{Ægteskabsstatistik.} p. 67. --- \textit{Statistisk Aarbog.} 8th year. p. 10---11.}. It
is then a matter of course that the number of unmarried adult women is very large. One will perhaps say to this, with a
German statistician, that a surplus of female population is ``a social evil''; but one does not get it out of the way
thereby, and one has no right to regard the individuals who make up this surplus as superfluous people. The matter must
then be to transform the social evil into a social good by putting the apparently superfluous forces to use. They are
superfluous only if one stops at the second division of labor and will not acknowledge the third. \textsc{Malthus} was
the first who demanded greater respect for unmarried women, and it was his investigation of the population question
that led him to it. Yet not long ago a German philosopher maintained that it is egoistic and ethically reprehensible
when a young girl, in order to preserve her freedom to choose her husband according to her inclination, trains herself
so that her life can become valuable even if she does not marry!\footnote[2]{A. \emph{Dorner}: \textit{Das menschliche
Handeln.} Berlin 1895. p. 422 f.}

3. It has been asserted that woman's physical and psychical nature lacks the conditions for such a training and
development as could make her the equal of man.

That she is as a rule physically weaker than the man has already been mentioned. But her past shows that her strength
suffices for more than one is accustomed to suppose. Civilized
% --- p. 319 ---
\opage{319}life in its one-sidedness has also developed her in a one-sided way. A sounder and more natural upbringing and
instruction will one day put the matter in another light. It must be especially noted that the direction in which all
sound plans for the reform of men's upbringing and instruction go will precisely make it easier to give both sexes the
same training. Greater weight on the development of the power of intuition, on the self-activity of the power of
thought, on the exercise of the muscular system, is precisely what woman in her way needs just as much as man. The
difference in physical strength that will no doubt always remain need not be greater than the one that can develop
among men themselves\footnote[1]{Yet it is said to have appeared that general paresis (a nervous debility bound up with
partial paralysis) has in recent years become commoner than before in women in the civilized states, where they take
part in the work of culture. (Clod-Hansen p. 181). Perhaps, however, this is only a phenomenon of transition.}.

When one holds that feeling is woman's chief psychological faculty, and that she is therefore not disposed to higher
and more independent intellectual development, one proceeds from an unpsychological opposition between feeling and
knowledge. The development of feeling does not necessarily hinder a development of knowledge. It is only feeling as
strong emotion, as affect, that stands in opposition to knowledge. But a feeling that has the character of intimacy more
than of violence need not hinder the development of the power of thought, but can even favor it. Lively sympathy thus
leads naturally to absorption in objects, to taking them up into oneself in their whole distinctiveness, and such a
sympathy is closely akin to the mood of inquiry. In order to gain clarity about itself, feeling on the whole drives one
to inquire into the causes that produce it, and so sets knowledge in motion. Even if that characterization of woman as
a creature of feeling held for all women without exception, it would therefore not necessarily exclude them from the
right to attempt a higher development of their intellectual capacity.

Many peculiarities which, with greater or
% --- p. 320 ---
\opage{320}less right, are cited as characteristic of the feminine nature are surely due only to the conditions under
which woman has lived for long ages, and could be changed when her conditions are changed. This holds of the
psychological peculiarity just mentioned. For the formerly usual feminine upbringing has not been fitted to develop
woman's intelligence and will, but has nourished the indefinite life of feeling at the expense of the other capacities.
With this it hangs together that women are more affected by religious cult than men. ``We read that among the Greeks
women were more easily brought to religious excitement than men. Sir Rutherford Alcock tells us of the Japanese that in
the temples it is very rare to see a congregation consisting of others than women and children; the men at any rate are
few and belong to the lower classes. Of the pilgrims to the temple of Juggernaut it is related that at least five-sixths,
and often nine-tenths, are women. And of the Sikhs we hear that the women believe in more gods than the men
do.''\footnote[1]{\emph{Spencer}: \textit{The Study of Sociology.} Chap. 15. --- This trait in the feminine nature is
very strongly stressed by \emph{Bachofen} as an important element in explaining the rule of women which he holds to have
taken place at a certain stage of the development of culture. \textit{(Das Mutterrecht.} Stuttgart 1861. p.
XIII---XVIII).} % translated from Høffding's Danish
That it is not only in Asia but also in Europe that the churches have their best supporters and most zealous attenders
in women is a well-known matter. Even the ``cult of reason'' under the French Revolution had its most zealous adherents
among women, and these continued to take part in it after the men had begun to stay away from it. An eyewitness
relates: ``Formerly one saw more women than men in the churches; the same is the case in the temples of
reason''\footnote[2]{A. \emph{Schmidt}: \textit{Pariser Zustände während der Revolutionszeit.} III. Jena 1876. p. 239.}.
In present-day France Taine\footnote[3]{\textit{La Régime Moderne.} II. Paris 1894. p. 148.} thought he could calculate
that while in Paris only 1 man in 50 and in the provinces only 1 man in 12 fulfill their Easter duty by confessing and
taking communion, this duty is fulfilled in Paris by 1 woman in 12 and in the provinces by 1 woman
% --- p. 321 ---
\opage{321}in 4. In Copenhagen in the year 1897 (according to a statement in a newspaper) 20,018 men --- but 49,505 women
--- took communion. In London (where the Daily News had church attendance investigated for half a year with the help of a
staff of investigators) there were likewise more women than men among the churchgoers (who made up 16 per cent of the
population). An American psychologist of religion has thought he found a qualitative difference between young men's and
women's religiosity within one and the same confession: in the former will and thought, in the latter fear and the
feeling of imperfection, played a predominant part\footnote[1]{\emph{Starbuck}: \textit{Psychology of religion.} p.
225.}. But even if this hangs together with the peculiarity of the feminine nature, this need not always express itself
in the same way; one cannot conclude from it that woman must always believe in authority. Eduard v.
\textsc{Hartmann}\footnote[2]{Ed. v. Hartmann: \textit{Die Phänomene des sittlichen Bewusstseins.} p. 521 ff. % PRINTED AS IS: Phänomene for Phänomenologie (read at the image)
--- \emph{Rousseau} already said that woman has first her mother's, then her husband's religion, and he found this in
order. \textit{(J. J. Rousseau og hans Filosofi.} p. 141).} holds that a woman must either believe in authority or have
a husband who is a freethinker. It is thus two kinds of relation of authority between which she gets the choice. But if
she chooses the latter (or chooses at all a husband of another faith than the one she was brought up in), her
intellectual activity will necessarily have been set in motion; she must then have the capacity to free herself from
what at first laid claim to her development; and why then should she not be able to arrive at an independent conviction
in religious respects, whether this goes in the one direction or the other? --- The same author holds that, because of
her lively sense for individual cases and individual personalities, she is not able to regard and treat a matter from
the standpoint of a strict and universal rule, and that she might therefore well be suited to be an advocate, but not to
be a judge. Even if this were so, it is always an advance when she is considered competent as an advocate. There are
many men who are not suited to be judges either, and there have been times when
% --- p. 322 ---
\opage{322}woman was not even considered fit to be a witness, or when her testimony counted only half against the
man's!\footnote[1]{\emph{Post}: \textit{Die Grundlagen des Rechts.} p. 451.} Woman's life has so long been passed in the
interior of the family, where the purely personal side of all questions and relations by the nature of the case stands
in the first rank, while the more distant, more impersonal considerations recede. In men whose development proceeds under
similar conditions one will find the same peculiarity. And the sense for the individual, whose shadow sides are used as
an objection against woman, has on the other hand great advantages that can become of much significance in practical
life. \textsc{Stuart Mill} says of his wife: ``In nothing was she of greater value to my mental development than by her
correct estimate of the relative importance of different considerations''\footnote[2]{\textit{Autobiography.} Danish
transl. p. 256.}. % translated from Høffding's Danish
Such a correct estimate is constantly needed in the various relations of life, and it is not frequent. ---

4. Experience has, after all, already borne its testimony, since women in fact work in a whole multitude of domains from
which they were formerly kept out. No one denies that they fill their place well. Whether, as some enthusiastic
adherents of women's emancipation think, a new era will thereby begin that is to reveal undreamed-of and unseen
wonders, time will show. But expectation need not be pitched so high. When it has often been said that no woman has
hitherto produced anything of the very highest rank in any domain, and this has been seen as an objection against
woman's access to independent activity, the answer is that this standard is surely not applied when the question is of a
young man's right to develop the capacities he feels in himself. Most men would fare badly if so ideal a standard were
to be applied to them, and the very great majority of men would have to be glad if their spiritual development, both as
regards intelligence and character, had reached a height like that of Sophie Germain, a George Sand or a George Eliot.

% --- p. 323 ---
\opage{323}Nor is there any reason to lay down with a German philosopher\footnote[1]{\emph{Lotze}: \textit{Grundzüge der
praktischen Philosophie.} § 49.} that what can be done without women ought not to be done by them. This consideration is
surely not applied to men either. It would be a long test that one would have to undertake if one were entitled to
strike out only on such an activity as could not be carried on by anyone else. The sound line of thought, surely, is
that every human individual, woman as well as man, seeks to become clear about his capacity and his impulse and chooses
his way accordingly. Experience must then show whether the choice was the right one. Every such choice is a venture,
but nothing ventured, nothing gained.

Natural differences that hang together with lasting and unchangeable conditions of life cannot be effaced. But it is not
likely that all the differences one thinks one can point out between the nature and gifts of man and woman are of this
kind. Which are the deepest-lying differences will first come to light when both are allowed to use their powers. It
will then perhaps appear that the likenesses are more numerous, and the differences other and finer, than we now
imagine.

A great advance has taken place in comparison with earlier times. It was then considered a dubious sign when a woman of
the middle class could read and write\footnote[2]{``Bei den virginibus,'' wrote an old schoolmaster in the year 1772,
``ist das Schreiben nur ein vehiculum zur Lüderlichkeit.'' Even Justus Möser thought that as a man of the people he
would not marry a girl who could read and write. (G. \emph{Schmoller}: \textit{Ueber einige Fragen des Rechts und der
Volkswirthschaft.} Berlin 1874. p. 120).}, and even \textsc{Goethe} (Zweite Epistel) wished the young girls brought up
in kitchen, cellar and sewing room, preferably with no other reading than the cookbook:
\begin{quote}
Wünscht sie dann endlich zu lesen, so wählt sie gewisslich ein Kochbuch.
\end{quote}
The movement has gone forward step by step, and woman's access to the single domains has not met with so much resistance
as the general principle of her ``emancipation.'' What happens slowly and gradually does not awaken so definite and
clear a consciousness, and therefore not so much
% --- p. 324 ---
\opage{324}resistance either, as what is new and sudden. Yet there are countries where --- especially because of the
prevailing militarism and bureaucratism --- only very small steps have been possible.

5. So far there has been talk only of the possibility and justification of woman's independent development and
activity. But here there can be talk not only of right, but also of \textit{duty.} Every single individual has the task
of making the most possible of his nature, in order thereby to fill his place as a member of the race in the best way.
What woman demands when she wants to be ``emancipated'' is really the right to be able to do her full duty in the
service of mankind, to work with others at the common tasks. In the North American women's movement this side of the
matter comes forward in a fine and interesting way. The American woman first demanded her right when it was necessary
for her in order that she might do her duty. The activity for woman's freer position, that is, developed historically
out of the activity for the emancipation of the Negroes. From the very beginning American women took eager part in this
activity. But their right to appear in public, even in so great and fine a cause as this, met with a bitter resistance,
which reached its height at the great antislavery meeting in London (1840), where the women delegates from America were
refused admission, because such public appearance by women conflicted ``with the customs of the country and the divine
law''\footnote[1]{\textit{History of Women Suffrage.} I, p. 55.}. % PRINTED AS IS: Women for Woman (read at the image)
Only then began the agitation for woman's liberation,
in which so many outstanding women took part. Another example is offered by a German lady, Frau Jeannette Schwerin, who
died in 1899 and worked zealously for women's legal independence. In her practical work among the poor she often stood
helpless, being cut off from helping women who were brutally treated by their husbands, who had the law to support their
power. It was such experiences that led her into the work for women's independent position\footnote[2]{\textit{Jeannette
Schwerin zum Gedächtniss.} Berlin 1899. p. 25.}.

6. When woman has both the possibility, the right and
% --- p. 325 ---
\opage{325}the duty to work at the common human tasks and to train her personality in an independent way, she cannot be
kept out of the municipal and political franchise either. The inner and outer conditions that this presupposes she can
possess just as well as most men, and she has no less interest than men in public affairs being well conducted. The
practice and experience she lacks she can by the nature of the case acquire only through practical participation in
public life. Already now she exercises a great influence in political respects, but since she is kept away from
practical experience, this influence becomes one-sided and is determined by narrow points of view. She also lacks the
sense of responsibility that the right and duty to vote give. When she gets this right, the men will be forced to give a
more serious account of their voting and will not so easily be able to compromise with their
convictions\footnote[1]{Stuart \emph{Mill}: \textit{Considerations on Representative Government.} Danish transl. p. 241
f. --- \emph{Herbert Spencer} fears that woman's franchise, because of her reverence for all authority and her tendency
to put philanthropy above justice, will have harmful consequences. Yet this hangs together with the conditions of our
time of transition, he thinks; the time will come when the female franchise will be beneficial. \textit{(Princ. of Eth.}
IV. § 108).}. Even if husband and wife should vote on different sides, there would lie in this no much greater
misfortune than may already be present now when the wife's political conviction is different from the husband's. The
greater inducement that political rights will give her to think independently will also free her from the authority of
priests and confessors, so that giving women the vote will not in the long run be the same as granting the priests more
votes than they already have. --- Where woman has already obtained this right, it seems to have exercised a good
influence on public life and not to have brought with it the improprieties that have often been feared.

% --- p. 326 ---
\opage{326}\grouphead{\textit{3. PARENTS AND CHILDREN.}}

\chapterhead{XX.}{SOCIOLOGICAL DATA.}

1. Even if maternal love stirs at the very lowest stages of human existence, even if we find parents' love for their
children a frequent trait among savages, --- it can nevertheless be said that the treatment and circumstances of the
child, just like those of woman, furnish a standard for ethical development. At lower stages of culture the child is
wholly in the power of the parents. The father of the family has power of life and death over it, can sell or kill it,
and is accountable to no one for it. What treatment the child receives will depend much on the outer conditions under
which the family or the tribe lives. When we find the exposure of children, especially of crippled children, of triplets
and of --- girls\footnote[1]{So it was early held that too many women are a ``social evil.''}, among almost all savage
and barbarian peoples, indeed even among some peoples civilized in other respects, it was hard want that at first led to
this custom. A roving band of savages, always in danger of famine or of attack by enemies, rids itself out of the urge
for self-preservation of children, the sick and the old, who hamper its movements and consume its stores. Among savages
and barbarians the killing of children is chiefly the work of fathers and husbands; but they take the decision about
the newborn's fate on behalf of the family or the band. In any case --- whether it is the mother or the father who kills
--- it is regard for the welfare of the tribe that is the predominant motive for infanticide, as also for girls' being
especially killed. ``Where infanticide is the custom, the struggle for existence will be less severe, and all the members
of the tribe will have an almost equally good chance of rearing their few surviving children. In most cases a larger
number of
% --- p. 327 ---
\opage{327}female than of male infants are destroyed; for it is obvious that the latter are of more value to the tribe,
as they will, when grown up, aid in defending it, and can support themselves. But the trouble experienced by the women
in rearing children, their consequent loss of beauty, and the higher estimation set on them and their happier fate, when
few in number, are assigned by the women themselves, and by various observers, as additional motives for
infanticide.''\footnote[1]{\emph{Darwin}: \textit{Menneskets Afstamning.} II. p. 360.} % after Darwin's English (Descent of Man, ch. 20), as far as Høffding's Danish allows
Under our ``civilized'' conditions it is most often the abandoned, lonely mother who commits infanticide, and her
motive is shame as much as want. Maternal love often proves stronger than want, while it succumbs to shame. It has been
thought that infanticides are rare in countries where the judgment on women who give birth outside marriage is very
mild\footnote[2]{\emph{Lecky}: \textit{History of European Morals.} II. p. 37.}. Yet sad experiences show that want and
overburdening with work can weaken and eradicate maternal love\footnote[3]{\emph{Jevons}: \textit{Married Women in
factories.} (Methods of social reform). p. 157 f.; 166.}. --- In many women maternal love first expresses itself when
they have seen the child and begun to show care for it. Despine cites (in his \textit{Étude psychologique sur les
personnes qui commettent l'infanticide,} in the 3rd volume of his ``Psychologie naturelle'') a series of examples of
this. It thereby becomes intelligible how the abandoned mother can take the life of her child immediately after her
confinement: the countermotive comes too late. --- For prostitutes it is an honor to have children; they therefore very
rarely commit infanticide. --- In the states of antiquity the notion of the father's absolute authority persisted, and
with it was bound up, in the case of the Greeks, the assumption that a state fared best with a very limited increase in
population. This explains why the ancient philosophers and lawgivers maintain the right, or even the duty, of depriving
the conceived or already born child of life. They thereby introduced nothing new, but confirmed a very old
% --- p. 328 ---
\opage{328}custom, which only Christianity, with its care for the soul of every single human being, led men to
overthrow. Only now was the right of the unborn and the newborn child maintained\footnote[1]{Cf., for the North: R.
\emph{Keyser}: \textit{Efterladte Skrifter.} II, p. 315. --- The strict Christian way of thinking went so far that when,
in a difficult childbirth, the mother or the child had to be sacrificed, the mother was sacrificed, so that the
unbaptized child should not go to hell. \emph{Lecky} ib. p. 25.}.

2. As in woman's case, so in the child's, it was partly economic reasons that led to better circumstances for it.
Before the thought of the value and significance of the single individual could assert itself, the help and use that
children could afford had to bring it about that they were treated gently. Children performed work like the thralls.
But to this were added, at least in the case of sons, other considerations too. In his son the father saw the one who
was to avenge him if he came to a violent end, and the one who was to perform the sacrifices at the family altar after
his death. Only if he left a son could he pay his debt to his forefathers by securing the continuance of the family cult.
But the relation nevertheless remained until the most recent times a relation of authority and subjection, the parents'
authority being acknowledged in a way that could still recall the primitive family. Since the state has set narrower
limits to parental authority, and since ideas of freedom and equality have become dominant in the rest of society, the
relation between parents and children has become far more a pure relation of sympathy than a strict relation of
authority and obedience. The family has thereby itself developed into being, to a far higher degree than at the
primitive stage, the model of an intimate and perfect society, in which no member is a mere means, but all are
independently cooperating members.

% --- p. 329 ---
\opage{329}\chapterhead{XXI.}{ETHICS AND PEDAGOGY.}

1. Physiology already gives us a clear hint about the ethical position of parents toward their children\footnote[1]{Cf.
with regard to what follows: \textit{Psykologi} VI C, 2---3.}. The parents ``put'' the child into the world, to be sure,
but they are themselves only intermediate links in the order of nature which, long before their consciousness and their
impulses developed, had laid the foundation for the first germ of the new individual. No right of property or of
dominion can therefore be grounded on their physical relation to the child. Their authority over the child is grounded,
on the contrary, on the office that rests on them as those nearest: to tend the germ and guide its growth. Their
authority is grounded on the necessity of education. Parental authority can, like all other authority, involuntarily or
voluntarily make itself an end instead of a means. This happens when the power placed in the parents' hands is enjoyed
as power, so that the children become means for the satisfaction of the craving for power. It happens also when the
parents give free expression to the feelings awakened by the children's conduct, without considering that both praise
and blame are difficult means to use, and that the main thing is not that the parents vent their feeling, but that the
judgments they pronounce can have the right motivating influence on the development of the children's character. The
right to exercise power and to pronounce judgments rests on the duty to be an educating providence. This duty is violated
in a still more involuntary way by every lack of self-control which, either directly, by crushing and tormenting the
children, or indirectly, through the example given them, can conflict with the end of education.

Under sound conditions that duty is borne up by one of the strongest natural feelings, which is nourished and
strengthened the more constant the opportunity for its satisfaction. Instinctive activity precedes the feeling proper,
the intimate feeling, and
% --- p. 330 ---
\opage{330}this grows with continued activity. Living together produces fellow-feeling. (Cf. the Aristotelian
principle). The common experience that benefactors love those to whom they do good more than they themselves are loved
by these is confirmed here. The feeling in the children that answers to the parents' love can hardly reach the same
degree of strength. Homer often uses the expression of young heroes who fall that they did not pay their parents the
``return for their rearing,'' the requital for what they had received from them. But however strong children's gratitude
and reverence may be, they will never be able to pay the full return for their rearing. The reason is that the child's
gaze is naturally directed forward. However much nourishment it may draw from the past, it is its impulse and its calling
to look forward and not back, and when the parents' sympathy is disinterested, they too will wish it so. They will know
that it is theirs to decline, the children's to rise. % cf. John 3:30
Yet where the right relation is present, the parents' sympathy with the children's progress will meet the children's
reverence for the parents, and thereby that harmony between the older and the younger generation will arise which is one
of the most significant sides of the family. (Cf. XIV, 2 and XII, 6).

Besides the bond that unites parents and children with one another, the one that unites brothers and sisters with one
another also has a great ethical significance. Community in origin, in conditions of life, in memories, in upbringing
and activity form --- besides the natural kinship in temperament and disposition --- a firm foundation for community and
understanding, which can hold even if the later experience and activity of life become most different. No better
expression has been found for the society borne by general love of mankind than universal brotherhood. ---

2. Even if it may now be said to be settled that the child is not merely the parents' property and a means for their
ends, it might yet seem as if the child must in another respect be treated as a means and not as an end. Childhood is,
after all, the beginning of life, a transition and preparation for full and proper human life. The child must then ---
it might seem --- be regarded as a means, since the point is to
% --- p. 331 ---
\opage{331}form it into as vigorous and perfect a human being as possible: so it must be a means for its own future
human being!

But just as little as the whole of life may be regarded as a mere means, just so little may any single part of life be
regarded as a mere means for another part, when there is no compelling necessity present. Childhood stands first and
foremost as an independent and distinctive period and form of life, which has its own conditions and is to be judged by
its own laws. It is no more a mere time of transition than every part of life is. The child must first and foremost be
regarded as a child, not merely as a grown-up in the making. Out of the larva, to be sure, an insect develops in the
fullness of time; but that does not mean one can treat the larva as an insect. The larva has its own peculiar
capacities and needs. Apart from ascetic tendencies that regard all healthy joy of life and natural activity with
suspicion, childhood has no worse enemies than an external doctrine of utility which asks only what the child \textit{would
have} use for, and not what it has use for.

First and foremost the point is that the child be allowed, and be helped (for it may need help), to be a real child, to
live through the time of childhood with complete devotion, undisturbed by anxiety, want and toil. It thereby gathers
bodily and spiritual strength for future development. --- And next the point is that the work required of the child
become, as far as possible, its own work, so that it can feel the joy of using its own powers. --- When these conditions
are fulfilled, the problem is thereby also solved: to treat childhood as an independent and distinctive period of life,
and yet to make it a preparation for the life that follows.

It is the great merit of \textsc{Jean Jacques Rousseau} to have set forth these truths in such a way that they cannot be
forgotten again\footnote[1]{See my book \textit{J. J. Rousseau og hans Filosofi.}\textsuperscript{2} p. 89---91;
133---148.}. He thereby became at once the evangelist of childhood and the reformer of pedagogy.

3. For the child, just as little as for the grown-up, is obedience
% --- p. 332 ---
\opage{332}the highest virtue. Obedience has no value in itself, but is only a means. It is necessary that the child
learn obedience, since otherwise it cannot have its full freedom. One can give the child free play only if there is an
elastic bond that holds it back when it goes too far. When the child approaches an abyss without suspecting the danger,
only obedience to the warning voice can save it; there is no time to give reasons and closer explanations, but trust in
the voice that is heard must be sufficient motive. Toward the child the parents represent mature experience and
independent judgment. Until the child attains independence, it is guided by the need to find approval of its conduct
from those it loves most; out of this need there then gradually develops the need for self-esteem and with it ethical
independence\footnote[1]{This is very finely shown by J. G. \emph{Fichte}: \textit{Reden an die deutsche Nation.}
Berlin 1808. p. 317---329. --- \emph{Preyer} \textit{(Die Seele des Kindes.} 3. Aufl. p. 227) remarks of his child's
eagerness to imitate what the grown-ups did: ``Sein Nachahmungstrieb erscheint hier fast wie Ehrgeiz.''}. The child has
its conscience outside itself before it can have it within itself. But obedience and trust in others lead to independent
conviction only by leading to action. We are reminded here once more of the Aristotelian principle that one becomes good
only by acting well. The involuntary, instinctive way in which the child follows the conduct of its elders and imitates
them, the eagerness with which it strikes out on the paths that are opened to its strong urge to use its fresh powers,
counts every bit as much as conscious obedience. Without the child itself noticing anything of it, precedents are created
which determine the future and furnish a better foundation than admonitions and threats, rewards and punishments.
Experiences it will later need it can in this way make on a small scale, in such a way that it feels joy in its own
activity. When the time comes that it understands the commandments and admonitions, it has already practiced them; and
that is really the only way in which it can properly understand them. Just as, according to newer methods in language
teaching, one does not begin with grammatical rules but with practical exercise in the language,
% --- p. 333 ---
\opage{333}so in practical ethics too one ought not to begin with moral rules, but with the exercise of the powers.

Unconscious education is every bit as important as conscious. There are thoughts and feelings that first become possible
when activity goes before. --- But by unconscious education one can also understand the education that parents carry on
without knowing it themselves. Their unguarded moments act on the child every bit as much as the moments in which they
appear before it with clear pedagogical consciousness; in one's unguarded moments, after all, one often speaks and acts
most energetically. It is therefore of no use to lay down the most excellent pedagogical system when numerous unguarded
moments tear down what the system has built up. In general a great part of the child's education consists in involuntary
imitation, which leads to training in activities and tasks it will later carry out with consciousness. Through active
imitation the child comes to know the various sides of life. In its play it repeats the conduct of the grown-ups,
appears in various dignities and situations in which it sees the grown-ups appear, is now master, now servant, now
giving, now receiving. It comes to know the contrasts of life, and it learns besides that even the grown-ups, who
otherwise seem freed from all barriers, have rules and considerations by which they direct themselves. Step by step it
is led into the great coherence of life, and moral development can here go hand in hand with intellectual. The capacity
for sympathy, for rejoicing and suffering with others in lively fellow-feeling, is nourished by this course of
experience called forth by the involuntary urge to imitate, which only in the most recent times has been more exactly
apprehended and described\footnote[1]{\emph{Josiah Royce}: \textit{On Certain Psychological Aspects of Moral Training.}
(International Journal of Ethics. 1893). --- \emph{Mark Baldwin}: \textit{Mental Development in the Child and the Race.}
Methods and Processes. New York 1895. --- The same: \textit{Social and ethical interpretations in mental development.} A
study in social psychology. New York 1897. (A treatise awarded a prize by the Royal Danish Academy of Sciences.)}, and
which is an important example of the Aristotelian principle (XIII, 4).

4. In the domain of intellectual education one sins
% --- p. 334 ---
\opage{334}against the proposition of the independent significance of the life of childhood when one wishes to develop
and exercise the child's capacities by means that have no interest or value in themselves. So when one justifies the
usefulness of grammatical rigmaroles by saying that they strengthen the memory. What is used for the development of the
formal capacities ought to be something that can also have interest through its content, so that satisfaction can be
felt at having appropriated it.

And in intellectual education it is of the greatest importance that self-activity be awakened. However different
childhood and grown-up age may be, they agree in this, that the highest that can be reached is to work with one's own
strength in the service of a significant cause. That is possible for the child in its way, as for the grown-up in his.
The material one requires the child to appropriate must therefore be arranged so that the child can become self-active,
whether it is the imagination or the understanding that it comes especially to work with. Through self-activity the
child comes to feel joy in the work, without suspecting that it is thereby at the same time laying the best foundation
for its own future. The work of \textit{learning} lies, as \textsc{Dühring} has aptly remarked\footnote[1]{\textit{Der
Werth des Lebens.} 2. Ausg. p. 95.}, between play, which overcomes self-created obstacles, and work proper, which
overcomes real, objective obstacles; the obstacles that the work of learning has to overcome are those that lie in
sluggishness and lack of practice.

Only an ideal pedagogical system could completely unite the two considerations: to give childhood its independent
significance and yet let it be a time of preparation. But all reforming thoughts in pedagogy point in that direction.
That there will be great difficulties in their full execution follows of itself. It lies already in the fact that it is
grown-ups who educate; the insect can never wholly understand the larva. For the time being it is a great good that
attention is so strongly awakened to the significance of education. If it is so extraordinarily difficult for grown-ups
to make capacities and tasks fit together,
% --- p. 335 ---
\opage{335}it must surely be possible to adjust them to one another for the child, since one is oneself master of what
tasks one will set it; and if it should once succeed to solve the problem of education, grown-ups too would thereby be
helped in the struggle with their problems.

5. The child's self-activity naturally has its limits. It cannot itself make anew all the experiences that the race has
made on its long way. Experience educates us to a great extent through disappointments, through expectations that are
not fulfilled. We thereby learn to fix the boundaries between the real and the merely possible\footnote[1]{Cf.
\textit{Psykologi} V B, 4.}. But it is the educator's duty to see to it that the disappointments do not become too
great for the child, so that they crush its courage. Too abrupt a transition from expectation or illusion to reality can
easily stifle the self-confidence and feeling of security that the child can least of all do without. A small
observation made by \textsc{Preyer} illustrates this. ``In the 62nd week the child could still not stand longer than a
moment without being supported or touched. This lack of ability no longer rests on the difficulty of keeping its balance,
but on lack of self-confidence; for it cannot stand alone only when it knows that one is not holding it. But when it does
not know that the father's supporting, ever less pressing hand has been removed from its back, it can stand for several
seconds without support''\footnote[2]{\textit{Die Seele des Kindes.} 3. Aufl. p. 216.}. % translated from Høffding's Danish
What holds of bodily walking holds also of spiritual development. The child's ideas and feelings develop in secure trust
in the provisional frame it has formed for its view of the world, partly involuntarily, partly by taking things up from
the utterances of those around it; and if this frame is burst too early, it can hamper development and bring in an
insecurity that can harm the soundness of life. The child begins by naively ascribing validity to all its ideas. It
begins with a sanguineness that anticipates experience. A brutal tearing apart of illusions, an all too sudden
inculcation of the difference between the world of imagination and that of reality, the child cannot bear. Training in
this
% --- p. 336 ---
\opage{336}difference extends, after all, through the whole of life and cannot be brought to an end all at once.

The art of education consists in letting the child's sanguineness and imagination have their free play as far as
possible, and yet exercising a criticism that cuts away what is pernicious and stresses the sound kernel.

Nowhere does this find a more important application than in religious ideas. The child's way of conceiving and
explaining things to itself has much in common with that of natural man. They both have a tendency to find the
explanation of what goes on around them in the intervention of personal beings. They personify nature and really
understand only personal causes. It does not always follow that a poetic view of nature comes out of this. Precisely the
tendency to personify can give the explanation an external and mechanical character. When the child explains the
appearance of the stars in the evening by saying that the Lord lights them, just as ``the man'' lights the gas lamps, or
when it thinks that in spring a man comes and puts the leaves on the trees, there is not so much poetry in this as in
the scientific explanation. But the child arranges things in its own way, and appeals to such causes as are best known
to it. That is why it also takes up with ease the biblical stories, which suit its whole way of imagining the world so
well. The child must have peace to go through this whole mythological period without being disturbed every moment by a
criticism proceeding from entirely different presuppositions. Both dogmatic and antidogmatic educators often offend
against this. Just as missionaries often mercilessly declare the faith of the heathen to be the work of the devil, so
the child's wildly growing mythological ideas are often torn apart by anxious parents because they do not agree with the
catechism. And freethinkers often think themselves obliged by love of truth to stunt the growth of all mythological ideas
in the child. In both cases the child's free natural development is hampered. One must allow it to run through a
mythological period similar to the one the race has gone through, just as in its early embryonic development it has
exhibited forms that recall lower species of animals. What one can do is only to shorten and ease the way; but it must
% --- p. 337 ---
\opage{337}work its own way out of this period if the development is to be sound. One can help it to draw all the
consequences of its ideas and thereby test these in particulars. One thereby develops its criticism and its capacity to
infer. One can give it the opportunity to make as many fruitful experiences as possible. One can impart to it an idea of
how difficult it is to reach exact knowledge, and how much it still has to learn in order to treat properly the
interesting questions it raises. And one can show it that the most important thing is to have a good heart and a clean
conscience, and that these can be found even where there are very different religious ideas.

Too early direct criticism will easily give the life of the understanding a preponderance over the other sides of the
life of the soul. Knowledge has on the whole a tendency to develop more quickly than feeling; and this tendency easily
brings with it a morbid one-sidedness when one strengthens it by forcing on the child what one considers truth, instead
of letting the child find it by its own way. There will then arise a psychical disharmony, a dividedness and paralysis in
the mind. And later there will then easily come a period of reaction of an unhealthy kind. By a kind of effect of
contrast the individual then later longs to experience a circle of feelings which it did not come to know at the
natural time.

If self-confidence and self-activity can be crushed by too early criticism, they are naturally crushed still more when
the child is systematically instilled with the idea that it dare not trust its own senses and its own understanding, but
must learn to take them captive, since what ought to be believed cannot be understood. But such conduct lies too far
removed from the whole standpoint from which we here proceed to need any closer criticism. An education conducted in an
ethical spirit cannot lay absolute weight on any belief or unbelief whatever. Its goal is the competence and independence
of character, the intimacy and soundness of the feelings, the strength and clarity of the understanding. All
confessional oppositions are of vanishing significance in relation to this goal. It is a \textit{human being} that is to
come out of education, not a believer or an unbeliever.

% --- p. 338 ---
\opage{338}\chapterhead{XXII.}{THE STATE AND THE CHILDREN.}

1. At the modern stage of culture the family is no wholly closed world, as the patriarchal family was, which stood in
connection with the larger society only through its head. The right of the single individuals, woman and child included,
to independent life has been acknowledged, and it is the business of the state to maintain this acknowledgment. The
demands the state makes on the family rest partly directly on its calling to protect the weak, partly on this, that the
members of the family are also members of the state and must fulfill the conditions thereby required.

Thus the state maintains the right of the unborn and newborn child to exist, and permits no one to declare any other
individual superfluous or unentitled or unfit to exist. Even where the prospects for the individual beginning its
existence seem ever so dark, such an authority to cut the thread of life once spun cannot be ascribed to anyone.

The state further takes care of abandoned children and holds parents to doing their duty where their own conscience and
parental feeling do not drive them to it. It protects children against their parents. For it turns out that maternal
and paternal feeling can be dulled and eradicated by want, shame and other causes, and where the family relation thus
does not afford sufficient motives to fulfill what the principle of welfare requires, the larger society steps in,
although it cannot solve the task as well as the narrower one can under sound conditions. It is by virtue of this
principle that the state supervises and restricts the employment of women and children in factories, although it has
hardly yet done enough ``to secure the child's right to its mother's breast.'' By virtue of this principle the state
ought also to compel parents to acknowledge their children, which, however, in the case of illegitimate children it for
the most part does not do. Toward illegitimate children the state follows a barbarous rule: mater semper certa, the
mother is certain, we hold to her. The burden
% --- p. 339 ---
\opage{339}is laid predominantly on the mother. The child's right to the father's name is not acknowledged; and even
where the father is identified, his contribution to the child's maintenance is fixed according to the mother's, not
according to his own, social position. Even in the new civil code of the German Empire it says (§ 1567): ``An
illegitimate child and its father are not deemed to be related''; --- in the original draft it even said: ``Between an
illegitimate child and its father no kinship exists.'' A striking example of the difference between juridical and moral
thinking! --- When such an arrangement is defended on the ground that otherwise loose unions would be encouraged and the
standing of marriage weakened\footnote[1]{\emph{Goos}: \textit{Almindelig Retslære.} I, p. 544---551.}, it is
overlooked that there is here an immediate and natural title on the child's part which cannot be pushed aside for more
remote and indefinite considerations. Here the child is in fact made to suffer for the parents' sins, --- once more a
barbarous principle which more developed ethical feeling has elsewhere rejected. The child's right to the father's name
will moreover make people feel greater misgivings about bringing children into the world outside
marriage\footnote[2]{Yet the statistical experiences here seem to contradict one another. Cf. \emph{Oettingen}:
\textit{Moralstatistik.} 3. Aufl. p. 323 with \emph{Rümelin}: \textit{Reden und Aufsätze.} Neue Folge. p. 620. Rümelin
maintains that the French rule: la recherche de la paternité est interdite, strengthens woman's power of resistance
against the seducer.}. An easier access to entering into lawful marriage will restrict the number of the unions from which
illegitimate children spring; the introduction of civil marriage is said in several countries to have brought about a
decrease in illegitimate births. It must not be forgotten, moreover, that such unions are far from being the worst, the
most degrading form in which human beings' lack of capacity to realize the ideal of sexual union, free monogamy, shows
itself.

2. The state must also protect the child's spiritual development, both because it is to protect the weak and because
the child cannot otherwise come to fulfill what the state requires of its citizens. Where the parents cannot or will not
provide for their children's instruction, the state therefore intervenes with force. Here too it must often protect the
children against the
% --- p. 340 ---
\opage{340}parents, when these keep them away from school in order to profit from the children's
work\footnote[1]{The English workers have seen the importance of this point both for the children and for the whole
working class. The English trade unions have worked zealously for the introduction of compulsory schooling. (L.
\emph{Brentano}: \textit{Die Arbeitergilden der Gegenwart.} II, p. 100).}. Where the parents of their own accord
provide for the children's instruction, the state has only to make sure that the minimum of knowledge it requires is
present.

When it has been held that the state thereby encroaches on the parents' right, the assumption is that the children have
no independent right which the parents could violate partly from indifference, partly from ignorance, partly from
self-interest. Experience shows that a control of parents in this respect cannot be dispensed with, and it lies without
doubt (which, however, only the theory of the state can ground more closely) within the circle of the state's tasks. The
state must have authority to require that every normal child have the knowledge that is necessary for being a citizen
of it.

That it is the child's right that is in question here is also expressed in the Danish Constitution, whose § 85 reads:
``Children whose parents have not the means to provide for their education shall receive free instruction in the common
school.'' In the first draft of this paragraph, on the other hand, there was talk of the parents' right to receive free
instruction for their children when they could not themselves pay for it. The revision was expressly justified in the
constituent assembly by this, that ``it is after all not the parents but the children themselves to whom the right here
in question belongs.'' --- In the corresponding provision in German constitutions (stemming from the German imperial
constitution of 1849) there is also talk of a right of the children. Parental power is restricted so that a general
human fundamental right may be protected\footnote[2]{Cf. \emph{Rümelin}: \textit{Ueber das Object des Schulzwanges.}
(Reden u. Aufsätze. Neue Folge). p. 475.}.

One need not fear that the parents' responsibility will thereby be diminished. \textsc{Spencer}, led by this misgiving,
has protested against all state instruction. He
% --- p. 341 ---
\opage{341}demands that family and state be kept wholly apart and is convinced that both can thrive only by standing as
far from one another as possible\footnote[1]{\textit{Principles of Sociology.} I. p. 738 ff.}. To this it must be said
that one surely does not take ethical responsibility away from a person by seeing to it that something gets done which
he is indeed the nearest to do, but which he nevertheless cannot or will not do. The state here carries on a helping,
supplementing activity. Family and state are, to be sure, very different in their nature, but they have in part common
ends, since the members of the family are also those of the state. The larger society steps in where the smaller for one
reason or another does not perform its office. What is required is a minimum, and it is always open to parents to go as
far beyond this as they can and will. ---

A further-reaching education the state can offer, and can make a condition for obtaining public office, but must not
force on anyone. The same holds of the religious view of life. The state can use its influence to work for the free
unfolding of the various spiritual tendencies in the people, but must leave it to the individual homes to decide in which
direction the child's religious life is to be guided. Its public teaching institutions ought to approach freedom from
confession as nearly as is at all possible.

% --- p. 342 ---
\opage{342}\grouphead{B. THE FREE CULTURAL COMMUNITY.}

\chapterhead{XXIII.}{FREEDOM AND CULTURE.}

1. Already in family life the standard of development had to be this: to what degree the single individual stood not
merely as a means but as an end. This standard must find still greater application in the free cultural community. The
individual belongs to the family even before it has reached its maturity; but the free cultural community, which comes
into being through the union of individuals around common or connected ends, presupposes that the individuals come with
mature and developed powers to work in the service of their ends.

Within the family too work is done for cultural ends, but there the connection between the individuals is the first
thing, what lies at the foundation, and the common endeavors are a consequence of this connection. In the cultural
community, conversely, it is the common ends that awaken common endeavors and thereby bring about connections between
the individuals.

Freedom and culture are the two concepts that in union make clear the nature of the free cultural community. They might
seem to stand in a certain opposition to one another: freedom isolates, culture unites; freedom stresses independence,
culture devotion. But they are very closely tied to one another.
% --- p. 343 ---
\opage{343}As soon as the course of life brings it about that human beings are subject to common conditions, have a
common fate, common danger, common enemies, common work or common play, there will develop between them a \textit{feeling
of solidarity,} which makes freedom and devotion go together as one. The independence of the individual is then not
hindered, but is made possible precisely by his working together with the others or for the others. The development of
culture is not possible without such a solidarity, since it sets tasks that can be solved only by a work exceeding the
individual's powers. But cultural tasks are also such as do not benefit the individual alone. Culture produces goods that
can directly or indirectly come to have significance for all, and through the work of culture the individual thus works
for something more comprehensive than his own little world. By sharing in the work of culture man is therefore developed
both through the solidarity that community of work produces, and through the personal entering into universal ends,
which is culture's best fruit.

Often there may be a long school to go through before community in work and in ends properly comes into being. A process
of education then goes on, which can be strict and hard. Sluggishness and coarseness, egoism and blindness must be
overcome, but can put up strong resistance. Life in the cultural community is, seen from this side, a continuation of
the education that takes place within the family. And it is always the most important significance of cultural life to
have an educating and developing effect on personal life. The whole history of mankind can be regarded as a process of
education that is never finished, but comes forward especially when new ends are to be worked out, and when community in
working for these ends must therefore first be produced. When a society has already formed around a certain cultural end
as its center, the task for the single individual is to live itself into the finished culture. It is then determined in
its development by tradition; imitation and suggestion play a great part without its being noticed; finished
associations of ideas and results of displacement benefit it and determine the level from which it begins its
development,
% --- p. 344 ---
\opage{344}and from which --- if independence and criticism awaken --- new processes of education and displacement can
then be initiated.

2. In the 18th century personal freedom was set up as a self-evident right of man\footnote[1]{In theory \emph{Kant} formulated
the right to personal freedom best: ``Freedom (independence from the constraining arbitrary will of another), insofar as
it can coexist with the freedom of every other according to a universal law, is the only original right belonging to
every human being as human being.'' \textit{(Rechtslehre,} 2. Ausg. Königsberg 1798. p. XLV). % translated from Høffding's Danish
--- In practice this right of man was proclaimed in the Declaration of Independence of the North American free states,
where it says: ``We hold these truths to be self-evident, that all men are created equal, that they are endowed by their
Creator with certain unalienable Rights, that among these are Life, Liberty and the pursuit of Happiness.'' % the Declaration's English
--- (It hardly needs remarking that freedom in the sense in which the word is used here has nothing to do with the
``freedom of the will'' discussed in ch. V).}. But freedom cannot be said to be a self-evident and absolute right. It is
a derived principle. It needs a justification, and this is given by inferences from the principle of welfare, from
which, as we have already seen (VIII, 6), the principle of the free personality springs by way of a psychological
consideration. Personal beings are the central places in the world where the value of life is felt; a violation of the
unhindered and harmonious unfolding of the life of such a being therefore conflicts with the principle of welfare. Every
intervention in the free development of personality must be justified by regard for the continued course of this very
development; it might be, after all, that a check at the single moment or on the single point was necessary in order
that the whole development might go on as freely and harmoniously as possible. Because of the close connection between
freedom and welfare, freedom is itself an end. But it is at the same time a means, prescribed by the principle of
welfare. For without an independent and harmonious life in individuals the life of the race as a whole cannot go
forward. Initiative, the new beginning of what is great and good, comes into being only when life is allowed to stir
freely in the single individuals. Compulsion and authority would here be insufficient; they may suffice for a time to
hold fast what has once been acquired,
% --- p. 345 ---
\opage{345}--- and yet independent appreciation of what has been handed down, and enthusiasm for maintaining it, will
require a spiritual freedom and an independence of material barriers which do not come into being where individuals are
subject to dominion and tutelage.

But although we can thus well distinguish between \textit{freedom as goal} and \textit{freedom as means,} they hang
closely together and condition one another. In order to have freedom as a means, one must in part have attained it as a
goal; in order to work as a center of force one must have gathered force within oneself. And conversely: only through
free use of one's powers does one become a real personality. We find here once again the Aristotelian principle: we
become free only by acting with freedom. Herein lies precisely the difficulty in applying the principle of freedom. Not
only is freedom a derived principle, but it can often --- both as goal and as means --- be applied only indirectly,
since the individuals who are to be made free may need support and care in order to get so far that they can act on
their own. The use of spiritual and physical power can be justified in the name of freedom; but the use of power or
compulsion, the principle of authority in any form, is nevertheless always derived relative to the principle of freedom,
just as the principle of freedom itself is relative to the principle of welfare. Welfare is goal; freedom is both goal
and means; but authority is really only means, a means that is to disappear when the goal is reached.

Only when authority takes the place thus assigned to it is it of an ethical nature. Between the principle of authority
and the principle of freedom stands the great struggle in history. They are the two forces whose interaction determines
social development. They each have their time and their domain, which are determined by the principle of welfare. --- It
makes a very great difference whether authority is used as a means of education or as an unconditional principle. The
whole spirit and manner become different where it is regarded as a means from where it is regarded as an end. It is
naturally presupposed here that the acknowledgment that the use of authority is only a means is sincere. For it is very
possible to have an abstract enthusiasm for freedom as a distant goal and yet in
% --- p. 346 ---
\opage{346}one's actual conduct to use power and authority in the most brutal way. That is one of the many forms in
which the always questionable and not always honest distinction between theory and practice appears. The standard for
the development of culture and of the cultural community consists in the degree to which freedom stands not merely as a
theoretical idea, but --- both as goal and as means --- lies at the foundation of the ordering of the single relations of
life.

By setting up the principle of freedom as the subordinate principle derived first and foremost from the general
principle of welfare, we part company with \textsc{Bentham}, for whom security and equality stood as the first derived
principles. But it is clear that security receives its value only as a condition of free possession, enjoyment and
activity. Insecurity hampers, crushes and divides. It is not freedom that is the condition of security, but the reverse.
And if one wants security at any price, it easily comes at the cost of freedom. Risks must be taken for development to go
forward, and in every such venture it is trust in the effects of free powers that keeps one up. By fixing one's gaze on
security in the first line one easily comes to stop up the source of development. Out of fear of the danger of continuing
life, one stops life. As regards equality too, it is clear that it is a subordinate condition of life. What matters is to
bring out personal life in individuals. But since these, the better their nature is known, prove to exhibit quantitative
and qualitative differences, they can be treated wholly alike only in purely external relations. Distributive justice
requires, as we have already seen (III, 9; IV, 2; XI, 9), a far-reaching individualization. Equality too, then, is a
condition subordinate to freedom. The distribution of material and spiritual goods, and the fixing of the demands to be
made on individuals with regard to material and spiritual culture, must take place with a view to their capacity to
satisfy the principle of freedom, that is, to their capacity to develop freely and harmoniously and to work with freedom
in the service of the ends of culture.

3. It was the great glory of the 18th century that it discovered the principle of freedom and thereby established one of
the most important
% --- p. 347 ---
\opage{347}conditions of a higher human social order. But in its enthusiasm for the new principle it overlooked not only
its derived nature and the limitations to which it is subject; it also set up as a self-evident and original right what
has only gradually worked its way forward in history and can always be realized only approximately through historical
development.

Unfreedom arises through the division of labor. One of the simplest forms of division of labor is this, that the
stronger loads all disagreeable work onto the shoulders of the weaker. Women and children were the first slaves. The
conquered enemy is made a slave when the cannibal instincts do not lead to his being eaten at once. Economic interest
works where ethical motives cannot yet work. While a roving people of hunters will often be obliged to get rid of the
conquered, a nomadic people or an agricultural tribe will find its advantage in making use of their labor. What is first
founded on egoistic motives and marks an economic advance becomes an ethical advance, inasmuch as the life together and
the care for others that egoistic and economic interest motivates develop sympathy, at any rate to a certain degree. It is
in the master's own interest that the slave's health and strength do not suffer. This relation comes forward especially
where serfdom or villeinage, not personal slavery, is found. This too is due to egoism and economic motives. Serfdom can
arise partly through conquest, partly through a transition from personal slavery to being bound to the soil, partly
through free men in dangerous times placing themselves under the protection of the more powerful, and this relation of
protection then --- perhaps through arbitrary compulsion on the protector's part --- passing over into being bound to the
estate. Its development presupposes that great tracts of untilled land lie waiting which can be made productive only in
this way\footnote[1]{The motive for the introduction of the Stavnsbaand [the binding of peasants to their estates] in
Denmark was the wish to prevent the peasants from moving from the infertile to the fertile districts, whereby many
estates would come to lack labor. V. \emph{Falbe Hansen}: \textit{Stavnsbaandsløsningen og Landboreformerne.} I, p. 6.}.
The
% --- p. 348 ---
\opage{348}serf has a greater personal scope than the slave, more motives for work and self-development. And the
relation between lord of the manor and serf need \textit{not merely} be a relation of power. The lord has the duty of
protecting his subjects, and he stands as an intermediate link between them and the supreme power of the state. Society
exhibits a series of stages, where the higher stage protects the material and spiritual well-being of the lower, while
the lower follows the higher with faithfulness and trust in struggle and work for material and spiritual ends. The
division of labor now becomes this, that some are to do the work of the sword and of the spirit, protect society and
think for it, while others in return are to carry out the material work. Such a division of labor was already found ---
with personal slavery as its foundation --- among the Greeks, and arose also --- with serfdom as its foundation --- in the
Christian Middle Ages, except that here the work of the sword and of the spirit was carried on by different estates. When
the knight ceased to fight, and when others than the priest began to think, the ethical justification of this social
order was gone. The principle of freedom had already been pronounced by Stoic philosophy\footnote[1]{The first opponent
of slavery on principle seems, however, to have been \emph{Antisthenes}, the founder of the Cynic school. He arrived at his
result chiefly by two paths. Man must limit his needs in order to preserve his spiritual freedom: but then he must also
be able to do without slaves! And spiritual freedom itself can be found in the slave; true slavery is dependence on
sensual lusts; the outward difference between slave and master has no essential significance. It is possibly the Cynics
above all whom Aristotle is attacking in his Politics. (\emph{Joël}: \textit{Der echte und der xenophontische Sokrates.}
II, 2. p. 569---574. Cf. I, p. 406 f.). The Cynic line of thought was continued by the Stoics, especially after Carneades
had shown that slavery conflicted with Stoic ethics. (Cf. \emph{Schmekel}: \textit{Die Philosophie der mittleren Stoa.}
p. 228. 378 f.). In a way analogous to Cynicism and Stoicism, the Vedanta doctrine and Buddhism worked in India.
(\emph{Deussen}: \textit{Geschichte der Philosophie.} I, 1. p. 163).} and by Christianity, without their drawing the
practical, social consequences of it (which again hung together with their regarding freedom only as a goal, not also as
a means). Now came the time when it could be proclaimed as a social principle. But that this could happen in a practical
way is due to the circumstance that freedom
% --- p. 349 ---
\opage{349}had long been stirring. In the domain of industry and trade men had been at work who were neither serfs nor
lords of the manor. It was their experiences, and the trust in free powers won through these experiences, that brought
about the great adherence to the principle of freedom. In the great series of emancipations that begins in the second
half of the 18th century, economic and ethical motives work together, just as at earlier stages of development.

Only at crude stages of human activity can slaves and serfs suffice. They do not have the great interest in the work
that free men have. From one of the districts in Denmark where the peasants stood highest in the 18th century, and where
the conditions of compulsory labor were best regulated, it was stated that ten laborers doing compulsory service do not
accomplish as much as two hired laborers. Slaves and compulsory laborers gain nothing by doing more than is just barely
demanded of them; that would only lead to the demands being raised for the future. They have no reason to show special
care for tools and material that do not belong to them. What is wasted is a loss only to the master. The emancipations
have therefore been every bit as much to the advantage of the employers as of the workers, and the emancipation of the
peasants has brought about a greater productivity of the land\footnote[1]{Cf. \emph{Roscher}: \textit{Die Grundlagen der
Nationalökonomie} § 71. --- v. d. \emph{Goltz}: \textit{Landwirthschaft.} (Schönberg's Handbuch. 1st ed. I). p. 580. ---
In a letter from the son of a former slaveholder (printed in \emph{Henry George}: \textit{Social Problems.} Ch. 15) it
says: ``The planters are pleased with the change. They say: `How stupid it was of us to go to war for the sake of
slavery! We get the work cheaper now than when we owned the slaves.' Why do they get it cheaper? Because in the form of
rents they take more of the Negro's labor than they could under slavery; for then they were obliged to give him
sufficient food, clothing and medical attention to keep him well, and they were compelled by their conscience and by
public opinion as well as by law to support him when he could no longer work. Now their interest and their responsibility
end when they have got all the profit out of him they can.'' % translated from Høffding's Danish
--- \emph{Falbe Hansen}: \textit{Stavnsbaandsløsningen og Landboreformerne.} I, p. 145: ``The fixing and commutation of
compulsory labor, the redistribution of land, the commutation of dues in kind, etc., were reforms that were just as much
in the landowners' interest as in the peasant's. They brought it about that the yield of the land rose and that its value
increased; but this was bound to benefit chiefly its owner, the landowner.'' (See also p. 68. 101. 149). It has gone the
same way in the industrial domain; in particular the shortening of the working day has also benefited the factory owners
(see below XXVI, 14).}.

% --- p. 350 ---
\opage{350}On the other hand, --- where the work was of such a nature that it exercised the workers' strength and
intelligence, there both the need for and the necessity of personal freedom came forward. Artisans have therefore as a
rule attained it before agricultural laborers. Spiritual independence is awakened by the feeling of self-activity and
trust in one's own powers. Free striving is carried over from one domain to another. If physical work goes best with
freedom, spiritual work must also go best with freedom.

But a certain spiritual development is necessary everywhere for the need for freedom to be felt. Unfreedom is felt as an
evil only by him who strives beyond the barriers set to him. The slave who has a mild master can often get his material
needs satisfied far better and more securely than the free man. The slave often feels himself a child and wishes for
freedom just as little as the child wishes to leave its home. He is kept away from every opportunity of making
comparisons\footnote[1]{In the North American slave states it was forbidden under severe penalty to teach a slave to read
and write, let alone to bring in books dealing with the emancipation of the slaves. \emph{Lyell}: \textit{Reisen in
Nordamerika.} Deutsche Uebers. p. 118. 120. --- Lyell also cites examples of how glad and proud Negro slaves could be of
their position. ``I confess,'' says the famous naturalist, ``that one may think and philosophize about this peculiar and
interesting kind of vanity until one finds in it the proof of the utmost social degradation; but the first impression it
made on me was very consoling, so that I could not possibly feel a painful pity for people who felt so extraordinarily
contented.'' % translated from Høffding's Danish
In a similar way one sometimes finds, in countries where polygamy prevails, that a woman is proud of her husband's many
wives: it is a sign that it is a mighty and rich man she is married to.}, and new wishes and impulses therefore do not so
easily arise in him. (Cf. VII, 3). Since the need for freedom is lacking, unfreedom is not felt as an evil. The thrall's
position at earlier stages of culture approaches the child's, and need not be bound up with
% --- p. 351 ---
\opage{351}the degradation found in this position at a later and higher stage of culture. On the contrary, especially
when the whole order is regarded as an order of nature, as something that cannot be otherwise, and when common experiences
in good fortune and ill come to work, a hearty relation can develop between the thrall and his master, as between
comrades and fellow combatants. Over against such a relation the proclamation of the principle of freedom would have a
dissolving effect. Nowhere does it show itself as here how the value of general principles is conditioned in each single
case by the special circumstances to which they are to be applied. Whether the state of unfreedom has that idyllic and
hearty character that was suggested, or becomes a state of oppression in which power and dominion are used not for
protection but to satisfy caprice, lust of power and lust of enjoyment, a long state of unfreedom --- for the individual
and for the race --- will not all at once be able to be replaced by a state in which freedom can express itself both as
means and as goal.

Since development is checked as long as unfreedom lasts, both impulse and means for it being lacking, one cannot expect
the mere proclamation of freedom at once to bring about an entirely new state. The first effects of freedom are seldom
wholly fortunate. There is, after all, a transition to entirely new conditions of life to which the individual's nature
has not yet been able to adapt itself. The freedmen had a bad reputation in Rome. Neither the Negro slaves nor the
Russian peasants could at once use their freedom in the right way. Of the effect of the agrarian reforms on the Danish
peasants it is said, on the basis of a historical investigation\footnote[1]{Falbe Hansen. p. 87.}: ``The most essential
significance of the reforms for the nation's economy was that which came about through the raising of the moral and
intellectual level of the peasant class. The agrarian reforms certainly also worked directly to improve agriculture\dots{}
but still greater, no doubt, was the mediate, the indirect influence that came about through the educating effect of the
reforms on the peasant class and on the population as a whole. It is evident, however, that this indirect effect could
% --- p. 352 ---
\opage{352}show itself only late. A generation of downtrodden, thrall-bound peasants is not changed all at once because
the mere declaration of freedom is issued; a wholly new generation is needed to see the effect of it show itself in
farming. And as regards, e.g., compulsory labor, one of its worst drawbacks, that the peasant grew used to working
sluggishly and the landlord to squandering labor, could disappear only late; it has even been thought that the
after-effects of compulsory labor could be traced in farmers' way of working more than half a century after its
abolition.'' --- Human nature, after all, does not change at one stroke. Only gradually, through transitional states that
perhaps last for several generations, can maturity be reached. And fully reached it can be only through the very use of
free powers. What matters is to find provisional forms within which self-activity can move until it is able to take
upon itself the greater tasks.

4. In proclaiming the rights of man, form was set above content; indeed, it was even thought that form could do the
whole. It was overlooked that the degree to which freedom can be realized and carried through at a given time depends on
a whole series of different social conditions. Within the old social order the boundaries were narrowly drawn, and there
were heavy burdens to bear. But those among the privileged who regarded their power as an office entrusted to them showed
a care and a fatherly feeling which brought it about that the dependents could never feel wholly forsaken. Now, on the
other hand, when duty fell away together with power, the emancipated individual was thrust out into the world with its
charter of freedom in its hand, without perhaps knowing what it was really to use its right of man and its freedom for. A
relation of protection and reverence was suddenly transformed into a legal relation. This came forward not least in the
abolition of the old guilds and the institutions attached to them. Every organization of labor was thereby abolished,
and individuals were left to make their way with their own powers. --- Characteristically enough, it was the same time
in which a sharp separation was made between ethics, political economy and jurisprudence, and it was thought that here
one had three wholly separate domains, each by itself.
% --- p. 353 ---
\opage{353}In political economy only the business of earning was to be considered; if only it were given free play,
men's economic interests would come into harmony of themselves. In jurisprudence only the conditions for such an outer,
mechanical order were to be sought as would make security and freedom possible for all; it was to keep to outer action
(legality) and wholly disregard disposition (morality).

At the same time the starting point was a view that regarded outer circumstances, upbringing and social conditions as
the sole cause of all differences between human beings. All human beings --- it was held --- were equal in nature and
gifts; it was only the outer circumstances under which they developed that brought about the differences. \textsc{Adam
Smith}, whose views gained so great an influence on the general conception of social questions, thus proceeded from the
assumption that differences in character and gifts do not derive so much from nature as from habit, way of life and
upbringing, are not so much causes of the division of labor in society as effects of it\footnote[1]{\textit{On the Nature
and Causes of the Wealth of Nations.} Danish transl. I, p. 25. --- Smith had this doctrine from \emph{Helvetius}.
(\textit{De l'esprit.} III, chap. 26---27). --- \emph{James} and \emph{Stuart Mill} still held it.}. When, then, all barriers were opened and all
compulsion kept away, there seemed to be a prospect that all could get ahead. Justice must require that the conditions
be equal for all, and nature too, he held, gave no one a head start over another.

One was led by this one step further. All free associations were even forbidden. Or rather, after what had been
experienced of the degeneration of the old guilds, one could not conceive of any association without compulsion and
tyranny. \textsc{Turgot} had the guilds abolished by an edict of 1776. He saw (as is stated in the preamble to this edict)
the source of all evils in the industrial domain in the right that artisans of the same trade had had to assemble and
unite in associations. Everyone now got the right to practice whatever trade or whatever commerce he wished; but it was
forbidden to all masters, journeymen
% --- p. 354 ---
\opage{354}and workers to form societies and associations under any pretext whatever. When Turgot was overthrown, this
reform was repealed, like all his other reforms; but the Revolution took it up again. A law of 14 July 1791, % PRINTED AS IS: Juli; the note says June (the law is of 14 June 1791)
that is,
from the first period of the Revolution, declared the annihilation of all corporations of those who had the same estate
and occupation to be one of the foundations of the French constitution, and forbade the restoration of such societies in
any form whatever. One was not allowed to assemble and deliberate in definite forms on one's ``pretended common
interests'' (interêts prétendus communs)!\footnote[1]{\emph{Léon Say}: \textit{Turgot.} Paris 1887. p. 150---159. ---
\emph{Taine}: \textit{La Révolution.} I. p. 221 f. --- It is surely due to an agitator's use of history when \emph{Karl Marx}
\textit{(Das Kapital.} 2. Udg. I. p. 772) presents the law of 14 June 1791 as an attack by the bourgeoisie on the
workers' freedom.} This step is characteristic of the abstract and overstrained conception of the principle of freedom.
Associations are feared because they can easily produce new differences, new peculiar social groups, and in one's zeal to
protect freedom against this danger one goes so far as to introduce one of the very worst restrictions of freedom by
forbidding free union with others. Indeed, the power of the state is even granted the authority to decide whether one has
interests in common with others at all! The French law mentioned is a telling example of the tendency to keep only the
danger lying nearest in view, and to show so great a zeal in averting it that great obstacles are laid in the way of
continued development. In the 18th century it was thought that everything was done when barriers and bonds were
abolished. It was not considered that because the old guilds had had harmful effects, it did not follow that all
associations were harmful. One got entangled in the great contradiction of forbidding, in the name of freedom, one of the
most important uses that can be made of freedom: entering into a free social relation with other individuals.

But this overstrained character does not prevent it from being an extraordinarily significant principle that was
proclaimed.

% --- p. 355 ---
\opage{355}What matters is that freedom become determinative as far as possible in the division of labor, so that this
does not take place, as at the primitive stages, through compulsion and in such a way that the disagreeable work is loaded
onto the weak. The individual must be able to follow his impulse and capacity, so far as he is conscious of them, and so
far as he does not thereby violate greater interests. Freedom ought to be the possibility of binding oneself to something
definite, of choosing a calling. In the old social order social differences prevailed that had almost the character of
caste differences, and each individual's place was assigned him from the first. In the freely chosen calling, on the
other hand, the individual determines his own place, the point from which he will and can work for the tasks of the race.

A free formation of society thereby becomes possible. All thoughts and ends must first arise in the consciousness of
single human beings before they have their effect in the rest of the world. They work to form societies by gathering
around them those who understand their significance; and when the thought has thus been tested and applied in smaller
circles, the time comes when the firmer organization of the state, resting in the last instance on means of compulsion,
can step in. In free work on culture the productive powers unfold; it is then the state's business, with its legal order
and its means of power, to protect where protection is possible and beneficial. The state's relation to the free cultural
communities is thus analogous to its relation to the family. In neither place does it produce the really operative
powers; in both places it stands over against something that is not of the world in which it itself belongs; --- but in
both places its activity can be necessary to organize and to help. And it is by no means so that organization is first
the state's business. Viable organizations must, as history shows, grow up through the union of free powers which, under
the influence of the conditions and interests of life, find one another and more or less consciously supplement one
another or work toward the same goal. The forms these organizations take can become perfect only by being tested against
real circumstances. Some of the most valuable experiences the human race has made stem from this involuntary or voluntary
adaptation of the forms of life together to the actual
% --- p. 356 ---
\opage{356}circumstances. They are experiences whose effects concern not only the outer forms, but branch out into mind
and thought. There is a constant interaction between the inner and the outer. One of the most important problems in
social ethics is how far social development can be left to free adaptation, of which free association is one of the most
significant forms, and where the point lies at which the state must intervene with its compelling authority. When the
Revolution proclaimed not only liberty and equality but also fraternity, it declared that mere emancipation is not
enough, but that positive union is needed. But the relation between these two sides has proved to be far more involved
than was seen in the first enthusiasm, and the social problems come forward precisely when one thinks over the relation
between the different members of the revolutionary formula. They present themselves, however, naturally in somewhat
different forms in the different domains of culture.

5. There are just as many different ends of culture as there are different domains where the struggle for life is
waged. Life rests not merely on physical self-preservation, the procuring of the material conditions of existence; at
higher stages the value of life is tied to ideal ends, the satisfaction of the needs of thought, imagination and feeling;
and there is besides a need to let material and spiritual goods fall to the lot of as many as possible. To these three
kinds of culture correspond three kinds of free cultural community.

\textit{Material culture} springs directly from self-preservation. --- If there were not the constant pressure on human
beings that ``the will to live'' causes, they would scarcely submit to the restless labor and the toilsome drudgery in
the service of material culture. This pressure has driven human beings forward from lower to higher stages even before
ideal motives could assert themselves. --- At the lowest stage of human existence nourishment consists of wild animals
and plants that are sought out. Perhaps even fire is not known. Only by the faculty of speech and by the use of simple
tools is man here distinguished from the animal. This stage has been called the \textit{savage} stage proper. --- A
somewhat higher stage is the \textit{barbarian,} where fire is known,
% --- p. 357 ---
\opage{357}and where one can make metal tools and practice agriculture and cattle breeding. --- \textit{Civilization}
proper begins when written language has been invented, and with it a sure and widespread memory and tradition in the race
becomes possible\footnote[1]{Cf. \emph{Tylor}: \textit{Anthropology.} p. 24. 179 f.}. Written language no doubt first
served purely practical ends, such as accounts, lists, messages. But through a displacement of value this means of
preserving and communicating experiences and ideas acquired a significance that completed the transition from material
to ideal culture which had already begun where there was interest in oral tradition. --- It is of great ethical interest
that the more material culture develops, the more it makes it possible to bring about and strengthen connection and
intercourse between human beings, and that it at the same time approaches ideal culture and leads over to it. In these
respects the standard must be sought for the development of material culture, and especially for the perfection of the
forms under which it develops. The more material culture has a uniting effect, and the more it leads over to ideal
culture, the higher it stands.

\textit{Ideal culture} comes into being as soon as ends arise that aim at something other and more than maintaining
life. Even if it is the same powers that are at work here and in material culture, they are here used for their own sake,
on account of the immediate satisfaction bound up with their use. This use of the powers for the sake of the activity
itself is a sign that a truly human stage of life has been reached. That it is reached is due to a displacement of motive
(see XIII, 4). Not only the ``necessary'' but also the ``beautiful'' (to use the Greek expressions)\footnote[2]{By the
``necessary'' the Greeks understood what is a means for something else, by the ``beautiful'' what is in itself an end.
Cf. \emph{Aristotle}: \textit{Rhet.} I, 9.} now becomes an object of interest and endeavor. Man now strives not merely in order to live,
but lives in order to strive. Now science and art arise, an aesthetic and a religious life of feeling.

\textit{Philanthropic culture} has for its object the satisfaction of the love
% --- p. 358 ---
\opage{358}of mankind. It aims especially at supporting those who struggle at the lowest stages of material and ideal
culture. It will remedy bodily and spiritual degradation where it is found. It will seek especially to prevent a
dualistic difference from developing between material and ideal culture. But it is directed against all suffering and
want, bodily and spiritual, wherever found. A helping hand may often be needed precisely where there is plenty of
material and spiritual goods. --- Philanthropic culture need not come forward as a special endeavor beside the two other
kinds of culture, but can be bound up with them and mark the spirit and direction in which they are cultivated. But
hidden or open it must be present if development is to be sound and vigorous.

\grouphead{\textit{1. MATERIAL CULTURE.}}

\chapterhead{XXIV.}{SOCIAL ANTAGONISMS.}

1. There is first talk of material culture when man must work in order to maintain life. Where he finds the means of
self-preservation in nature, e.g. fruits that thrive without cultivation, he remains at a half-animal stage. But when
material culture has run through its first stages, there comes a point at which the activities that aim at procuring the
first and most indispensable conditions of life are looked down upon with contempt. This hangs together with the division
of labor that sets in early. The strong choose for themselves the work they wish to carry out and leave the rest to the
weak. At the barbarian stage it is hunting and war that alone are considered worthy of a man.
% --- p. 359 ---
\opage{359}With more advanced culture, spiritual activity is set beside these pursuits. A dualism now asserts itself
between one circle of activities, chosen for their own sake and able to fill and develop personal life, and another
circle of activities, forced on men by power or by want, which make the single personality a mere means for others. It is
then considered a degradation to work for the necessary. This line of thought can be traced through classical antiquity
and the Middle Ages. It was partly the slight connection of mechanical work with personality, partly the dependence that
arises where one works for wages, that made handicraft be considered unworthy of a free man: ``Nothing noble,'' says
\textsc{Cicero}, ``can thrive in the workshop'' (nec enim qvicqvam ingenuum potest habere officina)\footnote[1]{\emph{Cicero}:
\textit{De officiis.} I, 150. --- Cicero follows here, as usual, Panaetius, who on this point was more of an
aristocrat than earlier and later Stoics. Cf. \emph{Bonhöffer}: \textit{Die Ethik des Stoikers Epictet.} Stuttgart 1894.
p. 73 f. 233---243.}. Within the Stoic school the dignity of work was maintained, inasmuch as the wise man's disposition
ennobles every human activity of life. And through Christianity this line of thought spread in wide circles. But medieval
society, built on aristocracy and hierarchy, came away from this conception again. As late as the year 1781 the Academy in
Madrid set the prize question: ``to show that the useful trades had nothing dishonorable about them.''

It is partly aesthetic and religious, partly ethical considerations on which this long-prevailing conception has relied.
--- For the aesthetic idealism of the Greeks and the religious idealism of the Middle Ages, work with material stuff was
degrading and drew one downward. It prevents man from making his personality a harmonious work of art, or from dying away
from the world of the senses. --- And since material production serves first and foremost the individual's own
self-preservation, it seemed to rest on egoism. Modern political economy has often favored this conception by strongly
stressing the egoistic motives for the business of earning.

% --- p. 360 ---
\opage{360}When in modern times a greater acknowledgment of material work develops, this hangs together with thoughts
that are peculiar to the modern view of life. --- This view is convinced of the close connection between spiritual and
material life. However one conceives the closer relation between spirit and matter, it is settled that the life of the
spirit is tied to the highest and finest form of material life. The brain of a highly gifted human being is the finest
material product we know. And material work is needed to produce as good brains as possible, and work for material
culture can thereby also become work for spiritual culture. --- Material work, further, comes more and more to stand as
the application of laws discovered by the work of thought. We work upon and master nature by virtue of the superiority
over it that natural science gives us. A bridge has thereby been built between the world of the hand and that of the
spirit which antiquity and the Middle Ages did not know. Cicero distinguished between labor (opera) and art (ars); the
former was unworthy, the latter worthy, of a free man. But for many kinds of work this sharp opposition gradually falls
away. If hunting, war and the work of thought stood as noble activities because personal powers could stir freely in
them, the possibility is now not excluded that work with material stuffs too can develop personality, the more it lays
claim to understanding and will and is not merely a mechanical activity imposed from outside. Finally, we have learned to
see that the individual's work, however insignificant it might seem, nevertheless plays its part in the great social
household. He thereby procures not only the means of maintaining his own life, but also contributes his mite to the life
of the race, increases the goods at society's disposal. His energy and his thrift can --- on account of the solidarity in
the individuals' struggle for existence --- benefit more people than himself. He can therefore also carry out his work in
a frame of mind directed toward a wider horizon than the one determined by his purely individual needs. Science has
taught us that nature, both in the spiritual and in the material domain, often reaches
% --- p. 361 ---
\opage{361}great results by the summing up of small activities. By being faithful in little the individual can work for
the great whole.% cf. Luke 16:10

2. But this altered view of work has posed a problem far more than it has reached a real result. The higher valuation of
material work presupposes an ideal relation between worker and work which is still far from present, but which is the
content of the demand that must be made on social development from the side of ethics. This demand furnishes the
standard for the valuation of the development of material culture at any given time. Before we go further into this,
however, we shall first bring forward the social antagonisms which testify that the primitive division of labor and its
evils are not yet overcome.

a. The antagonism between \textit{propertied} and \textit{merely working} individuals is a historical continuation of
the primitive division of labor between master and thrall. The propertied man often merely enjoys the fruit of past work
or of the work of his contemporaries; in that case there is a complete antagonism between him and the worker. But even
if he works, the antagonism is sharp enough. For he disposes both of the means of work and of the finished product. The
propertyless worker, on the other hand, who has only his two hands, must put them at the propertied man's disposal in
order to get material to work on and a share in the yield that is won. Without material no work can be carried out; but
even if labor cannot be had, the material can perhaps still be enjoyed and consumed by its owner. However indispensable
an intermediate link work is in the transition from raw material to product, it is nevertheless regarded from the
standpoint of production and exchange merely as a means. Wages are regarded essentially as an expense that must be
reduced to the least possible. From his standpoint the worker must naturally regard wages as an end. Their size marks his
social position, and in a way also his position as a human being. But he lacks the power to make this view prevail as
long as he is left to himself. Just as his work is essentially regarded as a mere means for the process of production, so
his person too is
% --- p. 362 ---
\opage{362}easily regarded as a mere means. A numerous working population is necessary in order that labor can be had
at as cheap a price as possible. For some to content themselves with the least possible wage, there must be others who
get no wage at all. The consistent carrying through of the view that wages are merely an expense thus leads to a mass of
human beings being pressed down to the lowest stage of existence.

This dependence, in which the worker, despite the personal freedom he officially possesses, comes to stand toward the
propertied man, extends further than to the conditions of his work and maintenance. Not only the wage, but also the
working time and the conditions under which the work is carried out, are determined predominantly without his will. It
sounds like mockery when he is told that he is a free man, and that it depends on himself whether he will take on a piece
of work or not\footnote[1]{When in an English government commission a demand was made for more exact inspection of the
coal mines on account of the frequent accidents, the employers' representative said: ``Does it not depend on the miners
themselves whether they will go down into the pits?'' A witness answered: ``Certainly. But it also depends on themselves
to starve to death if they do not go down.'' % translated from Høffding's Danish
L. \emph{Brentano}: \textit{Die Arbeitergilden der Gegenwart.} II, p. 17. --- Cf. p. 165 on an employer's joy over the
many unemployed workers who make wages cheap.}. The ware he has to offer --- his labor power --- is inseparably bound up
with his person, so that in selling it on terms he does not control he also sells his person. Dependence on the employers
easily extends from the outer, material domain to the inner, so that in social, political and religious questions too he
cannot act as a free man.

b. Very few workers have yet been able to benefit from what science has accomplished in the service of industry. The work
of thought and of muscle, \textit{intelligent and physical work,} are not yet united, but are for the most part the
business of different individuals. The application of the discoveries of science, the laying out of the plan of work and
the procuring of the means to
% --- p. 363 ---
\opage{363}carry it out require capacities and conditions that the physical workers do not command. Physical work is
carried on as good as blindly, without insight into the forces of nature it works with or into the significance of the
results reached. --- \textit{Within the physical workers' own circle} too a division of labor is carried through more and
more, which in the end restricts the single worker to the mechanical repetition of one and the same quite simple activity.
In the development of work there takes place partly a dissolution and parceling out of a compound work among several
different workers, partly a combination of several different branches of work for the production of one and the same
object. These are the same two forms of development --- isolation and combination --- that appear in all domains of human
activity. The manufacture of a pin was already in Adam Smith's time divided among 18 hands. In the year 1860 there were
only four workers to make a chair; but in the year 1895, 23 workers were needed to make the same kind of chair with the
help of machines. In such cooperation, whether it arises through parceling out or through combination, the single worker
loses the overview. He is trained in a single simple activity, and only a single side of his personality and his
capacities is developed. Human life is not developed in him in its full extent; seen from this side too he comes to stand
as a means, not as an end. The more mechanically he can carry on his task, the better it perhaps fits into the whole
enterprise. His one-sidedness makes him an all the better member of the whole machinery. And one-sidedness also makes him
more dependent, since he is not able to take on new, more compound work. To be sure, parceling out makes it easy to pass
from one simple work to another; but this advantage is outweighed by the fact that practically everyone can take on such
simple activities, so that competition becomes great and wages low. --- It then becomes of little use to the worker to
train his mind and his capacities, since the activity that falls to his lot requires no further training. Without joy in
his activity, without hope
% --- p. 364 ---
\opage{364}of advancement and without an impulse toward development, he easily sinks down to a half-animal existence.
---

The antagonisms that have here been portrayed in their strongest forms\footnote[1]{\emph{Rousseau} first brought forward with
energy the unfortunate sides of the division of labor and thereby posed the social question. Cf. my book \textit{J. J.
Rousseau og hans Filosofi.} p. 99---106; 125 f.; 130 f. --- Shortly after Rousseau, \emph{Adam Ferguson} \textit{(Essay on
the History of Civil Society.} Edinbourgh 1767. IV, 2: Of the subordination consequent to the separation of art and
professions) illuminated the same question from a culture-historical point of view. --- In his work \textit{Gemeinschaft
und Gesellschaft} (Leipzig 1887, 2nd ed. 1912), built on thorough sociological and psychological studies, \emph{Ferdinand
Tönnies} has given an instructive account of social pathology. Cf. my review of this work in the essay \textit{Social
Pessimisme} in ``Mindre Arbejder'' I.}, condition the social question.

\chapterhead{XXV.}{THE SOCIAL QUESTION.}

1. When one speaks of the social question, the misunderstanding must be kept away that it is a wholly single and simple
question, and that it could be solved at one time and once for all. It arises from the strong antagonisms that threaten
to split society; but these antagonisms are in turn conditioned by the interaction of many different circumstances. It is
not only material work whose ordering is in question here; ideal culture too, the development of spiritual life, becomes
significant; and finally the constitution of the state too will have an essential influence both on the shape of the
question and on its solution. There is just as close a connection between these different sides of the social question as
there is between brain, heart and bowel. Disease in a single one of these
% --- p. 365 ---
\opage{365}organs can cause the death of the whole body, and a fault in one organ will have consequences for the activity
of the others, which then in turn act back upon it. --- A solution of the question once for all is improbable, since it
hangs together with so many circumstances, and since, even if it comes forward most clearly and sharply in our time, it
has nevertheless existed in one form or another at all times.

The reason why it is only in our days that people speak of the social question must be sought in several circumstances.
It is hardly because there is more egoism and more envy in the world than there has been before. Nor is it because want
and suffering in themselves should be greater nowadays than formerly; rather the opposite might be asserted. One may
perhaps even say that if conditions had not improved, there would be no talk of a social question. For the greatest want
strikes to the ground, crushes thought and does not awaken the impulse to get further. It is to that extent precisely an
advance in the workers' circumstances that makes it possible for the social question to be posed by themselves. It is a
sign that their nature is after all not completely crushed. Conditions must be said to have grown worse only insofar as
the one-sidedness and dependence brought about by the ever further carried division of labor must be felt all the more
strongly the more the worker, by virtue of the principles of the 18th century, now stands as a free man. As long as he was
only a slave or servant, he did not feel these evils so much, especially since the old social order, besides the
dependence in which it held him, afforded him protection and help in many ways. To this is added, further, increased
enlightenment. The power of thought is awakened, and comparisons are made. One no longer bows before the given social
order and the difference in lots of life as before something that should be justified merely because it is. One asks
about its justification, and one sets up ideals. And the greater the opposition between ideal and reality, the more
strongly the imagination is set in motion.

But it is not only awakening comparison and reflection that pose the social question so sharply. It is, after all, posed
not only by those who suffer directly under the conditions; others too feel its sting. Sympathy and the feeling of justice
% --- p. 366 ---
\opage{366}have developed to greater fineness and greater extent. The unrest complained of in the present comes to a
great extent from sympathy. If one could be more egoistic and thoughtless, one would live more happily for one's own
person. The pessimism that is so common in our time springs in part, no doubt, from blasé indifference; but its noblest
source is lively sympathy with the great suffering and disharmony found in the world, which earlier times passed over
more lightly.

2. At the same time as the consciousness of the significance of the social question has been sharpened, and the impulse
to work at its solution has been awakened, we have also learned how deep a foundation it has in human nature and its
conditions. The childish philosophy of history that derives all great movements from the appearance of single individuals
we are getting beyond; it runs aground already on this, that it remains incomprehensible how the single individuals
themselves came to think and act as they did. Among conservatives and radicals alike, however, this childish conception
can still be traced. According to conservative optimism everything would be good if only there were no radical
agitators; according to radical optimism everything would be good if only there were no kings and priests. Both overlook
partly that in the social question we have a historical inheritance, a continuation of the social difficulties of earlier
times, partly that in human nature and the conditions of human life there lie causes that will perhaps always lead to the
social question's being posed, even if in new forms.

It was mentioned in the foregoing that the great numbers of the working population were a cause of labor's taking so
relatively subordinate a place among the factors of production. If there were not so many workers, the workers'
dependence on those who dispose of material and product would not become so great. It is the strong increase of
population that causes the struggle for existence in the social domain as in the whole of nature. In all domains of
organic life the species are maintained by far more germs being thrown out than can develop under the given conditions.
In order that a few germs may take root and develop, manifold ones must perish. It
% --- p. 367 ---
\opage{367}seems as if life can subsist only when a mighty force is put into producing a surplus of individuals. The
human race too has a tendency to propagate itself more quickly than the means of subsistence can always keep up with.
This proposition, which was put forward with exaggeration and one-sidedness by \textsc{Malthus} \textit{(Essay on
population.} 1798) against the prevailing optimism that derived all social evils from human institutions, has drawn
attention to an essential side of the social question. Social life now proved to hang together with natural life as a
whole. \textsc{Darwin} rediscovered throughout plant and animal life the conditions Malthus had assumed in the case of
human life. We do not understand the social question when we do not see it in its connection with the whole life of
nature. When forces and tendencies lying so deep are at work, one will not so easily think that the question could be
solved in one single way, by one single formula and for all time.

Malthus' one-sidedness consisted, in the first place --- as \textsc{Auguste Comte} already remarked\footnote[1]{\textit{Cours
de philosophie positive.} IV. 2. éd. p. 455.} --- in his not stressing that a numerous population is necessary for social
life to be able to unfold in its various forms. Division of labor, lively exchange, social organization are possible only
with a numerous population. The more numerous population is also a denser population, and thereby life together is
furthered, and everything to which it in turn leads. Against outer dangers too numbers can be of importance. --- In the
second place, Malthus' one-sidedness consisted in his not seeing that the disproportion between the size of the
population and the means of subsistence available may also be due to an imperfect development of human activity. Greater
energy or cleverness might perhaps be able to find use for the apparently superfluous labor power, either with the help of
new inventions or by a better use of capital. Or new and more easily accessible means of subsistence might be procured.
The disproportion can thus have different causes in the single cases, now a defect in production, now a defect in
distribution, now a defect in consumption.
% --- p. 368 ---
\opage{368}--- But what Malthus is right about is that each time a certain harmony between population and means of
subsistence has been produced, the strong tendency of population to multiply will always anew upset the balance and make
new exertions necessary in order to preserve the stage of life that has been reached. It is this that always anew drives
human beings onto the path of work and culture. Malthus says, and one need not be a pessimist to agree with him in this,
that if the means of subsistence increased of themselves in the same degree as the human race, he does not see what
mainspring could be strong enough to overcome man's indolence and carry him further on the path of
culture\footnote[1]{\textit{Versuch über die Bedingungen und die Folgen der Volksvermehrung.} Aus dem Engl. von
\emph{Hegewisch}. Altona 1807. II. p. 152.}. Sloth, in which we have found the germ of all evil (VI, 1), needs a strong
counterweight to be overcome. As long as higher motives have no greater power in human nature than has so far been the
case, the elementary motive stressed by Malthus will always anew be necessary. A strong pressure is needed to awaken
impulse and power. Only when it has had its effect is a stage reached where more ideal motives can work.

There is no reason to feel admiration for the order of nature we stand before here. But it must be taken as something
given. And on the other hand it must not be forgotten that the mighty force that leads to the posing of the social
question is the same that, by leading to the founding of the family, lays the foundation of the most intimate human
society, the source of all sympathy in the human race, and the same that, through its influence on the imagination and by
awakening enthusiasm, sets the capacity for devotion and sacrifice in activity.

3. The social question is an \textit{ethical} question. This lies already in what was said above, that it is only on
account of awakened sympathy that it comes forward so sharply and so pointedly in recent times. Think away the standard
of ethics, the ethical thought of an ideal society (III, 10; XIII, 6) --- and the ``question'' will at once fall away, or
rather, it will be reduced to a pure question of power; the individual will merely be set
% --- p. 369 ---
\opage{369}the task of saving his skin or feathering his nest as well as possible in the general struggle of all against
all. In the thought of the ideal society, of the kingdom of humanity, there lies, on the contrary, the demand that every
human being be more than a means, that he take his peculiar and independent place in the great realm of
mankind\footnote[1]{In my essay \textit{Forestillingen om Retfærdighed i dens historiske Hovedformer} (Nationaløkonomisk
Tidsskrift 1873) I have developed this thought and its significance for the social question.}. It conflicts with the
ideal of a human society when a greater or smaller number of human beings stand merely as a passive mass, as subordinate
means and tools, whose enjoyment and suffering are not taken into the reckoning when the account of social progress or
regress is to be made up. It must be well remembered here that enjoyment or suffering is conditioned not only by the goods
or evils that fall to the individuals' lot from outside, but also by the tasks they can set themselves and work at, and
the duties they are able to take on. What matters is not only the distribution of the yield, but also the distribution of
the work itself. The proletarian's striving to attain a higher level of life very often finds its justification in this,
that only when such a level is reached can he come to work at the tasks and fulfill the duties that correspond to his
capacities. So too woman's wish for ``emancipation'' often appeared as grounded in the need for active participation in
the great work of mankind. (Cf. XIX, 5).

As long as the concept of \textit{mass} can still be applied to human beings, the goal is not reached and more or less
serious disharmonies are present. Within a society, mass is what is not taken up into the organization, is no living and
active component of society. In the mass the single personalities are effaced and disappear and do not stand as
distinctive points of force. Where there is mass, personality does not come into its own. Neither the unity nor the
multiplicity that a perfect society presupposes can therefore be present: not unity, because society is split
% --- p. 370 ---
\opage{370}into parts that do not hang together organically, but perhaps even stand hostile to one another; --- not
multiplicity, because the independence of the single personalities disappears. The principle of the free personality,
the most important proposition that can be derived from the principle of welfare (VIII, 6), is violated, personal beings
being reduced to means without at the same time being regarded as ends.

4. But can it now be otherwise? --- There are those who answer this question with No. Their line of thought is roughly
the following. --- In social life, as in natural life, definite laws prevail that cannot be defied. Such a law is the law
of supply and demand. It is of no use to summon up the greatest exertion when the result one attains is not in demand;
for our work we get only what can be given in the market. Under this law those must suffer especially who come to market
with mere physical labor power. They will not be able to reach a secure material existence, and they will have no time
and strength left over to share in spiritual culture. This cannot be otherwise, and it has never been otherwise. On the
other hand it is necessary that material work be carried on. Society's material needs must be procured by the great mass,
in order that a small band can work at spiritual culture. Higher culture has at all times been accessible only to a small
circle, favored by fortune, which has been able to develop harmoniously and many-sidedly without needing to slave for
material needs. In order that the light may fall clearly on this elite of the human race, the great mass must lie in
darkness. It will not even be fortunate if too much light is spread from the narrow circle out into the great one. It is a
natural division of labor that some think, others work, and thinking will easily make work less good, or will awaken
discontent and unrest and thereby disturb the sound order. Nevertheless the chosen give the multitude more than they
receive from it, so that one cannot speak of heartlessness and inhumanity in them. And insofar as those who belong to the
great mass do not find the life they are bound to lead worth living, they must be referred to
% --- p. 371 ---
\opage{371}the faith of the Church and console themselves with the hope of another and better life, in which they can
find redress for what they have suffered here.

As a counterpart the following way of thinking is often put forward. --- The so-called social laws of nature appealed to
are only the expression of a social order which in part has even been brought about by violence and encroachment. That
some individuals have so great a head start in the great competition is due to an injustice which it must be possible to
correct. With what right does one really speak here of a social order? In the competition between human individuals there
really prevails only the animal struggle for existence, merely in a somewhat masked form. Present society is not at all
built on ethical and humane principles, but on power and egoism. It must be wholly dissolved. Bare ground must be cleared;
then we can always talk about what is to be built in its place. While material work now stands as a subordinate, serving
force, the insight must break through that work is the source of all wealth and culture. ``Over against the working
class,'' it says in a program issued by German Social Democrats (at Gotha, May 1875), ``all other classes are only one
reactionary mass.''

As you call into the forest, so it answers. From both sides ``Mass!'' is called out to the opponent. From both sides people
meet in the assertion that a sharp antagonism, a dualism, exists within society, only that from the one side it is held
to be beneficial or at any rate necessary, from the other side to be inhuman and easy to abolish if only the will is good.

5. According to the general view I have earlier maintained (III, 19---20; XIII, 1---4), ethical development is to carry
further what natural, more or less unconscious development has begun. The human will cannot create out of nothing, but can
only build further on the foundation that nature has laid. The point, then, is to find points of departure and germs of
the future in what is given. Where such are wholly lacking, the state of things is hopeless. But according to both of the
views just brought forward we end in hopelessness. For the first-mentioned view
% --- p. 372 ---
\opage{372}this is consistent; for since it considers dualism in the human race a normal state, there is for it no reason
to long to get beyond it. Over against the second view, on the other hand, the question must be put: if in present
society there is nothing whatever of value, whence then are the powers to be got for working on the new building of
society? The powers we use have always been procured by the past. It is therefore not easy to see how an absolute
condemnation of what history has handed down can be united with a great hope for the future.

There is, moreover, one thing that both parties overlook. The dualism is in reality not complete. When the social
question is posed more sharply in our time than before, it is, according to what we have seen, to a great extent because
consciousness of its significance has awakened. This consciousness is not confined to any single class or circle. It is,
as \textsc{Stuart Mill} has said, characteristic of our time that all classes take part in the debate on the important
questions, and that more than ever before the suffering classes have an important voice in it\footnote[1]{\textit{Principles
of political Economy.} II, 1, 2. --- \emph{Lassalle} maintained that his agitation inaugurated a work of reconciliation,
in that it gave the propertied and the working class the possibility of meeting in discussion of the social question.
\textit{(Arbeiterlesebuch.} p. 54 f.). --- In the part that Positivists and ``Christian Socialists'' have played in the
English workers' struggle for the rights of their class (see \emph{Sidney} and \emph{Beatrice Webb}: \textit{History of
Trade Unionism.} London 1894. p. 197, 228 f., 246 f.), there is also a testimony that the working class does not wage its
struggle without help and sympathy from other classes, something which the Marxists' constant talk of the class struggle
conceals. The socialist ideas themselves were, moreover, first put forward by men from the ``higher'' classes. --- That
the Positivists naturally work together with the proletarians is seen from what is cited in \textit{Den nyere Filosofis
Historie}\textsuperscript{2} II. p. 350 and 359.}. The two opposing parts of society are thus, after all, not wholly
isolated, but must have a common ground, since they can discuss the problem with one another. The spiritual unity of
society is thereby maintained, and at the same time an essential means of getting further is provided. For history shows
clearly that no reform and no improvement leads to the goal if it is not actively
% --- p. 373 ---
\opage{373}supported by those who are to be helped. They know best where the shoe pinches, and without their
cooperation the pressure cannot be relieved. Neither theoretically nor practically can the social question be treated
fruitfully without the workers themselves taking part in the discussion and the decision. And a beginning has now been
made with this. ---

But it is a long way from discussion to decision, and the consolation drawn merely from everyone's being able to join in
the discussion might rightly seem somewhat questionable. We must then examine what practical possibilities of development
present themselves.

\chapterhead{XXVI.}{POSSIBILITIES OF DEVELOPMENT.}

1. What makes a crowd of human beings a mere mass is that the single personalities do not attain independent significance,
do not each develop in a distinctive way. But with this effacing of personal peculiarities there also follows a lack of
inner connection, of organized social life. The single parts of the mass stand indifferent toward one another. Instead of
a real \textit{community,} in which the individuals are bound together by common ends, we then get a mere
society\footnote[1]{This contrast between the concepts \textit{community} (Samfund, Gemeinschaft) and \textit{society}
(Selskab, Gesellschaft) is developed by \emph{Tönnies} in his work mentioned above (XXIV, 2).}, in which the one individual
regards the other as a mere means. Outer interests bind human beings together without their therefore standing in any
essential inner connection. The one wants, with the least possible sacrifice on his part, to obtain as much as possible of
what is desirable to him and is possessed by the other. When the one affords the other means for his ends, it is only in
order to get means for
% --- p. 374 ---
\opage{374}his own ends in return. The relation between the individuals takes on an ever more impersonal character; for
what is impersonal we use and regard merely as a means for our own ends. Each one follows his instincts and his ends. But
since each one follows his interest without there existing any real community with the interests of others, where these
cannot be made means for him, collision necessarily comes about. Each one wants to live without seeing the necessity of
others living too. It is especially in the domain of material goods that the collisions must come. A material good can
belong to one alone. What satisfies my hunger does not satisfy my neighbor's. I must push aside his instinct of
self-preservation in order to satisfy my own, and the more material goods I appropriate, the fewer must remain for him.
One and the same atom cannot at the same time be a component of my organism and of his. Therefore a struggle of all
against all easily arises here, and the task presents itself of distributing goods as the welfare of all requires. The
demand of distributive justice (III, 9) is a demand for a social organization through which individuals can supplement
one another instead of fighting one another. The great question is by what way such an organization is to be brought
about, and what possibilities for it experience points out to us. What powers does the human race command for the
solution of this problem?

We discuss this question here, in social ethics, from a somewhat different side than in political economy. We have to
undertake an ethical valuation of the different ways in which a social distribution and organization can be produced.
Political economy, on the other hand, keeps more to an investigation of whether the possibilities are really present; it
discusses rather the technical side of the matter. Yet, as was already remarked earlier (III, 14), no sharp boundary can
be drawn between ethics and political economy, since the latter is not only a doctrine of the production and exchange of
material goods, but also a doctrine of their distribution.

Now there will be, as experience shows, two groups of forces that can work in the direction of a distribution in accord
with the principle
% --- p. 375 ---
\opage{375}of welfare, and thereby also in the direction of social organization in the domain of material culture. On
the one hand individual forces can work together in free associations, called forth by community of interest and by
sympathy; on the other hand the state with its compelling central power can intervene and lay down an ordering of the
work of material culture. We treat these two groups separately. They both in turn present different forms.

\begin{center}a. The ordering of labor through free association.\end{center}

2. The first effect of the principle of freedom was dissolving and isolating. Its setting up worked like a mine that blew
up the old building of society, and the world has not yet come to rest after this catastrophe. As we have seen, in
revolutionary zeal isolation was even regarded as a necessary condition for the preservation of freedom. But freedom shows
itself just as well in the choice of union with others as in the choice of separation from others. It means only this,
that no arbitrary obstacles are set to the individual's activity. Free association must precisely be said to be the
natural continuation of emancipation. If freedom has brought misfortunes and evils, it is --- in part at least --- able to
heal the wounds it has struck itself. The proclamation of freedom takes place by virtue of the principle of welfare,
because the free use of powers has value both as end and as means, and it is likewise by virtue of the principle of
welfare that emancipation is continued by association. This will come out clearly when we consider the main forms of
workers' associations that have appeared historically.

3. First of all the point is to abolish isolation and strife within the workers' own circle. The worker's dependence on
him who buys his labor hangs together, after all, chiefly with the strong competition among the workers themselves. Wages
sink on account of the strong supply of labor. By appearing in common before the buyers of labor the workers become able
to set their conditions, as regards wages, working time and conditions of work alike. Only thereby do
% --- p. 376 ---
\opage{376}they become really free men, since they have the possibility of getting their demands acknowledged. As long as
freedom is not bound up with the least power, it is only a word; and power is reached only through union and organization.
As so often in history, power must be unfolded in order that right may be acknowledged. Long after thraldom and serfdom
had ceased, and after the guilds had been abolished, employers regarded themselves as absolute masters and authorities
over their workers and considered unconditional authority over them a basic condition of industrial organization. They
laid claim to the rights of the old regime after the corresponding duties had fallen away. And not only this: they
appealed at the same time to the principle of freedom that the Revolution had proclaimed, and demanded only to enter into
a free contractual relation with the single worker; they used freedom to isolate. All the more, then, was it a matter for
the workers to unite; it became a basic condition for them in the struggle for existence.

This called forth the \textit{trade unions,} which especially in England have developed in a distinctive and significant
way\footnote[1]{In what follows I keep especially to the English trade unions, whose history is the most interesting and
best known, and here rely on \emph{Lujo Brentano}: \textit{Die Arbeitergilden der Gegenwart} (Leipzig 1871---1872) and on
\emph{Sidney} and \emph{Beatrice Webb}: \textit{History of Trade Unionism.} (London 1894). For Denmark reference is made to
J. \emph{Jensen} and C. M. \emph{Olsen}: \textit{Oversigt over Fagforeningsbevægelsen i Danmark 1871---1900.} (Copenhagen
1901).}. The very history of their development offers great ethical interest. They had a hard struggle to go through,
partly with the power of the state, which would not acknowledge them, partly with the employers, partly with the
coarseness of their own adherents. The workers were denied the right to unite around their common interests, and the more
they were placed outside the law, the more they resorted to deeds of violence, destruction of machines and murder. By the
prudence and moderation they increasingly displayed, however, they have gradually managed to procure acknowledgment; and
the more this acknowledgment of their ends and their activity made its way, the more the coarseness of their adherents
ceased.
% --- p. 377 ---
\opage{377}Their striving aims especially at obtaining an even and secure rise in wages and at any rate at preventing a
decline in the workers' circumstances. They see the importance of the workers' becoming accustomed to make certain demands
on life, accustomed to fight for a certain way of living (standard of life), so that they do not let themselves be pressed
down to mere animal self-preservation. Such a level is produced not by sudden and strong fluctuations but by regular and
lasting conditions. They therefore also lay greater weight on the shortening of working time, the procuring of healthier
outer conditions, safety measures, etc., than on raising wages. To this belongs not least the endeavor to let unhealthy
work at home be replaced by work in workshops. This endeavor has rightly been cited as testimony to the workers'
intellectual and moral development. ``A working class sunk in ignorance and social-political dependence will sacrifice
nothing to free its home from the squalor of the workshop, but will direct its endeavors exclusively to getting a little
more pay for its work and be content with that''\footnote[1]{\emph{Jensen} and \emph{Olsen}: \textit{Fagforeningsbevægelsen
i Danmark} p. 69.}. --- Through the overview of the market and of working conditions in different places which they manage
to procure by their widespread organization, they become able to decide which demands they have a prospect of carrying
through. The struggle for the workers' advance is now no longer waged blindly, but supported by clear insight into the
real circumstances. This organization on the workers' side has in turn made possible the establishment of boards of
conciliation and negotiation, so that representatives of workers and employers can meet to discuss common affairs.
Something has happened here within the single nation like what happens when arbitration between contending peoples takes
the place of war. And only through the struggle of the trade unions was the workers' civic independence acknowledged. It
is characteristic that
the English law regulating the relation between employers and workers was called until 1875 the ``Law of Masters and
Servants'' (Masters and Servants Act); only in that year was this
% --- p. 378 ---
\opage{378}name replaced by ``Law of Employers and Workmen'' (Employers and Workmen Act)\footnote[1]{S. and B.
\emph{Webb}: \textit{History of Trade Unionism.} p. 274 f.}. A proof that freedom is first carried through with the help
of association, although it is itself the condition for association's arising! In their struggle for freedom of
association the workers therefore also stand together with the adherents of radical individualism. But while these
thought that everything was attained when freedom --- freedom of association included --- was carried through,
association was for the workers the foundation of a further-reaching development.

Association brings with it effects of great ethical significance. --- The individual's feeling of independence is kept up
when he sees that by his participation in the trade union he can put a weight into the scale, while by standing isolated
he is able to accomplish nothing. Demands are also made on his competence, for the trade unions require that every member
shall have learned the work and really be able to earn an ordinary day's wage. Incompetent workers are not wanted, since
they make wages fall. --- Comradely feeling, sympathy with others, grows. The individual sees that his conduct at and
outside work, his competence, diligence, self-control and thrift, become of significance for far more people than himself.
An ethical sphere forms, so to speak, around him, a great family of which he feels himself a member. He comes to know
brotherhood, which was before almost forgotten for ``liberty and equality.'' He learns to subordinate his own interests to
the common ones. He feels solidarity with his fellow tradesmen and --- through the federation of the different trade
unions --- with other workers, indeed with the workers of other countries. His horizon widens; he gains the capacity to
set himself greater ends, and he grows through his relation to these greater ends. And through the common experiences and
the overview of commercial and factory conditions that determines the trade unions' policy he gains a clearer conception
of the workers' position relative to other classes of society, comes to know better both his rights and his duties as a
member of the race. --- It is, in short, a whole education from egoism to sympathy, from
% --- p. 379 ---
\opage{379}blind coarseness to clear-sighted strength, from struggle to peaceful negotiation, that goes on here. And all
this happens by way of freedom. There is no better answer to those who regard our time as a mere time of dissolution
than to point them to the formation of society and the ethical development that go on here. And it proceeds according to
the psychological laws brought forward earlier (XIII, 4). Fellow-feeling is developed through living and working
together, through common fate and common work. The Aristotelian principle comes clearly to light here. A peculiar
displacement of motive appears in this, that although it is often the purely individual struggle for self-preservation
that leads the worker to seek association, he can nevertheless so enter into the common interests of life that these
become ends for him directly. And finally, through these new forms of society effects are produced of significance far
beyond the workers' circle. The other classes of society are educated by this new element that actively intervenes in
development; they come to know their limits better, and they also get new ends and tasks, inasmuch as the new element
brings with it a successive transformation of social life as a whole.

Not all trade unions have reached the same stage of development. Both in organization and in procedure great differences
assert themselves. We have here kept to the picture presented by the most advanced unions. --- For the individual worker
a serious ethical conflict can arise during a strike decided on by the union, since he may be placed between his starving
family and what he must regard as the honor and welfare of his class. The unions' conduct toward the so-called
``blacklegs'' has often been harsh; but one must remember that an ethical conflict is present here. If the strike really
is in the interest of the whole class, it is undoubtedly the duty of the individual --- a duty that the feeling of
solidarity enjoins wherever it stirs --- to hold out as long as possible. Even if he does not belong to the union, he will
reap the benefit of its victory, and he will therefore not be able to separate his cause from its cause in the time of
struggle. It is a great responsibility that those take upon themselves who declare the war; but once war is
% --- p. 380 ---
\opage{380}declared, the individual must put up with the unavoidable suffering. And in these struggles qualities have
undoubtedly been displayed in narrow and hidden circumstances which on a larger stage would have brought historical fame
with them. An economic writer who, far from being an unconditional admirer of the trade unions, criticizes them very
sharply, \textsc{Stanley Jevons}, says of them: ``I do not doubt that if the history of strikes and labor disputes were
fully written, it would present as many examples of faithfulness and heroism, and of fearless suffering of misery, indeed
even of death, as many a war described in history''\footnote[1]{\textit{Trades Societies, their objects and their
policy.} (In the work: Methods of Social Reform). p. 115. Cf. with this a similar statement by \emph{Kropotkin}:
\textit{Mutual Aid.} London 1902. p. 270.}.% translated from Høffding's Danish

It has been objected against the trade unions that by striving for higher wages and more favorable conditions of work
they raise the price of goods that workers in other trades need. Here, then, the workers are required to take a regard
for consumers which employers do not take when they set the price of their goods. They are held up to an ideal which no
one thinks of holding up to merchants and manufacturers. But even apart from this the objection is
unfounded\footnote[2]{Cf. \emph{Brentano}: \textit{Arbeitergilden.} II, p. 233---244.}. When wages rise in one trade,
it also benefits workers in other trades, since the better-placed workers have greater purchasing power. The consequence
is thus only that a greater part of the national income benefits the working class. Distributive justice has taken a step
forward.

With greater right it may be objected against them that they found an aristocracy within the working class by wishing to
admit only competent and trained workers. A certain aristocratic tendency can also be traced among the different classes
of workers in the same factory, e.g. when the skilled workers demand a dining room for themselves and do not permit
unskilled workers to take a place among them. But all social development
% --- p. 381 ---
\opage{381}takes place, after all, in layers. The whole mass cannot be organized at once. It is a great advance when for
the time being the uppermost layers come along into the development. And the history of the trade unions shows that,
even if they should have a tendency to close themselves off and form an aristocracy, the principle to which the
association movement is due works in ever new layers. The ``unskilled'' workers too are beginning to organize; the strike
of the London dock workers (1889) marks an epoch in the history of association in this respect. Shortly afterward the
laborers in Copenhagen were organized, and their union was acknowledged after some struggle on the employers'
side\footnote[1]{\emph{Jensen} and \emph{Olsen}: \textit{Fagforeningsbevægelsen i Danmark,} p. 51. --- Cf., with regard to
the most recent development abroad: \emph{Harald Scavenius}: \textit{Fagforeningsbevægelsen i Frankrig} (1911) and
\textit{Af den nyere engelske Arbejderbevægelses Historie} (1912).}.

4. With all their advantages the trade unions nevertheless represent the interests of only a single class. And when by
their organization of the workers they make mediation between these and the employers possible, mediation always
presupposes preceding strife and opposition. Since, on the other hand, mediation can succeed only when common interests can
be pointed out, it is natural to ask whether there could not be free associations between workers and employers, just as
well as between workers among themselves. And such we do find in the most recent times. Prudent and sympathetically
disposed employers, who are not merely annoyed at the losses from strikes and lockouts but grieve over the bitterness and
discord that the present ordering of labor so often brings with it, have introduced giving the workers a share in the
profits, in order thereby to bind them more closely to their work\footnote[2]{Stanley \emph{Jevons}: \textit{On Industrial
Partnerships.} (Methods of Social Reform. p. 122---155). --- \emph{Stuart Mill}: \textit{Principles of political Economy.}
IV, 7, 5. --- In her criticism of the industrial partnerships \emph{Beatrice Webb} shows that they can become harmful to
the trade unions' endeavors. \textit{Die britische Genossenschaftsbewegung.} Leipzig 1893. p. 139---145.}. A certain fixed
part of the profit is reserved for the employer; from the rest a sum is then deducted for the repair of machines and the
extension of the business; and what then remains is divided into two parts,
% --- p. 382 ---
\opage{382}one of which falls to the employer, while the other is distributed among the workers in proportion to the wages
they draw. In addition the workers are given access to buy shares in the business. --- Such at least is the arrangement in
some of these associations.

Such \textit{industrial partnerships} are due to the initiative of employers; but their rise is nevertheless not a matter
of grace. In trades that are suited to this arrangement --- and these are especially those where the workers'
reliability in carrying out the work and in handling machines and tools is of great importance --- in such trades the
employers too will profit by forming a community with their workers. Just as the abolition of serfdom turned out to the
advantage of the owners, so progress will pay here too. The workers' increased diligence and care, the way in which they
check one another, the peace and trust that reign among all, the avoidance of stoppages of work, --- all this brings
material and spiritual gain with it. --- It is not only the workers, it is also the employers who need to be educated. The
social question can be brought nearer its solution only when employers regard their position as a social task that
brings social duties with it. To these duties belongs not only that of delivering certain products as well and cheaply as
possible, but also that of leading the little society at whose head they stand forward to greater material and
spiritual welfare. Employers can practice a kind of social robber-exploitation of a pernicious sort. So when they set up
factories, draw a mass of workers to them, reap the greatest possible yield from the business --- and then stop the
whole thing, either because they have earned enough or because it no longer pays. The workers are dismissed and now
stand, perhaps with large families, without bread and without work. The whole spirit in which the employers' activity is
carried on is therefore of the greatest importance, and the future of social conditions depends to a great extent on
their ethical qualities, not only, as is so often pharisaically said, on the workers'. It is not charity and great
sacrifices that are required. The most important thing here will be a keen eye, guided by sympathy, for common interests.
It may be a long prospect before
% --- p. 383 ---
\opage{383}such qualities develop. It is, after all, the case that the discontented improve sooner than the contented. And
it would in any case not be fortunate if the workers were to confine themselves to waiting patiently until it happens.
Their own development and their independence would suffer by it. Perhaps \textsc{Stuart Mill} too is right when he thinks
that ``long before the superior classes could be sufficiently improved to govern in the tutelary manner supposed, the
inferior classes would be too much improved to be so governed.''% Mill's English (Principles of Political Economy IV, 7, 1)

5. There is a kind of free association that aims still more than those so far named at abolishing the difference
between employer and worker, and between intelligent and physical work. It is the \textit{producers' associations,}
unions of workers who with money they have saved up or been lent themselves buy the means of production and afterward
divide the yield among themselves\footnote[1]{Stuart \emph{Mill}: \textit{Principles of political Economy.} IV, 7, 6. ---
L. \emph{Brentano}: \textit{Die christlich-soziale Bewegung in England.} Leipzig 1883.}. These associations bear witness to
the energy, intelligence and capacity for sacrifice that can be found among the workers, and already thereby contain a
good omen for the future. In order to reach their goal the workers who found these associations submit to privation and
toil, and moreover to a compulsion and a discipline which men working in the service of others would not put up with.
They have often managed to work their way forward in spite of all the resistance the authorities put in their way. In the
members of these associations an interest in enlightenment has been awakened, as well as a morality that could hardly be
produced by sermons or through temperance societies. --- Producers' associations are often founded by trade unions or
consumers' associations.

They presuppose, however, qualities that will for the time being be present only in very few. It has turned out that
they thrive only where the start is made with self-saved means, while support from the state or from private persons has
not worked well. It will --- for the time being at least --- be only a small elite that works its way forward in this
way. And often the successful founders of producers' associations,
% --- p. 384 ---
\opage{384}when they are not animated by great and lasting enthusiasm, end as capitalists and employers who take other
workers into their service on the usual terms.

6. Through trade unions, industrial partnerships and producers' associations the workers seek to procure their class a
share in the yield of production. But there is another way by which union can lead to raising the condition of the class.
For what matters is not only how much money one earns by one's work, but also --- and every bit as much --- how much one
can get for one's money. And the working population as a whole is, after all, society's greatest buyer; it thus has a
means of power in its very consumption, which is by no means --- as Karl Marx thought --- vanishing in comparison with
that of the capitalists. The point is only to draw advantage from this very consumption. By the formation of
\textit{consumers' associations} it becomes possible to get a share in the profit that otherwise falls to the middlemen.
The consumers form an association whose first end is to obtain cheaper provisions by having the turnover of goods led as
directly as possible from producers to consumers. But as it goes with the other associations, so it goes with these: new
valuable qualities develop as a consequence of working together, and higher, more comprehensive ends are set besides the
original and elementary ones. Where the consumers' associations have properly developed (especially in England and
Switzerland)\footnote[1]{Beatrice \emph{Webb}: \textit{Die britische Genossenschaftsbewegung.} Leipzig 1893. --- \emph{Hans
Müller}: \textit{Die schweizerischen Konsumgenossenschaften.} Basel 1896. --- Only after the repeal of the Socialist Laws
did the German workers get a proper eye for the significance that consumers' associations could have for the development
of the working class. Only from the beginning of the nineties have they properly taken an upswing, partly under violent
opposition on the part of the retail traders. Cf. with this R. Riehn: \textit{Das Konsumvereinswesen in Deutschland.}
Stuttgart und Berlin 1902.}, they have been mighty causes of the development of the feeling of solidarity. They have
moreover required honesty and unselfishness, prudence and energy, a democratic disposition and capacity for
self-government. Beside the trade unions they have become the most important school of the working class, the most
important means of social organization of this class, which after
% --- p. 385 ---
\opage{385}the abolition of the Stavnsbaand and the guilds stood as a chaotic mass. It has become of special significance
that work in the service of such associations has produced a staff of people with administrative capacity and with
understanding of comprehensive interests. It is a development that has taken place from below, through the exercise of
powers in the smaller circles and often on apparently insignificant and material tasks. Yet it is not merely the
procuring of cheap provisions that stands as the end. Ideal ends too are kept in view, especially the intellectual and
aesthetic culture of the members, and many consumers' associations have lecture halls and libraries in their buildings.
With the help of common capital, moreover, they often go over to common production of necessities, and it has turned out
that producers' associations of this kind lead a sounder existence than those formed on bare ground, and that they do not
so easily degenerate into mere joint-stock companies. This whole movement will perhaps become the most important
counterweight to the extreme division of labor to which the development of productive activity has led. It unites just
as much as this has divided. And it will be able more and more to embrace all domains, the ideal as well as the material.
The free-church movement has rightly been pointed to as the model for the consumers' associations, and the common meal as
the old image of intimate community between human beings\footnote[1]{Ch. \emph{Gide}: \textit{La solidarité économique.}
(In the work ``Philosophie de la solidarité''). Paris 1902. p. 229---231.}.

It is not necessary to set out more closely how here, as with the associations mentioned earlier, the three main laws of
ethical development (XIII, 4) come to light. Whatever future may be in store for these and similar associations, their
history surely contains the most significant historical experiences the human race has made in the last century. The
forms of society that develop on the ground of freedom have the great advantage that they are due to demands that assert
themselves of their own accord and are not artificially coaxed forth. They are tested first in smaller circles before they
spread in larger. They show us
% --- p. 386 ---
\opage{386}how self-control, sympathy and a sense for ideal goods awaken as soon as human beings struggle for life in
common and do not stand each by himself with his blind urge for self-preservation.

\begin{center}b. The ordering of labor by the intervention of the state and the municipality.\end{center}

7. The social phenomena described in the foregoing belong decidedly under the concept of the free cultural community.
They consisted in free union, conditioned by the striving after common ends. Their ethical significance lay in the
psychological changes and displacements that life in these associations calls forth, and in the new duties, tasks and
conflicts that arise through the relation to them. But the question naturally arises how these free cultural unions
relate to the forms of society whose ordering is not left to the involuntary or voluntary unions of free forces, but is
if necessary carried through with power and compulsion. If by \textit{socialism} is understood the view that the state or
the municipality ought to dispose of all means of production and determine their use as well as the distribution of the
yield of production, then it is the relation of the free cultural community to socialism that we stand before here. But it
must be well noted that the expression socialism is used of very different social views, something which especially those
who are frightened by the mere sound of the word must not forget. Under all its forms socialism designates the
maintenance of, and admiration for, a social ideal of the future; but the single traits in this ideal picture, and still
more the relation of the whole ideal of the future to the present state as well as the way in which the transition is to
be made from the present state to the future state, can appear in most different ways in the different views. We must
therefore briefly characterize the different main forms of socialism. There appears, when we set them in historical
order, a rhythmically progressing course of development in socialist ideas, called forth by interaction with historical
circumstances.

a. \textsc{Plato} in the \textit{Republic} has portrayed an ideal society in which private property was abolished for the
ruling classes, so that they could live for public interests and
% --- p. 387 ---
\opage{387}intellectual endeavors. Influenced by him, but above all under the influence of social conditions in England at
the end of the Middle Ages, \textsc{Thomas More} wrote his \textit{Utopia} (1516). And a century later \textsc{Campanella},
in the prison to which his part in the social and political movements of southern Italy had brought him, wrote his City of
the Sun \textit{(Civitas solis.} 1623)\footnote[1]{See on this work \textit{Den nyere Filosofis Historie}\textsuperscript{2}.
I. p. 152 f.}. In the Platonic state, in Utopia and in the City of the Sun alike, the social order is laid down by
compelling laws. No distinction is yet made between society and state, and the power of free forces is not known. They
thereby differ from the later utopias, such as \textsc{Fr. Chr. Sibbern}'s (in \textit{Meddelelser af Indholdet af et Skrift
fra Aaret 2135}), which lets complete community come about only after all political life and all compelled authority have
been abolished\footnote[2]{Cf. my \textit{Mindre Arbejder.} p. 108---112.}. --- The interest of these ideal constructions
lies chiefly in the criticism of actual social conditions that they presuppose. It is injustice and want that by effect of
contrast call forth the social ideas. At the foundation lies here the same standard that we have set up in the foregoing,
although this standard is not applied with full consistency. Most characteristic of the form of socialism appearing in
the works named is that the picture of the future is set without further ado in opposition to the actual state. No attempt
is made to show how the transition from the present to the ideal future can be made. At all times the critical attitude
toward the actual state and the enthusiastic faith in the ideal have been the marks of socialism, so that now the
critical, now the ideal element had the preponderance. And it owes its great influence partly to criticism and
indignation at what is given, partly to the attracting power of the picture of the future. Human beings need not only
criticism, but great pictures that can fill their minds and give their need for hope a definite content, and these
utopias satisfy this need in relation to their historical conditions. Other forms of socialism seek in various ways to
find intermediate links between doubt and
% --- p. 388 ---
\opage{388}faith, between the dark picture of the present and the bright picture of the future.

β. A series of men from the first decades of the 19th century expected great things from the cooperation of free forces.
\textsc{Saint-Simon}\footnote[1]{See on this man, whose ideas became of great significance for the development of philosophy:
\textit{Den nyere Filosofis Historie}\textsuperscript{2}. II. p. 313---317.}, \textsc{Charles Fourier} and \textsc{Robert
Owen} met in the fundamental thought that the want and injustice arising from men's mutual exploitation of one another
would be remedied if they united in order to exploit nature in common. Here, then, a definite way to the attainment of the
ideal goal is pointed to. But the limitation of these men lies in this, that they in turn know no means of getting onto the
way, --- of forming the free unions and brotherhoods. They think chiefly of associations of people who in philanthropic
enthusiasm would directly realize the social ideal. This \textit{philanthropic socialism} thus comes nearer to reality, to
be sure, than that older tendency which one could call \textit{utopian socialism,} but it still lacks several intermediate
links. Nevertheless --- and this is not the least interesting thing about it --- its enthusiasm was not without use, but
became an important operative cause within the course of social development. It was to a great extent the disciples whom
Saint-Simon, Fourier and Owen won within the working class whose enthusiasm and practical sense led to the founding of
trade unions, consumers' associations and producers' associations. The producers' associations are especially the work of
the Saint-Simonians; from them the English ``Christian Socialists'' took them over. In the history of the English trade
unions and consumers' associations Owen's ideas and Owen's disciples play an essential part. The Swiss consumers'
association movement is inspired by Fourierist ideas, later influenced by the English movement. It was thus shown that it
is not enough to begin with ideas, but that the involuntary need and unfolding of life must meet the movement of ideas.
Construction and experience must work together. This became possible because --- as a consequence of the reforms and
discoveries of the 18th century ---
% --- p. 389 ---
\opage{389}a distinction could be made between society and state. A free social development could now take place, by which
new ideas could be tested from below, in the smaller circles of life, so that the way from above, through law and
compulsion, was not the only possible one.

γ. While the philanthropic socialists thought chiefly of a free organization of the forces liberated by the great
Revolution, the third main form of socialism marks a return to utopian socialism, insofar as it teaches the state's taking
over of all means of production. But it thinks it has got beyond utopianism in that what in this stood only as an object
of faith and a picture of fancy is now derived and proved by way of science. This tendency stands toward utopianism
roughly as speculative philosophy in the period of Romanticism stood toward positive religion: agreed in content, but
different through its scientific form\footnote[1]{Cf. \textit{Den nyere Filosofis Historie}\textsuperscript{2}. II. p.
183---187.}. It is hardly an accident that the founders of this tendency came out of the Hegelian philosophy. It likes to
call itself scientific socialism, but ought rather to be called \textit{speculative socialism.} The scientific significance
that \textsc{Friedrich Engels} and \textsc{Karl Marx}, the founders of this tendency, undoubtedly have consists chiefly in
their having maintained the historical character of the concepts and laws of political economy, whereas there had
otherwise been a tendency to see eternal ideas and truths in them. This hangs together with Engels and Marx bringing the
social question into connection with the whole of historical development. But they conceive this connection in such a
way that it is the economic questions around which the whole of history turns. All social order, all morality and law,
all political and religious ideas depend according to Marx on economic conditions. The whole of ideal culture consists
only in the way in which human beings become conscious of their economic conditions of existence. In particular, social
science itself is a product of the historical movement in the direction of new economic conditions, not --- as with
utopians and philanthropists --- a more or less well thought out system. True social
% --- p. 390 ---
\opage{390}science arises when one becomes conscious of what has happened or is in the process of happening. And what has
happened is that, as a consequence of the historical distribution of power, a division of labor has come about by which
the worker was separated from the means of work (the land and the tools), capital being concentrated in relatively few
hands. With the introduction of machines the worker's dependence on the employer was completed. But after the capitalists
as a class have expropriated the workers, the small capitalists were expropriated by the great. We now stand at a turning
point where there is a decided disproportion between productive power, which is to be sought in labor alone, and the
ordering of production, which lets by far the greatest part of the yield fall to those who do not themselves work, namely
the capitalists. In this disproportion, which keeps growing greater and greater, the social question consists. The
solution can come only when the possibilities of the present order are exhausted. This will happen when the concentration
of capital has advanced so far that the state can easily expropriate the few capitalists. The reversal then takes place
by a negation of the negation: the expropriators hitherto are now themselves expropriated\footnote[1]{``The transformation
of scattered private property, arising from individual labour, into capitalist private property is, naturally, a process,
incomparably more protracted, violent, and difficult, than the transformation of capitalistic private property, already
practically resting on socialised production, into socialised property. In the former case, we had the expropriation of
the mass of the people by a few usurpers; in the latter, we have the expropriation of a few usurpers by the mass of the
people.'' % the English translation of Das Kapital (Moore and Aveling), ch. 32; Høffding translates the German
Karl Marx: \textit{Das Kapital,} I. 2. Aufl. Hamburg 1872. p. 793.}. But the condition for this is that the workers have
gathered in great associations, ready to take over power, and also that great technical advances have made common
operation both possible and --- because of the costliness of the means of work --- necessary\footnote[2]{In one of his
earlier writings Marx expresses himself clearly and interestingly about his view of the task of social science. ``The
socialists and communists are the theorists of the proletariat. As long as the proletariat is not sufficiently developed
to constitute itself as a class, and the struggle of the proletariat with the bourgeoisie therefore has no political
character, and as long as the productive forces within the bourgeoisie itself are not yet sufficiently developed to
reveal the material conditions necessary for the liberation of the proletariat and the formation of a new society, so
long are these theorists only utopians who think up systems and seek a regenerating science to relieve the needs of the
oppressed classes. But in the measure that history advances, and with it the struggle of the proletariat stands out more
clearly, the theorists no longer need to seek science in their own heads; they need only take account of what is happening
before their eyes and make themselves its organ\dots{}. From this moment science becomes the conscious product of the
historical movement; it is no longer doctrinaire, but has become revolutionary.'' % translated from Høffding's Danish
\textit{Misère de la philosophie} (1847). German transl. Stuttgart 1885. p. 121. (A polemic directed against
Proudhon).}. What matters first and foremost is that the workers
% --- p. 391 ---
\opage{391}awaken to consciousness of their mission and their power, and therefore Marx's call was: ``Proletarians of
all countries, unite!'' Only through the class struggle can power be wrested from those who hold it. And here, then, is the
point where Marx's speculative line of thought issues in agitation.

The most characteristic thing about this whole view is the purely deductive way in which it is grounded. It is derived
from a general theory of the philosophy of history according to which economic conditions lie at the foundation of all
culture. Ideal culture is only an effect or reflection of material culture. Therefore labor, by which is meant material
labor, is regarded as the source of all value, and therefore it is denied that historical development is in any way
determined by ideal factors. It is a struggle for food that alone is in question, in however many ``ideological''
disguises this original motive may appear.

The relation between economics and ethics, and between material and ideal culture, is not so simple. It is right that
material culture is a presupposition of ideal culture, insofar as one must live in order to be able to live ideally, ---
and insofar as one can live without living ideally. But it does not follow that ideal culture should be merely the effect
and reflection of material culture. When the material demands of life are reasonably satisfied, the energy not used for
them can be used for ideal pursuits, for thought and mood, for art, religion and science. But this use is not a matter of
course, and there
% --- p. 392 ---
\opage{392}are problems here which Marx has not heeded. Moreover, ideal culture, once it has developed, acts back on
material culture. This happens already under such primitive conditions that we know no people whose struggle for existence
is not determined by religion, tradition and morality. ``Ideology'' very early becomes a factor in development that no
theory of the philosophy of history dare overlook. Even if the thoughts that the struggle for existence calls forth at
first stand merely as means and detours for reaching the material goal, they very soon become ends themselves, value being
ascribed to them for their own sake. It may be, for example, that the individual worker first joins his comrades in order
to obtain a definite advantage thereby; but the honor and advance of his class can later become an end that he pursues
without egoistic ulterior motive. Such displacements of motive make the relation more involved than Marx's purely deductive
account can acknowledge.

Marx's theory would portray social development wholly independently of the influence of all ideal motives. But the history
of the social problem and of socialism itself proves that such motives do in fact play their part. And in the foregoing
(IV, 6) we saw that conscience itself is after all also a force, a link in the causal series that determines development.
According to Marx all theory, all ``ideology,'' is to be only consciousness of what happens; but what value has that which
thus stirs in the brains of more or fewer individuals, if it has no practical effect? Marx indeed thinks that it is of more
than theoretical interest to find a society's economic law of motion. ``Even if a society,'' he says in the preface to
``Das Kapital,'' ``has got on the track of the natural law of its movement, it can neither leap over nor decree away natural
phases of development. But it can shorten and lessen the birth-pangs.'' % translated from Høffding's Danish
Already such a shortening and lessening is more than Marx could consistently concede. But Marx has altogether more
presuppositions than he will admit. His theory is at bottom an ethical theory; the result he arrives at rests on an ethical
postulate, which at least in one single place he comes to betray, the postulate namely that a human being is not to be
treated merely as
% --- p. 393 ---
\opage{393}a means, but always at the same time as an end. This postulate, which contains the justification of the
indignation that can be traced through his account, however learned and heavy this may be, is expressed in the following
proposition: ``In the capitalist mode of production the worker exists in order that existing values may find employment,
instead of, conversely, objective wealth existing to serve the worker's need for development.'' (Das Kapital. 2. Aufl. I.
p. 646). % translated from Høffding's Danish
Whence, one wonders, does Marx know what wealth exists for? By a purely historical route, with deduction attached, he can
know nothing about it, only something about what wealth will one day be used for according to a necessary law of
development. From the worker's having only a slight share in the culture he helps to produce it by no means follows in
itself that he is to have a greater share, even if that is possible. This follows only if one proceeds from the ethical
principle that every personality is an end, not merely a means. It is then no wonder that this principle involuntarily
forces its way forward in Marx, although he will acknowledge no special ethical presuppositions, but in ``scientific''
superiority leaves such presuppositions to utopians and philanthropists. In an utterance like the one cited he expressly
demands an intervention in the course of things from ethical motives. That proposition contains a whole ethics in germ.
When the practical valuation that thus lies at the foundation is taken into account, one better understands how knowledge
of the law of development can have a ``shortening and lessening'' effect. It is indignation at the distance between ideal
and reality that has a goading and strengthening effect in the struggle against resistance. Marx's own ``ideology'' is a
mighty weapon in his and his adherents' hands, and it perhaps came into being from the first because it could serve as a
weapon, not merely from purely theoretical interest. In vain does Marx seek to deny or conceal the idealism that lies
behind his own appearance, and which is at the same time a necessary condition for such great things to be accomplished
in history as the production of a new social order. He stands on an ethical standpoint without being willing to own it.

In his demonstration of the way to reach the goal Marx lays the weight on the class struggle. It is for him the most
important phenomenon
% --- p. 394 ---
\opage{394}in modern social development\footnote[1]{In the account of Marxism that \emph{Werner Sombart} gave in his Zürich lectures, this side of the matter was quite especially stressed. The lectures have appeared under the title:
\textit{Socialismus und sociale Bewegung im 19. Jahrhundert.} Bern 1897. In the discussion that followed these lectures I
had occasion to express myself on Marxism as Sombart had presented it. (See the work just cited, p. 56---59). It has
interested me to see that Sombart later reproached Marx with the trade unions' being for him only means to revolution,
without independent significance for the workers' development. Sombart has further sought later to show ``that capitalism
and socialism are not opposites that exclude one another, but that their ideals can to a certain degree very well be
realized in one and the same society.'' The brilliant pioneers, the royal merchants, cannot be done without if economic
progress is to be possible. But the labor relation is to pass over from the domain of private law to that of public law,
and the working class is to be secured co-determination in economic and political decisions. \textit{(Dennoch!} Aus
Theorie und Geschichte der gewerbschaftlichen % PRINTED AS IS (read at the image)
 Arbeiterbewegung. Jena 1900. p. 92---94). In his great work \textit{Der
Kapitalismus} (1901) he distinguishes --- under the evident influence of a more critical philosophy than the one that lay
at the foundation in Marx --- between three different points of view from which economic phenomena must be regarded: 1) the
causal point of view, which seeks out the actually operative causes and motives; 2) the teleological [better: ethical]
point of view, which proceeds from definite social-political ends; 3) the critical point of view, which leads to a
discussion of the fundamental concepts that the causal and the teleological points of view use, with regard to their
justification and their limitation. \textit{(Der Kapitalismus.} I. p. XXXII). --- With this view the transition has really
been made from speculative to empirical socialism.}. Yet it is hardly right to keep to the negative side of the matter.
Where a new social layer is to work its way forward, it naturally has a hard struggle to go through with the layers that
hitherto held the whole of social power, a struggle that is not always bloodless. But the weight must not be laid
exclusively on the opposition to the other classes. New qualities are acquired through the union, the brotherhood among
the workers themselves, that the struggle makes necessary. The association makes it possible for the individual, through
the culture and subordination required in the service of common ends, to come to know better his duties and rights as a
member in the service of mankind. Through this growing
% --- p. 395 ---
\opage{395}self-esteem, as well as through the material power the associations dispose of, and through their political
influence, the workers come more and more to stand as ends, not merely as means in the social process. That Marx does not
stress this side of the matter more decidedly comes from this, that according to his theory, with the growing concentration
of capital, ``misery, oppression, tyranny, degradation and plunder'' are to grow, and thereby the catastrophe set in in
which ``the expropriators are expropriated.'' But the educating influence the associations exercise on the workers, and the
steady improvement of the material conditions of life that these same associations make possible, do not agree well with
the catastrophe theory. Had Marx ascribed a positive significance to the associations, it would have been impossible for
him to deduce the catastrophe.

In spite of the scientific character Marx thinks he has given socialism, he has nevertheless not wholly broken with
utopia. This shows itself not only in the certainty with which he foretells the catastrophe, but also in the hints ---
few and short, to be sure --- that he gives of the social state that will set in after the catastrophe. After the
expropriation of the few capitalists by the whole mass of the people there is to be no more class difference. ``Will
there,'' asks Marx in one of his earlier writings\footnote[1]{\textit{Misère de la Philosophie.} German transl. 1885. p.
181 f.}, ``after the overthrow of the old society, be a new class domination culminating in a new political power? No. The
condition for the liberation of the working class is the abolition of every class\dots{}. There will no longer be any
political power properly so called, since political power is precisely the official expression of class antagonism within
bourgeois society.'' % translated from Høffding's Danish
A society without ``any political power properly so called'' is evidently a utopia, and indeed a utopia that surpasses
those of Plato, More and Campanella; for in their ideal states there is an organized political power. In his later
writings Marx expresses himself only vaguely and negatively about the future social state. He and his adherents think that
there will always be time enough to speak of it when the present order of things is past. This too
% --- p. 396 ---
\opage{396}hangs together with the catastrophe theory, according to which the ``reversal'' or the ``revolution'' is to lead
to something wholly opposed to what exists; but it conflicts with the experience that the future is prepared in what
precedes it not only negatively but also positively. Marx, moreover, does not have the naive trust in the constructions of
imagination that the old utopians had, and the vague hints he allows himself come forth rather for use in agitation than
as necessary links in the theory.

δ. In opposition to utopian, philanthropic and speculative socialism alike, there asserts itself in the most recent
times, especially in England, an endeavor that can rightly bear the name of \textit{empirical socialism.} It does not aim
at constructing a picture of the future, and it sees that philanthropic wishes are not enough, but that one must soberly
ask and test what can be attained under the given conditions; nor does it rely on deductions from the philosophy of
history, but will proceed inductively, testing the possibilities on the basis of definite experiences. And it imposes this
limitation on itself not only from critical deliberation and from aversion to demagogic
declamation\footnote[1]{Cf. \textit{Fabian Tracts} No. 41. (\emph{Bernard Shaw}: The Fabian Society. What it has done, and
how it has done it). p. 5. ``Our preference for practical suggestions and criticisms, and our impatience with all general
expressions of sympathy with working-class aspirations, not to mention our way of chaffing our opponents instead of
denouncing them as enemies of the human race, repelled from us some earnest and eloquent socialists\dots{}. There was too
much equality and intimacy among the Fabians for any member to stand up and preach to the others in the way the working
classes still put up with from their leaders.'' % after Shaw's English, as far as Høffding's Danish allows
--- The Fabian Society is named after the Roman general Fabius, who accomplished more by waiting for the right moment and
preparing it than many hotheads who thought the thing was to strike at once.}, but also and above all because it, more than
any of the other forms of socialism, acknowledges freedom both as means and as goal. A social organization has real value
only when it consists in the union of free forces. Therefore great weight is laid on the free workers' organizations which
% --- p. 397 ---
\opage{397}have been mentioned in the foregoing. The lower degrees and forms of freedom are to be used to produce the
higher. The capitalist system is reproached with hindering the development of personality in so many by pressing down the
conditions of existence and making them insecure. The point is to maintain a certain economic level (standard of life) for
the workers, so that they can acquire and maintain a higher spiritual level. By socialism (or, as is often said in the
most recent times, collectivism) is here understood a doctrine according to which it is society's business to maintain
such conditions for labor that the physical and spiritual development of the workers is not hampered. Regard is also to be
taken for the different powers of individuals; the weak is to work in proportion to his slight power, just as the strong
in proportion to his great power. Steadiness in economic conditions is to be provided for, so that the paralyzing feeling
of insecurity can fall away. The outer, physical and social conditions under which man lives determine his character to a
great extent. Therefore care ought to be taken that the social machinery be shaped as far as possible so that it can have
fortunate effects on character. Here it is of no use to rely on the philanthropic disposition of employers. And the
workers themselves are able only through experience to appreciate the significance of healthy, pure and noble conditions
of life. Therefore it is the business of society, of the state and the municipality, to procure as favorable hygienic and
moral conditions for the great population as possible. The working class and its spokesmen must use the civic rights they
owe to radical individualism to gain growing influence on public life through school boards, municipal councils and
parliament. Thereby experiences can be harvested and experiments made that can guide further development. State socialism
begins from above and would work through a bureaucratic regime and on a doctrinaire foundation. Empirical socialism lays
weight on self-government in trade unions and consumers' associations, in municipality and state; it moves through the
small circles to the large ones and seeks to develop the people's powers, through work on small tasks, to cope with the
great
% --- p. 398 ---
\opage{398}tasks. It is characteristic that empirical socialism\footnote[1]{The best account of it that I know is
\emph{Sidney Ball}'s essay: \textit{The Moral Aspect of Socialism} (International Journal of Ethics VI. Cf. the discussion
called forth by this essay in the same journal, vols. 6 and 7). The expression ``empirical socialism'' I have from
\emph{Sidney} and \emph{Beatrice Webb}: \textit{History of Trade Unionism.} p. 403. These authors' historical works show
how empirical socialism has developed. Cf. also \emph{Hans Müller}'s work on the Swiss consumers' associations. --- Since
1883 the \emph{Fabian Society} in London has worked through writings and lectures in the same direction. \textit{(Fabian
Tracts,} a series of small pamphlets, clearly and concisely written, can be had by application to The Fabian Office. 276
Strand. London, W. C.). --- The standpoint of empirical socialism is not very different from the one \emph{John Stuart Mill} took
toward the social question. \textit{(Den nyere Filosofis Historie}\textsuperscript{2}. II. p. 427---430). Nor does
\emph{Eugen Dühring}'s standpoint seem to be very different (work cited, p. 564 f.).} has its home in England, while
speculative socialism has its home in Germany\footnote[2]{Marxism is an offshoot of German speculative philosophy. As late
as 1891 Fr. \emph{Engels} wrote under his portrait: ``We German socialists are proud of being descended --- not only from
Saint-Simon, Fourier and Owen --- but also from Kant, Fichte and Hegel. The German labor movement is the heir of German
classical philosophy.'' (\emph{Werner Sombart}: \textit{Friedrich Engels.} Ein Blatt zur Entwickelungsgeschichte des
Sozialismus. Berlin 1895. p. 13). --- Cf. on the relation of Marxism to the Hegelian philosophy \textit{Den nyere Filosofis
Historie}\textsuperscript{2}. II. p. 266. 282 f.}. While speculative socialism is most reminiscent of the utopians,
empirical socialism hangs together most closely, as regards both character and historical origin, with the philanthropic
socialists.

8. From the standpoint from which we here regard social development, we must sympathize in advance with socialism on two
points: in its most important fundamental thought and in the main features of its criticism of the present social order.
And these two points are common to all forms of socialism.

Often, to be sure, socialism appears as a dream in which the mind has found rest when it was tormented and frightened by
the misfortunes of the time. But even then it contains a real ethical fundamental thought: the idea of a distributive
justice, of a perfect society in which everyone's capacity and need comes
% --- p. 399 ---
\opage{399}into its own. It appears especially in our century as a useful counterweight to the one-sided individualism
that splits society into isolated individuals. It in reality maintains the very fundamental thought of social ethics: that
the individual's position in society ought to be determined by the good of the whole society (the individual's own
included). And it is from this fundamental thought that socialism subjects the present social order to sharp criticism.
What self-examination and self-knowledge are for the single human being, such a criticism is for society. It lays bare
the evils and sufferings, and that is the first condition for their being remedied.

But one can acknowledge both the fundamental thought and the criticism without acknowledging the means and ways by which
the thought is to be realized and the evils remedied. It is one thing to point out the disease, another to prescribe the
remedy. Against such forms of socialism as proceed by testing and experimenting and work with the help of free
associations or of the existing municipal and political organizations, to which the workers are getting more and more
access, nothing can be objected. Unless, that is, one --- as is by no means unusual --- regards the present distribution of
property and labor as the sum of all wisdom. To what degree this order can be replaced by another, history must
show\footnote[1]{Cf. what so expert and conservative an observer as \emph{Alexis de Tocqueville} states as the result of
his inquiry and his experience: Ce qu'on appelle les institutions nécessaires, ne sont souvent que les institutions
auxquelles on est accoutumé, et en matiére de constitution sociale, le champ du possible est bien plus vaste que les hommes
qui vivent dans chaque société ne se l'imaginent. \textit{Souvenirs.} Paris 1893. p. 111 f.}. The possibilities we cannot
see through in advance --- all the less since the constant displacements (XIII, 4) can bring it about that something
wholly other often comes out of conscious endeavors than could be guessed in advance. The displacements of motive and of
value make interventions in social development leaps in the dark. But to do nothing is here also to leap. Every
intervention or non-intervention in the course of things must be motivated by the best-grounded conviction that can be
won.

% --- p. 400 ---
\opage{400}Empirical socialism sets itself the goal of altering the existing order of labor and property in a social
spirit; it leaves it undecided how far development will be able to go. An energetic worker in the Swiss consumers'
association movement has defined socialism as ``the doctrine of such an arrangement of human society that the disproportion
between excessive wealth and mass poverty is remedied''\footnote[1]{See the quotation in \emph{Hans Müller}: \textit{Die
schweizerischen Konsumgenossenschaften,} p. 453.}. A task is set here that can be worked at step by step. In contrast to
this stands Marxism's definition of socialism: ``Abolition of the capitalist order of labor by socialization of the means
of production, and class struggle as the way to this goal''\footnote[2]{Werner \emph{Sombart}: \textit{Friedrich Engels.}
p. 25. It is added: ``Diese programmatische Kernpunkte werden mehr und mehr Gemeingut des kämpfenden Proletariats.''}. It
is clear that the two views set the free cultural community in a different relation to the state; according to the first,
the spirit of the free cultural community is to be gradually grafted into the mechanism of the state; according to the
second, the point is to get control of the mechanism of the state as quickly as possible in order to be able to determine
the ordering of the free cultural community.

In the closer critical examination of socialism that follows we have chiefly this last, speculative or Marxist view in
mind. It will appear that the criticism of it leads to laying all the more weight on empirical socialism.

9. The power of the state exercises under all circumstances a great influence on the distribution of the means of work and
of consumption, and one can to that extent say that every constitution is a relative socialism. It will on the whole prove
impossible to stake out absolute boundaries for the activity of the state. The nature of the state is not unchangeable; it
develops with human nature and with historical circumstances, and no one can say how the state will one day be constituted
in the future. But by handing over all distribution to the power of the state one presupposes qualities in the human beings
who exercise the power of the state (and it is, after all, always definite human beings
% --- p. 401 ---
\opage{401}who exercise it) which these cannot so far be said to have shown themselves to possess. How could human beings
become so perfect that they would not misuse so enormous a power, when they already, as history amply shows, misuse the
smaller power that has hitherto been entrusted to those who govern? It is characteristic that in England, the country where
the sharpest control can be exercised over the use of the power of the state, there is the greatest disinclination to
extend its domain. State socialism presupposes a superstition about the state, a trust in what can be done from above
downward, which overlooks that the helm of the state is after all always held by human beings, not by gods. And this will
not be changed even when the majority of the people determines the composition of the government. When human beings become
as good and capable as socialism presupposes, the social question will already be out of the world.

Not only moral perfection, but omniscience too would have to be possessed by the bearers of the power of the state in the
socialist state. In order to be able to distribute work and yield, they would have to know the capacities and needs of the
different individuals. But both capacities and needs are in constant development, and the individual is the one who can
best discover them, if only they are allowed to unfold as freely as possible, so that they can be tested. The state, to be
sure, has the task even now of choosing the individuals best fitted for certain activities and of meeting various needs.
But there are surely few who think that it exercises this office to such a degree of perfection that one should wish it
extended to everything without compelling reasons. And the state, as it is now placed, has the free impulse to
development, the free initiative of single individuals, to rely on: it can choose among those who out of their own need
have developed in a certain direction. It has, moreover, private undertakings to compete with in many domains. However
extended one may imagine the activity of the state, it will nevertheless not be able to do without the competition of free
initiative and private activity, if it is not to stiffen into doctrinairism and slovenliness. This holds in the domain of
material and of ideal culture alike.

As regards the distribution of the yield in particular, the great question arises here what just distribution really
means.
% --- p. 402 ---
\opage{402}The socialist writers divide here into two groups, some wishing to have each one's share determined by the
work he has performed, while others wish to have distribution determined by each one's need. The first standpoint is
expressed in \textsc{St. Simon}'s proposition: ``Chacun doit être classé selon sa capacité et rétribué selon ses
oeuvres!'' --- the second in the words of the Social Democratic \textit{Gotha Program:} ``Jedem nach seinen
vernunftgemässen Bedürfnissen.'' --- In the first case one comes up not only against the difficulty of dividing equally
between the different kinds of material and spiritual work that society needs, but also against the great problem of how
one can make sure that the value of the yield corresponds to the work spent. For it depends, after all, on whether the
product of labor has sufficient use value, that is, whether it satisfies real needs in society. The value of a product is
not determined merely by the work its production has cost, or the time this work has taken. It depends on whether there is
use for what has been produced, thus whether a want, a need, is present that can be satisfied by it. Therefore one can
regulate work only if one can regulate needs. The first theory, then, like the second, must grant the state authority to
determine the needs of individuals. And one must also see to it that nothing other and nothing more is produced than can be
consumed in the country itself. For over foreign needs the single state is not master. Therefore world trade, which leads to
more being produced than the country itself needs, and thereby makes the country dependent on foreign countries, must be
restricted and controlled by the government. This consequence J. G. \textsc{Fichte} saw and was not afraid to draw it (in
his work \textit{Geschlossener Handelsstaat.} 1800). This would, to be sure, also stop up the source of the modern social
disharmonies; for these go back historically to the 13th and 14th centuries, when a world market began to form, so that
production came to be determined by other than the domestic, easily surveyed needs\footnote[1]{Lujo \emph{Brentano}:
\textit{Ueber die Ursachen der heutigen socialen Noth.} Leipzig 1889. Cf. the same author's \textit{Die Arbeitergilden der
Gegenwart.} I. p. 58 ff., II. p. 320.}. ---
% --- p. 403 ---
\opage{403}When the second theory would have distribution determined by each one's reasonable needs, it is clear that it
cannot be the single individual's own reason that is to decide whether his needs are ``reasonable.'' It must be the reason
of those who exercise the power of the state, and the single individuals are thus made minors. Even if a regulation of needs
is necessary, since the state must after all determine its officials' standard of life, it will nevertheless be just as
little fortunate here as in production if this regulating does not have a free unfolding and adaptation beside it, so that
comparison and choice can take place on as large a scale as possible.

10. A collection of individuals who have the right to decide neither what capacities nor what needs they possess, but who
must put up with being cut to a scheme laid down by authorities, --- such a collection of individuals is a mere
\textit{mass,} not an organized society. And it is indifferent whether one imagines the future leaders of society as
geniuses or as idiots, --- whether one imagines them as absolute monarchs and dictators or as issuing from the election of
the whole people. What gives life its value, the free unfolding of capacities and impulses, has in all cases fallen away.

The need to decide for oneself what capacities and needs one possesses, and which of them deserve to be developed and
satisfied, is not merely an egoistic need. It is, as we have sought to show, the condition for there being productive
forces in society that bring forth something new and do not go on moving in the old ruts. Even customary activity itself
is carried on rightly only when the individual determines himself to it. But when it is a matter of striking out on new
paths, the satisfaction of following one's own initiative toward what one considers good and useful is often the only
reward given. Great inventors often care nothing at all for the advantage they themselves could have from their inventions,
and they are often enough cheated of it too. What embitters their life and makes their fate hard is the obstacles that
stand in the way of the free use of their powers in the direction they wish. How will it go, then,
% --- p. 404 ---
\opage{404}when \textit{all} means of work are seized by the power of the state! Whence are the means then to come for the
private experiments to which culture owes so much, if it does not owe them everything?

--- \textsc{Schäffle}, to be sure, has maintained that one does socialism an injustice in thinking that it necessarily
abolishes all free movement and all free disposal of material goods. He distinguishes sharply between means of enjoyment
and means of production and seeks to show that socialism abolishes property only as a means of production, not as a means
of enjoyment. Over the means of enjoyment that fall to us in the socialist society we can freely dispose. We can accumulate
income for use in our own private purposes or to give gifts and help to others\footnote[1]{\textit{Socialismens
Kvintessens.} Danish transl. p. 84---97. --- Cf. also \emph{Gustav Bang}: \textit{Den socialistiske Fremtidsstat.}
Copenhagen 1903. p. 78.}. But it is surely no great scope that free movement gets here. The consumption of what is saved
takes place with great speed when I am not allowed to make it fruitful by letting others who lack it have it. If these, in
return for being allowed to use what I had saved, gave me a certain compensation for the time during which I let them have
the use of it, then I could use this time for pursuits that were not immediately productive, although they demand serious
and lengthy work. And this time will be all the better used since it can perhaps be used for works whose value no one else,
and especially not the authorities ruling in society, acknowledges, either because it is not seen that they answer to real
needs, or because they answer to needs that are first to be awakened, and whose awakening will make life richer. In the
socialist state, which forbids all interest, there will be room only for such pursuits as the power of the state will
favor. Not only will individual production fall away, but individual enjoyment will be kept within certain narrow limits.
Freer and distinctive tendencies and endeavors will not be able to form when one lives at the discretion of the power of
the state. Interest is, to be sure, an income that the owner gets without himself producing it at the moment;
% --- p. 405 ---
\opage{405}but it has the social significance of making activities other than the immediately useful possible and of
favoring thrift. The thrifty man now knows that he can manage not only to secure his own existence, but also to further, in
a lasting way, interests and endeavors he values. Interest can naturally be misused for idleness; but every right of
disposal can be misused: the state's right of disposal can, after all, be misused too. --- Just as socialism presupposes a
mankind that can furnish perfect rulers, so too it presupposes a mankind whose zest for work and inventiveness are not
weakened because its initiative is checked and because its needs are determined by others.

Through private initiative fresh blood comes into social life. The individual must therefore not be absolutely excluded
from producing. It would be absurd if he were to be allowed to consume his savings in debauchery, but not to use them as
means of production. This barrier to the freedom to try and to venture will never be endurable. And it will be not only a
barrier to the individual's freedom, but also to society's own development. Bureaucratic and parliamentary complacency will
check progress in an intolerable way. Not only material but also ideal culture will suffer by it. Through the law of
inheritance, and through the state's taking over private undertakings of a generally useful character when they have for a
certain number of years benefited their founders and when they are suited to be run by the public, it can be prevented
that out of the private enterprise and new formation so important for the whole of society there should develop a
hereditary nobility. The matter could be arranged by analogy with the newer arrangements of authors', artists' and
inventors' rights.

Speculative socialism would stop up the source of goods so that distribution can be all the better. But it thereby gets
entangled in a contradiction, since in the end there will be no goods to distribute if the source of progress is stopped
up. Utopian socialism was more consistent: it presupposed not only a limited or closed state; it also required needs to be
directly regulated. So for example Plato and Campanella.
% --- p. 406 ---
\opage{406}Newer socialism makes individual freedom the concession of permitting private property in means of consumption
and enjoyment, and would only forbid their use as means of production. The further social experience advances, the more
surely will the adherents of strict socialism too come to see that social ethics must for its own sake demand and further
individual freedom, not merely for consumption and enjoyment, but also for work and production. At any rate the abolition
of freedom of production as a solution of the social problem presupposes psychological and social conditions so different
from those now present that it would be dogmatic presumption to pronounce decisively on what they would make possible or
not possible.

11. It is a main socialist proposition that labor is the source of all value and all culture. The Gotha Program begins with
this proposition. But it contains a certain ambiguity, which also lies in the word labor. There is, after all, not only
material but also spiritual labor. And even if in considering material culture we think especially of material labor, it
has nevertheless appeared that the altered and improved position that labor in the service of material culture occupies in
modern times hangs together with modern industry's being so largely an application of the results of modern science.
Material labor thus presupposes spiritual labor here; the work of the muscle presupposes that of the brain. The
experiments, thoughts and plans to which spiritual labor has led have given countless material workers plenty to do. One
will therefore not be able to carry through a social order if one begins, like the Social Democratic programs of Gotha
(1875) and Ghent (1877), by setting the working class (by which is meant the material workers: the Ghent program calls it
the proletariat) in opposition to all other classes. This may be justified in the heat of battle; but if the feeling of
opposition and separation from the rest of society is nourished too strongly, the only way to better conditions is blocked.

Not physical labor alone is the source of wealth and culture. The greatest work of culture has been carried on in the
spiritual domain.
% --- p. 407 ---
\opage{407}Naturally material needs must be satisfied in order that one may work. And much spiritual work needs material
help. The organist would be helpless without the bellows-treader. But it is not the bellows-treader who is the ``source'' of
the music. --- Socialism too puts free thinking and inquiry and public instruction on its program. It must then acknowledge
that the whole spiritual atmosphere in which the worker lives is of the greatest significance for him. Modern philosophy and
social science led first to the abolition of the old guilds and then of the provisions that prevented the workers from
forming lawful and legally protected associations. The most solitary investigator can send out thoughts into the world
which, through their influence on the general view of life and on public opinion, can determine the course of culture to a
far higher degree than the material labor of many thousands. This is simply so, and no programs can prevent it. If future
development is to go in a sound direction, one must strive to diminish the distance between material and spiritual work.
But such an endeavor cannot succeed when the working class's opposition to \textit{all} other classes is stressed as sharply
as usually happens. --- It must be conceded that the ``other'' classes, the classes that until now have almost exclusively
been so placed that they could undertake spiritual work, have not always shown the conduct toward the working class that
their duty enjoined. Prejudices of various kinds, intellectual arrogance and lack of sympathy have prevented them from
acknowledging its right. Here lies the most important guilt for the disproportion. But here we do not ask first and
foremost about guilt; we examine how great the opposition between the classes in society actually is and, from an
ethical-social point of view, ought to be.

Those Social Democratic programs really conflict with \textsc{Marx}'s doctrine, according to which all class difference is
to be brought to an end, just as the difference of estate between nobility and burghers fell away with the Revolution. The
working class does not yet stand over against the other classes as the third estate at the outbreak of the Revolution stood
over against the ``higher'' estates of the time. The class that is now scornfully called the bourgeoisie had created
% --- p. 408 ---
\opage{408}modern industry and commerce, science and art, and out of it, in the Germanic countries, the Protestant movement
of freedom had come. In spite of all the admiration that must be felt for the development of the working class in the last
century, it cannot be said that it has now reached a point like the one the third estate had reached a hundred years ago.
Its development is not yet complete. That is not its own fault; for the conditions have only recently come to be present.
But all the less is it justified, then, to stress the class opposition as strongly as Marxism does. In any case the opposed
classes have much to learn from one another. And the working class that will, it is to be hoped, one day come to embrace
all will not be the present, imperfectly developed working class, which not only because of its interests, but also because
of its less perfect development, stands in opposition to the other classes of society. Imperfectly developed it will remain
as long as it still --- because of the class struggle --- must appear as a mass, without any free individualization being
able to take place within it. As long as, for example, proportional representation is regarded as a purely aristocratic
institution, it is clear that understanding of the significance of personal independence is lacking. The dogmatic quarrels
within German Social Democracy, which have led to the expulsion of personalities whose intellectual independence could not
bow to the one saving truths of the program, point in the same direction. If it is true that the Social Democratic party
leadership in Germany literally has greater power and influence over its fellow members than the state authorities over
their subjects,\footnote[1]{B. \emph{Friedländer}: \textit{Die vier Hauptrichtungen der modernen socialen Bewegung.}
Berlin 1901. I, p. 302. --- Cf. also \emph{Bruno Wille}: \textit{Die Philosophie der Befreiung.} Ch. 13 (Parteiherrschaft).}
this bears witness to great spiritual immaturity within the class. The theoretical anarchists do a useful work here by
their criticism; in return they are regarded as apostates and counted among the ``bourgeoisie.''

12. The problem, however, will always be posed anew, on account of a circumstance that socialists as a rule push aside,
% --- p. 409 ---
\opage{409}namely the natural instinct that brings about a tendency to an increase of population stronger than always
corresponds to the increase of the means of subsistence. Even if Malthus should have exaggerated this tendency, there is
evidently a force at work here that will always anew lead beyond any state of equilibrium one could imagine brought about.
(Cf. XXV, 2). If an arrangement were successfully made by which everyone could look the future confidently in the face, this
confidence would show itself among other things in numerous families being founded, so that there would soon again be
superfluous labor, more hands than there was perhaps work for, and more mouths demanding food. This is a simple consequence
of the same cause that brings it about that even now the number of marriages rises, not only when the price of grain falls,
but already when there is merely a prospect of it, or when a hopeful mood prevails\footnote[1]{``After a good harvest the
number of weddings and births usually increases considerably, and likewise conversely decreases after a bad harvest. In
the first case it is still more hope than actual possession that incites to the founding of new families, which is why the
strongest increase is found not at the absolutely lowest prices of grain, but at those that stand in the most striking
contrast to a preceding bad year.'' \emph{Roscher}: \textit{Die Grundlagen der Nationalökonomie.} § 240.}.

The social state of things is determined at any time by the relation between the growth of the population and the degree
to which increased cleverness and energy can increase the productivity of the earth. Where different forces work together,
a rhythmic movement will set in, now the one, now the other tendency gaining the upper hand. An existence not subject to
such a rhythmic alternation we do not know, and do not understand. It is the different forces, and the way in which they
break against one another, that make life a struggle and cause pain, especially each time the amplitude of the
oscillations grows greater. It is not said that human nature will constantly bring with it the same strong tendency to
increase; but a decisive change in this respect lies too remote for it to acquire any ethical significance for us. Through
the hard struggle with the given conditions
% --- p. 410 ---
\opage{410}human nature has developed to its present stage; it is this that makes the strict demands of self-control,
prudence and sympathy necessary; and the conditions will, as far as we can see, not change for the time being. (Cf. with
what is developed here XI, 10; XVII, 2 and XXV, 2).

We therefore also see that when a social layer has worked its way forward to more favorable conditions, to a higher
standard of living, a new layer forms beneath it whose standard of living needs to be raised. After the agrarian reforms in
Denmark at the end of the 18th century, the number of cottagers and agricultural laborers grew on account of the parceling
out of land and the increased need for labor --- but grew to a degree that posed a new social question in place of the one
that had just been solved. ``Even for the cottagers who had land the position was by no means always good. And the cottagers
who were landless or were lodgers, and that was about half of the whole number, nearly always lived in very poor
circumstances; they had gone back rather than forward''\footnote[1]{\emph{Falbe Hansen}: \textit{Stavnsbaandsløsningen og
Landboreformerne.} I. p. 139.}. In a similar way the trade unions and consumers' associations at first embraced only the
most intelligent workers; the trades that require no further special training were not organized. The labor movement
threatened for a time to end in the formation of a labor aristocracy which --- as has been said --- was regarded by the
``unskilled'' workers with feelings like those for the House of Lords in the parliamentary world. A movement then had to get
going to organize the ``unskilled'' workers\footnote[2]{\emph{Sidney} and \emph{Beatrice Webb}: \textit{History of Trade
Unionism.} p. 373---393.}. Moreover, unfavorably placed workers will naturally seek their way to districts and countries
where a more favorable position has already been attained for the working class. There thus takes place in Europe a constant
migration of crowds of workers from east to west and from south to north; --- and far out on the eastern horizon the
industrious and frugal workers of the yellow race are a threat. --- The problem will follow the race on its way for long
ages; a solution once for all is improbable.
% --- p. 411 ---
\opage{411}13. The socialist theory is a decidedly idealistic theory, insofar as it rests on the conviction that the human
will is able to clear all difficulties out of the way of the formation of a harmonious human society\footnote[1]{The ideal
of the future appears in an especially idealistic form in the theory of \emph{syndicalism}. Not only capitalism, the state too is
to fall away. An economic federation of free workshops is to replace the state. Joy in work, fulfillment of duty and
sympathy are to take the place of authority and compulsion. Through artistic and moral development the workers are to be
matured to be able to realize this ideal. Syndicalism places itself in opposition both to the usual trade unions and to
socialism, which would make the state the sole employer. Cf. (besides the works of H. Scavenius named above (4))
\emph{Challaye}: \textit{Le Syndicalisme revolutionnaire.} (Revue de Métaphysique et de Morale. 1907). \emph{Paul Louis}:
\textit{Le Syndicalisme contre l'état.} 1910. --- In its rejection of the state syndicalism recalls utopias of Sibbern's
kind (in contrast to those of Plato, More and Campanella); in its stress on sympathy it recalls philanthropic socialism; in
the significance it ascribes to the strike (in the end the general strike) as a means of action it is, in spite of all
opposition, akin to Marxism. (1913).}. It regards labor, after all, as a source of all wealth and culture and disregards
all natural causes that could further or check the flow of the source. Through this idealism, which does not let its hope
and its enthusiasm be weakened by anxious regard for the naturally and historically given conditions, socialism is akin to
the best in the thought of the 18th century, although on the other side by its anti-individualism it marks a reaction
against the 18th century. However many theoretical mistakes and illusions one may be able to point out in socialism, it
nevertheless appears in practice as one of the most significant ethical-social movements of the present. It has known how
to awaken and inspire the workers; it has directed their thoughts toward ideals and tasks that embrace far more than the
narrow circle in which the isolated individual drive for self-preservation moves. It has set millions of human beings a
great common ideal that lifts the individual up above the narrow circle of interests, has directed the gaze toward a
picture of the future that shows the golden age before us, not behind us; --- and it has at the same time managed to get a
firm hold on purely practical and material interests, to use these as
% --- p. 412 ---
\opage{412}motives that only lead the individual out beyond himself. Without a great picture of the future no social
movement can go forward. Sometimes --- so a few years ago in England --- the trade union movement was coming to a
standstill, and it was then the enthusiastic faith in the socialist ideal that awakened the impulse to advance again.
Development cannot be guided from above; even if fruitful germs of thought and impulses can come from the other social
layers, self-activity must above all be awakened, and for that ideals are needed that stand in a natural relation to the
need that is felt and to the intellectual and moral horizon that lies open. If the other social layers do not like these
ideals, and if they find the need too low and the horizon too narrow, then it is their task to change conditions of life
and social circumstances so that the need can be raised, the horizon widened, and with it the ideal too take on another
character. The socialist movement is a continuation of the great movement of freedom which many liberals expected to be
completed with the Declaration of the Rights of Man, with parliamentarism and free trade. It is the positive side of the
great process of emancipation that begins in the 18th century. The harmony of interests that older economists and
utilitarians believed in, it seeks to make a fact. And at the same time it would republicanize and democratize the feeling
of humanity which the old society knew essentially only in the form of the master's patriarchal relation to his dependents.
It would develop love of mankind into justice. The shadow sides and illusions in the picture of the future can be
corrected in the course of further development. Socialism must learn to distinguish between ideal and reality. The ideal
need not lose its driving power because one does not veil and overlook the real circumstances. And it must learn above all
that if it has a future, this comes essentially through free labor and free association, even though the state may very
well, to a far higher degree than we are now used to thinking, come in to protect, help, equalize and educate.

14. As already touched on (9), the power of the state intervenes uninterruptedly in social development, even where one
does not become
% --- p. 413 ---
\opage{413}clearly conscious of it. There is not a single side of the political order (constitution, administration,
finance, law and the military, church and education) that does not in one way or another become of determining influence
on the ordering of the conditions of labor. It would already be a great advance if one became more clearly conscious of
this, so that this very essential side of political questions were drawn forward more. It is, after all, the social
significance of political questions that first gives them real interest. Where this significance is wholly lost from
sight, political strife is either a purely personal strife for power or a strife over mere formalities. --- We shall here
name some points where the state, without violating the principle of freedom, can accomplish much as regards the ordering
of labor and the distribution of its yield, and where it has also already intervened --- by an unconscious or involuntary
socialism.

a) First of all the demand must be brought forward that simple justice be shown. The great thieves are still in many
respects more fortunately placed than the small ones. Movements in the workers' circles are looked upon with suspicious
eyes which in other circles of society one would not put obstacles in the way of. The resistance that has been offered to
the workers' right to personal freedom and their right of assembly and association has not been able to dispose them
kindly toward the ruling classes. Adam Smith complained in his time that masters might combine to lower wages, while workers
might not combine not to work below a certain payment\footnote[1]{\textit{Om Nationalvelstand.} Danish transl. I, p. 202
f.}. German legislation likewise until quite recently acknowledged the employers' guild associations, but not the workers'
corresponding associations, and the much talk of restoring the old guilds really aims at making the workers dependent on
the employers\footnote[2]{L. \emph{Brentano}: \textit{Die gewerbliche Arbeiterfrage.} (Schönberg's Handbuch. 1st ed. I). p.
931 f. 970.}.

b) The state has recognized it as its duty to protect the workers' freedom, health, safety and morality against the
arbitrariness of
% --- p. 414 ---
\opage{414}the employers. After a hard struggle --- first in England --- the so-called factory legislation went through,
which regulates safety and health measures in factories and mines, determines the length of working time, fixes limits and
closer provisions for the work of women and children, forbids the payment of wages in taverns and requires them paid in
cash\footnote[1]{\emph{Gneist}: \textit{Das Self-Government in England.} 3. Aufl. p. 314 f. --- L. \emph{Brentano}:
\textit{Die gewerbl. Arbeiterfrage.} (Schönberg's Handbuch I). p. 973. --- K. \emph{Marx}: \textit{Das Kapital.} 2. Aufl.
I. p. 224---314. --- Characteristic is the resistance the utilitarians offered to factory legislation, except insofar as it
concerned the protection of children. \emph{Leslie Stephen}: \textit{The Utilitarians.} III. p. 175---177.}. --- Just as it
has turned out that free labor is more productive than slave labor, so too it has turned out that the more favorable
conditions which this legislation, carried through on the English model in other countries too, has procured the workers,
far from harming productivity, have precisely increased it. It has not only, as even \textsc{Karl Marx} admits, brought
about ``the physical and moral regeneration of the factory workers''; but the greater freshness and strength with which
work could now be done often even increased the productivity of labor. The employers themselves found their advantage in
it and gave up their earlier resistance to what they --- characteristically enough --- had called an encroachment on their
personal freedom! --- The restriction of working time in particular has been a great good, and that for employers and
workers alike. As the English factory inspectors say in a report: formerly the master had no time to think of anything but
money, and the workers no time to think of anything but work. The long working time made the workers purely physical
beings, all their free time going to sleep and rest to get strength to set to again. With the successive limitation of
working time one approaches \textsc{More}'s Utopia, where ``the chief end of the constitution is to regulate labor by the
needs of the people, so that time may be left over for the development of the mind, in which the Utopians think the
happiness of life consists.'' % translated from Høffding's Danish
\textsc{Fichte} has rightly said that a people's true wealth
% --- p. 415 ---
\opage{415}consists in the free time at its disposal\footnote[1]{\textit{System der Rechtslehre} (1812). Nachgelassene
Werke. II, p. 543: ``Das Eigenthum des Ganzen [ist] die Musse, die allen nach vollbrachter Arbeit bleibt.'' --- Cf. L.
\emph{Brentano}: \textit{Die Arbeitergilden der Gegenwart.} II. p. 356: ``Die Frage nach der Länge des Arbeitstages ist
eine Frage nach dem Stand der Civilisation.''}. For the time that can be spared from work for material needs can be used
for a higher and freer development, for noble enjoyment of life, for working in the service of ideal culture and feeling
oneself to be more than a wheel in a great machine. --- It depends, naturally, on how free time is used, and one cannot
wonder that it is not straight away or always used as well as possible. The impulse must be awakened and means procured for
higher development. Here education, both that directed by the state and that provided by free forces, becomes of the
greatest significance, and not least everything that can be done to spread and satisfy the sense for the beauty of art and
of nature. Through participation in political life the feeling of civic duties and rights, and of working in the service of
the whole of society, is developed. One cannot wait to limit working time until a sense for using free time well has been
attained; for this sense first arises when the free time is there. And when it is complained that the workers do not know
how to use their free time, there is far more reason to complain of the way in which the well-to-do classes use theirs. When
one is not used to having free time, it is no wonder that one has not learned to use it; but it is sad to see how empty and
low is often that with which the so-called cultivated and higher classes fill their ample and often undeserved free time.
The workers must precisely have more free time in order that the way in which they now often misuse free time may fall
away. The eight-hour working day is demanded in order that ``Blue Monday'' may fall away\footnote[2]{Robert \emph{Seidel}:
\textit{Der achtstundige Arbeitstag.} Zürich 1896. p. 4. --- The division of eight hours' work, eight hours' sleep and eight
hours' free time was already proposed by \emph{Comenius}, and later by \emph{Hufeland}.}.

c) The state can further exercise an equalizing influence on social antagonisms through the ordering of finance and
% --- p. 416 ---
\opage{416}taxation. A new period in the history of taxation has even been dated\footnote[1]{Adolf \emph{Wagner}:
\textit{Direkte Steuern.} (Schönberg's Handbuch. 1st ed. II). p. 169. 259. Cf. also \emph{Jhering}: \textit{Der Zweck im
Recht.} 2. Aufl. I, p. 533.} from the time when this equalizing and distributing effect of the ordering of taxes began to
be stressed as a main point of view. It is the ``social-political'' theory of taxation that applies this point of view
beside the merely financial one. \textsc{Henry George} and his adherents even expect the complete solution of the social
question through the taxation of ground rent. Without the state, ground rent would not exist, and it is therefore the
proper object of taxation\footnote[2]{Cf., however, \emph{Westergaard}: \textit{Nationaløkonomien i Hovedtræk.} (1908). p.
84---89.}. --- To this belongs also a heavier taxation of inheritance and the abolition of inheritance through collateral
lines.

d) A step further leads the question whether the state is to set up insurance institutions for the workers. If it is to be
compulsory for the worker to make use of such institutions, he is easily robbed of his personal independence. For since the
state cannot guarantee him steady work at a certain wage, a fluctuation in conditions can make him unemployed and thereby
unable to pay contributions, whereby he in turn loses what he has paid in before. It will then be in the employer's power
to impose conditions on him which he would not put up with if he did not fear losing his earlier
contributions\footnote[3]{L. \emph{Brentano}: \textit{Die gewerbl. Arbeiterfrage.} (Schönberg's Handbuch. 1st ed. I). p.
985 f.}. The state will surely here, as on so many other points, work best by working indirectly, so that it supports and
controls the organizations that have developed through the workers' own free union.

e) There are many activities, in the domain of production as well as of exchange and communication, that state and
municipality have already taken over and will be able to take over gradually. Where the limit lies in this respect only
advancing experience can decide. Here precisely empirical
% --- p. 417 ---
\opage{417}socialism has its great significance. But in order to be able to use experience one must have sure facts. And
here is a domain where the state can do incalculable good: by procuring an exact statistics. The North American free
states have led the way in this respect. Both in the single states and for the Union as a whole institutions for labor
statistics have been set up. In the law by which in the year 1888 \textit{the United States Department of Labour} was set
up, its task is stated to be to collect, and to spread among the population of the United States, information ``upon the
subjects connected with labor, in the most general and comprehensive sense of that word, and especially upon its relation
to capital, the hours of labor, the earnings of laboring men and women, and the means of promoting their material, social,
intellectual, and moral prosperity.'' % the Act's English, as far as Høffding's Danish allows
The social legislation and the public debate in the United States have already reaped benefit from this
measure\footnote[1]{\emph{Ernst Beckman}: \textit{Den sociala frågan och stastistiken.} (Nordisk Tidsskrift. 1893).}.

The closer development and justification of the points named here (and of others that could be added) belong to political
economy. They are merely cited here because they all presuppose an ethical-social conception of the state and its
activity, and a recognition that state help and self-help do not exclude one another. What matters is only to establish
the right relation between them. The domain of free forces is the really productive one, the one where initiatives thrive;
the state can only afford protection, form and material support to what has worked its way forward of itself. The relation
between state help and self-help ought to be precisely the reverse of the one that speculative socialism, in remarkable
agreement with bureaucracy and absolutism, lays down. And since experiences are best made in smaller circles, municipal
ordering of production\footnote[2]{Cf. F. \emph{Pio}: \textit{Kommunesocialisme.} (Tilskueren. 1903).} --- as empirical
socialism too maintains --- will have many advantages over ``state socialism'' proper. On the socialist side too
% --- p. 418 ---
\opage{418}people have begun to get away from thinking of the state --- in the sense in which we now take this word --- as
the form of society that is to take over the whole of production. Organizations of the most different kind and extent are
envisaged, according to the kind of task, and a cooperation between these different associations\footnote[1]{Gustav
\emph{Bang}: \textit{Den socialistiske Fremtidsstat.} p. 64.}. This view, which --- be it said in passing --- leads away from
Marxism proper, must probably be widened so that free individual initiative gets its place beside the self-government
granted to the small and great organizations. There must be a constant interaction between individuals, the free cultural
communities and the state if development is to be sound and vigorous.

15. Speculative socialism does not absolutely abolish private property, but restricts it to means of use and enjoyment.
There would be a complete contradiction between private property and socialism only if one took the right of property to be
unconditional. But an unconditional right of property cannot be pointed out anywhere in history. The idea of the
unconditional right of property is a fantasy that individualism set up against the whims of an arbitrary power of the state.
At all times the power of the state has intervened in the individual's possessions when common interests seemed to require
it. The limits of these interventions are shifting, change with historical circumstances.

The very concept of \textit{private property} in the extent in which we now take it is relatively new. At the primitive
stages of culture community prevails. The hunting grounds are common to the whole tribe, and among agricultural tribes each
family is assigned its lot yearly, which the members of the family cultivate in common. It would naturally be wrong to think
of this community as a communism in the proper sense; it means only that the need for a carried-through distribution has
not yet arisen, and that the individual works together immediately with the family or the horde, has as yet no occasion to
try to separate his existence from it. On the other hand, by
% --- p. 419 ---
\opage{419}characterizing ``primitive communism'' as an insurance, primitive community as ``a spatial being-together of
individual isolated interests''\footnote[1]{C. N. \emph{Starcke}: \textit{Samvittighedslivet.} I. p. 235.}, one would carry
over modern concepts, which owe their origin to the mechanics of individualism, to the old social conditions. One cannot
compare the old community with a modern joint-stock company, which is precisely the expression of hitherto isolated
interests meeting in order to work together in a purely external and impersonal way. As regards landed property in
particular, it is of significance that social organization in clan and family precedes in time settlement and permanent
possession of land\footnote[2]{B. W. \emph{Leist}: \textit{Alt-Arisches Jus Civile.} I. Jena 1892. p. 512 f.}. The first
private possession proper is the game the individual brings down, the fruit he gathers, the tools and clothes he procures
for himself. Movable property thus precedes immovable. Only later does land pass over into private property. The necessity
of regular cultivation of the soil carried through according to a comprehensive plan, the division of labor, the awakening
self-esteem and urge for independence in the single individuals (cf. XI, 14 and XXIII, 3), all these were causes that were
bound to bring about an individualization of property. Primitive culture most often knows only a common right of use within
the tribe or the family, and a right of property only where tribes or families come into collision. Only with the
development of the state does individual property arise as the possession of material goods acknowledged and protected by
the power of the state. But this very acknowledgment and protection, which is an indispensable moment in the concept of
property, shows that the right of property is not unconditional. The power of the state always attaches definite conditions
to its acknowledgment and its protection. It sets limits to the power of disposing of possessions as one likes. One may not
build on one's ground, cultivate one's land or bequeath one's money as one oneself wishes, and one contributes greater or
smaller amounts of what one possesses to common ends. The power that by its acknowledgment and protection makes mere
possession into property is the same that imposes limits and duties;
% --- p. 420 ---
\opage{420}these thus by no means appear as encroachments on an originally absolute right of property.

It is in the end ethical considerations that ground both the acknowledgment and the restriction. --- A determination of the
limits of the single individuals' disposal of material goods is, in the first place, necessary in order that the state may
be able to solve its most elementary task, namely to maintain peace and security. For this it is not enough to fix the
boundary between the scope of tribes or families. As development advances, the state has to do with the single individuals
and their mutual relations, and must have means of establishing where encroachment has taken place between them. If, for
example, agriculture begins to develop in a tribe that has hitherto been nomadic or lived by raids, it will at first be the
affair of single persons. But the more people strike out on this path, the more easily collisions will occur, several
wishing to cultivate the same piece of land. The mightiest would seek to seize the largest pieces of land; but against this
the interest of the great multitude would then rise, since they were used to using the common land as pasture for the whole
tribe's cattle. Only a definite distribution and demarcation on the part of the power of the state could here procure peace
and security\footnote[1]{Cf. D. M. \emph{Wallace}: \textit{Rusland.} Danish adaptation. II. p. 57---61.}. --- Next,
acknowledged and protected possession is a condition of steady, independent and self-sacrificing activity for the ends of
culture. To possess property is to hold a social office of trust. The individual's feeling of independence and power is,
to be sure, also satisfied by it; but social ethics cannot regard as sufficient the justification that merely regards
property as the extension of personality out into the surrounding world. Such an extension or spreading is naturally
justified where it does not lead to strife and collision. But its positive, social significance first comes forward when
the individual's feeling of self and of power is taken into the service of society. Not from the concept of the individual
personality in and for itself, but from the social necessity of having as many free, active points of departure as
possible for
% --- p. 421 ---
\opage{421}the activity of culture, does the ethical character of private property follow. The individual has not even a
self-evident right to what he produces and acquires by his own work; for this work is not a creation out of nothing; it is
not conditioned merely by his own capacity and will; --- that it could be carried on, and could become fruitful for him, is
due to the protection and help of society. As he must thus feel with gratitude and reverence what he owes to society, so
too, when he sees things from the ethical point of view, he must regard his possessions as a means of carrying on an
activity in the service of society. He is just as much an official as those whom the state specially appoints, and has his
ethical responsibility for the use of what he owns. Ethically regarded, therefore, there are also narrower limits to the
use of property than there are from the point of view of the outer legal order. The material goods in the individual's
possession he is to make as fruitful as possible --- not only for himself but for the race. They belong to him now, but
they may later fall to the share of others. He therefore has, for example --- even if legislation cannot always prevent it
--- no right to exhaust the soil by robber-farming and thereby diminish its productive capacity for the future. That is an
egoistic recklessness toward the following generations. For the sake of a momentary personal advantage the lasting good of
the race is pushed aside. Often, however, it is plain ignorance that lies at the bottom of it, and we then have an
interesting example of how increased insight can increase ethical duties. It was only Liebig's doctrine of plants' need for
mineral substances that proved the harmfulness of robber-farming and demonstrated the solidarity of the race with regard to
the use of the soil.

Much that is now private property can be conceived one day to become common property. It might, for example, be conceived
that it turned out expedient and possible to bring all landed property under the state, so that ground rent could benefit
the whole of society. If that happened, it would be no encroachment; it would be in consistency with the line of thought on
which the ethical justification of private property rests. What speaks against this thought is that a greater yield is
probably got out of the land by private cultivation than by state operation.

% --- p. 422 ---
\opage{422}16. Private property in material goods makes \textit{commerce,} exchanging activity, necessary. In the socialist
society there would be no commerce: from the seat of the central power (in the whole state or in the single municipality)
the products and means of enjoyment would be distributed round to the single individuals, but there would be no exchanging
activity going on among the countless single individuals. A centralization of exchanging activity will indeed also be in the
interest of society, since superfluous middlemen can thereby be avoided. The strong flowering of the consumers' associations
(XXVI, 6) shows that there are far more superfluous middlemen than one would have thought in advance. Since society does not
exist for the sake of commerce, but commerce for the sake of society, an endeavor like that of the consumers' associations is
fully justified and is an essential member of empirical socialism. But one ought not to judge the significance of commerce
by the superfluous middlemen. Commerce rests on the activity of free, individual forces, which here show themselves every bit
as valuable as in production. For the one who holds goods he does not himself use, the point is to find a place where there
is need of them. The commercial spirit consists far from merely in the art of selling for more what one has bought for less;
its most essential side is the capacity to find where there are needs to satisfy. Productive work is carried out when a
place is found where a product can be of use, and when the product is moved there from a place where it is of no use.
Socialist theorists have not been willing to acknowledge this. And yet it is just as clear as that it is productive work the
farmer does when he drives his grain to the merchant instead of letting it lie in the barn or stand in the field. It is
naturally self-interest that at first sets that capacity in activity, and the unfortunate sides of commerce and of those
engaged in it hang together with this. When, especially in earlier times, the Greek philosophers, the Christian Fathers of
the Church and the scholastics looked down on commerce as an immoral way of earning, it was because it was thought to
produce pettiness, greed for money, faithlessness and
% --- p. 423 ---
\opage{423}heartlessness\footnote[1]{\emph{Plato}: \textit{Laws.} Fourth Book. p. 705 f.}. The great ethical-social significance
that commerce has in binding human beings together through their interests was overlooked. Commerce is an expression of
the fact that the single human being and the single people is not sufficient to itself, which is why they seek
supplementation from free impulse. The socialist state, which is to exclude commerce, must either (as \textsc{Fichte}
conceived it) be a closed state, economically sufficient to itself, or (as \textsc{Rodbertus} conceived it) a world state
embracing all peoples. Until such a world state has developed, commerce must be a necessity, and commerce will even be a
main means of its development. And as regards the closed state, this was according to \textsc{Fichte} to stand only in
intellectual, not in economic, intercourse with other states. Science, not commerce, was to condition the connection between
human beings. But the ideal connection depends here precisely on the material; the latter paves the way for the former. When
the bond has once been tied for economic reasons, a further-reaching connection can thereby have been prepared. Already a
lasting commercial connection presupposes a relation of trust that is built not only on egoistic interests but also on
characters. Lively interaction and firm bonds between individuals of the same people and between different peoples are tied
in this way. Commerce has therefore played a great part in the history of culture. It has initiated connections and brought
about mutual acquaintance between human beings who would otherwise not have come into contact with one another. It has often
widened the narrow horizon within which the gaze moves, before the need for and the possibility of interaction with more
distant parts of mankind had arisen. Especially where there is sea trade with distant lands, a wider view of life and its
conditions is furthered; manifold new thoughts and plans stir; new customs and institutions become known, and a liberation
from the accustomed and the handed-down is initiated. With the ships and the caravans travel the invisible passengers: through
the material
% --- p. 424 ---
\opage{424}connection, human beings' feelings and thoughts too are brought into connection. Material exchange thereby also
gains significance for spiritual life. (Cf. XXIV, 1).

\grouphead{\textit{2. IDEAL CULTURE.}}

\chapterhead{XXVII.}{MATERIAL AND IDEAL CULTURE.}

1. Material culture aims chiefly at procuring a great system of means. But means belong with ends. Material culture
therefore always points beyond itself. The social question arises, after all, precisely because the means seem to spread
at the expense of the ends, a great part of mankind seeming to be bound to keep life going but to do without what first
makes life worth living. We seem to groan under the weight of an apparatus that was brought about in order to make it
easier to live. The ethical inquiry into material culture concerned the possibility of getting beyond this evil and of
maintaining the fundamental thought that a human personality may never be regarded merely as a means. Already the
application of this fundamental thought leads beyond material culture, since it shows us that all culture is to serve
the development of personal life. In ideal culture, which consists in the free development of thought, imagination and
feeling, personality is more than a means; it is the personality's own powers that are here set in activity, and that for
no other reason than that there is an immediate need to use them. By virtue of a higher urge for self-preservation,
thought, imagination and feeling unfold and create distinctive forms in which they find
% --- p. 425 ---
\opage{425}vent and expression. Here there is no more distant goal to be reached: clear understanding, the living image of
imagination, intimate and deep feeling have their value in themselves.

To that extent the relation between material and ideal culture might seem quite simple: the former seems to stand to the
latter as means to goal. But where is the boundary to be set? Is there any development whatever of thought, imagination
and feeling that cannot in turn act back on material culture and ease the production of material goods directly or
indirectly? A constant circulation takes place between ideal and material culture. A higher intelligence, a more lively
imagination, a more intimate feeling changes the conditions of work and determines its direction. Free movements in the
domain of spiritual life awaken hope, boldness and energy to go forward in the domain of material culture too. And on the
other side material work is not always merely a means of producing a product; it can become a school for the will, a form
of exercise of the powers. If there is an advance in the conditions of material work, it must consist in this, that work,
even if it is not play, is nevertheless bound up with an immediate satisfaction in the use of the powers. Man is not only a
being of nerves, but also a being of muscles; he has just as much a need to use his bodily powers as to use his spiritual
powers, and the satisfaction of this need is more than merely a means, just as the satisfaction of the need for knowledge
is. If material work could always be bound up with such an immediate satisfaction, the opposition between material and
ideal culture would fall away on an essential point. History shows that at every stage of human development both ideal and
material culture are present, although they may be developed in different degrees and stand at different times in
different relations to one another.

In the end there is only one single valuation. Both material and ideal culture receive their value by being conditions or
forms of advance toward the greatest welfare for as many as possible. Everything that directly or indirectly works in this
direction is, ethically regarded, productive. The soundness of the development of culture
% --- p. 426 ---
\opage{426}depends on the right relation between material and ideal culture, between the work of securing the material
foundation of life and the work of developing thought, feeling and imagination. The ethical standard of wealth and culture
is free time, i.e. the time that can be used, and that really is used, for ideal work of culture. (Cf. XXVI, 14).

2. We have, however, not yet got further than that by far the greatest part of work aims at procuring the means to live.
The great labor spent on this is only in a very slight degree immediately satisfying and developing for the worker's
personality. And the free time we have, and which far too few of us have, we do not yet know how to use in the right way.

At the lowest stage of existence almost all activity aims at the maintenance of the individual and of the race. The time
left over from this is essentially time of rest, used to gather strength for new work of maintenance. It marks a
significant turning point when the time of rest can be used for activity that is not necessary. Life can then get a
natural rhythm, not merely of physical work and physical rest, but of physical work and spiritual work. On this rests the
significance for the history of culture of the festivals of the Greeks and the Sabbath of the Jews. The latter in
particular is of great significance, since it seems to have fixed the length of the rhythm natural for most
people\footnote[1]{During the French Revolution an attempt was made to replace the rhythm of seven days by one of ten (the
week by the decade), which was found more rational. But it foundered for various reasons, not only religious ones. The
women in the suburb of St. Marcel declared that a worker could not work nine days in a row without resting. --- Ad.
\emph{Schmidt}: \textit{Pariser Zustände während der Revolutionszeit von 1789---1800.} III, p. 247.}. But the time of rest
is not always used in the best way. Just as the size of free time is a standard of national wealth, so the way in which
this free time is used is a standard of national culture. \textsc{Aristotle} says of the Spartans that their state
declined because they did not know how to use their free time rightly; they had procured power and wealth, but for lack of
culture of mind they had nothing to use
% --- p. 427 ---
\opage{427}these means for\footnote[1]{Polit. II, 9.}. Of a great many in the so-called higher classes the same still
holds. Precisely in these classes people often give themselves up to sensual enjoyments of the lowest kind, because the
impulse toward nobler enjoyments is lacking. One cannot wonder, then, that the same is the case to so high a degree in the
class of physical workers, where the strong exertion of working hours often leaves no energy over for the time of rest. On
the whole, however, a comparison between the kind of pleasures in older and in newer times seems to lead to the result that
coarseness has become less than in earlier times\footnote[2]{L. \emph{Felix}: \textit{Der Einfluss der Sitten und Gebräuche
auf die Entwicklung des Eigenthums.} p. 197.}.

3. For two reasons closely connected with one another ideal culture acquires a higher value than material culture. ---
It stands in \textit{more intimate connection with man's personality.} There is not the difference between work, means of
work and products of work that there is in material culture. We work here with our own mind, and what we produce cannot be
separated from our own mind, but belongs to it for ever. The social question, on the other hand, arose because labor power
was not always bound up with the means of work or in control of the product\footnote[3]{Cf. the closer development of this
in the essay \textit{Aandelig Kultur} (Mindre Arbejder. Third Series).}. --- And although work here is more bound up with
the personality itself, \textit{the individual nevertheless does not work for himself alone, but for the whole race.} If he
succeeds in producing new thoughts, new images or new forms for the life of feeling, he has thereby increased the race's
spiritual capital, and he does not himself become poorer because all share it with him. A material particle can at one and
the same time be part of only one single organism; that is why the struggle over material goods is so fierce. But a
thought can be produced in as many consciousnesses as you like. No isolated right of property holds here. And not only is
community more easily possible in the spiritual than in the material domain; the need for it is also greater. Without
interaction and common striving
% --- p. 428 ---
\opage{428}the higher spiritual life does not develop. It is not only finished thoughts, fancies and feelings that lead to
interaction between individuals; it holds every bit as much of those in the making. New thoughts and feelings can be
developed only when different consciousnesses influence one another. It is a matter, however, not only of the production
of new thoughts and forms of the spirit, but also of appropriation, independent appropriation. Through that too spiritual
life goes forward, and there is no sharp separation between production and appropriation, activity and passivity. Everyone
who has the possibility, that is to say: capacity, strength and time, to develop his life of thought and mood has not only
an individual but a social obligation to do so, whether appropriation or production is the predominant thing in him.
There is an inner art in the mind that is the essential thing in the great masters, and that can also be practiced by
those who learn from them. ``To discover something oneself is fine,'' says \textsc{Goethe}, and he adds: ``but when you
have grasped and valued what others have happily found, is that any the less to be called your own?'' % translated from Høffding's Danish (Goethe's distich "Erfinden, das ist schön...")
Open sense and ideal enthusiasm can pave the way for spiritual goods in every circle. No one is excluded. ---

Yet ideal culture too has its shadow sides. It often develops in a one-sided and unhealthy way; egoistic passions often
romp about on its preserves; and work in its service is often bound up with great sufferings. --- We consider each of
these three points by itself.

Where thought, imagination and feeling do not develop hand in hand with the will and muscular strength, ideal culture
easily takes on the character of comfortable enjoyment, play with thoughts and images, wallowing in fancies and feelings,
sentimentality and a morbid love of reflection. There is a kind of spiritual gourmandise that values thoughts and feelings
merely by their momentary taste, not by their real nutritive value. And there is a dead learning and an absorption in
specialties that can no less hamper one's own personal development. Such evils are to be compared with the stunted and
degraded figure that bodily life takes on under the disproportionate pressure of material
% --- p. 429 ---
\opage{429}work. Ideal culture too has its thralls, on whom it stamps the mark of one-sidedness.

Ideal culture, like material culture, develops through struggle, and struggle sets the passions in motion here too.
Ambition, lust for domination and jealousy have a fertile soil here. Sects and parties flourish, and the oppositions often
acquire here a depth they do not have in the struggle over material goods. For it is a matter here of something which is
bound up with the personality itself far more intimately than material goods can be; the judgment on a person's thoughts,
fancies and feelings is in an entirely different sense a judgment on himself than the judgment on his material work.
Ideal culture moreover develops the differences between human characters to a far higher degree than material culture.
The production of material goods lays claim only to elementary powers that are on the whole alike in all; but in the free
unfolding of spiritual life differences and nuances come forward that were formerly imperceptible. The differences of
individuality grow with growing ideal culture. But with them grow also the possibilities of strife and disharmony. There is
only one truth; the strife must then become fierce between those who, in spite of their great differences, each think they
possess it. And if only it were the truth itself that was at stake. But the honor of finding the truth often plays a
greater part than the truth itself, and that only one can get. --- If there is much in ideal culture that makes it a
mighty uniting force, there is also the possibility that precisely here great oppositions may arise.

Finally, from the intimate relation between spiritual work and the worker's personality there follows the possibility of a
kind of suffering that only the spiritual worker knows. It is not at all times that the inner work succeeds. There may be
inner resistance to overcome. There may be a disharmony between need and capacity, which is felt the more painfully the
more spiritual work is felt as a personal matter, and the more ideal a standard one applies to one's striving. It is a
matter not only of putting something into the world that can satisfy the demands of
% --- p. 430 ---
\opage{430}the market, but of satisfying one's genius. Often the spiritual worker may envy the material worker who leaves
his workshop with an easy mind when he has performed a certain quantity of mechanical work in the set time. The spiritual
worker has no such security, especially when that disharmony asserts itself. Many dark hours and times are called forth
thereby; an anguish can be felt at the stopping of spiritual life in us, which recalls the organic feeling of anguish when
breathing or the circulation of the blood is hampered. Thought will not clear, and feeling seems to have dried up. The
happiness of the good times is dearly bought with the anguish and doubt of the dark times. When the mind has experienced
what it is to be filled with a great ideal interest, it feels the slackness and incapacity of the bad times far more than
if it had not known such interests at all. To this is added the outer resistance, misunderstanding, mockery and coldness
that spiritual work so often meets, not least when it is distinctive and original. --- This suffering falls not only to
the lot of the great geniuses, but can also be the share of those who appropriate in a distinctive and independent way what
they have produced. They too can experience both the inner and the outer resistance.

\grouphead{\textit{A. INTELLECTUAL CULTURE.}}

\chapterhead{XXVIII.}{THE ETHICAL SIGNIFICANCE OF SCIENTIFIC KNOWLEDGE.}

1. The assumption that everything in the world is subject to a rhythmic movement, and that if there is progress, it at
any rate does not go in a straight line, is perhaps nowhere more clearly confirmed than by the judgments on the
significance of intellectual development.
% --- p. 431 ---
\opage{431}--- For the Greeks thought or reason stood as man's patent of nobility. Only he was a true human being who knew
the beautiful and the good. For the Greek philosophers thinking and knowing stood as the highest of all activities. In the
Christian Middle Ages, on the contrary, the point was to take reason captive. The natural understanding was a heathen that
had to bow to the authority of faith. When later the break with the principle of authority had been made, a complete
transformation, a perfecting of human life, was expected from the development of science and enlightenment. Already in Bacon
and Descartes this great expectation can be felt; but in many of the writers of the 18th century it became a passion, a
fanatical faith. Revolutionary cruelty finds its explanation in part in the fact that people could imagine nothing but an
evil will as the cause of not everyone's bowing to the new gospel: perfection lay so near, after all, if only one would
open one's eyes! And yet already during the revolutionary period itself an entirely opposite current makes itself felt,
partly the after-effect of Rousseau's struggle for the right of feeling against the understanding, partly the fear that a
new aristocracy would form if science were favored. The reaction in the first part of the 19th century would not, like
the revolutionary fanatics, abolish science; but it demanded what was still more questionable, that science should ``turn
back,'' i.e. deny its own principles, and it regarded enlightenment, moreover, as a dangerous thing. --- In the thinkers of
the most recent times one finds attempts at a more unbiased and many-sided investigation of the relation between knowledge
and the other sides of the life of the soul. In \textsc{Stuart Mill}'s autobiography one can see how much effort it cost
a noble and able mind to free itself from the one-sidedly intellectualist view in which he had been brought up.

These fluctuations in the valuation of the ethical significance of intellectual culture will be understood when one looks
partly at the psychological relation between knowledge and the other sides of the life of the soul, partly at the
historical experiences we have of the effects of intellectual development.

% --- p. 432 ---
\opage{432}2. Psychology teaches us that knowledge develops more quickly than feeling and will. Our observations and ideas
can arise and unfold in consciousness without corresponding feelings and impulses at once stirring in their full strength.
We can encompass far more in thought and imagination than we can ever embrace with living feeling or realize through the
work of the will. This comes partly from our disposing at any time of only a certain definite sum of energy, and when
intellectual activity lays claim to a great part of this sum, all the less remains for the other activities of the soul;
--- partly from feeling and will being more conservative in their nature than knowledge, holding more firmly to the
direction they have once entered upon. It takes time before the results of knowledge settle into flesh and blood, so that a
harmonious relation is brought about between the different sides of the life of the soul.

There therefore often sets in, as a provisional consequence of intellectual development, a state of disharmony and
discord in consciousness. The widened frame that thought has brought about cannot be filled by feeling and will. And the
instinctive security with which life was led as long as the horizon was still limited gives place to doubt and unrest. ---
Where this disharmony is not present, one often sees instead a weakening and blunting of the life of feeling set in. A
certain languor spreads; life loses in warmth what it gains in clarity.

Revolutionary and reactionary fanaticism can to this extent appeal to psychology. The tree of knowledge is not without
further ado the tree of life. But it does not follow, after all, that one should at once cut it down or prune it.

3. These psychological results are confirmed by history insofar as no sure and general positive ethical effect of
increasing enlightenment can be pointed out. It makes itself more noticeable by the unrest it awakens than by its educating
influence. Insanity and suicide are especially frequent in countries where education stands relatively high, and one can at
any rate not show that crimes decrease in such countries. The hope that has been
% --- p. 433 ---
\opage{433}entertained, that one would be able to close a prison for every school one opened, has not been fulfilled. Even
if one would now\footnote[1]{Cf. \emph{Rümelin}: \textit{Ueber den Zusammenhang der sittlichen und intellektuellen
Bildung.} (Reden und Aufsätze. Neue Folge). p. 4. --- \emph{Mondière} in ``Dictionnaire des sciences anthropologiques.''
Art. \textit{Instruction.} --- R. \emph{Garofalo}: \textit{La Criminologie.} Paris 1888. p. 135 f. 162. --- \emph{Buckle}
goes a step further in his History of Civilization in England, denying that any moral development (though there is an
intellectual one) takes place at all. And \emph{Oettingen} goes yet a step further in (\textit{Moralstatistik.} 3. Aufl. p.
601 f. 762 ff.) blaming ``half-education'' for the increase of crime. --- \emph{Rubin} has shown in ``Tilskueren'' (1884) p.
466 ff. on how uncertain a foundation \emph{Oettingen} has constructed this concept of half-education.} confine oneself to
saying that no causal connection between intellectual and moral development can be shown \textit{statistically}, this must
already arouse wonder and disappointment when one considers the great work spent on instruction and enlightenment. And if
one would say that it is only the hitherto so imperfect enlightenment and instruction that is the cause of the evils and of
the slight ethical effects, a misgiving of another kind arises. If only the knowledge that is won and appropriated in an
independent way can have good effects in ethical respects, then it seems dangerous to give instruction to others than those
who have the inner and outer conditions for being able to reach the goal completely. But how many would not then be
excluded! And what an opposition between knowing and ignorant would not then arise! Society would be split in a way no less
questionable than the single individual's being is split by a knowledge that has not taken root in his nature. It will be
only a poor consolation that history always shows us such an opposition between knowing and ignorant in various forms. The
savages had their ``medicine men'' and rainmakers; the ancient Aryans had their singers, who could lure the gods down to the
sacrifices before battle; the Indians and Egyptians had their priestly caste, the Chinese have their mandarins. Everywhere
there is a knowledge that is held in high esteem, and that requires
% --- p. 434 ---
\opage{434}special initiation and practice, so that it can be found only in a few, who exercise a spiritual dominion over
all the rest. Will it not always go so with our scientific knowledge, the more special and deep-going it is? Science, so it
seems, is and must be aristocratic.

There might seem in this respect to have been a step backward since the Middle Ages. At that time too, no doubt, there was a
learned class; but the principles of the doctrine of faith were common to the simplest and to the most learned scholastic.
Every village could be called an Athens: the highest truths could be proclaimed and understood there. Is anything
corresponding now possible? --- Science has specialized itself to such a degree that the single investigator can master only
a very small domain in an independent way. How then can there be talk of a common, generally spread intellectual culture? ---

The questions touched on here are the most important that social ethics has to treat with regard to intellectual
development.

4. We have surely in recent times come into an overvaluation of the intellectual capacities. We develop them without making
sure whether the development takes place in harmony with the other sides of the life of the soul. It is the natural part of
knowledge to determine the content and direction of feeling and will. But only that knowledge can guide us in face of
reality which has itself sprung from reality. Complaints about mere enlightenment of the understanding and its uselessness
or harmfulness ought therefore to be directed especially at its standing too far from real life both in its origin and in
its application. That was also in many ways the case with the enlightenment of the 18th century. It was limited to a small
circle; the great multitude was to borrow its light from this small circle, and there could therefore be no talk of an
independent acquisition. And the small circle itself of really enlightened and independently thinking men stood too far from
real life. Absolutism kept them out of all practical participation in the affairs of the people. The whole literature was
thereby bound to take on a one-sided, in part fantastic stamp. The people became dependent on account of the prevailing
centralization and lacked both the need for enlightenment and the capacity to use it. Popular
% --- p. 435 ---
\opage{435}enlightenment is even said to have been greater in France in the 16th than in the 17th and 18th
centuries\footnote[1]{Cf. \emph{Tocqueville}: \textit{L'ancien Régime et la Révolution.} 7. éd. p. 207 f. Adolf
\emph{Schmidt}: \textit{Pariser Zustände.} III, p. 335---337. --- Already \emph{Rousseau} stressed the harm the life of
knowledge suffers from standing too far from practical life: ``Tant que la puissance sera seule d'un côté, les lumières et la
sagesse seules d'un autre, les savans penseront rarement de grandes choses, les princes en feront rarement de belles, et les
peuples continueront d'être vils, corrompus et malheureux.'' \textit{(Discours sur les sciences et les arts).}}.

Science grows forth ever anew out of life. Artificial means are able neither to prevent it when a need for it is stirring,
nor to keep it up when such a need is lacking. Natural science grows out of the demands of practical life, out of the need
to master outer nature; philosophy comes into being (as in modern times the history of philosophy in England especially
shows) through questions that life itself leads one to ask; history comes into being out of the need to preserve in memory
the life of the people and of the race. The knowledge that remains without effect is finished knowledge, communicated from
outside. Where one can get a seed to sprout, one cannot always get a whole tree to take root.

It is the ideal of the pedagogical art to awaken the need for knowledge before knowledge is communicated, and to communicate
knowledge only to the extent to which a need for it is felt. The more this ideal is successfully realized, the more will the
disharmony described above disappear. There is no reason to turn back; it is only a matter of continuing the process of
intellectual development on as good conditions as possible. This process of development is no artificial product, but has
come forth from life's own need. ---

The striving for enlightenment that judges a human being merely by the development of his understanding, and whose whole
psychology is confined to distinguishing between obscurity and clarity, received its deathblow from
\textsc{Rousseau}\footnote[2]{Cf. my book: \textit{Rousseau og hans Filosofi}\textsuperscript{2}, especially p. 49---78 (the break
with the Encyclopedists and with Voltaire).}. However great the significance we ascribe to
% --- p. 436 ---
\opage{436}science and knowledge, there are nevertheless other sides of man's being to which we look first and foremost in
judging an individual. Over against the life that unfolds in feeling and will, knowledge is only a forecourt in which we
can move without entering the sanctuary proper. It is here that kinship between the simplest human being and the greatest
thinker can constantly be present in spite of all difference in intellectual culture. It is no less a man than
\textsc{Kant} who has said: ``I am myself an investigator, and am so by inclination. I feel the whole thirst for knowledge
and the eager unrest to get further in it, as well as the satisfaction at every advance. There was a time when I believed
that all this could constitute the honor of mankind, and I despised the rabble who know nothing. Rousseau has set me
right. This blinding superiority vanishes. I learn to honor human beings, and would consider myself far more useless than
the common laborer if I did not believe that this very consideration could work to establish the rights of
mankind''\footnote[1]{Fragmente. Werke ed. Rosenkrantz. XI, p. 240.}. % translated from Høffding's Danish
In his first writings Kant had consoled himself for the want and struggle in the world with the thought that
enlightenment was after all going forward in a small chosen band. The study of Rousseau's writings produced a whole
revolution in his views and made him seek what is properly human deeper than in the understanding. ---

When science grows out of life and always preserves its connection with life, and when, moreover, the center of spiritual
life lies not in the intellectual domain but in feeling and will, the danger is diminished of the disharmony in the
individual and the split in the race to which intellectual development so easily seemed to lead. And no less will the
danger thereby be avoided that lies in the development of an intellectual proletariat. Where intellectual development
stands in close relation to a definite scientific work, or to a pronounced need for clarity in the view of life, or to
activity for definite practical ends, there it is sound and necessary. But it can become a misfortune where the conditions
of life,
% --- p. 437 ---
\opage{437}inner and outer, do not allow this connection to subsist. It is a task for the pedagogical art to preserve
harmony in this respect, as far as possible, in the individual and in the people.

5. It cannot be otherwise than that there will be an opposition between different stages of intellectual development. Even
if it is not between knowing and ignorant, it is between those who know more and those who know less. This follows from
the division of labor. If science is to be cultivated in proper fashion, there must be those who devote their whole lives
to it. A learned circle or class will form in the people. Two things then matter here if development is to be sound. In
the first place, there must not form a learned caste recruited from some classes of society but not from all. Lower and
higher instruction must be so ordered that there is an even transition from the lowest to the highest
stages\footnote[1]{Such an even transition was present when theology was the main subject at the university, as the
catechism was in the common school. Now theology is reduced to a minimum in the grammar school and to a specialty at the
university, but the catechism still gives the common school its stamp. \emph{Martensen} maintains that so it must be: ``It
is religion that makes the common school a common school'' \textit{(Social Etik,} p. 355), an utterance that is striking even
for his own standpoint. A dualism is here fixed that will have very questionable consequences if it is not remedied in time.
And it will then hardly be science that comes to ``turn back.''}, and that the difficulties of running through all the
stages do not become insuperable for him in whom serious impulse and capacity stir. And next the learned class must regard
research not merely as the satisfaction of its own impulse, but as a social work, a work carried on on behalf of the whole
of society. The single investigator is placed at his lookout post to observe everything he can catch sight of there. When
this disposition animates him, he preserves his unity with the race even if he feels isolated at his solitary post, and
even if he long remains unrecognized and misunderstood because the others cannot catch sight of what he has seen. His faith
in truth is at the same time a faith in the significance of truth for the race.

% --- p. 438 ---
\opage{438}The progressive specialization of science surely makes it more and more difficult, even for its own
cultivators, to obtain an independent overview of its results. But intellectual culture does not presuppose that everyone
knows everything. There can be an independent intellectual culture in wide circles, if only the power of thought itself is
developed. In spite of the differences between the sciences, it is after all one and the same power of thought with which
they all work, and it is one and the same world that they all, each from its own side, seek to penetrate. Therefore he who
has been led in a reasonable way into a single domain can easily be brought to understand the tasks and difficulties that
present themselves in other domains. There are, moreover, certain main thoughts that recur wherever human knowledge works,
and which it is the task of philosophy to bring forward and develop. Such a thought is the idea of existence as a
law-governed connection. The different sciences each seek to penetrate their piece or their side of one single great causal
connection. To the idea of causal connection is joined the idea of development, which in its main features is found again in
the different domains of experience. In a small part of the world we can therefore study the whole world and win an
understanding of it. There forms in our consciousness a general picture of the world, whose gaps the single sciences are
working uninterruptedly to fill. Even he who cannot devote himself to scientific research will nevertheless be able to form
an idea of the main laws of the existence of which he is a member, and get guidance for taking a more universal view of
himself and his life than is possible as long as he believes that everything revolves around him and the narrow horizon of
his interests. By recognizing his place in existence he will also gain in self-knowledge. He will be strengthened, besides,
in his faith in the unity of the human race when he gets an insight into the great common work of forming the human race's
picture of the world, a work that generation after generation lays claim to the noblest powers, the most persevering and
self-sacrificing striving. The individual does not here work for himself
% --- p. 439 ---
\opage{439}alone; his own research would not lead him far. But he can lay his stone to the great building that is the
common world view of the human race, which slowly develops through the ages. And he can perhaps from his workplace survey
so much of the building that he can form an inkling of how it would look if it were wholly finished\footnote[1]{Cf. my
lecture \textit{Videnskab og Folkeliv} (Tilskueren 1903), and a series of lectures by French historians and philosophers,
published under the title: \textit{L'éducation et la démocratie.} Paris 1903.}.

6. Science works to found community through the necessity of cooperation on the race's picture of the world. Not only are
contemporary investigators thereby bound together, but the thought of the debt in which we stand to the past becomes alive.
Nowhere can progressive development be shown so clearly as in the domain of scientific knowledge. Nowhere is it easier to
show how the one generation stands on the shoulders of the other and thereby gets a wider view than it. But since science is
in constant development, cooperation cannot always be quite peaceful. Schools and parties stand against one another and
fight one another. Here, then, not one great community forms, but several smaller ones in strife with one another. But such
strife can be a means of advance when it is more than personal squabbling. It then arises partly from the matter itself
presenting several different sides, partly from the investigators coming with different presuppositions. In both respects
the formation of parties can be fruitful. When one joins with others in order to exhaust everything that can be found out
about the matter from the side from which one sees it, or from the presuppositions one proceeds from, it will happen with
greater thoroughness and zeal. The very affect that party opposition brings with it can sharpen the gaze. Even purely
personal quarrels can in this way become fruitful. The common work is furthered not only through the disinterested devotion
of individuals, but also through the clash of their egoisms. They work
% --- p. 440 ---
\opage{440}restlessly to be proved right, to maintain their honor, and can thereby often further an end far more
comprehensive than the one they have in view in the heat of strife.

\chapterhead{XXIX.}{THE FREEDOM AND INDEPENDENCE OF INTELLECTUAL CULTURE.}

1. \textsc{Freedom} is the most important condition of intellectual culture. Science knows no boundaries staked out from
outside, nothing that is too high or too low to fall under its treatment. It stakes out its boundaries itself, deciding,
from the very nature and mode of working of knowledge, which questions it is able to discuss and on which it is unable to
decide anything. And even with regard to these last questions, insoluble for science, science claims the right to examine
the answers that have been given to them, --- to inquire into the origin of these answers and the reason why they have found
adherence. Where the theory of knowledge ends, historical criticism and psychology take over.

Freedom is here, as everywhere (see VIII, 6; XXIII, 2), both end and means. There is an immediate satisfaction in the
unhindered use of the power of thought, so that an unnecessary intervention here is an offense against the principle of
welfare itself. But freedom is also a means. Only free inquiry can see the matter from all sides, discover the different
possibilities and draw all the consequences. Yet we meet here the peculiar phenomenon that even those who abhor every
justification of ethical rules by regard for utility often appeal to it on this point. ``Free inquiry,'' it is then said,
``shakes everything, even opinions that human beings cannot do without. Truth is not harmful; but doubt and inquiry without
firm ground under one's feet
% --- p. 441 ---
\opage{441}have a harmful effect. The individual may be allowed to doubt and inquire if he can defend it before his
conscience; but let him keep his doubt to himself and not draw others onto his tottering ground!'' --- The justification
such a line of thought can have under certain circumstances we have already acknowledged in another connection (XII, 4---6,
cf. also VII, 4 and XXI, 5). But here, where the talk is of the significance of intellectual culture for the life of the
race as a whole, regard for the race's lasting conditions of welfare must be placed above regard for individuals'
sensitivity. It will be pernicious to take regard for the momentary pain caused\footnote[1]{Stuart Mill has rightly
stressed it as characteristic of the most recent times that people speak more of the utility of opinions than of their
truth. ``In the present age,'' he says \textit{(On Liberty.} Danish transl. p. 39), people ``are not so much convinced of the
truth of their opinions as that they do not know what to do without them.'' % translated from Høffding's Danish
Mill, the zealous utilitarian, feels misgivings at this. The point is that truth is for him a condition of life and welfare
of the first rank, and what he is fighting is the short-sighted and limited conception of what is useful.}, and thereby to
forget the consequences that the withholding of truth --- and doubt too is a truth when it is a well-grounded doubt --- can
have for wide circles and for long ages. Pedagogical considerations it is chiefly the individual's business to take toward
other individuals; in society they must often recede altogether. --- Those who from a social standpoint have misgivings about
free inquiry, moreover, as a rule think that they themselves possess the truth, and instead of negotiating with them about
the justification of \textit{uttering} doubt, it would therefore be more to the point to negotiate with them about the
justification of doubt itself. They take cover behind the consequences of uttering doubt, but it is really doubt itself they
want to get at. In any case they must admit that the utility of uttering doubt must be able to be discussed freely and
openly. They cannot, after all, think that they hold an infallible knowledge of what is useful or not. And can the utility
or harm of doubt be discussed publicly without doubt itself thereby being publicly uttered and made known? --- --- However one
% --- p. 442 ---
\opage{442}twists and turns, one will not be able to avoid this danger, if it is a danger.

Free inquiry surely brings dangers with it. But the way to truth goes only through errors. ``A living error,'' it has been
said, ``is better than a dead truth.'' What is living in the error is a thought that can work its way off the wrong path;
but in a truth long since settled, which can no longer set the powers of the mind in motion, there are no germs of life.
When we remember those who found the truths on which we build, we must also remember those who developed the ``living''
errors and thereby built intermediate stations on the way to truth, even if this way is a detour. --- A living truth is
naturally the very best of all. ---

2. Intellectual culture works to found community not only by bringing investigating individuals together for common
inquiry and communication of thought, but also by leading to the formation of institutions of learning. History shows us
here some of the finest examples of free formation of society. In Greek antiquity the philosophical schools were formed by
young men joining around a thinker who had awakened their interest. In the Middle Ages schools arose at the monasteries,
where young men studied in common. The universities arose in the Middle Ages through the free union of men eager to learn;
the word university means, as is well known, precisely a community or corporation of teachers and students. Only later did
church and state gain influence over these scientific institutions, which at that time had an all the greater significance
since they were the only means to higher culture; books and other scientific means were as good as wholly lacking. In our
days we have seen a similar phenomenon here at home in the folk high school movement.

Where such a free union takes place, it is the state's duty not to put obstacles in its way but to support it according to
its ability. The state has, as was mentioned earlier (XXII), the right and the duty to secure a minimum of knowledge in every
one of its future citizens. But it is also one of the state's tasks to use its means and its organization for the benefit
% --- p. 443 ---
\opage{443}of the progress of higher intellectual culture, apart from its having to require a certain scientific culture
of those who wish to be its officials. Free union easily acquires a certain aristocratic character. There are many who are
bound to the soil by need and cannot without help join in the appropriation and continuation of intellectual culture. And it
will not be fortunate if only the rich and well-to-do classes carry intellectual culture. The state has here, as in the
domain of material culture (XXVI, 14), a \textit{distributing} activity to carry on. It is not only to seek to favor the
formation of a capital of knowledge as a common good for the whole of society, but also to further the best possible
distribution of this knowledge. And it does not happen here, as it so easily happens in the distribution of material goods,
that nothing can be given to the one without being taken from the other.

The danger will always lie near, however, that the state makes science its servant instead of being itself science's
servant. The state is, after all, at any given time represented by the holders of the power of the state, whose party
interests or personal interests can lead them to favor or hinder certain tendencies or persons without regard for the real
good of intellectual culture. It has even been wished to conceive the matter so that endeavors for intellectual culture are
always to be subordinated to other considerations. It has been said: the real human societies are home, church and state;
the school (the society for the furthering and appropriation of science) is to serve them, not conversely; home, church and
state have in union to determine how the school is to be\footnote[1]{``The school is no independent institution beside the
home, the state and the church, but a helper dependent on them. That is the school's essential position, which nature and
religion have assigned it, and therein the groundless untruth shows itself of the modern endeavors that would place the
school independently of home and church. Home, state and church must have schools that answer to their spirit and their
demands. They cannot be prevented from this without great injustice.'' % translated from Høffding's Danish
Bishop \emph{Ketteler}: \textit{Freiheit, Autorität und Kirche.} Mainz 1862. p. 209. The Catholic Church demands, by
virtue of its divine origin, the right to supervise all schools, from the highest to the lowest. The right of parents and of
the state is subordinate to the Church. And in the school itself (according to a pronouncement of the College of the
Inquisition of 30 June 1775) religious instruction is the main thing; all other subjects are to be regarded only as
supplements. Non-confessional schools (scholae neutrae) are an abomination: Jesus himself said, after all, that he who is not
with him is against him! Cf. \emph{Cathrein}: \textit{Philosophia moralis,} p. 396. \emph{Gury}: \textit{Theologia moralis.} I. p. 314. ---
Not only in Catholic writers does one meet such a line of thought.}.

% --- p. 444 ---
\opage{444}Here the school, the society founded on the cultivation of science, is denied all initiative. But what a
self-contradiction one thereby gets entangled in! For what is it that state, church and home wish to attain through the
school? It must surely be truth. But they prescribe beforehand themselves how truth is to be constituted. They must then
think they are beforehand in possession of the rule of truth, so that the school is merely to apply it in particulars. But
from that come only opinions of authority, not scientific truths. And it is overlooked besides that science, among other
things, also studies the origin, nature and significance of home, church and state. What significance has science if it does
not in the end, perhaps by many detours, but nevertheless unavoidably, lead to changes in the life that is led in home,
church and state? An intellectual culture deprived of the possibility of acting back on life is, as we have seen (XXVIII,
4), unhealthy and useless. One is therefore forced to grant intellectual culture a place as an independently cooperating
cause in the development of human life. If home, church and state cannot use independent science, that becomes their own
affair; as surely as truth is an indispensable condition of life, they thereby show that they are on a pernicious path; ---
but it is meaningless to imagine that one is supporting science when one does not give it complete freedom. Specialties and
curiosities can thrive, --- bastards begotten by fantastic mysticism and sober thinking can romp about, --- old doctrines can
be kept up without scientific freedom\footnote[1]{According to an agreement concluded in the year 1857 between the Pope and
Württemberg, the Catholic theological faculty in Tübingen was placed under the bishop's direction and supervision. The bishop
was to give the professors authority to hold theological lectures and could take this authority back again. He was to
examine their notebooks and their textbooks! The academic senate in Tübingen set up a committee to examine whether the
theological faculty could under these circumstances still be said to be a member of the university. This committee came to
the result that the professors of Catholic theology could from now on not be regarded as representatives of free science and
therefore not be elected members of the academic senate. Bishop Ketteler \textit{(Freiheit, Autorität und Kirche.} p. 23)
is much offended at this; but it is not easy to see how there can be talk of free science under the constant censorship of a
bishop! --- The difference, for the rest, is not great whether it is a living bishop or a dead book to which one is to
subordinate oneself. Where there is not unconditional freedom for thought, there can be no science. A constant difficulty
asserts itself here in the position of the theological faculties, congregations and churches demanding control over what
they teach, thus determining the results to which they are to come. The Pope, for example, refuses to acknowledge
theological faculties that are not under such control. In France the Catholic theological faculties at the state
universities were therefore abolished in the year 1885, since they were not found useful without papal acknowledgment.}; but
the great
% --- p. 445 ---
\opage{445}and bold thoughts that carry our knowledge and thereby in turn life forward will be sought in vain, and without
them even the coziest life in home, church and state is after all only a meager thing. If the school is to fill its place in
society in the right way, the teaching profession must be acknowledged as an independent profession, of equal standing with
the other organs of society and independent of ecclesiastical and political authorities. Its task is to work for the culture
of the people, and its obligation to take regard for what exists in family life, church and state must be subordinate to
this purely human task. Such an independent position for the teaching profession will be more important than school laws and
school plans. At the highest institutions of learning this goal has in essentials been reached; the point is to carry it
through for the common school too.

To the freedom and independence of science belongs not only that it be allowed to develop within its own domains without
being hampered by outer authorities, but also that it not be prevented from exercising the influence on the general view of
life and of the world that it lies in its nature to exercise.
% --- p. 446 ---
\opage{446}The scientific picture of the world necessarily determines our view of life and our faith. For it is in the
real world that our convictions must stand their test, prove their strength; and this real world we simply cannot come to
know in any other way than with the help of science. A faith that shuns truth is a mere dream. The difficulty lies in this,
that the changes the results of scientific research produce in a view of life as a rule take place slowly and imperceptibly,
and that it cannot always be decided with certainty what the definitive consequences of a scientific result are. Great
demands are made here on love of truth and honesty. The Socratic ``Know thyself!'' finds application here as a summons to
intellectual self-examination, in order to come to a decision as to where the center of gravity of our view of life lies,
and to what degree the inherited religious ideas have been changed by the results of knowledge that have been appropriated.
Here there is the possibility of a great self-deception and of many illusions. What alone can help here is a real
intellectual culture, a practiced scientific way of thinking that with right tact is able to hold fast the essential points
of view, to distinguish between hypotheses and settled truths, to venture where venturing is needed, and to withhold its
judgment where nothing can be known. Intellectual culture is an art that is far from found in all men of science. One can
bore one's way down into a specialty and work there with competence, and yet be devoid of all capacity and tact when one
stands before other matters. Either one then throws oneself, when one comes out of one's cave, into the arms of the first
faith that comes along, or one exaggerates the results won in the special domain by leading everything back to them, or one
ends up in a more or less blasé doubt. One maintains the omnipotence of science, or one declares that, apart from a few
special domains, it can teach us nothing about existence. And one perhaps swings between both assertions, so that for a
while there is talk as if all riddles were solved, and when this has been carried to excess, science is suddenly declared
bankrupt. Over against such oscillations sound intellectual culture --- which
% --- p. 447 ---
\opage{447}consists in a union of critical thinking and personal art --- is alone able to preserve continuity in
spiritual development\footnote[1]{Cf. my essay \textit{Filosofi som Kunst.} (Mindre Arbejder. I.).}. It conditions the
capacity to lead the results of thought over into personal and social life in the right way.

3. History shows that encroachments by home, church or state on the school do not reach their goal. They rest partly on an
overvaluation of the effects of intellectual culture, partly on a failure to recognize that direct influence often even
leads to the wholly opposite result of the one intended. --- Enlightenment of the understanding does not at once determine
the whole direction of life. Interests, feelings and impulses are not at once changed by the knowledge and ideas that are
inculcated. The ``danger'' is thus not so great and overwhelming as is thought. --- And on the other side the zealous
inculcation of a certain series of ideas can precisely awaken a need to strike out on series of ideas going in the directly
opposite direction. Such an effect of contrast often makes the one generation go in the direction opposite to the other.
Capacities and impulses that found no nourishment in the ideas with which one's consciousness was filled during upbringing
then later force their way forward all the more strongly and violently.

In France the changing governments of the last century have all sought to use the school for their ends, without their
duration having become any greater for it\footnote[2]{Adolf \emph{Schmidt}: \textit{Pariser Zustände.} III. p. 391.}. The
only thing attained is to offend against the good of intellectual culture. The basic condition of all higher culture is that
it must be sought for its own sake. ---

Although it is necessary that the state support science, it would nevertheless not be fortunate if all scientific
investigators in a people worked in the service of the state. Just as where the state appears as a distributing power in the
domain of material culture, it is necessary here too that it have competitors in private endeavors. When scientific research
is sound and flourishing, it grows out of life and is not merely called forth by the helping hand the state can afford.

% --- p. 448 ---
\opage{448}The state will, however, as a rule be able to show greater impartiality than private and ecclesiastical
institutions. Precisely because the state embraces the \textit{whole} people, its gaze will be directed toward the universal
points of view for culture, and the prejudices with which single circles and sects confront science will not so easily come
to rule in it. The private school depends on ``the public,'' the state school on the people.

When the state school is ordered and directed in a true scientific spirit, it will exercise a uniting and binding influence.
Youth from all estates and classes will be able to meet there and form a brotherhood. The presupposition for this is that no
particular political or religious tendency come to rule in the school. At the point where political and religious
differences begin, it is the home's business to educate. The school can afford only the education that scientific culture
brings about. It may be difficult to draw the boundaries between the common life of thought that has its firm foundation in
science and the political-religious views that are essentially determined by feelings and traditions. But only by a serious
endeavor to fix such boundaries can intellectual culture be maintained as an independent element in the development of
mankind.

\grouphead{\textit{B. AESTHETIC CULTURE.}}

\chapterhead{XXX.}{ART AND LIFE.}

1. We understand a thing when we see it in its definite place in the series of things, borne by other things and itself in
turn bearing other things. Science dwells on the single and individual only in order to find the laws of its connection with
% --- p. 449 ---
\opage{449}the rest of the world. When, on the other hand, we relate ourselves aesthetically to things, we see each of them
by itself as a distinctive and complete whole, rejoice in the image simply because it gives us something distinctive and
characteristic. Where we cannot ourselves find or form distinctive and characteristic images, art helps us to them.

Both intellectual and aesthetic feeling are akin to the sympathetic feelings, since our pleasure and displeasure in them is
determined not by our physical urge for self-preservation, but by devotion to something that lies beyond ourselves. But
aesthetic feeling bears the stamp of sympathy more than intellectual feeling does. The latter is determined more by the
relation between things than by the single thing itself in its distinctive nature. Aesthetic feeling is determined by our
seeking to make ourselves one with things, to live their life, to live life itself over again in our imagination. There is
an aesthetic pleasure already in the free use of the limbs, in play, in which we imitate pursuits otherwise carried on for
serious and practical purposes. Indeed, all art can be called a play, insofar as it brings life or parts of life before us
in an image, in an ideal reproduction, and invites us to be moved by it as if it were a real life.

2. The development of aesthetic feeling acquires ethical significance already in that it accustoms us to regard things
without egoistic ulterior motive, accustoms us to forget ourselves over the value of the sight. There works in the truly
beautiful a power that compels acknowledgment, and not merely acknowledgment but love and admiration. Both in him who is
able to produce beautiful works and in him who stands as aesthetic beholder before the works of art or of nature, a
disinterested (i.e. unegoistic) need expresses itself, at the same time as the power of intuition, attention and the power
of production are set in vigorous and harmonious activity. And the feeling thereby awakened will not be confined to the
moments of contemplation or of production, but will influence both the view of life and the conduct of life and there
exercise a harmonizing and idealizing influence. Where purely moral judgment and discipline would be too
% --- p. 450 ---
\opage{450}clumsy an instrument --- in many of life's most delicate and most individual relations --- there aesthetic or
artistic tact can often carry one over difficulties and avert catastrophes. The aesthetic education of man, for which
\textsc{Schiller} spoke up\footnote[1]{Cf. \textit{Den nyere Filosofis Historie}\textsuperscript{2}. II. p. 128 f.}, is,
however, more than a mere means for moral life; it leads to a free and many-sided use of the powers instead of the
dividedness that the conflicting elements in human nature bring with them, and instead of the one-sidedness that the division
of labor brings with it for the single working human being. --- With the disinterested character of aesthetic goods it hangs
together that they are common goods, that they can be shared by many without their value being diminished. The community they
make possible has its great ethical-social significance. --- But it is of significance especially that aesthetic pleasure is
conditioned by the use of powers that real life too lays claim to. On this rests the seriousness in aesthetic play.
Aesthetic enjoyment is more than mere passive entertainment. It is attained only where we feel ourselves gripped by a real
image of life. The ethical effect is here not different from the perfect aesthetic effect. \textsc{Aristotle} defined tragedy
as an imitation of a significant action, an imitation that awakens our fear and pity and thereby purifies these feelings. By
``purification'' he probably understood that the feelings are awakened and shaped in a harmonious way with the help of the
artistic presentation, and that they are thereby freed from what is wild and chaotic. Just as music at once releases and
harmonizes the feelings, so too dramatic art is able to do this\footnote[2]{Cf. Erwin \emph{Rohde}:
\textit{Psyche}\textsuperscript{2} II. p. 48 f.}. What \textsc{Aristotle} said of tragedy we can extend to all art. In play
or through the image the same feelings are awakened as by the real circumstances presented; but they are awakened in such a
way that what is painful and egoistic falls away. Real experiences come upon us without our always being able to find
coherence in them; chance seems to rule. They present, moreover, several different sides and therefore also set several
different
% --- p. 451 ---
\opage{451}feelings in motion. Our mood in the experience is therefore not ``pure,'' i.e. different feelings not in
harmony with one another assert themselves. In the artistic image, on the contrary, everything is laid out so that a total
mood is awakened which answers to the characteristic main feature of the thing or event. Without robbing the image of its
concrete-individual character, artistic treatment brings out for us what is essential. We are thereby freed from what is
incoherent and get a collected impression, and can therefore also far more easily make ourselves one in our feeling with
what is presented. With this it also hangs together that artistic reproduction teaches us to ``understand'' the thing better.
It gives us no scientific explanation of it, to be sure, but by showing it to us in its whole distinctiveness art makes clear
what justification this thing has for existing. Whatever else it may be, it is a characteristic part of the world, or rather
is itself a little world. And this little world is as it were moved out of time by being fixed in an image; what passed by
fleetingly as an intermediate link between past and future lives an eternal life in art. With the help of art we are freed
not only from dispersion and chance, but also from transience. --- While in tragedy it is fear and pity that are
``purified,'' in comedy it is the feeling of power and self-esteem. It is not personal scorn that finds vent in the laughter
that comic poetry procures us, but a feeling of liberation at seeing the finite, the self-contradictory and the foolish in
its whole nakedness and yet as something that, after all, also belongs to life\footnote[1]{Cf. \textit{Psykologi} VI E, 9
c.}.

3. When the question is whether an aesthetic presentation agrees with what ethics demands, many will at once think of the
question whether strongly sensual pictures and descriptions are permissible or not. There is, however, according to the way
in which we have here conceived the ethical significance of aesthetic culture, no reason to raise this question. As regards
content and manner of presentation there can be no conflict between what aesthetics really demands and what is ethically
% --- p. 452 ---
\opage{452}justified. Everything that really has aesthetic value must also be ethically justified. We stand here before
values of life that belong to the ends all ethics is to serve. The beautiful belongs to what makes life worth living; it
belongs to the domain free of conscience. (VIII, 1). --- The pedagogical use of art is another matter. One cannot put every
work of poetry into everyone's hand. But that has nothing to do with aesthetic value. He whose sensual impulses are set in
strong motion by a description does not come to relate himself aesthetically to it: his feeling is not ``purified,'' but
incited. The young man in one of Lucian's tales who let himself be locked in for the night in the temple of the Cnidian
Aphrodite in order to embrace the statue was not exactly moved by \textit{aesthetic} feeling.

It follows of itself that when art is an ideal life, the value of art will answer to the value of life, and the value of a
work of art therefore rests not only on the competence or the genius with which the subject is grasped and carried out, but
also on the value of the life it presents. Here lies what I have called the open place of art (Psykologi V B, 12), since the
valuation of life plays its part here too. At this border point a struggle will always be waged. The moralists will here
easily come to take the matter too heavily. They easily do the artist an injustice by overlooking that the artistic treatment
itself, the artistic point of view, places the artist in a different relation to a subject than the purely practical one.
--- Art is not to undertake any ethical valuation, is not to moralize directly. But genuine art will not demoralize either by
directing attention in a one-sided way to certain sides of life. If the spectator can transgress against a work of art, the
artist can also transgress against life. Naturally there can be precisely different views about how life is and what
attention is especially to be directed to. Every significant artist has a struggle to wage here to make his view of things
prevail. An artist's genius is inseparable from his personality. The boundary cannot be drawn definitely and is always
unnatural. Therefore an involuntary valuation enters into his work, even where he --- from conviction, or from
``distinction,'' or from doctrinairism --- strives
% --- p. 453 ---
\opage{453}to keep all valuation out. It is this hidden and involuntary valuation (which need not at all presuppose any
``Absicht'') that can rightly be made an object of ethical criticism. But it is by no means easy to discover and establish.
It will appear especially as an unjustified generalizing. Art can by its nature give only the single, individual, special;
for only that can be united into a living image. But whether the single image gives an example of a general rule or merely
presents an exception due to special circumstances,
--- about that art as art can say nothing. Therefore it cannot put problems up for debate, except insofar as it can prove by
its presentation that something such as is presented does after all exist in the world, lies within the possibilities of
life. A real posing of a problem presupposes a large and rich material and much reflection.

Much in the resistance to modern aesthetic realism surely derives from many people's seeking in art only entertainment and
rest and therefore not wishing to see the bitter and dark sides of existence. They do not want their fear and pity awakened
at all. The old tragedies people were used to; they had the necessary guidance for hearing the heavy steps of fate in them;
but in a modern tragedy like ``Ghosts'' they were unable to hear them. The subject was here too close to life. And yet modern
realism in its finest works has really done nothing but carry the great points of view deeper down into real life and show
us seriousness where one was used to seeing only everyday life and everyday stories\footnote[1]{In the preface to ``Germinie
Lacerteux'' the brothers \emph{Goncourt} declare that they wished here to test whether tragedy was to be definitively dead.
They wished to test ``whether the misfortunes of the small and the poor would speak as loudly to interest, feeling and pity as
those of the great and the rich.''}. Art here has an educating effect by making us open our eyes, by strengthening our feeling
of the seriousness of life and awakening our sympathy with its suffering. It teaches us especially, in great, vivid images,
to see the deep connection in human life. Without direct moralizing the poet, merely by pointing with a sure hand
% --- p. 454 ---
\opage{454}to the decisive causes of the development of characters or of the course of fate, can give us a deeper
instruction and a better guide for valuing actions and orders of life than can be got from any philosophical ethics. What
has been said of Goethe's ``Wahlverwandtschaften,'' that it ``is not written for the sake of morality, but has a
morality''\footnote[1]{Richard \emph{Meyer}: \textit{Goethe.} Berlin 1895. p. 286.}, holds of every significant work of
poetry. And we may add: it holds of it the more, the more it is permeated by the spirit of realism and determinism, for the
more definitely can the point be indicated in the series of events from which the fateful effects spring. In a gripping way
Mrs. Alving in ``Ghosts'' thus points back to the point where her husband's misfortune came into being and the ground of his
decline, with all its consequences, was laid. Thereby a light emerges in the midst of the darkness into which this tragedy
leads us\footnote[2]{Ghosts (Act 3): ``You spoke before of the joy of life; and then it was as if a new light rose for me
over all the things in my life\dots{} You should have known your father when he was quite a young lieutenant. The joy of life
was in him!\dots{} Your poor father never found any outlet for the overpowering joy of life that was in him. Nor did I bring
any holiday weather into his home.'' % translated from Ibsen's Dano-Norwegian
}. The old Greek tragedian's word, that mortals must win knowledge through suffering (πάθει μάθος), is confirmed here. The
tragic thing is that the suffering can be so overwhelming that the knowledge won becomes for the sufferer himself only a
light that falls back on the past. But for him who by the artist's hand follows the sufferer's fate ``in fear and pity,''
this light can come to shine forward. --- Works of poetry of an ``idealistic'' kind often have very little ethical value,
because they do not acknowledge the fixed causal connection of life.

4. However much art can and should be for life, it must nevertheless not step into the place of life, and real life must not
be treated merely as an aesthetic object. --- This will not so easily happen with the serious artist, who relates himself not
enjoyingly but working, and for whom art is a serious task of life, a social task, since he, like the man of science in his
own
% --- p. 455 ---
\opage{455}way, feels himself placed at a post to observe life and teach others to see it. He feels what \textsc{Michel
Angelo} expressed, that ``genuine art is of itself pious and noble through the spirit in which it works; for nothing makes the
soul so pious and pure as the endeavor to produce something perfect.'' % translated from Høffding's Danish
And he feels keenly enough the resistance that reality, the refractory material, offers to what stands before him in his
imagination. He perhaps perishes by this resistance. The danger lies far nearer for the predominantly receptive natures and
the dilettantes. They easily come to treat life as an aesthetic play. \textsc{Schiller} said that man is truly man only where
he plays. His meaning was that the world of imagination and play is man's own work, and that a free heart and an energetic
will are required to withdraw from the pressure of reality and hold this ideal world fast. But his words have been used as a
motto for an aesthetic view of life that relates itself jestingly and ironically to every practical relation, without always
taking art itself any more seriously for that. Work is left to the philistines; the geniuses in their own conceit look down
from their ideal height on the world as on a droll thing that does not concern them. --- The realistic tendency in art can
just as well as Romanticism lead to living in an imagined world and becoming a stranger to reality. The realist moves in
imagination just as much as the idealist, and the danger may for him perhaps even lie nearer than for the latter. The
overstrained idealist will as a rule have the consciousness that he lives in two worlds, the world of dreams and the prosaic
world; he scorns the latter, but can for that very reason understand quite well how to take it as it is. For the realist,
whose imagination seeks to sate itself with impressions from the real world, the temptation will be greater to treat all
relations of life in an aesthetic way.

This danger appears in history at every time when art and aesthetic interest stand in the foreground. In any case it becomes
the inheritance the next generation receives; this comes forward ``with the demands of culture, but without originality
% --- p. 456 ---
\opage{456}and without living creative power''\footnote[1]{With these words \emph{Hettner} (\textit{Italienische Studien.
Zur Geschichte der Renaissance.} p. 175) characterizes the generation that followed the time of Raphael and Michelangelo. But
he finds the reason for the abrupt fall in the very nature of the Italian Renaissance. It had burst the medieval chains, but
was able neither by way of ethical disposition nor by that of working thought to shape a new ideal of man. It was an age
without firm moral concepts and models, an easy prey for the ecclesiastical reaction.}. Even a great and significant art (like
the Italian in the time of the Renaissance) can stand without connection with the real powers struggling in the time and
without connection with the life of the whole people.

5. Art is to give the content of life form and clarity, to widen the gaze and sympathy, to point to the way along which
development is to go. The great artist is at the same time something of a prophet. Art can never become a merely private
matter. The whole people and the whole age are to come to know themselves in it. Every people and every age must therefore
have its own art and cannot live wholly on the art of other peoples and other ages, however great the significance of coming
to know it may be. Even where what is universally human is expressed and portrayed, it must nevertheless be done in the
particular form that answers to the single people's experiences. The demand for a national and contemporary art can be
grounded in two different ways. The artist sees and works with his best powers only where he is animated by the feelings and
ideas of his own home and his own time; only then can he appropriate what he sees and reproduce it. And on the other side
every people and every age properly understands only itself; what is to affect them aesthetically must be flesh of their flesh
and blood of their blood. --- Naturally it is of no use for art to be national and contemporary if it is not real art.
\textsc{Julius Lange} has aptly shown how unfortunate consequences it can have when the national demands on art gain the upper
hand over art's own demands. ``The condition for an outstanding and excellent art rests,'' he says\footnote[2]{\textit{Sergel
og Thorvaldsen.} Studier i den nordiske Klassicismes Fremstilling af Mennesket. Copenhagen 1886. p. 77.}, ``more on the
nation's learning to enter into
% --- p. 457 ---
\opage{457}the distinctive tasks of art than on the demand that art shall enter into those of the nation. The nation must not
merely place itself in relation to art as its master; from that it in fact draws no advantage; it must first and foremost make
itself art's apprentice, learn from it to see the spirit in form, light and color, learn the immediate intuition of man and
nature.'' % translated from Lange's Danish
There is a common content of thought and feeling into which every people and every age must live itself and draw nourishment
from. Through Homer, Dante, Shakespeare and Goethe we live the spiritual life of European mankind over again.

As regards the future of art, there is --- on account of the connection between art and life --- reason to fear for it only if
life itself becomes poorer, --- if it is led in a smaller style and without the struggle and tension that great tasks and ideas
bring with them. Genius, which is at once the great child and the great discoverer, will again and again open the depths of
life in spite of the philistines of the future, as it has done in spite of those of the past. What matters essentially is to
keep the senses open and the heart alive; the world of the spirit will not be closed. And hand in hand with art, the serious
conduct of life that takes life as the struggle it is will work to maintain the conditions for the great poetry of life, which
is and remains the only form in which we are able to sum up the content and value of life. Here the ethical reckoning and the
aesthetic need become one.

\grouphead{\textit{C. RELIGIOUS CULTURE.}}

\chapterhead{XXXI.}{ETHICS AND RELIGIOUS FEELING.}

1. \textsc{Ethics} must, as was stated earlier (II, 3), build on as few presuppositions as possible. It must not demand
any special position in science and must not try to shake the
% --- p. 458 ---
\opage{458}principles, results and hypotheses that other sciences set up. But as religious life has developed in the human
race, it has been tied to assumptions and dogmas that have again and again led to conflict with science. Ethics might then
seem bound to enter into relation with religious culture as with a purely historical phenomenon, a foreign power that it
must reckon with, that it can criticize and evaluate, but with which it has, by its nature and its presuppositions, no
kinship. That would, however, be a hasty conclusion.

For, in the first place, the religions that have appeared in history, at any rate their higher forms, have each in its time
been ethical powers. Ethical ideas have made up an essential part of the content of the positive religions. Seen from this
side, then, there must after all be a certain kinship between ethics and positive religion. --- In the second place, it is
a question whether, on the standpoint ethics takes, there might not be the possibility of a feeling which --- even if one
will not \textit{call} it religious feeling --- is nevertheless in its psychological nature akin to religious feeling as this
appears in the higher positive religions. I shall try in what immediately follows to describe such a feeling, while I postpone
the investigation of the ethical significance of positive religion to the next chapter.

2. By the \textit{feeling of life} psychology understands the feeling of pleasure or displeasure that answers to the course
of organic life in us. It hangs together, then, with the ease and strength with which breathing, the circulation of the blood
and the whole activity of nutrition go on in us\footnote[1]{Cf. \textit{Psykologi} VI A, 3 a. --- The following description of
religious feeling in the sense of a cosmic feeling of life is a further elaboration of what is briefly set out in
\textit{Psykologi} VI C, 8 b. In my \textit{Religionsfilosofi} I have attempted a discussion of the religious problem from the
epistemological, the psychological and the ethical point of view. And in \textit{Den menneskelige Tanke} (p. 369---392) I
have treated the religious problem in connection with the other main problems.}. As consciousness develops, feeling is
determined not only by our organic state, but
% --- p. 459 ---
\opage{459}also by a greater or smaller sum of ideas. Our feeling attaches itself to many things that lie beyond the
boundaries of our organism. In intellectual and aesthetic feeling our pleasure and displeasure are determined by the activity
of knowing and by the images that nature or art makes our imagination form. But when, with the help of thought, we have
reflected on our position in existence, have seen that with all our longing, all our plans and all our ideals we stand as a
single link in the great, incalculable series of causes and effects, --- then there arises a feeling at life, not merely at
the life that stirs in our own organism, but at the life that stirs in the whole universe of which we are members. Our
feeling of life is widened and determined by the course of life and development in the world, so far as we can form any idea
of it. There arises a \textit{cosmic feeling of life,} which forms an analogue to the organic feeling of life. From the organic
feeling of life the cosmic feeling of life differs by its content of thought, which is made up of all the experiences we have
won and all the ideas we have formed of the course of things. From intellectual and aesthetic feeling the cosmic feeling of
life differs by its personal and real character. In the activity of thought and imagination we forget ourselves over what we
think and intuit; in the organic feeling of life we have to do merely with ourselves as organic beings; but in the cosmic
feeling of life our pleasure or displeasure is determined by the position of our whole personality and our highest values of
life within the development of the world.

3. The cosmic feeling of life presupposes a view of the world. It need not, however, have any great apparatus of speculative
hypotheses. Where such hypotheses are developed, they will often, on closer investigation, even prove to be effects rather
than causes of the feeling. It can content itself with the two main thoughts that were stressed above (XXVIII, 5) as the most
important for intellectual culture, namely the thought of existence as a law-ordered connection and the thought of existence
as a great process of development.

Every link in the great causal connection is borne by
% --- p. 460 ---
\opage{460}other links and in turn bears others. The strength with which it fills its place it cannot have given itself
from the first; it has received it, and it must at every moment receive it anew in order to be able to maintain itself. Even
in my most energetic act of will, where my personal distinctiveness shows itself most strongly, and where my self-esteem seems
to have the firmest ground to build on, even there I am using up a capital that I did not myself bring about from the first.
I am dependent precisely in my highest activity, and the more so the more strength it lays claim to. But dependence appears
not only in the energy I dispose of being \textit{given} me, but also in the energy I dispose of being \textit{limited.} My
fate is interwoven in the great process of development, and my will can intervene determiningly in it only in part. I am
tossed about by hope and fear in the great swell of development and dissolution, arising and passing away, that nature
presents. Who is right, and on whose side is victory, --- with the forward-striving, developing tendency in the world or with
the hampering and dissolving tendency?

The relation between my will and my fate thus determines my feeling at life. This feeling can be of a purely or predominantly
egoistic kind. But when I do not think of my fate as this single individual, --- when my sympathy with others and my ethical
feeling are awakened, the cosmic feeling of life takes on another character. These feelings, which lead me to set up and work
for ideals lying far beyond my individual self-preservation, have nevertheless developed according to definite laws of
nature; and that such a development has been possible gives us a testimony that valuable forces are at work in existence.
(Cf. IV, 5). Nature stands before us --- precisely when we apply a naturalistic explanation to the origin and development of
conscience --- as a home of ideal forces. Whatever else development has brought with it, \textit{this} it has at any rate
also brought with it! There has developed an urge and an impulse of life of another kind than the mere physical urge for
self-preservation. What I feel in me, in my
% --- p. 461 ---
\opage{461}conscience, is just as much a world force as the forces that express themselves in the interaction of material
masses.

But here a significant difficulty arises once more. We cannot explain existence from ethical feeling. The world is not
constructed according to the principles our conscience sets before us as the highest. What is the highest end of our ethical
life of will we cannot declare to be the end of the development of the world. Not only is the very thought of an end of the
world just as impossible to carry through scientifically as the thought of a first cause. But the glaring disharmonies in the
world, the suffering and cruelty, the sum of misfortune and crime with which development is bought and progress --- insofar
as any progress can be shown --- is won, all this makes it logically and ethically impossible to maintain an ethical principle
as the source of the development of the world. Every theological and philosophical attempt to overcome this difficulty has
proved fruitless. Orthodox theology has merely pushed the question further back, and speculative philosophy has volatilized
and explained away the difficulties. The only way one can escape these difficulties is to refrain from thinking about them,
and this way out does not come equally easy to all individuals\footnote[1]{Several theological critics have been so kind as
to interpret this sentence as meaning that I myself was among those to whom it came easy to refrain from thinking about the
problem mentioned. They have not understood that the sentence is meant ironically. I was aiming partly at those whose
interest in such questions is dulled, so that they do not occupy themselves with them at all, partly at those theologians who
seemingly occupy themselves with them, but help themselves over the difficulties in so easy a way and with such superficial
arguments that this occupation cannot be called thinking. --- In my \textit{Religionsfilosofi}\textsuperscript{2}, p.
243---249, I have developed the thoughts that are for me the most essential with regard to the problem we stand at here.}.
Not out of pride, but precisely out of conscience and out of clear insight into the limits of our knowledge, we must leave
this question open. The philosopher sees no reason to take reason captive in favor of supposed solutions which, to unbiased
criticism, contain insuperable self-contradictions or come into conflict
% --- p. 462 ---
\opage{462}with what has been won by way of experience or thought, or which perhaps even present great ethical misgivings.
An honest doubt is better than a thoughtless faith. Life is full of great oppositions; but they first become contradictions
when one would make one member of the opposition into the whole and explain existence from it. From a good world principle
the evil in the world cannot be explained; from an evil world principle, such as absolute pessimism sets up, the good in the
world cannot be explained.

But our sympathy with the suffering in the world and our need to hold fast the ideal way of regarding things need not be
shaken either by our not knowing the end of the world drama, or even knowing whether it has any end. The good and valuable in
existence is a struggling power, and with every serious struggle uncertainty and tension are bound up. The cosmic feeling of
life will therefore swing between fear and hope even in those in whom it on the whole bears the stamp of confidence. ---

A feeling such as the one described here quite certainly exists. It marks a standpoint where religion's being a matter of
feeling has been taken seriously. It encloses no dogma, but arises quite simply from the relation of ethical feeling to the
real world. Those who feel a need for it may try with the help of speculations and symbols to win a further understanding;
there will be nothing questionable in this, provided only that two things are not shaken: the idea of existence as a
law-bound connection, and the independence of ethics from dogmatic assumptions. Religious thought can rise above scientific
knowledge only by setting up the demand for a completion of our view of the world, which it constantly
progressing experience makes impossible, and by interpreting the picture of the world that experience gives us by analogy
with, and as the expression of, the inner states of the life of the soul. But beyond speculative hypotheses one cannot get by
these ways\footnote[1]{Cf. \textit{Den nyere Filosofis Historie.} (The index under ``Ide'' and ``Analogi''). --- \textit{Den
menneskelige Tanke} p. 306---345.}. And far stronger than
% --- p. 463 ---
\opage{463}the theoretical interest in understanding the world will be, under the influence of religious feeling, the need for
self-understanding. The more strongly the feeling determined by the course and conditions of life stirs, and the more obscure
its origin is, because so many small and great experiences of life determine it, the more will man seek to win such
expressions for it that his inner life can become intelligible and intuitable to him. And not only to himself, but also to
others. For he will not be able to rest content with what stirs in him being wholly peculiar to him. Through signs and
symbols he then seeks to make known to others how he is attuned to life, in order to see whether something similar might not
be stirring in them. Through the involuntary need for self-understanding and for spiritual community, religious ideas
originally come into being. And because of the incompleteness of knowledge they cannot become anything other than symbols.

Religious feeling sets no ethical tasks that would not be set without it. It expresses the most intimate and highest
self-understanding man can win with regard to his position in the world and with regard to the fate that befalls that for
which he lives. It ``purifies'' the mind, as aesthetic feeling does; but the purification is here more far-reaching, because
it is man's own real position in the real life of the world that determines the feeling. Through its close connection with
ethical feeling the cosmic feeling of life is free of sentimentality and quietism. We climb a mountaintop, not in order to
build and dwell up there and withdraw from real life, but to breathe the pure and strong, perhaps also sharp, air up there,
and to widen and clear our gaze, so that we can better find our way in the narrower circumstances to which our practical
tasks bind us. There is an element of will closely bound up with religious feeling, which keeps it from becoming speculative
or sentimental rapture. The mood swinging between hope and fear, called forth by the changing fate of man himself and of the
highest values of life, leads to the peculiar concentration and deep self-reflection in which religious feeling consists; but
out of this same concentration springs
% --- p. 464 ---
\opage{464}the will to hold out on the side of the good cause, to work together with all the powers in existence that serve
those values of life. Religious faith, which now has more the character of patience, now more that of confidence, is a
willing that holds fast the thought of the ideal in spite of all inner and outer checks and resistance and in spite of all
uncertainty and doubt. There is here an inner field of work of a central kind, where endurance and courage can be shown no
less than in the domain of outer life. From the purely ethical side, however, the tasks that lie here have all been mentioned
already earlier (see X, 3---4).

4. The feeling described does not stir equally strongly in all individuals. Some natures are especially receptive to it, as
other natures are to intellectual or aesthetic feeling. And it is of so compound a nature that it cannot be surprising if its
elements are found in most divergent relations to one another in single individuals, so that it appears with very different
timbre. --- Already the relation between the cosmic and the organic feeling of life can be most different. They can often
stand in decided opposition to one another. Temperament does not unconditionally govern the way we feel the world. He who is
melancholic in temperament, and whose organic feeling of life thus predominantly has a dark and heavy character, can therefore
well be an optimist in his view of the course of things in nature and history. The darkness in the world may be confined to
the small place he himself fills, and by an effect of contrast the rest of the world may precisely come to stand in a bright
light for him. Conversely there are those whose happy organic basic mood is interrupted and disturbed when they cast their gaze
around them in the world. --- Some will lay great weight on the enigmatic and unfathomable in existence, others on the lawful
connection that science is gradually able to show in existence. --- In some it will be this, in others that class of
experiences that especially determines the feeling, --- in some, for example, especially the life of nature, in others human
history. --- In some (e.g. Spinoza) the cosmic feeling of life has essentially the character of resignation, in others (e.g.
Bruno and Fichte) it has the character of enthusiasm and
% --- p. 465 ---
\opage{465}the impulse to activity. --- A decisive difference lies between optimism and pessimism. The course of life proves
both partly right, and they are not incompatible, except when each is formulated as a speculative system. In Plato and in
Christian theology on the one side, and in Schopenhauer on the other, the two opposed forms of cosmic feeling of life are
developed into diametrically opposed structures of thought. --- A considerable difference will come to light in the
originality and independence with which individuals are able to shape and express the ideas in which their religious feeling
finds vent. In some these ideas will depend essentially on the inherited religious forms of the family or the nation; in
others a life of feeling similar to the one from which the inherited form of religion sprang will rise again in an
independent way, except that the changed culture-historical circumstances will bring new nuances or a special stress on
single sides; in others again the life of feeling, in interaction with the experiences of life, will unfold in so peculiar a
way that they cannot find understanding among those around them, and also cannot acknowledge the current symbols as an
expression of the way in which life stirs within them. People of this last kind must then go their solitary ways, insofar as
they do not have the symbol-creating capacity and the depth of personality that can make them religious models and
leaders\footnote[1]{Cf. the fuller account of differences in type and disposition in \textit{Religionsfilosofi}\textsuperscript{2},
p. 114---128; 258---282.}.

But in spite of these differences there is nevertheless here a common spiritual domain, where in the main features a common
development too may be possible, and where the one need not stand in the other's way. Where the feeling described is genuine,
it is determined by the real world and its conditions, and it is, after all, the same world that appears to all, even if their
original dispositions of feeling (temperaments) and their experiences are different. When they properly familiarize themselves
with how the world is, they must also take into consideration the very circumstance
% --- p. 466 ---
\opage{466}that it can and must be seen from so many different standpoints. We know the world, after all, only through the
images of it that form in each one's consciousness; and in these images it is not the same features that everywhere appear
most clearly and vividly; nor are they all seen in the same light. The world thereby becomes precisely richer and more
manifold, and sympathy is widened, the more one knows how to enter into others' ways of conceiving and feeling the world. In
every more intimate life together between human beings a community, or at least familiarity and understanding, must assert
itself in this respect. The spiritual help that one individual can bring another consists especially in furthering the truth
and clarity of his feeling at life. For that it is not necessary that the one individual stand on quite the same standpoint
or belong to quite the same religious type as the other; where there is a sense for personality and for personal truth, one
will be able to help another to full clarity and to complete expression of his being, even if one's own being is of a wholly
different kind. Some individuals are especially suited to afford such cure of souls, and often practice it without knowing
it. But no one is wholly cut off from it. The universal priesthood that Protestantism proclaimed, but never took fully
seriously, can become truth the more the life of feeling is allowed to develop freely and distinctively in single individuals,
and the more sympathy and psychological sense grow at the same time. Physicians and pastors are the nearest to exercise a
special influence here. But the former have still too many physiological problems to solve to be able properly to develop
their psychological sense; they are too inclined to regard the life of the soul from its outer, material side. And the latter
must, in accordance with their standpoint, first and foremost make sure under which rubric of dogmatics they are to place the
psychical phenomena before them; they lack the free human gaze. ---

Even if one wishes to call the cosmic feeling of life religious feeling, it in any case needs to found no church and to
establish no cult, just as little as it needs to rest on any dogma. The community it leads to is
% --- p. 467 ---
\opage{467}the freest of all; it expresses itself only in mutual understanding and help. Its church is great nature itself,
and its cult is work, life together with human beings and with nature, life in science and in art. Quite peculiar conditions
are in any case required, as we shall see, for a common symbolism and a common circle of outer forms to form.

\chapterhead{XXXII.}{THE SOCIAL-ETHICAL SIGNIFICANCE OF THE POSITIVE RELIGIONS.}

1. There are exceedingly many nuances and forms of religious life and religious faith, and when freedom once comes in
earnest to rule in the spiritual domain, there will be still more than now. Historically \textit{the positive religions}
have so far had the greatest ethical and social significance. Two things characterize positive religion: cult and dogma.
Religion is not from the first merely feeling or a purely subjective relation to the powers of existence; it is a way by
which man thinks to attain definite practical ends, to provide for his bodily and spiritual salvation\footnote[1]{As I have
sought to show in my \textit{Religionsfilosofi}\textsuperscript{2} (p. 231---234), the decisive thing in the concept of
\textit{positive} religion is that it builds on historical tradition, and that there is a community in the forms and
expressions of religious life. ``Positive'' religion is to be understood on the analogy of ``positive'' law. Cult and dogma
are thus only special, historically conditioned expressions of positive religion; they are derived relative to what is
essential in this concept. A religious society could be conceived with historical tradition and common symbols without cult
or dogmas in the strict sense of these words. --- I have, however, not wished to change the account here in the Ethics,
since the usage is not of decisive importance.}. \textit{The act of cult,} the oldest and most constant element in a
positive religion, has a magical
% --- p. 468 ---
\opage{468}character, since through it man thinks to intervene in a supernatural order of things, or himself to stand as
the object of influence from a supernatural order of things. In the act of cult man meets his deity directly; therefore it
stands as the central element in positive religion. Community in cult is the original religious community, and a community
of great significance, since what binds is the highest and most decisive experiences known. \textit{Dogma} is an assumption
that builds on a supernatural revelation, which unveils the mystery of existence, an authentic explanation from the source of
existence of what the meaning of the whole of existence is. In its cult and in its dogmas, then, positive religion has not
merely human thoughts about existence and human feelings at it, but the deity's own thoughts and present working. As soon as
dogma is made a symbolic concept and the act of cult a beautiful custom, we are outside positive religion. This stands and
falls with the punctual presence of the deity at a definite time and place and under a definite form. The positively
religious consciousness does not distinguish between symbolic ideas and objective truths that unveil the innermost of
existence, nor between symbolic customs and actions that intervene in the innermost of existence. The mighty symbolic forms
that positive religions have created became possible precisely because human beings did not yet know the difference between
symbol and reality, which is the work of reflection and analysis. When this work has been done, the great images can still
keep their value as symbolic expressions of human experiences and feelings that are constantly being had anew; but it is a
question whether the symbol-forming activity would have been released from the first with so high a degree of energy if it
had not sprung from the firm assurance of having to do with objective realities. The undivided state, unweakened by doubt and
reflection, has been a condition for the production of some of the mightiest images and forms the human race commands. It is
an essential side of the religious problem whether the great productive power can be preserved when the distinction
% --- p. 469 ---
\opage{469}between symbol and reality has come into force. But this much is certain, that some of the human race's most
significant experiences are woven into cult and dogmas. In the worship of common gods the community of the family, the tribe,
the people or mankind is felt. The god is offended when an offense is done to one of his worshipers. In the loftiness of the
deity the feeling at what is great and mighty in life and in nature finds its expression. And that even the gods suffer, and
perhaps suffer willingly, hallows suffering and makes it easier to bear. Constantly and involuntarily the thoughts and actions
that were thought to have validity and significance for another world were formed on the basis of influences, experiences and
moods from real life. It is not least on this that the ethical significance of the positive religions rests.

The history of the development of the positive religions shows us that they stand in a \textit{relation of interaction with
practical ethics} (positive morality, see I, 2). They have exercised great influence on it, but have also themselves been
influenced by it. The adherents of the positive religions can from their standpoint assume only the first, not the second. In
their religion absolute truth is revealed to them; from it human life is to learn, but it has nothing to learn from human
life. The historical view, on the other hand, for which everything that appears in history is also itself a result of
historical development, holds that the ethical influence of religions first becomes possible through their taking up into
themselves ethical ideas that have developed to a certain degree independently of religious influence. Man must know from his
own experience the qualities he ascribes to the deity. In human life love and justice must have proved their significance
before the thought of a loving and just deity can arise. How else should man be able to connect any meaning with these words
when they are used of the deity? A potentiation, an idealization takes place: the qualities are ascribed to the deity in a
degree and fullness that passes all understanding. The positively religious feeling is determined by the idea of one or more
beings far exalted above human nature and human
% --- p. 470 ---
\opage{470}conditions, whose capacity far surpasses what human understanding can encompass, but who are nevertheless thought
of on the analogy of human beings. All the single qualities ascribed to them must therefore undergo an extension and
heightening, but experiences must lie at the bottom that can be extended and heightened. The world of the gods is endowed with
the qualities that stand before human consciousness as the most excellent. We thereby understand that as human beings are
improved, or at any rate as customs are softened, the character of the gods too is improved and softened. In place of the
wild and cruel gods of war, gods of love and mercy gradually come. The history of religions shows us a progressive
humanization of dogma and cult. Religion and morality hardly have a common origin. They probably joined with one another only
at a later stage of development: ``the religious relation has in the course of time been moralized''\footnote[1]{\emph{Chantepie
de la Saussaye}: \textit{Religionsgeschichte.} Freiburg 1887. I. p. 35.}. --- At lower stages the gods' thirst for vengeance
can be quenched only by blood. But real human sacrifices gradually fall away or are replaced by symbolic customs, whose
original meaning believers seldom properly picture to their imagination. In Greek religion this development can be traced
especially clearly. In the Jewish religion sin offerings played a great part; but already the Jehovah of the prophets does not
thirst for blood, and love is better to him than many sacrifices. Through Paul's influence the idea of the sin offering to the
angry God nevertheless passed over into ecclesiastical Christianity, and it is often still held fast there in a form that quite
recalls the primitive ideas\footnote[2]{A clergy meeting in Lund (September 1864) declared: ``No Christianity without an
atonement --- and an atonement in blood!'' and referred expressly to the fact that ``the history of religions speaks of an
atonement in blood through the bloody sacrifices that occur in so many religions.''}. --- In the idea of hell and eternal
damnation we have another example. Even if this idea is still held fast by orthodox church doctrine, it is far from being held
with the glowing lust for vengeance and the lively
% --- p. 471 ---
\opage{471}imagination of the Christians of antiquity\footnote[1]{Tertullian and Cyprian looked forward to seeing their
persecutors tormented in the flames of hell while they themselves sat saved at God's right hand. In Augustine and Thomas
Aquinas a softening has already set in, the sight of the torments of the damned being really only a contrast that is to give
the saved an all the stronger feeling of the grace that has befallen them. --- In our days zealous pastors are said sometimes to
console those whose nearest ones they take to be damned with the thought that all memory of these damned will be blotted out in
the saved, so that the latter's own bliss will suffer no harm from the former's torments. Many, however, will think that no
very great advance in humanity is marked by this assumption. --- One is expressed, on the other hand, in Schleiermacher's
objection to eternal damnation: the joy of the blessed would be spoiled by the thought of the torment of their fellow men!}.
--- The foundation has changed; the religious ideas formed under the influence of the life of feeling of earlier times still
hold their ground, no doubt, but another meaning is put into them, or they are no longer brought before consciousness in their
whole sharpness and clarity. Sympathy, love of mankind, has grown, and therefore imagination too has lost its energy in this
direction.

2. It is not only the results of \textit{ethical} development that are thus translated into dogma and cult. We first
understand the significance of the positive religions when we see them as \textit{condensations} of \textit{all} sides of
spiritual life. The influence of intellectual and aesthetic development can be traced no less than that of ethical-social
development when we compare higher forms of dogma and cult with lower. --- The higher religions have taken up more rational
knowledge into themselves than the lower. The idea of the deity is always brought to a certain degree into harmony with the
growing knowledge of the connection of nature. The development from fetishism to polytheism and from that again to monotheism
is in part conditioned by a growing understanding of the connection in nature, which compels a limiting of the intervention of
supernatural powers. The higher religion is rationalistic relative to the lower. --- The aesthetic capacities too are laid claim
to by religions. Imagination seeks to form as living, intuitable and appealing images as possible of the beings that are the
objects of faith.
% --- p. 472 ---
\opage{472}It strains to picture what is great and exalted.

Positive religion is originally the only form of ideal culture. Ethical, intellectual and aesthetic life are led essentially in
the form of religion. All sides and directions of spiritual life are condensed and given shape in dogma and cult. Therefore
positive religion will not satisfy any single side of human nature by itself alone, neither thought, imagination, feeling nor
will; but it turns to them all and is in \textit{its classical times} everything for man in spiritual respects. No special
science, no special art, no special ethics is then known, and therefore religion itself does not become a special thing beside
other endeavors either. The significance of religion for the history of culture rests on its remarkable power of condensing so
much of what stands or stood before human beings as true, beautiful and good. With the collected force won by this
condensation, religion works on human beings. It nourishes and educates them with elements it has itself absorbed and
transformed. --- It is the task of the history of religions to show the different processes of condensation that have taken
place here, and the elements that have been taken up at different times. A great religion is not shaped all at once or by a
single human being. Many generations and personalities of various kinds contribute to cult and dogma through the involuntary
projection of their experiences of life. Religious genius shows itself in the capacity to initiate a new condensation of the
elements of spiritual life at a given time; and religious sense consists in the need for, and receptivity to, spiritual
nourishment in this form.

3. With the continued development of culture a division of labor sets in in the spiritual domain. Dispersion takes the place
of condensation, analysis the place of synthesis. Each single side and direction of spiritual life now lays claim to special
attention and energy. A special science and a special art arise, and ethics seeks its independent foundation, independent of
dogma and cult. Then strife arises between faith and knowledge, between theological and philosophical
% --- p. 473 ---
\opage{473}ethics, between church and state. At this stage of development positive religion comes to suffer from a
difficulty, a self-contradiction, called forth by what was its strength at the earlier stages: it is to express and determine
the \textit{whole} of spiritual life and yet comes itself to stand as a \textit{special} expression of life \textit{beside}
science, art, ethical life and the activity of culture.

A positive religion's ethical significance will rest partly on the definite way in which it unites the different
(intellectual, aesthetic, ethical) components, partly especially on what ethical ideas it has taken up into itself. With its
dissolution a significant ethical force may perish. Partly there are natures that are receptive to spiritual nourishment only
in condensed form, and that therefore now feel split and empty. Partly it is not certain that the single elements are able to
work as strongly outside the peculiar religious
condensation as within it. Oxygen, after all, can as a component of carbonic acid help to call forth effects it is unable to
call forth in free form. Just as little as a positive religion's ethical significance follows of itself, just so little does
it therefore follow of itself that its dissolution is an advance.

The components of a positive religion's content \textit{need} not fall away with its dissolution. The human experiences that
find their potentiated and idealized expression in religious ideas can always be had anew. Therefore religious motives form no
complete opposition to other motives, but can include them within themselves. Love of mankind is an ethical motive; faith in
God as the one whose being is love is a religious motive; but this faith presupposes that one knows love from one's own
experience. This relation between religious and ethical motives makes a comparison of the ethical significance of the different
religions so difficult. It is not easy to decide with what degree of independence the single element works, and how much
depends on the form and connection in which the religion lets it appear.

Even if the forces that work in a concentrated way in positive religion could very well accomplish the same by working each
by itself,
% --- p. 474 ---
\opage{474}there nevertheless remains the drawback that a collected expression for ideal culture is lacking. A natural need
will always assert itself to receive a collected influence of life, instead of always having to do only with a single side of
it. The division of labor leads to dividedness and disharmony when new processes of condensation cannot take place. It will be
of essential significance for the future of the human race whether such condensations can be brought about without their
needing to take on the shape of dogma or of the act of cult. This would require an intimate interaction of science and art, of
theory and practice, of knowledge and feeling. The great basic experiences and basic moods of human life would have to find
their expression in symbols which, without being fixed dogmatically, speak just as much to feeling as to understanding, just
as much to will as to imagination. Whether this task can be solved the future must show. Human consciousness is not yet ripe
for it. We live for the time being in the age of criticism and analysis, and bear its marks in our whole spiritual life. Under
the overwhelming impression of the religious traditions so much work must be spent on spiritual liberation that energy is not
easily left over for the positive shaping of a new view of life. And this struggle for liberation must moreover be waged by
individuals each by himself. Indignation at the way in which independent religious development is crushed and hampered then
easily comes to play a greater part than the need to relieve the spiritual distress in the race. An energetic religious
thinker, \textsc{Christoph Schrempf}, has --- on the occasion of the fiftieth anniversary of the founding of a free
congregation in Königsberg --- expressed himself on this\footnote[1]{\textit{Die Wahrheit.} Halbmonatschrift zur Vertiefung in
die Fragen und Aufgaben des Menschenlebens. V. Stuttgart 1896. p. 261.} in the following way: ``The soul of this [the
free-church] movement was that the duty and the right of independence forced themselves forward with power. That has given me
something to think about. Me too self-assertion has brought into conflict with the church, and where religious tension
% --- p. 475 ---
\opage{475}is at the moment great, one may suppose that it is caused especially by resentment at religious
subjugation\dots{} I see in this a sign that we are still in the preparatory stages, not yet in the decisive movement of the
religious consciousness. This movement will not proceed from a striving for independence,\dots{} but from a newly awakened
feeling of brotherhood\dots{} In particular the leading spirits will not be moved so much by care for their own religious self
as by compassion for the people, who are now again as sheep without a shepherd. Naturally this compassion, if it is of the
right kind, will in turn enjoin greater respect for every human being's self. It will not find man's true distress in
calamities, nor in social inequality, but in the danger of losing his self. It will thus see its main task in education to
independence. But the real driving force of religious reform will nevertheless not be resentment at religious subjugation, but
--- as in Jesus and Luther --- compassion.'' % translated from Høffding's Danish
It is aptly said here that the need for spiritual self-assertion in those who free themselves from religious tradition is
nowadays easily developed at the expense of the capacity for devotion and sympathy. --- Attention has rightly been drawn to yet
another consequence of religious emancipation. \textsc{Stanton Coit}, one of the most energetic champions of an ethics
independent of religion, says in an address on ``The Dangers of Radicalism in Religion'': ``The mere casting away of inherited
doctrines and striving after ideas of one's own gives man a certain impulse to constant change. He must always have something
new\dots{} Newly fledged freethinkers are like birds which, when they have escaped from the cage, fly restlessly on from branch
to branch, from valley to mountain and from mountain to lake, without ever stopping to build a homely nest for their heart,
but always hasten on and on, until at last their wings are paralyzed and they fall to the ground in a desolate
desert''\footnote[1]{\emph{Stanton Coit}: \textit{Die ethische Bewegung in der Religion.} Uebers. von Gizycki. Leipzig 1890.
p. 93 f.}. % translated from Høffding's Danish
This restlessness too, which makes it difficult for many to attain a new inwardness, a new spiritual home after giving up
% --- p. 476 ---
\opage{476}the old, snug home in the shelter of tradition, is a moment that hinders positive new formations in the domain of
views of life.

4. Yet another very important side of the positive religions must be brought forward. The dogmas and forms of cult that the
religious process of condensation forms are common to greater or smaller circles of human beings. Positive religion and
\textit{church} cannot be separated. It is not in the isolated individual that the religious process goes on; what the
individual utters is only what more or less stirs in many, an adherence to or a carrying forward of what is given in
\textit{tradition.} In the domain of positive religion even the most productive individuals feel themselves not as those who
give, but as those who receive. Individuals are held together around a common tradition; and religious quarrels arise when the
question arises in what direction tradition is to be carried on. A positively religious society, a church, arises not by the
free union of individuals who seek one another, but by the same word sounding to all. --- In its simplest form this society
consists of a single family, within which the spirits of the dead are worshiped from generation to generation. The hearth fire
is kept up in their honor, and their commands are respected as highest law. Already here there is a sacred history. In the
practice of such a \textit{family religion} only those take part who belong to the kindred; the father of the family, who leads
it, stands as the highest living authority in religious as in other respects. Even today this kind of cult is the most
widespread on earth, existing under, beside or as a member of the national religions. Only in Christianity and Mohammedanism has
this worship wholly vanished\footnote[1]{\emph{Henry Maine}: \textit{Early Law and Custom.} p. 57 ff.}. \textit{National
religion} arose through the whole state having its common gods and heroes and its common hearth. The state, like the family,
was originally both a religious and a political society. There was no difference between state and church, as little as
between religious, juridical and ethical laws. It was every citizen's duty to take part in the public cult, but strangers were
excluded
% --- p. 477 ---
\opage{477}from it. Each state had its gods, as each family had its own. One fought for one's fathers' gods as for one's
fathers' land. --- Already here the religious society can begin to appear different from the political when a special priestly
class develops whose office becomes the preservation of the old traditions and customs. But the difference first comes forward
decisively when \textit{the idea of the church as a universal religious society} arises. Differences of family and national
barriers now acquire only subordinate significance. There is no longer to be Jew or Gentile, Greek or barbarian. A purely
spiritual society is to be founded. Here too, however, the religious society, the spiritual unity, is conditioned by the common
tradition of a sacred history, from which individuals draw the highest law for their life and firm assurance for their fate.
For European mankind Christianity made this idea prevail. Greek philosophy had, no doubt, by its own way and taught by the
historical experiences after the time of Alexander the Great, arrived at the idea of the equality and common nature of all men
and of universal love of mankind. And without this whole change in ways of thinking a world religion would hardly have
developed out of Judaism either. Here development in the Greek world prepared Christianity. But precisely here the significance
of the religious process of condensation shows itself: for only as elements in a positive religion have community and love of
mankind been able to win a wider acknowledgment.

Thus positive religion has proved its power of founding and preserving society in an ever wider extent. The high point is
reached in the idea of the church as a universal human society. We shall then undertake a closer examination of this idea.

5. The universal human society that positive religion proclaims is not supposed to have developed in a natural way; its
founding is a supernatural act that must be repeated each time a new individual is to be received into it. The individual
cannot gain access to it by his own endeavor. From outside, as historical tradition, the account of its founding comes
% --- p. 478 ---
\opage{478}to him. The society is to be universal, to be sure, but the tradition into which one must live oneself in order to
be a member of it flows nevertheless through a narrow channel. He who does not meet this stream on his way is lost. Everything
then depends on whether tradition reaches us and whether it is genuine. For this we must have guarantees, and faith in the
guarantees therefore necessarily becomes the most important thing. Positive religious faith therefore becomes with logical
necessity faith in the church, the channel through which religious tradition comes to us. This transition from faith in what is
guaranteed to faith in the guarantee is seen in the recent history of both Catholicism and Protestantism. The dogma of papal
infallibility, which was proclaimed only a millennium and a half after the other dogmas, is especially characteristic here.

Just as positive religion, as regards its content, came at further advanced stages of culture to suffer from the contradiction
that it is to be everything and yet becomes a special expression of life beside others, so here, as regards its foundation, it
comes to suffer from the contradiction that as the condition of the universal ethical human society it sets up a faith which can
be lacking without any of the noblest ethical qualities therefore needing to be lacking, and which contains very much whose
value for life is difficult to show. Here too (as with material culture, see XXVII, 1) we groan under the weight of an apparatus
that was after all brought about
to make life easier to live! He whose reason does not allow him to accept dogma and do homage to the authority of the church
can therefore very well ``hunger and thirst after righteousness,'' and be ``meek and lowly in heart.'' % Matt. 5:6; 11:29 (KJV)
Indeed, it may be that it is precisely hunger and thirst after righteousness that makes it impossible to accept dogma. (Cf.
XXXI, 3). There may be points in dogmatic doctrine that conscience, in spite of the most honest examination, cannot
acknowledge. To such things the church has only the answer that it is self-righteousness and pride. It demands obedience as the
first condition. (Cf. X, 4).

The contradiction that comes to light here is a
% --- p. 479 ---
\opage{479}contradiction between faith and love. Faith sets barriers, love abolishes them. Faith separates, love unites.
This contradiction runs through the whole history of the church. It was love of mankind that burst the national barriers; it
will also be able to burst the dogmatic barriers. Its influence has helped to bring about \textit{religious freedom.}

6. The principle of religious freedom abolishes all outer, civil and political differences that were formerly tied to
differences of faith. It abolishes first and foremost the punishments formerly set for transgression of religious duties and
for unbelief. Slowly in this way a series of crimes that formerly stood in the first line has vanished from the civil law
books\footnote[1]{\emph{Durckheim}: \textit{De la division du travail social.} Paris 1893. p. 172---177.}. Yet it is not to
be conceived as a purely negative principle. It is not a proclamation of indifference, but a natural consequence of widened
community and of love of mankind. It not only makes it possible for everyone to withdraw undisturbed by others into his own
inner self, but rests on the presupposition of a common human foundation behind all differences of faith. This view is of
importance; for only by stressing it can one prevent the great social force that lay in the positive religions from being
wholly lost. There is a great force in the union that rests on a common view of the greatest problems of life. When we are
allowed each to be saved after his own fashion, it can easily lead to each of us going his own way and losing mutual
understanding. But even if indifference and blasé superiority also find their account in religious freedom, that is
nevertheless not its essential side.

In religious freedom there lies, in the first place, the presupposition that ethics is independent of religion. Difference in
confession of faith, the opposition between accepting and rejecting a dogma, cannot be one with the difference between good and
evil. On a standpoint where faith is the only principle of all good,
% --- p. 480 ---
\opage{480}unbelief of all evil, religious freedom is impermissible. However great a difference one may make between civil life
and the inner life of conscience, this difference can never become so great that one could give people who lack the first
conditions of an ethical conduct of life full right to active participation in civil and political life. In civil and
political life there is both room and use for the highest ethical qualities, and no one can be acknowledged as an independent
member of his people when the possibility of his displaying such qualities is wholly excluded. According to strict orthodoxy
the unbeliever has, after all, even committed the greatest sin of all, the mother of all other sin: that of not being willing to
take reason captive. With such a view the acknowledgment of religious freedom is evidently incompatible.

In the second place, religious freedom rests on the presupposition that each individual has his personal distinctiveness,
which it is of importance to get developed as fully and independently as possible. Whatever he believes or does not believe, he
is to be treated as an independent personality, and he is to win truth by his own distinctive way. Every personal and
distinctive expression of life has its value when it does not stand in the way of more comprehensive considerations. Even in
the simplest spiritual activities individual differences assert themselves, to which attention has only properly been drawn in
the most recent times. How much more, then, must the higher spiritual life take different forms, and how careful one must be in
drawing inferences from outer likenesses to inner ones! Word, symbol and action can be the same, and yet there can be the
greatest difference within. Only he who has the capacity to distinguish between shell and kernel discovers these differences,
this richness. He who stops at the positive or negative confession of faith gets no further than the surface. The flower that
grows in the tall grass is discovered only when one bends the long stalks aside. Spiritual kinship proves to be other than
dogmatic kinship. And not only can a most different inner life of feeling lie behind the same outer confession of faith; but
behind different
% --- p. 481 ---
\opage{481}confessions of faith a closely kindred life of feeling can also lie. It may be difficult to penetrate to it; we
are still far behind in psychological sense for personal life in all its forms, and dogmatic quarrels will still for a long
time be a hindrance to the full development of this sense. But with religious freedom the way is paved for its development.

7. Common interests and common life in civil society and in the service of love of mankind will gradually lead both to a
fuller acknowledgment of the independence of ethics and to a better understanding of personal life in spite of the different
forms of faith. The more domains there are where work can be done without differences of faith playing any part, the more such
acknowledgment and understanding will develop, even if it happens half unconsciously. Here as elsewhere (cf. XIII, 4) common
activity precedes sympathy, practice precedes theory.

The church from its standpoint naturally cannot accept these consequences of the principle of religious freedom. Just as the
Roman Church continues to protest against the abolition of the church's secular power and against the doctrine that church
discipline may not entail temporal punishments, so every orthodox church must hold fast the difference between faith and
unbelief as the greatest opposition of life and regard the ``relative'' virtues that can be found in unbelievers as vanishing in
relation to that absolute opposition. But life goes its way in spite of all dogmas and unfolds in practice the consequences that
dogmatic prejudices prevent from being drawn in theory.

On the other hand there is nothing that prevents the freethinker from drawing the consequences mentioned at once. The time of
criticism and polemic against the positive religions is not over, to be sure. Here is a struggle that must for a long time yet be
taken up again and again. But the freethinker, who cannot be hindered by dogmatic prejudices, ought to feel that he has a far
greater obligation than the believer to show love of mankind and understanding toward those of other faith. Only if our thought
and our feeling have really been widened has liberation from the doctrines of positive religion any ethical significance. Not to
believe anything is about the least that can be said of a human being, and cannot in
% --- p. 482 ---
\opage{482}itself give him any value, especially since a negative confession of faith can be adopted in just as dependent and
thoughtless a way as a positive one can. In the freethinker who takes life seriously it is not merely ``cold understanding'' that
rules. It is free conscience every bit as much as free inquiry that animates him. \textsc{George Eliot} has expressed this for her
part in the following way: ``I say it now, and I say it once for all, that I am now determined in my conduct by far higher
considerations and by a far nobler idea of duty than was ever the case while I held evangelical views''\footnote[1]{\textit{Life of
George Eliot.} I, p. 148. (Tauchnitz Edition).}. % translated from Høffding's Danish

Of what use is all our historical research, all our absorption in different periods of culture, if a greater and more thorough
understanding is not thereby developed of the human life that is led around us, even if it clothes itself in forms and symbols
that do not agree with our views? --- If the freethinker does not surpass the believer in sympathy and in historical-psychological
understanding of spiritual life, mere negative criticism will not be sufficient to give him any real superiority.

Whether the religious strife itself will ever end, --- whether at all times believers and freethinkers will stand against one
another, about that we can know nothing, and it does not concern ethics either. As long as the opposition lasts, a constant
interaction will take place between the two tendencies. Some feelings and ideas are favored by the one, others by the other
tendency; and if these feelings and ideas are of essential significance for human life, they must be appropriated also by the
tendency that has not itself fostered them. How much, for example, has not the free inquiry of our time learned from Romanticism
and the religious reaction at the beginning of the century? And how great an influence has not the influence of experience and
science constantly had on religious consciousness, not only through their results, but especially through the method of thought
that they slowly train human beings to use? --- In this way development
% --- p. 483 ---
\opage{483}goes forward --- slowly and through many oscillations. But whether a lasting harmony will ever be reached, only a
very distant future can show.

\chapterhead{XXXIII.}{STATE AND CHURCH.\footnote[1]{Cf. with this chapter my essay \textit{Stat og Kirke} in the journal ``Det ny Aarhundrede'' (1904).}}

1. The church not only has been in the past, but still is and will for the time being continue to be one of the greatest powers of
culture at work in human life. This was clearly shown by the religious movement in the nineteenth century, when the church
blossomed forth into inner life and outward power\footnote[2]{Cf., e.g., \emph{Taine's} interesting account of the development of
the Catholic church in France after the Revolution. \textit{Le régime nouveau.} II. p. 89---151.}, although shortly before it had
seemed to be succumbing to outward enemies and to be forsaken by all spiritual forces. On account of its great influence and its
great vitality it falls within the purview of social ethics. Quite apart from the objective validity of the foundation on which the
church builds, it is of decisive significance that it should come to make --- within the right limits --- the contribution to
development that it is able to make.

The church has a great ethical responsibility, which its representatives, with a strange blindness, seem precisely in the most
recent time not to realize. The more enlightenment and knowledge spread in wider circles, the more doubt and criticism of the
dogmatic assertions of church doctrine will inevitably spread as well. The intellectual difficulties of church faith will be felt
by ever more people. But precisely in the most recent time the church has sharpened its position: without faith no ethics! From
this position the church must find it the right consequence when he who rejects faith
% --- p. 484 ---
\opage{484}thereupon throws himself into bestiality. The apostle Paul, after all, already spoke the fateful words: ``If Christ be not
risen, let us eat and drink; for to morrow we die.'' % Høffding's paraphrase of 1 Cor. 15:14, 32; KJV wording
If this line of thought gained currency --- and certain preachers do all they can to that end --- the prospect for the future would
be grim. The church does not see what a game of hazard it is entering upon here. In blind confidence that its dogmas will
triumph, it stakes everything on this one card. Whatever merits it may otherwise have earned from mankind, it will have a heavy
account to settle if the orthodox-pietistic tendency continues to determine its conduct toward the great masses, whose most
important spiritual leader and educator it still is. This account will be demanded of it in the name of mankind. For every
church, let it call itself by the proudest names, is after all only a sect in relation to mankind. The church exists for the sake
of human beings, and human beings do not exist for the sake of the church.

It is the church's tendency to make ethical questions dependent on dogmatic principles that above all makes it necessary to set
definite limits to its activity. This is not to deny that the church exercises a great educating influence on those who attach
themselves to it. Where ethical ideas would not be able to work by their own strength, they can work as elements of religion, and
where intellectual and aesthetic culture could not otherwise reach, they become possible in the form of religion. For many people
of all strata of society religion still stands as the epitome of all spiritual culture. It is by no means the case that it is only
``the great multitude'' that cannot do without religion. Among peasants and workers there are critics and doubters just as well as
among the so-called educated, and the criticism and doubt are there often precisely the most genuine and most serious, because
they have grown out of uninfluenced thinking and experience. There are far more people in ``the great multitude'' than is usually
suspected whose view of life is formed in an independent way. The culture of the ``educated'' does not always bear the stamp of
independence; for many it consists in sharing in a certain intellectual tradition, a tradition that is often adopted too early and
with too little
% --- p. 485 ---
\opage{485}work of one's own. And even among the truly educated there are undoubtedly always some for whom it is a spiritual necessity
that all ideal culture, all understanding of life and all feeling for life should in the end condense into great, richly colored
symbols, in which the distinction between image and reality is no longer to hold. It is this that makes the religious problem so
complicated.

2. The connection between church and state is no accident. On the one side, the church had to demand to permeate the whole of
society if it was to provide the epitome of all spiritual culture and be the universal human community. It had to regard the state
as its servant. On the other side, the state --- as long as ecclesiastical ideas held undisputed dominion over minds --- had to
believe that it had a religious mission and to regard it as its highest task to spread religion and protect it. Yet the state did
not regard itself merely as a means for the church. It also regarded the church and religion as its means. It considered religious
faith a necessary condition for its being able to have good citizens, and for security and peace to prevail in the land. Heresy was
forbidden not only because it was an affront to God, but also because it was harmful to the peace and welfare of the
realm\footnote[1]{Cf. F. \emph{Pollock}: \textit{The Theory of Persecution.} (Essays in Jurisprudence and Ethics. London 1882). p.
160.}.

As long as positive religion was the only spiritual power of culture, the dominion of the church over the state, or at any rate in
the state --- as church-state or state church --- had its justification. But when the times changed, the church itself began to
see that religious life fared best in freedom. And since modern states now comprise people with different religious standpoints,
--- since many domains that could formerly be looked after by the church demanded independent treatment, --- religious freedom has
more and more been acknowledged or demanded both on the part of the church and of the state. A person's relation to church and
religion is then to have no consequence whatever for his position in the state, is not to deprive him of any of his rights or free
him from any
% --- p. 486 ---
\opage{486}of his duties as a citizen of the state. In this respect, however, we still live in a state of transition. Civil life is
still by no means freed from ecclesiastical interference. Many forms and institutions remain from the time when the state regarded
either itself as the servant of religion or religion as its means. A couple of examples of this may be mentioned.

By means of the oath the state sought to obtain the greatest possible guarantee that an individual was speaking the truth or
promising honestly. Such a guarantee it found by appealing to what must be most sacred to the individual. The oath is properly
speaking\footnote[1]{\emph{Post}: \textit{Die Grundlagen des Rechts.} p. 441 f.} a remnant of the ``judgment of God'' so much
used among savages and among the barbarians of the Middle Ages: he who swears calls down a supernatural punishment upon himself.
But the oath was not meant to be a confession of faith. It occurred to no one that any individual could fail to believe in a
supernatural punitive authority. The individual was confidently referred to presuppositions that everyone was assumed to share. If
the oath were a confession of faith, its continuance as an institution would be a gross violation of the principle of religious
freedom, as this is laid down in the constitutions of most modern states. If the institution of the oath continues after the
recognition of religious freedom, it must be legitimate to understand it as a summons to make one's declaration or one's promise
after the most serious and conscientious examination and self-reflection, each person fixing his thought on what stands for him as
most sacred. One then takes the oath in the same spirit in which it was originally instituted. But as long as the oath is used in
its old theological form, the door stands open to legal chicanery. It is asserted that the oath contains a confession of faith, and
one thereby obtains a means of excluding those who are honest enough to acknowledge their free religious conviction from the full
enjoyment of their civil rights\footnote[2]{An English example of such misuse of traditional oath formulas is given in F.
\emph{Pollock}: \textit{The Oath of Allegiance.} (Essays p. 187. 197). --- In the Danish Rigsdag it has been stated by a
high-ranking jurist that if the oath is to contain a heightened guarantee, it must be more than a solemn assurance, it must be an
invocation of a power that can ``intervene in human conditions in this or the coming life,'' that can in general ``do one
something.'' (Sitting of the Folketing, 20 December 1880). There seems to lie at the basis here an ethical standpoint of a kind
similar to that of the young girl in Poul Møller's ``En dansk Students Eventyr.'' She was so afraid of hell, but was reassured on
learning that it was nothing other than --- a bad conscience. What can that ``do one''?}.
% --- p. 487 ---
\opage{487}People do not see to what consequences this will lead. Bishop \textsc{Martensen}, who is greatly pleased that the state,
which otherwise ``in so many respects has given way to irreligious humanity,'' admits, by employing the institution of the oath,
that it cannot do without the religious guarantee, points out with full justice that this single guarantee, the oath, becomes
meaningless if one does not take better care of the religiosity of the people in general\footnote[1]{\textit{Individuel Etik.} p.
278.}. % translated from Høffding's Danish
So: the state must see to it that all citizens have the requisite dogmatic presuppositions; these must therefore be inculcated in
everyone with all possible force. Here one ends, consistently, with the Inquisition. --- What best agrees with the principle of
religious freedom is for the formula to be changed so that it contains no theological expressions. One thereby also avoids the
danger to which \textsc{Kant} already drew attention, namely that of giving people the belief that a lie is not so bad as long as
one does not swear to it.

Another example is offered by the form for the public contracting of marriage. Here the only arrangement fully in accord with
religious freedom is that there be a common, purely civil form for contracting it, whereas now in many countries it is only those
who do not wish to be married in church who are married civilly. The natural thing --- and at the same time the thing in accord with
the religious spirit --- is, after all, that those who wish to be married in church can add this to the civil ceremony, not that
those who do not wish it should have to get themselves exempted from it. By leaving the choice open between civil and church
marriage, one moreover\footnote[2]{Cf. E. \emph{Zeller}: \textit{Staat und Kirche.} Berlin 1873. p. 226.} establishes that they are
opposites that exclude one another, and thereby nourishes religious prejudices. --- In the discussion of this
% --- p. 488 ---
\opage{488}question it has been far too much forgotten that it is neither the ecclesiastical nor the legal sanction that ethically
founds marriage. Marriage is founded by a free covenant between man and woman, and the recognition of church or state is of
subordinate significance. (Cf. XVII, 4).

Similar examples could be drawn from burial customs, from the dominion of the church in the elementary school, etc. But the
examples given are sufficient to show the necessity of bringing about an arrangement of the relations of civil life that is
untouched by theological presuppositions.

3. The religious life that the church protects will become truer and more intense the more the church gives up ruling in the civil
world. The historical development of the last hundred years shows clearly enough how deep a root the church still has in the
peoples. At the same time as its outward privileges gradually fall away, its inner life has grown stronger. Where it is wholly or
partly separated from the state, it seems to have the greatest influence on minds. The exasperation that was aroused against the
church in the eighteenth century concerned its outward dominion far more than its dogmas and cult. The history of the Catholic
church after the Reformation, the position of the church in North America, the position of the German church after the civil
status law of 1875 show that the church itself can fare very well under the carrying out of religious freedom. On the ecclesiastical
side too great a fear has been entertained, and on the part of many freethinkers too great hopes have been built, of the separation
of the church from the state.

Strictly speaking it is rather the church's withdrawal from the state than a divorce between church and state that has been
demanded in the foregoing. All that is contained in the foregoing is that civil life in the state is to be ordered according to
general human presuppositions, not according to particular religious ones. The state must presuppose that its citizens, even if
they belong to no church, have the requisite ethical presuppositions for exercising the rights and duties that fall to them. But
this does not yet settle anything about what the future position of the church is to be. Because it is not to rule in the state
and in general civil relations, it does not follow that the state has nothing at all to do with
% --- p. 489 ---
\opage{489}it. As long as the church still lives in the people, it is a power of culture of great significance, toward which the state
cannot be entirely indifferent.

Even if a strong religious movement arose within the church, an enthusiasm like that of the first congregation, so that all
support and help from the state was renounced, --- even if the church thus became a purely private association, the state would
not on that account be able to lose sight of it entirely. The state would have to exercise a certain control over it to prevent
dangers from issuing from it to public peace and to the ethical foundation on which the state rests. It would have to prevent
fanatical encroachments on those of other faiths and the oppression of the church's own adherents. Here, as with respect to the
family, it would have the task of protecting the weak against encroachments, even within private and very intimate relations. The
state represents the general ethical principle over against the special religious-ethical principles that the various religious
communities profess. ``Morality stands as a more general law above the dogmas and the religious
tenets''\footnote[1]{This statement was made in the Danish constituent Rigsdag (at the sitting of 3 May 1849) by the then Minister
of Public Worship (Madvig) during the debate on § 76 of the present Constitution.}. % translated from Høffding's Danish
Religious associations or orders that disturb the peace between the different confessions and attack the very presuppositions of
the state the state will be compelled to keep a watchful eye on, although it will best maintain its dignity by leaving it to free
development to get beyond the abuses that may arise here. --- I disregard here the fact that the state must decide whether private
religious associations are to be recognized as legal persons, whether they are to have the right to own real property, whether
they are to have the right to teach, etc. ---

In countries where the church has taken root in the people and for centuries has permeated its life, it will be natural for the
state to continue to give it support as a community within which the greater part of the people finds the most important
nourishment for its spiritual need. The state thus regards the church as a historically given power of culture, and the support it
gives
% --- p. 490 ---
\opage{490}the church rests on the same motives as that which is given to science and art. The state cannot directly produce culture
of any kind whatever. Its activity in the service of culture is always indirect. It can neither produce nor eradicate religion; but
it can give material support and legal forms to the religious community. Whether, besides the church historically handed down in
the people, it will also support other religious communities will depend on the significance it attaches to them. --- As
conditions for the material and legal support the church receives from the state, the latter must demand the right to exercise a
more searching control than can be exercised over a purely private association. There are above all two demands that the state
must definitely make in the interest of spiritual culture and of freedom of the spirit.

It cannot be a matter of indifference to the state what constitution the church has. It must work toward a church constitution
that can serve to further the whole life and development of the people. A constitution that created a hierarchy exercising an
unlimited spiritual dominion over the members of the church the state will not be able to accept. It cannot favor spiritual tyranny
and the systematic narrowing of the people's spiritual horizon. It will as far as possible favor freedom of teaching and work
toward the people's being able to gather round the religious teachers who win its free adherence. It will not in the long run be
able to allow the church to demand that a definite system of doctrine be made the basis for all religious teachers, despite the
great improbability that they should all have appropriated it with full and independent conviction. Ecclesiastical unity and
ecclesiastical fellowship are bought too dearly when those who are to lead religious development are, involuntarily or
deliberately, induced to refrain from testing what they are to teach as impartially as love of truth
demands\footnote[1]{An illuminating example from the most recent time is offered by a ruling of the Württemberg ``Disciplinary Court
for Evangelical Clergy'' of 21 Feb.\ 1896, by which a clergyman (Pastor Steudel) was removed from office because on ``essential''
points he departed from the Augsburg Confession and the liturgy. In the judgment this man was reproached with the fact that ``the
one-sided elaboration and cultivation of his ethical scruples had its source in part in a certain overestimation of himself, by
which he was led to make \textit{a conflict that hundreds beside him fight out in silence in one way or another} the object of a
public martyrdom.'' % translated from Høffding's Danish
The high ecclesiastical authority thus here admits that hundreds of clergymen find themselves in a serious conflict with church
doctrine, a conflict in which not only intellectual but also ethical scruples play a part. It is very improbable that this should
be something that occurs only in the Württemberg church. --- Cf. on this case \emph{Christoph Schrempf}: Ein kirchliches
Ereigniss. (In the journal \textit{Die Wahrheit} edited by Schrempf. VI. Stuttgart 1896).}. --- The state will also demand that
those who are to be
% --- p. 491 ---
\opage{491}teachers in the church receive a certain scientific education. It is of no small significance that the young men who are to
enter the service of the church pursue their studies at a scientific university together with the rest of the country's student
youth. They come to know life from several sides. The very life of comradeship makes all too great one-sidedness impossible, and
they acquire a certain scientific conscience besides the religious conscience.

The influence that the state can exercise on the church in this way is supported by its work for culture in other domains. The
better popular education is, and the more luxuriantly scientific life flourishes in the people, the greater the demands that will
be made on the clergy, and the less will people put up with the dominion of a hierarchy. The state can support the individual
powers of culture, the intellectual, aesthetic and religious; but their mutual relation, whether it becomes harmonious or
disharmonious, it must leave to the free interplay of forces. It
can procure means and a material basis for the waging of a spiritual struggle, but cannot itself intervene decisively in this
struggle. To the people belong all, believers and unbelievers, and the life of the people unfolds under the constant clash between
the different spiritual tendencies.

4. How long such a positive relation between state and church will continue to be natural, the practical statesmen must decide.
The decision of this question will differ from country to country; it depends on the character of the people and on its
traditions. There are, however, several different circumstances that will more and more work in the direction of a withdrawal of
the church from
% --- p. 492 ---
\opage{492}the state as the natural course. The carrying out of religious freedom, and in connection with it the carrying out of purely
human forms of life, constantly work in the direction of isolating the church. Within the church itself the strict emphasis on the
doctrines of faith in their literal interpretation will produce a growing opposition to cultural life in other domains. And on the
other side, the church's endeavor to work as far as possible by voluntary forces and voluntary contributions will make it more and
more independent, so that it does not need to further the cause of the faith (or of the faithful) by means that the state levies
from unbelievers who for one reason or another have not withdrawn from the state church.

The question is whether a complete withdrawal of the church would be a good for the life of the people. Even if it becomes natural
and necessary, it is not thereby given that it signifies progress. If the church gets its own seminaries and its own clerical
constitution, fanaticism and narrow-mindedness will spread. The split in the people in religious respects will become greater, and
religion will come into the sharpest opposition to the rest of culture. The greatest dangers to peace and culture in the people
could spring from this. As experience shows, narrow, barbarous and life-hostile doctrines can gain entrance into the church despite
the moderating and humanizing influence that the connection with the state brings with it; but it will surely be worse if the
clergy are wholly replaced by lay preachers, and if the connection with the spiritual life of the people in other domains is
broken off. Only if the church were once and for all to acknowledge and adopt such tendencies would it become absolutely hazardous
for the state to maintain the connection; then the church would no longer be able to exercise any true power of educating the
people and would no longer be one of the representatives of ideal culture.

On the part of the church there appears in the most recent time an inclination, natural in itself, to secure the advantages that
religious freedom brings and yet keep the advantages that the position of state church brings. It demands full freedom to order its
own affairs, full self-government, and yet at the same time as much influence as possible on civil life and as much
% --- p. 493 ---
\opage{493}support as possible from the state. The church, as has rightly been remarked\footnote[1]{\emph{Zeller}: \textit{Staat und
Kirche.} p. 132.}, nowadays often comes forward in the name of religious freedom with demands that in earlier times it derived from
its divine mission. It is not only Catholicism that makes use of the name of freedom in this way\footnote[2]{\emph{Ketteler}:
\textit{Freiheit, Autorität und Kirche.} p. 159 ff.}. The Protestant state church too has raised the complaint that whereas the
sects have the right to order their affairs independently, it stands under the authority of the state and is thus the only
religious community in the land that does not have religious freedom!\footnote[3]{Cf. the article ``Religionsfrihed'' in the 2nd
edition of the ``Nordisk Konversationslexikon,'' and \emph{Martensen}: \textit{Social Etik.} p. 417 ff.}. In other words: on the
ecclesiastical side one would prefer to have it both ways, ``both in the bag and in the sack,'' and does not see that the more one
gets into the sack, the less one can get into the bag. An either--or is in general what more recent theologians seem to have the greatest difficulty in understanding.
Only if the church is willing to return to the position of a purely private association will it be able to obtain a virtually
complete self-government. When it accepts support from the state, the relation is altered.

Whether a complete divorce between church and state is a good will depend on the circumstances under which it takes place. It is
not, as many freethinkers in particular have believed, a magic remedy that does away with the difficulties of the religious
question. The exaggerated significance that has been attached to the divorce of church and state is connected with the exaggerated
significance that is attached to the state in general. The religious problems are not got out of the world because the state
withdraws from the church while the latter still has roots in the people. It is not the priesthood and other institutions that
produce and maintain religions; it is the need and adherence of the people that make the existence of the priesthood and the church
possible\footnote[4]{Frederick the Great protested against the assumption that the princes had allied themselves with the priests
in order to exploit the people jointly: on the contrary, it was the superstition of the people that gave the priests a power by
which the hands of governments too were tied! (E. \emph{Zeller}: \textit{Friedrich als Philosoph.} p. 32).}. In all probability the
bond between
% --- p. 494 ---
\opage{494}state and church will become looser and looser. Already now voluntary forces are working zealously and in great numbers in
the service of the church. Even the English state church, the stiffest of them all, now relies to a very large extent on
``voluntarism''\footnote[1]{\emph{Arthur Elliot}: \textit{The State and the Church.} London 1882. p. 127 f.}. And in the North
American free states there is a remarkable example of how the church (or the churches), after separation from the state, has on
the basis of voluntariness preserved and increased its spiritual influence and its material means. ``In the course of only three
generations American democratic society has carried out the complete separation of church and state, a reform that no other people
has ever attempted.'' ``A great part of the moral energy that three generations have been able to spare from labor for material
livelihood has been applied to carrying out the voluntary system in religion''\footnote[2]{The American \emph{Charles W. Eliot},
quoted in \emph{James Bryce}: \textit{The American Commonwealth.} London 1888. III. p. 347 f.}. % translated from Høffding's Danish
The North American community has from of old been permeated by the spirit of religious freedom; free, independent thinking is in
general so highly esteemed, and public opinion guards so strongly against all encroachments, that the humanizing influence of the
state is not missed\footnote[3]{Yet \textit{Moncure Conway} maintains that orthodox Christianity is, in spite of everything, still
a national American institution, because all church property is free from taxation, --- because Sunday rest is determined by law,
--- and because national chaplains are still appointed and paid. \textit{(Autobiography.} 1904. II. p. 291).}. Here, as on so many
other points, conditions in the European countries, which for such long times have lived under hierarchical and aristocratic
government, differ greatly from the American. If in Europe too the church is one day completely detached from the state, it is to
be hoped that religious freedom will to such a degree have permeated the whole way of thinking of the people, and that the people's
demands on its religious teachers will be so great, that the dangers that such a solution would now bring with it will not arise.

% --- p. 495 ---
\opage{495}\grouphead{\textit{3. PHILANTHROPIC CULTURE.}}

\chapterhead{XXXIV.}{THE NATURE AND SIGNIFICANCE OF PHILANTHROPY.}

1. \textsc{If} culture means the working-up and carrying further of something given by nature, there must also be a special
philanthropic culture. The need to relieve the suffering of others and to heighten their joy is just as much a human need as the
need to sustain life and to develop knowledge and imagination. Philanthropic culture, like material and ideal culture, aims at
developing and organizing an essential side of human nature.

The task of philanthropic culture is to work toward making the goods that material and ideal culture produce fall to the share of
as many as possible. Both in the domain of ideal and in that of material culture there can, as we have seen, arise a narrowing
tendency, so that the goods of culture fall only to the share of limited circles. In philanthropic endeavors an expansive, outgoing
tendency asserts itself, which lays the weight on the widest possible distribution of goods, not merely on their production.
Already in the foregoing, in connection with the great social problems in the domain of material and ideal culture (see especially
XXV, 3; XXVIII, 5; XXXII, 5), we have met this expansive tendency. It was indeed this tendency that first made the problems of
culture into \textit{ethical} problems.

The need that is satisfied through philanthropic culture is therefore also of a more immediate and direct ethical-social nature
than the needs that material and ideal culture satisfy. Sympathy itself can, after all, bind the individual together with others
even before he enters into connection with them through common work of culture. But it also comes to work indirectly as a founder
of community, and it is especially this indirect influence that is to be discussed here. Sympathy
% --- p. 496 ---
\opage{496}founds a community not only between giver and receiver, but also between all those who give in common, between all those whose
sympathy leads them to work in the same direction toward relieving the material and spiritual distress in the world. There arises a
special class of ends and tasks that can gather forces together and through common activity develop the community further.

2. Along two quite different lines, however, one might seek to ground the assertion that philanthropic endeavors do not rightly
form any special group within social ethics.

Seen from the one side, it might seem as if \textit{the whole} of ethics were really a doctrine of philanthropic culture. If
disinterested and universal sympathy is the psychological foundation presupposed in ethical valuation, then there seems to be no
part of the content of ethics that is not an expression of this sympathy and does not satisfy the need to relieve and help. Why then
set up a \textit{special} section on philanthropic culture? --- A line of thought such as this would come to a point in the
proposition: ``Nothing is needful but love of mankind!'' It sounds fine, but it is not true. For precisely where love of mankind is to
be able to work and accomplish something, there is need of many different forces and abilities if welfare in the world is to be
furthered. Welfare, after all, itself consists precisely in abilities and impulses being developed in as many as possible and in as
great fullness and harmony as possible. Therefore material and ideal culture, the family and the state, retain their independent
significance. One cannot root out all forces and impulses in human nature with the exception of love of mankind itself. One would
thereby obliterate personal differences and make human life poor. Love of mankind itself would in any case have to demand that they
be produced anew. It is to be a force that ennobles the work of culture, but it is not to take the place of this work or to make it
a mere means. Work, art and science have value first and foremost as expressions of deep-lying interests, and only \textit{because}
they thus have an independent value do they also become of ethical significance.
% --- p. 497 ---
\opage{497}There are philanthropists who work with great self-sacrifice and energy, but who disdain and reject all endeavors of culture
that do not directly and immediately ``improve'' human beings and their conditions. They regard science and art as expressions of
egoistic enjoyment and vanity. Of the Irish apostle of temperance, Father Matthew, it is told that he feared knowledge and
intellectual and moral freedom almost as much as brandy. % ``Matthew'' PRINTED AS IS (Father Theobald Mathew)
Many philanthropists also lack a sense of the significance of the
independent and active struggle for existence of the less favorably placed classes. They are glad to give help; but they are
offended that the class to which the needy belong organizes itself in order to be able to carry its demands through in society,
e.g., by strikes. A confident love of mankind, on the other hand, will believe precisely in the significance of the free unfolding of
life and rejoice in all the different forms in which human powers move, even where responsibility and danger are bound up with them.
---

On the other side, one might assert that it would be most beneficial if sympathy were not allowed to express itself immediately at
all. ``The welfare of the race and of individuals,'' it is then said, ``is increased by work in the service of culture. The more
material and ideal goods are produced, the greater the capital mankind has at its disposal, and that, after all, is what matters.
So no direct intervention, no immediate activity in a philanthropic direction, but energetic activity in the domain of material and
ideal culture!'' --- This view overlooks first and foremost that the production of a capital for the race is not enough; it also
matters how this capital is distributed among the members of the race. If there are those who go away empty-handed in the
distribution, then immediate help and intervention cannot be dispensed with. And it is not only the order of society that constantly
causes distress and misfortune through the unequal distribution of the goods produced. Distress and misfortune also arise through
the course of nature, independently of human relations and human wills. However high culture may be imagined to rise, it will
hardly ever be able to stop up the source of distress and misfortune that lies in
% --- p. 498 ---
\opage{498}the conditions of human beings as natural beings. Finally, not all have the conditions for taking part in the tasks of
culture. It is a favor and a good fortune to belong to the workers of culture, and it becomes a special task to work toward this
good fortune falling to the lot of as many as possible.

3. The two opposed views just criticized are closely connected with the two different sides that the struggle for existence
presents.

The gaze of the first is held captive by the sufferings that the struggle for existence brings with it. The beings whose nature and
strength do not correspond to the demands that circumstances make go under after first having passed through a longer or shorter
time of suffering. The weak are trampled down in the heat of battle, and the weary sink down before the fight is ended. When love of
mankind first awakens, it naturally turns with great fervor and with partiality toward these suffering, cowed and oppressed ones. It
does not confine itself to seeking them out and sacrificing itself for them; but it looks with suspicion on the happy, on those who
laugh, and at last finds the true and proper mood of life in sorrow and weeping. It is not content to comfort them for what they
must suffer, but even tells them that in suffering they have found the real truth of life. It thereby entangles itself in a strange
contradiction, since its constant aim is to remedy suffering, and it thus seeks to bring the human being out of the state that by its
own account should be the normal one.

The other view dwells chiefly on the fact that through the struggle for life new forms of life are developed. It is, once and for
all, the case that circumstances make their demands, and that the individuals and races that cannot satisfy these demands must sooner
or later perish. Even the most well-meant intervention only prolongs the sufferings of those doomed to destruction and prevents the
race from drawing the benefit from its experiences that it ought to draw. One relieves the distress of the moment at this particular
spot, but does not see that one thereby often produces misfortune and disproportion at other points and delays or halts development
for those who have the ability to advance. Just as a band of Indians on its hunting or war expedition must often abandon children,
% --- p. 499 ---
\opage{499}the weak and the old to their fate so as not to be hampered in its free movement, so the human race on its great forward
march must not be encumbered with those who cannot keep up. The more forces are spent on the ambulance, the fewer one has available
for the battle.

The two considerations can very well be combined with one another, and are in fact combined in the modern theory of evolution.
Every living being is required to be able to adapt itself to the conditions of life. This adaptation is brought about in two ways,
partly indirectly, partly directly. Indirect adaptation takes place through the weakening and perishing of the individuals and races
that cannot master the conditions. Direct adaptation takes place through the awakening and developing of new abilities that can
master the conditions. The higher the stage of development on which an individual stands, the greater the role that this direct
adaptation plays, and human culture is its highest form. All material and ideal culture is a form of the activity by which man makes
himself master of the conditions of life. But indirect adaptation constantly does its work, and it is precisely the task of
philanthropic culture to prevent it from assuming the brutal form that it has at lower stages of development. Nowhere perhaps does
what is properly human come out so characteristically as in care for and compassion with germs of life and forms of life that would
otherwise perish. To soften the consequences of the struggle for existence is a task that man sets himself as soon as he has got
beyond the very rawest stage; and it is a task in which he only continues what already expresses itself in the instincts that make
animals work and sacrifice themselves for their offspring. The suffering belong to the race just as much as the healthy and strong.
The gaze must not be turned one-sidedly to the one side or the other; we should thereby get a dualism in the race that conflicts with
universal love. Though it may be said that Buddhism and Christianity stared themselves blind at the suffering and misfortune in the
world and lacked an eye for the significance of the active work of culture, --- and though there was compassion and love of mankind
in
% --- p. 500 ---
\opage{500}the world before their time\footnote[1]{Orthodox theologians are not content to show that the virtues of the heathen are
splendid vices, but also deny to the heathen the virtues they really have. Thus \emph{Martensen} says: ``In heathendom mercy was
extinguished and unknown'' \textit{(Individuel Etik.} p. 304), and he sets out how ``the many hospitals and charitable foundations
suddenly appear as something new and unknown to heathendom, according as Christianity spreads in the Roman empire. For among all the
rich heathen and among all the wise heathen none ever thought of such things to help their suffering brothers, but only left them to
their fate.'' (\textit{Social Etik.} p. 160 f.). % translated from Høffding's Danish
This is historically incorrect. That the spirit of mercy was not unknown to the heathen, every reader of the Odyssey will already
know: the poor, strangers and suppliants stand there under the protection of Zeus, as is stated in several places with great
fervor. In Athens there was a special altar to Pity; the state supported those who on account of bodily weakness could not provide
for their own maintenance, and well-to-do private men considered it their duty to help the poor. Thucydides tells in his description
of the plague in Athens how many, without letting themselves be deterred by the danger of infection or by the horror of the disease,
went into the houses to nurse the sick. In Plato's ``Laws'' (p. 729---730) very great weight is laid on the duties toward strangers,
the friendless and suppliants. Such commandments passed over from the ancient Aryan ethics into the Greek (see \emph{Leist}:
\textit{Alt-Arisches Jus gentium).} --- Nor were hospitals and charitable foundations something altogether new and unknown in
heathendom. In particular there arose in the Roman imperial period, essentially under Stoic influence, a whole series of measures and
foundations for the benefit of slaves, children and the helpless. --- Cf. L. \emph{Schmidt}: \textit{Die Ethik der alten
Griechen.} II, p. 289 f. --- \emph{Schiern}: \textit{Om Humaniteten i den ældre romerske Kejserlovgivning.} (Historiske Studier.
II).}, --- it still stands as their immortal merit to have asserted \textit{the human rights of the suffering.} They have turned
men's eyes toward sides of life from which the eye has a natural inclination to keep away. They thereby first truly taught men to
know life. Despite the high-strung idealism in these world religions, they therefore nevertheless have a decidedly realistic
character. The love of mankind they proclaimed often took on, no doubt, a passive and sentimental character and was not always fitted
to ennoble the work of culture; but in its real nature it was nevertheless a new form of spiritual strength, the outcome of the
fullness of an inner life that was able to bear more than what the individual's personal fate laid upon it. (Cf.
% --- p. 501 ---
\opage{501}IX, 1). It does not fear that the race will become worse through the suffering not being left to their fate. It will continue
the struggle for life as long as possible; it does not easily capitulate before distress, and it is able to find and use
possibilities where the indifferent eye does not become aware of them. It does not fear that development will be delayed through the
help given to those who cannot keep up. The goal, after all, does not lie only in the future. The single generation is not merely a
means or a transitional link for future generations, but its distress and suffering weigh as an independent weight in the scale
when the value of life is to be judged. On this rest the justification and significance of philanthropic culture.

4. Love is often blind, and it then becomes a matter of chance what benefit it does. Insight into the actual circumstances, first
and foremost into the nature and position of the receiver, is a necessary condition for love not to bring about pernicious effects.
(Cf. XII, 2). Love must develop so as to become one with justice. (III, 9; X, 4). It rests on the conviction that in the one who is
to be helped there is a peculiar center, a force that is to be awakened to independent activity. It therefore fixes no distance
between giver and receiver; the gift is not handed down from above; the independence of the receiver is acknowledged. Love of mankind
thereby differs from ``grace and mercy,'' which are unethical concepts. In the concept of ``grace'' we have a remarkable example of
how wholly opposed standpoints can come to determine one another precisely through their opposition. When love is conceived as
grace, it is put in the greatest possible opposition to the obligingness that is a return for services rendered or springs from
feeling one's own power to be limited. Grace is shown by him who has complete power over others and stands so high above them that
he cannot receive anything at all from them. This concept is taken from the world of war; for thus the victorious warrior stands
over against the disarmed enemy. Grace is the victor's renunciation of using his power. But ethically regarded there can never occur
a situation in which the independent right of the receiver should have wholly
% --- p. 502 ---
\opage{502}vanished, and in which the giver stands as absolute starting point. They both stand as members of the race, subject to a common
law that gives each of them his right\footnote[1]{Only when the concept of grace is interpreted as a higher justice (see
\textit{Etiske Undersøgelser} p. 63) do these misgivings fall away.}.

That the giver is never an absolute starting point is easy to show. It may even be said to him: ``What hast thou that thou didst not
receive?'' % 1 Cor. 4:7, KJV
The favorable conditions under which he began the struggle for life and has got as far as he has, he did not himself produce. He
stands here in a debt to the race which he can pay off only by the help he gives others in their struggle for life. Besides, he has
perhaps, consciously or unconsciously, been the cause of the conditions for others not becoming as favorable as they would otherwise
have been. Or the favorable conditions that have benefited him are connected with an order of society that does not fully satisfy the
demands of justice. \textsc{Kant} raises the question whether the ability to do good is not for the most part a result of the
injustice of the state, which brings about an inequality in prosperity that in turn makes beneficence necessary, --- and whether under
such circumstances the assistance the rich man gives the needy deserves the name of beneficence at
all\footnote[2]{\textit{Tugendlehre} § 31.}. --- History shows that poverty increases especially when traditional institutions
dissolve without being replaced by others that can afford a support like that which the former gave. Thus poverty developed in the
Middle Ages after the cessation of personal slavery, later with the cessation of serfdom and of the guild system, in general, then,
with the transition from bound labor to free labor\footnote[3]{Cf. \emph{Lecky}: \textit{History of European Morals.} II, p.
78---87. --- \emph{Gneist}: \textit{Das Selfgovernment in England.} 3. Aufl. p. 260.}. Slavery, serfdom and guild compulsion did not
accord with the demands of justice; but just as little was there full justice in their falling away without new forms being formed
that could afford the help they had given. (Cf. XXIII, 4). The individual here inherits the guilt the race has brought upon
% --- p. 503 ---
\opage{503}itself, inherits it by continuing it and enjoying the benefit of its consequences. Seen from this side, philanthropic culture
stands as an endeavor to restore what has been neglected and incurred. \textsc{Friedrich Nietzsche} sought to derive the great misery
and distress found in the European world from the ``slave morality'' that has kept the weak and the sick alive at the expense of the
strong and the healthy. In opposition to this assertion it has surely been maintained with full justice that if an increase in the
number of sick and unhappy individuals has taken place, it is to a great extent to be derived precisely from the exploitation of the
weak by the strong which Nietzsche advocates in his ``master morality.'' The ``masters'' whom Nietzsche admires sacrifice with a
``good'' conscience the health and honor of others for their own enjoyment or advantage\footnote[1]{F. \emph{Tönnies}: \textit{Der
Nietzsche-Kultus.} Leipzig 1897. p. 109.}.

To this it may further be added that the individual right of property --- as shown earlier (XXVI, 15) --- can be grounded ethically
only when property is conceived as a means entrusted to the individual so that by it he may solve a social task. So-called
beneficence thus does not stand as something wholly new, as a quite special virtue, but as a natural consequence, a duty.

5. The great significance that Buddhism and Christianity have had and have in the history of philanthropy is limited not only by
the blindness of the compassion and sympathy they awaken\footnote[2]{It has even been possible to assert that the church by its
irrational almsgiving produced a great part of the evils it sought to remedy. F. \emph{Raumer}: \textit{Geschichte der
Hohenstaufen.} VI, p. 578. \emph{Lecky}: \textit{History of European Morals.} II, p. 101.}, but also by the fact that to a very
great extent it was religious motives that lay at the basis precisely of the most magnificent works of love. Where love of mankind
receives a theological motivation, there is an intermediate link between giver and receiver, and these do not stand in an immediate
relation to one another. The work of love does not spring immediately from fellow feeling, but takes on the character of a work of
penance. What one gave to the poor, one gave not so much for their sake as in order to win one's own salvation. Alms --- it was said
in the Middle Ages
% --- p. 504 ---
\opage{504}--- bear fruit a hundredfold and blot out sins, just as water quenches fire\footnote[1]{Raumer op. cit. p. 575. ---
\emph{Lecky} op. cit. p. 85. 99.}. The reason why the great zeal for charity cooled for a time at the end of the Middle Ages was
partly that belief in the power of good works to procure remission of sins fell away with the Reformation. The poor were merely means
or objects on which penitence was to exercise itself: what was essential was that one deprived oneself of something, not that the
others received something. From this standpoint, therefore, endeavors aiming to make almsgiving superfluous are often viewed with
suspicious eyes. The Irish politician O'Connell, a zealous Catholic, for this reason had misgivings about modern poor-law
legislation: begging he regarded as something holy in itself, which it was impious to abolish. ``The medieval conception of the
meritoriousness of alms necessarily entailed that beggars were looked upon with a particularly mild eye. They were by no means
regarded as the outcasts of the world; the beggar's life was on the contrary regarded as a life especially well-pleasing to God. When
all was said and done, it might perhaps be open to doubt whether it was the givers who were the real benefactors, and not rather the
beggars. These, after all, gave people occasion to give alms and thereby furthered the eternal well-being of the givers; for temporal
gifts one received eternal ones in return. And one did not look with critical eyes at the beggars to see whether they were worthy of
help or not; one poured out indiscriminately to worthy and unworthy''\footnote[2]{M. H. \emph{Nielsen}: \textit{Fattigplejen i
Danmark før Reformationen.} (Aarbog for dansk Kulturhistorie 1895). p. 85.}. % translated from Høffding's Danish

The theological motivation has, moreover, the drawback that it naturally comes to make worthiness to receive help dependent on
dogmatic faith. Philanthropy becomes confessionally limited. A limitation is thereby set that is foreign to the nature of love of
mankind and prevents it from unfolding freely. Besides, hypocrisy is furthered when a particular confession, or perhaps even
participation in certain ecclesiastical customs, becomes a condition, or at any rate a favoring circumstance, for enjoying the benefit
of the good deed.

% --- p. 505 ---
\opage{505}Therefore the emancipation of philanthropy from confessional barriers is of extraordinarily great ethical significance. Yet most
people still seem to need theological motives if beneficence is to have any real vigor. It is not only medieval Catholics who know
how to trade with heaven and who give nothing away without a satisfactory draft on compensation in another world. Among Protestants
too there are traces enough of such ``other-worldly worldliness'' (other-worldliness). As for those, on the other hand, whose faith
breaks through the confessional barriers, the capacity for fellow feeling and self-sacrifice does not always grow along with the
critical sense. Knowledge often thrives at the expense of feeling, and it may look as if theoretical narrowness were the condition of
practical expansion. --- Love of mankind, which has overcome so many difficulties, will surely overcome this one too. It was born in a
manger, but it will grow, and it will try to permeate the whole world.

\chapterhead{XXXV.}{THE ORGANIZATION OF PHILANTHROPY.}

1. Already the individual must organize his beneficence. He must manage the means he can and will use to help others in such a way
that they fulfill their purpose in the best manner. He must not let himself be guided by mere whim, so that it becomes a matter of
chance that he gives to the one and refuses the other. Whim will naturally prevail especially the less personal the relation in which
one places oneself to those who are helped, as with beggars. When humane men like James Mill and Archbishop Whately prided themselves
on never having given a beggar anything, it was a right principle, since it was supplemented by active help to those whose need one
knew, and toward whom one could therefore also find the best way
% --- p. 506 ---
\opage{506}of helping. It is an easy matter to earn a daily wage of 3 to 4 kroner by going from door to door in Copenhagen; this is learned
from the statements of the many beggars brought before the criminal court\footnote[1]{\textit{Københavns Understøttelsesforening
1874---1899.} By \emph{Jul. Rüdinger}. p. 4.}. Begging often pays better than real work. But the beneficence that furthers such
things has gone astray and needs organization. It has even been said that he who gives blindly, ignorantly and thoughtlessly acts like
one who fires a cannon into a crowd of people\footnote[2]{\emph{Henderson}: \textit{Introduction to the study of the dependent,
defective and delinquent classes.} Boston 1901. p. 139.}. --- The individual's means, experience and insight are limited. His mite
can often have the best effect by being added to that of others, and a certain division of labor naturally asserts itself, since some
individuals are especially suited to work as investigators, directors or distributors in the service of philanthropic associations.
A free bureaucracy can develop here, a kind of honorary posts, which may one day carry no less esteem than the offices of the state.
Yet it will not be for the good if the personal relation between giver and receiver fell away entirely. It is the living kernel of
philanthropic culture, the element of freedom that this culture can no more do without than any other kind of culture. Individual
beneficence and comradely beneficence will always be necessary presuppositions for, and transitional links to, the systematic
activity of larger associations. Individual and comradely philanthropy will strive precisely that the systematic machinery be set in
motion only where it is absolutely necessary. The old neighborly relation, which so easily loses its significance in the course of
culture, especially in the large cities, must always stand as the ideal one seeks to hold fast or restore under the changed
conditions. \textsc{Bosanquet} has therefore rightly said that charity ought to mean a service that one neighbor renders another, not
blind almsgiving. The philanthropic associations themselves ought to work with as great decentralization as can be united with the
overview and
% --- p. 507 ---
\opage{507}the concentration that it is precisely the purpose of these associations to bring about. Otherwise the contribution to the
philanthropic association easily becomes a kind of tax that one half unwillingly takes upon oneself because it is, after all, the
custom; and those who receive the help also easily come to feel that they stand before a machine. The art will be to combine
individualizing activity with sufficient regularity and uniformity in the activity as a whole. Without organization of philanthropy
the work is done blindly and in a haphazard way. The poor are accustomed to lie in order to obtain help from several places at once,
and the help that can come from a single place is also frequently so small that it accomplishes nothing but enabling the receiver to
obtain work by underbidding his comrades. Some years ago it was found in London that there was no lack of funds for the relief of
distress, nor as a rule of persons and institutions willing to take part in the work of relief, but that for lack of organization,
not least for lack of sufficient knowledge of the circumstances, the whole often did more harm than good\footnote[1]{B.
\emph{Bosanquet}: \textit{The principles and chief dangers of the administration of charity.} (Journal of Ethics. III). --- On
philanthropic organizations see, besides the report on Københavns Understøttelsesforening cited above, the description of the great
English and American organizations in \emph{Henderson}: \textit{Introduction.} Part. II, chap. 9: The charity organization society.}.

Free philanthropic association is still only in its beginning. It surely has a great future before it; in any case it will not come
to lack tasks.

2. Neither individual nor organized philanthropy, however, will be able to master all tasks or provide a guarantee against the worst
distress. For that a firmer organization and greater means are required, which only the state commands. In poor relief, as in
schooling, it is the task of the state to see to it that the most elementary conditions are present. This is not only because the
state thereby indirectly provides for security. Just as little as it supports ideal culture for the sake of security, just as little
is it from mere prudent calculation that it organizes
% --- p. 508 ---
\opage{508}help for the poor. It steps in where the most extreme distress is abroad, even if public security is not threatened. And just
as little as the help of private individuals is the help the state gives to be a matter of grace; it is to be regarded as an
ethical-social act of justice. By setting up a system of poor relief the state admits that it has something to make good. In an ideal
state there would be no poor; that there are those who lack what is necessary for the maintenance of life, although they have the
ability and the will to work, will always be connected with faults in the social organization. It does not follow, however, as
\textsc{Fichte} thought, that support of the poor is first and foremost the business of the state and only in the second place the
business of private individuals. The intervention of the state must here, as at other points, be only the last resort, the necessary
supplement to the activity of free forces. Free philanthropy has precisely the task of preventing as long as possible the state from
having to intervene with its heavy apparatus.

From two quite different sides the philanthropic activity of the state has been attacked. --- In state poor relief people have found
a bad communism, in which the state gives its gift without having sufficient influence on the receiver. When those who receive poor
relief cannot be forbidden to marry, poverty is of course increased by state help. Instead of poor relief, therefore, a system of state
insurance has been demanded, in which the workers were to be compelled to take part\footnote[1]{So the German economists \emph{Wagner}
and \emph{Schäffle}. (Schönbergs Handbuch der polit. Oekonomie. 1. Ausg. II, p. 598. 615). --- Cf. above XXVI, 14.}. But it is not
easy to understand how people can be got to insure their future when they have nothing to live on in the present. This proposal thus
presupposes that the worst distress has already been relieved in one way or another. But how this result is to be reached is
precisely the difficulty. --- From quite another side it is desired to keep the state altogether away from this domain. As great and
fine a thing as private beneficence is, so pernicious a thing is state beneficence held to be. Only free love of mankind ---
% --- p. 509 ---
\opage{509}it is thought --- allows beneficence to be exercised with the right sense of responsibility, so that by supporting the unworthy
one does not lead the race to regression instead of progress. A bureaucracy cannot have the right sense of responsibility; it does
not come into personal relation with the suffering. Besides, it is dangerous to increase the power of the state; one thereby easily
creates dangers for individual freedom. A beneficent despot nevertheless remains a despot. The state has only the task of providing
for right and security; in the state, unlike in the family, sympathy cannot rule immediately\footnote[1]{Cf. \emph{Herbert Spencer}:
\textit{The man versus the state} (London 1884), where this line of thought is carried through in consistent form. In the second
volume of \textit{Principles of Ethics} (1893) (Chap. 7: Relief of the poor), where the difficulties of and misgivings about all
philanthropy are sharply emphasized, Spencer nevertheless thinks that state support of the poor has a just basis, since the right
that the serfs of former times had to use the land was taken from them in the course of time by the landowners: the poor law is in a
way an acknowledgment of the primeval right to the land. But state help is and remains for Spencer ``a kind of social
opium-eating,'' which soothes for the moment but makes things worse later. A period of weaning is needed in order to pass over to
``the healthy state of self-help and personal beneficence.'' % translated from Høffding's Danish
}. To this it must be answered that hardly such sharp boundaries can be drawn between beneficence and justice as the line of thought
cited presupposes. True beneficence is itself an act of justice, and it is therefore right that the Danish Constitution (§ 84)
acknowledges that the needy person has a right to public support. Care for security is indeed the most elementary task of the state;
but that does not exclude the state from stepping in to help with its organization and its means where important humane ends cannot
be reached merely by the activity of free forces. If we could provide for our security on our own, we would prefer to do without the
power of the state; the reason why we call the state to our help is in all domains essentially the same: the insufficiency of free
forces, an insufficiency due not least to imperfect organization. The danger, to be sure, is also everywhere the same: namely that
the system and the power that are set on foot should spread too much at the expense of the real ends. This
% --- p. 510 ---
\opage{510}danger exists, however, also, though naturally in lesser degree, in organized free philanthropy; it arises wherever the
immediate relation between giver and receiver gives way to a mediated one.

3. Where the state is at work, compulsion always stands in the background, and a systematic procedure is employed that does not
always allow the individual, actual circumstances to come into their own. On a large scale this was shown in England by the effect of
the older English poor-law legislation dating from the time of Queen Elizabeth. Partly through the easygoingness of the
administration, partly through too strong centralization and the tendency to treat all cases according to an abstract scheme, partly
through lack of control on the part of the taxpayers, the most pernicious conditions arose. Since relief was given without closer
investigation, the poor in the population were accustomed to improvidence and indifference. They did not think of the morrow; they and
their children would, after all, in case of need be maintained by the public, perhaps even be better off than many of those who worked
hard to be able to pay their share of the poor rate. Not only indifference and idleness, but also ingratitude and egoism were bred.
Children refused to help their old parents and left it to the public. The demoralization became so great that Lord Chancellor
Brougham could even declare in the English House of Lords that the English poor laws were the chief cause of the moral corruption of
the population and the increase of crime. It then became common (from the end of the eighteenth century) for the parishes to grant a
supplement to wages where these were insufficient. Thereby poor relief was really transformed into a subsidy to the employers, who
could now press wages down quite below what was necessary for the maintenance of the workers!\footnote[1]{\textit{Das Armenwesen. Frei
nach Duchatel und Naville.} Von einem deutschen Staatsbeamten. Weimar 1837. p. 502 ff. --- \emph{Gneist}: \textit{Das Selfgovernment
in England.} 3. Aufl. p. 700 f.} The more recent legislation (from 1834 on) has sought to remedy this. But it has constantly proved
difficult to decide where
% --- p. 511 ---
\opage{511}need is really present and where not. In England the so-called ``workhouse test'' was introduced, it being thought that
able-bodied people who would not put up with the work the public imposed on them as a condition of relief thereby showed that they
were not really in need. But so mechanical a means only became necessary because sufficient decentralization was not
employed\footnote[1]{Cf. \emph{Gneist} op. cit. p. 740: ``When in the assessment of distress and the determination of the right means
of relieving it one lacks the understanding that a neighbor can have, and that presupposes personal and local acquaintance such as
can be found only in living communal bonds, one arrives at the schematic workhouse test at the expense of humanity and at the moral
degradation of the working classes. The factual question of the cause of distress can be settled only in a living communal bond by
the members of the commune; if anything of the kind is lacking, the providing power of the state arrives at the mechanical workhouse
test for the same reason that the punishing power of the state once arrived at the rack.'' % translated from Høffding's Danish
Cf. with this Bosanquet's sentence mentioned above (1): Charity for us means neighbourly service.}. In smaller circles (communes and
still narrower neighborly relations) there is the possibility of personal acquaintance with all the single individuals, and the state
ought therefore as far as possible to leave the distribution of relief to such circles. The more recent arrangements of this matter
also go in this direction. Here, even more than in other domains, self-government is the only arrangement that can lead to the goal.
Where the smaller circles have self-government, each single case can be treated according to its individual character. In the
``Elberfeld system''\footnote[2]{\emph{Löning}: \textit{Armenflege und Armenpolizei.} (Schönbergs Handbuch. 1. Ausg. II), p. 588.
601 f. % ``Armenflege'' PRINTED AS IS
--- According to \emph{Henderson} (\textit{Introduction.} p. 164) the Elberfeld system arose in part on the model of the poor relief
founded by Chalmers in Scotland. ``The friendly visitor'' is in Chalmers' system the good neighbor who brings help and counsel in
distress and sets life going where it threatens to come to a standstill.} decentralization and individualization have been carried
so far that each helper has to do with only four families. Their circumstances he can then learn to know thoroughly, and it becomes
possible at the same time for every
% --- p. 512 ---
\opage{512}citizen to take on the honorary communal office of helper, since his time is not so heavily taken up by it. By such an
arrangement state philanthropy comes nearer and nearer to the private philanthropic organizations. It works best where it cooperates
with them. Through a series of transitional forms the great distance between the personal relation and impersonal state intervention
is shortened.

% --- p. 513 ---
\opage{513}\grouphead{C. THE STATE.}

\chapterhead{XXXVI.}{PEOPLE AND STATE.}

1. \textsc{If} one wishes to set the state in sharp opposition to the family and the free cultural community, it is the coercive power
that most clearly characterizes it. But it is also only where that opposition is to be presented in its extreme sharpness that this
factor stands out so strongly. Power always stands in the background; but where it is the only bond holding things together, the state
is near its dissolution. The nature of the state first comes before us in its whole scope when we see that it is able to awaken
feelings and motives that can even overcome those awakened by family relations and the ends of culture: as when the individual
sacrifices not only his own life but also regard for the welfare of his family in the service of the state, and when the artist and the
man of science subordinate the interests of their field to the welfare of the fatherland. Such feelings the state can awaken only
because its existence and its activity are necessary conditions for the attainment of great universally human ends. This has already
appeared in the foregoing, inasmuch as the investigation both of the family and of the free cultural community led to pointing out
tasks that only the state could solve. To that extent we have already in the foregoing presupposed a definite concept of the state,
which it is now a matter of developing more closely in such a way that both the
% --- p. 514 ---
\opage{514}opposition between the state and the other forms of community and the close connection between them can find their grounding.
An investigation of the relation between people and state will serve this purpose.

2. In discussing the family and the free cultural community as independent social forms of life, we have disregarded the fact that in
reality they always occur in a definite historical and national context. At the most primitive stages of development there is no sharp
boundary to be drawn between family, cultural community and state. In the horde or the tribe, as it appears at early stages of
development, the bond of blood prevails, even if it is not the only thing that binds. The primitive horde is at the same time a
cultural community, insofar as there can be any talk of culture at all at this stage. Work is done in common, and what is possessed
is owned by the horde in common. Ideal culture is represented by belief in the gods of the tribe. And the horde is finally a little
state, inasmuch as it gathers itself for common resistance against outward enemies, and inasmuch as the father of the family or the
chieftain also exercises power inwardly in order to keep the refractory in check. --- It goes here as everywhere: the sharp and
distinct differences only emerge in the course of development and do not make themselves felt from the start.

A horde is not a people, but first becomes one when there develops a more conscious feeling of community than can be present where
people split up into small groups just as quickly as they gather for mutual assistance in the hour of need. A coherent memory and
tradition is a presupposition for the coming into being of a people. Through common fate and common activity every feeling of
community arises (cf. XIII, 3 and elsewhere), national feeling too. Common fate and common activity are conditioned first and
foremost by all dwelling and building in the same land. Its character conditions the common nature of the work as well as common
physical vicissitudes, common moods toward nature, common fear and common joy\footnote[1]{On the conditions for the formation of a
people see \emph{Spencer}: \textit{Principles of Sociology.} II. p. 265---287 (Political integration).}. Even a
% --- p. 515 ---
\opage{515}nomadic people has, after all, its places to which it returns from time to time, so that there is coherence in its memories.
If it has no other fixed possession, it still has the graves of its fathers, to which it feels itself bound and around which it
circles as around a center\footnote[1]{Cf. \emph{Herodotus} IV, 127, where the Scythians refuse to engage in battle with the
Persians, since they have no cities or cultivated land to defend; they will fight only if the Persians find the graves of their
fathers and destroy them. --- In Genesis, chap. 23, Abraham, who otherwise lives as a nomad, buys a piece of land in order to bury
Sarah in it.}. It is indifferent what originally brought the horde together, --- whether it arose by descent from the same family, or
by the joining together of several wandering individuals, or by a stronger horde subjugating a weaker. The main thing is that such a
life in common gets going that a common life of thought and feeling develops under common fate and common activity. The common
language is itself often first a product of this life in common, and it is only at higher stages of development that it can come to
play a decisive role. Language is not always the decisive thing. The Swiss people is an example of how the boundaries of national
feeling and of language need not coincide; the main thing is that a common consciousness develops, and for that only a common history
is needed.

Through common fate and common activity the habits and customs common to a human group develop. Life takes shape involuntarily in a
certain definite way and in certain definite directions. These involuntarily formed forms then in turn work determiningly for the
future, and that not only in an unconscious way but also through more or less conscious need and feeling. A need is felt to perform
the same action in the same way under the same circumstances. Tilling the soil, hunting and war, clothing and weapons, the contracting
of marriage, upbringing, worship of the gods, revenge and enjoyments --- all are ordered and practiced not so much according to what
on reflection is found expedient and reasonable as according to the way in which the fathers ordered and practiced them. A breach of
this or a
% --- p. 516 ---
\opage{516}hindrance to it can awaken a passion like that which an animal feels when it is prevented from performing an instinctive
action. Indignation is felt at him who would introduce a new custom, and an astonishment mixed with indignation is the feeling with
which the morals and customs of other human groups are observed. --- It is --- so we may also put it --- the common positive morality
(see I, 2) that, in connection with the common memories, makes a human group into a people.

3. National feeling first comes to full and clear consciousness through the clash with, and indignation at, what is foreign. It is a
psychological law that a more or less strong contrast is required for a state of consciousness to acquire its definite, pronounced
character. History accordingly shows that national feeling has been able to work unifyingly only by working dividingly: it has united
through the consciousness of common peculiarity that was aroused especially by a common feeling of indignation, contempt or enmity
toward what is foreign.

From the start it is more than an aesthetic antipathy that is at work. Everything foreign and unknown can at primitive stages awaken
fear or hatred, because it can bring danger to life, possessions and freedom. One puts oneself in any case, for safety's sake, in a
state of defense. But even if fear falls away, antipathy maintains itself. From the standpoint of national self-esteem a sharp
judgment is passed on what is foreign. It is two ethical systems that here stand against one another in flesh and blood, and a
mediation or reconciliation can take place only through a historical development that brings with it community in fate and work. A
longer time of peaceful life together is needed to discover that what is foreign has its peculiar value. It is a widening of sympathy
and a widening of ideas that must take place.

National feeling can have its full ardor only where it appears with the strength and blindness of instinct. It has been a mighty
historical force. It has made history into one great history of war, but precisely by occasioning war it has worked to hold together
and concentrate. War brings with it the most far-reaching common experiences for the whole group. Common danger and fear, common
suspense and
% --- p. 517 ---
\opage{517}hope, common defeat and common victories, all this binds together with the strongest cement. And to this is added that at the
stage of instinctive national feeling everything is staked on one card: in the struggle all the goods of life are at stake at once;
if one is defeated, one must be prepared to lose freedom, property, perhaps life itself; indeed one even loses the gods, since the gods
of the victors displace them. It is a question of to be or not to be. In \textsc{Aeschylus}' tragedy \textit{The Persians} the cry goes
up on the Greek side as the battle of Salamis begins: ``Sons of the Hellenes, free your fatherland, your children and wives, the
dwellings of your fathers' gods and the graves of your forefathers! Now the struggle is for all!'' % translated from Høffding's Danish
--- It is otherwise when the time of wars of extermination is past. War itself is waged with greater humanity, and it is no longer all
the goods of life that are at stake. There is a universally human social life that reaches out beyond national differences, and the
opposition between our own nationality and the foreign no longer coincides with the opposition between good and evil. Will not
national feeling, then, one day have played out its role as a mighty historical driving force, and will not the ethical significance
it has had be at an end? It seems bound to lack sufficient nourishment and therefore also to be unable to retain sufficient energy.

4. National feeling has in common with every feeling that approaches instinct that it is not called forth by the value that can in and
for itself be attributed to that to which it is attached. We do not love our land and our people because we consider them the best,
most beautiful, richest and best endowed; but we love them first and foremost because it is \textit{our} land and \textit{our} people.
It is a purely subjective, not an objective valuation that is made in national feeling. In a naive way this finds expression in the
superlatives with which we adorn what we love; these superlatives are expressions of the fervor and strength of the feeling, but not
of the result of an objective investigation. We may even quite readily admit that other lands and peoples possess a greater objective
excellence; yet they always have the defect that they are not \textit{our} land and \textit{our} people: ``fairer they may be, perhaps,
alas, but they are not it.'' % translated from Høffding's Danish (a line of Danish verse)
It is here as with self-preservation: we preserve life first and foremost because it is
% --- p. 518 ---
\opage{518}our life, not because we have examined it and found it to be an excellent life. On this half unconscious, instinctive character
of national feeling it depends that it can be preserved even if the horizon is widened and antipathy toward what is foreign falls
away. With the sharp contrast to what is foreign, the fervent feeling for what is native does not necessarily fall away. And only in
what is native do we have a firm hold and a foundation from which we can orient ourselves in the world\footnote[1]{In my essay
\textit{Nationalitet og Kultur} (Tilskueren 1895) I have described national feeling as a feeling of recognition and as a feeling of
community, and have also emphasized that in every work that goes at all deep the nationally determined elements in the personality
must take part.}.

Advancing culture does indeed increase unity and community in the human race, but it also opens the eye to the peculiar nuances and
differences. Each people takes life in its own way, and thereby life becomes richer and more manifold for him who looks at it with
sympathetic understanding. The barbarian tramples down what is foreign because he does not understand it, and if there is danger for
national peculiarities, it is rather because there is too much barbarism in the world than because there is too much culture. One
nation has not yet properly seen the necessity that the other should exist. Advancing culture both opens the eye to the peculiarity of
the different peoples and gives the single people greater opportunity to make its contribution to common development.

Just as the family and the free cultural community are natural intermediate links between the individual and the state, so the
nationalities are natural intermediate links between single individuals and the whole race. What is universally human is not thereby
hindered in its development; on the contrary, it thereby becomes richer and more manifold.

5. The state is the organized people. The state presupposes such an organization that the people can act as a unity outwardly, and
that the conduct and ordering of its inner affairs also present a certain unity. One cannot point to any definite point at which the
state comes into being. Already the family and the horde can have a certain organization, so that they both resist outward
% --- p. 519 ---
\opage{519}attacks and order the relation among their members. A large family is already a little state.

Through the firmer organization that the state presupposes, the people gains the same advantage that the animal organism gains by
having a nervous system. In the plant a luxuriant organic life is led; metabolism, growth and reproduction go on with great energy.
But only where a nervous system connects the parts of the organism can it act actively and with concentrated force toward the
surrounding world, direct its activity toward definite points and be determined by regard for what lies far beyond its immediate
surroundings. Without the organization of a state the people leads a plant life. A certain culture can develop; material activity and
the life of thought and feeling can unfold to a certain degree; but the whole bears the stamp of chance and dispersion. Further
development is conditioned by concentration and firm coherence of everything that belongs to the people. In the struggle for existence
such a concentration becomes of the greatest significance. A kingdom divided against itself cannot stand. % cf. Mark 3:24, KJV
War has therefore at all times been a cause of firmer organization. Yet it is only at the most elementary stages that concentration
in the face of outward danger is the ruling consideration. The organization that is first brought about to meet outward resistance is
naturally employed to help and protect all the interests and ends that the life of the people comprises. There are two types of
political life, one that is predominantly determined by the necessity of displaying power outwardly, and another that is essentially
determined by regard for the inner life of the people itself\footnote[1]{\emph{Spencer}: \textit{Principles of Ethics.} II. p.
181---187.}. As long as a state must fight for existence against other states, the first type will predominate; as soon as more
peaceful conditions set in, the second type will develop. It will be easy to see which of the two types will be the more fortunate for
a sound development of material and ideal culture.

Already at the primitive stages of development there is something that is stronger than mere power and takes it into its service. It
is the unwritten laws contained in the customs of the people.
% --- p. 520 ---
\opage{520}(Cf. 2). The ordering of the relations of the people can never take place merely through commands of power. The holders of power
themselves are inwardly more or less possessed by fear or awe of what has been handed down. No holder of power stands so far above his
time that he does not have his models and precedents, which he follows. In such models and precedents we have the primitive form of
\textit{law.}

We thus have different factors in the concept of the state. Power is the last means and foundation of the state; law is the norm
that determines the use of power; and culture is the end that the use of power guided by law is to serve\footnote[1]{Syndicalism in
its criticism of the state takes account only of the factor of power, not of law and culture. And in power it sees only the means of
a ruling minority. Even Bakunin's order built on the communes is rejected, because the commune is a political institution. It does not,
however, go into the question of how the state will collapse and how the new order can be organized. (\emph{Poul Louis}: Le
Syndicalisme contre l'État). (1913). % ``Poul'' PRINTED AS IS (the author is Paul Louis)
}. And with all these factors the question of the ethical significance of the state still remains to be raised.

\chapterhead{XXXVII.}{LAW AND MORALITY.}

1. Where law consists only in the customs and traditions that involuntarily govern life, the opposition between morality and law
does not make itself felt. By law we understand, from our modern standpoint, the sum of the rules for the use of power expressed in
definite promulgations, whereas by morality we understand what is made known to us by conscience within us. But such a sharp
opposition between an outer and an inner, between outer power and inner feeling, is only the fruit of a long course of development.
Originally there is no difference between
% --- p. 521 ---
\opage{521}custom, tradition, law, morality and religion. A need asserts itself to treat like cases in like ways. That something has
once happened in a certain way is reason enough for it to happen the next time in the same way. \textsc{Tocqueville} even thought he
had observed that among the Parisians a certain usage and custom had developed according to which they conduct themselves during and
after the revolutions they make.\footnote[1]{\textit{Souvenirs} (Paris 1893) p. 105: Il s'est formée parmi nous une espèce de moralité
particulière ou désordre, et un code spécial pour les jours d'émeute. D'après ces lois exceptionnelles, le meurtre est toléré, la
devastation est permise, mais le vol est sévèrement défendu} % French as printed
So even here repetition and habit have been at work! --- As already remarked, it is not a purely unconscious habit that prevails
here. In the need to follow the paths once laid down there is also at work reverence for the past and for the obscure powers from
which the prevailing customs and traditions\footnote[2]{We lack in Danish a word that fully expresses what lies in the German
``Sitte.'' ``Sitte'' differs from habit in that it appears as normative for consciousness, whereas with habit no distinct
consciousness need assert itself. There is in ``Sitte'' an inner, subjective factor that may be lacking in habit, which approaches
pure reflex movement. --- Cf. the excellent development of the concept of ``Sitte'' in \emph{Jhering}: \textit{Der Zweck im Recht.}
II, p. 19---41. It is defined there as ``verpflichtende Gewohnheit.''} are derived. Their first origin is hidden in the unknown;
thereby they acquire a mysterious character. And since they contain the sum of the experiences the forefathers have had in the
various relations of life, they often afford a better guide than the individual's own more limited insight.

Where these unwritten laws are recorded and made known, the power of the state comes forward more definitely as that which maintains
and enforces the laws, and these now acquire the character of express declarations of will, of general rules for action. Law is now
no longer the sum of the involuntary norms of the people, but is the will of the ruling power. Over against the power of the state
stand the single individuals with their personal conviction, and here it may come to a conflict between law and morality. They are
no longer united in tradition,
% --- p. 522 ---
\opage{522}but what law commands or permits, morality perhaps forbids, and what morality commands, law perhaps forbids. A conflict may
arise here that often has a tragic outcome. Society must maintain its legal order, and the individual must maintain his conviction,
must ``obey God rather than men''\footnote[1]{Cf. \textit{Plato's Apology.} Chap. 17. --- \textit{The Acts of the Apostles} IV, 19;
V, 29. --- Jurists do not like to contemplate the possibility that the individual can have an ethical right against ``the law.'' Even
A. S. \emph{Ørsted} says: ``Everyone is bound to subordinate his private opinion to the one that lies at the basis of the laws. At
least he must be regarded as bound to do so by every public authority.'' \textit{(Om de første Grundsætninger for
Straffelovgivningen.} --- Eunomia II, p. 202). --- In \textit{the Danish civil Penal Code} § 42 there is talk of ``the
\textit{mistaken} opinion that the act forbidden by the law is permitted or even commanded by the dictates of conscience or of
religion,'' and so it seems as if it were laid down that the wrong must always be on the side of the individual. Ørsted is more
cautious, then, when he merely maintains that the public authority must necessarily regard the matter in this way. The framer of the
law seems to have wished to veil the possibility of a real conflict. But this possibility is always present as soon as any
distinction at all can be made between law and morality.}. % Acts 5:29, KJV; Ørsted and the Code translated from Høffding's Danish
The individual has here nothing else to do than to appeal from the present legal order to the legal order of the future, built on a
better foundation. --- But such cases of conflict must not be determinative for fixing the relation between law and morality in
general. Law and morality have developed historically from a common root, and only in exceptional cases do they become wholly
foreign to one another.

2. When the primitive stage of development, where the difference between law and morality does not yet appear, is left behind, and
when law is developed as a special domain of its own, there are especially four points at which it proves to be different from
morality: it relies on outer compulsion; it concerns only the outer action; it is more universal; it is more elementary.

a) A moral action must have sprung from the agent's own inner feeling and will. But the legal action can be enforced by outer power,
regardless of whether the agent's mind agrees with it or not. \textit{The possibility of compulsion} is the outer mark of law. The
legal order makes itself known by physical
% --- p. 523 ---
\opage{523}power, whereas the moral order of the world is an inner one. The sanction of law is an outer one, that of morality an inner
one. (Cf. IV, 5---6).

b) Closely connected with this is the fact that law demands merely \textit{outer action.} Disposition and will it does not care about;
these it cannot reach. For purely ethical valuation, on the other hand, the outer action may often become of vanishing significance
in relation to the motives that have led to it. These motives can be highly questionable, although the outer action in itself is not
pernicious; and they can be praiseworthy, although the outer action in itself is reprehensible.

c) Therefore law can also concern only such relations in which \textit{the same thing can be demanded of all with as close an
approximation as possible.} The great individual differences come out more in disposition and motives than in outer actions. People
of the same nation resemble one another more in their conduct than in their way of thinking. The action is not so individual as the
disposition.

d) And what can and must thus be demanded of all, and can be done whatever the motives, and can therefore also be enforced by power,
--- that can only be certain elementary ethical actions, that is, actions aiming at the maintenance of the most necessary conditions
of human life together. Law constantly tends --- or ought constantly to tend --- toward a minimum. This lies already in the fact that
compulsion is always an evil, and therefore, by virtue of the principle of welfare, ought to be limited to the least possible, the
most indispensable. Law is to fix by its strict provisions the limits below which human action must not sink, but above which it can
rise as high as it is able.

3. The legal order is a natural order brought about by human beings partly involuntarily, partly deliberately. Although it is
different from the ethical order, it nevertheless works as a preparation and foundation for it. It exercises partly an indirect,
partly a direct education. Indirectly it works by crushing all attempts to set oneself above the most necessary conditions of human
life together. Thereby the tendencies that conflict
% --- p. 524 ---
\opage{524}with these conditions are gradually checked. Directly it educates by providing a firm framework within which free forces can
move securely. It conditions a feeling of security and stability without which every higher endeavor becomes impossible.

But the connection between law and morality is far from exhausted by this. The task of maintaining law by power falls naturally to
those who command the collected, concentrated power of society. The ethical justification and necessity of this use of power rests
on the fact that the good of the whole stands above the desires of the individual. It is the same presupposition that holds for
every ethics that does not profess the standpoint of the moment or the individualistic standpoint. It is an ethical task to maintain
the fundamental conditions without which social life cannot develop. Limits must necessarily be set to the need of single individuals
to spread themselves in existence, if a social life is to be possible. --- Not only by its effects, then, but also by the
presupposition on which it builds, law maintained by power acquires an ethical character.

Despite the constant possibility of a conflict between law and morality, law is nevertheless always ethically conditioned and the
legal order a part of the ethical order. It is precisely this that gives such a conflict its tragic character. It is a clash of
ethical powers, or rather of powers each of which must lay claim to ethical authority, --- for when a clash takes place, one of the
parties must have lost its ethical character. Which of them in the particular case has forfeited its authority, that is precisely
the great question. Is it my conscience that, when it conflicts with a valid law, is a mere figment of the brain, or is it the
command of the law that, when it conflicts with the clear verdict of my conscience, is a mere arbitrary and indefensible demand? The
final decision will in all cases lie in the individual's conscience; for conscience judges everything, itself included (IV, 3). If I
choose to bow to positive law, I must be able to defend it before my conscience just as well as if I choose to break positive law.

Goos, whose development of the relation between law and morality
% --- p. 525 ---
\opage{525}has been very instructive for me, declares\footnote[1]{\textit{Almindelig Retslære} I. p. 47.} that positive law only forfeits
its ethical justification ``when a despotism ruling in the name of law rose to such a height that it could be said from the
standpoint of the idea of morality that no law was better than such a law.'' % translated from Høffding's Danish
The limit of the wrong of law in relation to morality seems, however, to be drawn too narrowly here. Not only where despotism, but
also where inertia rules in the name of law, can a positive legal provision lose its ethical character. A positive law can stand in
the way of an arrangement that the ethical law demands, and it does not follow that the fulfillment of the ethical demand can wait
until the valid law is replaced by a better arrangement in the way prescribed by the laws of the state. In every beneficial
reformation and revolution an ethical breakthrough takes place that bursts the barriers of positive law. Under ordinary conditions
ethical demands can find their outlet into the forms of positive law in a quiet and even way; but such an ethical tension may have
accumulated that a sudden breakthrough becomes necessary. Not only, then, when positive law stands very low, but also when the ethical
demand stands very high, the conflict will come, and come as an ethical necessity.

4. How definitely the difference between law and morality asserts itself when a higher stage of development has been reached can be
seen from the fact that it is precisely from a serious ethical standpoint that some have wished to maintain that law and morality
have no common root at all, and that law has no ethical character. \textsc{Kant} and \textsc{Fichte} conceived law as a purely outer
order that must necessarily be present when independent beings are to live together. The freedom of the one person must be limited in
such a way that the equally great freedom of the other person becomes possible. But law fixes only an \textit{outer} order, forbids
\textit{actual} encroachments on the freedom of others. Law demands merely legality, outer conformity with the rules valid for the life
together of several persons. Ethics, on the other hand, demands \textit{morality,} inner adherence to the recognized duty. Just as
little as the good will has anything to do in the domain of law,
% --- p. 526 ---
\opage{526}just as little does legality as such have any ethical significance.

\textsc{Fichte} made the most characteristic attempt to ground the doctrine of law as an independent science, wholly independent of
ethics. The principle of law presupposes merely the resolve of an individual to live together with other individuals. \textit{Why}
the individual makes this resolve, whether for ethical or non-ethical reasons, does not concern the doctrine of law. With him who will
not live together with other individuals the doctrine of law has nothing to do; none of its arguments holds against him. But if there
is a resolve of the content named, the doctrine of law shows that it can be carried out only if the individual is willing to limit his
freedom. If one has said A, one must also say B, and if one will not say B, one must also refrain from saying A: he who will not limit
his freedom cannot live together with other individuals. We thus arrive --- says \textsc{Fichte} --- at law without in the least
presupposing morality. --- But to this is further added that morality and law can conflict with one another. The moral law forbids me
to take the poor man's last lamb when he cannot pay his debt to me; but the juridical law permits it. How can such a permission have
any ethical grounding? From a common presupposition two mutually contradictory conclusions cannot spring\footnote[1]{J. G.
\emph{Fichte}: \textit{Grundlage des Naturrechts.} Jena und Leipzig 1796. Einleitung § 2 and Erstes Hauptstück § 4. --- Before Fichte,
Anselm \emph{Feuerbach} had put forward a similar theory, though not in so thoroughgoing a form. (Cf. A. v. \emph{Feuerbach}:
\textit{Biographischer Nachlass,} I. p. 23).}.

To this it must first be said that the legal order can very well have ethical significance because it is maintained by outer
compulsion, and because the individual is led already by purely egoistic motives to acknowledge it. When law maintains the most
elementary ethical demands, it is no wonder that it lays down something that can be recognized as necessary from standpoints quite
other than the properly ethical one as well. Security, order and firm organization can well be ethical necessities because they
% --- p. 527 ---
\opage{527}are also egoistic necessities. They acquire ethical value through what can unfold on their basis, and they spring from the
principle of welfare like all ethical demands.

As regards, next, the possibility of a conflict, there is nothing strange in the fact that one can arise, since the purely elementary
ethical demands must necessarily be roomier and lower than the higher demands that cannot be maintained by outer power. My duty sets
narrower limits to my freedom of action than my juridical right. Law is, as remarked earlier (2), more universal than duty, and
ethical justification (cf. VIII, 1) does not coincide with juridical justification. If I remain standing on the juridical standpoint
and push the stricter ethical demands aside, it is no wonder that a conflict arises. --- That the lower circle of rights must exist is
conditioned by the fact that otherwise worse drawbacks would arise than those that follow from a conflict between law and morality.
The creditor can use the law for inhuman conduct toward the debtor; but if harsh provisions did not hold here, levity in borrowing and
consuming would be fostered, since one could count on always keeping something back. To such unethical conduct on the debtor's side
there would be far greater temptation than there is to inhuman use of the creditor's juridical right. Naturally the provisions ought
not to be harsher than strictly necessary, and there is surely no question that here, as at so many other points, great advances in
humanity remain to be made.

Law, it is true, does not count on more than legality, outer conformity with the valid laws. But no legal community can in reality
exist if motives other than purely egoistic ones do not lead to obedience to the legal order. If the legal order is to be more than an
ordered state of war, it must have its subjective foundation in more than mere self-interest\footnote[1]{Cf. \textit{Om Grundlaget for
den humane Etik.} p. 56 ff.}. And this it will get when it answers to its end, to guard and protect the development of life in its
various directions. People will then look up to the legal order as the condition of the
% --- p. 528 ---
\opage{528}greatest goods. It will inspire in them not only fear and a feeling of dependence, but also gratitude and admiration. They will
see in it a framework within which everything they wish and strive for can develop securely, and even when they see the necessity of
criticizing and changing it, indeed perhaps even of breaking entirely with the track on which it has got, it is because they want to
make it better able to perform its office within the ethical world.

It can sometimes be precisely an ethical duty to make ruthless use of the right one possesses under the valid order, in order thereby
to maintain its validity and right understanding and to further respect for the freedom and independence of single individuals. The
struggle for law is a struggle for one of the most important goods of human life, even if \textsc{Jhering}\footnote[1]{\textit{Kampen
for Retten.} Danish transl. p. 48 ff.} goes too far when he asserts that of the two commandments ``do no wrong!'' and ``suffer no
wrong!'' the latter is the more important. It marked a great ethical advance when the proposition was first put forward that it is
better to suffer wrong than to do wrong (see XII, 2). Jhering's assertion is right only where people put up with wrong out of
easygoingness and indifference, or where they have lost sight of the practical significance of law for human life.

5. There are two main points in the history of the development of law that are of special interest in an ethical respect. The one
concerns the authority that fixes and maintains law, the other the subject that possesses the right, or that is to satisfy the demand
of law.

a) As long as law is not yet separated out from usage and \textit{custom,} but is contained in these in close union with morality and
religion, it cannot be formulated in clear and definite propositions. The way stands open to much arbitrariness and chance. It is
therefore a great advance when the power ruling in society acknowledges \textit{definite propositions as the norm} for its conduct.
Power thereby subjects itself to a higher law and can now be judged by whether it satisfies this law. It can be measured with a
yardstick it has itself acknowledged. It is an important turning
% --- p. 529 ---
\opage{529}point in the history of authority\footnote[1]{\textit{Om Grundlaget for den humane Etik.} p. 75---79.}. The inconsistencies of
power will in the long run come down on its own head. Even if originally it only put on law as a kind of disguise, it has thereby
challenged a dangerous enemy. It has challenged \textit{criticism,} and it will not be able to stop this again except by
systematically annihilating every freethinking consciousness. That it is not able to do this is shown by the fact that as a rule it
--- even if it is done by ever so many distortions and sophistries --- seeks to procure for itself a semblance of legal grounding for
its conduct. Once criticism has begun at one point, it soon goes further. It does not confine itself to testing the agreement between
the particular use of power and the law that power has itself recognized as valid; but it turns against the valid law itself and
demands that its inequalities and drawbacks be removed. Already through public criticism and debate the thoughts and feelings of the
people gain influence on the fixing and application of law. The holders of power themselves will in the long run not be able to
withdraw from this influence; in their own consciousness these thoughts and feelings assert themselves to a certain degree, just as
already at the stage of customary law the holders of power also feel the power of tradition (XXXVI, 5). --- That has even been
praised as the true \textit{free} constitution in which the final decision of affairs lies, to be sure, with the absolute monarch, but
in which unconditional freedom of discussion prevails, so that a control of all over all is exercised. The monarch then merely lays
down what has developed in the mind of the people as the fruit of the many-sided debate\footnote[2]{Cf. F. C. \emph{Sibbern}:
\textit{Dikaiosyne.} Kbhvn. 1843. See my essay on Sibbern in ``Mindre Arbejder'' I. p. 99---103. --- In a similar way, according to
Cardinal \emph{Newman} \textit{(Apologia pro vita sua.} 1864. p. 271 ff.), the infallible Pope himself is to have only the task of
``defining'' what the church already believes beforehand: he is merely to lay down the result of the development of faith within the
church. Catholics, said Newman, do not believe in the Immaculate Conception of the Virgin Mary because it has been laid down by the
Pope; but it was laid down because Catholics believed in it!}. But here there is still too great a distance between the debate and
the decision, and no guarantee that the latter really comes into as
% --- p. 530 ---
\opage{530}close a connection with the former as possible. Law therefore does not yet come into the intimate connection with the
consciousness of the people that is necessary for its secure existence. That happens only where \textit{the people} has political
freedom, and with it \textit{direct influence on the fixing of law} through legislation. Only thereby will it also be possible for the
particular interests of the holders of power to be more and more subordinated to the common interests of society. --- But the living
sense of justice in the people is --- whatever constitution one may imagine --- always the last bulwark of the legal
order\footnote[1]{Cf. \emph{Jhering}: \textit{Der Zweck im Recht.} I. 2. Aufl. p. 381: ``On the moral power of the national sense of
justice rests in the last instance the whole security of law. \textit{Not} on the constitution --- devise it as artfully as one will;
none can be conceived that would actually deprive the power of the state of the possibility of trampling the law underfoot.
\textit{Not} on the oaths by which one thinks to secure it --- experience shows how often it is broken. \textit{Not} on the nimbus of
sanctity and inviolability with which theory clothes the law --- it does not impress arbitrariness.'' % translated from Høffding's Danish
}, just as it is the source from which the legal order develops in the first place. Where this source dries up, life is over, even if
the wheels of the state mechanism may go on clattering for a time.

b) At primitive stages it is not the single individuals but the kindreds, the family groups, that the legal order has to do with. The
individual is regarded only as a member of his family and his tribe. Within these groups community of property prevails, perhaps even
with regard to women and children. Marriage is not an affair between two individuals but between two kindreds, which thereby come into
connection with one another. If a murder or a robbery takes place, the whole kindred to which the perpetrator belongs is struck by the
retribution. There exist no individual obligations and no individual guilt. Only gradually, as instinctive custom develops into a
proper sense of justice, does an \textit{emancipation of single individuals} take place, by which each of them becomes a legal person, a
center of rights and obligations, a subject that can incur guilt and pay the penalty for it. Traces of the primitive way of thinking
are still found in our legal order (cf., e.g., XXII, 1).

6. An intermediate link between positive law and ethical conviction
% --- p. 531 ---
\opage{531}is formed by \textit{public opinion.} It can be defined as the popular expression of positive morality. It is a ``moral police''
that exercises a mighty influence by its praise and its blame. It draws somewhat narrower limits for what is justified than law does.
Law keeps to the outer action; public opinion goes further and passes judgment on character. Circumstances come into consideration
before its tribunal that the juridical tribunal cannot take into account. Therefore it sometimes acquits those who are condemned
juridically, although still more often it accuses those who cannot be brought before a juridical tribunal.

It can form in very different ways and be of very different value. Its origin can often be difficult to trace, because it is the
expression of a long series of constant experiences that have not each drawn attention to themselves, but whose result only then
emerges clearly. In other cases it can be sudden experiences, had in common, that at once set the same mood in motion in all. In both
ways opinions come to be ``in the air,'' and in both ways they can assert themselves with great force. The unknown in their origin
gives them a certain mysterious character and stamps them with the seal of the self-evident. Often it can be single eminent men who
determine public opinion by their authority, especially when they are able to clothe their thoughts in striking words. But it can
also be the persistent repetitions of unknown and insignificant persons that make opinions come to be ``in the air.'' A constant
interaction takes place between the narrow circle of people who have new experiences and get new ideas and the great circle that for
the most part receives their opinions in a passive way. No one, however active his life of thought and feeling may be, will be able to
avoid being determined by the presuppositions and traditions of his time and people, or of his class and family; but one can place
oneself in a more or less active relation to these and have a more or less wakeful eye for new possibilities. In different countries
the proportion between those who actively form the prevailing opinions, those who passively receive them, and those who --- have no
opinion, will be different. \textsc{James}
% --- p. 532 ---
\opage{532}\textsc{Bryce}, who in his work on the American republic undertakes an interesting investigation of the rise and significance
of public opinion, thinks that in the United States the number of those who belong to the first class is very small, whereas in England
it is comparatively large, but that on the other hand the third class is over there comparatively far smaller than in England --- and
would be far smaller still if immigrants and Negroes did not lower the proportion. It is therefore more difficult in the United States
than elsewhere to point out the successive development of public opinion; its rise and development have there, even more than
elsewhere, the character of a growth. There are few superior intellects, but there is a great multitude of people who are able to
take up ideas, and these spread extraordinarily quickly in great circles\footnote[1]{James \emph{Bryce}: \textit{The American
Commonwealth.} London 1885. III. p. 10---13; 99---105. % ``1885'' PRINTED AS IS (the first edition is of 1888; cf. p. 494)
--- Bryce found, especially in North America, a more critical view of the press among the public in general; the American public is
cooler and more cautious than the European and is not so easily impressed by the mysterious ``we.'' (p. 37 f.).}. --- In the spreading,
reflection and conviction often play only a small part, but the impulse to imitate plays a main role, and not infrequently a public
opinion comes into being merely because it already seems to exist!\footnote[2]{\emph{Holtzendorff}: \textit{Wesen und Werth der
öffentlichen Meinung.} München 1879. p. 93.}

Public opinion can often, ethically regarded, stand higher than any of the persons who contribute to it and profess it. This is in
itself no stranger than that the legal order also stands higher than the actual condition of the country, which the necessity of penal
provisions already shows, --- or than that the individual's ethical conviction too can have an ideality that his actual willing and
acting are far from possessing. But it is surely also connected in part with the fact that each of us blames the motes in others and
forgets the beams in ourselves; there results from this in public opinion a strictness that the individuals each take good care not to
apply in their self-criticism\footnote[3]{Cf. \emph{Rümelin}: \textit{Ueber den Zusammenhang der sittlichen und intellectuellen
Bildung.} (Reden und Aufsätze. Neue Folge. p. 24).}.

% --- p. 533 ---
\opage{533}Just as the legal order, so too public opinion ought to be limited to a necessary minimum. It is more intrusive than law, since
it also judges character; but nonetheless it does not command the thorough methods that are at the service of the legal order when it
is a matter of throwing light on the nature of an action. Character is more difficult to judge than the outer action. The course of
development of the will will always be more or less hidden from outsiders. Public opinion is nevertheless not inclined to recognize its
limits; it easily credits itself with omniscience and becomes intolerant. It lumps everyone together. It is often governed by
prejudices of class, race, party and religion. Despite all this it plays an important role both toward the power of society and toward
single individuals. It is in its way, like the legal order, a kind of natural order that sets limits and conditions. It exercises a
control that cannot be dispensed with. But just as the legal order in the end springs from the people's sense of justice and exists
only by its connection with it, so it is also of the greatest significance that the stream of public opinion be fed from the serious
conviction of individuals, so that it does not become fleeting and superficial. And far more frequently than conflict arises between
law and morality, conflict will arise between public opinion and the individual's conscience.

It was especially fear that the growing power of public opinion would become dangerous to the independence of character that in his
time led \textsc{Stuart Mill} to write his book ``On Liberty.'' But the solution he gave of the question was, as I have tried to show
above (VIII, 5), an unfortunate one. One cannot agree with \textsc{Mill} that the individual's inner self-development concerns only
himself and not others. There is no trait of character in a human being that cannot gain significance for his relation to others and
for his position in society. Nor will one be able to prevent others from making a valuation of the picture of his character that forms
in their consciousness. It is not the very justification for judging, but the uncertainty of the data by which one judges, on which the
weight must be laid when public opinion is to be pushed back within its proper
% --- p. 534 ---
\opage{534}limits. On the other hand, the individual is often himself to blame for the public misjudgment that falls to his share, when he
does not do enough so that his conduct may come to stand in the right light for others. One has no right to let oneself be made a martyr
when one has not done everything possible to set forth the real connection of one's cause as clearly as possible. Even the ironist
Socrates spoke out plainly when it became serious, and when it was a matter of preventing his countrymen from committing a crime.

\chapterhead{XXXVIII.}{THE ETHICAL SIGNIFICANCE OF THE STATE.}

1. \textsc{Most} directly and conspicuously the ethical significance of the state would appear if those were right who conceive the
state as an immediate revelation of the idea of the good, as an ethical authority that stands independent of and raised above
individual consciences. The power of the state would then be the ethical guardian of individuals, to which they had to bow.

This view appears in its most extreme form where it is asserted that the state is to be built on a theological foundation. The state
is here not only an ethical power, but --- it is added --- since all ethics rests on religious authority, the state must in the end
build on a religious foundation. The state needs --- it is said --- ``a supreme, reliable standard'' in order to be able to make a
valuation of the human ends that it is its task to further. Outward justice cannot be carried through without inward justice, and
this in turn presupposes a religious disposition that can be won only through Christianity. The state must therefore be a Christian
state. This presupposes ``not that living personal Christianity is to be found in everyone in the people, but that the people bows
under the authority of the Christian tradition.'' The Christianity of the state shows itself first and foremost in its supporting the
church; and
% --- p. 535 ---
\opage{535}next in its providing for Christian usage and custom in the people and for Christian schools, and in its impressing Christian
principles in general on its laws and institutions\footnote[1]{\emph{Martensen}: \textit{Den sociale Etik.} p. 125 ff. Martensen's
concept of a ``Christian state'' nevertheless means a great abatement in relation to \emph{Stahl's}, who demanded exclusive protection
for the Christian church and a Christian confession of faith as the condition for all offices and for a seat in the representative
assembly. And Stahl's doctrine is in turn an abatement in relation to what the Holy Alliance demanded. (Cf. \emph{Bluntschli}:
\textit{Politik.} p. 226---228). The Catholic view naturally goes further still, since it demands that the supremacy of the church be
acknowledged in all questions concerning both church and state. The Pope protested against the Holy Alliance because the princes in
it wished to act as the vicars of Jesus.}. % Martensen translated from Høffding's Danish

There is something ambiguous in this doctrine. It does not tell us in plain words where the reliable standard is really to be found.
Christian tradition and Christian principles are loose concepts, especially when it is expressly admitted that it is not presupposed
that everyone in the people has a living, personal faith in Christianity. No one can deny that Christianity has made significant
contributions to the development of ethical consciousness in the human race. But these contributions must be supplemented by
contributions from other sides and are neither in form nor in content sufficient to serve as a reliable standard. When the
theological side boasts of having a reliable standard, there is sense in it only if the principle of authority is consistently held
fast. So the authority of church faith (whether this be the Pope, the Bible or the creed) is to be the foundation of the state. What
follows from this is a theocracy or a church-state. The state becomes not a purely human community, but dependent on a supernatural
authority. There can be no question of anything of the kind from the standpoint we have taken here. Just as we have maintained the
freedom of inquiry and of conscience, so we must also maintain the freedom and independence of the state over against the
theological principle of authority. The relation is even, as we have seen in the foregoing (XXXIII, 3), the reverse of what the
adherents of ``the Christian state'' imagine: it is the state that in the end decides whether the various
% --- p. 536 ---
\opage{536}religious confessions agree with ``morality and public order''; it cannot itself take its concept of morality from any of
them.

From the standpoint of Christianity itself there is sense in the idea of a Christian state only if one can start from the assumption
that everyone in the people is a Christian. But it was expressly said that one did not start from that! Christianity is thus to rule
over those who do not believe in it! --- All this talk of a Christian state is unclarity from first to last.

Over against the unclarities of High Church theologians it is refreshing to see \textsc{Grundtvig's} sound and clear conception of the
nature of the state. However deeply he was personally seized by Christianity, he nevertheless saw that ``where civil conditions are
truly to be reborn and gain firmness, there one must leave Christianity, as something altogether free and incalculable, altogether
out of the reckoning \dots{} and keep to the human nature that history reveals in every given place''\footnote[1]{\textit{Mands
Minde.} p. 127 cf. p. 569 f.}. % translated from Høffding's Danish

2. Even if one will not build the state on a theological foundation, one may still think that the state is the expression of a
higher morality, that its laws teach the individual what is good and right. The individual, no doubt, has his subjective ideals, but in
the state the ideal meets him as an actual power. The ancient conception of the state, as it was developed by Plato and Aristotle, went
in this direction. In modern times \textsc{Hegel} in particular has asserted this view in an extreme form. ``The state is the
actuality of the ethical idea.'' ``The individual has truth and ethical life only insofar as he is a member of the state.'' The power
of the state is ``reason actualizing itself as will''\footnote[2]{\emph{Hegel}: \textit{Philosophie des Rechts.} § 257---258. --- Cf.
\textit{Den nyere Filosofis Historie}\textsuperscript{2} II. p. 180---183.}. % Hegel translated from Høffding's Danish

This view has in common with the preceding one that it abolishes the independence and free conviction of the single personality.
Single individuals are regarded as merely dependent and receptive in relation to the state, not each for himself as constituent
members of it. The state is given a kind of
% --- p. 537 ---
\opage{537}mystical life outside or above its citizens. But the state lives only in its citizens. (Cf. VIII, 4). The spiritual and
physical power that the state commands is that which its citizens are able to supply: the thoughts and feelings, the strength that
they possess. The morality that is in the state must be in its citizens. Or could perhaps a moral state consist of immoral citizens?
That is surely just as impossible as that a state can be called Christian when it does not consist of believing Christians.

The view discussed fits best the stage of development at which an opposition between the power of the state and individual freedom
does not yet make itself felt, --- at which the difference between law, morality, usage and custom has not yet arisen. Ethical ideas
are not seen here as produced by the single individual; they are components of a venerable tradition on which individuals feel
themselves dependent and to which they bow. But among the civilized nations this instinctive adherence to what has been handed down is
supplemented and controlled by public criticism, and single individuals do not merely behave receptively but react actively on the
life of the state.

By attributing direct ethical authority to the state one is in reality compelled in the end to appeal to physical power. The power of
the state cannot let it depend on whether the individual feels convinced of its ethical justification. It then uses its sharp means,
and ``the actualization of the ethical idea'' then advances with police and cannon. The overstrained ethics reveals itself as physics
in disguise.

3. Now the nature of the state has been found precisely in mere power. It is especially pessimists who emphasize this factor in the
concept of the state. Thomas \textsc{Hobbes} conceived the state as the unconditional will over against which individuals give up their
individual wills so that peace and security may prevail. \textsc{Schopenhauer} described the state as an institution of compulsion
whose only purpose is to protect individuals against one another and the whole society against outward enemies. \textsc{Taine} finds as
the real kernel behind the whole machinery of the state ``the gendarme armed against the savage, the robber and the madman whom each of
% --- p. 538 ---
\opage{538}us hides in the depths of his heart, slumbering or chained, but still always alive''\footnote[1]{\emph{Hobbes}: \textit{De
cive} V, 9. --- \emph{Schopenhauer}: \textit{Die beiden Grundprobleme der Ethik.} 2. Aufl. p. 217. --- \emph{Taine}: \textit{L'ancien
régime.} p. 316.}. % Taine translated from Høffding's Danish

This view can appeal not only to the fact that the last means of the state is always the use of power, but also to the fact that most
states have historically been founded by conquest. By the sword the state came into being, and by the sword it endures. But it has
always proved that mere power is too uncertain a bond. The realms founded by conquest first gained firmness when common customs, forms
of life and traditions could form. Through common fate and activity the human groups brought together grow together into a people.
(Cf. XXXVI, 2). By mere power one can gather a horde; but a real state presupposes that there is a people. Where the power of the
state and the life of the people go in quite different directions, there is an unstable equilibrium, and a small shock may be enough to
overturn the whole.

Mere power can maintain security and unity in the people. No one murmurs, and no one leaves the ranks. But security and unity are goods
only insofar as they are conditions for a free and secure unfolding of the life of the people. One must beware here of making the means
the end. The more one lets the means of security spread, the more posts one sets up to forestall dangers and misfortunes, the more one
hinders the free movement that might perhaps, in its own way, be able to keep misfortunes away. A nervous staring at possible dangers
can destroy the whole life of the people, just as the single individual can be prevented from living by fear of death. There can of
course be times for a people when it ought to be especially or exclusively occupied with providing for its security and its unity. But
it is wrong to find the nature of the state exhausted in the role it is compelled to play in such times. Centralization and violent
concentration are an exertion of strength that may be necessary and beneficial, but that would stifle all natural and independent
development if it became constant and unlimited.

% --- p. 539 ---
\opage{539}Besides, power cannot and must not be blind. It must be guided by rules and principles, and these cannot be taken from power
itself, but must be derived from the ends that power is to protect. What power can produce, security and unity, is only a means for
higher ends that power cannot produce, but that free development produces, and to which power can only then come into relation. From
mere power neither culture nor law can be derived.

4. Both when the state is conceived as an immediate revelation of the idea of the good, and when it is conceived as the power ruling
over all and everything, the main weight is laid on the \textit{unity} in the state, and it is made all in all. Over against such views
it is rightly urged that what is really alive in the state is, after all, the single individuals. The unity consists in their
community and union, but is only a form, whose significance depends on the content with which it is filled. It might then even seem
that one must take the living components, the elements, the single individuals, as the basis in considering the nature of the state,
and conceive the nature of the state as conditioned by their union. The unity then becomes not original but derived. We thereby arrive
at the \textit{individualistic} view, according to which the state rests on an agreement or contract between the single individuals.

This theory has often been wronged by the assumption that its meaning must necessarily be that the state had \textit{historically}
arisen through a contract. That is in any case not the meaning of all individualistic theorists of the state\footnote[1]{Cf.
\textit{Den nyere Filosofis Historie.} The index under ``Naturret.''}. The theory uses the concept of contract in order to make clear
how the community of the state would have to be ordered if it answered to its ideal. It is thus used as a standard in the valuation of
the actual state, not as a means of explaining the historical origin of the actual state. The perfect state --- it is thought --- would
be one in which everything was so arranged as the single individuals could wish it by free agreement, --- so arranged, that is,
% --- p. 540 ---
\opage{540}that account was taken of the interests of all, insofar as they could be brought into harmony with one another.

What can be traced of the historical development of the relation between state and individual does not give the idea either that the
state historically arises through the union of independent individuals. Apart from conquests, which of course already presuppose
states as given, the formation of states seems to have taken place through the fusion of hordes, groups or single families, sometimes
no doubt also through simple growth and development within one and the same family. Of single independent individuals there is at
earlier stages really no question; they are regarded only as parts of the family or the tribe. But in the course of historical
development a dissolution or ``crumbling'' of the original groups takes place. The individual is emancipated step by step. The
principle of freedom slowly works its way forward. (Cf. XXIII). This process of emancipation comes out more and more distinctly as one
goes from east to west: it is less prominent in Asia, where, e.g., the Indian village communities have held their ground to this day,
than in Europe, and less prominent in eastern Europe (the Slavic countries) than in the western (England and France)\footnote[1]{\emph{Henry
Maine}: \textit{Ancient Law.} p. 311. --- \textit{Early History of Institutions.} p. 386 f.}. A thoroughgoing individualistic view
imagines this process of emancipation completed and measures the historically given forms of state by the degree to which they approach
this ideal.

The great significance of individualism lies first and foremost in its emphasis on the principle of personality, on the value of each
single individual. It is the great task of the state, as of every other community, to work --- directly and indirectly --- toward the
development of a free personal life in as many individuals as possible. The perfect state would be one that satisfies the personal life
of all, each one's need for the development of his abilities and impulses. Individualism has therefore grasped an essential side of
the ideal of the state.

Next, individualism is also right in that the deliberations and conscious choices of single individuals in fact gain more and more
% --- p. 541 ---
\opage{541}influence on the ordering of the state and on the decision of affairs of state. Where political freedom prevails, the individual
has access to letting his will help determine public affairs. Not only at the founding of new states or federations of
states\footnote[1]{\emph{Alexander Hamilton}, during the deliberations on the North American federal constitution, held up to his
countrymen that ``it seems to have been reserved to the people of this country, by their conduct and example, to decide the important
question, whether societies of men are really capable or not of establishing good government from reflection and choice, or whether
they are forever destined to depend for their political constitutions on accident and force.'' \textit{The Federalist.} 1787. No. 1.
(New edition, Boston 1882, p. 49). % Hamilton's English (Federalist No. 1)
}, or at the introduction of new constitutions, but also in the single decisions during the ordinary course of political life, it
comes more and more to depend on bringing about agreement among the various individual wills taking part. Along this path everything
that lives and develops within individual personalities and in the smaller circles of society can benefit the life of the state.
Individualism thus points not only to the last end of political life, but also to one of the most important means of its progress.

But in its extreme forms individualism lays the weight one-sidedly on the \textit{separateness} of individuals, regards each of them as
absolutely independent and sets up the concept of the sovereignty of the individual. It is one thing that every single individual is a
peculiar and independent starting point, another to attribute to it an absolute significance apart from its relation to other
individuals. If the word sovereignty is to be taken in the strict sense\footnote[2]{Which seldom happens. Cf. \emph{Stuart Mill}:
\textit{Levned.} Danish transl. p. 259. \emph{Dühring}: \textit{Cursus der Philosophie.} p. 268. When it is said here that only those
encroachments on sovereignty are permissible that are necessary to maintain the equal sovereignty of all, one obviously cannot derive
these limitations from the sovereignty of the individual; why should it limit itself? The limitation presupposes a standpoint that at
once embraces all individuals and where separateness therefore does not hold.}, every community will only be able to become a crowd or
a heap of individuals; there will be no inner
% --- p. 542 ---
\opage{542}connection among them. They come to terms with one another but do not live a common life. A really common life becomes
possible only when the individual from the start regards himself, and is regarded, as a member of the race and of society.
Sovereignty --- absolute independence and plenitude of power --- cannot lie with any single person or with all single persons regarded
as each isolated. It can lie only with \textit{the people,} to which individuals belong through community of fate and activity.

Consistently, individualism must go a step further. No form of state can let individual wills rule without interruption. Even where
the greatest political freedom prevails, it is only on single, especially important matters that the whole people is consulted, as in
the popular votes (referenda) in Switzerland and some other countries. In the interval between such votes individual wills recede,
and the central organs of the state govern the course of things. Strict individualism must regard such an interval as a period of
unfreedom. It is therefore consistent when \textsc{Rousseau} writes: ``The English people thinks it is free. It is greatly mistaken. It
is free only while members of Parliament are being elected; as soon as they are elected, it is a slave, it is nothing!''\footnote[1]{\textit{Contrat
social.} III, 13.}. % translated from Høffding's Danish
Rousseau therefore rejected the representative system: sovereignty cannot be represented! But he overlooked that within the single
individual, too, the choosing will does not stand isolated from the other elements in man's nature. It has significance only by being
the outcome of the individual's whole nature and character, and when the individual takes the consequences of his choice, he therefore
does not bow to any foreign power, but to his own work. Thus it is no breach of the voters' freedom either that they are bound by their
own choice. Without such boundness there is neither character in the single individual nor national character in the whole nation.
---

5. That such different conceptions of the state as those now discussed can assert themselves comes first and foremost from the fact
that the state is a many-sided being and that its concept
% --- p. 543 ---
\opage{543}contains elements whose mutual relation is difficult to determine. But it is also connected with the fact that the state is
still in its becoming, since the different elements in its nature have not yet been brought into firm connection with one another, and
since in the course of development now one, now another element pushes into the foreground. Power, law and culture stand at different
times in different relations to one another. The question now is whether we can state a normal relation among them, i.e., a relation
that best answers to the demand of the principle of welfare.

Of great significance here is the recognition that the involuntary always precedes the deliberate. The institutions and arrangements
that human beings make with full consciousness and deliberation always presuppose that a content has developed involuntarily that can
be ordered and applied in a certain way. In our institutions and arrangements we never begin absolutely afresh. This holds
psychologically for the relation of the choosing will to the other sides of mental life, and it holds sociologically for the relation
of the state to the people. The state is to order and protect the life of the people, but it does not produce it in the first place.
--- First and foremost, \textit{ethical life} in its highest form and development is independent of the state. It is due to feelings
that develop in the life of individuals together but that cannot be enforced by power. The state can inspire fear by its mere power;
but it does not get far with that. Reverence it first inspires when the individual finds that the state is guided by the principles to
which he himself feels bound inwardly. When the state demands submission on ethical grounds, it thus appeals to a power that it has
not itself produced, but that it presupposes to be present. The ethical principles that the state lays down and applies in its
constitution and its laws have first developed in the consciousness of the people. --- Nor does the state produce \textit{family life.}
This has its root in human nature; its forces and its forms are subject to the intervention of the state only in their outer appearance
and expression, not in their origin and inmost nature. --- Nor is \textit{culture} the work of the state. It develops in free society.
The origin both of material and of spiritual culture always leads us back to single
% --- p. 544 ---
\opage{544}individuals in whose consciousness the new thoughts came into being. Within the individual's consciousness these new thoughts
have their first struggle for existence to undergo. When they go out into the world, they must first gather small circles or
congregations round them, within which they can find their further nurture and development, and thereby become fitted to influence
larger circles. The significant processes of development in history usually take place sporadically, from scattered points, not from
the center. They often take their beginning where one would least expect it. People then ask in astonishment: Can any good thing come
out of Galilee? % cf. John 1:46 (KJV ``Nazareth''), 7:52
In a remote corner of the great Roman empire a force was born that was to live longer and work further than the powers that people in
Rome regarded as eternal. Who in the twelfth century could have imagined that the movements that expressed themselves in the French
``communes'' were the prelude to a complete reordering of the whole of public and social life? In dark workshops and in despised
circles of workers one of the most significant social phenomena of the nineteenth century, the English trade unions, came into being.
What has thus developed in secret can later benefit the life of the state in many ways, but the latter cannot itself produce it in the
first place. --- The same holds even, as we have seen, of law. It develops as an involuntary form of life in the people; the state can
only fix, systematize and maintain it.

The state is just as little productive as the choosing will is in the single human being. The possibilities between which a choice is
made must always be given beforehand in human nature. Choice naturally has its great significance, even if it is no creative act. What
is chosen was there beforehand, but now passes over into a firmer and more definite form.

6. The proposition that the state is not productive holds, however, only with regard to the \textit{content} of life. The state unfolds
its productivity where it is a matter of \textit{ordering this content in the form of law.} Here a series of significant tasks is set at
all times, in the treatment of which the genius of the statesman is to show itself. The gift of the great statesman forms a counterpart
to that of the great founder of a religion. The latter's significance too does not lie
% --- p. 545 ---
\opage{545}(cf. XXXII, 1---2) in producing wholly new thoughts, but in ``condensing'' in a peculiar way spiritual elements and forces that
are already stirring. The great statesman is he who with a gifted and sympathetic eye discovers the new germs of life in the people and
--- directly or indirectly --- gives them the help and support they need. When egoism and ambition do not dull his gaze, he knows that he
can do nothing but help forward what has sprouted without him. His productivity consists in finding the form and means through which
this help can be given. --- This task will be essentially lightened where political freedom prevails. A more intimate relation arises
between the state and free society, between the form of life and the content of life, when the people through its representatives
takes part in deciding how the power of the state is to be used. Only through political freedom does the state really become the
organized people, since the cultural community and the state can now come into a living relation of interaction with one another.

\textsc{Gneist}, who has described the development and organization of self-government in England in so thorough and interesting a way,
regards the whole recent social and political movement (for England from the first parliamentary reform on) as a great process of
dissolution. Society --- he thinks --- has forced its way into the state and seized its power for its own interests. How he conceives
the relation between free society and the state can be seen from the following statement: ``What the strife between duties and
impulses means in the life of the individual, the eternal struggle between state and society means in the life of
peoples''\footnote[1]{\textit{Das Selfgovernment in England.} 3. Aufl. p. 1017.}. % translated from Høffding's Danish
But even if, with Gneist, one means by society especially economic society, in which material culture is cultivated, this is far too
one-sided a judgment. In the domain of material culture forces and feelings of just as much value as those fostered by political life
can be developed. Material culture furthers not only self-interest and egoistic calculation; it develops the will, the practical eye
and the ability to cooperate with others. And it is no wonder that those who feel
% --- p. 546 ---
\opage{546}these forces in themselves also think they have a right to put in a word in the public affairs of the country. That is the
condition for the state's being able to continue to be the expression of what is stirring in the people. --- But it is one-sided here to
think merely of economic and material social life. The movement that has led to the extension of the franchise is closely connected
with intellectual and religious development, and of this one surely cannot say that it stands to the activity of the state as impulse
to duty\footnote[1]{In remarkable contrast to Gneist's repeated complaints about ``the dissolution of the old society'' stands a single
statement of his (p. 996 Note). ``The idea of a dissolution of society,'' he says, ``rests only on the fact that accustomed ideas are no
longer able to discover the accustomed barriers.'' --- When it is merely the old rubrics that are missed, there seems not to be so much
reason for complaint.}. --- Moreover, Gneist admits that the old order of things had its great dark sides: no account was taken of the
welfare of the middle class; no care was taken for the moral and spiritual elevation of the masses or for the cultural purpose of the
state\footnote[2]{Op. cit. p. 891. 938.}. An order of state against which such reproaches can be raised surely deserves to be replaced
by another arrangement! Gneist's great admiration for the old English self-government is not well compatible with this sharp judgment,
which makes the process of dissolution come as a Nemesis. ---

It is no underestimation of the significance of the state when its productivity is limited to finding and maintaining forms for a
content that is given without it. But this view, if it is correct, gives us a warning against laying too great weight on political life
in its separation from the other sides of life. In the struggles of political life it is only too often forgotten that the forces of
the people must first and foremost develop in other domains, in family life and in the free cultural community, and that only when this
has happened, and constantly happens anew, can they be applied fruitfully in the political domain. In politics one easily comes to set
form above content.

7. In the firm forms by which the life of the people is ordered, and by which it becomes able to act with concentrated force
% --- p. 547 ---
\opage{547}outwardly, the unity of the people shows itself most clearly. Single individuals, families and the free cultural community
present the picture of a multiplicity of changing phenomena. The state represents not only the unity of the people, but also the most
important conditions for its lasting existence, reaching beyond the time of single individuals and generations. In the state the people
has its most concentrated expression.

With this nature of the state the general rule is given for which tasks it will be able to solve. Since the state is the centralized
power of the people, what it is to be able to help forward must be of such a nature that it can be directed systematically from one
point, and that the use of power can be employed as a last means to keep it up. Where these two conditions are lacking, the state can do
nothing. But no domain can be named where the conditions named cannot be present at some stage of development and in some degree. ---
What and how much the state draws under its activity will differ under different conditions. What one age cannot imagine the state
concerning itself with will perhaps at another time become its main task. Either it takes upon itself the whole activity, as with the
administration of justice and defense, or it sets up institutions with which it allows private institutions to compete, as with
education, or it supports private institutions with material assistance, or it confines itself to securing for free endeavors the
support that the general legal order affords.

The real strength of the state consists in the adherence that the power of the state is able to win among the whole people. The greater
the freedom the state is able to leave to individual, family and cultural community, without its unity suffering by it, the stronger it
is. The extent of personal and political freedom is therefore one of the most important means of judging a state's stage of
development.

Full harmony between state and people, power and freedom, unity and multiplicity is a very distant goal. In reality every state is a
state of necessity, answering only very imperfectly to its concept. --- On the one side, power still always acts in a more or less brutal
way. To exercise power is bound up with far too great a gratification for human
% --- p. 548 ---
\opage{548}egoism for abuse not to appear. To this are added fear and a blind impulse of self-preservation, which can lead the power of the
state to spread at the expense of the ends it was to serve. A long course of development will still be needed before we get so far that
compulsion is used only where it is absolutely necessary by virtue of the principle of welfare. --- On the other side, freedom ---
not infrequently occasioned by encroachments of the power of the state --- often expresses itself as arbitrariness, censoriousness,
suspicion and egoistic individualism. --- These two extremes provoke one another. A front must now be made to the one side, now to the
other, so that development may through rhythmic oscillations slowly carry us forward in the direction of the goal indicated.

\chapterhead{XXXIX.}{THE PUNITIVE AUTHORITY OF THE STATE.}

1. Historically punishment has developed out of the instinct of revenge or retribution\footnote[1]{According to more recent
investigations, however, the need for retribution is preceded by the need for purification from the defilement that the criminal,
especially the murderer, brought upon the whole social group where his deed had been done. The reaction of the group was then at once
a disinfection and an exorcism. The transition to punishment consisted in the purification being accomplished through the criminal's
exile or death. See on this \emph{Frazer}: \textit{Psyche's Task.} A discourse concerning the influence of superstition on the growth
of institutions. (1909). p. 79 f. (1913).}. In its simplest form this instinct expresses itself in a movement directed toward the
quarter from which one has received a strong and painful impression. One flies at the real or supposed attacker. This instinct is
partly a special form of the instinct of self-preservation, insofar as the reacting movement leads to removing
% --- p. 549 ---
\opage{549}the cause of the pain, partly it is connected with the need for movement, for getting relief, that accompanies every excited
state of consciousness. It works blindly. It knows no other limits than the exhaustion of its strength. A comparatively slight
influence can call forth the most violent reaction when the state is already excited, just as a little spark can call forth the
greatest explosion. But the blindness shows itself not only in the fact that there is often no reasonable proportion between
influence and reaction. It shows itself also in the fact that no distinction is made between the real and the apparent cause of the
painful influence, and in the fact that the further consequences of the reaction do not become an object of attention. The essential
thing is to get relief for the inner unrest and pain; --- what objects it falls upon, and how much it falls upon them, plays no role.

The development that the instinct of revenge undergoes in the course of the development of social life and of culture consists partly
in the object of the revenging reaction being individualized, so that no one is struck but precisely the one from whom the influence
proceeded, --- partly in the strength and kind of the revenging reaction being determined in proportion to the strength and kind of the
influence and to the degree of consciousness with which it was exercised, --- partly in account being taken also of the further
consequences that the revenging reaction has for the individual it strikes and for society in general. All this presupposes that the
instinct is stopped in its course, --- that an interval is inserted between influence and reaction, that interval which is in general
so significant in the development of consciousness and especially of the will\footnote[1]{\textit{Psykologi} IV, 4. VII B, 1---3.}, ---
and finally that the reaction becomes not the affair of single individuals or kindreds, but of the state. Only through this series of
thoroughgoing changes does revenge become punishment. We shall consider some of these points more closely.

2. At primitive stages of culture the single individual is not regarded as isolated, either as regards rights or duties. It is
families and kindreds, not independent individuals, that
% --- p. 550 ---
\opage{550}stand over against one another. Revenge is therefore exercised by \textit{any} member of the kindred to which the injured
individual belongs against \textit{any} individual of the kindred to which the injurer belongs. Blood must flow when blood has flowed;
that is the main thing; \textit{whose} blood is the subordinate matter. Where the old family organization still prevails, as in
Albania, this way of looking at things still prevails too. The Albanian therefore does not receive a stranger before he has made sure
by inquiry that there is no blood between their kindreds. In the Middle Ages the guild brothers had the obligation to avenge one
another and to pay compensation for one another that members of the same kindred formerly had. ``All shall bear it,'' says an old guild
statute, ``when one transgresses, and all suffer the same.'' % translated from Høffding's Danish
It is due to the advancing emancipation of personalities that the idea of the limitation of guilt to the single individual gradually
develops\footnote[1]{Cf. on this \emph{Spencer}: \textit{Political Institutions.} p. 549 ff. --- \emph{Post}: \textit{Die Grundlagen des
Rechts.} p. 354 ff. --- L. \emph{Brentano}: \textit{Die Arbeitergilden.} I, p. 2. --- R. \emph{Keyser}: \textit{Efterladte Skrifter.} II,
p. 359.}.

Just as little as one originally distinguished between the acting individual and his whole kindred, just as little did one distinguish
between the outer action and the will that showed itself in it. One kept to the effect of the action. The instinct of revenge does not
give itself time to enter into the state of the agent; besides, it also lacks the ability to do so. One keeps to the fact that a
painful effect has proceeded from him, that he is a being that causes pain; anything else about him one does not care to know. --- No
distinction is made between intentional and unintentional injuries, between the harm caused with deliberate intent (dolus) and that due
to indifference and negligence (culpa). Involuntary and willed murder defiled in equal degree. And one punishes --- still in the Middle
Ages and later --- animals as well as human beings!\footnote[2]{\emph{Post} op. cit. p. 356 ff. --- \emph{Fustel de Coulanges}:
\textit{La cité antique.} 4 éd. p. 108.}.

That revenge bred new revenge in return, so that an endless series of actions and reactions between the kindreds
% --- p. 551 ---
\opage{551}was initiated, --- that the whole society suffered from this strife, --- that rawness and cruelty were nourished in minds, ---
did not come into consideration where the unconditionally commanding instinct prevailed.

3. When the recognition arises that the contending individuals and kindreds belong to one society, which suffers from their strife, and
when they themselves at the same time feel the drawbacks of the endless struggles of revenge, the interval mentioned before can set in.
There arises a possibility of mediation, and a deliberation is made before proceeding to the act of retribution. The first courts were
courts of mediation. If the guilty party had been found, and no settlement could be brought about, he was left to the revenge of the
injured party, either without more ado or on certain conditions. If the injured party was not able to revenge himself, society helped
him in it. That the intervention of the power of the state was nevertheless not regarded for the time being as the intervention of a
highest instance is seen from the fact that he who was not satisfied with the judgment had in France (before the legal reform of Louis
the Saint) the right to challenge the judge to single combat: it was an appeal from the judgment of men to ``the judgment of God through
the sword'' (jugement de Dieu par l'épée). --- The next step is that society itself takes the whole of retribution into its hands. It
must then command sufficient physical power both to hold back the instinct of revenge and to carry out the judgment. This power
naturally coincides with that which is raised against outward enemies, so that the authority that maintains law and the defending
authority come to rest in the same hand\footnote[1]{\emph{Maine}: \textit{Early Law and Custom.} p. 170 ff. --- \emph{Spencer}:
\textit{Political Institutions.} p. 618 f. --- \emph{Post}: \textit{Die Grundlagen des Rechts.} p. 401 ff. --- In an interesting way
the transfer of the authority to punish to the power of the state appears in the history of France in the Middle Ages, especially in
the legal reforms of Louis the Saint. \emph{Henri Martin}: \textit{Histoire de France.} 4. éd. IV p. 290---303.}. It was surely not
only reasons of expediency that brought about this passing of the right of retribution from individuals and kindreds to the highest
power of society. The rulers' own interest, their influence and their revenue were increased by their intervening as mediators and
judges in the disputes of single individuals and
% --- p. 552 ---
\opage{552}kindreds, and by the fines, one of the most important means of retribution, going into their coffers. But as so often in
history, what was mainly due to egoistic motives here became the condition of an essential advance in an ethical-social respect. Now a
strife for justice with reasons and testimony could replace the strife with naked weapons. The whole matter was now regarded from a
higher standpoint. It came to be a matter of something more than merely the satisfaction of individual instincts. Society wishes to
preserve its peace and its security, one of its first conditions of life, and to this consideration both the injured and the injuring
party, both avenger and attacker, must bow. There now come penal statutes and penal laws that regulate retribution.

In place of the blind haste with which revenge was carried out as long as instinct prevailed, there now comes a weighing of the
circumstances under which the injury was done, and a measuring of the degree of suffering to be inflicted on the injurer. Even if this
measuring is still done in a rather external way, so that eye for eye and tooth for tooth is demanded, it is nevertheless an advance in
relation to the fury that knew no other limit than its own exhaustion or the annihilation of the victim. There is now at least a
possibility of letting one's gaze reach further than merely to the momentary satisfaction of the collective or individual thirst for
revenge. One can consider whether the suffering inflicted is really necessary, and whether, when it is inflicted in the measure that the
need for retribution demands, it does not bring with it drawbacks greater than the drawback that lies in the need for retribution having
to be checked. One grows accustomed to regarding the whole matter from the point of view of society and to keeping to the significance
that the suffering inflicted on the injurer has for the whole society --- the injurer himself included. One discovers that the stronger
the power of the state is, the milder punishments can be without the peace and security of society suffering from it, that on the other
hand cruel punishments awaken and nourish savagery and rawness in the minds of the people and can thereby themselves become dangerous to
the peace of society, --- and that it depends just as much on the inevitability of punishment
% --- p. 553 ---
\opage{553}as on its severity\footnote[1]{This was experienced in the eighteenth century, when in several countries people tried to
work by ever severer punishments. Of Bavaria, e.g., which distinguished itself especially in this respect, it is said: ``Crimes
increased the more gallows and wheels were set up along the highroads.'' (A. v. \emph{Feuerbach}: \textit{Biogr. Nachlass.} I. p. 432).
% translated from Høffding's Danish
The very cruelty with which punishment was inflicted surely contributed to increasing the rawness. --- \emph{Romilly} justified his
reform of the severe English penal laws by the fact that crimes were on the increase.}. It is acknowledged that because a man is
accused he is not thereby deprived of his human rights, and care is taken that his case can be put in the most favorable light
possible.

The original starting point in the thirst for revenge is thereby definitively left behind. For the thirst for revenge and for the
doctrine of retribution the suffering inflicted on the injurer is the end. Punishment is willed for its own sake, and reflection does not
embrace what lies beyond the moment when retribution has struck its victim. By the checking of the immediate thirst for revenge and
retribution, real justice, which is something other and more than mere ``like for like,'' can come to prevail. The main weight is now
laid not on the instinct that demands satisfaction, but on how the good of society demands that the injurer (who himself continues to
belong to society) be treated. What the primitive instinct wholly disregarded now becomes the guiding point of view. It will appear that
only thereby too is it possible to give an ethical grounding of the state's authority to punish.

4. When the question is raised on what the state's authority to punish is grounded, the relation developed in the foregoing between
society and the state on the one side and the individual on the other will guide our decision. (Cf. III, VIII and XXXVIII).

It is the task of the state to maintain the legal order as one of the fundamental conditions of the development of human society.
Without security and peace neither the single individual, the family nor the free cultural communities can thrive. The state therefore
employs its coercive power when an individual refuses to fulfill the duties that according to the valid legal order are incumbent
% --- p. 554 ---
\opage{554}upon him. And where the individual at some point or other breaks the legal order, the power of the state steps in to maintain
it. This cannot be done merely by stopping and repelling the attack. In the breach of the legal order a will makes itself known that,
if it were allowed to proceed consistently as it has begun in the deed committed, would burst the community of the state. Therefore the
state subjects the transgressor to such a treatment that the relation between his will and the legal order can become a harmonious and
not, as in the action committed, a disharmonious relation. It forces on him an education by which he can become able to keep within
the limits the legal order prescribes. He is trained in self-control, and he is accustomed to work. The state does not thereby act as a
wholly foreign power toward him. It must not for a moment forget that he himself is a member of the race and of society. It does not
treat him as a \textit{mere} means for its ends, does not sacrifice him for the sake of society. What matters is, as far as possible,
to make the interest of society and that of the transgressor coincide, so that social security and order are maintained through an
education that becomes a benefit to the very person who has violated them.

When the state as the upholder of the legal order stands over against an individual who has broken the legal order, it is really two
ethical systems that stand over against one another. \textsc{Grotius} in his famous work rightly presented this relation as a relation
of war, a war between the individual and society, in analogy with the wars between individuals among themselves and between societies
among themselves\footnote[1]{See \textit{Den nyere Filosofis Historie}\textsuperscript{2}. I, p. 53 f.}. Principle stands against
principle, and no negotiation is possible as long as the individual holds fast to the presuppositions that led him to the breach of
law. Where a clash between different principles and systems takes place, practical-psychological influence is, as we have seen
(III, 13), the only means that can lead to an understanding. It is a matter of changing the psychological foundation, such a change as
all education aims at. Before education is completed,
% --- p. 555 ---
\opage{555}the individual naturally cannot acknowledge the justification of the treatment he is subjected to. But the justification of
the state rests on its wishing to make him fit to live in society.

The education the individual is subjected to is first and foremost a political or juridical education to legality. The individual
learns that it fares best by conforming to the law in its actions, whatever thoughts it may otherwise entertain. It is the impulse of
self-preservation that is appealed to first and foremost. With the elementary demands that the legal order makes, it depends, after
all, if it comes to the worst, only on the outer performance. But there is no reason why the state, in its treatment of the individual
who has offended, should always stop at outer, juridical education. This ought to be determinative for the kind and duration of the
punishment. But that does not exclude trying to influence the individual in an ethical respect so as to root out the evil. Only when
the individual's whole disposition and way of thinking have developed so that it bows to the legal order not merely from egoistic
interest but from inner conviction, only then is the legal order completely restored after the breach, and more than restored, since in
place of an enemy it has won a friend. Work will be the most important means of education. Very often the cause of the transgression is
that the individual seeks gain through fraud and encroachment instead of through honest and upright work. But a more direct influence
is not excluded. The ethical education must naturally not be exercised in such a way that religious freedom is thereby violated. It
must be different in form and method for different individuals according to their standpoint. It would be an extra punishment if the
individual were to be compelled to be instructed in a confession other than its own. If the individual belongs to no confession, the
education must be led by a teacher or another personality fitted for serious influence. Where the breach of the legal order springs,
as is very often the case, from bodily and spiritual degradation, something can be accomplished by the individual --- perhaps for the
first time in its life --- coming to stand face to face with people whom it feels it can trust fully. This alone, without much
% --- p. 556 ---
\opage{556}dogmatizing and moralizing, can bring about important ethical changes in the individual's nature. This ethical education, it is
true, belongs rather to the manner in which punishment is carried out than to punishment itself. It can take place only where the
individual shows itself receptive or even feels a need for influence in the direction mentioned. Some have even wished to deny the
state's right to intervene in the individual's ethical development, because no human being has the right to make himself judge of the
inner disposition of others, and the state has no right to use power to improve the citizen who is of age\footnote[1]{\emph{Fichte}:
\textit{Grundlage des Naturrechts.} II, p. 114. --- \emph{Goos}: \textit{Indledning til den danske Strafferet.} p. 63.}. But from the
point of view of social ethics there is no reason to distinguish so sharply between punishment and the carrying out of punishment as may
perhaps be necessary for technical reasons in jurisprudence. He who is sentenced to a punishment is, after all, also sentenced to its
carrying out and to everything it brings with it, and in many cases ethical improvement is either a natural continuation of, or a
necessary condition for, the juridical one. It will prove impossible to maintain a sharp distinction between them. The relation between
the juridical and the ethical element in the unlawful action itself will, after all, also be different in each single case. In some
breaches of law (such as political crimes and police offenses) the ethical element may be as good as lacking, or there may even be one
of those conflicts between morality and law that cannot be avoided. The greatest and most far-reaching breaches of the legal order need
not also be those that are most questionable in an ethical respect. ``The entanglement of circumstances,'' says A. S.
\textsc{Ørsted}\footnote[2]{\textit{Om de første Grundregler for Straffelovgivningen.} (Eunomia II). p. 8.}, ``can often lead a man who is
far from dead to the good and noble into the most dangerous and harmful crimes, while a crime of slight consequences can betray a
disposition corrupt to the core. Who will assert that murder or crimes against the state unconditionally betray greater moral corruption
than fraud?'' % translated from Høffding's Danish
To this question of Ørsted's it must be answered, according to the experience of several experts, that
% --- p. 557 ---
\opage{557}murderers very often stand morally higher than other criminals. (See below in 6). --- One therefore cannot ground the state's right
to punish without demanding that punishment be \textit{individualized} as much as possible. It is especially the constantly differing
relation between the juridical and the ethical element that makes this demand necessary. Without individualization the individual
punished is treated as a mere means for society's need to maintain its legal order. But with the individualizing education that we
demand, there is no point at which the activity by which the legal order is maintained is not at the same time a pedagogical activity
toward the individual who has transgressed it. The penal activity of the state acquires the greater ethical grounding and justification
the more it approaches this ideal. The state's right to subject the transgressor to an education that makes him fit to satisfy the
elementary conditions of life together with others in an ordered society rests essentially on the same foundation as its right to see
to the upbringing and instruction of children. In both places the state maintains the level below which its members must not sink, and
seeks to help up those who have not yet reached it, or who have sunk below it\footnote[1]{The development given here of the ideal
grounding of the authority to punish attaches most closely to a single side of \emph{Fichte's} theory. Fichte's theory as a whole is a
theory of penal threat (like A. Feuerbach's and A. S. Ørsted's) and untenable, since it must consistently make the individual punished
a mere means for society. But in the closer development of his theory he comes upon thoughts that could be used as the foundation for
an entirely different theory. For the rest, the pedagogical theory of punishment is of very old origin. It stems from \emph{Plato} as a
necessary consequence of this great thinker's break with the old doctrine of retribution. (Cf. XII, 2).}.

The grounding of the state's authority to punish that builds on \textit{the pedagogical effect of punishment} has its great significance
in setting up the ideal that the ordering of the penal system is to strive to approach. It is no objection to it that this ideal is not
reached in reality, indeed that the so-called houses of correction very often are institutions of deterioration\footnote[2]{\emph{Garofalo}:
\textit{La criminologie.} Paris 1888. p. 98; 129 ff.}. This only shows
% --- p. 558 ---
\opage{558}that an ideal is needed that can serve as a standard and show the direction in which work is to be done. And against the
consideration on which that grounding rests, namely that the single individual is and remains a member of society even if it violates
society's laws, just as sick individuals and those unable to work belong to society, no objections can be raised. The individual that
breaks the legal order has itself come into being and developed within society. Its character and whole way of thinking are to a great
extent determined by the spirit that prevailed in society and by the conditions society had arranged for it. In the breach of law we can
often see a simple effect of defects and faults in society. In judging the transgressor, society to that extent judges itself. When
society must thus see in the breach of law, to a greater or lesser extent, its own work, its own guilt, the right to educate the
transgressor changes into an obligation, and the necessity of the individualizing mode of treatment becomes all the more evident. Only a
view that regards the individual as the absolute starting point of all its actions could maintain that by transgressing the legal order
it separates itself altogether from society, so that the latter no longer has any obligation toward it. It is not, however, only
indeterministic views that can arrive at this result. The so-called Italian school in criminology (\textsc{Lombroso} and his school),
which explains habitual criminals as atavistic revivals of a human type that civilized societies have essentially got beyond, likewise
thinks it can maintain that society has no complicity in or obligation toward criminals\footnote[1]{Cf. \emph{Garofalo}: \textit{La
criminologie.} p. 55 ff., 117 ff.}. It is clear, however, that it is already a sign of imperfection in society that hereditary
dispositions constantly maintain themselves and only wait for a good opportunity to pass over into actuality. It shows that the social
spirit is not really able to permeate the nature of single individuals, and that upbringing and social conditions are not as they ought
to be\footnote[2]{See my criticism of the doctrine of the Italian school in \textit{Etiske Undersøgelser.} p. 73---81.}.

% --- p. 559 ---
\opage{559}The more one's eyes are opened to the influence of social conditions on the character and conduct of individuals, the less
one will expect great effects from the penal system alone. This must be seen as a link in the whole great process of education that the
race undergoes through the institutions of society. As in ordinary upbringing of children, so also in social education an indirect and
preventive influence will more and more take the place of the direct intervention by which one seeks in punishment to remedy the harm
that has been done. Already \textsc{Thomas More} in his \textit{Utopia} brought against society the charge that it first produced
thieves and afterward punished them. As in medicine, so too in the doctrine of law hygiene and prophylaxis will surely take the place of
the curing that can only be applied when the misfortune has happened, and is then often applied in vain.

--- The pedagogical theory of punishment gives the most ideal grounding of the state's authority to punish. But it is not sufficient.
When it lays weight solely and exclusively on the fact that the single individual grows out of society and always remains a member of
it, it overlooks that he who transgresses society's law acts as society's enemy and sets an example that may easily find imitation. Crime
establishes a state of war between individual and society and thereby sets society another task than the merely educating one, namely
that of maintaining the legal order and preventing the single breach from becoming a precedent. Therefore the relation of opposition
between society and the single individual must also be emphasized. This consideration is made the basis in \textit{the theories of
deterrence and of penal threat.} Here the end of punishment is seen in its working as a deterrent, since an example is made, as when a
kite is nailed to the barn door, or the matter is conceived in such a way that the power of the state in its penal laws utters threats
against those who perform certain actions, and that the punishment that falls on the single criminal stands as testimony that the
threats are seriously meant. The right to deter and threaten is derived from the state's right to exist, from its significance for
human welfare.

The theory of deterrence, which is professed especially by jurists and
% --- p. 560 ---
\opage{560}statesmen (among philosophers by \textsc{Schopenhauer})\footnote[1]{\textit{Die Welt als Wille und Vorstellung.} II. Kap. 47. --- Cf.
in Danish literature A. S. \emph{Ørsted}: \textit{Om de første Grundregler for Straffelovgivningen.} (Eunomia II). 1817. --- In most
recent times F. \emph{Tønnies} has spoken for the theory of penal threat. \textit{(Probleme des Verbrechens und der Strafe ---} in the
journal ``Deutschland.'' Oct.---Nov.\ 1902).}, stands in an interesting relation to the theory of retribution and the theory of
education. It regards punishment not as an end but as a means. Thereby it parts from the doctrine of retribution and acquires, like the
pedagogical theory, a teleological character. But whereas the pedagogical theory regards the very individual punished as an end, the
theory of deterrence regards it as a means, since it is punished in order to inspire terror in others or to show that the penal threats
are really carried out. The end of punishment lies for it in the effect on the rest of society, especially in the security obtained. It
does not lay so much weight on how punishment works on those punished as on how it works on those not punished. What matters is that
the hope of escaping punishment be struck down. And above all it is a matter of preventing the first, dangerous step on the path of
crime. This theory resembles the theory of retribution in not presupposing any interest in how the individual fares during and after
the carrying out of the punishment. Consistently, it will not lay any weight either on the degree of consciousness with which the
violation of the law took place; the decisive point of view becomes the dangerousness of the action, which makes it necessary through
deterrence to put a bar to all attempts in a similar direction. As regards its effects in experience, it has in common
with the theory of education that it is not completely confirmed by experience. Although it is easier to call forth terror than to bring
about a real change of character, experience nevertheless shows that crimes are very often committed by individuals who have been
punished several times, and the criminological problem lies precisely in the great number of relapsers (recidivists). Not even the
punishment that might especially be supposed to work as a deterrent on others, the death penalty, has this effect according to
experience.

% --- p. 561 ---
\opage{561}But both the theory of education and that of deterrence state ends that must necessarily be striven for with regard to those
who violate the law. In a complete theory of punishment both must be united. Punishment will then at once work to change the character
of the one punished and stand as an example that the legal order cannot be violated. The individual punished will then stand at once as
end and as means. To carry through such a conception there is admittedly required an art and a knowledge of human nature that are not
yet available. Only in recent years has the study of prisoners and of the prison system on a scientific basis begun. But \textit{the
decisive standard of the perfection or imperfection of the penal system will nevertheless have to be taken from the degree to which it
has succeeded in uniting education and deterrence.} This standard is only a special application of the principle of free personality
(VIII, 6), and the theory of punishment is thus connected with the fundamental thoughts of ethics. \textit{A philosophical theory of
punishment can give nothing other and more than such a standard, agreeing with the fundamental ethical thoughts, for use in the
valuation of the penal system that actually exists.} Experience also shows that it is more and more the principle of the close union of
education and deterrence that determines the judgment on the penal system as it appears at a given place and a given time. Now education,
now deterrence gains the upper hand as the guiding point of view\footnote[1]{Cf. \emph{Francis Hagerup}: \textit{Et Blad af Straffens
Historie.} (Nordisk Tidsskrift, udg. af den Letterstedtske Forening. 1893).}. In opposition to the theories of retribution and deterrence
that had prevailed earlier, the philanthropists of the eighteenth century asserted the pedagogical theory. They demanded a humane
treatment of those who were to be punished, and set up their improvement as the end of punishment. Even if this philanthropic movement
often had a sentimental character and often led to laxity and untimely leniency in the penal system, it undoubtedly did a good work by
asserting the right of the individual punished. It appealed to a higher concept of justice than the earlier modes of view knew, a concept
of justice that it admittedly
% --- p. 562 ---
\opage{562}did not succeed in setting forth in clear and definite forms. In the most recent time there is a tendency to bring forward the
point of view of deterrence again; but since the philanthropic movement had led to putting penalties of deprivation of liberty in place
of the formerly so frequent corporal punishments, and since penalties of deprivation of liberty, especially the short ones, have proved not
to work sufficiently as deterrents, one stands here before great practical and theoretical difficulties. At the same time criminal-anthropological
 investigations have begun to give an insight into the nature of criminals and have in particular shown that there are
great differences between the individuals who transgress the laws of the state, so that the treatment that befalls them must also be
different and cannot be determined merely by the transgression itself as an outer fact. The theories of retribution, of deterrence and of
improvement all overlooked individual differences and were inclined to regard all human beings as alike. In the most recent time attempts
have been made to take account of individual differences by introducing conditional sentences, i.e., such as are carried out only if the
individual offends once more, and indeterminate sentences, i.e., such in which the length of the term of punishment depends on how the
individual's character expresses itself and develops during the carrying out of the punishment. One has also become more and more clear
that crime is a social phenomenon connected with many other social phenomena. The criminological problem does not stand isolated, but is
closely connected with the other social problems. One thereby comes to the already mentioned conception of punishment as a means of
combating crime regarded as a social evil. The ``International Criminalistic Association'' has made this conception its basis, and has
also insisted that punishment must not be torn loose from its connection with the other means of combating, especially of preventing,
crime. The indefinite expression \textit{combating} has thus been put in place of the more definite expressions of the older theories,
\textit{improvement, deterrence, retribution.} Criminal policy takes the place of the theory of punishment. It will be a matter for
future experience to decide how much can be reached along the preventive path, especially by improving education and
% --- p. 563 ---
\opage{563}economic conditions, and how much can be reached only through the infliction of penal suffering. Greater demands will surely
come to be made on society when account is to be taken of all the varieties of violators of the law than when one contented oneself ---
as soon as a violation of the law was thought to be established --- with putting the person concerned in the lockup or on the galleys or
cutting off his head. Yet the standard in valuing the way in which criminals are treated will always be the degree to which the single
individual comes to stand as an end and not merely as a means. In one form or another the problem that came forward in the struggle
between the theories of retribution, of deterrence and of improvement will always appear anew.

When punishment is conceived as one means beside several others for combating crime as a social phenomenon, it may become a question
whether one then still takes the word punishment in the same meaning as that in which it occurred in the older conceptions. (Cf. above
V, 5 f.). An acute thinker has denied this. ``The concept of punishment,'' says F. \textsc{Tönnies}, ``must be transformed into a higher
concept, namely into the concept of an expedient treatment of the criminal.'' % translated from Høffding's Danish
To such an expedient treatment suffering need not in itself belong. When the individual has been compelled, as far as possible, to make
good the harm it has caused others, its further treatment must consistently be of a healing and educating kind\footnote[1]{F.
\emph{Tönnies}: \textit{Die Verhütung des Verbrechens.} (Deutsche Worte. Wien 1891). --- Tönnies describes the theory set forth in the
first edition of this work as a pure theory of improvement. It was, however, my meaning (as I have sought to bring out more clearly in the
later editions) that deterrence too stands as a necessary end, and that the standard for the perfection of a given penal system lies in
\textit{the degree to which the two ends are united,} especially united in such a way that deterrence does not spread at the expense of
education. (See 1st ed. p. 379).}.

5. According to the view developed here, punishment is inflicted not only because a transgression has taken place, but also so that no
transgression shall take place (not only: qvia peccatum
% --- p. 564 ---
\opage{564}est, but also: ne peccetur). It is the future, not the past, that makes punishment necessary. That no transgression takes place
means, after all, that the legal order stands firm, and that is what is to be attained. It is thus here a clear and definite end that is
attained, or at any rate can be attained, by punishment, --- if not by punishment alone, then by it in connection with other means
against the social evil that crime is. And this end in turn finds its grounding in the general principle of welfare.

By the assumption that punishment must serve an end and does not have its validity ``in itself,'' the view developed stands in definite
opposition to the theory of retribution, according to which it is supposed to be a degradation both of the authority to punish and of the
personality punished when punishment is regarded as a means. And since the theory of retribution still has so great an influence on
common ideas of crime and punishment that it must perhaps even be said to be one of the most essential obstacles to a rational ordering
of the penal system, we must examine this theory somewhat more closely and test its justification. The theory of retribution maintains
it as an eternal, self-evident ethical truth that whoever does evil, upon his head evil shall come. But is it really self-evident? And is
it an ethical thought? Why should evil, why should suffering in the world be increased because evil has once been done? That is, after
all, simply to make evil worse, as long as one does not point out a definite good that is attained by the evil one lets follow upon the
evil once done. But such an end one will not be able to state without appealing to the principle of welfare and thus getting into the
``morality of utility.'' Kant regarded it as a ``categorical imperative'' that a man is to be treated as he himself treats others. But
that is high-handed language without grounding. On the other hand it follows with great clarity from the principle of welfare that the
infliction of suffering is not permissible, no matter against whom it is directed, when a correspondingly greater good is not thereby
attained for the individual himself or for society or for both. Evil must not be repaid with evil when this latter evil is not a good in
% --- p. 565 ---
\opage{565}disguise. And it follows likewise from the principle of welfare that no one may be treated as a mere means --- not even, or
least of all, as a means for the satisfaction of blind instincts or the application of abstract maxims. Retribution is either an outcome
of the instinct of revenge or an abstract principle of reason that is asserted to be self-evident.

``But'' --- it will be said --- ``it was also a mistake of Kant's to regard that proposition as a pure proposition of reason. It first
gets its significance when it is regarded as the expression of society's ethical feeling, which is outraged by the crime and is first
satisfied by the criminal's suffering. The criminal must atone for his guilt; only after that can he again count as a member of
society.'' --- This does not make matters better. For by what right is the individual to suffer in order to satisfy the feeling of
others? One must establish the value and justification of this feeling before one can demand a satisfaction of it that falls upon a
human being. People's feelings can be set in motion by many different things, and one must therefore establish that this particular
feeling has the right to have its demand fulfilled. If one tries to establish its justification, one will once more end up in the
principle of welfare. Public indignation and outrage are justified only when they are motivated by a violation of the conditions of life
of society and demand what is requisite for maintaining these. It is with the suffering that the state inflicts in punishment as with that
which is felt in repentance. This too was not (see V, 3) of ethical value ``in itself,'' but only by being a necessary psychological
transition.

If the outraged feeling is really of a sympathetic nature and not a flaring up of the primitive instinct of revenge, it will no doubt
demand an atonement, a definite suffering that is to come upon the transgressor for his deed; but it will not stop at this as the last
goal. If it does so, only the need to get relief is at work, without regard to the consequences, the need that so often transforms the
sympathetic feelings into egoism\footnote[1]{\textit{Psykologi} VI C, 7.}. True
% --- p. 566 ---
\opage{566}sympathy is not concerned merely with getting relief; it will increase pleasure in the world and diminish suffering; it would
then be a self-contradiction if at any point it found its goal reached by the infliction of suffering. --- The transgressor himself, too,
often wishes for atonement. This need contains nothing mysterious, and no proof of the doctrine of retribution either. It is
psychologically intelligible that he in whom repentance begins to work feels a relief in physical suffering. To submit to it is a
discharge, is a way of getting relief for the strong spiritual tension. By submitting to suffering he shows, next, his earnest will to
restore the normal relation to the legal order; by suffering the punishment it prescribes he already satisfies, after all, one of its
demands. It is with this as with the painful self-knowledge that the individual who has really broken with a reprehensible direction of
will willingly goes through. (See V, 3). A murderer can, as \textsc{Fichte} has remarked, feel it to be just that he is killed, even if
it happens by the hand of another criminal or by a power that knows nothing of his crime.

Behind the ``reason'' and the ``feeling'' that are appealed to in favor of the self-evident character of the thought of retribution
lies the old instinct of revenge, deeply rooted in human nature. When we go back to it, we understand well that the grounding is
lacking. For instinct is not guided by the thought of any end, but consists merely in the need to perform certain actions. And since in
this case the actions have been, and in part still are, of great significance for the self-preservation of the individual and the race,
we understand that the instinct maintains itself and is difficult to overcome. It was therefore so remarkable an advance in the ethical
development of mankind when people began to doubt the self-evidence of the proposition that evil is to be repaid with evil (XII, 2). ---
It is most strange that even \textsc{Kant} could think he had an eternal ethical truth in the principle of retribution, and that
\textsc{Fichte} and A. S. \textsc{Ørsted}, who rejected this principle in the doctrine of law, could hold to it in
ethics!\footnote[1]{\emph{Fichte}: \textit{Grundlage des Naturrechts.} II, p. 128. --- A. S. \emph{Ørsted}: \textit{Om de første
Grundregler for Straffelovgivningen.} (Eunomia II). p. 4. --- In an interesting essay \textit{Vergeltung und Zurechnung} (in vols. 5
and 6 of ``Vierteljahrsschrift für wissenschaftliche Philosophie'') E. \emph{Laas} has sought to maintain the ethical significance of
retribution and to conceive punishment as ``ethicized revenge.'' But I cannot find that he puts forward anything that answers the doubt
put forward above, and it seems to me unjustified to use the word revenge when every egoistic need is to be thought away. (See
especially vol. 5, p. 152 f.). In a similar direction to Laas two thinkers so different from each other as \emph{Lotze}
(\textit{Grundzüge der praktischen Philosophie} § 64---65) and \emph{Dühring} (\textit{Cursus der Philosophie.} p. 226) express
themselves. According to \emph{Westermarck} (Origin and development of moral ideas. I. p. 122) revenge and punishment are not mother
and daughter, but stem from a common mother, the feeling of social resentment (public resentment). Perhaps this resentment in turn stems
from the feeling of defilement that Frazer takes as basis. --- From the fact that punishment has its historical origin in kin revenge
and private revenge it does not follow that its ethical justification is grounded on the instinct of revenge. (Cf. I, 4; XIII, 4).}.

% --- p. 567 ---
\opage{567}The energy that shows itself in the instinct of revenge and retribution need not be lost because the instinct must undergo an
essential metamorphosis. Strong and living feelings of indignation and resentment are constantly needed in order to maintain the
violated legal order and to enter the lists for it. It is only what is purposeless in these feelings that we speak against.
Because the infliction of suffering is not the last goal, there can very well be room and use for a living sense of justice. By
``nemesis'' \textsc{Aristotle} understood the feeling we have when we see good fortune fall to the lot of the unworthy\footnote[1]{\textit{Eth.
Nic.} II, 7.}. Such indignation at a disproportion between inner and outer happiness is an important ethical force. But this indignation
need not stop at the wish for suffering for those who have done evil; it is not weakened by this suffering being regarded as a means to
a change of character in the doer.

Where the need for retribution does not directly work in the service of self-preservation, its satisfaction rests on an illusion. What
does one really attain by retribution? The evil that has been done cannot be undone, and how can it be a relief to us that the
perpetrator comes to suffer? The illusion here is similar to that in envy and malicious joy. One has attained that one has got relief;
but nothing whatever in
% --- p. 568 ---
\opage{568}the order of things has been changed, except insofar as the sum of suffering in the world has been increased. This illusion in
revenge \textsc{Bjørnson} has depicted in ``Bergljot'':

\begin{quote}
Revenge? --- Who speaks of revenge?\\
Can revenge awaken my dead\\
or cover me against the cold?\\
Is there in it a widow's seat shut fast\\
or comfort for a childless mother?
\end{quote}
% translated from Høffding's Danish (Bjørnson's poem; a plain rendering)

6. It might seem as if the principle of retribution were excellently suited to be applied in measuring out punishment. It seems to be
a clear and reasonable thought that one is to be treated just as one treats others. One finds it self-evident that he who has killed
another is himself to be killed. But one will already (at our present stage of ethical development) become somewhat hesitant about
putting out the eye of him who has put out the eye of another. And now he who has committed rape, how is he to be punished according to
the principle of retribution? Kant thought by castration. But that is already a strong deviation from the principle of like for like.
This cannot be carried through, and where it seems to be applicable, it is rather a certain symbolism that comes into force than a real,
reasonable connection between deed and punishment. --- If one will carry through the principle of retribution consistently, one must go
as far as some savage peoples. A Basuto, e.g., whose son had been wounded in the head by a blow with a stick, wanted to get hold of the
perpetrator in order to strike him with the same stick on the same place on the head, standing on the same spot of ground on which the
perpetrator had stood. Only then did he feel himself completely satisfied.

The doctrine of retribution cannot ground the fact that actions springing from intent (dolus) are punished more severely than those
springing from negligence (culpa). It is the action that is to be repaid; how can the retribution differ according as the action has
been done with will or through negligence? It would surely be absurd to expose the man who through negligence had caused another's
death to a similar
% --- p. 569 ---
\opage{569}danger to the one he had brought the other into! One would then have to see to it that it too brought death with it. Yet the
retribution would not be complete, since death would precisely not be caused through negligence\footnote[1]{Cf. A. S. \emph{Ørsted's}
apt criticism of retribution as the standard of penal law: \textit{Om de første Grundregler} (Eunomia II) p. 12---24.}. But if one cannot
take account of the difference between intentional and negligent action, it becomes clear that retribution is in reality the outcome of
a blind instinct that makes us strike out without investigating the circumstances. That is a strange fate for so highly strung an
ethical view as the doctrine of retribution, which likes to look down with contempt on the ``morality of utility'' and ``humanity.''

The theory of retribution will finally not be able to ground what most recent penal laws lay down about the limitation of criminal
liability. According to the grounding of the authority to punish given above, it is self-evident that breaches of law lying far back in
time and only now discovered are not punished. The motive that sets the power to punish in motion is here lacking. Just as ethics looks
forward and looks back only in order to be able to look the better forward (V, 3 f.), so the state has to do not with the past but with
the present and the future; the action committed far back in time is neither a sure mark of what the perpetrator's character is now
like, nor does it contain any danger to the legal order now existing. The practical interest has fallen away. But for the doctrine of
retribution there must here be a hole in the order of the world, and no lapse of time however long can ground the punishment's not being
inflicted. Why should the blood cease to cry out?

--- According to the grounding of the right to punish given above, there is a twofold end to be reached by punishment: the restoration
of the legal order and the change of the perpetrator's character. In no mode of punishment may these two considerations be separated. Of
the punishments employed nowadays there are two that conflict with this, namely the death penalty and imprisonment for life. They cannot
therefore be justified, and we find too that they are used more and more seldom. Only a fifth of
% --- p. 570 ---
\opage{570}the death sentences pronounced in Europe are carried out. These punishments are an expression of the fact that we still stand at
a barbarous stage of development. We live under a state of war that makes the factor of deterrence come to play a greater role than can be
defended ethically. But the reason is also especially that we are still so extraordinarily far behind in the psychology and pedagogy of
punishment. Therefore we are not yet able to work on the transgressor's character in such a way that punishment can be concluded without
ending in death. To let the one who is to be educated ``sit out'' for life is a strange pedagogical method, and no less strange is it to
take his life in order to get the better of his character. It is easy enough to declare a man incorrigible; but where does one get the
proof that \textit{all} means have been exhausted? We are not yet so advanced in psychological insight and pedagogical practice that we
dare to issue such a declaration of outlawry. A number of experiences are on record showing that murderers condemned to death, but then
pardoned, have led a blameless life afterward. In some cases this was perhaps because the bad example and the outer opportunity were kept
away; but in other cases it was surely because a change of character had taken place\footnote[1]{\emph{Holtzendorff}: \textit{Das
Verbrechen des Mordes.} Berlin 1875. p. 177 f. --- \emph{Prosper Despine}: \textit{Psychologie naturelle.} Paris 1868. II. p. 260.}.

After the investigation made earlier (V) into ``the freedom of the will,'' there is no reason to go into this question here. If ethics can
be reconciled with determinism, the doctrine of law will be able to be so too. But for all that it will nevertheless --- precisely
according to the grounding of the right to punish given here --- be of great importance that the subjective conditions of will in those on
whom punishment is to be inflicted become an object of attention. --- It depends in the first place on how far the action is really
willed, how far it is the fruit of the individual's clear deliberation and conscious resolve, or has only been the object of intent, not
of a resolve proper\footnote[2]{\textit{Psykologi} VII B, 1.}, or has only been caused by the individual through lack of attention and
reflection,
% --- p. 571 ---
\opage{571}or finally has been carried out in a state of unusual excitement. In these different cases there is a very different degree of
disharmony between the individual's will and the legal order, and the punishment ought necessarily to vary accordingly. Yet too great
weight has often been laid on the degree of consciousness and deliberation with which the action was carried out. A long and conscious
deliberation may indicate that there have been great obstacles in the individual's mind to overcome before the criminal resolve could
come into being. In such cases deliberation is then an extenuating circumstance. A character that can perform an action only after
overcoming inner resistance stands above the character that feels no misgiving at all about such an action. --- In the second place,
punishment is, after all, to influence the individual's will as a motive that has proved to be requisite. The individual is to be
educated. But it then depends on whether the individual is normal, so that the influences one brings it under can have any effect. If
the individual is decidedly insane, educating influence in the usual form will be useless or even harmful; it must then be referred to an
asylum for the insane. If the individual is not insane in the usual sense, but (even if the action that brings it before the court is
comparatively insignificant) reveals strong and dangerous tendencies in its nature that it will probably not be able to overcome by the
capacity for self-control at its disposal, it must be referred to a ``prison asylum,'' an intermediate form between prison and asylum for
the insane\footnote[1]{Cf. A. \emph{Forel}: \textit{Du traitement des causes pathologiques du crime} (Bulletin de l'union international de
droit pénal. Vol. X, p. 390).}. And the more the nature of such individuals is studied, the more such intermediate forms will be required.

% --- p. 572 ---
\opage{572}\chapterhead{XL.}{THE CONSTITUTION OF THE STATE.}

1. When in the year 1789 the French were engaged in making themselves an entirely new constitution, an American statesman,
\textsc{Morris}, who was staying in Paris, wrote that they wanted an American constitution without considering that they had no American
citizens, and that at least a generation would be needed before people could become used to the new constitution and really live under
it. At the same time Wilhelm v. \textsc{Humboldt} predicted that the new constitution would not last, because it was a work of reason.
Reason --- he taught --- can, no doubt, order a given material, but cannot itself produce it. It can guide and awaken by setting up
models, but the living forces develop slowly through historical experiences\footnote[1]{\emph{Taine}: \textit{La Révolution.} I, p. 158.
183. --- W. v. \emph{Humboldt}: \textit{Ideen über Staatsverfassung, durch die neue französische Constitution veranlasst.} (1791). (Werke
I. p. 301 ff.).}. --- What a practical statesman and a philosophical thinker thus agreed on, history has sufficiently confirmed. Forms of
state and constitutions cannot be made all at once and impressed on the people from outside or from above. Since the beginning of the
nineteenth century 350 new constitutions are said to have been made\footnote[2]{\emph{Maine}: \textit{Popular government.} London 1885.
p. 174.}. The number alone shows that very many of them must have disappeared again because they were unable to strike root.
Constitutions must grow up. For us here this is no new truth. The form of state or the constitution is, after all, only the part of the
legal order that concerns the government of the state, and it holds of all law that only when it proceeds from or leads to customary
activity does it acquire a firm existence as the form of a real life. Everything is not done by formulating principles and schemes. As a
comparison between the history of the continental peoples and that of the English and North American peoples shows,
% --- p. 573 ---
\opage{573}constitutional life develops more soundly and vigorously where the constitution comes into being through a long series of small
expedients or improvements applied to single given cases than when a constitution is constructed and introduced all at once. That the
North American federal constitution was the conscious work of single men and was introduced all at once does not conflict with this. For
what was gifted in this constitution was precisely the way in which it made use of the possibilities lying in the conditions and
institutions of the former English colonial states\footnote[1]{Cf. my essay \textit{Alexander Hamilton og den nordamerikanske
Unionsforfatning.} (Mindre Arbejder).}. This constitution has accordingly had the good fortune, unique for an ``introduced''
constitution, that its centenary could be celebrated with good prospects of continued existence. --- Reason, conscious deliberation,
does not therefore lose its significance, as Humboldt too indicates. It works back in many ways, ordering, guiding and awakening, on what
is wholly or half unconscious. It sharpens attention to the essential conditions, warns against the dangers that threaten, and discovers
new possibilities of development. (Cf. III, 19---20).

From this it follows that it cannot be the task of ethics to construct an ideal constitution. It is no use making the frame when one
does not know whether there is anything to fill it with. What is decisive is the possibilities and starting points that the people
presents according to its nature and its historical development. Long before one can begin to discuss constitutions, an inner
constitution has in fact developed, which shows itself in the people's usage and custom, in its character and its culture, and to this
the outer, official ordering of the relation between the various elements of the people must attach itself if it is to have existence
and significance. The frame is naturally not indifferent. If it is too narrow, it presses on the content and hinders its unfolding; if it
is too wide, the content is not ordered in a firm form. And a bad frame need not only be one that is arbitrarily put on from outside; it
can also have developed involuntarily. But one can remedy
% --- p. 574 ---
\opage{574}such defects only by relying on possibilities and starting points that are given in the people as it now is. All sound
development goes from within outward, from below upward.

The proposition that constitutions cannot be arbitrarily chosen and introduced does not, however, leave us altogether powerless before
political life. We may perhaps be unable to accomplish anything against the finished results of development, and we cannot all at once
set going something that requires a long time of development to thrive. But we can act upon the causes that determine future
development. We can learn from experience and after mature deliberation arrange conditions better than they have lain hitherto. It is
only for the earliest periods of social and political development that it holds that outer conditions determine the ordering of society
and state without conscious insight and choice playing any role whatever. Where we find expedient arrangements existing, we have not
always the right to assume that their origin is due to wise foresight and calculation. But it does not follow from this that one can
(with \textsc{Spencer}) set up as a general rule that it is the \textit{given conditions} and not \textit{the conscious intentions} that
determine the form of state\footnote[1]{Conditions and not intentions determine. \emph{Spencer}: \textit{Political Institutions.} p.
395.}. When a higher degree of spiritual development has been reached, there arises, after all, precisely insight into the given
conditions, and this insight will then be able to become a guide for our action. When historical experience and social and political
insight have developed, it would conflict with all psychology if they did not also to a certain degree determine our practical conduct.
None of our ideas is altogether without influence on our feeling and will. The \textsc{Spencerian} proposition marks the high point of the
reaction against the enthusiastic but naive confidence of the eighteenth century in reason. The reaction begins with no longer believing
that it is ``ideas'' that rule the world: not thoughts but feelings and passions determine human actions, and thoughts are themselves only
by-products of feelings. The next step is that these feelings and passions in their turn find their explanation in the outer
% --- p. 575 ---
\opage{575}conditions under which human beings live and act. Thought and conscious intention are regarded only as products of a whole
causal series whose first link is sought in the outer conditions of life. But once thought has been produced, it can nevertheless work
back on the conditions to which it owes its coming into being. What is effect is itself in turn cause. That is what the Spencerian theory
overlooks. --- When insight is sufficiently thorough, it will also bring with it an understanding of how extraordinarily complex a human
social life is, and how many-sided and ramified the effects that any intervention at a single point can therefore have. It will
therefore lead to caution and to the conviction that indirect intervention reaches the goal sooner than direct. It will seek to serve life,
not to lord it over it.

2. What always brings an irrational element even into the most rationally devised political order is the role that \textit{power}
must necessarily play. That all use of power is necessary to preserve security and peace outwardly and inwardly is a testimony that
impulses and passions are still at work in human nature that are not adapted to the demands of a society built on a firm legal order. It
is to that extent necessary, to use Taine's expression, to have the gendarme on guard against the savage. But the command of power and
the exercise of power also appeal to the impulses and passions of human nature. All egoism is really a kind of feeling of power. The
brutal side of man's nature can express itself just as well in the raw use of power as in defiance of the order power is to maintain.
The gendarme himself has something of the savage in him. Men must perhaps be governed with tight reins; but those who hold the reins are
also men.

``If men were angels,'' said \textsc{Alexander Hamilton}, ``no government would be necessary. If angels were to govern men, neither
external nor internal controls on government would be necessary''\footnote[1]{\textit{The Federalist.} (1787). Boston 1882. p. 399.}.
% Hamilton's English (Federalist No. 51)
The problem to be solved in the constitution of the state is precisely this: a control supported by power is needed to govern the
impulses in human nature that
% --- p. 576 ---
\opage{576}threaten to burst the legal order; but next there is also needed a control over those who exercise that control; and how can
this latter control obtain the requisite authority without weakening the power that those who exercise the former control must
necessarily command?

3. At lower stages of political development it is in reality often assumed that it is angels who govern. People look with blind
confidence to those who govern as the epitome of all wisdom. So in the theocratic state, and in every state where the government stands
as the guardian of the people. Here only the first kind of control then seems to exist, not the second. Yet on closer consideration this
too can be pointed out. It consists in the people's very belief in the wisdom and ability of the government. Only as far as this belief
reaches will the government, even the most patriarchal, be able to go safely. To this are added national feeling and whatever other
feelings the government may be able to awaken. Mere power will not in the long run be sufficient where inner adherence is lacking. Those
who govern become, especially where active and concerted conduct is in question, in reality always dependent on the mood and way of
thinking of the people, which as a rule they themselves to a great extent share. Moreover the government cannot immediately see to
everything; it must have intermediaries, employ a bureaucratic system, which it can control only imperfectly and on which it easily
becomes dependent\footnote[1]{The Danish philosophers of law in the middle of the eighteenth century even conceived absolutism as a free
constitution, because it relied on public opinion and the official class. See E. \emph{Holm}: \textit{Om det Syn paa Kongemagt, Folk og
borgerlig Frihed, der udviklede sig i den dansk-norske Stat i Midten af 18. Aarhundrede.} (University programme). 1883. p. 66. 85. 95.}.

Public opinion, which already plays an important role under the regime of guardianship, will, where the development of the people is
sound and vigorous, pass over into becoming \textit{a public will} that demands active participation in the decision of affairs (Cf.
XXXVII, 5 and XXXVIII, 5). The presupposition for this is that the people has become conscious of its strength through activity within
free society and in the smaller circles. Political freedom in modern
% --- p. 577 ---
\opage{577}Europe is a result of processes of development that have taken place in the domain of material and ideal culture. Under the
luxuriant flourishing of trade, industry and agriculture there unfolded an intelligence and a capacity for action that could not be
kept within such narrow limits as the old forms of state laid down. Through the free development of religious life --- especially among
the Dissenters --- self-esteem and the feeling of freedom rose to a degree that made it impossible to exclude the people from influence on
the decision of matters concerning its own weal and woe. The English sects have acquired world-historical significance through their
influence on English-American political life. Modern science brought about not only by its results an entirely new view of nature and
history, but called forth, especially by its method and by the desire for inquiry and criticism that it awakened, a free and bold way of
looking at things that had to extend to the political domain as well. People could no longer put up with being controlled without
controlling. Political freedom is an expression of the fact that human life must be guided by the same laws in all domains; he who would
get the need for political freedom out of the world root and branch must stop all independent work in the domain of material and ideal
culture.

This is not to be understood as if the people could first be educated to political freedom and then be made a partaker in it. As long as
it does not have freedom, it cannot acquire the abilities and feelings that active participation in political life presupposes. Here
again the Aristotelian principle holds, for only \textit{by having political freedom can a people become politically free.} Freedom is here
too both means and end. (Cf. XXIII, 2). A regime of compulsion and guardianship cannot educate to freedom. Only through practice and
habit are abilities developed, and practice in using freedom one cannot get when one does not have freedom. --- This consideration does
not, to be sure, remove all difficulty; it moves it further back, since it must now be a matter of when a people has come so far in
development that political freedom can become a means for its further development. But not a little has been gained by this changed
% --- p. 578 ---
\opage{578}statement of the problem, since it can more easily be decided when a means is applicable than when a perfect proficiency is
present.

Nor is it to be understood as if political freedom were merely a right that the people acquired by its proficiency in other domains. Its
possession and exercise can become a \textit{duty} no less than a right. The state first really comes into being when as many as possible
take active part in political life. The state, after all, has no existence outside or above single individuals, but exists in them, in
their minds and in their will. The more cooperating wills, the firmer the foundation for the edifice of the state. The free and living
adherence of single individuals is the true strength of the state. Political freedom binds single individuals to the state through their
own activity for the purposes of the state, and it is the duty of individuals to show themselves active members of the state. Great
statesmen have therefore regarded universal suffrage as a necessary condition for the state's being able to be built on as deep and firm a
foundation as possible\footnote[1]{\emph{Alexander Hamilton}: \textit{The Federalist.} (1787). Boston 1882. p. LIX. Cf. my \textit{Mindre
Arbejder.} First Series. p. 128. --- \emph{Bismarck}: Cf. \textit{Fürst Bismarck als Redner. Vollständige Sammlung der parlamentarischen
Reden Bismarcks.} III. p. 194 ff. For Bismarck universal suffrage had political significance in several respects: partly as a national
German tradition from 1849, partly as a counterweight to dynastic interests in the individual states, partly as a contrast to the Prussian
Chamber of Deputies, founded on class voting and indirect election, partly finally as a means of frightening the reactionary governments
abroad. Cf. \emph{Sybel}: \textit{Die Begründung des deutschen Reichs.} II. p. 451. 540. V. p. 379.}. The extension of political freedom
is no dissolution of the state, but precisely the perfecting of political life itself.

Political freedom is \textit{one of the most important means of education} for a people\footnote[2]{Cf. \emph{James Bryce}: \textit{The
American Commonwealth.} London 1888. II, p. 316: ``Universal suffrage, as it exists in the United States, is not only a great means of
security in the present, but is perhaps the mightiest educating force to which the masses of mankind have ever been subjected.'' % translated from Høffding's Danish
--- See also H. F. \emph{Feilberg}: \textit{Dansk Bondeliv.} København 1889. p. 130. 137.}. It calls upon the individual to procure for
himself all the enlightenment and knowledge he can in order to be able to perform
% --- p. 579 ---
\opage{579}his office as an active member of the political community. But above all he learns to cast his gaze beyond the narrow horizon to
which his life's work confines him, and he feels his human worth as one who has his mite to contribute to the life of the people. ---
Besides being a natural fruit of activity in the domain of culture, political freedom is thus also a means of strengthening the state by
binding individuals more closely to it, and a spur to the energetic development of personality and its abilities.

Even if one is no adherent of the regime of guardianship or of bureaucracy, one can readily admit that political freedom has its great dark
sides. It can be misused, which is what the best is always exposed to. But here it is not a matter of a historical account of it; it is a
matter, rather, of showing its necessity and the value that must be attributed to it from an ethical point of view. We must submit to the
dangers and drawbacks of freedom in order to share in the indispensable goods of which it is the condition.

4. Political freedom we have provisionally defined as the right to active participation in political life. But what more precisely is
understood by this? We cannot, after all, all govern immediately. Something of the kind could approximately be done in the small states of
antiquity, where state and city coincided; but it worked there too only as long as they coincided. For larger realms with many cities and
communes it does not fit. In modern times, therefore, the participation of single individuals in political life takes place through the
election of representatives. The election is the act of will by which everyone --- besides through participation in public debate, which
the regime of guardianship does not exclude either --- can intervene in the course of political life. Where nothing can become law, ---
where no tax can be levied, --- and where no decisive provision concerning the inner and outer fate of the state can be made without the
consent of the people's elected representatives, there a free \textit{constitution} exists, whether it is written down on paper or not;
and there a relation of confidence between government and people can prevail of a deeper and firmer kind than that which can be present
under a government of guardianship, where the people follows passively because it is perhaps allowed to express its opinion, but not to
have any will.

% --- p. 580 ---
\opage{580}``But,'' one might ask, ``is it not after all the most important thing that the wisest and best men govern, whether they
enjoy the people's confidence or not?'' --- To this it must be answered that it is precisely the great question how one is to settle
which men are the wisest and best. Spiritual stature is not so easy to measure as bodily stature. The great strife in the human world
is, after all, precisely about the standard of wisdom and goodness, and political strife is only a branch of this great strife. But this
much is certain, that the really wisest and best statesmen (one can be wise and good without being a statesman) will be able to find out
not only \textit{what} serves the good of the people, but also \textit{how} it can be realized in the best way. They will see that only
that serves the good of the people which agrees with its nature and its forces, and they will know how to awaken and guide these forces.
They do not regard the people as a dead material, a mass that is passively to take the form one wishes to give it. They will work not
only for the people but also with the people, and they will be able to win the confidence of the people, in which they will see one of
the very most important conditions for being able to accomplish anything.

5. Even the freest constitution, however, cannot be wholly simple and uncompounded if it is to last. There are two mutually opposed
dangers to which it is exposed, and which must be forestalled by special arrangements.

The one danger consists in the involuntary promptings, the prejudices and passions that prevail in the great multitude, passing over
immediately into action without being tested and purified through calm and many-sided deliberation. The state then comes to pursue a
policy of instinct or impulse, not a policy of will. Here again (cf., e.g., XXXIX, 1) the importance appears of an interval being brought
about between prompting and action. The people, like the single individual, must be able to show self-control, must be able to check
involuntary promptings, so that deliberation can take place. Such self-control is an essential quality of every character, and that the
people has character is to show itself precisely in the fact that not the fancies and promptings of the single moment, but constant and
% --- p. 581 ---
\opage{581}firmly grounded thoughts and feelings determine its course. (Cf. III, 4---7 and XI, 8---9). Because nothing can become valid in
the state without the will of the people, the people's momentary promptings need not therefore pass over into reality at once. An
interval of time between impulse and resolve is necessary. The people must be protected by its constitution against its own rash acts.
The experience and expert knowledge to be found within the people must get the opportunity to assert themselves before the decision is
made. How this goal can be reached in the constitution of the single people will for the most part depend on the history of the people.
As a single example it need only be mentioned that the North American free states, by the way in which the relation between executive
and legislative power and between the two chambers is ordered, and by the independent position given to the supreme court, have developed
a constitution that --- on account of the guarantees it offers in the respect mentioned --- has won recognition even from those who look
with very great misgiving on the development of ``popular government''\footnote[1]{\emph{Henry Maine}: \textit{Popular Government.}
London 1885 --- It is characteristic that the author who in recent years has spoken against democracy as a form of government with the
greatest expert knowledge and talent is not far from regarding the North American constitution as an ideal.}. --- It must be well noted
that what needs to be checked with the help of reflection and deliberation may be not only impatience and an uncontrollable urge for
change, but also conservative inertia. There may be prejudices, egoism and narrowness in the popular way of looking at things that only
gradually give way to clearer insight and freer thoughts. Important and well-grounded reforms may meet with resistance from the great
multitude, which in its antipathy to the new ideas perhaps even prefers, at any rate for a time, a regime hostile to freedom. The history
of plebiscites in France affords illuminating examples of this. The Jesuits in Lucerne in the forties found their account in appealing to
a general popular vote. The appeal to the people (referendum) introduced by the Swiss federal constitution of 1874 has proved to be a very
conservative institution. According to the Swiss federal laws a popular vote is to be held on a law when 30,000
% --- p. 582 ---
\opage{582}voters demand it. In the period from 1874 to 1893 the referendum was employed for 19 laws (out of 196). Of these 19 laws the
people approved only 6 but rejected 13. In the canton of Zürich, where every law must be submitted to a popular vote, in the period from
1869 to 1893 only 97 laws were approved out of 128 laws that the cantonal council had adopted. On the whole, in the cantons where a
compulsory referendum exists, a fourth or more of the laws passed by the legislative power are rejected\footnote[1]{\emph{Lawrence
Lowell}: \textit{The Referendum and Initiative.} (Journal of Ethics VI). p. 52. --- The referendum existed earlier in single cantons,
just as it already existed in the eighteenth century in some of the North American free states and later in Australia. Cf.
\emph{Stewardson} in ``Journal of Ethics'' XIII. p. 134.}.

Another, quite opposite danger is that under a constitution founded on political freedom there should be too much and too long
deliberating, negotiating and talking, so that there is no willing and acting with the energy necessary for the existence and progress
of the state. The interval that has proved so necessary thus spreads too far here\footnote[2]{Here too we have a psychological parallel.
Cf. \textit{Psykologi} VIII B, 3.}. This danger, however, like the preceding one, will be lessened the more a relation of confidence
prevails between the people and its government. Only as long as suspicion and fear of encroachments on the part of the rulers assert
themselves will one seek to limit the executive power as much as possible. Where on the other hand a relation of confidence develops,
there will also be readiness to put as much power at the disposal of the government as is necessary to solve the tasks at hand. Without
such a relation of confidence no free state can be a strong state, and really no state at all can be a strong state. --- But it is, as
already remarked earlier (XXXVIII, 7, end), precisely here that it appears that every state is a state of necessity. The right relation
between power and freedom can only be approximately brought about through many oscillations and struggles, --- struggles in which there
may be guilt on both sides.

6. It is, where political freedom prevails, the majority of the voters
% --- p. 583 ---
\opage{583}in the country that comes to determine the decision of affairs. This cannot be otherwise. When a decision is to be made, it is
more reasonable that the minority conform to the majority than the reverse. The will of the majority is the greatest moving force, and
it alone is fitted to represent the will of the whole society when absolute unanimity is not present. The demand made on the minority to
bow to the will of the majority has as its presupposition that the minority too wills that society shall continue to exist; from this
presupposed will there follows with necessity the acknowledgment of the strongest will in society as the deciding one\footnote[1]{Cf.
\emph{John Locke}: \textit{Of Civil Government.} II. 96.}. The only thing the minority can fairly demand is to get the opportunity freely
\textit{to express its views and put forward its reasons.} The minority must be allowed to try to convince the majority and thereby itself
become the majority. Therefore it is of great significance that the minorities be represented, % the print has ``Majoriteterne'' (PRINTED AS IS; the sense is minorities)
and the idea of proportional representation, which shows the most practical way to this, will intervene mightily in future political
development. A real representation of the people will only thereby be brought about.

Everything great and good has first proceeded from single men and found its first nurture in smaller circles. (See XXXVIII, 5). But
therein lies no objection to the majority's right to determine the final decision. No single man and no minority can demand immediate
and instantaneous dominion for its ideas. They can do so all the less since not even the majority has the right to carry its will through
\textit{instantaneously}. A dominion of the minority would, after all, be a worse tyranny than a dominion of the majority. For it is a
matter of chance who the minority comes to consist of, so that one has no right to assume that it would always be the best and wisest.
And the minority can just as well be radical as conservative. But in all cases there will be more who must bend their will and their
reason when the minority governs. --- Naturally the minority \textit{can} have right on its side. This kind of tragic conflict can just as
little be avoided here as it can be avoided, e.g., in
% --- p. 584 ---
\opage{584}the relation between law and morality (XXXVII, 3). The right of the minority cannot go further than freely to develop its
thoughts and assert them before the decision, and to criticize the decision itself.

It is a frequent charge against the democratic constitution, which lets the will of the majority prevail, that it leads to tyranny of the
majority, and that it is unfavorable to the rise of original men and ideas. Already in the remarkable work that played so great a role in
the coming into being of the American federal constitution, \textsc{Alexander Hamilton} discussed how this danger could be avoided. He
saw a barrier to the tyranny of the majority in the great multitude of interests and tendencies that a state that is not all too small
will comprise. Just as religious freedom is secured by the multiplicity of sects, so political freedom will be secured by the
multiplicity of interests. From this Hamilton drew precisely an argument in favor of carrying through the Union, since he thought that
freedom would be less in the North American states if each of them were to form a small, closed whole by itself\footnote[1]{\textit{The
Federalist.} (1787). Nr. 51. (Boston 1882. p. 401).}. When \textsc{Tocqueville} half a century later visited America and then wrote his
famous book on democracy in America, he found in the tyranny of the majority a chief fault of the American people and its form of
government. ``I know,'' he said, ``no country where, on the whole, less independence of mind and true freedom of discussion prevail
than in America. The majority draws a formidable circle around thought.'' % translated from Høffding's Danish
The most recent investigation of American conditions, however, has come to another result. \textsc{James Bryce} is inclined to believe
that \textsc{Tocqueville} here, as on other points, was somewhat too hasty in his judgments. He at any rate found no tyranny of the
majority, either in a political or in a religious respect. Just as little does he agree with Tocqueville in the assertion, repeated again
and again since, that democracy is supposed to be an unfavorable soil for original men and original ideas. He found that the Americans had
precisely a tendency to hero worship, and that
% --- p. 585 ---
\opage{585}he who has achieved something eminent in one of the domains of spiritual life is more highly esteemed in America than in most
European countries. In general, according to \textsc{Bryce's} experiences, the higher classes in America had not lost in culture through
democratic development, whereas the lower classes had gained in independence. Tocqueville seems in doctrinaire fashion to have attributed
far too great a significance for the life of the people to the form of government\footnote[1]{\emph{James Bryce}: \textit{The American
Commonwealth.} London 1888. I. p. 14; III. p. 134---143; 169; 534---549.}.

7. Political freedom brings with it a peculiar kind of community formation that is just as free as community formation in the domain of
culture and yet is most closely bound up with political life and intervenes in it. It leads naturally to union among those who agree on
the great political questions, that is, to the formation of \textit{political} parties. That different parties form has its ground partly
in the fact that the great political questions can be seen from different sides, partly in the fact that the ways of thinking and
dispositions of single individuals are different and with advancing civilization will surely become more and more different. The party
system is necessary wherever freedom prevails, and there is hardly anyone who declaims against it without having some party or other that
he consciously or unconsciously excepts. The drawbacks it brings with it can be remedied only by better and more energetic use of freedom,
that is, by the very cause that has called it to life.

It is of the greatest significance what motives and points of view guide the formation of parties. History shows us now racial
oppositions, now religious oppositions, now class and estate oppositions as determining in the formation of political parties. But each of
them brings great drawbacks with it.

With \textit{racial opposition} the danger of a division of the state lies plain to view, and often there is even work directly toward this
goal, so that the party opposition is only the first intimation of complete dismemberment. From 1830 to 1864 the most important party
opposition within the Danish state of that time was determined more and more in this way, and the result was indeed a real splitting
apart.

% --- p. 586 ---
\opage{586}Where \textit{religious differences} form the basis of political party division, it is to the injury both of religion, whose
domain is the inward man, and whose inwardness is violated by its being used as weapon and banner in the struggles of outer life, --- and
of the state, whose affairs cannot be ordered according to religious points of view. The more the church is withdrawn from the state, the
more this basis for party formation must fall away. From a standpoint that would have ecclesiastical ideas rule the world, it is consistent
when one\footnote[1]{\emph{Ketteler}: \textit{Freiheit, Autorität und Kirche.} p. 99.} thinks that the difference of political parties
must in the end be determined by the acceptance or rejection of the supernatural. And in countries where the population is divided into
two great ecclesiastical communities, each of which jealously guards its position and its interests, there a religious grounding of party
division will not easily be kept away. It can, when it is driven to the extreme, lead just as well as racial opposition to a splitting of
the state. But it is then at least natural and honest. In countries, on the other hand, where religious freedom prevails, and where people
of the same religious faith can belong to different political parties and people of different religious faith to the same political party,
--- there an appeal to religion in party struggle can only have as its purpose to disparage and cast suspicion on opponents.

It cannot be otherwise than that \textit{class and estate oppositions} exercise great influence on party division. Behind every great
political question there has hitherto stood a social question. The political struggle is often only the form under which social
oppositions show themselves. This follows already from the fact that progress in history consists to a great extent in ever new strata of
society coming forward to get political freedom and the right to have a say. Such was the progress that took place in France in 1789, in
England in 1832, 1867 and 1884, and with us in 1848. The great goal that historical development aims at is to put an end to all dead and
passive mass within society and to give as many as possible access to the free use of their forces in the material and spiritual domain,
and
% --- p. 587 ---
\opage{587}political freedom is partly a consequence of, partly a cause of this free use of forces. --- But where the opposition of estates and
classes is one-sidedly determining, narrow and irrelevant considerations come in just as when races and sects determine party division.
Particular interests come to play the greatest role at the expense of the common good of society, and the people is split into different
strata that in the end no longer understand one another. Class egoism stands out especially clearly where a stratum of society that has
for a time been the favored and ruling one does everything to keep the next one down, although the nature of things and the importance of
a sound and peaceful development seem to impress the necessity of meeting the new strata halfway in such a way that their political
education can take place without too great difficulties.

Every class of society, like every single individual, has its definite calling within society as a whole. Party division determined by
class opposition must therefore, if a sound development is to take place, be subordinated and gradually give place to a party division
that springs from the very fundamental conditions of \textit{the existence of the whole society} and its development. The ideal for party
formation would be that each party gathered round one of these fundamental conditions. It is here a matter especially of two principles,
\textit{order and progress:} order as the basis for the calm course of life and for securing what has been attained, and progress, because
conditions change, and he who does not go forward goes back. To this opposition in the fundamental conditions of society corresponds a
main difference in the gifts and character of single individuals. Some are disposed to preserve, others to acquire; the minds of some are
turned with reverence back toward the good old, those of others look with hope forward toward the good new. The ideal party division would
be present where individuals of all different classes of society united round the one or the other of these two great principles. It is
no good sign when one can at once infer a man's political views from his social position. The two great political parties grounded in
nature ought to be represented in all classes of society from the highest to the lowest. The smaller
% --- p. 588 ---
\opage{588}groups would then be able to find their place on the one side or the other. In the political party relations of England and
North America it has hitherto gone in this way, and on this rests not least the sound character that political development has on the
whole had in these countries.

With such a basis for party division the parties would naturally regard one another as supplements. The one party will feel that it
needs the criticism of the other. And the criticism will essentially consist in pointing out that the opposing party does not answer to
its name and its idea. The conservative party is not criticized from a wholly foreign point of view when one draws its attention to the
fact that its task is to preserve what has value and to see to it that the continuity in the life of the people is not broken, but not
blindly and stubbornly to oppose every change. A superior conservative statesman will precisely lead his party to carry out reforms,
since they can then be carried through with reverence for what has been handed down and without all too abrupt transitions. The party of
progress fails its own idea when in its zeal to change what has been handed down it violates the lasting and deeper-lying conditions of
progress, which consist in the seriousness and proficiency of the national character. Restless haste and merely negative criticism are not
good conditions for the development of character. Character is developed through \textit{constant} work; the extraordinary significance of
the work of reform must not make one blind to the fact that without continuity in the conditions of life the strong and deep feelings
cannot develop that are the driving force in all action, whether it aims at reforming or at conserving.

The opposition between the conservative party and the party of progress combines in different ways with two other oppositions, whereby
the most various combinations become possible. The opposition between centralization and decentralization does not without more ado
coincide with the opposition between order and progress, and just as little is this the case with the opposition between the display
of power and the work of culture. At each particular time the party relations in a country are determined by the varying relation among
all three pairs of opposites. ---

Party formation is an organization without which political
% --- p. 589 ---
\opage{589}life cannot have coherence and clarity. The various wishes and plans that move in the mind of the people are by it brought
back to a few main thoughts. Like every community, the party too has its egoism; it can set, and very often does set, its special
interests above the interests of the public. But it will, as said, in such cases come into contradiction with its own idea and will
sooner or later come to feel the consequences of this. The sense of justice and public opinion in the people will besides always
preserve a certain independence of the party system and exercise their influence on its development. In North America professional
politicians have obtained a great influence within the individual parties and have often misused it in the most shameful way, ---
something whose ground must be sought partly in the all too numerous elections, partly in the political raw material that immigration
constantly brings and that easily falls into the hands of the professionals. But the sound sense of the people and the power of public
opinion set definite limits to the abuse of the party system and are able to instill in egoistic partisans the requisite
respect\footnote[1]{\emph{James Bryce}. III, p. 151 ff., 333 ff.}. A party may just as well need reform as the whole state may. And the
source of reform is in both cases the same: from the free development of life in the smaller circles come new problems and tasks on
which the existing parties are compelled to take a position. Parties are means for the people and are no more to lord it over it than
the state is. In themselves they are just as little really productive as the state. They are to be compared with the rings that form
when a stone has been thrown into the water; what is really productive is the throw of the stone; the rings are the forms under which
its effect spreads over the surface of the water.

The relation between the individual and the party he has joined, or whose representative he is, has the character of an ethical
relation. Conflicts can arise here of the same painful and tragic character as everywhere one must break with a smaller whole for the
sake of the larger whole. In the famous speech in which Sir \textsc{Robert Peel} took leave of political life he said, among other
things: ``In relinquishing power,
% --- p. 590 ---
\opage{590}I shall leave a name, severely censured I fear by many who, on public grounds, deeply regret the severance of party ties --- deeply regret that severance, not from interested or personal motives, but from the firm conviction
that fidelity to party engagements --- the existence and maintenance of a great party --- constitutes a powerful instrument of
government.'' % Peel's English (House of Commons, 29 June 1846); Høffding's Danish abridges it
He thus recognized that by carrying through the repeal of the Corn Laws against the will of his party (the Tory party) he had broken a
bond of an in itself ethical nature. But he held fast to the view that a higher ethical consideration compelled him to it: the necessity
of removing an injustice that made the poor worker's daily bread dearer and more bitter to him\footnote[1]{\emph{Molesworth}:
\textit{History of England from the year} 1830---1874. Abridged Edition. London 1877. p. 367.}. In this statement the whole ethics of the
party system is really contained.

8. However perfect one imagines the organization of parties, and however harmonious the relation between government and people, the
state and its constitution will nevertheless hang in the air if there is no firmer foundation for political life than that discussed so
far. It is not enough that what stirs in single individuals, in the family and in the free cultural communities can through political
freedom and through the struggle of parties gain influence on the development of political life. There is here still too great a leap
from individual life and free social life on the one side to the state on the other. This interval can be filled only by
\textit{self-government,} i.e., by as many as possible of the public affairs being seen to by the citizens' own activity and not merely by
a bureaucracy appointed by the government. Self-government forms an intermediate link between activity in the free cultural community and
the activity of the state proper. In the cultural community the weight is laid on common ends; in self-government it is \textit{local
belonging together} that makes common activity natural. And by this cooperation conditioned by local belonging together (neighborly
relations) self-government differs from the political life proper that unites the population spread over the whole territory
% --- p. 591 ---
\opage{591}of the country. The small local circles (communes, counties) each have interests and tasks whose satisfaction requires
organization, but which cannot be grasped and treated in their full peculiarity from the standpoint of the more remote central
government. It is common material, ideal and philanthropic tasks that are to be solved, and in solving them the small local circles have
the advantage that they can more easily collect exact experiences, --- that all individuals can more or less know one another, --- and
that the machinery can be limited to the least possible. Connection with the state cannot, however, be dispensed with. Partly the state
must exercise a control in order to prevent encroachments and arbitrariness toward single individuals in the circle, partly the
experiences from the other local circles must be made fruitful. There is just as little any necessary hostile relation between
self-government and the government of the state as between political freedom and the order of the state. Just as the state is precisely
strengthened by all citizens having the right and duty to take part in the decision of affairs of state, so it is also strengthened the
more independently and vigorously life stirs in the local circles. We then do not get a bureaucratic machinery of state, culminating in
the government, on the one side and all the single, isolated individuals on the other. Absolute monarchy had led to such a dualism in
many European countries\footnote[1]{Cf. \emph{Tocqueville}: \textit{L'ancien régime et la révolution.} 7 éd. p. 100. ---
\emph{Gneist}: \textit{Das Selfgovernment in England.} 3 Aufl. p. 978. --- \emph{Wallace}: \textit{Rusland.} Danish transl. I, p.
168---174. --- E. \emph{Holm}: \textit{Danmarks og Norges indre Historie under Enevælden fra 1660 til 1720.} II. p. 407. --- The French
National Convention was hostilely disposed toward l'esprit de localité in departments and communes. By the law of 24 August 1793 the
juridical existence of the communes was abolished, their property and their debts passing to the state. Thereby political freedom lost
its points of support. \emph{Arthur Desjardins}: \textit{De la liberté politique dans l'état moderne.} p. 12. --- The National
Convention thus followed in the footsteps of the absolute rulers.}. And it could not be otherwise. All kinds of absolutism set up an
either--or: either state or individual. The more power is laid on the one side, the less remains on the other side. If one demands an
absolute concentration of all public activity, then individual
% --- p. 592 ---
\opage{592}and civic self-activity must disappear. The central power then looks with suspicion on every independent movement, and with a
similar suspicion the population in turn will naturally come to regard the measures of the government, whether they go in a reforming or
a conserving direction. Through self-government, on the other hand, individuals are accustomed to take an interest in and work for common
and public affairs. Through activity in the smaller circle the sense for the larger circle is nourished. Self-government is therefore a
necessary foundation for a free political life. --- In this as in many other respects England is a model. The English constitution rests,
as \textsc{Gneist} in particular has shown, on the firm foundation that self-government in the smaller, local circles affords. The English
constitution has grown out of administration. Thereby there has been provided at once a good school for participation in political life
and also a domain where common activity can take place without necessarily being touched by the great political oppositions.

Only by resting on such a foundation does the state really become what by its definition it was to be: the organized people.

\chapterhead{XLI.}{CONCLUSION.}

1. The ethical world consists of a series of smaller worlds, each of which has its own center. Each single individual is a little world
of its own, even if the forces of a larger world reach into it. It is similar with family, cultural community and state. They are all in
the end embraced by the great kingdom of humanity, which extends as far as the human race, and which at some points even widens itself
% --- p. 593 ---
\opage{593}to embrace all beings with the capacity to feel pleasure and pain. (Cf. XII, 3). Several times (e.g., III, 17; XII, 1; XIV, 5;
XXXVI, 4) it has been shown that it would be a misunderstanding to think that one worked best for the good of the great world by skipping
over the smaller worlds. There is an inverse relation between the strength and the scope of feeling, and it is therefore of decisive
significance that interest be concentrated in the smaller circles and grow strong there before it passes into the larger circles.
Faithfulness in small things is the condition for being able to accomplish anything in great things. In the ethical world a law of
continuity prevails just as in the physical world; one cannot skip anything. The idea of the universal kingdom of humanity has for the
time being only the significance that it forbids us to set up arbitrary barriers and declares that \textit{every} human being has a center
in itself, which prevents it from ever being regarded as a \textit{mere} means. The kingdom of humanity does not lie \textit{outside} and
\textit{beyond} family, cultural community and state, but provides precisely --- as has constantly appeared in the foregoing --- the
principles for the ethical ordering of these circles. We are not first members of families, workers of culture and citizens and then
human beings; but we are precisely to live as human beings and treat one another as human beings in all relations within family, cultural
community and state. Not only for the \textit{scope} of our endeavors, but also for their character, the idea of universal human society
thus gains significance, and it was therefore made the basis at every point in the foregoing account. The kingdom of humanity can be where
only two or three are present. % cf. Matt. 18:20

But this does not exclude that there may be a question of a special organization of universal human society as a society that embraces all
nations and states. It would be a natural continuation of development within the single people. Just as the people gets its organization in
the state, might not an organization then be conceived that bound all states together into one great whole?

The different states still live in the state of nature in relation to one another. It is power that stands against power, and the
% --- p. 594 ---
\opage{594}ruling need is the need to spread at the expense of others. The idea of perpetual peace stands, no doubt, as a distant goal; but each
single state would prefer to enter this perpetual peace as great and rich as possible. And it is not even said that it is admitted that this
goal can ever be reached. A great German general has even called war a necessary element in the divine order of the world. It is a legacy
that seems now to have passed from the German philosophers of an earlier time to the generals, to credit oneself with an intimate
knowledge of the divine order of the world. From a more natural point of view war stands as a remnant of the brutality that is as a rule
characteristic of lower stages of human development. It by no means follows with necessity from this that war will ever cease; for it is
not at all said that we shall reach the highest degree of development we can imagine. But in ethics we have nothing to do with fantasies
of the future either. On the other hand it is an ethical consideration how far there is a possibility of realizing the kingdom of humanity
and peace and of pushing back brutality to a higher degree than hitherto. Thereby a whole series of endeavors is pointed to from a new
side, some of which have already been discussed from other points of view.

2. Such a state of war as now prevails among states once prevailed within the single states between the different individuals, the
different kindreds and classes. The modern state first came into being by stopping this inner war. One of the first forms under which it
succeeded in pushing back the lust for attack and the thirst for revenge was mediation. The first courts were courts of mediation. What
has been attained in the smaller sphere one will also with reason try to carry through in the larger. Courts of arbitration are, after all,
already not so rare between states, even if the questions referred to them are not the really burning and decisive ones; these one prefers
to reserve for decision by force --- when one thinks one possesses the force.

Greater hope is given by another consideration. --- Wars themselves are waged more and more humanely. A generation ago the philanthropic
endeavors of the Genevan \textsc{Henry Dunant}
% --- p. 595 ---
\opage{595}succeeded in bringing it about that the wounded and the whole military or voluntary medical staff whose task is to nurse them were
declared neutral in the midst of battle. After Dunant had with great self-sacrifice prepared the matter and negotiated about it with
various governments, the Swiss Federal Council invited to a congress in Geneva in August 1864, where the so-called Geneva Convention was
concluded. From that time on the red cross, the symbol of the Geneva Convention, marks a place under protection of peace, even if the
battle rages round about\footnote[1]{\emph{Rudolf Müller}: \textit{Henry Dunant, der Begründer des Roten Kreuzes und der Genfer
Konvention.} Stuttgart 1896.}. --- And although the great powers stand on a constant war footing toward one another, their intercourse is
nevertheless far more conditioned by the interests of peace than by those of war. Whereas international law in the time of
\textsc{Grotius} consisted essentially only of rules for the conduct of war and the conclusion of peace, as is also indicated in the title of
his work\footnote[2]{On the Law of War and Peace \textit{(De jure belli et pacis).} The work appeared in 1625. At the beginning of the
first chapter Grotius says that it is war he will treat of, but that he also comes to speak of peace, since war springs from it and in
turn leads to it.}, it is now especially the relations and tasks of peace that provide the material for the provisions of international law
observed by the different states\footnote[3]{``Modern international law consists for the most part of rules whose existence is due to
physical and social conditions that did not exist two centuries ago. On account of the ease of travel, on account of postal and telegraph
communication as well as on account of enlarged moral conceptions and more enlightened principles of commerce, the intercourse of the
nations with one another and the intercourse between the citizens of the different nations is far more considerable and important in time
of peace than in time of war.'' \emph{Sheldon Amos}: \textit{The science of Law.} 2. ed. London 1874. p. 341. % translated from Høffding's Danish
}. It is the tasks of culture (of material, ideal and philanthropic culture) that bring the peoples into contact, and however frequent wars
may be in some periods, peaceful contacts nevertheless increase in frequency and multiplicity to such a degree that they far outweigh them.
We have in another connection discussed the great community-founding significance of trade. Ideal culture follows in the footsteps of
material culture, and sometimes even goes ahead of it. Both egoistic and ideal interests
% --- p. 596 ---
\opage{596}thus work toward binding together people of different nations. There exists in fact a human society for which national boundaries
have no decisive significance. --- In his remarkable essay \textit{Zum ewigen Frieden} (1795) \textsc{Kant} laid special weight on the
fact that egoistic interests and the commercial spirit must bring about a state of peace and a universal federation of states, although it
was naturally his conviction that this first gets its proper grounding when its acknowledgment springs from an ethical disposition.
Unfortunately it proves that commercial interests do not always unite, but often just as much bring into strife. European and Asiatic
powers contend for markets to dispose of their enormous production. They hunt for customers, and they come into strife over the hunting
grounds, --- just like the Indians who, following the tracks of the buffaloes, come into strife over the buffalo districts.

Peaceful development will be favored the more perfectly it succeeds within the single state to push back the brutal element of power, the
more, that is, political freedom really comes to hold. The use of power outwardly and inwardly are closely connected.

Ethically regarded, it was laid down already when the idea of love of mankind first came forward (III, 8; 13) that there is a universal
kingdom of man. If love of mankind now grows in the race, it will also lead to the practical acknowledgment of this kingdom and limit the
state of war to the least possible. --- It is not only among civilized peoples that this practical acknowledgment is still greatly lacking.
It holds no less as regards the relation between the European peoples and the inhabitants of the other parts of the world. The Europeans
have everywhere appeared as the born masters of the earth and have without more ado regarded the other races as without rights, their land
as without an owner. Like Thomas \textsc{More's} Utopians, they think it just to make war on a people that does not exploit the resources
of its own land when one needs these oneself for one's existence. This principle often enough leads in practice to savage peoples being
regarded as without rights. But conscience has begun to stir here. Not only have philosophers long ago protested
% --- p. 597 ---
\opage{597}against this conduct\footnote[1]{\emph{Kant}: \textit{Zum ewigen Frieden.} Königsberg 1795. p. 42.}, but ``societies for the
protection of the uncivilized aboriginal peoples'' have been formed\footnote[2]{L. \emph{Felix}: \textit{Der Einfluss der Sitten und
Gebräuche auf die Entwickelung des Eigenthums.} Leipzig 1886. p. 395. --- For the rest, it is reported here that the Chinese too oppress the
aboriginal inhabitants of the Celestial Empire and regard them as without rights.}, and the great statesman who led the founding of France's
colonial empire energetically maintained the duties of the higher race toward the lower\footnote[3]{A. \emph{Rambaud}: \textit{Jules
Ferry.} Paris 1903. p. 317. 516 f. (Ferry maintained among other things that one had no right to punish savages for ``breach of treaty'' as
long as one had not taught them what a treaty is).}.

Even if perpetual peace should never come to prevail, there is thus both a possibility of and an ethical necessity for an advance in the
direction of the ideal it indicates. The endeavors that can lead in this direction are for the most part merely a continuation of
endeavors that already appear within the smaller circles.

3. Almost a hundred years ago \textsc{Kant} raised the question whether the human race was constantly progressing in a moral respect.
Experience, he thought, was uncertain and ambiguous in this respect, unless a fact could be pointed out that decidedly testified to a moral
disposition, a tendency toward the good in human nature. On such a fact belief in progress could rely. Such a fact he found in the general
enthusiasm that the attempt of the French people to found a state of society more in accord with reason and justice awakened round about
in Europe. ``The revolution of a gifted nation,'' he says, ``which we have seen taking place in our day, may succeed or miscarry, --- it may
be filled with misery and atrocities to such a degree that a right-thinking man, even if he could hope, undertaking it a second time, to
carry it out successfully, would nevertheless not make the experiment at such a cost, --- but this revolution nevertheless finds in the
minds of all spectators (who are not themselves caught up in the game) a sympathy bordering on enthusiasm, the very expression of which is
% --- p. 598 ---
\opage{598}bound up with danger, so that it can only stem from a moral disposition in the human race.'' ``Such a phenomenon in human history
\textit{is not forgotten again,} because it has revealed a disposition and a capacity for progress in human nature that no politician could
have figured out from the previous course of things''\footnote[1]{\textit{Streit der Fakultäten.} 1798. (Vermischte Schriften. Halle
1799. III) p. 440 f. 445.}. % translated from Høffding's Danish

What Kant sought there is at all times reason to seek. And he is surely right that the essential thing is whether there is an inner force
and tendency to build on. Only if our experience gives us facts like the one Kant pointed to is there reason to share his hope. The whole
account of ethics given here rests precisely on the presupposition that there is in human nature a capacity for universal and disinterested
sympathy, and the valuation of human actions and forms of life that has here been attempted rests on this foundation.

But in our day we would surely demand more than the eminent minds of the eighteenth century demanded. It was a grand idealism that marked
the close of that century. A belief in the powers of the human spirit, a belief all the more justified since some of the most significant
works of the human spirit stem from that time, prevailed in various forms among the spokesmen of the age. No wonder that the conviction was
held that if only all barriers and hindrances to the free unfolding of the human spirit were removed, there were no limits to progress. At
no time perhaps has a purer and nobler enthusiasm gone through the world.

Much has been experienced since that time. We have become much wiser and must see to it that we do not become wise beyond our years. We no
longer have the firm belief in the inner power of the human spirit to build and order the relations of life according to its ideals. We
have gained an eye for the many complicated conditions that assert themselves here. We have learned a good deal of history, not only that
which has taken place in the last century and which
% --- p. 599 ---
\opage{599}lies so near to us; earlier times too now stand before us in another light. We stand more or less doubting before reformers and
improvers of society. The toughness and endurance of historical powers, the deep-lying vein that brings the blood of the past still to
pulse in the organs of the present, the firm bonds by which the present, even with those who think they represent the newest of the new,
hangs together with the old, the forgotten or even the despised, --- all this we have learned to open our eyes to. We admire, no doubt, the
mightily bursting stream, but then set ourselves to seek out the many small brooks to whose confluence it owes its force. The best, the most
honest and the most energetic will can do nothing when the conditions are not there. And often it seems to us that the conditions for an
advance are so many that one must despair whether they can really come about. History thereby often takes on a tragic character for us,
since we have so often seen the best forces founder or lead to misfortunes because reality offered too strong a resistance, and the small
brooks could not be united into one great stream.

Yet the historical mode of view, which with us has taken the place of the more abstract tendency of thought of the eighteenth century, need
not stifle idealistic faith, however severe a test this may be put to. We must only make ourselves familiar with the fact that everything
great takes time, and that the beginning of the great is often inconspicuous. That is a thought that has repeatedly been brought forward in
the foregoing account as well. And connected with it is another thing: the involuntary goes before the deliberate; we must live and gather
experiences before we can form the thoughts that can work back on life; it is life itself that produces the forms that we can fix and
develop and use to carry life further. With the help of the intimate connection between the small and the great that the modern theory of
evolution has shown, it is possible for us to unite our realism with the idealism of our predecessors, by professing the rule: ``to be
enthusiastic only for what is great, but to be faithful in small things too.''

% ---------------------------------------------------------------------
%  INDEX (pp. 601--606): Høffding's Register, translated: the print's order (Danish alphabet) and page
%  numbers kept (said in one bracketed line); subject entries carry the Danish term in italics; names
%  as printed, with English forms where the English reader has one (Aischylos -> Aeschylus, Platon ->
%  Plato); the print's misprints kept with a comment.  Generated from the Register by
%  ~/oddspots/parts9/mkindex.py (a one-off; edit the entries here from now on).
% ---------------------------------------------------------------------
\clearpage
\chapter*{Index}
\addcontentsline{toc}{chapter}{Index}
\fancyhead[RE,LO]{\textit{Index}}
\opage{601}%
{\small\noindent [Høffding's \textit{Register}, translated. The entries keep the order of the Danish alphabet, and the Danish term follows each subject entry in italics; the page numbers are those of the 1913 edition, given in the margin of this translation.]\par}\medskip
\begin{multicols}{2}
\small\setstretch{1.1}\raggedright
\regletter{A.}
\entry Causal freedom (\textit{Aarsagsfrihed}) 107.
\entry The law of causality and ethics (\textit{Aarsagsloven og Etiken}) 106 ff. 122 f.
\entry The temperance movement (\textit{Afholdsbevægelsen}) 191.
\entry Deterrence (\textit{Afskrækkelse}) 556 ff.
\entry Aeschylus 126. 517.
\entry Public-mindedness (\textit{Alsind}) 249.
\entry Altruism (\textit{Altruisme}) 166.
\entry Ambrose 203.
\entry Anarchism (\textit{Anarkisme}) 164. 185.
\entry Repentance (\textit{Anger}) 117 ff. 565.
\entry Responsibility (\textit{Ansvar}) 117.
\entry Antisthenes 221. 348.
\entry Division of labor (\textit{Arbejdets Deling}) 361 ff.
\entry Aristippus 30.
\entry Aristotle 33. 84. 91. 98. 188. 204. 223. 231. 426. 450. 567.
\entry Aristotelian principle (\textit{Aristotelisk Princip}) 84 f. 267 f. 330. 343. 379. 481. 514 f. 543 f.
\entry Arnauld 137.
\entry Asceticism (\textit{Askese}) 137. 189 f.
\entry Augustine 109. 111. 129. 203. 218.
\entry Authority (\textit{Autoritet}) 14 ff. 59. 81 ff. 344.
\regletter{B.}
\entry Bain 105.
\entry Ball (Sidney) 399.
\entry Bang (Gustav) 419.
\entry Bang (N. H.) 69---72. 170.
\entry Childhood (\textit{Barndommen}) 329 f.
\entry Infanticide (\textit{Barnemord}) 327 f.
\entry Beale (L.) 228.
\entry Beneke 92.
\entry Bentham (Jeremy) 26. 44. 49. 55. 64. 94. 170. 248. 346.
\entry Birckner 241.
\entry Bismarck 578.
\entry Bjørnson 568.
\entry Boström 159.
\entry Buckle 433.
\entry Buddhism 499. 503.
\entry Burkhardt 296. 314.
\entry Butler 196.
\entry Brentano (Jujo) 376. 380. 402. 413.% PRINTED AS IS (Lujo)
\entry Bröchner 121.
\entry Bryce (James) 532. 578. 585. 589.
\regletter{C.}
\entry Campanella 387. 405.
\entry Carlyle 200. 254. 300.
\entry Carus (Paul) 168.
\entry Cathrein (Victor) 80. 82. 105. 153.
\entry Christensen (S.) 72.
\entry Cicero 90. 188. 202. 359.
\entry Clarke 247.
\entry Clifford 254.
\entry Coit (Stanton) 475.
\entry Coleridge 256
\entry Comte (Aug.) 166. 367.
\regletter{D.}
\entry Dante 120. 127. 141.
% [Register p. 602 begins]
\opage{602}\entry Darwin (Charles) 177. 293. 367.
\entry Democritus 85. 246.
\entry Despine 327. 571.
\entry Determinism (\textit{Determinisme}) 106 ff.
\entry Dhammapada 98.
\entry Dorner (A.) 318.
\entry Dostoevsky 133.
\entry Dunant (H.) 595.
\entry Duns Scotes 16.% PRINTED AS IS (Scotus)
\entry Dühring (Eugen) 334. 398. 541. 567.
\entry Virtue (\textit{Dyd}) 156. 159. 201 ff.
\entry Animals (duties toward) (\textit{Dyrene (Pligter mod)}) 247 f.
\entry Death (\textit{Døden}) 197.
\entry Death penalty (\textit{Dødsstraf}) 569 f.
\regletter{E.}
\entry Oath (\textit{Ed}) 486 f.
\entry Egoism (\textit{Egoisme}) 34 f.
\entry Ehrenfells 271.
\entry Eilschow 247 f.
\entry Property (\textit{Ejendom}) 234 f 418 ff.
\entry Eliot (George) 322. 482.
\entry Eliot (Charles W.) 494.
\entry Enfantin 289.
\entry Engels (Fr.) 389. 398.
\entry Epictetus 239.
\entry Esquirol 228.
\entry Ethics (\textit{Etik}) 3 f. 7---9. 22. --- Historical and philosophical ethics 9---13. --- Theological and philosophical ethics 13 ff. --- Ethics and metaphysics 18---20. --- Individual and social ethics 160 ff. --- Ethics and political economy 52 f. 417. Ethics and religion 457 ff. Ethical judgments 1---7. 22. 49 f. 75.
\regletter{F.}
\entry Trade unions (\textit{Fagforeninger}) 376 ff.
\entry Falbe-Hansen 235. 347. 349.
\entry Family (\textit{Familie}) 274 f. 278 ff. 593.
\entry Poor relief (\textit{Fattigvæsen}) 503 ff.
\entry Feilberg (Ludvig) 144.
\entry Ferguson 564.
\entry Feuerbach (Anselm von) 115. 526. 553.
\entry Ferry (Jules) 597.
\entry Fichte (J. G.) 93. 127. 166. 200. 332. 402. 415. 508. 526 f. 556 ff. 566.
\entry Philanthropy (\textit{Filantropi}) 495 ff.
\entry Philosophy (\textit{Filosofi}) 9 ff. 18 f. 437 f. 462.
\entry People (\textit{Folk}) 513 f.
\entry International law (\textit{Folkeret}) 594 ff.
\entry Model (\textit{Forbillede}) 198---201.
\entry Consumers' associations (\textit{Forbrugsforeninger}) 384 f.
\entry Constitution (\textit{Forfatning}) 572 f.
\entry Obduracy (\textit{Forhærdelse}) 134 f.
\entry Merit (\textit{Fortjeneste}) 104 f.
\entry Parents and children (\textit{Forældre og Børn}) 326 f.
\entry Fourier 388.
\entry Frederick the Great 493.
\entry Freedom (of the will) (\textit{Frihed (Villiens)}) 106 ff.
\entry Freedom (personal) (\textit{Frihed (personlig)}) 344 ff.
\entry Freedom (political) (\textit{Frihed (politisk)}) 576 ff.
\entry Free love (\textit{Fri Kærlighed}) 289.
\entry Free monogamy (\textit{Frit Monogami}) 286 ff.
\entry Free time (\textit{Fritid}) 414. 425 f.
\entry Industrial partnerships (\textit{Fællesforeninger}) 382.
\entry Community (\textit{Fællesskab}) 235.
\regletter{G.}
\entry Retribution (\textit{Gengældelse}) 117 ff. 549 ff. 564 ff.
\entry George (Henry) 416.
\entry Gneist 510. 545. 591.
\entry Good (a) (\textit{Gode (et)}) 155.
\entry Goethe 151. 200. 230. 425.
\entry Goos 339. 524 f.
\entry Grotius (Hugo) 554. 595.
\entry Grundtvig 536.
\entry Fundamental value (\textit{Grundværdi}) 3. 9. 51.
\entry Guyau 187.
\regletter{H.}
\entry Hafstrøm 285. 292.
\entry Hagerup 561.
\entry Hamilton (Alexander) 541. 573. 575. 578. 584.
\entry Hammurabi 302.
\entry Commerce (\textit{Handel}) 422 f.
% [Register p. 603 begins]
\opage{603}\entry Motive of action (\textit{Handlingsmotiv}) 60 f.
\entry Hansen (S.) 72 f.
\entry Harnack 38.
\entry Hartley 268.
\entry Hartmann (E. v.) 321.
\entry Heegaard (S.) 112. 115. 159. 162.
\entry Hegel 64. 94. 168. 269. 536.
\entry Devotion (\textit{Hengivelse}) 185 f. 205. 236. 260.
\entry Hensel 150.
\entry Herbart 158.
\entry Hettner 456.
\entry Hetaerism (\textit{Hetærisme}) 285.
\entry Hobbes 64. 184. 272. 277. 537.
\entry Holberg 92. 215.
\entry Holtzendorff 532.
\entry Hugo a St. Victore 129.
\entry Humboldt (W. v.) 309. 572.
\entry Hume 64. 184.
\entry Hutcheson 44. 64. 91. 248.
\entry Hägerström 57. 152.
\entry High-mindedness (\textit{Højsindethed}) 242.
\regletter{I.}
\entry Ibsen (Henrik) 133. 162. 454.
\entry Ideal 62 f. 64 f. 96 f. 596. --- See ``Model.''
\entry Ideal culture (\textit{Ideel Kultur}) 424 ff.
\entry Indeterminism (\textit{Indeterminisme}) 111 ff.
\entry Individualism (\textit{Individualisme}) 34. 35. 161 ff. 539 ff.
\entry Individual relativity (\textit{Individuel Relativitet}) 89 ff. 178 f. 222 f.
\entry Individual differences (\textit{Individuelle Forskelligheder}) 88 ff. 222 f. 464 f.
\entry Individual ethics (\textit{Individuel Etik}) 60. 183 ff.
\entry Sovereignty of the individual (\textit{Individets Suverænitet}) 541.
\entry Intellectual culture (\textit{Intellektuel Kultur}) 430 ff.
\entry Intellectual probity (\textit{Intellektuel Redelighed}) 254 f.
\regletter{J.}
\entry Jacobi (F. H.) 5. 91. 254.
\entry James (William) 255.
\entry Jevons 380. 381.
\entry Jhering 108. 236. 521.
\entry Johnson (Dr.) 259.
\entry Jouffroy 187.
\regletter{K.}
\entry Kant 11. 23. 26. 41. 98. 104. 116. 157. 159 f. 174---176. 204. 205. 226. 436. 502. 525. 597 f.
\entry Categorical imperative (\textit{Kategorisk Imperativ}) 41.
\entry Ketteler (Bishop) 443. 445. 586.
\entry Kierkegaard (Søren) 30. 92. 164. 192. 251. 253 f.
\entry The church (\textit{Kirken}) 476 ff.
\entry Church and state (\textit{Kirke og Stat}) 483 ff.
\entry Non-confessional school (\textit{Konfessionsløs Skole}) 339 f. 443.
\entry Consumption societies (\textit{Konsumforeninger}). See Consumers' associations.
\entry Cosmic feeling of life (\textit{Kosmisk Livsfølelse}) 459.
\entry Kropotkin 187. 380.
\entry Christianity (\textit{Kristendommen}) 20 f. 38. 161. 186. 203. 499 ff.
\entry Kroman 76. 112. 115. 116. 118.
\entry Culture (\textit{Kultur}) 139. 145---152. 340 ff. 354---356. 421---427.
\entry Cultural community (\textit{Kultursamfund}) 275 f. 342---512.
\entry The position of woman (\textit{Kvindens Stilling}) 310 ff.
\entry Love (\textit{Kærlighed}). See Devotion.
\entry The sexual relation (\textit{Kønsforholdet}) 225---229. 281 ff.
\regletter{L.}
\entry Lange (Julius) 456 f.
\entry Lasalle 372.% PRINTED AS IS (Lassalle)
\entry Liguori (Alphonso di) 82.
\entry Locke 583.
\entry Lombroso 558.
\entry Lotze 323. 567.
\entry Luther 134.
\entry Obedience (\textit{Lydighed}) 203. 332.
\entry Lyell (Charles) 350.
\regletter{M.}
\entry Maimonides 303.
\entry Maine (Henry) 581.
\entry Majority and minority (\textit{Majoritet og Minoritet}) 582 ff.
\entry Malebranche 137.
\entry Malthus 244. 367. 409.
\entry Mantegazza 295.
\entry Marcus Aurelius 186. 230. 232. 239.
% [Register p. 604 begins]
\opage{604}\entry Martensen (H. L.) 220. 437. 487. 493. 500. 535.
\entry Marx (Karl) 354. 389 ff. 392 f. 414.
\entry Marxism (\textit{Marxisme}) 389 ff.
\entry Mass (\textit{Masse}) 369 f.
\entry Material culture (\textit{Materiel Kultur}) 358 f. 424 ff.
\entry Maudsley 233. 238.
\entry Meinong 92.
\entry The human race (\textit{Menneskeslægten}) 38. 237 ff. 492 ff.
\entry Metchnikoff 227.
\entry Mill (James) 94. 268. 505.
\entry Mill (John Stuart) 105. 140. 170. 176. 316. 322. 372. 383. 398. 431. 441. 533. 541.
\entry Monogamy (\textit{Monogami}) 286 ff.
\entry Montaigne 283.
\entry Montesquieu 214.
\entry More (Thomas) 216. 387. 414. 559.
\entry Displacement of motive (\textit{Motivforskydning}) 10 f. 269 ff. 379. 392.
\entry Müller (P. E.) 138.
\entry Møller (Poul) 253.
\regletter{N.}
\entry Grace (\textit{Naade}) 501.
\entry National feeling (\textit{Nationalfølelse}) 514---518.
\entry Political economy (\textit{Nationaløkonomi}) 52. 374. 416.
\entry Newman (J. H.) 529.
\entry Nietzsche 163. 186. 503.
\entry White lie (\textit{Nødløgn}) 257 f.
\regletter{O.}
\entry Public opinion (\textit{Offentlig Mening}) 70 f. 529 ff. 576.
\entry Evil (the ethical) (\textit{det Onde}) 125 ff. (cf. 35. 42). 497.
\entry Deliberation (\textit{Overlæg}) 570 f.
\entry Owen 388.
\regletter{P.}
\entry Panaetius 90. 202.
\entry Parties (\textit{Partier}) 585 ff.
\entry Pascal 154. 254.
\entry Peel (Robert) 589 f.
\entry Pelagius 113.
\entry The principle of personality (\textit{Personlighedsprincipet}) 18. 173 ff. 392 f. 560.
\entry Personal truth (\textit{Personlig Sandhed}) 253.
\entry Reverence (\textit{Pietet}) 259.
\entry Plato 15. 33. 137. 140. 199. 201. 209. 246. 272. 387. 557.
\entry Duty (\textit{Pligt}) 39. 104. 155. 159.
\entry Poincaré 68.
\entry Polygamy (\textit{Polygami}) 282 ff. 297 f.
\entry Positive morality (\textit{Positiv Moralitet}) 6.
\entry Positive religion (\textit{Positiv Religion}) 467 ff.
\entry Preyer 332. 335.
\entry Principle (\textit{Princip}) 22.
\entry Private property (\textit{Privatejendom}) 234 f. 418 f.
\entry Probabilism (\textit{Probabilisme}) 82 f.
\entry Producers' associations (\textit{Produktionsforeninger}) 383.
\entry Promiscuity (\textit{Promiskuitet}) 283 f.
\entry Pedagogy and ethics (\textit{Pædagogik og Etik}) 329 ff.
\regletter{R.}
\entry Real self (\textit{Realt Selv}) 32 f. 39 f. 80.
\entry Referendum 582.
\entry Relativity (\textit{Relativitet}) (psychological r. 48 f. Historical r. 264 f. Individual r. 88 ff. 178 f.).
\entry Religious freedom (\textit{Religionsfrihed}) 479 ff.
\entry Religious culture (\textit{Religiøs Kultur}) 457 ff.
\entry Right, law (\textit{Ret}) (ethical right 156 f. Juridical right 522 ff.).
\entry Justice (\textit{Retfærdighed}) 40. 187 f. 206. 369 f. 374.
\entry Doctrine of law (\textit{Retslære}) 52. 522 f.
\entry Revolution 524 f.
\entry Rohde (E.) 450.
\entry Ritschl (Alb.) 95.
\entry Rousseau 186 f. 195 f. 217. 248. 254. 331. 364. 431. 435. 436. 542.
\entry Robber-farming (\textit{Rovdrift}) 421.
\regletter{S.}
\entry Love of truth (\textit{Sandhedskærlighed}) 249 ff.
\entry Saint-Simon 388.
% [Register p. 605 begins]
\opage{605}\entry Society, community (\textit{Samfund}) 168 f. 272---277.
\entry Conscience (\textit{Samvittighed}) 32. 39 f. 79 ff.
\entry Sanction (\textit{Sanktion}) 96 f
\entry Scharling (H.) 112
\entry Schiller 50. 142. 455.
\entry Schleiermacher 23. 88. 91. 158. 192. 300.
\entry Schmidt (L.) 12. 85.
\entry Schäffle 404.
\entry Schopenhauer 243 f. 465. 538.
\entry Schrempf 255. 491.
\entry Schwerin (Jeannette) 324.
\entry Self-control (\textit{Selvbeherskelse}) 220 f.
\entry Self-assertion (\textit{Selvhævdelse}) 207 ff.
\entry Suicide (\textit{Selvmord}) 210 ff.
\entry Self-preservation (\textit{Selvopholdelse}) 207 ff.
\entry Self-government (\textit{Selvstyrelse}) 590 f.
\entry Independence (\textit{Selvstændighed}) 231 f.
\entry Seneca 204. 239.
\entry Sentimentality (\textit{Sentimentalitet}) 187.
\entry Shakespeare 119. 133. 457.
\entry Sharp 89.
\entry Sibbern (F. C.) 37. 81. 103. 127. 130. 159. 230. 250. 268. 387. 529.
\entry Sibbern (G.) 250.
\entry ``Sitte'' 521.
\entry Divorce (\textit{Skilsmisse}) 308---310.
\entry Smith (Adam) 37. 64. 243. 353.
\entry Social ethics (\textit{Social Etik}) 169 ff. 263 ff.
\entry Socialism (utopian, philanthropic, speculative, empirical) (\textit{Socialisme (utopisk, filantropisk, spekulativ, empirisk)}) 386 ff.
\entry The social question (\textit{Socialt Spørgsmaal}) 364 ff.
\entry Sociology (\textit{Sociologi}) 263.
\entry Socrates 18. 50. 85. 90. 111. 127. 132. 199. 245.
\entry Sombart (W.) 394.
\entry Spencer (H.) 103 f. 167. 272. 312. 340 f. 509. 514.
\entry Spinoza 97. 127. 137. 186. 197. 232. 268. 464.
\entry Starcke 69 f. 282. 302.
\entry Stahl 535.
\entry State (\textit{Stat}) 276 f. 513 f.
\entry Steinthal 107.
\entry Stephen (Leslie) 171. 255.
\entry The Stoics 85. 90. 161. 214. 239. 293.
\entry Punishment (\textit{Straf}) 548 ff.
\entry Strauss (David) 209.
\entry Subjective thinkers (\textit{Subjektive Tænkere}) 254.
\entry Suneson (Andreas) 129.
\entry Symbols (\textit{Symboler}) 463. 468. 474.
\entry Sympathy (\textit{Sympati}) 37. 57. 236 ff.
\entry Syndicalism (\textit{Syndikalisme}) 411. 520.
\entry Synonymy (\textit{Synonymik}) 138.
\regletter{T.}
\entry Taine 320. 537 f.
\entry Theological ethics (\textit{Teologisk Etik}) 13 ff. 56.
\entry Thomas Aquinas 130. 204.
\entry Permissible (the) (\textit{Tilladelige (det)}) 158.
\entry Imputation (\textit{Tilregnelse}) 117 ff.
\entry Tocqueville 399. 521.
\entry Tolerance (\textit{Tolerance}) 258 f.
\entry Treschow 37.
\entry Defiance (\textit{Trods}) 128.
\entry Turgot 353 f.
\entry Tønnies 364. 373. 503. 560. 563.
\regletter{U.}
\entry Hypothesis of evolution (\textit{Udviklingshypotese}) 101.
\entry Utilitarianism (\textit{Utilitarianisme}) 46. 47. 49. 56. 150. 441.
\entry Utopia 387. 414.
\entry Illegitimate children (\textit{Uægte Børn}) 339.
\regletter{V.}
\entry Freedom of choice (\textit{Valgfrihed}) 110.
\entry Welfare (principle of welfare) (\textit{Velfærd (Velfærdsprincip)}) 43---46. 54 f. 94. 135 ff.
\entry Webb (Beatrice) 381.
\entry Wilckens 112.
\entry Wille (Bruno) 85. 152 164.
\entry Voltaire 140.
\entry Wundt (W.) 37. 168. 269.
\entry Valuation (\textit{Vurdering}) 3. 9 f. 25 f. 32 f. 61 f.
% [Register p. 606 begins]
\opage{606}\entry Motive of valuation (\textit{Vurderingsmotiv}) 3. 27. 62.
\entry Value (\textit{Værdi}) 3. 23. 32 f. 169.
\entry Displacement of value (\textit{Værdiforskydning}) 269.
\regletter{Y.}
\entry Humility (\textit{Ydmyghed}) 231.
\regletter{Æ.}
\entry Marriage (\textit{Ægteskab}) 281 ff.
\entry Honor (\textit{Ære}) 234.
\entry Honesty (\textit{Ærlighed}) 252.
\entry Aesthetic culture (\textit{Æstetisk Kultur}) 449 ff.
\regletter{Ø.}
\entry Standpoint of the moment (\textit{Øjebliksstandpunkt}) 30.
\entry Ørsted (A. S.) 92. 170 f. 522. 558. 560. 566. 569.
\entry Ørsted (H. C.) 162.
\end{multicols}

\end{document}
